% From True-Concurrent Games to a Categorical Logic of Agency
% Self-contained research manuscript; main constructions and detailed appendices.
\documentclass[11pt,a4paper]{article}
\usepackage[T1]{fontenc}
\usepackage[utf8]{inputenc}
\usepackage{lmodern}
\usepackage[margin=29mm]{geometry}
\usepackage{amsmath,amssymb,amsthm,mathtools}
\usepackage{microtype}
\usepackage{enumitem,booktabs,array,longtable}
\usepackage{tikz,tikz-cd}
\usetikzlibrary{arrows.meta,positioning,calc,decorations.pathreplacing}
\usepackage[authoryear,round]{natbib}
\usepackage[colorlinks=true,linkcolor=blue!45!black,citecolor=blue!45!black,urlcolor=blue!45!black]{hyperref}
\usepackage{aliascnt}
\usepackage[nameinlink,noabbrev,capitalise]{cleveref}
\hypersetup{pdftitle={From True-Concurrent Games to a Categorical Logic of Agency},pdfauthor={Conor Sheridan}}
\setlist{itemsep=2pt,topsep=5pt}
\setlength{\emergencystretch}{2em}
\numberwithin{equation}{section}
\newtheorem{theorem}{Theorem}[section]
\newtheorem{proposition}[theorem]{Proposition}
\newtheorem{lemma}[theorem]{Lemma}
\newtheorem{corollary}[theorem]{Corollary}
\theoremstyle{definition}
\newtheorem{definition}[theorem]{Definition}
\newtheorem{example}[theorem]{Example}
\newaliascnt{construction}{theorem}
\newtheorem{construction}[construction]{Construction}
\aliascntresetthe{construction}
\newtheorem{convention}[theorem]{Convention}
\newtheorem{notation}[theorem]{Notation}
\theoremstyle{remark}
\newtheorem{remark}[theorem]{Remark}
\theoremstyle{plain}
\newtheorem{maintheorem}{Theorem}
\renewcommand{\themaintheorem}{\Alph{maintheorem}}
\crefname{maintheorem}{Theorem}{Theorems}
\Crefname{maintheorem}{Theorem}{Theorems}
\crefname{construction}{Construction}{Constructions}
\Crefname{construction}{Construction}{Constructions}
\setcounter{tocdepth}{2}
\makeatletter
\renewcommand*{\l@subsection}{\@dottedtocline{2}{1.5em}{3.1em}}
\makeatother
\newcolumntype{L}[1]{>{\raggedright\arraybackslash}p{#1}}

\providecommand{\Ag}{\mathsf{Ag}}
\providecommand{\sig}{\vec{\Sigma}}
\providecommand{\taus}{\vec{T}}
\providecommand{\ths}{\vec{\Theta}}
\providecommand{\Dec}{\mathsf{Dec}}
\providecommand{\AnnPol}{\mathsf{AnnPol}}
\providecommand{\Pol}{\mathsf{Pol}}
\providecommand{\BR}{\mathsf{BR}}
\providecommand{\NE}{\mathsf{NE}}
\providecommand{\DK}{\mathsf{DK}}
\providecommand{\CK}{\mathsf{CK}}
\providecommand{\ctx}{\mathsf{Ctx}}
\providecommand{\Inst}{\mathsf{Inst}}
\providecommand{\Lin}{\mathsf{Lin}}
\providecommand{\Val}{\mathsf{Val}}
\providecommand{\I}{\mathsf{I}}
\providecommand{\Top}{\top}
\providecommand{\botell}{\mathbf{0}}
\providecommand{\ordturn}{\vdash}
\providecommand{\linturn}{\vdash}
\providecommand{\ltype}{\mathsf{LType}}
\providecommand{\type}{\mathsf{Type}}
\providecommand{\cC}{\mathcal{C}}
\providecommand{\cV}{\mathcal{V}}
\providecommand{\cT}{\mathcal{T}}
\providecommand{\cL}{\mathcal{L}}
\providecommand{\cP}{\mathcal{P}}
\providecommand{\Set}{\mathbf{Set}}
\providecommand{\Sub}{\mathop{\mathrm{Sub}}}
\providecommand{\llbracket}{\mathopen{[\![}}
\providecommand{\rrbracket}{\mathclose{]\!]}}
\providecommand{\sem}[1]{\llbracket #1\rrbracket}
\providecommand{\id}{\mathop{\mathrm{id}}}
\providecommand{\im}{\mathop{\mathrm{im}}}
\providecommand{\kerpair}{\mathop{\mathrm{ker}}}
\providecommand{\ev}{\mathop{\mathrm{ev}}}
\providecommand{\pow}{\mathcal{P}}
\providecommand{\seq}{\mathbin{;}}
\providecommand{\parcomp}{\mathbin{\Vert}}
\providecommand{\End}{\mathsf{End}}
\providecommand{\Send}{\mathsf{Send}}
\providecommand{\Recv}{\mathsf{Recv}}
\title{From True-Concurrent Games to a Categorical Logic of Agency}
\author{
  Conor Sheridan\\[0.5ex]
  \href{mailto:siridconor@gmail.com}
       {\texttt{siridconor@gmail.com}}
}
\date{\today}
\begin{document}
\maketitle
\begin{abstract}
Strategic games assign complete contingent policies to agents but normally
abstract from the causal structure of their execution. True-concurrent games
describe causal interaction without determining which economic agent controls
a choice. We connect these descriptions compositionally. Starting from the
compact-closed category of receptive courteous event-structure strategies,
we construct an indexed dependent linear--nonlinear completion and a
displayed, owner-indexed agent type. Its states are naturally equivalent to
complete-policy families of concurrent strategies, compatibly with
composition and ownerwise reparametrization.

Owner projections yield coalition-deviation kernel pairs. Predicate feedback,
observation descent, and legal comprehension then define equilibrium,
informed choice, and knowledge without identifying economic control with
operational polarity. We construct witnessed institutions, finite causal
decision protocols, and binary sessions that retain off-path policies and
secured execution witnesses. Evaluation of complete truthful communication
policies into current public messages is a split surjection with exact
strategic and epistemic successor lifting; this transports the designated
modal fixed points. Residual controllers and visible refinements extend the
construction to continuing agency. We prove reconstruction and soundness,
and establish first-layer equational completeness for the class of indexed
models. In the fixed semantics, original concurrent arenas separate open
agent morphisms, and finite retained frames admit characteristic formulas
for every bisimulation-invariant predicate.
\end{abstract}
\tableofcontents
\clearpage

\section*{Introduction}\label{sec:paper-introduction}
\addcontentsline{toc}{section}{Introduction}

Strategic and operational descriptions answer different questions about
agency. A game can say which player controls a complete contingent policy
and which other policies count as deviations, without describing how those
policies are executed. A concurrent strategy can record causal prerequisites,
conflict, and independent events, without declaring who controls the
resulting choices. Identifying these descriptions would conflate economic
control with operational polarity. We instead ask how they can be connected
compositionally while retaining both kinds of information. True concurrency
provides the operational starting point: independent moves are represented
without choosing a total interleaving.

The first task is to represent an owner-indexed family of complete policies
as an interactive object whose executions are concurrent strategies. Once
that representation is available, owner projections determine controlled
alternatives, while predicate feedback evaluates them without building a
preference into the operational strategy. Observation quotients determine
which responses are informed; legal comprehensions and execution witnesses
determine which actions are permitted or realized. Communication must then
be implemented, not merely postulated as an information change, and its
complete policies must retain responses at histories not actually reached.
Residual strategies extend this account beyond a single interaction. The
paper constructs these links categorically and gives the associated formal
calculus and verification results.

Four principal theorems organize these tasks. The first establishes the
representation of complete-policy agents; the second equips those agents
with controlled, evaluated, and institutional alternatives; the third gives
operational communication and continuation; and the fourth establishes
reconstruction and the stated logical adequacy. They concern pure policies
and qualitative information. Preferences, permissions, and signals not
determined by a protocol remain separately typed inputs.

\begin{maintheorem}[Concurrent agents and compositional representation]
\label{main:A}
Receptive courteous strategies on polarized event structures form a compact
closed category $\mathcal C$ under secured interaction and hiding.
Its restricted Yoneda completion $\mathcal V$ is symmetric monoidal closed
and supports the indexed dependent linear--nonlinear structure.  For every
finite owner list $N$, ordinary complete-policy family
$\vec\Sigma=(\Sigma_i)_{i\in N}$ and original arenas $A,B$,
there is a natural bijection
\[
\mathcal V\bigl(I,[\mathsf{Lin}(\prod_i\Sigma_i)\otimes yA,yB]\bigr)
\cong\mathbf{Set}(\prod_i\Sigma_i,\mathcal C(A,B)).
\]
Under the displayed ownership signature, this representation preserves
sequential and parallel composition, ownerwise reparametrizations, and
the realization-preserving $2$-cells of the owned-open-arena bicategory.
\end{maintheorem}

The representation in Theorem~A rests on the secured operational
category, Day completion, and dependent model of
\cref{thm:egt-base,thm:egt-day-preservation,thm:egt-dependent-model}.
Its decisive step is the dependent agent abstraction and the recovery of
owned presentations in
\cref{agt-thm-agent-dependent-product,agt-thm-representation,agt-thm-owned-equivalence}.

\begin{maintheorem}[Controlled, situated, and institutional agency]
\label{main:B}
Owner-coordinate projections determine compositional coalition
counterfactuals.  Owner regrouping and associative delegation preserve the
realized concurrent interaction and admit exact deviation comparisons on
the implemented policy image.  Predicate feedback along realization
constructs canonical selection, best response, and equilibrium.
Observation quotients construct admissible policy descent and knowledge;
dependent universal images construct observation-uniform legal policies.
Witnessed record-changing institutions form a symmetric monoidal bicategory,
and their dependent sums and products determine proof-relevant dynamic
modalities compatible with sequential execution.
\end{maintheorem}

Theorem~B first turns owner projections into alternatives and then
transports outcome comparisons back to those alternatives. Regrouping,
delegation, legal choice, and institutional execution extend the same
structure; the precise construction and comparison results appear in
\cref{thm:egt-counterfactual,thm:egt-owner-regrouping,thm:egt-delegation-substitution,thm:cgt-coplay,agt-thm-uniform,agt-thm-institutions,agt-thm-dynamics}.
The crucial distinction is between a controlled change and the comparison
by which that change is assessed.

\begin{maintheorem}[Communication, epistemic change, and continuation]
\label{main:C}
Finite owned causal decision protocols and finite binary sessions admit
true-concurrent realizations preserving complete off-path policies.
Their residual operations satisfy subject reduction, and secured interaction
is adequate for the completed communication witnesses.  Actual delivery
records determine epistemic action-event quotients.  For complete truthful
message policies, evaluation into the truthful public-message model is a
split surjection with the designated strategic and epistemic lifting
properties.  At the completed delivery boundary, collective disclosure
transports distributed knowledge to the stated common-knowledge and
announcement-ability predicates.  Complete local policies generate their
residual controllers, with the constructed criterion for finite-state
implementation; witnessed residual transitions support temporal fixed-point
agency and exact visible refinement.
\end{maintheorem}

Theorem~C moves from these abstract interfaces to causal realizations.
The protocol and session arguments in
\cref{thm:egt-causal-protocol,agt-thm-session,agt-lem-session-input-residual,agt-thm-message-comparison}
retain off-path responses and operational witnesses. The pooling and
disclosure comparisons of \cref{agt-thm-pooling,agt-thm-disclosure}
then identify the information created by delivery. Section~\ref{sec:paper-12}
uses residuals and visible refinements to pass from completed interactions
to continuing ones.

\begin{maintheorem}[Soundness, recovery, and finite behavioural completeness]
\label{main:D}
The ordinary, linear, displayed-agent, and contextual judgments of the
agency calculus are sound in the completed true-concurrent semantics.
The owner-indexed and enriched situated presentations are recovered by the
explicit reconstruction maps, while reductions relative to designated
tests retain their exhibited alternative-lifting properties.
The first-layer equational calculus is complete for its class of split
indexed linear--nonlinear models by the syntactic classifying model.
In the fixed concurrent completion, original arenas separate open agent
morphisms. On image-finite retained frames, the exact bisimulation
quotient identifies precisely the states satisfying the same designated
Boolean modal tests; on finite frames, characteristic formulas define
all bisimulation-invariant predicates. Explicit finite causal-agent
interpretations have decidable semantic equality. 
\end{maintheorem}

Theorem~D closes the comparison between presentations and semantics.
Section~\ref{sec:paper-14} reconstructs the displayed data, establishes
exact reductions, and proves operational separation and finite-state
modal expressivity in the fixed semantics. Appendix~\ref{app:metatheory}
proves reconstruction, syntactic soundness, and classifying-model
equational completeness. The fixed-model tests and the syntactic
completeness theorem have distinct conclusions.

\paragraph{Construction-to-theorem guide.}
The table separates the imported categorical foundations from the
constructions established here and points to the comparisons on which the
four theorems depend.
\begin{center}
\small
\begin{tabular}{@{}L{.12\textwidth}L{.38\textwidth}L{.41\textwidth}@{}}
\toprule
Claim & Construction and comparison & Proof location \\
\midrule
A & Imported secured interaction and Day convolution; constructed
$\Lin\dashv\Val$ and the open-context agent representation. &
\Cref{thm:egt-base,thm:egt-day-preservation,thm:egt-dependent-model,agt-thm-representation,agt-thm-owned-equivalence}. \\
B & Ownerwise kernel-pair deviations, legal alternative lifting,
comparison feedback, and matched institutional witnesses. &
\Cref{lem:ownerwise-legal-lifting,thm:egt-delegation-image,thm:cgt-selection-equilibrium,thm:egt-institution-bicategory}. \\
C & Finite causal wiring, session-prefix residuals, truthful-message
back-lifts, and the greatest-bisimulation quotient. &
\Cref{thm:egt-causal-protocol,agt-thm-session,thm:egt-announcement-policy-comparison,lem:egt-announcement-successor-fibres,thm:egt-exact-visible-refinement}. \\
D & Reconstruction and the first-layer classifying model;
original-arena separation and finite-state modal characterization. &
\Cref{agt-thm-reconstruction,agt-thm-completeness,thm:agt-concurrent-probe-separation,thm:egt-hennessy-milner-agency,thm:egt-finite-modal-completeness}. \\
\bottomrule
\end{tabular}
\end{center}
The mathematical proofs of the constructed operational, strategic,
institutional, and epistemic results appear alongside their statements
in Parts~I--V. Appendix~\ref{app:rules} presents the natural-deduction
calculus and assertion language; Appendix~\ref{app:metatheory} contains
the remaining reconstruction, soundness, and equational-completeness proofs.
The compact-closed strategy category and Day convolution are imported
results; their connection to displayed ownership, dependent policy
interfaces, legal and epistemic alternatives, and causal implementations
is constructed here.

\paragraph{Three semantic levels.}
Ordinary policies and their dependent families live in $\mathbf{Set}$;
linear interaction interfaces live in presheaves over the concurrent
strategy category; and assertions classify properties of states and
witnesses. Policy parameters specify which interactions can be chosen, not
which resources can be copied. Comparison maps supply economic evaluation,
while proof-relevant execution data remain available even when existential
images record only feasibility. The adjunctions and representation maps
below give the precise passage between these levels.

\paragraph{What the operational semantics determines.}
Causality, enabled moves, residuals, and the signals of displayed wires
come from execution. Preferences, permission rules, and observations not
realized by a protocol do not. They enter as explicit maps or subobjects,
from which the paper derives knowledge, strategic selection, coalition
ability, and updates. In particular, economic rationality is not read off
an event's polarity.

\paragraph{A running finite instance.}
The recurring example has two owners and a private bit for each: its world
carrier is $W=2\times2$, with observations $o_1(x,y)=x$ and
$o_2(x,y)=y$. At separate decision nodes owner $1$ applies
$\pi_1:2\to2$ to $x$ and owner $2$ applies $\pi_2:2\to2$ to $y$.
Thus $\Sigma_1=2^2$, $\Sigma_2=2^2$, and the complete profile carrier
has $16$ elements, including all off-path responses. The operational
realization is the tensor of the two response strategies, with their
specified private inputs. An institutional rule can couple the policies:
we use the concrete legal domain
\[
 L=\{(\pi_1,\pi_2):\pi_1(1)\leq\pi_2(1)\}\hookrightarrow
 \Sigma_1\times\Sigma_2,
\]
which has $12$ elements. The order is the usual order on $2=\{0,1\}$.
The current-world payoff comparison below uses
$u_1(a,b)=2ab-a$ and $u_2(a,b)=2ab-b$, where
$a=\pi_1(x)$ and $b=\pi_2(y)$. Later, a collector publicly delivers
$(x,y)$, and an institutional authorization event may guard a final
execution. These are different interfaces: the legal domain and utilities
are supplied, while the response, collection, and guarded execution have
the exhibited causal realizations. The concrete deviations, epistemic
updates, and authorization witnesses are calculated in
\cref{sec:agt-running-institution,sec:agt-running-example}.

\paragraph{Route through the paper.}
Part~I constructs the concurrent and dependent semantics and represents
agents within it. Part~II shows how complete policies compose, which
alternatives owners control, and how comparison feedback selects them.
Part~III treats information, legal choice, institutions, and epistemic
update. Part~IV constructs the finite communication processes needed for
those updates and compares their complete policies with public messages.
Part~V develops residual agency, temporal and epistemic rationality,
reconstruction, and abstraction. Appendix~\ref{app:rules} specifies the natural-deduction
rules, and Appendix~\ref{app:metatheory} proves reconstruction, soundness,
and equational completeness.

\subsection*{Relation to the literature}

Two existing lines of work clarify why agency requires more than an
action modality. The type-theoretic action calculi of Barber, Gardner,
Hasegawa, and Plotkin relate interaction to linear logic
\citep{agt-BarberEtAl}; resource-sensitive agency logics likewise connect
modal action with resources \citep{agt-PorelloTroquard}. Our question is more
specific: an agent must itself be a typeable, composable resource, while
its owner-indexed control structure remains visible even when its
underlying linear object is passed to another process.

Parameterized open games already distinguish controlled policy
coordinates from evaluation \citep{CapucciEtAl2022Agency}. We retain that
distinction over a true-concurrent operational category, representing the
policy family internally while keeping utility in a comparison interface.
Session calculi relate channel interaction to cut
\citep{agt-CairesEtAl}; our finite sessions give a causal realization with
explicit complete policies and secured interaction witnesses. The
operational implementation is needed to identify information actually
delivered, rather than information merely described by an abstract update.

The contextual construction also draws on epistemic actions treated as
resources \citep{agt-BaltagCoeckeSadrzadeh} and multi-type presentations of
dynamic epistemic logic \citep{agt-FrittellaEtAl}. Here observation
adjunctions produce knowledge, comprehension produces public restriction,
and witnessed transitions produce general update. Keeping these operations
at distinct interfaces makes it possible to compare abstract epistemic
actions with their causal executions.

For the logical characterization of finite branching systems, the
Hennessy--Milner separation theorem \citep{HennessyMilner1980,HennessyMilner1985} supplies
the modal precedent. The contribution below is its application to the
retained observation, deviation, and witnessed transition relations,
together with characteristic tests for the exact refinement already
constructed in the agency semantics.

The representation uses standard copower, internal-hom, and Yoneda
identities. The comparison established here connects their displayed
owner-indexed interpretation to receptive courteous execution and preserves
strategic alternatives and epistemic successors. Its principal constructions
are the source-fixed legal lift of \cref{lem:ownerwise-legal-lifting}, the
secured implementation of \cref{thm:egt-causal-protocol}, and the simultaneous
epistemic and strategic back-lifts of
\cref{thm:egt-announcement-policy-comparison}. Compact closure and Day
completion supply the operational semantics; the displayed ownership,
comparison, legal, and observation interfaces determine the strategic and
epistemic properties preserved by these constructions.

\noindent\textbf{Use of AI-assisted tools.} Generative AI systems were used during the research and development of this manuscript. Responsibility for the mathematical content, arguments, and conclusions rests with the author.

\part{Concurrent semantics and first-class agents}

\section{True-Concurrent Games and Dependent Linear Semantics}\label{sec:paper-1}

A concurrent strategy must preserve the causal structure of an
interaction without imposing an order on independent events. The
polarized event structures and secured composition recalled from
\citet{CastellanEtAl2017} provide that operational category. They do not,
however, provide the dependent sums, products, and ordinary policy
parameters required for agent types. We first construct finite value and
collection protocols, exhibiting which data operations impose causal
barriers. We then use restricted Yoneda and Day convolution to obtain the
dependent linear--nonlinear interfaces while recovering every original
concurrent arrow. This supplies the operational and categorical inputs to
the agent representation in the next section.

\subsection{Operational event structures and concurrent interaction}
\label{sec:egt-interaction}

\subsubsection{Concurrent histories}

\begin{definition}[Event structures and configurations]
\label{def:egt-event-structures-and-configurations}
An \textbf{event} is one occurrence of an action or communication, so two
occurrences of the same kind of action remain distinct events.
\textbf{True concurrency} records prerequisites and compatibility directly,
without first choosing a total schedule for independent events. In the
formal data below, a down-closed set contains every prerequisite of each
of its events. An event is \textbf{enabled} after a configuration $x$ when
$e\notin x$ and $x\cup\{e\}$ is again a configuration. Two distinct events are concurrent
when they can occur together and neither causally precedes the other;
incompatible collections are in conflict. The consistency family also
records incompatibilities involving more than two events.

An \textbf{event structure} is a triple $E=(|E|,\leq_E,\operatorname{Con}_E)$.
The relation $\leq_E$ is a partial order with finite principal downsets.
The family $\operatorname{Con}_E$ consists of finite subsets of $|E|$, contains the
empty set and every singleton, is closed under subsets, and satisfies
\[
 X\in\operatorname{Con}_E,\quad e\leq_E e'\in X
 \quad\Longrightarrow\quad X\cup\{e\}\in\operatorname{Con}_E.
\]
A finite \textbf{configuration} is a down-closed consistent subset.
Write $\operatorname{Conf}(E)$ for the set of finite configurations.
\end{definition}

The order records prerequisites. Consistency records which finite collections
can occur together. Two events can be consistent without either causing the
other. A configuration therefore records a partially ordered history rather
than an arbitrary linear ordering of its independent events.

\begin{construction}[Enabled-event span]
\label{con:egt-enabled-event-span}
For an event structure $E$ put $W_E=\operatorname{Conf}(E)$ and define
\[
 \mathsf{Step}(E)
 =\{(x,e)\in W_E\times|E|:e\notin x,\ x\cup\{e\}\in W_E\}.
\]
It has projections $p(x,e)=x$ and $\ell(x,e)=e$, and a target
$t(x,e)=x\cup\{e\}$. Thus
\[
 W_E\xleftarrow{p}\mathsf{Step}(E)\xrightarrow{t}W_E
\]
is a relation span: $(p,t)$ is jointly injective, since
$e=t(x,e)\setminus x$ is determined by its two endpoints. The inclusion
$\mathsf{Step}(E)\hookrightarrow W_E\times|E|$ is the comprehension
of the event-enablement predicate. Every event of a finite configuration
can be executed in some linear extension of that configuration's causal
order, but the span does not choose a linear extension.
\end{construction}

\noindent The incidence inclusion and source projection give the enabled-event diagram.
\begin{center}
\begin{tikzcd}[column sep=large,row sep=large]
\mathsf{Step}(E) \arrow[r,hook,"\iota"] \arrow[dr,"p"']
 & W_E\times |E| \arrow[d,"\pi_1"] \\
 & W_E
\end{tikzcd}
\end{center}

\begin{figure}[htbp]
\centering
\begin{tikzpicture}[>=Stealth, every node/.style={font=\small},
  event/.style={draw,rounded corners,inner sep=6pt}]
\node[event] (a) at (0,0) {commit $A$};
\node[event] (b) at (3.6,0) {commit $B$};
\node[event] (c) at (1.8,-1.3) {settle};
\draw[->] (a)--(c); \draw[->] (b)--(c);
\node[align=center,text width=9cm] at (1.8,-2.35)
  {Neither commitment orders the other. Settlement depends on both.};
\end{tikzpicture}
\caption{Causal dependence without a scheduler-imposed order.}
\label{fig:egt-concurrency}
\end{figure}

\begin{definition}[Total event-structure map]
\label{def:egt-total-event-structure-map}
For event structures $E,F$, a \textbf{total map of event structures} is a
function $f:|E|\to|F|$ satisfying, for every finite configuration
$x\in\operatorname{Conf}(E)$,
\[
 f[x]\in\operatorname{Conf}(F),\qquad
 f|_x:x\longrightarrow f[x]\text{ is injective}.
\]
It thus preserves configurations and is locally injective. It need not
preserve every source causal order as an order relation between event
labels: a strategy may add dependencies specifying when it responds.
\end{definition}
\begin{definition}[Polarized arenas]
\label{def:egt-polarized-arenas}
The polarity is relative to an interaction boundary: negative events are
moves the surrounding interaction may supply, and positive events are moves
the strategy may make. Reversing an interface reverses these roles. Economic ownership is recorded by the separate owner-indexed policy signature. For example,
several players' chosen messages can all be positive outputs of separate
components, while becoming negative inputs of the component that receives
them.

An \textbf{arena} is an event structure with a polarity
$\operatorname{pol}:|A|\to\{-,+\}$. Its dual $A^\perp$ reverses polarities.
The parallel arena $A\parallel B$ is the disjoint union, with componentwise
causality and consistency. The empty arena is denoted $I$.
\end{definition}

\subsubsection{Strategies and their composition}

To compose strategies, occurrences at the common interface must be
matched without creating a cycle of causal prerequisites. The
synchronization pullback retains the secured histories supporting each
match, including distinct incompatible histories that have the same
visible event labels.

We use the secured-bijection pullback of \citet{CastellanEtAl2017}.
Its universal property supplies the synchronization object subsequently
used in interaction and hiding.

\begin{definition}[Secured matching]
\label{def:egt-secured-matching}
Let $f:S\to A$ and $g:T\to A$ be total maps of unpolarized event
structures. A matching is a finite graph $\varphi\subseteq |S|\times|T|$
whose projections are configurations, which is the graph of a bijection
between those configurations, and for which $f(s)=g(t)$ whenever
$(s,t)\in\varphi$. Order its pairs by the transitive closure of
\[
 (s,t)\mathrel{\trianglelefteq}(s',t')
 \quad\text{if}\quad s\leq_S s'\text{ or }t\leq_T t'.
\]
The matching is secured when this closure is antisymmetric. It is prime
when this finite partial order has a greatest pair. Write
$\mathcal M(f,g)$ and $\mathcal M_\bullet(f,g)$ for the secured and
prime secured matchings.
\end{definition}

\begin{construction}[Synchronization object]
\label{con:egt-pullback}
The events of $\mathsf{Sync}(f,g)$ are $\mathcal M_\bullet(f,g)$, ordered by
inclusion. A finite set of primes is consistent precisely when its
union is a secured matching. A prime is projected to the coordinates
of its greatest pair, giving maps $p:\mathsf{Sync}(f,g)\to S$ and $q:\mathsf{Sync}(f,g)\to T$.
\end{construction}

\noindent Synchronization is determined by its projections into the shared event-structure interface.
\begin{center}
\begin{tikzcd}[column sep=large,row sep=large]
& \mathsf{Sync}(f,g) \arrow[dl,"q"']\arrow[dr,"p"] & \\
 T \arrow[dr,"g"'] && S \arrow[dl,"f"] \\
 & A &
\end{tikzcd}
\end{center}

\begin{theorem}[Synchronization universal property; recalled]
\label{thm:egt-pullback}
The construction is an event structure, with a bijection
\[
 \operatorname{Conf}(\mathsf{Sync}(f,g))\cong\mathcal M(f,g),
 \qquad X\longmapsto\bigcup X.
\]
The maps $p,q$ exhibit it as the pullback of $f,g$ in the category of
event structures and total maps.
\end{theorem}
\begin{proof}[Source and identification]
\leavevmode
\begin{enumerate}[label=\textup{(\arabic*)},leftmargin=*,itemsep=4pt,topsep=3pt]
  \item \textbf{Secured operational interaction.} This is the secured-bijection pullback theorem of
\citet{CastellanEtAl2017}. In the notation fixed above, a secured
matching $\varphi$ determines the configuration
\[
 X_\varphi=\{\downarrow_\varphi u:u\in\varphi\}
\]
of prime secured matchings, and a configuration $X$ of primes recovers the
secured matching $\bigcup X$. These operations are inverse.
  \item \textbf{Lifting through the comparison.} The maps $p$
and $q$ take a prime to the coordinates of its greatest pair, and the
cited pullback theorem identifies the unique lift of a cone with the prime
secured histories generated by its configurations. We use that universal
property below; no different synchronization construction is introduced
here.
\end{enumerate}
\end{proof}

\begin{figure}[htbp]
\centering
\begin{tikzcd}[column sep=large,row sep=large]
V \arrow[dr,dashed,"\ell"] \arrow[drr,bend left=16,"h"]
  \arrow[ddr,bend right=16,"k"'] & & \\
& \mathsf{Sync}(f,g) \arrow[r,"p"] \arrow[d,"q"'] & S \arrow[d,"f"] \\
& T \arrow[r,"g"'] & A
\end{tikzcd}
\caption{Prime secured histories supply the synchronization pullback.}
\end{figure}

\begin{definition}[Concurrent strategy]
\label{def:egt-concurrent-strategy}
Receptivity requires a unique extension in the strategy for the arrival of
each enabled negative event. Courtesy permits the strategy to wait for an
input before producing an output, but prevents it from inserting arbitrary
new dependencies between moves whose ordering is already controlled by the
arena. Here $e\lessdot e'$ means an immediate causal dependency, with no
event strictly between them. Thus a protocol may wait for both submitted
messages before announcing their joint outcome without imposing an order
between the two submissions.

A concurrent strategy from $A$ to $B$ is a polarity-preserving total map
\[
 s:S\longrightarrow A^\perp\parallel B
\]
with two conditions. \textbf{Receptivity} says that every enabled negative move
has a unique extension in the strategy. More explicitly, if $x$ is a
configuration of $S$ and $s(x)\cup\{a\}$ is a configuration extending $s(x)$ by
a negative event, there is a unique $e$ with $s(e)=a$ and
$x\cup\{e\}\in\operatorname{Conf}(S)$. \textbf{Courtesy} says that a new immediate
dependency may only run from a negative event to a positive event. Thus, if
$e\mathrel{\lessdot}_S e'$ and either $\operatorname{pol}(e)=+$ or
$\operatorname{pol}(e')=-$, then $s(e)\mathrel{\lessdot}s(e')$ in the arena.
\end{definition}

\begin{construction}[Secured interaction and hiding]
For strategies $S\to A^\perp\parallel B$ and
$T\to B^\perp\parallel C$, interaction synchronizes the occurrences on $B$.
Only \textbf{secured} synchronizations are admitted: the combined causal orders
of the matched histories must be acyclic. Composition then hides the matched
$B$-events while retaining the induced order on visible events. The
acyclicity condition prevents a circular collection of unfulfilled
prerequisites from being mistaken for a possible interaction.
\end{construction}

\begin{figure}[htbp]
\centering
\begin{tikzcd}[column sep=large,row sep=large]
& S\mathbin{\circledast_B}T \arrow[dl,dashed] \arrow[dr,dashed]
  \arrow[d,dashed,"\text{hide }B"] & \\
S \arrow[d] & T\odot S \arrow[d] & T \arrow[d] \\
A^\perp\parallel B & A^\perp\parallel C & B^\perp\parallel C
\end{tikzcd}
\caption{Interaction precedes hiding. Dashed arrows are partial maps of
event structures: projection forgets events private to the other component,
and hiding forgets internal events. Solid arrows are total strategy maps.}
\end{figure}

\begin{theorem}[Concurrent composition; imported]
\label{thm:egt-base}
Polarized event structures and receptive courteous strategies, taken up to
isomorphism over their arenas, form a compact closed category $\mathcal C$.
Composition is secured interaction followed by hiding, identity is copycat,
tensor is parallel composition, and duality reverses polarity.
\end{theorem}

\begin{proof}[Reference and size verification]
\leavevmode
\begin{enumerate}[label=\textup{(\arabic*)},leftmargin=*,itemsep=4pt,topsep=3pt]
  \item \textbf{Composition.} For composition and compact closure, see
\citet{CastellanEtAl2017}. Parallel composition, finite secured histories,
and hiding preserve countability. On a fixed countable carrier, orders,
consistency families, polarities, and labellings range over sets.
  \item \textbf{Causal execution.} Taking a
set of countable carriers and quotienting strategy presentations by
isomorphism therefore gives a set $\mathcal C(A,B)$ for every $A,B$.
\end{enumerate}
\end{proof}

\begin{notation}[Compact structure]
We write $\otimes$ for parallel composition in $\mathcal C$ and $I$ for the
empty arena. Arena-preserving strategy isomorphisms represent equality of
arrows. Put
\[
 \widetilde\eta_A=
 I\xrightarrow{\eta_A}A^\perp\otimes A
 \xrightarrow{\gamma}A\otimes A^\perp.
\]
Interaction polarities are independent of economic ownership.
\end{notation}

\subsubsection{The linear account of interaction}

\begin{construction}[Composition and copycat in the chosen convention]
The preceding pullback describes the interaction used by
Theorem~\ref{thm:egt-base}. Forget the polarities on the shared interface
while forming the pullback of
\[
 S\parallel C\xrightarrow{s\parallel\operatorname{id}}
 A^\perp\parallel B\parallel C
 \xleftarrow{\operatorname{id}\parallel t} A^\perp\parallel T.
\]
The additional factors carry the events private to the other component.
The pullback has a map to the three-interface arena. Keep just its events
labelled in $A^\perp\parallel C$, restrict order and consistency, and
restrict the labelling map. This gives the composite. Restricting the
transitive order retains the visible consequence of an internal causal
chain; it does not simply erase the chain's dependencies.

The copycat strategy of \citet{CastellanEtAl2017} is constructed as follows: start with the two interface copies of $A$, add
the forwarding edge from the negative occurrence of each arena event to its
positive occurrence, take transitive closure, and define consistency through
the causal closure in the arena. We use the cited verification that this gives
a receptive courteous strategy and that it is the identity for secured
interaction followed by hiding. Thus no separate copycat correctness claim is
being introduced here; the concrete description only fixes the convention
used by the imported compact-closed structure.
\end{construction}

\begin{construction}[The closed structure of concurrent interaction]
The tensor is $A\otimes B=A\parallel B$, with unit the empty arena.
Compact closure constructs the internal hom
$A\multimap B=A^\perp\parallel B$ and the natural bijection
\[
 \mathcal C(X\otimes A,B)\cong\mathcal C(X,A\multimap B).
\]
Evaluation and transposition are obtained from the displayed cups, caps,
and symmetry. The snake equations prove the inverse laws for this
bijection. Identity is copycat, composition is interaction and hiding,
and tensor places interactions in parallel.
\end{construction}

Compact closure supplies the tensor and internal hom used for linear
interaction; compare \citet{Girard1987}. It does not yet make an ordinary
policy parameter reusable or allow a linear interface to depend on that
parameter. Section~\ref{sec:egt-dependent-semantics} constructs precisely
those additional operations.

\begin{remark}[Related concurrent semantics]
The extension by symmetry and the construction of concurrent Hyland--Ong
games are developed in \citet{CastellanClairambaultWinskel2019}.
\citet{Winskel2023} relates concurrent strategies to functional models,
including lenses and optics. These works situate the chosen operational
category and its functional fragments; the constructions used here are the
ones explicitly specified above and below.
\end{remark}

\subsubsection{Ordinary functions as response protocols}

\begin{construction}[Finite value arenas and response strategies]
For a finite set $X$, let $D(X)$ be the positive arena with one event for each
$x\in X$, no nontrivial causal order, and configurations $\varnothing$ and
the singletons. Distinct values conflict. A function $f:X\to Y$ determines
a strategy $\widehat f:D(X)\to D(Y)$. For each $x$, its corresponding branch
contains a negative input event $x$ and a positive output event labelled
$f(x)$, with the input preceding the output; different input branches
conflict.
\end{construction}

\begin{proposition}
\label{prop:egt-functions}
The assignments $X\mapsto D(X)$ and $f\mapsto\widehat f$ form a faithful
functor from finite sets and functions to $\mathcal C$.
\end{proposition}
\begin{proof}
\leavevmode
\begin{enumerate}[label=\textup{(\arabic*)},leftmargin=*,itemsep=4pt,topsep=3pt]
  \item \textbf{Secured operational interaction.} The only added dependency is negative-to-positive, so courtesy holds. At the empty configuration every input is accepted once; after an input,
all other inputs conflict with it in the arena.
  \item \textbf{Receptivity and functional composition.} This proves receptivity. Synchronizing $\widehat f$ and $\widehat g$ produces, on the branch indexed
by $x$, the chain $x\to f(x)\to g(f(x))$. Hiding the middle event gives
$\widehat{g\circ f}$.
  \item \textbf{Secured operational interaction.} For the identity function the construction is copycat
on $D(X)$. Distinct functions differ on an input and therefore on an output
label of the corresponding strategy; they cannot be isomorphic over the
arena.
\end{enumerate}
\end{proof}

\noindent The point strategies form a natural family with respect to finite response maps.
\begin{center}
\begin{tikzcd}[column sep=large,row sep=large]
I \arrow[r,"\mathsf{pt}_X(x)"] \arrow[d,equal]
 & D(X)\arrow[d,"\widehat f"] \\
 I\arrow[r,"\mathsf{pt}_Y(f(x))"'] & D(Y)
\end{tikzcd}
\end{center}

\begin{notation}[Value states]
For $x\in X$, write
\[
 \mathsf{pt}_X(x):I\longrightarrow D(X)
\]
for the one-event positive strategy selecting $x$. Equivalently, if
$\bar x:1\to X$ is the corresponding function and
$\upsilon:I\to D(1)$ is the unique one-event state, then
\[
 \mathsf{pt}_X(x)=\widehat{\bar x}\circ\upsilon.
\]
Functoriality gives the naturality equation
\begin{equation}
\label{eq:egt-value-state-naturality}
 \widehat f\circ\mathsf{pt}_X(x)
 =\mathsf{pt}_Y(f(x)).
\end{equation}
Equivalently, under $X\cong\mathbf{FinSet}(1,X)$, the functions
$\mathsf{pt}_X:X\to\mathcal C(I,D(X))$ are the components of the
natural transformation
\[
 \mathsf{pt}:\mathbf{FinSet}(1,-)
 \Longrightarrow \mathcal C(I,D(-)),
\]
between functors $\mathbf{FinSet}\to\mathbf{Set}$.
Its construction is the hom-functor action of $D$ on
$X\cong\mathbf{FinSet}(1,X)$, followed by precomposition with
$\upsilon$. The symbol denotes a selected point of a data arena and is
separate from both compact coevaluation $\eta_A$ and later quotation.
\end{notation}

The response functor is faithful, but it does not identify Cartesian
products of values with parallel channels. The arena $D(X\times Y)$ reports
one paired value, whereas $D(X)\otimes D(Y)$ admits independent outputs.
Collecting those outputs must therefore record their causal prerequisites.
The resulting barrier will prevent a later policy from using a joint report
before every required input has arrived.

\subsubsection{Communication of data and synchronization barriers}

\begin{construction}[Copying and disposal of data]
For data that may be reported to several recipients there are explicit
copying and disposal strategies
\[
 \Delta_X:D(X)\to D(X)\otimes D(X),\qquad
 \epsilon_X:D(X)\to I.
\]
The source of $\Delta_X$ has, for each value $x$, one negative input
and two independent positive outputs, each labelled $x$ and each
dependent on the input. Different input branches conflict. The source
of $\epsilon_X$ is just the negative input arena. These are receptive
and courteous by direct inspection: every input is accepted uniquely,
and the only new dependencies are negative-to-positive.
\end{construction}

\begin{proposition}[Duplicable data]
\label{prop:egt-data-comonoid}
The triple $(D(X),\Delta_X,\epsilon_X)$ is a commutative comonoid
in $\mathcal C$. For every finite function $f:X\to Y$,
\[
 (\widehat f\otimes\widehat f)\Delta_X
   =\Delta_Y\widehat f,\qquad
 \epsilon_Y\widehat f=\epsilon_X.
\]
\end{proposition}
\begin{proof}
\leavevmode
\begin{enumerate}[label=\textup{(\arabic*)},leftmargin=*,itemsep=4pt,topsep=3pt]
  \item \textbf{Associativity.} Both bracketings of the iterated copy strategy have, after interaction
and hiding, one negative input event labelled $x$ and three independent
positive output events labelled $x$, each with that input as its
prerequisite. Matching corresponding events gives an isomorphism over
the arena.
  \item \textbf{Copying comonoid laws.} This proves coassociativity, including its prescribed
associators. Swapping the two outputs is a branchwise isomorphism,
proving cocommutativity. Composing either output with disposal leaves
one input and its single copied output, exactly copycat. This proves
the two counit equations.
  \item \textbf{Causal execution.} For the naturality equation, each side has an input branch for $x$
and two independent outputs labelled $f(x)$. Their only visible causal
prerequisite is the input. The branchwise identification is again an
isomorphism.
  \item \textbf{Secured operational interaction.} Discarding the output of $\widehat f$ leaves just its
receptive input behavior, which is $\epsilon_X$.
\end{enumerate}
\end{proof}

\noindent Naturality of copying is expressed by the following commuting square.
\begin{center}
\begin{tikzcd}[column sep=large,row sep=large]
D(X)\arrow[r,"\Delta_X"]\arrow[d,"\widehat f"']
 & D(X)\otimes D(X)\arrow[d,"\widehat f\otimes\widehat f"] \\
 D(Y)\arrow[r,"\Delta_Y"'] & D(Y)\otimes D(Y)
\end{tikzcd}
\end{center}

Copying and disposal are designated operations on these finite data
interfaces, not structural rules of the linear syntax. An isomorphism in
the ambient concurrent category can transport a comonoid structure, so
material nonduplicability cannot be inferred from an object's isomorphism
class alone. Material feasibility and authorization are instead tracked by
exhibited certificates and witnessed institutional executions, as constructed
in Section~\ref{sec:egt-witnesses}; an ambient morphism alone does not
establish either property. Thus communication may copy a signal without
thereby certifying another authorization or resource occurrence.

\begin{construction}[Distribution of a paired value]
There is also a distributor
$d_{X,Y}:D(X\times Y)\to D(X)\otimes D(Y)$. For each pair $(x,y)$ its
source has one negative input labelled $(x,y)$ and two independent positive
outputs labelled $x$ and $y$, both following that input. Different paired
inputs conflict. It is receptive and courteous by the same direct check as the
response protocols above.
\end{construction}

A collector must accommodate partially received tuples. Completed
tuples sharing a channel value therefore share the corresponding negative
input event, rather than generating incompatible copies of that input.
The canonical collector below makes this sharing explicit.

\begin{construction}[Canonical finite collector]
\label{con:egt-finite-collector}
Let $J$ be a finite index set and let $(X_j)_{j\in J}$ be finite sets. Put
$X_J=\prod_{j\in J}X_j$. Define an event structure $C_{\mathbf X}$ with
negative events
\[
 e_{j,x}\qquad(j\in J,\ x\in X_j)
\]
and positive events
\[
 o_{\mathbf x}\qquad(\mathbf x=(x_j)_{j\in J}\in X_J).
\]
Generate its causal order by
\[
 e_{j,x_j}<o_{\mathbf x}\qquad(j\in J).
\]
A finite set of events is consistent precisely when it contains at most one
input value on each component, at most one tuple output, and every input that
occurs together with $o_{\mathbf x}$ has the value prescribed by
$\mathbf x$. Label $e_{j,x}$ by the input occurrence $x$ on component $j$
and label $o_{\mathbf x}$ by the tuple $\mathbf x$. Denote the resulting map
by
\[
 c_{\mathbf X}:C_{\mathbf X}\longrightarrow
 \left(\bigotimes_{j\in J}D(X_j)\right)^\perp
 \parallel D(X_J).
\]
For $J=\{1,2\}$ write $c_{X,Y}$, and for $J=\varnothing$ write
$c_0:I\to D(1)$.
\end{construction}

\begin{proposition}[Finite collectors]
\label{prop:egt-finite-collectors}
The map $c_{\mathbf X}$ is a receptive courteous strategy
\[
 c_{X_j\,|\,j\in J}:
 \bigotimes_{j\in J}D(X_j)
 \longrightarrow
 D\!\left(\prod_{j\in J}X_j\right).
\]
A configuration containing the tuple output $o_{\mathbf x}$ contains exactly
its required input events $e_{j,x_j}$ in their causal past. Every iterated
binary collector, followed by the canonical reassociation of the Cartesian
product, is isomorphic to this strategy over its arena. The nullary collector
$c_0$ is the one-event positive state selecting the empty tuple.
\end{proposition}

\begin{proof}
\leavevmode
\begin{enumerate}[label=\textup{(\arabic*)},leftmargin=*,itemsep=4pt,topsep=3pt]
  \item \textbf{Causal execution.} A principal downset of an input event is a singleton, while
\[
 \downarrow o_{\mathbf x}
 =\{o_{\mathbf x}\}\cup\{e_{j,x_j}:j\in J\},
\]
so principal downsets are finite. Consistency is downward closed. For the
cause-extension axiom, the only nontrivial case is
$e_{j,x_j}\leq o_{\mathbf x}$ with $o_{\mathbf x}$ already in a consistent
set.
  \item \textbf{Causal execution: second comparison.} Such a set cannot contain a different input value on component $j$;
therefore adjoining $e_{j,x_j}$ preserves consistency. Thus
$C_{\mathbf X}$ is an event structure.
  \item \textbf{Causal execution.} A configuration containing $o_{\mathbf x}$ is down-closed, hence contains all
$e_{j,x_j}$, and its image is a configuration of the arrow arena. A
configuration without an output contains at most one value on each input
component. The labelling is injective on either kind of configuration, so it
is a total map of event structures.
  \item \textbf{Secured operational interaction.} For receptivity, let $z$ be a source configuration and suppose that the arrow
arena enables a negative input $x\in X_j$. Then $z$ contains no input on
component $j$. It cannot contain a tuple output either, because down-closure of
such an output would already contain an input on component $j$.
  \item \textbf{Collector receptivity and courtesy.} Hence
$z\cup\{e_{j,x}\}$ is a configuration. The source has exactly one event with
that arena label, so the lift is unique. Every added causal edge is of the form
$e_{j,x_j}<o_{\mathbf x}$, hence is negative-to-positive; courtesy follows.
  \item \textbf{Secured operational interaction.} Now expand any iterated binary collector before hiding its intermediate tuple
events. On a completed tuple $\mathbf x$, the visible source events are the
original inputs $e_{j,x_j}$ and the final tuple output. The internal tuple
events form a directed tree from those inputs to the final output, so secured
interaction is acyclic.
  \item \textbf{Secured operational interaction: second comparison.} Hiding the internal tuple events leaves
$e_{j,x_j}<o_{\mathbf x}$ for every $j$ and adds no order between distinct
inputs. Its visible consistency condition is exactly the one in
Construction~\ref{con:egt-finite-collector}. Matching equally labelled
visible events therefore gives an isomorphism over the arena.
  \item \textbf{Causal execution.} Cartesian
coherence identifies every parenthesization with the same product tuple. The
nullary case has no negative events and one positive event labelled by the
unique element of $1$.
\end{enumerate}
\end{proof}


\begin{proposition}[Lax symmetric-monoidal data interface]
\label{prop:egt-data-lax-monoidal}
The faithful functor $D:\mathbf{FinSet}\to\mathcal C$ of
Proposition~\ref{prop:egt-functions}, together with
\[
 c_{X,Y}:D(X)\otimes D(Y)\to D(X\times Y),
 \qquad c_0:I\to D(1),
\]
is lax symmetric monoidal. Its canonical $J$-ary comparison maps are the
finite collectors of Proposition~\ref{prop:egt-finite-collectors}.
\end{proposition}
\begin{proof}
\leavevmode
\begin{enumerate}[label=\textup{(\arabic*)},leftmargin=*,itemsep=4pt,topsep=3pt]
  \item \textbf{Coherence of finite collectors.} The associativity and unit diagrams compare two iterated collectors with the
same finite collector; Proposition~\ref{prop:egt-finite-collectors} supplies
the required isomorphisms, and their coherence is the coherence of Cartesian
reassociation together with symmetric-monoidal coherence in $\mathcal C$. Symmetry exchanges the input components and sends a tuple to its permuted
tuple, giving the symmetry diagram branchwise.
  \item \textbf{Causal execution.} For functions
$f_j:X_j\to Y_j$, both composites
\[
 \widehat{\prod_jf_j}\circ c_{X_j\,|\,j\in J}
 \quad\text{and}\quad
 c_{Y_j\,|\,j\in J}\circ\bigotimes_j\widehat f_j
\]
have one input event $e_{j,x}$ for each supplied coordinate and one final
output labelled $(f_j(x_j))_j$, depending on all of those inputs. After
hiding the intermediate response events they are isomorphic over the arena. This proves naturality and the lax symmetric-monoidal laws.
\end{enumerate}
\end{proof}

\noindent The lax monoidal structure records the naturality of finite collection.
\begin{center}
\begin{tikzcd}[column sep=large,row sep=large]
D(X)\otimes D(Y)\arrow[r,"c_{X,Y}"]\arrow[d,"\widehat f\otimes\widehat g"']
 & D(X\times Y)\arrow[d,"\widehat{f\times g}"] \\
 D(X')\otimes D(Y')\arrow[r,"c_{X',Y'}"'] & D(X'\times Y')
\end{tikzcd}
\end{center}

\begin{proposition}[The synchronization barrier]
\label{prop:egt-barrier}
One has
\[
 c_{X,Y}\circ d_{X,Y}=\operatorname{id}_{D(X\times Y)}.
\]
For nonempty $X,Y$, the reverse composite is not the identity on
$D(X)\otimes D(Y)$: each of its outputs depends on both inputs.
\end{proposition}
\begin{proof}
\leavevmode
\begin{enumerate}[label=\textup{(\arabic*)},leftmargin=*,itemsep=4pt,topsep=3pt]
  \item \textbf{Secured operational interaction.} In the first composite there is one negative input pair and, after hiding its
two distributed occurrences, one positive output carrying the same pair. Its
source is copycat on $D(X\times Y)$.
  \item \textbf{Reverse implication.} In the reverse composite the collector's
tuple output requires both independent input events. Each distributed output
therefore inherits both prerequisites after hiding. In identity copycat on
$D(X)\otimes D(Y)$, the first positive output can occur after its own negative
input alone, before the second input.
  \item \textbf{Secured operational interaction.} Choosing any values in the two nonempty
sets exhibits a configuration present in copycat and absent from
$d_{X,Y}\circ c_{X,Y}$, so the two strategies are not isomorphic over their
arena.
\end{enumerate}
\end{proof}

\noindent Collection and distribution determine a retraction onto paired data.
\begin{center}
\begin{tikzcd}[column sep=large,row sep=large]
D(X\times Y) \arrow[r,"d_{X,Y}"]\arrow[dr,"\id"']
 & D(X)\otimes D(Y)\arrow[d,"c_{X,Y}"] \\
 & D(X\times Y)
\end{tikzcd}
\end{center}

The retraction identifies one paired report with the completed
collection of two channels. Its reverse retains the synchronization
barrier: neither output can be forwarded before both inputs have arrived.
This causal difference persists in protocols that make decisions between
individual submissions and completion of a joint report.

\subsection{Dependent completion of concurrent interfaces}
\label{sec:egt-dependent-semantics}

An owned interaction starts from the policy sets $\Sigma_i$ and a
realization $g:\prod_i\Sigma_i\to\mathcal C(A,B)$ in the concurrent
category $\mathcal C$. These form its policy family and concurrent
realization, but these alone do not interpret a policy-dependent boundary
or quantified alternatives. The completion below adds those interfaces
without changing what an arrow between original arenas means. We first
index the operational category by ordinary parameters and then embed it
fully faithfully into a presheaf category with Day tensor. The resulting
indexed linear--nonlinear model will supply agent abstraction and the
logical operations used later.

\subsubsection{Indexed families of true-concurrent interactions}

\begin{construction}[Indexed operational category]
The ordinary parameters of an economic model are kept separate from its
linear interfaces. For a set $\Gamma$, let $\mathcal C_\Gamma$ have the same
objects as $\mathcal C$ and hom-sets
\[
 \mathcal C_\Gamma(A,B)
   =\mathbf{Set}\bigl(\Gamma,\mathcal C(A,B)\bigr).
\]
Composition, tensor, identity, symmetry, cups, and caps are defined
pointwise. A function $f:\Delta\to\Gamma$ induces reindexing
\[
 f^*:\mathcal C_\Gamma\longrightarrow\mathcal C_\Delta,
 \qquad (f^*t)_\delta=t_{f(\delta)}.
\]
\end{construction}

\begin{proposition}[Indexed operational semantics]
\label{prop:egt-indexed-operational}
The assignment $\Gamma\mapsto\mathcal C_\Gamma$ and
$f\mapsto f^*$ is contravariantly functorial. Every fiber
$\mathcal C_\Gamma$ is compact closed with the pointwise structure inherited
from $\mathcal C$, and every reindexing functor preserves that structure.
Thus ordinary parameter substitution and linear interaction form an indexed
compact-closed semantics over $\mathbf{Set}$.
\end{proposition}
\begin{proof}
\leavevmode
\begin{enumerate}[label=\textup{(\arabic*)},leftmargin=*,itemsep=4pt,topsep=3pt]
  \item \textbf{Monoidal comparison.} Identity and successive reindexing are the identity and associativity laws
for function composition. Every categorical operation is evaluated at one
parameter value at a time, so reindexing commutes with composition and tensor
and sends the constant identity, symmetry, cup, and cap families to the same
constant families. The compact-closed equations therefore hold pointwise in
every fiber because they hold in $\mathcal C$.
\end{enumerate}
\end{proof}

\noindent Successive parameter substitutions commute with indexed realization.
\begin{center}
\begin{tikzcd}[column sep=large,row sep=large]
\mathcal C_\Gamma \arrow[r,"f^*"]\arrow[dr,"(fg)^*"']
 &\mathcal C_\Delta\arrow[d,"g^*"]\\
 & \mathcal C_\Theta
\end{tikzcd}
\end{center}

\begin{definition}[Owner-indexed parameter category]
\label{def:egt-owner-indexed-parameter-category}
Fix a finite set $N$ of economic agents and put
\[
 \mathbf{Par}_N=\mathbf{Set}^N,
 \qquad
 \Pi_N(\vec\Sigma)=\prod_{i\in N}\Sigma_i.
\]
The product functor $\Pi_N:\mathbf{Par}_N\to\mathbf{Set}$ forgets which
factor is controlled by which owner while retaining the total profile space.
Pulling the indexed operational semantics back along $\Pi_N$ leaves the
true-concurrent operational category unchanged and records exactly the
factorization of ordinary alternatives that will determine deviations.
\end{definition}

\subsubsection{The canonical completion of concurrent interfaces}

Economically, the completion supplies spaces in which policy families,
their dependent alternatives, and predicates can coexist with the original
operational interfaces. Its new objects organize those families; they do
not add new economic decisions to a given arena. Full faithfulness of Yoneda
will ensure that an arrow between original interfaces still denotes exactly
one original concurrent strategy.

\begin{lemma}[Scalars and the obstruction to empty linear limits]
\label{lem:egt-native-empty-obstruction}
The scalar set $\mathcal C(I,I)$ is a singleton. Every arena $A$ has a
state $I\to A$, but $\mathcal C$ has neither an initial nor a terminal
object.
\end{lemma}

\begin{proof}
\leavevmode
\begin{enumerate}[label=\textup{(\arabic*)},leftmargin=*,itemsep=4pt,topsep=3pt]
  \item \textbf{Causal execution.} A total event-structure map to the empty arena has empty source, proving
the scalar assertion. In $A$ take the induced substructure
\[
 A_- = \{e\in A: d\leq_A e\ \Longrightarrow\ \operatorname{pol}(d)=-\}.
\]
It is downward closed. Its inclusion into $A$ is a total event-structure
map and is courteous: an immediate dependency in this downward-closed
substructure is immediate in $A$.
  \item \textbf{Lifting through the comparison.} If a negative event is enabled after
an image configuration, all its strict predecessors belong to that
configuration, so it belongs to $A_-$ and has a unique lift. The inclusion
is therefore a receptive courteous strategy $I\to A$.
  \item \textbf{Causal execution.} If $O$ were initial, its state $s:I\to O$ and the unique $r:O\to I$
would satisfy $rs=1_I$ by scalar uniqueness and $sr=1_O$ by initiality. Then $I$ would be initial. The one-positive-event arena, however, has
distinct empty-source and one-event-source strategies from $I$.
  \item \textbf{Secured operational interaction.} Neither
violates receptivity, and their sources are not isomorphic. This is a
contradiction. Compact duality exchanges initial and terminal objects,
so a terminal object is impossible as well.
\end{enumerate}
\end{proof}


\begin{convention}[Base, size, and finite ownership]
\label{conv:egt-base}
$\mathbf{Set}$ denotes the ordinary set-theoretic parameter category.
The ambient logic is classical. Predicate calculations use its subobject
doctrine, written with Heyting operations; arguments that use classical order
or finite maximization are identified as such. Set-indexed meets, joins, and
power objects are taken in this base.  The operational category $\mathcal C$ has arbitrary countable
arena carriers. For each fixed pair of arenas, its arrows form the set of
countable strategy presentations whose source carriers are subsets of
$\mathbb N$, modulo arena-preserving isomorphism. Every countable source
admits such a presentation, so this convention retains every operational
arrow. Let $\mathcal C_0$ be the full subcategory on arenas whose event
carriers are subsets of $\mathbb N$, and write
$J:\mathcal C_0\hookrightarrow\mathcal C$ for the inclusion.
The category $\mathcal C_0$ is small. Fixed bijections for finite tagged
carriers equip it with a tensor $\otimes_0$ and coherent comparisons
making $J$ strong symmetric monoidal. The empty arena is its unit $I$.
Every countable arena is isomorphic to some $Jc$ by transport of its
structure along an injection into $\mathbb N$. This is an objectwise
existence statement; no simultaneous choice of representatives is made.
The owner set is a fixed finite list $N$. A coalition $C\subseteq N$
inherits its order from this list, and $-C$ denotes the complementary list.
Thus $\Sigma_C=\prod_{i\in C}\Sigma_i$ retains exactly its members'
controlled coordinates. Empty products are singletons and empty conjunctions
are true; the empty coalition changes no coordinate. Collective improvement
predicates below explicitly quantify over nonempty coalitions.
\end{convention}

An empty dependent linear sum or product requires limits and colimits
not present in the original arena category. We use the small monoidal
category $\mathcal C_0$ and the fully faithful, essentially surjective
inclusion $J:\mathcal C_0\hookrightarrow\mathcal C$ from
Convention~\ref{conv:egt-base}. Presentations with carriers in
$\mathbb N$ retain finite data arenas $D(X)$ and residual arenas
as objects of $\mathcal C$. For each original arena, restricted
Yoneda recovers its concurrent arrows. The new colimits and limits permit
dependent interfaces that the operational arena category itself lacks.
Thus the completion enlarges the language of interfaces, not the possible
executions of any already specified original strategy.

\begin{construction}[Presheaves and Day convolution]
\label{con:egt-day-completion}
Put $\widehat{\mathcal C}=[\mathcal C_0^{\mathrm{op}},\mathbf{Set}]$ and
write $yA=\mathcal C(J(-),A)$ for restricted Yoneda. For presheaves
$P,Q$ and $c\in\mathcal C_0$ define
\[
 (P\otimes_D Q)(c)=
 \int^{a,b\in\mathcal C_0}\mathcal C_0(c,a\otimes_0 b)\times P(a)\times Q(b),
 \qquad \widehat I=yI,
\]
\[
 [P,Q]_D(c)=\operatorname{Nat}\bigl(P,Q(c\otimes_0-)\bigr).
\]
The coend is the quotient of the disjoint union over pairs $(a,b)$ by
the relations induced by arrows in either variable. For example, for
$u:a'\to a$, it identifies $(k,P(u)x,z)$ with
$((u\otimes_0 1)k,x,z)$; the other variable is treated identically.
Reindexing in $c$ is precomposition. Limits and colimits of presheaves
are formed objectwise.
\end{construction}

\begin{theorem}[Completion and exact recovery of concurrent arrows]
\label{thm:egt-day-preservation}
The category $\widehat{\mathcal C}$ is complete, cocomplete, and
symmetric monoidal closed under the displayed operations. Tensor
preserves colimits separately. Yoneda is fully faithful and strong
symmetric monoidal, and preserves the original internal hom:
\[
 \begin{aligned}
 \widehat{\mathcal C}(yA,yB)&\cong\mathcal C(A,B),\\
 yA\otimes_D yB&\cong y(A\otimes B),\\
 [yA,yB]_D&\cong y(A\multimap B).
 \end{aligned}
\]
It preserves the compact duals of representable objects. A natural
transformation $\alpha:yA\to yB$ recovers the unique original strategy
\[
 \alpha_c(u)\circ u^{-1}:A\longrightarrow B,
 \qquad u:Jc\xrightarrow{\cong}A.
\]
The result is independent of the chosen isomorphism $u$.
\end{theorem}
\begin{proof}
\leavevmode
\begin{enumerate}[label=\textup{(\arabic*)},leftmargin=*,itemsep=4pt,topsep=3pt]
  \item \textbf{Day-convolution calculation.} Day's convolution theorem supplies the monoidal coherence and closed
structure \citep{Day1970ClosedCategories}. In the present formulas, the coend universal property gives directly
\[
 \widehat{\mathcal C}(R\otimes_D P,Q)
 \cong\widehat{\mathcal C}(R,[P,Q]_D).
\]
Products with fixed sets and coends preserve colimits, proving the
separate cocontinuity of tensor. Objectwise limits and colimits prove
completeness and cocompleteness.
  \item \textbf{Representation and Yoneda decoding.} We first verify restricted Yoneda directly. Choose, for this argument only,
an isomorphism $u:Jc\to A$. Put $f=\alpha_c(u)u^{-1}$.
  \item \textbf{Functoriality and substitution.} For any
$t:Jd\to A$, fullness of $J$ gives the unique arrow $v:d\to c$ with
$Jv=u^{-1}t$. Naturality yields
$\alpha_d(t)=\alpha_c(u)Jv=ft$. Thus $\alpha=yf$, and evaluation
at $u$ proves uniqueness. The recovered $f$ is therefore independent
of the local witness $u$.
  \item \textbf{Day-convolution calculation.} The strong monoidal comparison of $J$ induces canonical maps
$yA\otimes_D yB\to y(A\otimes B)$. For $A=Ja$ and $B=Jb$,
co-Yoneda makes these maps isomorphisms. Every $A,B$ is locally
isomorphic to such a pair, so the same canonical maps are isomorphisms
for all $A,B$.
  \item \textbf{Unit and identity laws.} The unit comparison is likewise an isomorphism;
coherence follows from these canonical maps and that of $J$. For the closed structure the same local argument gives
\[
 [yA,yB]_D(c)\cong\mathcal C(Jc\otimes A,B)
                  \cong\mathcal C(Jc,A\multimap B).
\]
The comparisons are natural and require no global choice of arena
presentations. They carry cups and caps to cups and caps, preserving
their snake equations.
  \item \textbf{Equation reflection under Yoneda.} Full faithfulness proves exact recovery and
reflection of equations between original strategies.
\end{enumerate}
\end{proof}

\noindent The Yoneda comparison intertwines the original tensor product and Day convolution.
\begin{center}
\begin{tikzcd}[column sep=large,row sep=large]
\mathcal C_0\times\mathcal C_0\arrow[r,"y\times y"]\arrow[d,"\otimes_0"']
 & \widehat{\mathcal C}\times\widehat{\mathcal C}\arrow[d,"\otimes_D"]\\
 \mathcal C_0\arrow[r,"y"'] &\widehat{\mathcal C}
\end{tikzcd}
\end{center}

Every presheaf is canonically a colimit of representables:
$P\cong\int^{c\in\mathcal C_0}P(c)\cdot y(Jc)$, where the dot denotes a copower. Thus the
new interfaces are generated from the concurrent ones by specified
categorical constructions. They need not themselves be single event
structures. General presheaves have the proved symmetric monoidal closed
structure; compact duality remains available on the original arena
fragment. The theorem is a preservation and recovery result for semantic
arrows.

\subsubsection{Dependent ordinary and linear fibres}

Dependence in this model is on ordinary parameters. The discrete index category $\Gamma_{\mathrm{disc}}$
indexes a family of presheaf interfaces, so substitution changes the
parameter at which a linear object is read. The linear context still
accounts for resource occurrences. This division allows dependent
quantification without turning a linear resource into an ordinary
variable.

\begin{construction}[The two indexed categories]
\label{con:egt-dependent-fibres}
For a set $\Gamma$ define
\[
 \mathsf T(\Gamma)=\mathbf{Set}^{\Gamma_{\mathrm{disc}}},
 \qquad
 \mathsf L(\Gamma)=\widehat{\mathcal C}^{\Gamma_{\mathrm{disc}}}.
\]
Objects of $\mathsf T(\Gamma)$ are ordinary dependent families; objects
of $\mathsf L(\Gamma)$ are linear dependent families. Their morphisms
are families of functions and natural transformations, respectively.
For $f:\Delta\to\Gamma$ put $(f^*A)_\delta=A_{f(\delta)}$.
Composition, tensor, internal hom, and their comparisons in the linear
fibres are pointwise. The unit is the constant family $\widehat I$.
An ordinary family $X$ has context extension
$p_X:\Gamma.X=\coprod_\gamma X_\gamma\to\Gamma$.
\end{construction}

\begin{proposition}[Substitution and the original indexed fragment]
\label{prop:egt-dependent-substitution}
These assignments are contravariantly functorial. Ordinary reindexing
preserves the cartesian closed structure, and linear reindexing preserves
the symmetric monoidal closed structure. Pointwise Yoneda embeds
$\mathcal C^{\Gamma_{\mathrm{disc}}}$ fully faithfully into
$\mathsf L(\Gamma)$, compatibly with substitution. Its constant-arena
part is exactly $\mathcal C_\Gamma$.
\end{proposition}
\begin{proof}
\leavevmode
\begin{enumerate}[label=\textup{(\arabic*)},leftmargin=*,itemsep=4pt,topsep=3pt]
  \item \textbf{Functoriality and substitution.} Both identity and composite substitution are identities of index
precomposition. Every listed structure is evaluated at a single index,
so reindexing preserves its defining universal property and comparison.
  \item \textbf{Representation and Yoneda decoding.} The hom-set of constant representables $yA,yB$ is
\[
 \prod_{\gamma\in\Gamma}\widehat{\mathcal C}(yA,yB)
 \cong\mathbf{Set}(\Gamma,\mathcal C(A,B)).
\]
Theorem~\ref{thm:egt-day-preservation}, applied at each index, proves
the remaining preservation and full-faithfulness assertions.
\end{enumerate}
\end{proof}

\begin{proposition}[Indexed operational equations]
\label{thm:egt-operational-soundness}
Composition, tensor, and currying of parameterized concurrent arrows
commute with parameter substitution. They satisfy associativity, units,
interchange, and the closed-structure equations
\[
 \operatorname{ev}\circ(\Lambda_A t\otimes1_A)=t,
 \qquad
 \Lambda_A\bigl(\operatorname{ev}\circ(u\otimes1_A)\bigr)=u.
\]
The same equations hold in every completed linear fibre.
\end{proposition}
\begin{proof}
\leavevmode
\begin{enumerate}[label=\textup{(\arabic*)},leftmargin=*,itemsep=4pt,topsep=3pt]
  \item \textbf{Adjunction and transposition.} In the original fragment each equation holds at every parameter by
compact closure, with the two displayed equations reducing to the snake
identities. In the completion they are the triangle identities for the
closed adjunction. Reindexing evaluates the identical map at the
substituted index, proving compatibility.
\end{enumerate}
\end{proof}

Ordinary parameters may be reused to describe or compare different
interactions; a linear occurrence must instead be consumed according to
its process. The functor $\mathsf{Val}$ extracts monoidal states of a
linear interface, not economic values, while $\mathsf{Lin}$ makes an
ordinary family of chosen states available linearly. Their adjunction is
the required bridge between policy descriptions and interactive resources;
it introduces no copying operation on arbitrary resources.

\begin{construction}[Values and linearization]
\label{con:egt-lnl}
Define functors between the two fibres by
\[
 (\mathsf{Val}_\Gamma A)_\gamma
    =\widehat{\mathcal C}(\widehat I,A_\gamma)=A_\gamma(I),
 \qquad
 (\mathsf{Lin}_\Gamma X)_\gamma
    =\coprod_{x\in X_\gamma}\widehat I.
\]
For a linear family its comprehension is
\[
 p_A:\Gamma.A=
       \coprod_{\gamma\in\Gamma}(\mathsf{Val}_\Gamma A)_\gamma
       \longrightarrow\Gamma.
\]
The generic element at $(\gamma,a)$ is the state $a:\widehat I\to A_\gamma$.
\end{construction}

\begin{theorem}[Constructed comprehension and linear--nonlinear adjunction]
\label{thm:egt-lnl-comprehension}
There is a substitution-stable adjunction
$\mathsf{Lin}_\Gamma\dashv\mathsf{Val}_\Gamma$ with
$\mathsf{Lin}_\Gamma$ strong symmetric monoidal. Linear comprehension
has the representing bijection
\[
 \mathsf L(\Delta)(\widehat I,f^*A)
 \cong\mathbf{Set}/\Gamma(f,p_A).
\]
Moreover
$\mathsf{Val}_\Gamma\mathsf{Lin}_\Gamma X\cong X$, so
$\mathsf{Lin}_\Gamma$ is fully faithful and
$\Gamma.\mathsf{Lin}_\Gamma X\cong\Gamma.X$ over $\Gamma$.
\end{theorem}

\noindent Comprehension transports linear states into ordinary sections of the displayed family.
\begin{center}
\begin{tikzcd}[column sep=large,row sep=large]
\mathsf L(\Delta)(\widehat I,f^*A)\arrow[r,"\cong"]\arrow[d,"g^*"']
 & \mathbf{Set}/\Gamma(f,p_A)\arrow[d,"(-)\circ g"]\\
 \mathsf L(\Theta)(\widehat I,(fg)^*A)\arrow[r,"\cong"']
 &\mathbf{Set}/\Gamma(fg,p_A)
\end{tikzcd}
\end{center}


\begin{proof}
\leavevmode
\begin{enumerate}[label=\textup{(\arabic*)},leftmargin=*,itemsep=4pt,topsep=3pt]
  \item \textbf{Frobenius and coproduct comparison.} The coproduct universal property gives
\[
 \mathsf L(\Gamma)(\mathsf{Lin}_\Gamma X,A)
 \cong\prod_\gamma\mathbf{Set}
              (X_\gamma,\widehat{\mathcal C}(\widehat I,A_\gamma))
 =\mathsf T(\Gamma)(X,\mathsf{Val}_\Gamma A).
\]
Separate cocontinuity of tensor gives
$\mathsf{Lin}(X\times Y)\cong\mathsf{Lin}(X)\otimes_D\mathsf{Lin}(Y)$
and $\mathsf{Lin}(1)\cong \widehat I$.
  \item \textbf{Monoidal comparison.} On summands the comparisons are unitors
and the bijection between pairs of indices; their coherence follows on
each summand from monoidal coherence. Evaluation at $I$ yields
\[
 (\mathsf{Val}\mathsf{Lin}X)_\gamma
 =\coprod_{x\in X_\gamma}\mathcal C(I,I)\cong X_\gamma
\]
by Lemma~\ref{lem:egt-native-empty-obstruction}. This is the adjunction
unit, proving full faithfulness.
  \item \textbf{Secured operational interaction.} For comprehension, a family of states
$a_\delta:\widehat I\to A_{f(\delta)}$ corresponds to the map
$\delta\mapsto(f(\delta),a_\delta)$; projection recovers the family. Pullback reindexes these same pairs. All the constructions commute with
index precomposition, establishing their substitution comparisons and
coherence.
\end{enumerate}
\end{proof}


\begin{construction}[The exponential]
\label{con:egt-exponential}
Put $!_\Gamma=\mathsf{Lin}_\Gamma\mathsf{Val}_\Gamma$. The counit
$\epsilon_A:!A\to A$ evaluates the state labelling each summand; the
comultiplication is $\mathsf{Lin}(\eta_{\mathsf{Val}A})$.
The adjunction triangles prove the comonad identities. The diagonal and
terminal map of $\mathsf{Val}A$, transported through the strong monoidal
functor $\mathsf{Lin}$, give
$!A\to !A\otimes_D !A$ and $!A\to \widehat I$ and their commutative-comonoid laws.
Thus these structural maps are constructed, and are stable under
substitution. Here $!$ is idempotent because $\eta$ is invertible.
This is the monoidal-adjunction account of linear--nonlinear
semantics \citep{Benton1994MixedLNL}.
For completeness, the lax monoidal comparison on values sends a pair
of states $a:\widehat I\to A$, $b:\widehat I\to B$ to
\[
 \widehat I\cong\widehat I\otimes_D\widehat I
       \xrightarrow{a\otimes_D b}A\otimes_D B;
\]
its unit is the identity state of $\widehat I$. Applying
$\mathsf{Lin}$ gives the monoidal comparison
$!A\otimes_D!B\to!(A\otimes_D B)$. Its coherence, and compatibility
with the comonad maps, can be checked on each pair of state-labelled
summands: both sides evaluate the same tensor of states. Also
\[
 !(A\mathbin{\&}B)\cong!A\otimes_D!B,
 \qquad !\top\cong\widehat I,
\]
because values preserve products and $\mathsf{Lin}$ is strong monoidal.
\end{construction}

The exponential renders named states reusable through the ordinary
interface. It does not turn the physical response arena $D(X)$ into a
copower: $\mathsf{Lin}(1)=\widehat I$, whereas $yD(1)$ has more than one
state and is not isomorphic to $\widehat I$. The copying and collection
strategies on finite data therefore retain the causal behaviour proved
above, independently of the abstract linear--nonlinear adjunction.

\subsubsection{Dependent adjoints, additives, and equality}

\begin{construction}[Dependent sums and products]
\label{con:egt-linear-dependent-adjoints}
For $f:\Delta\to\Gamma$ define
\[
 (\Sigma_f^\ell A)_\gamma
       =\coprod_{\delta:\,f(\delta)=\gamma}A_\delta,
 \qquad
 (\Pi_f^\ell A)_\gamma
       =\prod_{\delta:\,f(\delta)=\gamma}A_\delta.
\]
They act on arrows componentwise using coproduct injections and product
projections. The identical set formulas define ordinary functors
$\Sigma_f^{\mathrm c},\Pi_f^{\mathrm c}$ on $\mathsf T$.
\end{construction}

\begin{theorem}[Adjunctions, base change, and Frobenius]
\label{thm:egt-linear-dependent-adjoints}
For every $f$,
$\Sigma_f^\ell\dashv f^*\dashv\Pi_f^\ell$.
Both adjoints satisfy Beck--Chevalley for pullback squares. There is a
natural Frobenius isomorphism
\[
 \Sigma_f^\ell(f^*B\otimes_D A)
       \cong B\otimes_D\Sigma_f^\ell A.
\]
The comparisons are coherent with identity and composite substitution.
The ordinary adjoints have the corresponding properties, with cartesian
product in place of tensor. All empty fibres are included.
\end{theorem}
\begin{proof}
\leavevmode
\begin{enumerate}[label=\textup{(\arabic*)},leftmargin=*,itemsep=4pt,topsep=3pt]
  \item \textbf{Frobenius and coproduct comparison.} The coproduct universal property gives
\[
 \mathsf L(\Gamma)(\Sigma_f^\ell A,B)
 \cong\prod_\gamma\prod_{f(\delta)=\gamma}
             \widehat{\mathcal C}(A_\delta,B_\gamma)
 \cong\mathsf L(\Delta)(A,f^*B).
\]
Replacing the coproduct universal property by the product universal
property gives the right adjunction.
  \item \textbf{Base-change comparison.} For a pullback square
$\Delta'\xrightarrow{a}\Delta$, $\Gamma'\xrightarrow{b}\Gamma$
with vertical arrows $f',f$, the map $a$ is a bijection between the
fibre of $f'$ at $\gamma'$ and the fibre of $f$ at $b\gamma'$. It therefore induces
$b^*\Sigma_f^\ell\cong\Sigma_{f'}^\ell a^*$ and
$b^*\Pi_f^\ell\cong\Pi_{f'}^\ell a^*$. Successive pullback comparisons compose the same index bijections, so
their coherence follows.
  \item \textbf{Base-change comparison: completion.} Frobenius is, at each $\gamma$, the
distribution of $B_\gamma\otimes_D-$ over the coproduct indexed by that
fibre. Its comparison on each summand is the identity; this also proves
compatibility with base change. The set-valued proof is identical.
\end{enumerate}
\end{proof}

\noindent The pullback square supplies Beck--Chevalley comparisons for both dependent adjoints.
\begin{center}
\begin{tikzcd}[column sep=large,row sep=large]
\Delta'\arrow[r,"a"]\arrow[d,"f'"']\arrow[dr,phantom,"\lrcorner",very near start]
 &\Delta\arrow[d,"f"]\\
 \Gamma'\arrow[r,"b"']&\Gamma
\end{tikzcd}
\end{center}

\begin{construction}[Additives and equality]
\label{con:egt-linear-additives-equality}
Finite products and coproducts in each linear fibre give
$A\mathbin{\&}B$, $\top$, $A\oplus B$, and $0$. The object $\top$ is
the terminal presheaf; the tensor unit is $\widehat I=yI$.
For an ordinary family $X$ let
$d:\Gamma.X\to\Gamma.X.X$ be its diagonal and put
\[
 \operatorname{Id}^{\ell}_X=\Sigma_d^\ell \widehat I.
\]
Its fibre at $(\gamma,x,x')$ is $\widehat I$ if $x=x'$ and $0$ otherwise.
It equals $\mathsf{Lin}(\operatorname{Eq}_X)$, where
$\operatorname{Eq}_X$ is the ordinary equality family. Reflexivity is
the identity of $\widehat I$ on the diagonal. For any linear $E,D$ over
$\Gamma.X.X$, Frobenius and adjunction give
\[
 \mathsf L(\Gamma.X.X)(E\otimes_D\operatorname{Id}^{\ell}_X,D)
 \cong\mathsf L(\Gamma.X)(d^*E,d^*D).
\]
This is its equality-elimination universal property; the adjunction
identities give evaluation and uniqueness. All these constructions are
stable under substitution by Theorem~\ref{thm:egt-linear-dependent-adjoints}.
\end{construction}

\begin{proposition}[Recovery of ordinary dependent structure]
\label{prop:egt-ordinary-recovery}
Every ordinary context extension is the comprehension of a linearized
family. Moreover
\[
 \mathsf{Val}(\Sigma_f^\ell\mathsf{Lin}X)
       \cong\Sigma_f^{\mathrm c}X,
 \qquad
 \mathsf{Val}(\Pi_f^\ell\mathsf{Lin}X)
       \cong\Pi_f^{\mathrm c}X.
\]
For $p:\Gamma.X\to\Gamma$ and $B\in\mathsf L(\Gamma)$,
\[
 \Sigma_p^\ell p^*B\cong\mathsf{Lin}X\otimes_D B,
 \qquad
 \Pi_p^\ell p^*B\cong[\mathsf{Lin}X,B]_D.
\]
\end{proposition}
\begin{proof}
\leavevmode
\begin{enumerate}[label=\textup{(\arabic*)},leftmargin=*,itemsep=4pt,topsep=3pt]
  \item \textbf{Frobenius and coproduct comparison.} Evaluation at $I$ preserves objectwise products and coproducts, and
$\mathsf{Val}\mathsf{Lin}X\cong X$. This proves the first two
comparisons and, by Theorem~\ref{thm:egt-lnl-comprehension}, the context
claim. The constant-family sum distributes $B$ over one copy of $\widehat I$
per $x$; the product comparison follows by applying closedness to that
copower. These calculations are natural in all parameters.
\end{enumerate}
\end{proof}

\noindent Dependent ordinary families are transported along the pullback of context extension.
\begin{center}
\begin{tikzcd}[column sep=large,row sep=large]
\Delta\times_\Gamma(\Gamma.X)\arrow[r]\arrow[d]\arrow[dr,phantom,"\lrcorner",very near start]
 &\Gamma.X\arrow[d,"p"]\\
 \Delta\arrow[r,"f"']&\Gamma
\end{tikzcd}
\end{center}

\begin{theorem}[The categorical dependent linear model]
\label{thm:egt-dependent-model}
The constructed ordinary and linear indexed categories, comprehension,
and substitution-stable monoidal adjunction form a dependent
linear--nonlinear model over $\mathbf{Set}$. They carry ordinary
dependent sums, products, and extensional equality; linear
multiplicatives and additives; the exponential of
Construction~\ref{con:egt-exponential}; and the dependent linear sums,
products, and equality just constructed. Dependence is on ordinary
parameters. Original concurrent arena families embed fully faithfully,
preserving their operational structure, and all original ordinary
families are recovered by $\mathsf{Val}\mathsf{Lin}$.
\end{theorem}
\begin{proof}
\leavevmode
\begin{enumerate}[label=\textup{(\arabic*)},leftmargin=*,itemsep=4pt,topsep=3pt]
  \item \textbf{Adjunction and transposition.} Proposition~\ref{prop:egt-dependent-substitution} supplies the indexed
monoidal structure and its substitution laws. Theorem~\ref{thm:egt-lnl-comprehension} constructs the comprehension
representing property and the monoidal adjunction. Theorem~\ref{thm:egt-linear-dependent-adjoints} proves the adjoints,
Beck--Chevalley, and Frobenius, including their coherence.
  \item \textbf{Assembly of dependent linear structure.} Construction~\ref{con:egt-linear-additives-equality} gives equality and
additives. These are the structural requirements in the indexed
comprehension account of dependent linear semantics
\citep[Sections~3--4]{Vakar2015LinearFrameworks}. The preservation assertions are
Theorem~\ref{thm:egt-day-preservation} and
Proposition~\ref{prop:egt-ordinary-recovery}.
\end{enumerate}
\end{proof}

\begin{proposition}[Transport of dependent linear realizations]
\label{prop:egt-dependent-realization-transport}
For $A,B\in\mathsf L(W)$ and any $t:U\to W$, an ordinary section
$\rho:1_W\to\mathsf{Val}_W[A,B]_D$ determines a section
\[
 t^*\rho:1_U\to\mathsf{Val}_U[t^*A,t^*B]_D.
\]
Under the linear--nonlinear adjunction, its linear transpose is
$t^*\widetilde\rho:t^*A\to t^*B$. These comparisons commute with
composition of state maps, and, for representable fibres, recover the
original concurrent strategies pointwise by Yoneda.
\end{proposition}
\begin{proof}
\leavevmode
\begin{enumerate}[label=\textup{(\arabic*)},leftmargin=*,itemsep=4pt,topsep=3pt]
  \item \textbf{Fibrewise calculation.} At $u\in U$ the transported state is the arrow
$\rho_{t(u)}:A_{t(u)}\to B_{t(u)}$. The adjunction is pointwise, so
transposing before or after reindexing gives that same component. Successive substitutions select the same component. For representables,
$\widehat{\mathcal C}(yA,yB)\cong\mathcal C(A,B)$ proves recovery.
\end{enumerate}
\end{proof}

The transport result will be needed when an epistemic action changes
the ordinary state index. Its reindexed family still denotes the same
pointwise concurrent arrow; later residualization can therefore track an
interaction across the update without replacing it by an abstract scalar.

\subsubsection{Policy-parameterized interactions and their composition}

\begin{construction}[Dependent owned realization]
\label{con:egt-dependent-owned-realization}
Fix $A,B\in\mathsf L(\Gamma)$ and ordinary policy families
$\vec\Sigma=(\Sigma_i)_{i\in N}$; put $\Sigma=\prod_i\Sigma_i$.
A realization is an ordinary morphism
\[
 g:\Sigma\longrightarrow\mathsf{Val}_\Gamma[A,B]_D.
\]
Since $\mathsf{Val}[A,B]_D\cong
(\widehat{\mathcal C}(A_\gamma,B_\gamma))_\gamma$, it specifies
$g_{\gamma,\sigma}:A_\gamma\to B_\gamma$. Adjunction and closedness
give the equivalent single linear morphism
\[
 \widetilde g:\mathsf{Lin}_\Gamma\Sigma\otimes_D A\longrightarrow B.
\]
An ownerwise map $r_i:\Theta_i\to\Sigma_i$ reparametrizes $g$ by
$\prod_i r_i$; its linear form is precomposition with
$\mathsf{Lin}(\prod_i r_i)\otimes_D1_A$.
\end{construction}

\begin{proposition}[Categorical construction of dependent game composition]
\label{prop:egt-dependent-game-composition}
Dependent owned realizations have sequential and parallel composition,
identity, and ownerwise reparametrizations. With the canonical ownerwise
product comparisons they form a symmetric monoidal bicategory at each
$\Gamma$, stable under substitution of $\Gamma$.
On original representable interfaces the construction recovers precisely
the concurrent open arenas and their compositions of
Section~\ref{sec:egt-players}.
\end{proposition}

\noindent The realization of a composite is computed by composing the policy-indexed arrows.
\begin{center}
\begin{tikzcd}[column sep=large,row sep=large]
\Sigma\times T\arrow[r,"g\times h"]\arrow[d,"\cong"']
 & \mathcal C(A,B)\times\mathcal C(B,C)\arrow[d,"\circ"]\\
 \prod_i(\Sigma_i\times T_i)\arrow[r,"h\circ g"']
 & \mathcal C(A,C)
\end{tikzcd}
\end{center}


\begin{proof}
\leavevmode
\begin{enumerate}[label=\textup{(\arabic*)},leftmargin=*,itemsep=4pt,topsep=3pt]
  \item \textbf{Composition through linearized policy sums.} For $\widetilde g:\mathsf{Lin}\Sigma\otimes_D A\to B$ and
$\widetilde h:\mathsf{Lin}\mathcal T\otimes_D B\to C$, form
\[
 \mathsf{Lin}(\Sigma\times\mathcal T)\otimes_D A
 \cong\mathsf{Lin}\mathcal T\otimes_D\mathsf{Lin}\Sigma\otimes_D A
 \xrightarrow{1\otimes\widetilde g}
 \mathsf{Lin}\mathcal T\otimes_D B
 \xrightarrow{\widetilde h}C.
\]
Its policy component is $h_{\gamma,\tau}g_{\gamma,\sigma}$.
  \item \textbf{Ownerwise alternatives.} Parallel composition uses the strong monoidal comparison for
$\mathsf{Lin}$ and symmetry to apply the two realized arrows to their
respective inputs; its component is
$g_{\gamma,\sigma}\otimes_D k_{\gamma,\rho}$. For each owner use the paired policy family, so the total-profile
comparison is the canonical rearrangement of those pairs. Identity uses
singleton policies.
  \item \textbf{Ownerwise alternatives: second comparison.} A $2$-cell is an ownerwise map preserving the
realization, and horizontal and parallel composition take products of
these maps. Associators, unitors, and interchange rearrange the same
named policy coordinates; their realization equations are the category
and monoidal equations in $\mathsf L(\Gamma)$. Thus every coherence
diagram has identical coordinate maps and identical realizations.
  \item \textbf{Representation and Yoneda decoding.} Substitution evaluates these formulas at the substituted index. Finally $\mathsf{Val}[yA,yB]_D\cong\mathcal C(A,B)$ and Yoneda
preserves composition and tensor, proving exact recovery.
\end{enumerate}
\end{proof}


The owner-indexed factorization remains displayed after the dependent
realization is curried. A diagonal can require two components to use one
policy description, but changing the reference parameter does not restrict
which alternative policies exist. The next section makes agent types from
these policy families; subsequent state-context and comparison constructions
will distinguish reference substitution from changes of controlled
alternatives.

\subsection{The indexed model underlying agent types}
\label{agt-sec-model}

\subsubsection{The indexed model used by the agent calculus}

The operational category, its restricted Yoneda completion, and its
substitution-stable dependent linear--nonlinear structure have already been
constructed in \cref{thm:egt-base,thm:egt-day-preservation,thm:egt-dependent-model}.
In the present notation the canonical comparisons are
\begin{equation}
  \cV(yA,yB)\cong\cC(A,B),\qquad
  yA\otimes yB\cong y(A\otimes B).
  \label{agt-eq-yoneda}
\end{equation}
Put $\cT(W)=\Set^{W_{\mathrm{disc}}}$,
$\cL(W)=\cV^{W_{\mathrm{disc}}}$ and $\cP(W)=\Sub(W)$.
The functors $\Lin X=\coprod_{x\in X}\I$ and
$\Val A=\cV(\I,A)$, and their indexed forms, are those of
Construction~\ref{con:egt-lnl}; in particular
$\Lin(X\times Y)\cong\Lin X\otimes\Lin Y$ and
$\Val\Lin X\cong X$ by \cref{thm:egt-lnl-comprehension}.
For $f:V\to W$ the previously constructed linear adjoints have fibres
\begin{equation}
 (\Sigma_f^\ell A)_w=\coprod_{f(v)=w}A_v,
 \qquad(\Pi_f^\ell A)_w=\prod_{f(v)=w}A_v.
 \label{agt-eq-indexed-adjoints}
\end{equation}
The ordinary adjoints use set sums and products; the predicate adjoints
use existential and universal quantification along the same fibres.

\begin{theorem}[Indexed adjunctions and substitution]
\label{agt-thm-indexed}
There are adjunctions $\Sigma_f^\ell\dashv f^*\dashv\Pi_f^\ell$,
Beck--Chevalley isomorphisms for pullback squares, and a coherent
Frobenius isomorphism
\[
 \Sigma_f^\ell(f^*B\otimes A)\cong B\otimes\Sigma_f^\ell A.
\]
These include empty fibres and commute with successive substitutions.
Ordinary and predicate versions hold as well.
\end{theorem}
\begin{proof}
\leavevmode
\begin{enumerate}[label=\textup{(\arabic*)},leftmargin=*,itemsep=4pt,topsep=3pt]
  \item \textbf{Base-change comparison.} This is \cref{thm:egt-linear-dependent-adjoints}, applied to the same
pointwise fibres. Its proof constructs the adjoints from their universal
properties, Beck--Chevalley by the pullback-induced bijections of index
fibres, and Frobenius by separate cocontinuity of Day tensor. The ordinary comparisons use the same fibrewise set formulas.
  \item \textbf{Base-change comparison: completion.} For
predicates, universal quantification is the product of subterminal
families, whereas existential quantification is the support (image)
of their indexed coproduct; their Beck--Chevalley and Frobenius laws
follow from the same fibre bijections and the predicate adjunctions
proved in \cref{thm:fib-comprehension}.
\end{enumerate}
\end{proof}

\begin{proposition}[Linearization as a dependent sum of units]
\label{agt-prop-lin-dependent-sum}
For an ordinary family $X$ over $\Gamma$ and its comprehension
$p:\Gamma.X\to\Gamma$ there are substitution-natural isomorphisms
\[
 \Lin_\Gamma X\cong\Sigma_p^\ell\I_{\Gamma.X},\qquad
 \Lin_\Gamma X\otimes B\cong\Sigma_p^\ell p^*B.
\]
They interpret $\Lin$-elimination with an independent linear continuation.
\end{proposition}
\begin{proof}
\leavevmode
\begin{enumerate}[label=\textup{(\arabic*)},leftmargin=*,itemsep=4pt,topsep=3pt]
  \item \textbf{Frobenius and coproduct comparison.} The first fibre is the coproduct of one tensor unit for each
$x\in X_\gamma$, by the definitions of $\Lin$ and $\Sigma_p^\ell$;
the second comparison is the Frobenius isomorphism of
\cref{thm:egt-linear-dependent-adjoints}. Both comparisons are
natural in $\gamma$ and preserve substitution by Beck--Chevalley.
\end{enumerate}
\end{proof}

\subsubsection{Closed reindexing and dependent-product coherence}

\begin{lemma}[Closed reindexing in the concrete indexed model]
\label{agt-lem-closed-reindexing}
For every ordinary substitution $f:V\to W$, the canonical comparison
$f^*[A,B]\to[f^*A,f^*B]$ is a natural isomorphism, coherent under
successive substitutions.
\end{lemma}
\begin{proof}
\leavevmode
\begin{enumerate}[label=\textup{(\arabic*)},leftmargin=*,itemsep=4pt,topsep=3pt]
  \item \textbf{Fibrewise calculation.} At $v\in V$ both sides are canonically the same internal hom:
\[
 (f^*[A,B])_v=[A,B]_{f(v)}
 =[A_{f(v)},B_{f(v)}]
 =[(f^*A)_v,(f^*B)_v].
\]
The comparison and its inverse are the identities on these fibres. Both orders of successive reindexing select $A_{f(g(u))}$ and
$B_{f(g(u))}$, proving naturality and coherence, including empty
indexing sets. 
\end{enumerate}
\end{proof}

\noindent Reindexing commutes with the internal hom by the identity on each fibre.
\begin{center}
\begin{tikzcd}[column sep=large,row sep=large]
f^*[A,B]\arrow[r,"\cong"]\arrow[d,"g^*"']
 &\left[f^*A,f^*B\right]\arrow[d,"g^*"]\\
 (fg)^*[A,B]\arrow[r,"\cong"']
 &\left[(fg)^*A,(fg)^*B\right]
\end{tikzcd}
\end{center}

For the dependent-product Beck--Chevalley comparison in a pullback
$V'=W'\times_WV$, both sides at $w'\in W'$ are products over the
same fibre $\{v:f(v)=g(w')\}$. The comparison is uniquely specified by
its product projections; repeated base change selects the same factors.
The corresponding sum comparison uses the same bijection and coproduct
injections. Thus the explicit indexed formulas verify the coherence
already established in \cref{thm:egt-linear-dependent-adjoints}.

\section{Agents as First-Class Linear Objects}\label{sec:paper-2}

\subsection{Agent types and distinct levels of control}
\label{agt-sec-orientation}

The dependent model can now represent an agent as a linear interaction
parameterized by an owner-indexed family of complete policies. This
follows the distinction between strategy parameters and realization in
\citet{CapucciEtAl2022Agency}. Ownership determines deviations; operational
ports determine how the process interacts; observations determine what can
be known; and comparison interfaces determine what is preferred. None of
these roles is recoverable merely by treating an individual as a label in a
modal assertion.

\paragraph{Four distinctions in the agent calculus.} An \emph{owner identity} indexes control and can occur at several decision
positions. A \emph{linear occurrence} of an agent is a resource in a typing
context; its use must be accounted for. A \emph{complete policy} is an ordinary
parameter specifying responses at all relevant decision positions, including
positions not reached during a particular execution. A \emph{situated agent}
is an agent presentation together with a legal domain, a reference-state
space, observation maps, and comparison data.

In particular, two decisions with the same owner contribute two coordinates
to that owner's complete policy. They do not become two independent owners.
Likewise, a set of alternative policies does not supply multiple concurrent
copies of the agent implementing them. The former is an ordinary family of
possibilities; the latter would require additional linear resources.

\paragraph{Constructed and supplied interfaces.}
The calculus interprets ordinary and linear constructors, but a bare
concurrent arrow does not determine an economic preference, an
institutional permission rule, or an observation never delivered by the
process. Such data enter as typed maps or subobjects. Where a protocol does
produce a signal, record, or certificate, its observation and legal
interfaces are constructed by images, pullbacks, and dependent
quantification. This distinguishes the information supplied to a model
from the information derived within it.

\paragraph{The agent-type construction.}
\begin{equation}
\bigl|\Ag_{\sig}(A,B)\bigr|=[\Lin\Sigma\otimes A,B].
\label{agt-eq-agent-overview}
\end{equation}
The signature keeps ownership displayed rather than attempting to
recover it from the underlying linear object. In this sense
``first-class'' concerns the linear occurrence and higher-order use of
the agent resource; ownership is a displayed index of its typing judgment,
not a property of its underlying object alone.
\Cref{agt-thm-agent-dependent-product,agt-thm-owned-equivalence}
identify dependent abstraction and owned presentations;
\cref{agt-thm-agent-coherence,agt-thm-regroup,agt-thm-delegation-image}
then compare composition, regrouping, and implemented choices. The
subsequent contextual and operational sections use precisely this
representation rather than adding an independent meaning for an agent.

Thus an agent is not merely an available policy set or an operational
strategy. Its linear type can be supplied to another process, while the
displayed signature records which owner controls each choice. Polarity
continues to describe interface behaviour; rationality belongs to the
later comparison predicate. The representation theorem establishes the
exact relation between states of this type and owned concurrent
realizations.

\subsection{Economic vocabulary: policies, outcomes, and information}
\label{sec:egt-economic-vocabulary}

To read the type of an agent economically, we must distinguish a
complete contingent policy from the concurrent arrow implementing it.
Ownership identifies who may replace a policy coordinate; comparison and
observation specify, separately, how that replacement is evaluated and
what information the owner possesses. The definitions below fix this
vocabulary before Section~\ref{sec:egt-agents} packages the resulting
data into a formal agent interface.

\begin{definition}[Agent identities, policies, and profiles]
\label{def:egt-economic-policy}
An \textbf{economic game} in ordinary strategic form consists of a finite
owner list $N$, sets $\Sigma_i$ of available policies, an outcome set $Z$,
a realization, and strict comparison relations:
\[
 \Sigma=\prod_{i\in N}\Sigma_i,\qquad
 z:\Sigma\longrightarrow Z,\qquad
 \prec_i\;\hookrightarrow\;Z\times Z\quad(i\in N).
\]
An \textbf{agent identity} $i$ controls the coordinate $\Sigma_i$.
The term \textbf{owner} describes this coordinate control; \textbf{player}
is the customary game-theoretic name for the agent when its incentives are
considered. The same identity may own several decisions in one process.
A \textbf{pure policy} $\sigma_i\in\Sigma_i$ is a complete contingent plan
for the decisions of $i$; a \textbf{profile} is
$\sigma=(\sigma_i)_{i\in N}\in\Sigma$. Write $\sigma_{-i}$ for the
profile of the other owners.

More explicitly, at a decision node $v$ with effective observation set
$O_v$ and action set $A_v$, a decision policy has type
$\pi_v:O_v\to A_v$ and chooses the \textbf{action}
$\pi_v(o)\in A_v$ upon observation $o$. For a protocol with node set
$V$ and ownership $i(v)$, the complete owner policy space is
\[
 \Sigma_i=\prod_{v\in V:\,i(v)=i}A_v^{O_v},
\]
including decisions at nodes not reached by realized play. The general
open-arena interface takes the sets $\Sigma_i$ as parameters; its
operational realization has type $g:\Sigma\to\mathcal C(A,B)$.
An economic policy is an element of $\Sigma_i$, whereas an
\textbf{operational concurrent strategy} is an event-structure morphism
representing an arrow $g(\sigma)$ of $\mathcal C$. These are distinct
kinds of objects.
\end{definition}
\begin{definition}[Coalitions and deviations]
\label{def:egt-economic-coalition}
A \textbf{coalition} is a subset $C\subseteq N$ of owners, with
complement $-C=N\setminus C$ and policy carriers
$\Sigma_C=\prod_{i\in C}\Sigma_i$ and
$\Sigma_{-C}=\prod_{i\notin C}\Sigma_i$. Its \textbf{deviation relation}
is the kernel pair of the outsider projection:
\[
 \Sigma\xleftarrow{p_1}
 D_C=\Sigma\times_{\Sigma_{-C}}\Sigma
 \xrightarrow{p_2}\Sigma,\qquad
 (\sigma,\tau)\in D_C\iff\sigma_{-C}=\tau_{-C}.
\]
A coalition deviation from a reference $\sigma$ is an element of the
fibre $\mathsf{Alt}_C(\sigma)=p_1^{-1}(\sigma)$; for the ordinary outcome
realization $z$ its paired outcome is
$\delta_C^z=(zp_1,zp_2):D_C\to Z^2$.
A \textbf{unilateral deviation} has $C=\{i\}$; the
\textbf{grand coalition} has $C=N$. For $C=\varnothing$, $D_C$ is the
diagonal, so no policy changes. The word deviation denotes a comparison
of available alternatives, not a temporal adjustment.

Joint control of these coordinates is distinct from a collective
preference, pooled access to observations, or common resource endowments.
Joint choice and guaranteed outcomes against arbitrary outsider choices
therefore have different quantifier structures. Control of material resources is represented separately through
occurrence-preserving processes and their certificates.
\end{definition}
\begin{definition}[Preferences, utilities, and valuations]
\label{def:egt-preferences-utility}
For an outcome set $Z$, an agent's \textbf{strict preference}
is a relation $\prec_i\hookrightarrow Z\times Z$; $z\prec_i z'$ means
that $z'$ is strictly better than $z$ for owner $i$. A real-valued
\textbf{utility} or \textbf{payoff} is a map $u_i:Z\to\mathbb R$ that
represents this relation when
\[
 z\prec_i z'\quad\Longleftrightarrow\quad u_i(z)<u_i(z').
\]
For an allocation carrier $X$, a numerical \textbf{valuation} is a map
$v_i:X\to\mathbb R$ assigning the value of an allocation before
separately specified payments. Alternatively, when owner $i$ receives a
bundle in $X_i$ through $a_i:X\to X_i$, use
$v_i:X_i\to\mathbb R$ and the composite $v_i a_i:X\to\mathbb R$.
The resulting \textbf{social welfare} is the map
\[
 \operatorname{SW}:X\longrightarrow\mathbb R,\qquad
 \operatorname{SW}(x)=\sum_{i\in N}v_i(a_i(x)),
\]
with $a_i=\operatorname{id}_X$ in the common-allocation presentation.
A valuation need not be an agent's final utility after payments,
and an allocation's resource feasibility is separately witnessed.
\end{definition}
The general theory uses the supplied comparison relation directly, including when no numerical utility representation is available. When numerical utilities are supplied, their
differences measure incentive gains. The displayed welfare map records the specified sum of valuations,
not an additional rationality law on agents.

\begin{definition}[Best responses and Nash equilibrium]
\label{def:egt-best-response-nash}
Let $z:\Sigma\to Z$ realize the outcome of a policy profile. Against fixed
outsider policies $\sigma_{-i}$, a policy $a_i\in\Sigma_i$ is a
\textbf{best response} if no available $b_i\in\Sigma_i$ satisfies
\[
 z(a_i,\sigma_{-i})\prec_i z(b_i,\sigma_{-i}).
\]
For a numerical utility this is equivalent to
$u_i(z(a_i,\sigma_{-i}))\geq u_i(z(b_i,\sigma_{-i}))$ for every $b_i$.
A \textbf{Nash equilibrium} is a profile $\sigma$ at which each $\sigma_i$
is a best response to $\sigma_{-i}$. With a legal domain, the alternatives
are restricted to those that complete the fixed outsiders to a legal profile.
\end{definition}

Equilibrium expresses stability of a profile under the designated deviations. Coalitional stability tests joint changes
using an explicitly stated improvement criterion. The selection construction
below recovers these predicates and retains all ties.

\begin{example}[Individual and joint incentives]
\label{ex:egt-economic-orientation}
Two agents choose whether to contribute to a project, with actions
$a_1,a_2\in\{0,1\}$. Let
\[
 u_1(a_1,a_2)=2a_1a_2-a_1,\qquad
 u_2(a_1,a_2)=2a_1a_2-a_2.
\]
At $(0,0)$ each receives $0$. Either unilateral contribution gives its
contributor $-1$, so $(0,0)$ is Nash. The joint change to $(1,1)$ gives both
agents $1$, so the same profile admits a strict improvement by the
grand coalition. At $(1,1)$ either unilateral withdrawal reduces the
withdrawing agent's payoff to $0$, so that profile is Nash as well.
This calculation separates individual stability from joint improvement.
The joint improvement is assessed in each agent's own utility; no
collective utility is needed.
\end{example}

\begin{definition}[Legality, observational admissibility, and resource feasibility]
\label{def:egt-choice-restrictions}
\textbf{Legality} of a policy profile is a subobject
$m:L\hookrightarrow\Sigma$, with exactly the elements of $L$ permitted.
For a history set $H$, an effective observation quotient
$o:H\twoheadrightarrow O$, and an action set $A$, a response $f:H\to A$
is \textbf{observationally admissible} precisely when there is a map
$\pi:O\to A$ with $f=\pi o$. Equivalently, for the kernel-pair
projections $p_1,p_2:H\times_OH\rightrightarrows H$, one has
\[
 fp_1=fp_2;
\]
the descended $\pi$ is unique by effective descent in $\mathbf{Set}$.

A \textbf{resource feasibility certificate} is retained in an exhibited
carrier $K$ together with its realization $w:K\to Z$ into the specified
allocation or outcome carrier. Its existential feasibility predicate is
\[
 \operatorname{Feas}_w=\operatorname{im}(w)\hookrightarrow Z.
\]
The image says which outcomes admit a certificate, while $K$ retains the
certificates themselves. Institutional legality, information-dependent
admissibility, and physical feasibility therefore have different carriers
and need not coincide.
\end{definition}
\paragraph{Words with more than one role.}
The following distinctions will be maintained when the constructions meet.
\begin{center}
\begin{tabular}{@{}p{.16\textwidth}p{.77\textwidth}@{}}
\toprule
Term & Meaning in this paper\\
\midrule
Agent & A persistent owner identity with a complete policy interface in an
interaction; its situated record includes observations and comparison feedback.\\
Policy & An ordinary complete contingent plan. Its realized concurrent strategy
is the corresponding arrow of the operational category.\\
State & A monoidal state is an arrow from the tensor unit.
An economic reference state indexes a policy profile; a public state
records a stage of a protocol.\\
Context & A state context indexes references, an operational context
surrounds an interaction by wiring, and a comparison context supplies
the outcome comparisons used for incentives.\\
Type & A dependent type is a family in the semantic model.
An economic type indexes an agent's valuation or utility.
A resource type names a kind of material occurrence.\\
Value & $\mathsf{Val}$ extracts states of a linear object.
An economic valuation or utility instead evaluates desirability.\\
History & A concurrent history is a configuration.
A tree history records a path of decisions or completed public rounds.\\
\bottomrule
\end{tabular}
\end{center}

\begin{definition}[Mixed strategies and probabilistic information]
\label{def:egt-probabilistic-vocabulary}
For a finite set $S$ write
\[
 \operatorname{Dist}(S)=
 \left\{\mu:S\to[0,1]:\sum_{s\in S}\mu(s)=1\right\}.
\]
When $\Sigma_i$ is finite, a \textbf{mixed strategy} of owner $i$
has type $\mu_i\in\operatorname{Dist}(\Sigma_i)$, as opposed
to a pure policy
$\sigma_i\in\Sigma_i$. A finite \textbf{prior} on a type space $T$
is $\mu\in\operatorname{Dist}(T)$. For an observation $o:T\to O$ and
$b\in O$ with $\sum_{o(t)=b}\mu(t)>0$, its \textbf{posterior} has type
$\mu_b\in\operatorname{Dist}(o^{-1}(b))$ and is given by
\[
 \mu_b(t)=\frac{\mu(t)}{\sum_{o(t')=b}\mu(t')}
 \quad(t\in o^{-1}(b)).
\]
For $u:T\to\mathbb R$, the \textbf{expected utility} is
$\mathbb E_\mu[u]=\sum_{t\in T}\mu(t)u(t)$. A finite
\textbf{chance kernel} is a map
$k:W\times\prod_iA_i\to\operatorname{Dist}(W)$.
A \textbf{Bayesian game} additionally indexes agents' information and
utilities by private types and a specified probability distribution;
a \textbf{stochastic game} uses a state-dependent chance kernel for its
transitions. For general measurable carriers the corresponding data are
probability measures and measurable kernels, not the displayed finite sums.
These probabilistic notions describe the boundary of the present pure,
qualitative model; no probability structure is used in the constructions
below.
\end{definition}
The paper's policies are pure, and its belief models record possibilities
without probabilities. The probabilistic structures in
Definition~\ref{def:egt-probabilistic-vocabulary} are vocabulary for the
theory's boundary, not additional data assumed by the constructions.
The carrier presentations in Convention~\ref{conv:egt-base} establish
smallness of the presheaf indexing category. Policy parameters retain their ordinary
set-theoretic meaning.



\begin{notation}[Operational and ordinary structure]
The symbol $A\multimap B$ denotes the internal hom of concurrent
interaction interfaces. Ordinary predicates have implication
$P\Rightarrow Q$. Policy parameters are ordinary objects; resource
occurrences belong to the linear interfaces. Their categorical relationship
is constructed in Section~\ref{sec:egt-dependent-semantics}.
\end{notation}

\subsection{Owned agents and their situated interfaces}
\label{sec:egt-agents}

The policy-parameterized concurrent realization is now available, but
an agent must also retain an identity across all of its controlled
decisions. This section packages those decisions into one complete policy
interface and connects it to the canonical selection relation. Legal
choices and observations are displayed with their state contexts, so that
neither evaluation nor information is mistaken for part of operational
polarity.

\subsubsection{Agent identities, complete policies, and situated interfaces}

Selection enters as a comparison interface, not as a property of a
concurrent event. The canonical relation below accepts a proposed policy
when no available alternative is designated as an improvement. It retains
all ties; the later feedback construction supplies the actual
improvement predicate.

\begin{definition}[Selection at a comparison interface]
\label{def:cgt-selection}
For a set $S$, a \textbf{selection relation} has the type
\[
 \varepsilon_S\hookrightarrow
 S\times\bigl(S\to\mathsf{Pred}(S)\bigr),\qquad
 \varepsilon_S(a,q)\iff\forall b\in S\;\neg q(a)(b).
\]
Here $a$ is a state of $\mathbb L(S)$, and
$q:S\to\mathsf{Pred}(S)$ is a comparison costate: $q(a)$ is the
predicate of improving alternatives to $a$. The selection relation
retains ties and does not choose a representative; for numerical
utilities it expresses maximization among all available choices.

In the terminology of Definition~\ref{def:egt-open-game-vocabulary},
a \textbf{concurrent open game} consists of the typed owned arena
$G=(\vec\Sigma,g):A\to B$ with the canonical tuple
$(\varepsilon_{\Sigma_i^G})_{i\in N}$ of such owner interfaces.
An \textbf{evaluated concurrent open game} is
$\bigl(G,(\varepsilon_{\Sigma_i^G})_i,P\bigr)$ with
$P\in\mathsf{Cmp}_N(A,B)$. Connecting the agent interfaces to the
pulled-back comparison costates constructs
$\mathsf{BR}_G(P)\hookrightarrow(\Sigma^G)^2$ in
Section~\ref{sec:egt-equilibrium}.
A numerical evaluation retains both its value maps and their induced comparison relations. This is the selection-relation interface
of \citet{CapucciEtAl2022Cybernetics} specialized to the paper's
canonical qualitative feedback.
\end{definition}

\begin{definition}[Agent interface and situated agent]
\label{def:egt-agent-interface}
In an owned open game $G=(\vec\Sigma,g):A\to B$, the
\textbf{agent interface} of identity $i\in N$ is
\[
 \mathsf{Ag}_i(G)=(i,\Sigma_i^G,\varepsilon_{\Sigma_i^G}).
\]
The complete policy carrier belongs to this occurrence of the interface;
the identity $i$ is retained when occurrences are composed. Its contextual
realization is
\[
 \alpha_i^G:\Sigma_{-i}^G\times\Sigma_i^G\longrightarrow\mathcal C(A,B),
 \qquad \alpha_i^G(u,a)=g_{(u,a)}.
\]
An \textbf{agent process} is the owned open game connecting these interfaces
through $g$. A \textbf{situated, evaluated agent presentation} consists of
the process together with a legal comprehension $m:L\hookrightarrow\Sigma^G$,
a reference map $s:W\to L$, observation maps $o_i:W\to O_i$, and a comparison
context $P_i\hookrightarrow\mathcal C(A,B)^2$. The unrestricted case uses
$L=\Sigma^G$. The ordinary reference context of the presentation is $ms$.
The corresponding situated agent is the record
\[
 \bigl(\mathsf{Ag}_i(G),\alpha_i^G,o_i,P_i,m,s\bigr).
\]
Its legal interface at $\sigma\in L$ is the dependent carrier
\[
 S_i^L(\sigma_{-i})
 =\{a\in\Sigma_i^G:(\sigma_{-i},a)\in L\},
\]
with its own canonical selection relation and the restriction of
$\alpha_i^G$. The outsider tuple indexes available alternatives; the
observation map specifies the information available to the agent.
\end{definition}

The owned policy parameter and the concurrent arrow $g_\sigma$ have
different types and different roles. The former supplies an alternative;
the latter executes that alternative. Selection, contextual comparison,
and execution are consequently interpreted by different constructions,
which are connected only after the evaluation context has been specified.

The triple $\mathsf{Ag}_i(G)$ records control and the canonical
selection interface, but it does not by itself determine observation,
objective, legality, or implementation. The situated presentation retains
those fields so that knowledge, rationality, ability, and commitment can be
posed as distinct predicates. In particular, failing a rationality or
knowledge test does not remove the agent's control interface.

Over an ordinary parameter context $\Gamma$, the dependent version
uses the same record in each fibre.
Its implementation is the linear map
$\widetilde g:\mathsf{Lin}_\Gamma\Sigma\otimes_D A\to B$;
world-indexed observations and comparisons remain ordinary dependent
terms. This makes the situated interface compatible with substitution and
with the original representable arena fragment.

\begin{proposition}[Recovery of the owned presentation from agent interfaces]
\label{prop:egt-agent-presentation}
Bundling an owned open game and its displayed state, observation, comparison,
and legal data into the preceding agent presentation, and then unpacking it,
recovers the original data. Conversely, unpacking and rebundling recovers the
agent presentation. These constructions preserve ownerwise maps, sequential
and parallel composition of the underlying processes, and substitution of
the ordinary base. Displayed state, observation, and comparison data are
retained under the same context transport.
\end{proposition}
\begin{proof}
\leavevmode
\begin{enumerate}[label=\textup{(\arabic*)},leftmargin=*,itemsep=4pt,topsep=3pt]
  \item \textbf{Secured operational interaction.} The interfaces retain the same indexed carriers $\Sigma_i$, and the process
retains $g$. Its contextual realization is obtained by the canonical
rearrangement $\Sigma_{-i}\times\Sigma_i\cong\Sigma$; evaluating it on the
corresponding coordinates recovers $g$. The selection relation is the
displayed predicate constructed from $\Sigma_i$.
  \item \textbf{Ownerwise alternatives.} All additional data remain
the same maps $m,s,o_i,P_i$. Thus the two constructions are inverse in indexed
family notation. An ownerwise map remains the same family of maps and
satisfies the same realization equation.
  \item \textbf{Policy-level comparison.} Products of its coordinates give
the same composition comparisons, and substitution selects the same fibre. In slice notation the comparison is the canonical equivalence between a
displayed map and the disjoint union of its fibres. No policy representative
is chosen.
\end{enumerate}
\end{proof}

\subsubsection{Displayed ownership and the underlying linear object}
\label{agt-sec-agenttypes}

\begin{definition}[Displayed agent signatures and judgments]
\label{agt-def-displayed-agents}
An agent signature over $\Gamma$ is a tuple
$\mathfrak s=(N,\sig,A,B)$ comprising a finite owner list, ordinary
families $\sig=(\Sigma_i)_{i\in N}$, and linear boundaries $A,B$.  Write
$\Sigma=\prod_i\Sigma_i$.  The signature is part of the judgment:
\[
 \Gamma\vdash\mathfrak s\ \mathsf{AgSig},\qquad
 \Gamma;\Delta\vdash a:\Ag_{\mathfrak s}.
\]
Its underlying linear type, obtained by forgetting the display, is
\[
 \bigl|\Ag_{\mathfrak s}\bigr|=H_{\mathfrak s}
 :=[\Lin\Sigma\otimes A,B].
\]
Formally a displayed agent term is a pair $(\mathfrak s,t)$ with
$\Gamma;|\Delta|\vdash t:H_{\mathfrak s}$.  The notation
$\Ag_{\sig}(A,B)$ suppresses the other entries of $\mathfrak s$.
Unpacking forgets the tag and returns $t$; packing $t$ with a specified
signature returns $(\mathfrak s,t)$.  Packing and unpacking are definitionally inverse, and their equations are those of the underlying linear calculus.  Every linear constructor acts on the underlying terms;
a resulting agent is tagged by an explicitly constructed signature.
\end{definition}

Ordinary substitution transports both an agent signature and its term,
whereas linear substitution transports the term with its signature fixed.
Consequently, conversion of displayed agent types must respect the owner
factors. An isomorphism of underlying hom objects cannot by itself say
whether two coordinates are controlled by one owner or by two: the
signatures $(2\times2)$ and $(2,2)$ have isomorphic total profiles but
different unilateral-deviation kernel pairs. Thus ownership cannot
be recovered from the underlying object $H_{\mathfrak s}$.

An agent variable $a:\Ag_{\mathfrak s}$ occurs linearly in
$\Delta$ and can be packed with its specified signature.
Because an agent term is linear, another process may receive it as an
input without acquiring permission to duplicate its occurrences. The
owner-indexed display has a different purpose: it specifies which
reparametrizations count as strategic maps. An underlying arrow between
agent objects need not preserve those coordinates, and repacking a term
with a different display is not a strategic morphism. The category of
elements in Theorem~\ref{agt-thm-owned-equivalence} characterizes the
realization-preserving ownerwise maps. Regrouping and delegation then give
explicit changes of display, while higher-order linear terms use agents
as resources.

\subsubsection{Introduction, execution, and extensionality}

The derived introduction and execution rules are
\[
 \frac{\Gamma,\sigma:\Sigma;\Delta,x:A\linturn t:B}
      {\Gamma;\Delta\linturn
        \Lambda\sigma.\lambda x.t:\Ag_{\sig}(A,B)},
\]
\[
 \frac{\Gamma;\Delta\linturn a:\Ag_{\sig}(A,B)\quad
       \Gamma\ordturn\sigma:\Sigma\quad
       \Gamma;\Theta\linturn x:A}
      {\Gamma;\Delta,\Theta\linturn a[\sigma](x):B}.
\]
In the first rule, $\Delta,A,B$ are pulled back along the policy parameter.
The introduction first forms the underlying linear term
\[
 \lambda z.\,\mathsf{let}\ u\otimes x=z\ \mathsf{in}\
 \bigl(\mathsf{let}\ \mathsf{return}(\sigma)=u\
       \ \mathsf{in}\ t\bigr),
\]
where $u:\Lin\Sigma$ and the ordinary variable $\sigma$ is bound in
the second elimination.  The specified signature tag is then attached.
Execution unpacks the agent and supplies
$\mathsf{return}(\sigma)\otimes x$ to the internal hom.  Neither
operation changes the ownership tag implicitly.
The equations are
\begin{align}
 (\Lambda\sigma.\lambda x.t)[\tau](u)&=t[\tau/\sigma,u/x],
 \label{agt-eq-agent-beta}\\
 \Lambda\sigma.\lambda x.a[\sigma](x)&=a.
 \label{agt-eq-agent-eta}
\end{align}
The second equation follows from the uniqueness equations for internal
hom, tensor, and $\Lin$ elimination applied to the displayed nested term.
The $\beta$-equation is elimination after introduction and the $\eta$-equation
is the corresponding universal-property uniqueness.  Both are equations of
open terms, so no separation of arrows by global points is invoked.

\begin{theorem}[Agent abstraction is a dependent linear product]
\label{agt-thm-agent-dependent-product}
Let $p:\Gamma.\Sigma\to\Gamma$ be the ordinary comprehension of the
complete policy family. In the indexed linear--nonlinear semantics of
Section~\ref{agt-sec-model}, there is a substitution-natural isomorphism
\begin{equation}
 [\Lin_\Gamma\Sigma\otimes A,B]
 \cong\Pi_p^\ell[p^*A,p^*B].
 \label{agt-eq-agent-dependent-product}
\end{equation}
Its transposition is precisely agent introduction, and its inverse is
policy-indexed application. Hence the displayed agent $\beta\eta$
equations hold on open linear contexts, without using global points to
separate arrows or postulating an independent agent abstraction rule.
\end{theorem}
\begin{proof}
\leavevmode
\begin{enumerate}[label=\textup{(\arabic*)},leftmargin=*,itemsep=4pt,topsep=3pt]
  \item \textbf{Frobenius and coproduct comparison.} For every test object $X$ of $\cL(\Gamma)$, closure,
Proposition~\ref{agt-prop-lin-dependent-sum}, Frobenius, the left indexed
adjunction, fibrewise closure, and the right indexed adjunction give a
natural chain
\begin{align*}
 &\cL(\Gamma)(X,[\Lin\Sigma\otimes A,B])\\
 &\ \cong\cL(\Gamma)(X\otimes\Lin\Sigma\otimes A,B)\\
 &\ \cong\cL(\Gamma)(\Sigma_p^\ell p^*(X\otimes A),B)\\
 &\ \cong\cL(\Gamma.\Sigma)(p^*(X\otimes A),p^*B)\\
 &\ \cong\cL(\Gamma.\Sigma)(p^*X,[p^*A,p^*B])\\
 &\ \cong\cL(\Gamma)(X,\Pi_p^\ell[p^*A,p^*B]).
\end{align*}
Reindexing preserves tensor by construction and internal hom by
Lemma~\ref{agt-lem-closed-reindexing}.
  \item \textbf{Representation and Yoneda decoding.} Yoneda therefore gives the
displayed isomorphism. The first and last transpositions send a
body $\Gamma,\sigma:\Sigma;\Delta,x:A\vdash t:B$ to its agent
abstraction and back to its application. The triangle identities and the
closed-structure $\beta\eta$ equations prove
\eqref{agt-eq-agent-beta}--\eqref{agt-eq-agent-eta} even for open terms.
  \item \textbf{Base-change comparison.} Each comparison is an indexed adjunction comparison, so
Beck--Chevalley makes the isomorphism natural under ordinary substitution.
\end{enumerate}
\end{proof}

\noindent The indexed adjunctions identify agent abstraction with dependent currying.
\begin{center}
\begin{tikzcd}[column sep=large,row sep=large]
\mathsf L(\Gamma)(X,[\Lin\Sigma\otimes A,B])
 \arrow[r,"\cong"]\arrow[d,"\cong"']
 &\mathsf L(\Gamma)(X\otimes\Lin\Sigma\otimes A,B)\arrow[d,"\cong"]\\
 \mathsf L(\Gamma)(X,\Pi_p^\ell[p^*A,p^*B])
 \arrow[r,"\cong"']
 &\mathsf L(\Gamma.\Sigma)(p^*X,[p^*A,p^*B])
\end{tikzcd}
\end{center}

\begin{proposition}[Agents as additive families of interactions]
\label{agt-prop-agent-product}
The following is the copower--closedness comparison specialized to the
canonical model, rather than a property peculiar to event structures.
In the canonical model constructed in Section~\ref{agt-sec-model}, there is
a natural isomorphism
\[
 [\Lin\Sigma\otimes A,B]
 \cong\prod_{\sigma\in\Sigma}[A,B].
\]
The product is an additive product in the linear category.
\end{proposition}
\begin{proof}
\leavevmode
\begin{enumerate}[label=\textup{(\arabic*)},leftmargin=*,itemsep=4pt,topsep=3pt]
  \item \textbf{Fibrewise additive-product representation.} By \cref{agt-thm-agent-dependent-product}, the left-hand object is
$\Pi_p^\ell[p^*A,p^*B]$. At $\gamma$ its fibre is, by the
constructed indexed product formula, precisely
$\prod_{\sigma\in\Sigma_\gamma}[A_\gamma,B_\gamma]$. Naturality and substitution stability are those of the dependent-product
comparison, so the result is a natural isomorphism of linear families.
\end{enumerate}
\end{proof}

\begin{remark}[Agent transformations and strategic maps]
A linear arrow between agent types is a transformation of agent resources.
An ownerwise reparametrization is a more specific piece of displayed data:
it changes the policy parameters in a way that identifies which owner
controls each change. The category of underlying linear types alone does
not recover this distinction. In particular, an isomorphism of total policy
sets does not identify their coalition structures.
\end{remark}

\subsubsection{Open agents}

An open term $\Gamma;\Delta\linturn a:\Ag_{\sig}(A,B)$ denotes, at each
ordinary parameter $\gamma$, an arrow
\[
 \bigotimes\sem\Delta_\gamma
 \longrightarrow[\Lin\Sigma_\gamma\otimes A_\gamma,B_\gamma].
\]
Its uncurried form is
\[
 \bigotimes\sem\Delta_\gamma\otimes\Lin\Sigma_\gamma\otimes A_\gamma
 \longrightarrow B_\gamma.
\]
The uncurried form accounts for the resources needed to supply an
agent separately from the policy and interface supplied when it executes.
Thus higher-order application does not erase the linear inputs on which
an open agent term depends.

\subsubsection{Exact representation of owned presentations}

\begin{theorem}[Representation of agents]
\label{agt-thm-representation}
There is a natural bijection
\[
  \cV\bigl(\I,\bigl|\Ag_{\sig}(yA,yB)\bigr|\bigr)
  \cong\Set\bigl(\Sigma,\cC(A,B)\bigr).
\]
More generally, for every linear test object $T\in\cV$ there is a
natural open-context bijection
\begin{equation}
\label{agt-eq-agent-open-context-representation}
 \cV\bigl(T,[\Lin\Sigma\otimes yA,yB]\bigr)
 \cong\prod_{\sigma\in\Sigma}\cV(T\otimes yA,yB).
\end{equation}
These comparisons identify execution, sequential and parallel composition,
ownerwise reparametrization, regrouping, and delegation with the
corresponding constructions on owned concurrent realizations.
\end{theorem}
\begin{proof}
\leavevmode
\begin{enumerate}[label=\textup{(\arabic*)},leftmargin=*,itemsep=4pt,topsep=3pt]
  \item \textbf{Frobenius and coproduct comparison.} Closedness and the coproduct description of $\Lin$ give, naturally in
$T$ and in the policy set $\Sigma$,
\begin{align*}
 \cV(T,[\Lin\Sigma\otimes yA,yB])
 &\cong\cV(T\otimes\Lin\Sigma\otimes yA,yB)\\
 &\cong\cV\Bigl(\coprod_{\sigma\in\Sigma}(T\otimes yA),yB\Bigr)\\
 &\cong\prod_{\sigma\in\Sigma}\cV(T\otimes yA,yB).
\end{align*}
The middle comparison is an isomorphism because Day tensor preserves
coproducts separately; the last is the coproduct universal property.
  \item \textbf{Representation and Yoneda decoding.} For $T=\I$, restricted Yoneda gives
$\cV(yA,yB)\cong\cC(A,B)$ and hence the first assertion. For $r:\Theta\to\Sigma$, precomposition by $\Lin r\otimes1$
sends the component at $\theta$ to the component at $r(\theta)$. Sequential and parallel composition send families $g_\sigma,h_\tau$
to $(h_\tau g_\sigma)_{\sigma,\tau}$ and
$(g_\sigma\otimes h_\tau)_{\sigma,\tau}$, respectively.
  \item \textbf{Ownerwise alternatives.} These calculations establish equality of the maps represented by the open-context bijection for arbitrary linear test objects. Ownerwise regrouping and delegation apply the same named-coordinate
substitutions on both sides; the comparison maps obey coherence by
ordinary-product and symmetric-monoidal coherence.
  \item \textbf{Policy-level comparison.} The equality
$g=h\prod_i r_i$ defining a realization-preserving $2$-cell becomes
exactly the equality between the corresponding reparametrized
agent states.
\end{enumerate}
\end{proof}

\noindent Policy indexing, currying, and Yoneda give an exact representation of realized agents.
\begin{center}
\begin{tikzcd}[column sep=large,row sep=large]
\mathcal V(\I,[\Lin\Sigma\otimes yA,yB])\arrow[r,"\cong"]\arrow[d,"\cong"']
 &\mathcal V(\Lin\Sigma\otimes yA,yB)\arrow[d,"\cong"]\\
 \Set(\Sigma,\mathcal C(A,B))\arrow[r,"\cong"']
 &\prod_{\sigma\in\Sigma}\mathcal V(yA,yB)
\end{tikzcd}
\end{center}

\begin{theorem}[Equivalence on owned presentations]
\label{agt-thm-owned-equivalence}
For fixed original arenas $A,B$ and owner list $N$, form the
presheaf of agent states on $\Set^N$,
\[
 F_{A,B}(\sig)=
 \cV\bigl(\I,[\Lin(\prod_i\Sigma_i)\otimes yA,yB]\bigr),
 \qquad F_{A,B}(r)(b)=b[r].
\]
Let $\mathsf{StateAg}_N(A,B)=\int F_{A,B}$ be its category of
elements. Its objects are $(\sig,a)$ with the retained owner display;
an arrow $(\sig,a)\to(\taus,b)$ is an ownerwise map
$r_i:\Sigma_i\to T_i$ satisfying $a=b[r]$.
Then execution and restricted Yoneda define an isomorphism of categories
\[
 \mathsf{StateAg}_N(A,B)\ \cong\
 \mathbf{Arena}_N(\cC)(A,B).
\]
These isomorphisms are compatible with sequential and parallel composition,
identity presentations, and their ownerwise associativity and symmetry
comparisons.  They also commute with ordinary base substitution.
\end{theorem}
\begin{proof}
\leavevmode
\begin{enumerate}[label=\textup{(\arabic*)},leftmargin=*,itemsep=4pt,topsep=3pt]
  \item \textbf{Representation and Yoneda decoding.} For $g:\Sigma\to\cC(A,B)$, copair the arrows $yg_\sigma:yA\to yB$:
\[
 \Lin\Sigma\otimes yA\cong\coprod_{\sigma\in\Sigma}yA
 \xrightarrow{[yg_\sigma]_\sigma}yB.
\]
Currying and tagging this map constructs the agent state $a_g$.
  \item \textbf{Reverse implication.} Conversely, evaluate an agent state $a$ on every coproduct injection,
then decode the resulting map $yA\to yB$ by Yoneda; call the family
$g_a$. The coproduct and closed-structure $\beta\eta$ equations give
$g_{a_g}=g$ and $a_{g_a}=a$, without relying on global-point
extensionality.
  \item \textbf{Ownerwise alternatives.} For an ownerwise map $r$, the equation $a=b[r]$
holds exactly when $g_a(\sigma)=g_b(r(\sigma))$ for every profile,
which is precisely the original presentation $2$-cell condition and the
category-of-elements arrow condition. In particular arbitrary arrows of
the underlying linear category are not thereby declared strategic maps. The execution equations~\eqref{agt-eq-seq}--\eqref{agt-eq-par} then
identify horizontal and tensor composition.
  \item \textbf{Ownerwise alternatives: completion.} Each higher coherence map
has the same ownerwise coordinate projections and the same concurrent
realizations on both sides. The ordinary and symmetric-monoidal coherence
laws prove compatibility with all comparison maps. Pointwise execution
also commutes with ordinary substitution.
\end{enumerate}
\end{proof}

\noindent Copairing and componentwise execution are inverse correspondences on owned presentations.
\begin{center}
\begin{tikzcd}[column sep=large,row sep=large]
(\vec\Sigma,g)\arrow[r,"g\mapsto a_g"]\arrow[dr,"\id"']
 &(\vec\Sigma,a_g)\arrow[d,"a\mapsto g_a"]\\
 &(\vec\Sigma,g)
\end{tikzcd}
\end{center}

\begin{proposition}[Interpretation of open terms]
\label{agt-prop-core-sound}
Every core derivation determines a morphism
$\bigotimes\sem\Delta\to\sem A$ in $\cL(\sem\Gamma)$.
Definitional equality gives equality of morphisms. Ordinary and linear
substitution are interpreted by reindexing and composition.
\end{proposition}
\begin{proof}
\leavevmode
\begin{enumerate}[label=\textup{(\arabic*)},leftmargin=*,itemsep=4pt,topsep=3pt]
  \item \textbf{Frobenius and coproduct comparison.} Interpret context splitting by tensor, abstraction and application by
closedness, additives by products and coproducts, the ordinary--linear
rules by $\Lin\dashv\Val$ together with
Proposition~\ref{agt-prop-lin-dependent-sum}, and indexed binders by
\eqref{agt-eq-indexed-adjoints}. Theorem~\ref{agt-thm-indexed} makes ordinary
substitution commute with these operations.
  \item \textbf{Fixed-point argument.} The induction of
Lemma~\ref{agt-lem-substitution} interprets linear substitution as
composition. Each $\beta$ equation is an evaluation equation of one of
these universal properties; each $\eta$ equation is its uniqueness
equation. Commuting conversions express naturality. Agent constructors
are the specified composites of these operations.
\end{enumerate}
\end{proof}

\begin{remark}[Linearization and physical data]
The functor $\Lin$ organizes ordinary alternatives in the completion.
The physical finite-data functor is $D:\mathbf{FinSet}\to\cC$.
Its collector $D(X)\otimes D(Y)\to D(X\times Y)$ waits for both inputs
and belongs to a lax monoidal structure; it is generally not an isomorphism.
We retain $D$ for communication implementations and $\Lin$ for the
ordinary--linear adjunction. Identifying them would erase a real
synchronization barrier.
\end{remark}

\part{Composition, control, and game-theoretic evaluation}

\section{The Compositional Calculus of Agents}\label{sec:paper-3}

Having represented agents by complete-policy families, we must show
that composition respects both their policies and their concurrent
executions. The two operations differ: ordinary products combine an
owner's contingent plans, while composition and tensor combine linear
interactions, as in parameterized open games
\citep{GhaniEtAl2018,AtkeyEtAl2021}. Their compatibility yields a
symmetric monoidal agent calculus and a bicategory of owned presentations.
The ownerwise factorization must survive this construction; forgetting it
would change the allowed unilateral deviations.

\subsection{Linear agent operations and ownerwise composition}
\label{agt-sec-composition}

\subsubsection{The two composition operations}

For signatures on a common owner set, define
$(\sig\times\taus)_i=\Sigma_i\times T_i$. The canonical rearrangement
identifies its total profile with $\Sigma\times T$. Unused owners have
singleton policy carriers; disjoint owner lists can therefore be combined
without identifying owners.

\begin{construction}[Agent composition]
The following are linear operations:
\begin{align*}
 \mathsf{seq}:&\ \Ag_{\sig}(A,B)\otimes\Ag_{\taus}(B,C)
       \longrightarrow\Ag_{\sig\times\taus}(A,C),\\
 \mathsf{par}:&\ \Ag_{\sig}(A,B)\otimes\Ag_{\taus}(A',B')
       \longrightarrow\Ag_{\sig\times\taus}(A\otimes A',B\otimes B').
\end{align*}
Their defining executions are
\begin{align}
 (a\seq b)[\sigma,\tau](x)&=b[\tau](a[\sigma](x)),
 \label{agt-eq-seq}\\
 (a\parcomp b)[\sigma,\tau](x\otimes y)
 &=a[\sigma](x)\otimes b[\tau](y).
 \label{agt-eq-par}
\end{align}
The identity agent has singleton policies and executes $x\mapsto x$.
\end{construction}

The first operation consumes each supplied agent once. The second splits
both the agent resources and their inputs. These are comparison maps between
types; they are not asserted to be isomorphisms. An arbitrary composite
agent need not admit a decomposition into two independently supplied agents.

\begin{theorem}[Coherence of agent composition]
\label{agt-thm-agent-coherence}
Sequential composition is associative and unital under the canonical
ownerwise product bijections. Parallel composition is coherently symmetric
monoidal. Sequential and parallel composition satisfy interchange:
\[
 (a\parcomp b)\seq(c\parcomp d)
 \cong(a\seq c)\parcomp(b\seq d),
\]
where the comparison regroups the complete policy coordinates.
\end{theorem}
\begin{proof}
\leavevmode
\begin{enumerate}[label=\textup{(\arabic*)},leftmargin=*,itemsep=4pt,topsep=3pt]
  \item \textbf{Associativity.} For associativity, agent abstraction followed by
\eqref{agt-eq-agent-beta} reduces both sides, at the same named
policy coordinates, to
\[
 c[\upsilon](b[\tau](a[\sigma](x))).
\]
The unit calculation reduces to $a[\sigma](x)$. The tensor comparison is
obtained by reassociating and exchanging the linear inputs and the ordinary
profile factors.
  \item \textbf{Interchange and higher coherence.} For interchange, both executions reduce to
\[
 c[\upsilon](a[\sigma](x))\otimes
 d[\omega](b[\tau](y)).
\]
Agent $\eta$ then gives the open-term equations. In every higher coherence diagram, the ordinary maps have the same
coordinate projections and the linear maps are the coherence maps of the
symmetric monoidal calculus. The product and monoidal coherence laws
therefore give the required commuting diagrams.
\end{enumerate}
\end{proof}

\subsubsection{Ownerwise policy substitution}

For ordinary maps $r_i:\Theta_i\to\Sigma_i$, write $r=\prod_i r_i$ and set
\[
 a[r][\theta](x)=a[r(\theta)](x).
\]
This defines a linear map
$\Ag_{\sig}(A,B)\to\Ag_{\ths}(A,B)$, implemented by precomposition with
$\Lin r\otimes1_A$. Computation proves
\[
 a[\id]=a,\qquad a[r][s]=a[r\circ s],
\]
and
\[
 (a\seq b)[r\times q]=a[r]\seq b[q],\qquad
 (a\parcomp b)[r\times q]=a[r]\parcomp b[q].
\]
Here $r\times q$ is formed ownerwise and then flattened.

At the level of declared presentations, an ownerwise map $r:a\Rightarrow b$
has components $r_i:\Sigma_i^a\to\Sigma_i^b$ satisfying
$a[\sigma]=b[r(\sigma)]$. Vertical composition is function composition;
horizontal and parallel composition take ownerwise products. Together with
the preceding coherence maps these presentations form the owned
symmetric monoidal bicategory recovered in Section~\ref{agt-sec-model}.

\begin{example}[Shared policy and duplicated execution]
Two occurrences controlled by one owner initially have policy carrier
$\Sigma\times\Sigma$. Precomposition with the ordinary diagonal
$\Sigma\to\Sigma\times\Sigma$ makes them use the same policy description.
Both linear occurrences must still be supplied. The diagonal ties parameters;
it does not construct a linear diagonal on the agent type.
\end{example}

\subsection{Open arenas, ownership and contextual wiring}
\label{sec:egt-players}

The linear agent operations were defined through internal homs.
To identify them with ordinary owned presentations, we now carry out their
composition on the original concurrent arenas. Full faithfulness of
Yoneda ensures that each component still denotes one genuine concurrent
strategy. State contexts will select reference profiles, and joint
comparison feedback will later test controlled alternatives without
breaking the composite into independently evaluated fragments.

The backward interface will be inverse image of predicates along the
complete wiring map. This construction retains comparisons between joint
alternatives until the owner's policy coordinates determine which changes
are allowed. It adapts the open-game and contextual approaches of
\citet{GhaniEtAl2018,AtkeyEtAl2021,CapucciEtAl2022Agency,Hedges2020}
to the concurrent category.

\subsubsection{The bicategory of open arenas}

\begin{definition}[Open arenas and open games]
\label{def:egt-open-game-vocabulary}
An \textbf{owned open arena} on a linear interaction boundary $(A,B)$ is
the data
\[
 G=(\vec\Sigma,g):A\to B,\qquad
 \vec\Sigma\in\mathbf{Par}_N,\qquad
 g:\Sigma^G=\prod_{i\in N}\Sigma_i^G\longrightarrow\mathcal C(A,B).
\]
Its owned policy families parameterize the concurrent realizations; at
this stage it carries no preference. A \textbf{concurrent open game}
adds the prescribed owner-selection interfaces
\[
 \varepsilon_{\Sigma_i^G}\hookrightarrow
 \Sigma_i^G\times
 \bigl(\Sigma_i^G\to\mathsf{Pred}(\Sigma_i^G)\bigr),
 \qquad
 \varepsilon_S(a,q)\iff\forall b\in S\;\neg q(a)(b),
\]
whose construction is developed in
Definition~\ref{def:cgt-selection}. An \textbf{evaluated concurrent
open game} has in addition a comparison context $P=(P_i)_{i\in N}$ of
type
\[
 P\in\mathsf{Cmp}_N(A,B)
 =\prod_{i\in N}\mathsf{Pred}\bigl(\mathcal C(A,B)^2\bigr).
\]
Connecting this context to the owner interfaces determines the
best-response relation $\mathsf{BR}_G(P)\hookrightarrow(\Sigma^G)^2$.
The outcome interface can retain the entire exposed concurrent interaction, with scalar evaluations and material allocations as particular observations.
\end{definition}
\begin{definition}[Reparametrisations of owned open arenas]
\label{def:egt-concurrent-open-arena-and-reparametrisation}
For owned open arenas $G=(\vec\Sigma^G,g)$ and
$G'=(\vec\Sigma^{G'},g')$ on the same linear interaction boundary
$(A,B)$, a \textbf{reparametrisation} or $2$-cell
$f:G\Rightarrow G'$ is an ownerwise map
$f_i:\Sigma_i^G\to\Sigma_i^{G'}$ satisfying $g=g'\Pi_Nf$.
An element of $\Sigma_i^G$ is a complete policy of owner $i$.
A presentation is \textbf{owned by $i$} if the other policy sets are
singletons, and \textbf{nonstrategic} if every policy set is a singleton.
The backward interface of an owned open arena is the canonical predicate
feedback constructed from its realization below, not additional
unspecified data.
\end{definition}

\begin{construction}[Horizontal and parallel composition]
For $G:A\to B$ and $H:B\to C$, put
\[
 \Sigma_i^{H\circ G}=\Sigma_i^G\times\Sigma_i^H,
 \qquad (h\circ g)_{(\sigma,\tau)}=h_\tau\circ g_\sigma.
\]
For $G:A\to B$ and $K:A'\to B'$, put
\[
 \Sigma_i^{G\otimes K}=\Sigma_i^G\times\Sigma_i^K,
 \qquad(g\otimes k)_{(\sigma,\rho)}=g_\sigma\otimes k_\rho.
\]
Use the canonical product rearrangements throughout. Both operations on
$2$-cells are ownerwise products. The identity at $A$ has singleton policy
sets and realizes $\operatorname{id}_A$.
\end{construction}

\begin{theorem}[Bicategory of concurrent open arenas]
\label{thm:egt-owned}
These operations define a symmetric monoidal bicategory
$\mathbf{Arena}_N(\mathcal C)$. Its hom-category at $(A,B)$ is the category of
elements of the presheaf
\[
 F_{A,B}:\mathbf{Par}_N^{\mathrm{op}}\to\mathbf{Set},
 \qquad F_{A,B}(\vec\Sigma)
 =\mathbf{Set}(\Pi_N\vec\Sigma,\mathcal C(A,B)).
\]
The quotient by invertible $2$-cells is a symmetric monoidal category
$\mathsf{ArenaIso}_N(\mathcal C)$.
\end{theorem}
\begin{proof}
\leavevmode
\begin{enumerate}[label=\textup{(\arabic*)},leftmargin=*,itemsep=4pt,topsep=3pt]
  \item \textbf{Functoriality and substitution.} The category-of-elements description is exactly the equation defining a
$2$-cell. The comparison
$\Pi_N(\vec\Sigma\times\vec\Theta)\cong
\Pi_N\vec\Sigma\times\Pi_N\vec\Theta$ is natural. It therefore makes the
displayed horizontal operations functorial on hom-categories.
  \item \textbf{Ownerwise alternatives.} The associator and unitors are the canonical ownerwise product bijections;
their realization equations are the associativity and unit equations in
$\mathcal C$. Their pentagon and triangle hold coordinatewise.
  \item \textbf{Ownerwise alternatives: second comparison.} The tensor compositor is the ownerwise permutation between the parameter
orders of $(H\otimes H')\circ(G\otimes G')$ and
$(H\circ G)\otimes(H'\circ G')$. Its realization equation is interchange
in $\mathcal C$. The monoidal associator, unitors, and symmetry have
singleton parameters and realize those of $\mathcal C$; their naturality
$2$-cells are the corresponding product rearrangements.
  \item \textbf{Monoidal comparison.} In each coherence
diagram both parameter maps send every named coordinate to the same
coordinate. The realized arrows agree by symmetric monoidal coherence in
$\mathcal C$ \citep{Selinger2011}. Since a $2$-cell is determined by its
parameter maps, this verifies the coherence diagrams themselves.
  \item \textbf{Passage to isomorphism classes.} Finally,
invertible $2$-cells are exactly ownerwise bijections; the coherence maps
become equalities on their isomorphism classes.
\end{enumerate}
\end{proof}

\noindent Ownerwise reparametrization commutes with evaluation of the composite realization.
\begin{center}
\begin{tikzcd}[column sep=large,row sep=large]
\Sigma^G\times\Sigma^H\arrow[r,"r\times s"]\arrow[d,"h\circ g"']
 &\Sigma^{G'}\times\Sigma^{H'}\arrow[d,"h'\circ g'"]\\
 \mathcal C(A,C)\arrow[r,equal]&\mathcal C(A,C)
\end{tikzcd}
\end{center}

\[
\begin{tikzcd}[column sep=large]
 \Sigma^G\times\Sigma^H \arrow[r,"g\times h"] \arrow[d,"\cong"'] &
 \mathcal C(A,B)\times\mathcal C(B,C) \arrow[d,"\circ"]\\
 \Sigma^{H\circ G} \arrow[r,"h\circ g"'] &\mathcal C(A,C)
\end{tikzcd}
\]

\begin{construction}[Reindexing and the nonstrategic embedding]
\label{con:egt-owned-reindexing}
For an ownerwise map $f:\vec\Theta\to\vec\Sigma$ define
\begin{equation}\label{eq:egt-owned-reindexing}
 f^*G=(\vec\Theta,g\Pi_Nf).
\end{equation}
It carries the canonical $2$-cell $f:f^*G\Rightarrow G$.
Define
\begin{equation}\label{eq:egt-nonstrategic-embedding}
 \mathbf j_N:\mathcal C\to\mathbf{Arena}_N(\mathcal C),
 \qquad u\longmapsto(\vec1,\operatorname{const}_u).
\end{equation}
Its structure maps are the ownerwise bijections $1\times1\cong1$.
\end{construction}

\begin{proposition}[Reindexing and nonstrategic recovery]
Reindexing is functorial and gives the cartesian arrows of each
category-of-elements projection to $\mathbf{Par}_N$. The functor
$\mathbf j_N$ is locally fully faithful and its image is biequivalent to
the nonstrategic sub-bicategory.
\end{proposition}
\begin{proof}
\leavevmode
\begin{enumerate}[label=\textup{(\arabic*)},leftmargin=*,itemsep=4pt,topsep=3pt]
  \item \textbf{Ownerwise alternatives.} The cartesian universal property is substitution of the equation
$g=g'\Pi_Nf$. Between two singleton-parameter presentations there is one
possible ownerwise map, and it is a $2$-cell exactly when their realized
arrows are equal. Every family of singleton policy sets is uniquely
isomorphic to $\vec1$.
\end{enumerate}
\end{proof}

\begin{remark}[Reference sharing and policy restriction]
Reindexing a family of reference states to a diagonal does not change its
space of alternative policies. To enforce reuse of a policy in two
components, use the ownerwise diagonal
$\Sigma_i\to\Sigma_i\times\Sigma_i$ in
Construction~\ref{con:egt-owned-reindexing}, or construct the corresponding
legal subobject. These are changes to the available policy presentation;
ordinary substitution of the reference state alone is not such a change.
\end{remark}

\subsubsection{Joint feedback for concurrent wiring}

Wiring joins the concurrent processes and combines the policy
coordinates belonging to each owner. Feedback must therefore be pulled
back along the whole wiring function, not computed independently at each
port. The resulting predicate remembers interactions between component
choices.

\begin{construction}[Multi-input coplay]
\label{con:cgt-joint-feedback}
Let $H_j=\mathcal C(A_j,B_j)$, and let
$m:\prod_{j=1}^rH_j\to H$ be a wiring function constructed using
composition, tensor, fixed concurrent arrows, and their coherent
reassociations. Apply $\mathbb L$ to $m$.
For a comparison $P\hookrightarrow H^2$ its transported costate is
the comparison
\begin{equation}\label{eq:cgt-joint-feedback}
 m^\sharp P((f_j)_j,(f'_j)_j)
 \quad\Longleftrightarrow\quad
 P\bigl(m((f_j)_j),m((f'_j)_j)\bigr).
\end{equation}
The input continuation interface is
$\mathbb L(\prod_jH_j)$: feedback is a predicate on joint alternatives.
For component arenas $G_j$ with plays $g_j$, their wired play is
$m\circ\prod_jg_j$, followed by the ownerwise parameter rearrangement.
Its coplay is the canonical coplay of that composite function.
\end{construction}

\begin{theorem}[Canonical coplay and open-arena coherence]
\label{thm:cgt-wiring-soundness}
Wiring substitution commutes with backward feedback. In particular,
\[
 (mn)^\sharp=n^\sharp m^\sharp,
 \qquad
 (m\circ\textstyle\prod_jg_j)^\sharp
   =(\textstyle\prod_jg_j)^\sharp m^\sharp.
\]
The sequential, parallel, unit, symmetry, and reparametrisation laws
of open arenas hold for both play and coplay. Every arena retains its
owner-indexed policy sets and its complete concurrent play map.
\end{theorem}
\begin{proof}
\leavevmode
\begin{enumerate}[label=\textup{(\arabic*)},leftmargin=*,itemsep=4pt,topsep=3pt]
  \item \textbf{Equality of the two constructions.} The displayed equations are
$((mn)\times(mn))^*=(n\times n)^*(m\times m)^*$. For sequential composition the two arguments of a transported comparison
are $(h_\tau g_\sigma,h_{\tau'}g_{\sigma'})$. For a triple composite the two bracketings give
\[
 P(k_\upsilon(h_\tau g_\sigma),
     k_{\upsilon'}(h_{\tau'}g_{\sigma'}))
 =P((k_\upsilon h_\tau)g_\sigma,
     (k_{\upsilon'}h_{\tau'})g_{\sigma'}).
\]
Associativity in $\mathcal C$ proves equality.
  \item \textbf{Ownerwise alternatives.} Interchange, units, and
symmetry are handled by their equations in $\mathcal C$ on both arguments. The ownerwise permutations are exactly those of
Theorem~\ref{thm:egt-owned}. For a reparametrisation $f$, the equation
$g=g'\Pi_Nf$ yields
\[
 (g\times g)^*P
  =(\Pi_Nf\times\Pi_Nf)^*(g'\times g')^*P.
\]
Thus all coherence diagrams commute also on costates.
  \item \textbf{Recovery of canonical coplay.} Finally, forgetting
the canonical backward map recovers the same $g$ and the same policy
coordinates, while Theorem~\ref{thm:cgt-coplay} reconstructs that backward
map uniquely by the displayed construction.
\end{enumerate}
\end{proof}

\noindent Ownerwise reparametrization commutes with joint wiring.
\begin{center}
\begin{tikzcd}[column sep=large,row sep=large]
\prod_j\Sigma_j\arrow[r,"\prod_jr_j"]\arrow[d,"m\circ\prod_jg_j"']
 &\prod_j\Theta_j\arrow[d,"m\circ\prod_jh_j"]\\
 H\arrow[r,equal]&H
\end{tikzcd}
\end{center}

\[
\begin{tikzcd}[column sep=large,row sep=large]
 (\Sigma^G\times\Sigma^H)^2
   \arrow[r,"{(g\times h)^2}"] \arrow[d,"{(h\circ g)^2}"'] &
 (H_{A,B}\times H_{B,C})^2 \arrow[d,"m^2"]\\
 H_{A,C}^2 \arrow[r,equal] & H_{A,C}^2
\end{tikzcd}
\]

The owner grading records precisely the policy factors preserved by
composition. When an owner acts in more than one component, its composite
coordinate is the product of those decisions. The graded-arrow structure
below packages parameter-indexed hom-families with their unit, composition,
strength, and reindexing; its coherence says that these operations commute
with substitution and concurrent tensor.

\begin{proposition}[The owner grading]
\label{prop:cgt-grading}
The family
\[
 \mathcal A_{\vec\Sigma}(A,B)
 =\mathbf{Set}(\Pi_N\vec\Sigma,\mathcal C(A,B))
\]
is a commutative $\mathbf{Par}_N^{\mathrm{op}}$-graded arrow over
$\mathcal C$: its unit is the singleton family, composition uses
ownerwise products, reparametrisation is precomposition, and strength
sends $g_\sigma$ to $g_\sigma\otimes\operatorname{id}_Z$.
\end{proposition}
\begin{proof}
\leavevmode
\begin{enumerate}[label=\textup{(\arabic*)},leftmargin=*,itemsep=4pt,topsep=3pt]
  \item \textbf{Naturality.} Naturality in the grade is substitution. Graded associativity and the
unit equations are the corresponding equations of $\mathcal C$ after
the canonical product comparisons of $\Pi_N$. The strength laws are
functoriality and associativity of tensor.
  \item \textbf{Ownerwise alternatives.} The two ways of forming
parallel composition agree by interchange. These are precisely the
graded-arrow equations of \citet{AtkeyEtAl2021}, verified here for
the constructed owner grading. Taking its category of elements retains
the reparametrisations and yields the bicategory already constructed.
\end{enumerate}
\end{proof}

\noindent Owner grading respects the two ways of inserting a fixed interaction.
\begin{center}
\begin{tikzcd}[column sep=large,row sep=large]
\mathcal A_{\vec\Sigma}(A,B)\arrow[r,"-\otimes\id_Z"]\arrow[d,"r^*"']
 &\mathcal A_{\vec\Sigma}(A\otimes Z,B\otimes Z)\arrow[d,"r^*"]\\
 \mathcal A_{\vec\Theta}(A,B)\arrow[r,"-\otimes\id_Z"']
 &\mathcal A_{\vec\Theta}(A\otimes Z,B\otimes Z)
\end{tikzcd}
\end{center}

The grading extends the agency and contextual-wiring constructions of
\citet{CapucciEtAl2022Agency,CapucciEtAl2022Cybernetics,Hedges2020}.
It matters because feedback must evaluate a repeated owner's 
\emph{complete} policy. For an owner acting in $G$ and $H$, the forward
carrier is $\Sigma_i^G\times\Sigma_i^H$ and the backward carrier is
$\mathsf{Pred}(\Sigma_i^G\times\Sigma_i^H)$. The next step fixes the
other owners' coordinates to obtain the individual costate, as in
Section~\ref{sec:egt-equilibrium}.

\subsubsection{Composition of contexts and their logical interfaces}

\begin{construction}[Sequential and parallel context composition]
\label{con:fib-context-composition}
For $G:A\to B$, $H:B\to C$, use the canonical owner rearrangement
$\Sigma^{H\circ G}\cong\Sigma^G\times\Sigma^H$ to define
\[
 \mathsf{St}(G)\times\mathsf{St}(H)\longrightarrow
 \mathsf{St}(H\circ G),\qquad
 (W\xrightarrow{s}\Sigma^G,V\xrightarrow{t}\Sigma^H)
 \longmapsto(W\times V\xrightarrow{s\times t}\Sigma^{H\circ G}).
\]
Its realization is $(w,v)\mapsto h_{t(v)}g_{s(w)}$.
For parallel arenas use the same product context and realization
$(w,v)\mapsto g_{s(w)}\otimes k_{t(v)}$.
Predicates have an external conjunction
$P\boxtimes Q=\pi_W^*P\wedge\pi_V^*Q$, while the complete fibre
is $\mathsf{Pred}(W\times V)$, including arbitrary joint predicates.
\end{construction}

\begin{theorem}[Fibred coherence of concurrent wiring]
\label{thm:fib-wiring}
The context-composition functors preserve cartesian lifts jointly in
their inputs. External conjunction commutes with these lifts. Their
associativity, units, symmetry, and interchange comparisons satisfy
the corresponding bicategorical coherence laws. For any concurrent
wiring $m:\prod_jH_j\to H$ and realizations $a_j:W_j\to H_j$,
\[
 (m\circ\textstyle\prod_ja_j)^*Q
   =(\textstyle\prod_ja_j)^*m^*Q,
 \qquad
 (m\circ\textstyle\prod_ja_j)^\sharp P
   =(\textstyle\prod_ja_j)^\sharp m^\sharp P.
\]
\end{theorem}
\begin{proof}
\leavevmode
\begin{enumerate}[label=\textup{(\arabic*)},leftmargin=*,itemsep=4pt,topsep=3pt]
  \item \textbf{Monoidal comparison.} For maps to $X$ and $Y$ there is a natural bijection
\[
 (S\times_XW)\times(T\times_YV)
 \cong (S\times T)\times_{X\times Y}(W\times V).
\]
Both sides consist of tuples $(s,t,w,v)$ satisfying the same two
equations. It identifies the product of the cartesian lifts in
Theorem~\ref{thm:fib-tower} with the cartesian lift of their composite. Inverse image commutes with conjunction and projections, proving the
predicate assertion.
  \item \textbf{Coherence.} The coherence maps rearrange the same named
coordinates. Their realizations agree by associativity and interchange
in $\mathcal C$, so the coherence diagrams commute in the tower itself. The final equations are functoriality of inverse image, applied once to
the realization and once to its square.
\end{enumerate}
\end{proof}

\noindent Cartesian lifts are compatible with products of independently wired state contexts.
\begin{center}
\begin{tikzcd}[column sep=large,row sep=large]
(S\times_XW)\times(T\times_YV)\arrow[r,"\cong"]\arrow[d]
 &(S\times T)\times_{X\times Y}(W\times V)\arrow[d]\\
 S\times T\arrow[r,equal]&S\times T
\end{tikzcd}
\end{center}

\begin{remark}[Whole joint feedback]
The source interface for a product of carriers is
$\mathbb L(X\times Y)=(X\times Y,\mathsf{Pred}(X\times Y))$.
The ordinary lens tensor instead has backward carrier
$\mathsf{Pred}(X)\times\mathsf{Pred}(Y)$; it is not identified with the
former. In particular $\mathbb L$ is faithful but not strong monoidal
for the cartesian tensor on sets and the Cartesian-product tensor on lenses.
The constructed context fibre keeps whole joint feedback until ownership
determines which alternatives are varied. This implements the distinction
between joint and partially evaluated feedback described by
\citet{Capucci2023}. The graded-arrow organization remains the one
proved in Proposition~\ref{prop:cgt-grading}; fibrational lenses and
related variants are developed in \citet[Section 6]{AtkeyEtAl2021}.
\end{remark}

\subsubsection{Composition of complete agent interfaces}

\begin{proposition}[The complete-policy composition law]
\label{prop:egt-agent-composition}
For composable agent processes $G:A\to B$ and $H:B\to C$, and for
$K:A'\to B'$, the agent of identity $i$ has interfaces
\[
 \begin{aligned}
 \mathsf{Ag}_i(H\circ G)
   &=\bigl(i,\Sigma_i^G\times\Sigma_i^H,
                   \varepsilon_{\Sigma_i^G\times\Sigma_i^H}\bigr),\\
 \mathsf{Ag}_i(G\otimes K)
   &=\bigl(i,\Sigma_i^G\times\Sigma_i^K,
                   \varepsilon_{\Sigma_i^G\times\Sigma_i^K}\bigr).
 \end{aligned}
\]
The realizations are $h_\tau g_\sigma$ and $g_\sigma\otimes k_\rho$.
The identity process has singleton agent policies. These operations and
ownerwise realization-preserving reparametrizations form the agent calculus
of $\mathbf{Arena}_N(\mathcal C)$, with its symmetric monoidal coherence.
They are stable under substitution in the dependent model.
\end{proposition}
\begin{proof}
\leavevmode
\begin{enumerate}[label=\textup{(\arabic*)},leftmargin=*,itemsep=4pt,topsep=3pt]
  \item \textbf{Ownerwise alternatives.} The parameter family of a composite is the ownerwise product by
Proposition~\ref{prop:egt-dependent-game-composition}. Applying
Definition~\ref{def:egt-agent-interface} gives the displayed interfaces. Three sequential components have the same flattened policy tuple
$(\sigma_i,\tau_i,\upsilon_i)$ for each $i$, and both bracketings realize
$k_\upsilon h_\tau g_\sigma$.
  \item \textbf{Ownerwise alternatives: second comparison.} Unit comparisons remove singleton coordinates;
tensor comparisons reassociate and permute named coordinates. Interchange
holds because it holds in $\mathcal C$. Products of ownerwise maps preserve
these equations.
  \item \textbf{Dependent substitution of the composite.} Thus every coherence diagram is the corresponding product
coordinate diagram together with the same realized concurrent arrow. The dependent construction evaluates these formulas at each base index, so
substitution commutes with all comparisons.
\end{enumerate}
\end{proof}

\noindent The composite realization is obtained by joining the two complete-policy profiles.
\begin{center}
\begin{tikzcd}[column sep=large,row sep=large]
\Sigma^{H\circ G}\arrow[r,"\cong"]\arrow[d,"g^{H\circ G}"']
 &\Sigma^G\times\Sigma^H\arrow[d,"g\times h"]\\
 \mathcal C(A,C)
 &\mathcal C(A,B)\times\mathcal C(B,C)\arrow[l,"\circ"]
\end{tikzcd}
\end{center}

The selection relation now acts on each owner's product policy
carrier, after the comparison has been pulled back through the entire
composite. Section~\ref{sec:egt-equilibrium} derives the resulting best
responses and equilibria. For instance, if one owner controls two bits
with value $u(x,y)=xy$, the policy $(0,0)$ has a profitable joint improvement to $(1,1)$,
although changing either bit alone does not improve its value.
Ownerwise composition retains this single-owner alternative.

Reparametrization and reference substitution must remain distinct.
The map $r_i:\Theta_i\to\Sigma_i$ changes the implemented policy
presentation by precomposition; substitution of a state reference changes
which policy is currently being tested but leaves the alternative carrier
intact. The transport laws in Theorem~\ref{thm:fib-strategic-transport}
make this difference precise.

\subsubsection{The semantic composition map on open agent resources}
\label{agt-app-coherence}

Let $G=[\Lin\Sigma\otimes A,B]$ and
$H=[\Lin T\otimes B,C]$. The uncurried sequential operation is
\begin{align*}
 G\otimes H\otimes\Lin(\Sigma\times T)\otimes A
 &\cong H\otimes\Lin T\otimes G\otimes\Lin\Sigma\otimes A\\
 &\xrightarrow{1\otimes1\otimes\ev}
       H\otimes\Lin T\otimes B
 \xrightarrow{\ev} C.
\end{align*}
Currying gives the map of Section~\ref{agt-sec-composition}.
For the parallel operation, rearrange
\[
 G\otimes H\otimes\Lin(\Sigma\times T)\otimes A\otimes A'
\]
into the tensor of $G\otimes\Lin\Sigma\otimes A$ and
$H\otimes\Lin T\otimes A'$, evaluate each, and curry. These descriptions
prove naturality on arbitrary test objects and open linear contexts.
They do not reduce equality of morphisms to equality only on global states.

\section{Control, Coalitions, and Delegation}\label{sec:paper-4}

Composition determines the complete policy coordinates, but agency
also requires a way to identify which of them a group may change.
For a coalition, the kernel pair of the projection onto outsiders fixes
exactly the policies beyond its control. We first prove that these
alternatives compose, then examine regrouping and delegation. Regrouping
changes the allocation of control; delegation implements a parent's
policies through subordinate decisions. The latter comparison is exact on
the implemented policy image. This distinction is related to additive and
multiplicative control in Cartesian frames
\citep{GarrabrantEtAl2021CartesianFrames}.

\subsection{Complete-policy alternatives and their composition}

\begin{definition}[Counterfactual policy alternatives]
\label{def:egt-counterfactual}
For a coalition $C\subseteq N$ of an open arena
$G=(\vec\Sigma,g):A\to B$, its \textbf{counterfactual alternative
span} is the kernel pair of the outsider projection:
\[
 \Sigma^G\xleftarrow{p_1}
 D_C=\Sigma^G\times_{\Sigma_{-C}^G}\Sigma^G
 \xrightarrow{p_2}\Sigma^G,\qquad
 \mathsf{Alt}_C(\sigma)=p_1^{-1}(\sigma).
\]
The counterfactual realization and its comparison tests have types
\[
 \delta_C^G=(gp_1,gp_2):D_C\to\mathcal C(A,B)^2,
 \qquad (\delta_C^G)^*P_i\in\mathsf{Pred}(D_C).
\]
Thus a reference and its alternative retain identical outsider policies.
The deviation relates a reference policy to its available counterfactual alternatives. The following construction gives the projections in ownerwise
coordinates and their compositional comparison maps.
\end{definition}
These kernel pairs support both existential and universal deviation
modalities. In either case the reference policy remains fixed while the
quantifier ranges over alternatives obtained by changing the coalition's
coordinates. No temporal move is implied by the word
``counterfactual.''

\begin{construction}[Deviation kernel pairs]
For $\Sigma=\prod_i\Sigma_i$, write $\Sigma_C=\prod_{i\in C}\Sigma_i$.
Define
\begin{equation}\label{eq:egt-deviation-kernel}
 D_C=\Sigma\times_{\Sigma_{-C}}\Sigma,
 \qquad p_1,p_2:D_C\rightrightarrows\Sigma,
\end{equation}
using the outsider projection $\pi_{-C}$. Put $D_i=D_{\{i\}}$.
The coordinate bijection is
\begin{equation}\label{eq:egt-deviation-coordinates}
 \chi_C:\Sigma_{-C}\times\Sigma_C\times\Sigma_C\xrightarrow{\cong}D_C,
 \quad (u,a,b)\longmapsto((u,a),(u,b)).
\end{equation}
For $G=(\vec\Sigma,g)$ its counterfactual realization is
\begin{equation}\label{eq:egt-delta-C}
 \delta_C^G=(gp_1,gp_2):D_C\to\mathcal C(A,B)^2.
\end{equation}
The first leg retains the reference profile and the second its alternative.
\end{construction}

\[
\begin{tikzcd}
 D_C\arrow[r,"p_2"]\arrow[d,"p_1"']
 \arrow[dr,phantom,"\lrcorner",very near start]&\Sigma\arrow[d,"\pi_{-C}"]\\
 \Sigma\arrow[r,"\pi_{-C}"']&\Sigma_{-C}
\end{tikzcd}
\qquad
\begin{tikzcd}
 D_C\arrow[r,"{(p_1,p_2)}"]\arrow[dr,"\delta_C^G"']&
 \Sigma^2\arrow[d,"g\times g"]\\
 &\mathcal C(A,B)^2
\end{tikzcd}
\]

\begin{construction}[Counterfactualization]
Define $\mathsf{Dup}_C:\mathbf{Par}_N\to\mathbf{Par}_N$ by squaring the
coordinates in $C$, leaving the others unchanged, and acting on maps by
$f_i\times f_i$ or $f_i$ respectively. Identify
$\Pi_N\mathsf{Dup}_C\vec\Sigma$ with $D_C$ through $\chi_C$.
On objects put $\mathsf{CF}_C(A)=(A,A)$; on open arenas put
\[
 \mathsf{CF}_C(G)=(\mathsf{Dup}_C\vec\Sigma,\delta_C^G)
 \quad\text{in }\mathbf{Arena}_N(\mathcal C\times\mathcal C).
\]
On $2$-cells use $\mathsf{Dup}_Cf$.
Write $\operatorname{Comp}((f_0,f_1),(h_0,h_1))
=(h_0f_0,h_1f_1)$ and define $\operatorname{Ten}$ coordinatewise.
\end{construction}

\begin{theorem}[Compositional counterfactualization]
\label{thm:egt-counterfactual}\label{prop:egt-composite-deviation}
$\mathsf{CF}_C$ is a symmetric monoidal pseudofunctor. In particular there
is a canonical bijection
$\rho_C:D_C^{H\circ G}\cong D_C^G\times D_C^H$ such that
\begin{equation}\label{eq:egt-counterfactual-composition}
 \delta_C^{H\circ G}
 =\operatorname{Comp}(\delta_C^G\times\delta_C^H)\rho_C.
\end{equation}
The tensor equation uses $\operatorname{Ten}$.
\end{theorem}
\begin{proof}
\leavevmode
\begin{enumerate}[label=\textup{(\arabic*)},leftmargin=*,itemsep=4pt,topsep=3pt]
  \item \textbf{Ownerwise alternatives.} The strong product comparison of $\mathsf{Dup}_C$ at a coalition coordinate
is the bijection
\[
 ((a,b),(a',b'))\longmapsto((a,a'),(b,b')).
\]
Outside the coalition it is the identity product comparison. These comparisons and $1\times1\cong1$ are
natural and satisfy symmetric monoidal coherence, since every composite
has the same action on named coordinates. They induce $\rho_C$.
  \item \textbf{Policy-level comparison.} Both realizations in the displayed equation send a tuple of reference and
alternative policies to
$(h_\tau g_\sigma,h_{\tau'}g_{\sigma'})$. A realization-preserving
$2$-cell stays realization-preserving after duplication. The pseudofunctor
and monoidal coherence equations reduce to these product comparisons and
the equations in $\mathcal C\times\mathcal C$.
\end{enumerate}
\end{proof}

\[
\begin{tikzcd}[column sep=large]
 D_C^{H\circ G}\arrow[r,"\rho_C", "\cong"']
 \arrow[d,"\delta_C^{H\circ G}"']&D_C^G\times D_C^H
 \arrow[d,"\delta_C^G\times\delta_C^H"]\\
 \mathcal C(A,C)^2&\mathcal C(A,B)^2\times\mathcal C(B,C)^2
 \arrow[l,"\operatorname{Comp}"]
\end{tikzcd}
\]

\begin{corollary}[Compositional counterfactual tests]
For a predicate $P\hookrightarrow\mathcal C(A,C)^2$,
\[
 (\delta_C^{H\circ G})^*P
 =\rho_C^*(\delta_C^G\times\delta_C^H)^*\operatorname{Comp}^*P.
\]
\end{corollary}
\begin{proof}
\leavevmode
\begin{enumerate}[label=\textup{(\arabic*)},leftmargin=*,itemsep=4pt,topsep=3pt]
  \item \textbf{Functoriality and substitution.} Apply functoriality of pullback to
Equation~\eqref{eq:egt-counterfactual-composition}.
\end{enumerate}
\end{proof}

The dependent reading of $D_C$ takes $p_1$ as the reference map and
the second projection as the alternative. Its paired realization is a term into a
constant comparison type. Thus the composition result above is also a
substitution law in the state-context doctrine of
Theorem~\ref{thm:fib-wiring}: changing a reference preserves the entire
dependent family of alternative policies.

\subsection{Collective control and regrouping of agent interfaces}

A coalition may be treated as one controller without altering the
component policies it controls. Regrouping collects its coordinates into
a single owner factor, leaving the realized interaction unchanged.
This changes the deviation relation, not the members' preferences:
collective evaluation is still supplied by a separate comparison. The
construction therefore distinguishes common control from a common
objective.

The coalition-deviation relation $D_C$ was defined in
Equation~\eqref{eq:egt-deviation-kernel} as the kernel pair of
$\pi_{-C}:\Sigma\to\Sigma_{-C}$. Regrouping will identify it
with the unilateral-deviation relation of the new collective owner,
without changing the underlying profile tuple.

\begin{proposition}
\label{prop:egt-coalition-composition}
On a full product profile space,
\[
 D_C;D_D=D_{C\cup D}.
\]
\end{proposition}
\begin{proof}
\leavevmode
\begin{enumerate}[label=\textup{(\arabic*)},leftmargin=*,itemsep=4pt,topsep=3pt]
  \item \textbf{Composition of coalition changes.} Two successive changes in $C$ and $D$ preserve every coordinate outside
$C\cup D$. Conversely, if $\sigma$ and $\tau$ agree outside $C\cup D$,
let $\rho$ have $\tau$'s coordinates on $C$ and $\sigma$'s coordinates
elsewhere. Then $\sigma D_C\rho$ and $\rho D_D\tau$. The construction
also covers overlap between $C$ and $D$.
\end{enumerate}
\end{proof}

\begin{construction}[Regrouping by right Kan extension]
Regard an owner
family $\vec\Sigma$ as a functor
\[
 \Sigma:N_{\mathrm{disc}}\longrightarrow\mathbf{Set}.
\]
For a surjection of owner sets $c:N\twoheadrightarrow J$, regarded as a
functor between the corresponding discrete categories, define
\[
 c_*\Sigma=\operatorname{Ran}_c\Sigma.
\]
Because the categories are discrete, the right Kan extension is the finite
product over each fiber:
\begin{equation}
\label{eq:egt-owner-kan}
 (c_*\Sigma)_j
 =\prod_{i\in c^{-1}(j)}\Sigma_i.
\end{equation}
There is a canonical flattening bijection
\begin{equation}
\label{eq:egt-owner-flattening}
 \kappa_\Sigma:
 \prod_{j\in J}(c_*\Sigma)_j
 \xrightarrow{\ \cong\ }
 \prod_{i\in N}\Sigma_i.
\end{equation}
\end{construction}

\begin{theorem}[Owner regrouping]
\label{thm:egt-owner-regrouping}
Every surjection $c:N\twoheadrightarrow J$ induces a symmetric-monoidal
pseudofunctor
\[
 c_*:\mathbf{Arena}_N(\mathcal C)
 \longrightarrow
 \mathbf{Arena}_J(\mathcal C)
\]
which is the identity on arenas. For an open arena
$G=(\vec\Sigma,g):A\to B$ its regrouping has parameter family
$c_*\vec\Sigma$ and realization
\[
 g^c=g\circ\kappa_\Sigma.
\]
For a $2$-cell $f:G\Rightarrow G'$, its $j$th component is
\[
 (c_*f)_j=\prod_{i\in c^{-1}(j)}f_i.
\]
Under the flattening bijection, deviation by the new owner $j$ is exactly
deviation by the old coalition $c^{-1}(j)$:
\[
 D_j^{\,c_*G}\cong D_{c^{-1}(j)}^{\,G}.
\]
For another surjection $d:J\twoheadrightarrow K$ there is a canonical
pseudonatural isomorphism
\[
 d_*c_*\cong(d\circ c)_*.
\]
\end{theorem}
\begin{proof}
\leavevmode
\begin{enumerate}[label=\textup{(\arabic*)},leftmargin=*,itemsep=4pt,topsep=3pt]
  \item \textbf{Naturality.} Naturality of the flattening bijections gives
\[
 \kappa_{\Sigma'}\circ\prod_j(c_*f)_j
 =\left(\prod_i f_i\right)\circ\kappa_\Sigma.
\]
Hence, if $g=g'\prod_i f_i$, then
\[
 g\kappa_\Sigma
 =g'\kappa_{\Sigma'}\prod_j(c_*f)_j,
\]
so $c_*f$ is realization-preserving.
  \item \textbf{Monoidal comparison.} The right Kan extension is strong monoidal for the pointwise Cartesian
products used in owned composition and tensor. Its component is the canonical
bijection
\[
 \prod_{i\in c^{-1}(j)}(\Sigma_i\times\Theta_i)
 \cong
 \left(\prod_{i\in c^{-1}(j)}\Sigma_i\right)
 \times
 \left(\prod_{i\in c^{-1}(j)}\Theta_i\right),
\]
and it sends the singleton owner family to the singleton owner family.
  \item \textbf{Monoidal comparison: second comparison.} These maps supply the compositor and unitor of the induced pseudofunctor. Their coherence is Cartesian-product coherence. Since composition and tensor
of realized interactions are evaluated after flattening the same parameter
tuple, the required realization equations are respectively associativity and
interchange in $\mathcal C$. The same strong-monoidal maps give compatibility
with tensor and symmetry.
  \item \textbf{Ownerwise alternatives.} For the deviation statement, forgetting the new $j$-coordinate from
$\prod_{k\in J}(c_*\Sigma)_k$ forgets, after $\kappa_\Sigma$, precisely the
old coordinates indexed by $c^{-1}(j)$. The two forgetful maps therefore
commute up to the canonical bijections, and their kernel pairs are
canonically isomorphic. This yields
$D_j^{c_*G}\cong D_{c^{-1}(j)}^G$.
  \item \textbf{Iterated right Kan extension.} Finally, the component at $k\in K$ of $d_*c_*\Sigma$ is
\[
 \prod_{j\in d^{-1}(k)}\prod_{i\in c^{-1}(j)}\Sigma_i,
\]
canonically isomorphic to
\[
 \prod_{i\in(d\circ c)^{-1}(k)}\Sigma_i
 =((d\circ c)_*\Sigma)_k.
\]
These reassociations are natural and satisfy the pseudofunctorial coherence
because every composite is the same flattening of the same owner-indexed
tuple.
\end{enumerate}
\end{proof}

\noindent Flattening preserves realized arrows while identifying a regrouped owner with its original coalition.
\begin{center}
\begin{tikzcd}[column sep=large,row sep=large]
\prod_{j\in J}(c_*\Sigma)_j\arrow[r,"\kappa_\Sigma","\cong"']\arrow[dr,"g^c"']
 &\prod_{i\in N}\Sigma_i\arrow[d,"g"]\\
 &\mathcal C(A,B)
\end{tikzcd}
\end{center}

Regrouping transfers strategic authority, not economic judgement.
A welfare order, transfer rule, or Pareto comparison supplies a collective
preference separately. The next construction, delegation, also changes
which choices implement that authority.

\subsection{Delegation and substitution of agent interfaces}
\label{sec:egt-delegation}

Regrouping reallocates control over existing coordinates; delegation
must additionally say how subordinate policies implement each parent
choice. The implementation maps therefore precede any comparison of
preferences. This follows the separation of agency from selection in
\citet{CapucciEtAl2022Agency} and makes the distinction from Cartesian
frame subagents explicit \citep{GarrabrantEtAl2021CartesianFrames}.
The main issue is which parent alternatives remain available after the
subordinate implementation is fixed.

\begin{construction}[Substitution of controlled policies]
\label{con:egt-agent-substitution}
Let $G=(\vec\Sigma,g):A\to B$ have finite owner set $N$.
A delegation presentation consists of a finite set $J$, a surjection
$c:J\twoheadrightarrow N$, subordinate policy carriers $\Theta_j$, and maps
\[
 h_i:\prod_{j\in c^{-1}(i)}\Theta_j\longrightarrow\Sigma_i.
\]
Write $\Theta=\prod_{j\in J}\Theta_j$. Regrouping and the product of the
$h_i$ give $h:\Theta\to\Sigma$. The substituted process
\[
 G[h]=(\vec\Theta,gh):A\to B
\]
is owned by $J$. If $L\hookrightarrow\Sigma$ is the legal domain, its
substituted legal domain is $L_h=h^{-1}(L)$.
The construction retains the implementation maps $h_i$: an unimplemented
parent policy is not silently made available to the delegates.
\end{construction}

\begin{theorem}[Associative delegation and compatibility with composition]
\label{thm:egt-delegation-substitution}
For $d:K\twoheadrightarrow J$ and
$k_j:\prod_{\ell\in d^{-1}(j)}\Xi_\ell\to\Theta_j$, define
\[
 (h\star k)_i
 =h_i\circ\left(\prod_{j\in c^{-1}(i)}k_j\right)\circ\kappa_i,
\]
where $\kappa_i$ groups the $\Xi$-coordinates over $(cd)^{-1}(i)$
by their intermediate owner $j$. Then
\[
 (G[h])[k]\cong G[h\star k],\qquad
 (L_h)_k=L_{h\star k}.
\]
Identity implementation maps are units, and three successive substitutions
agree under the canonical product reassociations.

Suppose also that $H:B\to C$ has owner set $N$, and that
$r_i:\prod_{j\in c^{-1}(i)}\Lambda_j\to\Sigma_i^H$ implements
its policies using the same owner map $c$. Define $h\boxtimes r$ at owner
$i$ by the canonical map
\[
 \begin{aligned}
 \prod_{j\in c^{-1}(i)}(\Theta_j\times\Lambda_j)
 &\longrightarrow
 \left(\prod_{j\in c^{-1}(i)}\Theta_j\right)
 \times\left(\prod_{j\in c^{-1}(i)}\Lambda_j\right)\\
 &\xrightarrow{h_i\times r_i}\Sigma_i^G\times\Sigma_i^H.
 \end{aligned}
\]
There are canonical realization-preserving comparisons
\[
 (H\circ G)[h\boxtimes r]\cong H[r]\circ G[h].
\]
The same assertion holds for tensor whenever the boundaries are those of
a tensor product. In particular, delegation respects the complete-policy
composition law without identifying the subordinate owners.
\end{theorem}
\begin{proof}
\leavevmode
\begin{enumerate}[label=\textup{(\arabic*)},leftmargin=*,itemsep=4pt,topsep=3pt]
  \item \textbf{Monoidal comparison.} On a subordinate tuple $\xi$, the $i$th parent coordinate of either
iterated implementation is
$h_i((k_j((\xi_\ell)_{\ell\in d^{-1}(j)}))_{j\in c^{-1}(i)})$. Thus the realizations agree after the displayed product bijection. Inverse images compose, giving the legal-domain equality.
  \item \textbf{Monoidal comparison: second comparison.} For three
levels, both bracketings evaluate the same nested coordinate functions;
the product reassociations are coherent because they have identical
coordinate projections. Identity maps leave every coordinate unchanged.
  \item \textbf{Composition.} For sequential composition, both sides have, after rearrangement, the
$J$-indexed carrier $(\Theta_j\times\Lambda_j)_j$. Their value at
$(\theta,\lambda)$ is
$r_H(\lambda)\circ r_G(\theta)$, where
$r_G=gh$ and $r_H$ is the realization of $H$ composed with $r$. For tensor the common value is $r_G(\theta)\otimes r_H(\lambda)$.
  \item \textbf{Witness and legality comparison.} The comparison maps are products of coordinate rearrangements, so preserve
realizations, identities, and the associativity and symmetry comparisons. If legal domains on a composite are specified by an additional
comprehension, pulling that comprehension back through either equal
implementation map gives the same legal domain. No independence of that
additional rule is required.
\end{enumerate}
\end{proof}

\noindent Associative delegation expresses equality of the iterated implementation maps.
\begin{center}
\begin{tikzcd}[column sep=large,row sep=large]
\Xi\arrow[r,"k"]\arrow[dr,"h\star k"']
 &\Theta\arrow[d,"h"]\\
 &\Sigma
\end{tikzcd}
\end{center}

\begin{lemma}[Ownerwise lifting of legal alternatives]
\label{lem:ownerwise-legal-lifting}
Let $r_i:T_i\twoheadrightarrow S_i$ be surjections over a finite owner
set $N$, put $r=\prod_i r_i$, and let $L\hookrightarrow S=\prod_iS_i$
be any legal domain. With $K=r^{-1}(L)$, for every coalition
$C\subseteq N$ the canonical map over $K$
\[
 \lambda_C:D_C^K\longrightarrow K\times_{L,p_1}D_C^L,
 \qquad(x,x')\longmapsto(x,(r(x),r(x')))
\]
is surjective. For any $g:S\to H$, it preserves both realized endpoints under
the factorized realization $g\circ r$. Writing $g_L=g|_L:L\to H$ and $(gr)_K=(gr)|_K:K\to H$,
and setting
$\delta_C^{g,L}=(g_Lp_1,g_Lp_2)$ and
$\delta_C^{gr,K}=((gr)_Kp_1,(gr)_Kp_2)$, for every comparison
$Q\hookrightarrow H^2$ one consequently has an equality of predicates on
$K$:
\begin{equation}
\label{eq:ownerwise-legal-comparison-pullback}
 r|_K^*\,\exists_{p_1:D_C^L\to L}
     (\delta_C^{g,L})^*Q
 =\exists_{p_1:D_C^K\to K}(\delta_C^{gr,K})^*Q.
\end{equation}
\end{lemma}
\begin{proof}
\leavevmode
\begin{enumerate}[label=\textup{(\arabic*)},leftmargin=*,itemsep=4pt,topsep=3pt]
  \item \textbf{Lifting along the surjection.} A point of the target is $x\in K$ and $y\in L$ agreeing with $r(x)$
outside $C$. Keep $x'_i=x_i$ for $i\notin C$. For each $i\in C$,
surjectivity of $r_i$ provides $x'_i\in T_i$ with $r_i(x'_i)=y_i$.
  \item \textbf{Witness and legality comparison.} The finitely many coordinate witnesses assemble to $x'\in T$. Then $r(x')=y\in L$, so $x'\in K$, and $x'$ agrees with $x$ outside
$C$. Hence $(x,x')\in D_C^K$ maps to $(x,(r(x),y))$.
  \item \textbf{Ownerwise alternatives.} The case $C=\varnothing$ uses $x'=x$, including empty owner lists. Both endpoints are preserved because the delegated realization is
$g\circ r$ by construction. To prove \eqref{eq:ownerwise-legal-comparison-pullback}, fix $x\in K$.
  \item \textbf{Witness and legality comparison.} The predicate on the right has a witness $x'$ exactly when
$xD_C^Kx'$ and $Q(grx,grx')$. Such a witness maps through $\lambda_C$
to the target alternative $(rx,rx')$, giving the left predicate.
  \item \textbf{Reverse implication.} Conversely a witness $y\in L$ for the left predicate is a target
point $(x,(rx,y))$ of $\lambda_C$; the explicit coordinate construction
above lifts it to $x'$ at that \emph{same} reference $x$. The two
realized endpoints agree, so $Q(grx,grx')$ holds. This gives equality
of the existential predicates; applying negation and intersecting the
specified coalitional tests gives exact stability transport.
\end{enumerate}
\end{proof}

\noindent Legal alternatives lift along the ownerwise surjection while keeping the reference fixed.
\begin{center}
\begin{tikzcd}[column sep=large,row sep=large]
D_C^K\arrow[r,"\lambda_C",two heads]\arrow[d,"p_1"']
 &K\times_L D_C^L\arrow[d,"\pi_1"]\\
 K\arrow[r,equal]&K
\end{tikzcd}
\end{center}

\begin{theorem}[Exact strategic comparison with implemented parent policies]
\label{thm:egt-delegation-image}
For the preceding delegation, put $T_i=\operatorname{im}(h_i)$,
$T=\prod_iT_i$, and $L_T=L\cap T$. Restrict the parent process to $T$
and regroup the subordinate process by $c$. Write
$\Theta^c=\prod_{i\in N}\prod_{j\in c^{-1}(i)}\Theta_j$
for the regrouped profile carrier, identified with $\Theta$
by flattening. The induced map
\[
 \bar h:\Theta^c\longrightarrow T
\]
is ownerwise surjective and takes $L_h$ onto $L_T$.
For every coalition $I\subseteq N$, every $\theta\in L_h$, and every
$t'\in L_T$ agreeing with $\bar h(\theta)$ outside $I$, there is
$\theta'\in L_h$ agreeing with $\theta$ outside the blocks indexed by $I$
and satisfying $\bar h(\theta')=t'$. Realizations of both the reference
and the alternative are preserved.

Consequently, writing $\bar h_L:L_h\twoheadrightarrow L_T$
for its legal restriction and pulling the same parent comparisons
back through realization,
\[
 \mathsf{NE}_{c_*G[h],\,L_h}
 =\bar h_L^{\,*}\mathsf{NE}_{G|_T,\,L_T}.
\]
The same equality holds for stability against any specified family of
parent coalitions, using the corresponding whole blocks of delegates and
the same coalition comparisons.
\end{theorem}
\begin{proof}
\leavevmode
\begin{enumerate}[label=\textup{(\arabic*)},leftmargin=*,itemsep=4pt,topsep=3pt]
  \item \textbf{Lifting along the surjection.} Corestricting each implementation map $h_i$ to its image $T_i$
produces an ownerwise surjection $\bar h:\Theta^c\twoheadrightarrow T$. Its inverse-image legal domain is exactly $L_h=\bar h^{-1}(L_T)$. \Cref{lem:ownerwise-legal-lifting} now gives the required lift of every
legal $I$-alternative, keeping the entire outsider blocks fixed and
preserving both concurrent realizations.
  \item \textbf{Reverse implication.} Conversely, every legal
subordinate block deviation maps to a legal parent deviation. Thus an improving deviation exists at a subordinate reference if and
only if one exists at its implemented parent reference. Universal
negation over singleton parent owners proves the displayed Nash equality;
the same comparison for each designated coalition proves the final
stability statement.
\end{enumerate}
\end{proof}

\noindent The image factorization retains the exact implemented parent policies.
\begin{center}
\begin{tikzcd}[column sep=large,row sep=large]
L_h\arrow[r,"\bar h_L",two heads]\arrow[d,hook]
 &L_T\arrow[d,hook]\\
 \Theta^c\arrow[r,"\bar h",two heads]
 &T
\end{tikzcd}
\end{center}

The image factorization answers that issue even when the delegation
maps $h_i$ are not surjective onto the original parent carriers: every
implemented parent alternative has a subordinate lift with the outsider
blocks fixed. Collective authority is nevertheless distinct from a
collective preference. If delegates share their parent's comparison,
whole-block stability implies individual stability, but not conversely.
For two delegate bits $x,y\in\{0,1\}$, with $h(x,y)=(x,y)$ and common
value $xy$, the profile $(0,0)$ has no profitable isolated change, while
the block improves to $(1,1)$. Thus exact parent comparison requires
whole-block deviations.

\begin{proposition}[The Cartesian frame underlying a control interface]
\label{prop:egt-cartesian-frame-comparison}
On an unrestricted profile carrier, every coalition $C$ determines the
Cartesian frame
\[
 \mathsf{Frame}_C(G)
 =\left(\Sigma_C,\Sigma_{-C},
       (a,e)\longmapsto g_{(a,e)}\right)
\]
with outcome set $\mathcal C(A,B)$. If $C\subseteq D$, write
$X=\Sigma_C$, $Y=\Sigma_{D\setminus C}$, and $Z=\Sigma_{-D}$.
The two frames have the form
\[
 (X,Y\times Z,p),\qquad (X\times Y,Z,p)
\]
for the same three-variable outcome function. Thus their relation is
exactly the multiplicative redistribution of control and environment
used in Cartesian frames. Restricting the controlled carrier to
$K\subseteq\Sigma_C$ instead gives the corresponding additive restriction.
\end{proposition}
\begin{proof}
\leavevmode
\begin{enumerate}[label=\textup{(\arabic*)},leftmargin=*,itemsep=4pt,topsep=3pt]
  \item \textbf{Policy-level comparison.} Reassociate the profile product as $X\times Y\times Z$. For $C$, the $Y$-coordinates belong to the environment; for $D$, they
belong to the controlled coordinates. Both evaluation maps are the
same $g$ after this reassociation. Restricting the first coordinate
of the evaluation to $K$ gives its additive restriction.
\end{enumerate}
\end{proof}

A legal domain may couple controlled and outsider coordinates. The
available choices are then a family $S_C^L(e)$ rather than a fixed
rectangular carrier, so the Cartesian frame comparison does not identify
all the paper's morphisms. In particular Chu-style frame maps reverse the
environment coordinate, whereas our owned maps preserve specified
implementations. Regrouping establishes control, while corporate
judgement and commitment still require aggregation or update data; compare
\citet{ListPettit2011GroupAgency}.

Legal comprehension turns this distinction into a dependent
interface and supplies the legal context $(L,m)$. For
$m:L\hookrightarrow\Sigma$, the kernel pair of the
outsider map gives
$D_C^L=L\times_{L_{-C}}L$; its two legs select a legal reference
and a legal alternative term in $\Delta_GL$. The following dependent
choice spaces are the
fibres of that span, so coupled legal opportunities are not replaced by
unrestricted products.

\subsection{Displayed coalition and delegation operations}
\label{agt-sec-ownership}

\subsubsection{Counterfactual alternatives in agent syntax}

For $C\subseteq N$, the unrestricted deviation is
$D_C=\Sigma\times_{\Sigma_{-C}}\Sigma$, as constructed in
\cref{def:egt-counterfactual,thm:egt-counterfactual}.
On a legal domain $L\hookrightarrow\Sigma$, its actual alternatives
are instead
\begin{equation}
 D_C^L=L\times_{\Sigma_{-C}}L.
 \label{agt-eq-legal-deviation}
\end{equation}
Their paired realization is $(a[\sigma],a[\tau])$.
The coordinate comparison in \cref{thm:egt-counterfactual} gives
$D_C^{\sig\times\taus}\cong D_C^{\sig}\times D_C^{\taus}$,
including the sequential and parallel agent presentations. Likewise
\cref{prop:egt-coalition-composition} proves
$D_C;D_D=D_{C\cup D}$ on the unrestricted profile product.
For legal carriers the intermediate profile of that proof need not be
legal; legal relations compose by their actual pullback and image,
not by replacing them with unrestricted products.

\subsubsection{Collective agents by regrouping}

\begin{construction}[Collective control]
Let $c:N\twoheadrightarrow J$ be a surjection. Define
\[
 (c_*\sig)_j=\prod_{i\in c^{-1}(j)}\Sigma_i,
 \qquad
 \kappa_{\sig}:\prod_{j\in J}(c_*\sig)_j\xrightarrow{\cong}\Sigma.
\]
The regrouped agent $c_*a$ has execution
$(c_*a)[\theta](x)=a[\kappa_{\sig}(\theta)](x)$.
Its underlying linear type is canonically isomorphic to that of $a$, while
its displayed control interface has owner set $J$.
\end{construction}

\begin{theorem}[Collective composition and coalition deviations]
\label{agt-thm-regroup}
Regrouping is compatible with sequential and parallel agent composition and
with ownerwise maps. It satisfies
\[
 d_*c_*a\cong(d\circ c)_*a,
 \qquad
 D_j^{c_*a}\cong D_{c^{-1}(j)}^a.
\]
The deviation comparison also holds on the transported legal domain.
\end{theorem}
\begin{proof}
\leavevmode
\begin{enumerate}[label=\textup{(\arabic*)},leftmargin=*,itemsep=4pt,topsep=3pt]
  \item \textbf{Ownerwise alternatives.} Apply the regrouping pseudofunctor of
\cref{thm:egt-owner-regrouping} through the owned-presentation
isomorphism of \cref{agt-thm-owned-equivalence}. Its compositor at new
owner $j$ is explicitly
\[
 \prod_{ci=j}(\Sigma_i\times T_i)
 \cong\left(\prod_{ci=j}\Sigma_i\right)
       \times\left(\prod_{ci=j}T_i\right).
\]
Both routes have identical policy projections and concurrent realizations.
  \item \textbf{Descent along the kernel pair.} Forgetting $j$ forgets exactly the old $c^{-1}(j)$ coordinates; the
kernel-pair comparisons restrict along any transported legal inclusion.
\end{enumerate}
\end{proof}

\noindent The kernel-pair comparison preserves the deviations controlled by each regrouped owner.
\begin{center}
\begin{tikzcd}[column sep=large,row sep=large]
D_j^{c_*a}\arrow[r,"\cong"]\arrow[d,shift left,"p_1"]\arrow[d,shift right,"p_2"']
 &D_{c^{-1}(j)}^a\arrow[d,shift left,"p_1"]\arrow[d,shift right,"p_2"']\\
 \Sigma^{c_*a}\arrow[r,"\kappa_\Sigma","\cong"']
 &\Sigma^a
\end{tikzcd}
\end{center}

A coalition $C$ is represented by the partition with block $C$ and singleton
outsiders. The empty coalition has unit policy carrier and identity
alternatives. Overlapping coalitions are represented by their different
projections $\Sigma\to\Sigma_{-C}$; they need not be blocks of a single
partition. Thus the calculus retains both nested collective agents and
overlapping coalitional tests.

Regrouping determines joint control. Collective observation is constructed
in Sections~\ref{agt-sec-perspective} and~\ref{agt-sec-actions}; collective
evaluation is constructed in Section~\ref{agt-sec-equilibrium}. The original
owner list and regrouping map remain displayed, so members' separate
observations and comparisons are not lost.

\subsubsection{Delegation as implementation of policies}

For a parent owner set $N$, take subordinate owners $J$, a surjection
$c:J\twoheadrightarrow N$, and maps
\[
 h_i:\prod_{cj=i}\Theta_j\longrightarrow\Sigma_i.
\]
Flattening gives $h:\Theta\to\Sigma$. The delegated agent is $a[h]$,
owned by $J$, with execution $a[h(\theta)]$ and legal domain $h^{-1}L$.
For a further implementation $k$, define $h\star k$ by composing the
coordinate maps after regrouping their subordinate inputs.

\begin{proposition}[Associative implementation]
\label{agt-prop-delegation}
There are coherent comparisons
\[
 (a[h])[k]\cong a[h\star k],
 \qquad (h^{-1}L)\text{ pulled back along }k=(h\star k)^{-1}L.
\]
Delegation commutes with sequential and parallel composition through the
ownerwise product of implementation maps.
\end{proposition}
\begin{proof}
\leavevmode
\begin{enumerate}[label=\textup{(\arabic*)},leftmargin=*,itemsep=4pt,topsep=3pt]
  \item \textbf{Functoriality and substitution.} The implementation maps $h\star k$ and the natural composition
comparisons are those constructed in
\cref{thm:egt-delegation-substitution}. The owned-presentation equivalence
\cref{agt-thm-owned-equivalence} preserves their explicit flattening
maps and realized arrows. Legal domains agree by functoriality of
inverse image, $(h\circ k)^{-1}L=k^{-1}(h^{-1}L)$.
\end{enumerate}
\end{proof}

\begin{theorem}[Exact implemented alternatives]
\label{agt-thm-delegation-image}
Put $T_i=\im(h_i)$, $T=\prod_iT_i$, and $L_T=L\cap T$. The regrouped
delegated legal carrier maps onto $L_T$. Every legal parent alternative in
a coalition $C\subseteq N$ lifts to a change of exactly the subordinate
blocks over $C$, with all other subordinate coordinates fixed.
\end{theorem}
\begin{proof}
\leavevmode
\begin{enumerate}[label=\textup{(\arabic*)},leftmargin=*,itemsep=4pt,topsep=3pt]
  \item \textbf{Lifting along the surjection.} After regrouping delegates into parent blocks, each implementation
corestricts to its image. The resulting ownerwise surjection has
inverse-image legal domain $h^{-1}L$. The explicit lift in
\cref{lem:ownerwise-legal-lifting} fixes the outsider blocks and realizes
the requested legal parent alternative.
  \item \textbf{Ownerwise alternatives.} Its converse takes the image
of any legal block deviation, as in \cref{thm:egt-delegation-image}.
\end{enumerate}
\end{proof}

\begin{example}[Delegated choices and collective authority]
Two delegates choose bits $x,y$, with common outcome value $xy$. At $(0,0)$
neither isolated delegate has a strictly improving change, but the whole
block can change to $(1,1)$. Consequently the exact strategic comparison in
Theorem~\ref{agt-thm-delegation-image} is between parent coalitions and whole
delegate blocks. The distinction is produced by the explicit projections
defining alternatives.
\end{example}

\section{Contextual Evaluation and Strategic Equilibrium}\label{sec:paper-5}

Ownerwise control specifies which policies can change, but not which
changes are profitable. To evaluate an agent, we must transport a
comparison on realized concurrent interactions back to its complete
policy space. We construct state contexts for reference profiles, derive
predicate feedback as the canonical backward action, and connect that
feedback to owner selection. The resulting equilibrium predicate uses
the same kernel pairs already defining control, and its finite
normal-form specialization recovers ordinary pure Nash equilibrium.
This is the concurrent interpretation of contextual open-game evaluation
\citep{CapucciEtAl2022Cybernetics}.

\subsection{Presentations, state contexts, and assertions}

The owned realization makes policies into ordinary parameters.
Reference profiles, however, can be indexed by richer state spaces whose
points also record histories or information. We first construct the
fibration of these state contexts and their assertion fibres. It will
supply the inverse-image operation used for coplay and keep reference
substitution distinct from restricting the policy space.

\begin{definition}[Economic state contexts]
\label{def:egt-economic-state-context}
A state context $s:W\to\Sigma^G$ is a family of reference profiles indexed
by $W$.
\end{definition}

Its points need not be policies themselves: an epistemic world may
also carry information about which other worlds an agent considers
possible. Several points of $W$ may consequently have the same profile.
A predicate on $W$ asks a question at those references, while the full policy
space still supplies alternatives. This use of state is distinct from
a monoidal state $I\to A$.

\begin{definition}[The tower over a concurrent interface]
\label{def:fib-tower}
Put $\mathcal A_{A,B}=\mathbf{Arena}_N(\mathcal C)(A,B)$.
The profile functor $U_{A,B}:\mathcal A_{A,B}\to\mathbf{Set}$ sends
$G$ to $\Sigma^G$ and an ownerwise reparametrisation $r$ to
$\bar r=\Pi_Nr$. The category $\mathsf{St}_{A,B}$ has objects
$(G,s:W\to\Sigma^G)$ and arrows
\[
 (r,u):(G,s:W\to\Sigma^G)\longrightarrow(H,t:V\to\Sigma^H),
 \qquad tu=\bar r s.
\]
An object of $\mathsf{Assn}_{A,B}$ additionally carries a predicate
$P\hookrightarrow W$. An arrow over $(r,u)$ from $P$ to $Q$ exists
exactly when $P\leq u^*Q$, and is then unique over that arrow.
Here $\mathsf{Pred}(W)=\operatorname{Sub}_{\mathbf{Set}}(W)$ and
$u^*Q=u^{-1}(Q)$. Composition uses composition of the two displayed
maps and functoriality of inverse image. The projections form the tower
\[
\begin{tikzcd}[column sep=large]
 \mathsf{Assn}_{A,B}\arrow[r,"p"]&
 \mathsf{St}_{A,B}\arrow[r,"q"]&
 \mathcal A_{A,B}\arrow[r,"\pi"]&\mathbf{Par}_N.
\end{tikzcd}
\]
The category $\mathsf{St}(G)=q^{-1}(G)$ is the category of
\textbf{state contexts} of $G$. Its objects retain their reference map
to the full policy space.
\end{definition}

\begin{theorem}[Cartesian lifts and the logical fibres]
\label{thm:fib-tower}
All three projections in Definition~\ref{def:fib-tower} are fibrations.
There are canonical identifications
\[
 \mathsf{St}(G)=\mathbf{Set}/\Sigma^G,
 \qquad
 p^{-1}(G,s:W\to\Sigma^G)=\mathsf{Pred}(W).
\]
More precisely, $p$ is a \emph{posetal} fibration, obtained by change of
base of the subobject fibration along the domain functor
$\mathsf{St}_{A,B}\to\mathbf{Set}$; $q$ is the change of base of
the codomain fibration along $U_{A,B}$; and $\pi$ is a
\emph{discrete} fibration, as the category-of-elements projection of
the presheaf $F_{A,B}$.
\end{theorem}
\begin{proof}
\leavevmode
\begin{enumerate}[label=\textup{(\arabic*)},leftmargin=*,itemsep=4pt,topsep=3pt]
  \item \textbf{Lifting through the comparison.} For $\pi$, lifting an owner map $r:\vec\Theta\to\vec\Sigma$
at $g$ gives $(\vec\Theta,g\Pi_Nr)$, and the factorization equation
is unique. This is the category-of-elements construction of
Theorem~\ref{thm:egt-owned}.
  \item \textbf{Lifting along the surjection: second comparison.} For $q$, lift $t:V\to\Sigma^H$ along $r:G\to H$ by the square
\[
\begin{tikzcd}
 \Sigma^G\times_{\Sigma^H}V\arrow[r,"a"]\arrow[d,"s"']
  \arrow[dr,phantom,"\lrcorner",very near start]&V\arrow[d,"t"]\\
 \Sigma^G\arrow[r,"\bar r"']&\Sigma^H.
\end{tikzcd}
\]
If $(r\ell,v):(K,z)\to(H,t)$ lies over $r\ell$, its unique
factorization through this square sends $w$ to $(\bar\ell z(w),v(w))$.
  \item \textbf{Cartesian lifting of assertion predicates.} This proves cartesianness. Finally the lift of $Q\hookrightarrow V$
along $(r,u)$ is $u^*Q\hookrightarrow W$. A factorization into it exists
exactly when its predicate is below the inverse image of $u^*Q$;
functoriality of inverse image gives precisely the required condition. The two change-of-base descriptions and the fibre identifications follow
from the definitions.
\end{enumerate}
\end{proof}

\noindent The state-context fibration is obtained by changing base along the policy projection.
\begin{center}
\begin{tikzcd}[column sep=large,row sep=large]
\Sigma^G\times_{\Sigma^H}V\arrow[r]\arrow[d]\arrow[dr,phantom,"\lrcorner",very near start]
 &V\arrow[d,"t"]\\
 \Sigma^G\arrow[r,"\bar r"']&\Sigma^H
\end{tikzcd}
\end{center}

\begin{construction}[Indexed presentation and constant types]
\label{con:fib-constants}
Write $\mathsf P_G:\mathsf{St}(G)^{\mathrm{op}}\to\mathbf{Pos}$
for the indexed predicate assignment $(W,s)\mapsto\mathsf{Pred}(W)$.
The terminal context is $1_G=(\Sigma^G\xrightarrow{\mathrm{id}}\Sigma^G)$.
For an ordinary set $X$, its constant type is
\[
 \Delta_GX=(\Sigma^G\times X\xrightarrow{\pi_1}\Sigma^G).
\]
For $s:W\to\Sigma^G$, a term $a:W\to X$ is represented by
$(s,a):(W,s)\to\Delta_GX$. A predicate $Q\hookrightarrow X$ is
represented by $\Delta_GQ=\Sigma^G\times Q\hookrightarrow\Sigma^G\times X$.
Consequently
\begin{equation}\label{eq:fib-constant-test}
 (s,a)^*\Delta_GQ=a^*Q.
\end{equation}
In particular play is the section
$\widehat g=(\mathrm{id},g):1_G\to\Delta_G\mathcal C(A,B)$.
Each owner projection likewise gives a section of
$\Delta_G\Sigma_i^G$.
\end{construction}

\begin{proposition}[Change of game and constant types]
\label{prop:fib-basechange}
A reparametrisation $r:G\to H$ induces pullback
$\bar r^*:\mathsf{St}(H)\to\mathsf{St}(G)$ and, on each predicate,
inverse image along the top map of its cartesian lift. There are natural
isomorphisms
\[
 \bar r^*\Delta_HX\cong\Delta_GX,
 \qquad \overline{tr}^{\,*}\cong\bar r^*\bar t^*.
\]
They satisfy the identity and associativity coherence equations. Under
the first isomorphism the pullback of $\widehat h$ is $\widehat g$.
\end{proposition}
\begin{proof}
\leavevmode
\begin{enumerate}[label=\textup{(\arabic*)},leftmargin=*,itemsep=4pt,topsep=3pt]
  \item \textbf{Pullback comparison isomorphisms.} The first map sends $(\sigma,(\bar r\sigma,x))$ to $(\sigma,x)$. The second forgets the repeated intermediate coordinate in an iterated
pullback. Their inverses insert the displayed coordinates.
  \item \textbf{Coherence.} Every
coherence composite therefore inserts or deletes the same coordinates. The final assertion is the realization equation $g=h\bar r$. Equivalently, represent a context as its family of fibres $(W_\sigma)$;
reindexing is $(W_\eta)\mapsto(W_{\bar r\sigma})$, with strict
substitution.
  \item \textbf{Equivalence with split families.} Passing between this split presentation and slices uses
the canonical bijection $W\cong\coprod_\sigma W_\sigma$.
\end{enumerate}
\end{proof}

\begin{construction}[Dependent fibres over state contexts]
\label{con:egt-state-dependent-fibres}
Pull back the indexed categories $\mathsf T$ and $\mathsf L$ along
$\operatorname{Dom}:\mathsf{St}_{A,B}\to\mathbf{Set}$, which sends
$(G,W,s)$ to $W$ and $(r,u)$ to $u$. Their fibres are
$\mathsf T(W)\simeq\mathbf{Set}/W$ and $\mathsf L(W)$, with
reindexing $u^*$. Their Grothendieck constructions are fibrations: the
cartesian arrow above $u$ is the identity $u^*E\to u^*E$ in the source
fibre, and every factorization is forced by composition of index maps.
The comparison, comprehension, and adjunction maps of
Theorem~\ref{thm:egt-dependent-model} pull back to these fibrations.
The original $\mathcal C_W$ is their constant representable linear
fragment, by Proposition~\ref{prop:egt-dependent-substitution}.
\end{construction}

The original assertion fibre is recovered as
$\mathsf{Pred}(W)=\operatorname{Sub}_{\mathsf T(W)}(1)$: a subterminal
family has a singleton or empty fibre at each state. These are the
ordinary subobject structures of
\citet[Chapters 1, 4, and 10]{Jacobs1999}, now situated inside the
constructed dependent linear--nonlinear model. For the original boundary,
the realization section is equivalently a section of
$\mathsf{Val}_W[yA,yB]_D$, since this family is constant with value
$\mathcal C(A,B)$.

\subsection{Predicate feedback, canonical lenses, and coplay}

The assertion fibration supplies a backward operation without adding
a new primitive: a forward map reindexes predicates by inverse image.
Representing this operation with power objects yields the canonical lens
used here as \emph{coplay}. Its forward value can be a complete concurrent
strategy, while its backward value is a predicate on alternatives. This is
a particular continuation semantics in the sense of
\citet{Capucci2023}; it does not claim that every game-theoretic coplay
is inverse image.

\begin{definition}[Play, predicate feedback, and canonical coplay]
\label{def:egt-play-continuation-coplay}
For a forward realization $f:X\to Y$, a \textbf{predicate
continuation} is $Q\in\mathsf{Pred}(Y)$. The subobject fibration
supplies its \textbf{predicate feedback} by reindexing. The
\textbf{canonical coplay} of the induced lens is this same map:
\[
 \mathsf{coplay}_f=f^*:\mathsf{Pred}(Y)
 \longrightarrow\mathsf{Pred}(X),\qquad Q\longmapsto f^{-1}(Q).
\]
A numerical or other $V$-valued continuation has type $k:Y\to V$,
and its backward transport has type
$f_V^-:V^Y\to V^X$, $k\mapsto kf$.
For a comparison predicate $P\hookrightarrow Y^2$, its transported
comparison is $(f\times f)^*P\hookrightarrow X^2$. These are the
backward components of the canonical lenses constructed below. Play
returns a result; coplay tests which alternatives produce desired
results, and does not execute a physical protocol backwards.
\end{definition}
For example, let $X$ encode keeping or exchanging a resource and let
$f$ return the resulting bundle. Pulling back the goal of possessing a
specified resource selects exactly the choices achieving it. The
calculation evaluates forward possibilities; it does not run the protocol
backwards.

\begin{definition}[Lenses and comparison interfaces]
\label{def:cgt-comparison-lens}
Write $\mathsf{Pred}(X)=\operatorname{Sub}_{\mathbf{Set}}(X)$.
A lens $(X,S)\to(Y,R)$ is a pair of maps
$f:X\to Y$ and $f^-:X\times R\to S$. Its composite with
$(h,h^-):(Y,R)\to(Z,T)$ is
\[
 (hf,\ (x,t)\longmapsto f^-(x,h^-(f(x),t))).
\]
The identity has backward map $(x,s)\mapsto s$. The tensor is the
componentwise product of forward and backward maps, with unit $(1,1)$.
These operations define $\mathsf{Lens}(\mathbf{Set})$.
The \textbf{comparison interface} on $X$ is
\[
 \mathbb L(X)=(X,\mathsf{Pred}(X)).
\]
A comparison $P\hookrightarrow X\times X$ is represented by the costate
$q_P:X\to\mathsf{Pred}(X)$, where
$q_P(x)=\{x':P(x,x')\}$. Thus its backward output at a reference $x$
is the predicate of alternatives compared favorably with that reference.
\end{definition}

\begin{definition}[Lens states and comparison costates]
\label{def:egt-lens-states-costates}
A lens state is an arrow $(1,1)\to\mathbb L(X)$, hence an element
$x\in X$. A lens costate is an arrow $\mathbb L(X)\to(1,1)$, hence a map
$q:X\to\mathsf{Pred}(X)$. The two arguments of $q(x)(x')$ have distinct
roles: $x$ is the reference and $x'$ is its candidate alternative.
\end{definition}

A costate must retain both arguments of an improvement comparison:
the policy currently evaluated and its candidate replacement. The map on
$X$ thus carries more information than a predicate that merely classifies
acceptable current outcomes.

\begin{construction}[Canonical relational and numerical feedback]
\label{con:cgt-canonical-coplay}
The cartesian lift of $Q\hookrightarrow Y$ along $f:X\to Y$ is
$f^*Q\hookrightarrow X$. Using the membership predicates of the power
objects to represent this operation, define
\[
 \mathbb L(f)=(f,f^-),\qquad f^-(x,Q)=f^*Q.
\]
For a value object $V$, put $\mathbb L_V(X)=(X,V^X)$ and define
$\mathbb L_V(f)$ with backward map $(x,k)\mapsto kf$.
Numerical values, their differences, and arbitrary comparison predicates
can consequently be transported along the same forward map. A value
presentation consists of the actual map $k:X\to V$ and its comparison
relation; the numerical map remains part of the evaluation data.
\end{construction}

\begin{theorem}[Functorial feedback and comparison costates]
\label{thm:cgt-coplay}
The assignments $\mathbb L$ and $\mathbb L_V$ are faithful functors
from $\mathbf{Set}$ to $\mathsf{Lens}(\mathbf{Set})$.
Transporting a comparison costate along $\mathbb L(f)$ gives
\begin{equation}\label{eq:cgt-costate-pullback}
 \bigl(\mathbb L(f)^*q_P\bigr)(x)(x')
 \quad\Longleftrightarrow\quad P(f(x),f(x')),
\end{equation}
so its uncurried comparison is $(f\times f)^*P$.
\end{theorem}
\begin{proof}
\leavevmode
\begin{enumerate}[label=\textup{(\arabic*)},leftmargin=*,itemsep=4pt,topsep=3pt]
  \item \textbf{Functoriality of predicate feedback.} For $f:X\to Y$ and $h:Y\to Z$, the composite backward map sends
$Q$ to $f^*(h^*Q)=(hf)^*Q$; identities give identity pullback. For numerical continuation interfaces the same calculation is
$(kh)f=k(hf)$. The forward component recovers the original map,
which proves faithfulness.
  \item \textbf{Pullback of comparison costates.} A costate of $(Y,\mathsf{Pred}(Y))$ sends
the reference $f(x)$ to $q_P(f(x))$. Pullback along $f$ then tests
$P(f(x),f(x'))$, proving the displayed equivalence.
\end{enumerate}
\end{proof}

\noindent The canonical coplay is functorial under composition of forward realizations.
\begin{center}
\begin{tikzcd}[column sep=large,row sep=large]
\mathsf{Pred}(Z)\arrow[r,"h^*"]\arrow[dr,"(hf)^*"']
 &\mathsf{Pred}(Y)\arrow[d,"f^*"]\\
 &\mathsf{Pred}(X)
\end{tikzcd}
\end{center}

\[
\begin{tikzcd}[column sep=large,row sep=large]
 X\times X \arrow[r,"f\times f"]
   \arrow[d,"{(x,x')\mapsto P(fx,fx')}"'] &
 Y\times Y \arrow[d,"P"]\\
 \mathsf{Pred}(1) \arrow[r,equal] & \mathsf{Pred}(1)
\end{tikzcd}
\]
Here a predicate is displayed as a truth-value map only to draw its
classification square; all subsequent constructions use subobjects and
pullback, as developed in Section~\ref{sec:egt-assertions}.

The forward carriers used below are hom-sets
$H_{A,B}=\mathcal C(A,B)$. Therefore a value of the play map is the
entire concurrent strategy of Section~\ref{sec:egt-interaction}. Applying $\mathbb L$ to an
ordinary map into such a hom-set constructs feedback on that strategy;
the linear interaction ports are still the original objects $A,B$.
For families $u,v:\Gamma\to H_{A,B}$, write
$\operatorname{Test}_P(u,v)=\langle u,v\rangle^*P$.
Section~\ref{sec:egt-assertions} derives the logical structure of this
testing operation.

Applied to the realization $g:\Sigma^G\to\mathcal C(A,B)$, this
construction equips the arena with $\mathbb L(g)$ and backward map
$g^-(\sigma,Q)=g^*Q$. The same operation is substitution along
$\widehat g$ into a constant predicate type in its state doctrine.
The equation $g=h\bar r$ for a reparametrisation therefore gives both
the forward equation and
$\mathbb L(g)=\mathbb L(h)\mathbb L(\bar r)$.

\subsection{Observable outcomes}

\begin{proposition}[Economic closure retains an outcome interface]
\label{prop:egt-scalars}
$\mathcal C(I,I)$ is a singleton. Therefore an evaluation depending only
on a fully hidden operational scalar is constant.
\end{proposition}
\begin{proof}
\leavevmode
\begin{enumerate}[label=\textup{(\arabic*)},leftmargin=*,itemsep=4pt,topsep=3pt]
  \item \textbf{Constancy of scalar evaluation.} This is the scalar calculation in Lemma~\ref{lem:egt-native-empty-obstruction}. Any function on that singleton is constant.
\end{enumerate}
\end{proof}

\begin{definition}[Economically closed realization]
\label{def:egt-economically-closed-realization}
An economically closed realization has type $I\to O$, with $O$ a retained
outcome interface. Economic closure supplies the decision inputs; it does
not hide $O$ before evaluation.
\end{definition}

\subsection{The category of operational contexts}

\begin{definition}[Operational contexts, comparison contexts, and environments]
\label{def:egt-economic-contexts}
An \textbf{operational context} from $(A,B)$ to $(A',B')$ is an
allowed wiring-induced function
\[
 T:\mathcal C(A,B)\longrightarrow\mathcal C(A',B'),
 \qquad T\in\mathsf{Ctx}(\mathcal C)((A,B),(A',B')),
\]
where the generators of $\mathsf{Ctx}(\mathcal C)$ are defined below.
A \textbf{comparison context} at $(A',B')$ is a tuple
$P\in\mathsf{Cmp}_N(A',B')$ of predicates on pairs of interactions.
An \textbf{environment} for $(A,B)$ is a typed pair $(T,P)$; its
induced comparison has type
$T^\sharp P\in\mathsf{Cmp}_N(A,B)$, with
\[
 (T^\sharp P)_i=(T\times T)^*P_i.
\]
A value-presented environment in particular may give
$a:I\to A$, $b:B\to O$, value maps
$u_i:\mathcal C(I,O)\to V_i$, and
$\prec_i\hookrightarrow V_i^2$. It induces
\[
 k_i:\mathcal C(A,B)\to V_i,\qquad
 k_i(f)=u_i(bfa),\qquad
 (k_i\times k_i)^*\prec_i\hookrightarrow\mathcal C(A,B)^2.
\]
A \textbf{state context} $s:W\to\Sigma^G$ instead indexes reference
profiles; it is distinct from these operational and comparison contexts.
\end{definition}
\begin{definition}[Operational context]
\label{def:egt-context}
An interface pair is an ordered pair $(A,B)$ of objects of $\mathcal C$.
The category $\mathsf{Ctx}(\mathcal C)$ has interface pairs as objects.
Its arrows $(A,B)\to(A',B')$ are functions
\[
 T:\mathcal C(A,B)\longrightarrow\mathcal C(A',B')
\]
generated under composition by identities, whiskering
$f\mapsto bfa$ for $a:A'\to A$, $b:B\to B'$, and tensor insertion
$f\mapsto f\otimes r$ or $f\mapsto r\otimes f$ for a fixed arrow $r$.
Generated functions are identified when they are equal. A parameterized
context is an ordinary family of such functions.
\end{definition}

\begin{construction}[Backward action on tests]
Put
\[
 \mathsf{Cmp}_N(A,B)=
 \prod_{i\in N}\mathsf{Pred}(\mathcal C(A,B)^2).
\]
For a context $T$, define
\[
 T^\sharp:\mathsf{Cmp}_N(A',B')\to\mathsf{Cmp}_N(A,B),
 \qquad (T^\sharp P)_i=(T\times T)^*P_i.
\]
For value maps $k_i:\mathcal C(A',B')\to V_i$, the corresponding action is
$k_i\mapsto k_iT$.
\end{construction}

\begin{theorem}[Contextual evaluation and coherence]
\label{thm:egt-context-action}
$\mathsf{Cmp}_N$ with this action is a contravariant functor
$\mathsf{Ctx}(\mathcal C)^{\mathrm{op}}\to\mathbf{Pos}$.
It preserves the pointwise predicate operations under reindexing, and
\[
 \operatorname{Test}_{T^\sharp P}(u,v)
 =\operatorname{Test}_{P}(Tu,Tv)
\]
for families $u,v:\Gamma\to\mathcal C(A,B)$.
The action respects every equation between operational contexts.
\end{theorem}
\begin{proof}
\leavevmode
\begin{enumerate}[label=\textup{(\arabic*)},leftmargin=*,itemsep=4pt,topsep=3pt]
  \item \textbf{Composition of contextual predicate pullback.} The costate action is the comparison part of the canonical lens
functor in \cref{thm:cgt-coplay}. In particular, for contexts $T,U$,
$((UT)\times(UT))^*=(T\times T)^*(U\times U)^*$. Testing follows by the same pullback identity:
$\langle u,v\rangle^*(T\times T)^*P
 =\langle Tu,Tv\rangle^*P$.
  \item \textbf{Functoriality and substitution.} All context equations are equalities of the generated functions, so
their backward actions agree by this functoriality.
\end{enumerate}
\end{proof}

\noindent The external test and its contextual pullback coincide.
\begin{center}
\begin{tikzcd}[column sep=large,row sep=large]
\Gamma\arrow[r,"{\langle u,v\rangle}"]\arrow[d,equal]
 &\mathcal C(A,B)^2\arrow[d,"T\times T"]\\
 \Gamma\arrow[r,"{\langle Tu,Tv\rangle}"']&\mathcal C(A',B')^2
\end{tikzcd}
\end{center}

\[
\begin{tikzcd}[column sep=large]
 \Gamma\arrow[r,"{(u,v)}"]\arrow[dr,"{(Tu,Tv)}"']&
 \mathcal C(A,B)^2\arrow[d,"T\times T"]\\
 &\mathcal C(A',B')^2
\end{tikzcd}
\qquad
\begin{tikzcd}[column sep=small]
 \mathsf{Cmp}_N(A,B)&\mathsf{Cmp}_N(A',B')\arrow[l,"T^\sharp"']
 &\mathsf{Cmp}_N(A'',B'')\arrow[l,"U^\sharp"']
 \arrow[ll,bend right=25,"(UT)^\sharp"']
\end{tikzcd}
\]

\begin{construction}[Evaluation supplied by an environment]
\label{con:egt-environment}
For $a:I\to A$, $b:B\to O$, and
$u_i:\mathcal C(I,O)\to V_i$, define
\[
 k_i(f)=u_i(bfa).
\]
More generally a comparison context is an element
$P\in\mathsf{Cmp}_N(A,B)$; it need not have a numerical representation.
If $\prec_i\subseteq V_i^2$, define
$P_i=(k_i\times k_i)^*\prec_i$.
For open arenas $G:A\to B$ and $H:B\to C$, the contexts obtained by fixing
$\tau\in\Sigma^H$ or $\sigma\in\Sigma^G$ are
\[
 T_\tau(f)=h_\tau f,
 \qquad S_\sigma(h)=hg_\sigma.
\]
Write $P^{G\mid\tau}=T_\tau^\sharp P$ and
$P^{H\mid\sigma}=S_\sigma^\sharp P$.
\end{construction}

\begin{proposition}[Context action for sequential and parallel composition]
\label{prop:egt-context-composition}
For fixed $h:B\to C$, $j:C\to D$,
$T_{jh}^\sharp=T_h^\sharp T_j^\sharp$.
Whiskering on opposite sides commutes. For fixed composable $r,r'$, tensor
insertion is compatible with composition by interchange. All these equations
hold also for parameterized contexts, pointwise in their parameters.
\end{proposition}
\begin{proof}
\leavevmode
\begin{enumerate}[label=\textup{(\arabic*)},leftmargin=*,itemsep=4pt,topsep=3pt]
  \item \textbf{Monoidal comparison.} The context functions satisfy $(jh)f=j(hf)$ and $b(fa)=(bf)a$. The tensor assertion follows from
$(h\otimes r')\circ(f\otimes r)=(hf)\otimes(r'r)$. Apply Theorem~\ref{thm:egt-context-action}.
\end{enumerate}
\end{proof}

The backward context action is the costate action of the canonical lenses
of Construction~\ref{con:cgt-canonical-coplay}. Its inputs and outputs are
complete concurrent strategies. Closing an evaluation context specifies
the comparison received by each agency interface; the retained output
arena continues to expose the economic outcome.

\subsection{Dependent assertions and semantic testing}
\label{sec:egt-assertions}

Comparisons of policies become assertions only after a reference
context has been chosen. The state-context fibration supplies ordinary
dependent objects for policies and witnesses and subterminal objects for
predicates. Comprehension restricts contexts to satisfying states;
quantifier adjoints express selection and existence of witnesses.
We now identify these operations with the standard dependent semantics
so that subsequent strategic and epistemic tests use the same logic.

\subsubsection{The predicate fragment and comprehension}

An allocation certificate and the assertion that an allocation is
feasible carry different information. An ordinary dependent family retains
the witnesses; its existential image records only their existence.
Comprehension selects the states satisfying a predicate, while its
subterminal fragment provides the truth conditions needed for agency.
The following constructions establish these interfaces without equating
witness-bearing types with propositions.

\begin{definition}[The assertion hyperdoctrine]
\label{def:fib-hyperdoctrine}
For an arena $G$, use the indexed assignment $\mathsf P_G$ of
Construction~\ref{con:fib-constants}. For $s:W\to\Sigma^G$, an assertion is a subobject of $W$;
entailment is inclusion. Its conjunction, disjunction, implication, and constants are
those of $\mathsf{Pred}(W)$, with $\neg P=(P\Rightarrow\bot)$.
An arrow $u:(W,s)\to(V,t)$ substitutes predicates by $u^*$.
The ambient notation $\mathsf{Pred}(X)$ permits the same calculation
before a reference context is chosen; constant types insert it into the
game's doctrine by Equation~\eqref{eq:fib-constant-test}.
\end{definition}

\begin{theorem}[Quantifiers, equality, and comprehension]
\label{prop:egt-quantifiers}\label{prop:egt-hyperdoctrine}
\label{thm:fib-comprehension}
The indexed assignment $\mathsf P_G$ is the subobject hyperdoctrine of
$\mathbf{Set}/\Sigma^G$. Every context map has adjoints
$\exists_u\dashv u^*\dashv\forall_u$, with
\[
 (\exists_uP)(v)\iff\exists w\,(u(w)=v\wedge P(w)),\qquad
 (\forall_uP)(v)\iff\forall w\,(u(w)=v\Rightarrow P(w)).
\]
Both quantifiers satisfy Beck--Chevalley for pullback squares. Moreover
$\exists_u(P\wedge u^*Q)=\exists_uP\wedge Q$.
Equality on a context $W\to\Sigma^G$ is its diagonal in
$W\times_{\Sigma^G}W$.

Comprehension of $P\hookrightarrow W$ is
$\{P\}=(P\hookrightarrow W\xrightarrow{s}\Sigma^G)$.
For the projection $p_G$ of the predicate fibration, the truth functor
$\top_G$ and comprehension satisfy
$p_G\dashv\top_G\dashv\{-\}$.
The same comprehension construction works in the total state-context
category and is stable under its cartesian substitutions.
\end{theorem}

\noindent Comprehension represents the predicates by inclusions of their satisfying contexts.
\begin{center}
\begin{tikzcd}[column sep=large,row sep=large]
\{P\}\arrow[r,hook,"m_P"]\arrow[d,"s|_P"']
 &W\arrow[d,"s"]\\
 \Sigma^G\arrow[r,equal]&\Sigma^G
\end{tikzcd}
\end{center}


\begin{proof}
\leavevmode
\begin{enumerate}[label=\textup{(\arabic*)},leftmargin=*,itemsep=4pt,topsep=3pt]
  \item \textbf{Predicate calculation.} A subobject of $(W,s)$ in the slice is exactly a subset of $W$ equipped
with the restricted map $s$. Slice pullbacks have the underlying
pullbacks of sets, proving the asserted identification of doctrines.
  \item \textbf{Base-change comparison.} The elementwise formulas give
\[
 \exists_uP\leq Q\iff P\leq u^*Q,
 \qquad u^*Q\leq P\iff Q\leq\forall_uP.
\]
In a pullback square
\[
\begin{tikzcd}
 W'\arrow[r,"a"]\arrow[d,"u'"']
  \arrow[dr,phantom,"\lrcorner",very near start]&W\arrow[d,"u"]\\
 V'\arrow[r,"b"']&V
\end{tikzcd}
\]
the fibre of $u'$ at $v'$ is canonically the fibre of $u$ at $b(v')$.
  \item \textbf{Frobenius and coproduct comparison.} The quantifier formulas therefore give
$b^*\exists_u=\exists_{u'}a^*$ and
$b^*\forall_u=\forall_{u'}a^*$. Frobenius follows because $Q(u(w))$ is constant on a fibre of $u$. The diagonal classifies equality in the slice product.
  \item \textbf{Predicate calculation.} An arrow from a predicate to truth is specified by its underlying context
map, proving $p_G\dashv\top_G$. An arrow from a true context to $P$
is a context map whose image lies in $P$; it factors uniquely through
$\{P\}$, proving $\top_G\dashv\{-\}$. In the total category the
same factorization keeps its arena reparametrisation unchanged. Inverse
image of a subset is its pullback, giving stability of comprehension.
\end{enumerate}
\end{proof}


\begin{construction}[Dependent types and predicate types]
\label{con:fib-dependent-types}
Over a state context $(W,s)$, a dependent type is a map $E\to W$;
its terms are sections. The slice
$(\mathbf{Set}/\Sigma^G)/(W,s)$ is canonically $\mathbf{Set}/W$.
These are the ordinary fibres of Construction~\ref{con:egt-dependent-fibres}.
For $u:W\to V$, their dependent sums and products have fibres
\[
 (\Sigma_uE)_v=\coprod_{u(w)=v}E_w,
 \qquad (\Pi_uE)_v=\prod_{u(w)=v}E_w.
\]
We write $\Sigma_u,\Pi_u$ for the ordinary adjoints
$\Sigma_u^{\mathrm c},\Pi_u^{\mathrm c}$ of
Theorem~\ref{thm:egt-linear-dependent-adjoints}; their adjunctions follow
from the same sum and product universal properties. The image of the dependent sum of a predicate gives
$\exists_u$; its dependent product gives $\forall_u$.
For $W\to\Sigma^G$, the power object has fibre
$\mathcal P_{\Sigma^G}(W)_\sigma=\mathcal P(W_\sigma)$, with fibrewise
membership. A subobject of a slice product is uniquely classified by
the map sending a parameter to its corresponding subset of $W_\sigma$.
The truth-value object is $\Sigma^G\times\{0,1\}\to\Sigma^G$.
Thus the construction also supplies higher-order predicate types.
\end{construction}

\begin{proposition}[Propositional reflection and recovery of quantifiers]
\label{prop:egt-propositional-recovery}
For $E\in\mathsf T(W)$ define $\|E\|$ by the image of its projection
$\coprod_w E_w\to W$. This is left adjoint to the inclusion of
subterminal families into ordinary families. It commutes with
substitution. For a subterminal $P$ over $W$ and $u:W\to V$,
\[
 \exists_uP=\|\Sigma_u^{\mathrm c}P\|,
 \qquad \forall_uP=\Pi_u^{\mathrm c}P.
\]
Thus the existing predicate doctrine and its quantifiers are recovered
inside the constructed dependent semantics.
\end{proposition}
\begin{proof}
\leavevmode
\begin{enumerate}[label=\textup{(\arabic*)},leftmargin=*,itemsep=4pt,topsep=3pt]
  \item \textbf{Characterization of propositional reflection.} The fibre $\|E\|_w$ is a singleton exactly when $E_w$ is inhabited. A map from $E_w$ to a subterminal $P_w$ exists exactly when this
inhabitation implies inhabitation of $P_w$, and is then unique.
  \item \textbf{Naturality and quantified reflection.} This proves the adjunction without choosing an element of $E_w$. Reindexing retains the same fibres, proving substitution stability. The coproduct is inhabited exactly when some summand is inhabited.
  \item \textbf{Monoidal comparison.} A product of subterminal sets is a singleton exactly when every factor
is inhabited: its element is the uniquely specified family of their
unique elements. Otherwise it is empty.
  \item \textbf{Quantifiers on empty and inhabited fibres.} This proves the quantifier
formulas, including empty indexing fibres, without a choice of witnesses
from arbitrary inhabited types.
\end{enumerate}
\end{proof}

\begin{proposition}[Logical change of base]
\label{prop:fib-logical-basechange}
The change of game in Proposition~\ref{prop:fib-basechange} preserves
the constructed finite-limit, comprehension, and predicate-type
structure up to its canonical comparisons. Quantifiers commute with
this change of base along the associated pullback squares. These
comparisons satisfy substitution coherence.
\end{proposition}
\begin{proof}
\leavevmode
\begin{enumerate}[label=\textup{(\arabic*)},leftmargin=*,itemsep=4pt,topsep=3pt]
  \item \textbf{Monoidal comparison.} In family notation a context has fibres $W_\eta$ and a dependent type
has fibres $E_{\eta,w}$. Pullback along $\bar r$ replaces $\eta$ by
$\bar r\sigma$ in these formulas. Products, equalities, subsets, and
power objects are consequently the identical constructions on identical
fibres.
  \item \textbf{Monoidal comparison: completion.} Sums and products indexed by a dependent context map range over
the same sets after this substitution. This proves their comparisons
and the quantifier equations; repeated substitution of the index proves
coherence. The set-theoretic constructions implement the fibred logical
semantics of \citet{Jacobs1999}.
\end{enumerate}
\end{proof}

The doctrine lives over state contexts, not over a Cartesian structure
on the category $\mathcal A_{A,B}$ of agent presentations itself. The latter indexes the
contexts and their changes of base; operational contexts in $\mathsf{Ctx}(\mathcal C)$ separately act
on hom-set values as in Section~\ref{sec:egt-players}. Keeping these roles
apart is necessary for treating owned reparametrizations and logical
substitutions correctly.

\subsubsection{Algebras and coalgebras of predicate adjunctions}

\begin{definition}[Posetal monads and comonads]
\label{def:egt-posetal-monads-and-comonads}
On a poset, a monad is a monotone endomap $M$ with
$\operatorname{id}\leq M$ and $M^2\leq M$; a comonad is a monotone
endomap $K$ with $K\leq\operatorname{id}$ and $K\leq K^2$.
These inequalities imply idempotence. Their algebra and coalgebra posets
are respectively the full subposets $MP=P$ and $KP=P$.
\end{definition}

\begin{theorem}[The two universal approximations to invariant assertions]
\label{thm:egt-invariant-adjunctions}
For a surjection $q:X\twoheadrightarrow Q$, let $\mathsf{Inv}(q)$ be
the full subposet of predicates constant on the fibers of $q$, with
inclusion $j:\mathsf{Inv}(q)\hookrightarrow\mathsf{Pred}(X)$. Then
\[
 M_q=q^*\exists_q\dashv j\dashv K_q=q^*\forall_q,
 \qquad
 \operatorname{Alg}(M_q)\cong\mathsf{Inv}(q)
 \cong\operatorname{Coalg}(K_q)\cong\mathsf{Pred}(Q).
\]
In the adjunction display, $M_q$ and $K_q$ are corestricted to
$\mathsf{Inv}(q)$; as endomaps they are the resulting idempotent
monad and comonad.
\end{theorem}
\begin{proof}
\leavevmode
\begin{enumerate}[label=\textup{(\arabic*)},leftmargin=*,itemsep=4pt,topsep=3pt]
  \item \textbf{Verification of the universal property.} Surjectivity gives $\exists_qq^*=\forall_qq^*=\operatorname{id}$. Thus $M_q$ and $K_q$ have invariant values, are idempotent, and the
adjunction units and counits give $P\leq M_qP$ and $K_qP\leq P$. If $Q'=q^*\bar Q$ is invariant, then
\[
 M_qP\leq Q'\iff P\leq Q',\qquad
 Q'\leq K_qP\iff Q'\leq P
\]
by the two quantifier adjunctions.
  \item \textbf{Predicate calculation.} These are the claimed universal
properties. Their fixed predicates are precisely those in the image of
$q^*$, and effective predicate descent identifies this image with
$\mathsf{Pred}(Q)$.
\end{enumerate}
\end{proof}

\noindent Invariant predicates descend uniquely along the observation quotient.
\begin{center}
\begin{tikzcd}[column sep=large,row sep=large]
\mathsf{Pred}(Q)\arrow[r,"q^*","\cong"']\arrow[d,equal]
 &\mathsf{Inv}(q)\arrow[d,hook,"j"]\\
 \mathsf{Pred}(Q)\arrow[r,"q^*"']
 &\mathsf{Pred}(X)
\end{tikzcd}
\end{center}

\[
\begin{tikzcd}[column sep=huge]
 \mathsf{Pred}(X)\arrow[r,shift left=1.2ex,"M_q"]
 \arrow[r,bend right=40,"K_q"']&
 \mathsf{Inv}(q)\arrow[l,shift left=1.2ex,"j"]
 \arrow[r,"\simeq"]&\mathsf{Pred}(Q)
\end{tikzcd}
\]

\begin{remark}[One construction, several economic meanings]
Observation quotients give epistemic invariant assertions. Owner-coordinate
quotients give assertions invariant under the corresponding deviations.
Quotients of redundant presentations give representation-independent
assertions. The adjunction is the same; the constructed quotient specifies
its economic meaning.
\end{remark}

\subsubsection{Realization terms and their assertion predicates}
\label{sec:egt-semantic-testing}

\begin{construction}[Semantic testing]
For $t:\Gamma\to X$ and $P\hookrightarrow X$, put
$\operatorname{Test}_P(t)=t^*P$.

In a state context $s:\Gamma\to\Sigma^G$, its typed expression is
$(s,t)^*\Delta_GP=t^*P$, by Equation~\eqref{eq:fib-constant-test}.
Thus the same test is an internal assertion on the reference context. For a finite tuple of indexed arrows use
$X=\prod_k\mathcal C(A_k,B_k)$ and
\[
 \operatorname{Test}_P(t_1,\ldots,t_n)
 =\langle t_1,\ldots,t_n\rangle^*P.
\]
\end{construction}

\begin{theorem}[Classification of semantic tests]
\label{thm:egt-semantic-test-yoneda}
Predicates on $X$ correspond bijectively to assignments
$\tau_\Gamma(t)\in\mathsf{Pred}(\Gamma)$ for $t:\Gamma\to X$
that are natural under substitution.
The inverse sends $\tau$ to $\tau_X(\operatorname{id}_X)$.
In particular,
\begin{equation}\label{eq:egt-semantic-test-reindexing}
 f^*\operatorname{Test}_P(t)=\operatorname{Test}_P(tf).
\end{equation}
For $F:X\to Y$ one has
$\operatorname{Test}_{F^*Q}(t)=\operatorname{Test}_Q(Ft)$.
\end{theorem}
\begin{proof}
\leavevmode
\begin{enumerate}[label=\textup{(\arabic*)},leftmargin=*,itemsep=4pt,topsep=3pt]
  \item \textbf{Naturality.} Naturality at $t:\Gamma\to X$ forces
$\tau_\Gamma(t)=t^*\tau_X(\operatorname{id}_X)$. Conversely pullback
is functorial, so every predicate gives such an assignment. These operations
are inverse. This is the Yoneda argument in its elementwise form.
\end{enumerate}
\end{proof}

\noindent Naturality of testing is the pullback square of an interpreted predicate.
\begin{center}
\begin{tikzcd}[column sep=large,row sep=large]
\Delta\arrow[r,"f"]\arrow[d,"tf"']
 &\Gamma\arrow[d,"t"]\\
 X\arrow[r,equal]&X
\end{tikzcd}
\end{center}

\[
\begin{tikzcd}
 \operatorname{Test}_P(t)\arrow[r]\arrow[d,hook]&P\arrow[d,hook]\\
 \Gamma\arrow[r,"t"']&X
\end{tikzcd}
\qquad
\begin{tikzcd}
 \Delta\arrow[r,"f"]\arrow[dr,"tf"']&\Gamma\arrow[d,"t"]\\
 &X
\end{tikzcd}
\]

The classification concerns tests of the specified semantic values.
Ownership, observations, and legality supply additional data; it does not
assert that every assertion about them depends only on the realized arrow.

\subsubsection{State contexts, legal comprehension, and operational substitution}

\begin{definition}[State realization]
\label{def:egt-state-realization}
A state realization is an object $(W,s)$ of $\mathsf{St}(G)$.
Its operational family is $gs:W\to\mathcal C(A,B)$, equivalently
a morphism $yA\to yB$ in $\mathsf L(W)$. Its realization section
$(s,gs)$ takes values in the ordinary family
$\mathsf{Val}_W[yA,yB]_D$, identified with the constant concurrent-arrow
family by Theorem~\ref{thm:egt-day-preservation}.
For a legal predicate with comprehension $m:L\hookrightarrow\Sigma^G$,
a legal state realization is a context over that comprehension,
$s:W\to L$, with operational family $gms$.
For $f:W'\to W$, substitution is $s\mapsto sf$ and $P\mapsto f^*P$.
\end{definition}

\begin{proposition}[Synchronized substitution]
For every semantic predicate $Q$,
\[
f^*\operatorname{Test}_Q(gms)=\operatorname{Test}_Q(gmsf).
\]
The intermediate state space need not contain representatives of every
available alternative policy.
\end{proposition}
\begin{proof}
\leavevmode
\begin{enumerate}[label=\textup{(\arabic*)},leftmargin=*,itemsep=4pt,topsep=3pt]
  \item \textbf{Representation and Yoneda decoding.} The equality is Theorem~\ref{thm:egt-semantic-test-yoneda}. The definitions require only the displayed map $s$ and impose no
surjectivity condition on it.
\end{enumerate}
\end{proof}

\begin{theorem}[Fibred compatibility of dependent realization and assertions]
\label{thm:fib-linear-logical-compatibility}
Over $\mathsf{St}_{A,B}$ the indexed ordinary and linear categories of
Construction~\ref{con:egt-state-dependent-fibres}, their comprehension,
monoidal adjunction, and dependent adjoints share coherent substitution.
The assertion fibration is their subterminal ordinary fragment. The
original indexed operational category embeds fully faithfully in the
linear fibration, preserving its concurrent structure.

If $t:W\to\prod_j\mathcal C(A_j,B_j)$ is a tuple of realized families,
$m$ is composition, tensor, or another exhibited operation on these
arrow carriers, and $Q$ is a predicate on its target, then
\[
 u^*\operatorname{Test}_Q(mt)
 =\operatorname{Test}_Q(m(tu))
 =\operatorname{Test}_{m^*Q}(tu)
\]
for every $u:V\to W$. The same equations hold for comparison predicates
on paired carriers. For dependent boundaries, replace the constant
carriers by $\mathsf{Val}[A,B]_D$ and use their fibrewise operations.
\end{theorem}
\begin{proof}
\leavevmode
\begin{enumerate}[label=\textup{(\arabic*)},leftmargin=*,itemsep=4pt,topsep=3pt]
  \item \textbf{Representation and Yoneda decoding.} The categorical assertions follow by change of base of
Theorem~\ref{thm:egt-dependent-model} along $\operatorname{Dom}$. Proposition~\ref{prop:egt-propositional-recovery} identifies the assertion
fibration, and pointwise Yoneda proves the embedding assertion. Each
cartesian lift reindexes the same ordinary parameter, so all the
comparison maps are those already proved coherent in the model theorem.
  \item \textbf{Secured operational interaction.} For the displayed equalities, $(mt)u=m(tu)$ and
$(m(tu))^*Q=(tu)^*m^*Q$ are inverse-image functoriality. The paired and
dependent cases are the identical calculation on the corresponding
fibres. 
\end{enumerate}
\end{proof}

\subsubsection{Effective descent of terms and predicates}

Observation descent gives a concrete example of this logical
architecture. If $q:H\to O$ identifies histories an agent cannot
distinguish, an implementable response $f:H\to A$ must be constant on
the fibres of $q$ and hence factor through $O$. It cannot prescribe different
actions at histories producing the same signal. The quotient thus has an
operational meaning before knowledge modalities are introduced, and its
dependent extension will retain complete families of such responses.

\begin{proposition}[Descent and finite products of quotients]
\label{prop:egt-kernel-descent}
\label{prop:egt-predicate-descent}
\label{prop:egt-finite-surjections}
Let $q:X\twoheadrightarrow Y$ have kernel pair
$p_1,p_2:X\times_YX\rightrightarrows X$.
A map $F:X\to Z$ descends uniquely to $Y$ exactly when $Fp_1=Fp_2$.
A predicate $P\hookrightarrow X$ descends uniquely exactly when
$p_1^*P=p_2^*P$; its descent is
\[
 \bar P(y)\iff\exists x\;(q(x)=y\wedge P(x)).
\]
Surjections are effective regular epimorphisms and are stable under
pullbacks and finite products, including the nullary product.
\end{proposition}
\begin{proof}
\leavevmode
\begin{enumerate}[label=\textup{(\arabic*)},leftmargin=*,itemsep=4pt,topsep=3pt]
  \item \textbf{Predicate calculation.} These are the effective-descent properties of sets
\citep[Chapters 4--5]{Jacobs1999}. Explicitly, the relation
$\bar F(y)=z\iff\exists x\,(q(x)=y\wedge F(x)=z)$ is total and
single-valued precisely when $F$ is fiber-constant. The displayed formula
for $\bar P$ pulls back to $P$ precisely when $P$ is fiber-invariant.
  \item \textbf{Verification of the universal property.} Surjectivity gives uniqueness in both cases. A pullback has the same fibers
as the original surjection. Products of two surjections are composites of
two such pullbacks; finite induction includes the identity on $1$.
\end{enumerate}
\end{proof}

\[
\begin{tikzcd}
 X\times_YX \arrow[r,shift left=.55ex,"p_1"]
 \arrow[r,shift right=.55ex,"p_2"'] &X\arrow[r,two heads,"q"]
 \arrow[dr,"F"'] &Y\arrow[d,dashed,"\bar F"]\\
 &&Z
\end{tikzcd}
\]

\subsubsection{Substituting the reference state}

Substitution changes the reference of an incentive question, not its
set of potential deviations. If $\Sigma_i=\{a,b\}$ and a context selects
$a$, then $b$ remains an alternative after reference pullback.
Restricting the policy carrier to $\{a\}$ instead removes that deviation.
The two operations have similar maps of reference states but different
counterfactual meanings, which the following span calculation separates.

\begin{construction}[Reference pullback]
For $X\xleftarrow pR\xrightarrow qY$ and $f:U\to X$, put
$R_f=U\times_XR$, with legs $p_f:R_f\to U$ and
$q_f=q\operatorname{pr}_R:R_f\to Y$.
For $s:W\to\Sigma$, the available $C$-alternatives are
$\mathsf{Alt}_C(s)=W\times_{\Sigma,p_1}D_C$.
Their paired realization sends $(w,\tau)$ to $(gs(w),g\tau)$.
\end{construction}

\begin{proposition}[Substitution of relational assertions]
\label{prop:egt-reference-bc}
$f^*[R]P=[R_f]P$ and
$f^*\langle R\rangle P=\langle R_f\rangle P$.
More generally, for a predicate $Q$ on the relation itself,
$f^*\forall_pQ=\forall_{p_f}\operatorname{pr}_R^*Q$, and similarly for
$\exists$. No target alternative has to lie in the image of $f$.
\end{proposition}
\begin{proof}
\leavevmode
\begin{enumerate}[label=\textup{(\arabic*)},leftmargin=*,itemsep=4pt,topsep=3pt]
  \item \textbf{Base-change comparison.} Apply both Beck--Chevalley equations of
Theorem~\ref{thm:fib-comprehension} to the displayed cartesian lift. Proposition~\ref{prop:fib-span-typing} types its target as a term into
$\Delta_GY$, so pulling back the reference context retains every target
alternative. This proves both modal identities and their quantified form.
\end{enumerate}
\end{proof}

\[
\begin{tikzcd}
 R_f\arrow[r,"\operatorname{pr}_R"]\arrow[d,"p_f"']
 \arrow[dr,phantom,"\lrcorner",very near start]&R\arrow[d,"p"]\arrow[r,"q"]&Y\\
 U\arrow[r,"f"']&X&
\end{tikzcd}
\]

\subsection{Selection and logical equilibrium}
\label{sec:egt-equilibrium}

The agent's canonical selection relation becomes economically meaningful
once a comparison costate has been supplied. Pullback transports joint
comparison feedback to each complete owner policy; universal
quantification tests whether an improving alternative exists. Pullback
along the diagonal then identifies best responses at the current profile
with equilibrium. We retain the off-diagonal best-response relation and
the paired concurrent realizations rather than only their equilibrium
truth values.

\subsubsection{Connecting agent interfaces to joint comparison feedback}

The complete owner interface, constructed in
Section~\ref{sec:egt-agents}, now receives the comparison on the wired
interaction. Fixing outsider coordinates turns that joint comparison into
a costate on the owner's own policy set. The selection relation then
classifies best responses, and its diagonal classifies equilibrium; the
same calculation is later restricted to legal alternatives.

\begin{construction}[Connecting an agent to joint feedback]
\label{con:cgt-player-feedback}
Fix an open arena $G$, comparison context $P=(P_i)_i$, and reference
profile $\sigma$. For owner $i$ put
\[
 j_{i,\sigma}:\Sigma_i\to\Sigma,
 \qquad j_{i,\sigma}(a)=(\sigma_{-i},a).
\]
The agent's costate is the costate transported through
$\mathbb L(gj_{i,\sigma})$:
\begin{equation}\label{eq:cgt-player-costate}
 q_{i,\sigma}(a)(b)
 \quad\Longleftrightarrow\quad
 P_i(g_{(\sigma_{-i},a)},g_{(\sigma_{-i},b)}).
\end{equation}
Define the best-response relation by
\[
 \mathsf{BR}_G(P)(\sigma,\tau)
 \quad\Longleftrightarrow\quad
 \bigwedge_i\varepsilon_{\Sigma_i}(\tau_i,q_{i,\sigma}).
\]
For arbitrary comparisons this means absence of an improving alternative;
for classical real-valued maximization it is the best-response condition of ordinary game theory.
\end{construction}

In $\mathsf{BR}_G(P)(\sigma,\tau)$, the first profile $\sigma$ determines the
outsiders against whom responses are evaluated. The second supplies the
candidate choices $\tau_i$, each compared with alternatives for that
same $\tau_i$ while
$\sigma_{-i}$ remains fixed. Evaluating against $\tau_{-i}$ would instead
change the reference game. Only on the diagonal $\tau=\sigma$ does the
relation say that each current choice is a best response to all the other
current choices.

\begin{proposition}[Selection and best response as internal predicates]
\label{prop:fib-selection-internal}
For $T_S=S\times(S\to\mathsf{Pred}(S))$, let
$B_S\hookrightarrow T_S\times S$ be the predicate
$B_S(a,q,b)\iff q(a)(b)$ and let $\rho:T_S\times S\to T_S$.
Then the agent's selection predicate is
\[
 \varepsilon_S=\forall_\rho\neg B_S\in\mathsf{Pred}(T_S).
\]
For an evaluated arena $G$, define
\[
 \kappa_i:\Sigma^G\times\Sigma^G\to T_{\Sigma_i^G},
 \qquad \kappa_i(\sigma,\tau)=(\tau_i,q_{i,\sigma}).
\]
Its complete best-response relation and equilibrium predicate are
\[
 \mathsf{BR}_G(P)=\bigwedge_i\kappa_i^*\varepsilon_{\Sigma_i^G},
 \qquad
 \mathsf{NE}_G=\Delta_{\Sigma^G}^{\,*}\mathsf{BR}_G(P).
\]
All these expressions are terms and predicates of $\mathsf P_G$ using
the constant types of Construction~\ref{con:fib-constants}.
\end{proposition}
\begin{proof}
\leavevmode
\begin{enumerate}[label=\textup{(\arabic*)},leftmargin=*,itemsep=4pt,topsep=3pt]
  \item \textbf{Strategic comparison.} The fibre of $\rho$ at $(a,q)$ is $S$, so its universal quantifier
is exactly $\forall b\,\neg q(a)(b)$. Substitution along $\kappa_i$
gives the agent condition defining best response, and substitution
along the diagonal makes its candidate response equal to the reference. The constant-type interpretation preserves each pullback by
Equation~\eqref{eq:fib-constant-test}.
  \item \textbf{Strategic comparison: completion.} This retains the off-diagonal
relation; passage to equilibrium is the displayed additional pullback.
\end{enumerate}
\end{proof}

The selection-relation approach of
\citet{CapucciEtAl2022Cybernetics,CapucciEtAl2022Agency} organizes
selection covariantly over lenses, with the Nash product giving a
monoidal opfibration. Here the corresponding predicates live in a
contravariant assertion fibration. Connecting feedback to selection is
the specified substitution, not an identification of those two
variances. A repeated owner retains one complete policy carrier before
this test is performed.

\begin{theorem}[Selection, counterfactual feedback, and logical equilibrium]
\label{thm:cgt-selection-equilibrium}
For every comparison context, the following predicates on $\Sigma$ agree:
\begin{enumerate}[label=\textup{(\roman*)},leftmargin=*]
\item the diagonal $\mathsf{BR}_G(P)(\sigma,\sigma)$;
\item satisfaction of every agent selection relation after connecting
the agent interfaces to the arena;
\item $\displaystyle\bigwedge_i
 \forall_{p_1:D_i\to\Sigma}\neg(\delta_i^G)^*P_i$.
\end{enumerate}
The equality retains the off-diagonal best-response relation and all paired
deviation realizations as separate, constructed semantic data.
\end{theorem}
\begin{proof}
\leavevmode
\begin{enumerate}[label=\textup{(\arabic*)},leftmargin=*,itemsep=4pt,topsep=3pt]
  \item \textbf{Policy-level comparison.} The first two assertions both expand to
$\bigwedge_i\forall b\in\Sigma_i\,
 \neg P_i(g_\sigma,g_{(\sigma_{-i},b)})$. The fiber of $p_1:D_i\to\Sigma$ at $\sigma$ is precisely that set of
alternative policies, and $\delta_i^G$ realizes its two endpoints. The adjunction $p_1^*\dashv\forall_{p_1}$ therefore proves equality
with the third assertion.
  \item \textbf{Typing of paired deviation tests.} The paired realization is typed into a constant
comparison type by Proposition~\ref{prop:fib-span-typing}, so this is
an equality in the assertion fibre over $1_G$. The proof uses no order
property of $P_i$. If $P_i(f,f')$ is $k_i(f)<k_i(f')$ for classical
real values, its negation is $k_i(f)\geq k_i(f')$, giving the argmax interpretation with every tie retained.
\end{enumerate}
\end{proof}

\noindent The two projections of a unilateral deviation determine the pair of realized interactions.
\begin{center}
\begin{tikzcd}[column sep=large,row sep=large]
&D_i\arrow[dl,"p_1"']\arrow[dr,"p_2"]&\\
 \Sigma\arrow[d,"g"']&&\Sigma\arrow[d,"g"]\\
 \mathcal C(A,B)&&\mathcal C(A,B)
\end{tikzcd}
\end{center}

\[
\begin{tikzcd}[column sep=large,row sep=large]
 \Sigma_i^2 \arrow[r,"j_{i,\sigma}\times j_{i,\sigma}"]
   \arrow[dr,"{(gj_{i,\sigma})^2}"'] &
 \Sigma^2 \arrow[d,"g\times g"]\\
 &\mathcal C(A,B)^2
\end{tikzcd}
\]

Independent owners recover the relational Nash product: every owner
receives the comparison obtained by fixing the other owners' policies.
A repeated owner has one larger carrier $\Sigma_i$ containing all its
decisions, and therefore must be tested against changes to that entire
carrier. This is the distinction on which equilibrium under composition
depends.

\begin{construction}[Maximum of a finite numerical family]
\label{con:egt-finite-max}
In the classical real-valued instance, a nonempty finite list
$(a_1,\ldots,a_m)$ has a maximum constructed by
$b_1=a_1$ and $b_{k+1}=\max(b_k,a_{k+1})$. To retain a maximizing
record from a displayed finite ordered domain, retain its first record
at a tie. Induction proves that $b_k$ is the maximum of the first $k$
values and that the retained record has that value. Thus this construction
provides both the maximum value and a specified first maximizer.
\end{construction}

\begin{definition}[Gain on a deviation]
\label{def:egt-deviation-gain}
For a numerical utility $u_i:\Sigma\to\mathbb R$ and a deviation
$(\sigma,\tau)\in D_i$, put
$\Delta_i u(\sigma,\tau)=u_i(\tau)-u_i(\sigma)$.
Strict profitable deviation is the inverse image of $\mathbb R_{>0}$.
For composable deviations of the same owner these gains add, and reversing
a deviation negates its gain, by cancellation of the intermediate utility.
\end{definition}

\subsubsection{Comparison contexts and equilibrium}

\begin{definition}[Preferences and value presentations]
\label{def:egt-preferences-and-value-presentations}
A strict comparison records an improvement, while a weak comparison
allows equally valued alternatives. For a preorder, two alternatives are
\textbf{indifferent} when each is weakly preferred to the other and
\textbf{incomparable} when neither is. Both fail to be strict improvements
over one another, for different reasons. The general comparison context can be an arbitrary relation; preference-order terminology applies
when the indicated preorder or numerical presentation is supplied.
The second argument of $\mathsf{Better}_i$ is always the candidate
improvement.

A comparison context for $G:A\to B$ is
$\mathsf{Better}\in\mathsf{Cmp}_N(A,B)$, with the second argument denoting
the improved interaction. No order law is part of this type. A value
presentation is a family $k_i:\mathcal C(A,B)\to V_i$ with relations
$\prec_i\subseteq V_i^2$ such that
$\mathsf{Better}_i=(k_i\times k_i)^*\prec_i$.
Write $v_i^G=k_ig$. When $V_i$ is preordered, its strict part is
$x\prec_i y\iff x\preceq_i y\wedge\neg(y\preceq_i x)$.
Every comparison has the tautological relational presentation
$V_i=\mathcal C(A,B)$, $k_i=\operatorname{id}$.
For an arbitrary comparison, the selection predicate expresses \emph{stability against the specified relation}; it has the Nash
interpretation when that relation is strict improvement. In particular,
reflexivity or irreflexivity is not inferred from the bare type
$\mathsf{Cmp}_N(A,B)$.
\end{definition}

\begin{construction}[Profitable deviations and equilibrium]
On $D_i$ put $\mathsf{Prof}_i=(\delta_i^G)^*\mathsf{Better}_i$ and define
\begin{equation}\label{eq:egt-ne-doctrine}
 \mathsf{NE}_G=\bigwedge_{i\in N}\forall_{p_1:D_i\to\Sigma}
 \neg\mathsf{Prof}_i.
\end{equation}
Write $\operatorname{NE}(G,\mathsf{Better})$ for this predicate, and
$\operatorname{NE}(G,k)$ when the value presentation is understood.
\end{construction}

\begin{theorem}[Derivation of pure strategic equilibrium]
For every profile $\sigma$,
\begin{equation}\label{eq:egt-ne}
 \mathsf{NE}_G(\sigma)
 \iff\bigwedge_i\forall a_i\in\Sigma_i\;
 \neg\mathsf{Better}_i(g_\sigma,g_{(\sigma_{-i},a_i)}).
\end{equation}
For classical real-valued utilities this is exactly
$v_i(\sigma)\geq v_i(\sigma_{-i},a_i)$ for all $i,a_i$.
For an incomplete preorder it expresses absence of a strictly better
alternative, without asserting comparability of all alternatives.
\end{theorem}
\begin{proof}
\leavevmode
\begin{enumerate}[label=\textup{(\arabic*)},leftmargin=*,itemsep=4pt,topsep=3pt]
  \item \textbf{Strategic comparison.} The fiber of $p_1$ at $\sigma$ is the set of profiles agreeing with
$\sigma$ outside $i$. Use the coordinate bijection
\eqref{eq:egt-deviation-coordinates}, then expand the quantifier and test. In a classical linear order, absence of a strict improvement is the
displayed weak inequality.
\end{enumerate}
\end{proof}

\begin{construction}[The ordinary dependent object of equilibrium proofs]
\label{con:egt-equilibrium-type}
Regard each $\mathsf{Prof}_i(\sigma,\tau)$ as its subterminal ordinary
family, and write $D_i(\sigma)$ for the fibre of the deviation projection.
Construct in $\mathsf T(\Sigma)$
\[
 \mathsf{Eqm}_G(\sigma)
 =\prod_{i\in N}\prod_{\tau\in D_i(\sigma)}
       \bigl(\mathsf{Prof}_i(\sigma,\tau)\to0\bigr).
\]
Each factor and therefore the product is subterminal. Its fibre is
inhabited precisely when there is no profitable deviation of any owner,
so it is the subterminal object classifying $\mathsf{NE}_G$.
This follows from Proposition~\ref{prop:egt-propositional-recovery},
applied to the displayed dependent products. Reference substitution
pulls back the complete deviation fibres, as in
Theorem~\ref{thm:fib-linear-logical-compatibility}.
\end{construction}

\begin{construction}[Improvement relation]
Put
$\mathsf{Cmp}_i^G=(g\times g)^*\mathsf{Better}_i$ on $\Sigma^2$ and
$\mathsf{PDev}_i=D_i\cap\mathsf{Cmp}_i^G$.
\end{construction}

\begin{proposition}[Equilibrium as modal safety]
\label{prop:egt-ne-modal-safety}
\begin{equation}\label{eq:egt-ne-modal-safety}
 \mathsf{NE}_G=\bigwedge_i\neg\langle\mathsf{PDev}_i\rangle\top.
\end{equation}
For a state map $s:W\to\Sigma$, testing and quantifying the reference
pullback $\mathsf{Alt}_i(s)$ yields exactly $s^*\mathsf{NE}_G$.
\end{proposition}
\begin{proof}
\leavevmode
\begin{enumerate}[label=\textup{(\arabic*)},leftmargin=*,itemsep=4pt,topsep=3pt]
  \item \textbf{Strategic comparison.} An existential successor of $\mathsf{PDev}_i$ is precisely a profitable
$i$-deviation. The state-indexed assertion follows from
Proposition~\ref{prop:egt-reference-bc} with the predicate
$\neg\mathsf{Prof}_i$ on $D_i$.
\end{enumerate}
\end{proof}

\begin{proposition}[Preferences are not determined by interaction]
No predicate depending only on the concurrent open-arena realization defines
equilibrium for all comparison contexts.
\end{proposition}
\begin{proof}
\leavevmode
\begin{enumerate}[label=\textup{(\arabic*)},leftmargin=*,itemsep=4pt,topsep=3pt]
  \item \textbf{Strategic comparison.} Use one owner and two distinct value-state realizations. Rank the first
strictly above the second, and then reverse the ranking. The two equilibrium
sets are different singletons, while the open-arena realization is unchanged.
\end{enumerate}
\end{proof}

\subsubsection{Contextual decisions}

\begin{proposition}[The open-decision calculation]
\label{prop:egt-open-decision}
For a decision presentation $\mathsf{Dec}_i(X,Y)$, fix an observed
input $x\in X$ and a continuation $k:Y\to\mathbb R$. Define
the comparison by
\[
 P_{x,k}(f,f')\iff\exists y,y'\;
 \bigl(f\mathsf{pt}(x)=\mathsf{pt}(y)\wedge
 f'\mathsf{pt}(x)=\mathsf{pt}(y')\wedge k(y)<k(y')\bigr).
\]
Its equilibrium condition is
\[
 \pi\in\operatorname{NE}(x,k)
 \iff\pi(x)\in\operatorname*{arg\,max}_{y\in Y}k(y).
\]
\end{proposition}
\begin{proof}
\leavevmode
\begin{enumerate}[label=\textup{(\arabic*)},leftmargin=*,itemsep=4pt,topsep=3pt]
  \item \textbf{Policy-level comparison.} The response-state equation gives
$\widehat\pi\mathsf{pt}(x)=\mathsf{pt}(\pi(x))$. Consequently the comparison on response policies is
$k(\pi(x))<k(\pi'(x))$. Every action $y$ occurs at $x$ under the constant
policy with value $y$.
  \item \textbf{Policy-level comparison: completion.} Quantifying over policies is therefore equivalent
to quantifying over actions. This recovers the contextual utility-maximizing
decision of \citet{GhaniEtAl2018}.
\end{enumerate}
\end{proof}

\subsubsection{Exact recovery of pure normal-form games}
\label{sec:egt-normal}

\begin{definition}[Normal-form games]
\label{def:egt-normal-form-game}
A \textbf{normal-form game} consists of an agent set $N$, a policy set
$\Sigma_i$ for each agent, and payoff maps
$u_i:\prod_{j\in N}\Sigma_j\to\mathbb R$. It specifies a payoff for each
joint policy profile without displaying the internal order of decisions.
A policy in this representation may itself be a contingent plan from a
larger protocol.
\end{definition}

The finite normal-form specialization checks the economic meaning of
the construction. Independent action selections give singleton realized
outcome tuples. The Hoare preorder extends the numerical comparison to
other strategy states; on the singleton outcomes generated by the game it
agrees exactly with the ordinary payoff order. Hence the general
selection predicate can be compared directly with pure Nash equilibrium.

\begin{construction}[Realizing a finite normal-form game]
Let $A_i$ be finite nonempty action sets and
$u_i:\prod_j A_j\to\mathbb R$ utilities. Player $i$ has a singly owned open game
$G_i:I\to D(A_i)$ whose parameter $a_i$ realizes the one-event strategy
$\mathsf{pt}_{A_i}(a_i)$ selecting $a_i$. Form
\[
 G=\bigotimes_{i\in N}G_i:I\longrightarrow O,
 \qquad O=\bigotimes_{i\in N}D(A_i).
\]
At a profile $a$, the realization has one event for each agent, all
independent. Its unique maximal configuration contains the tuple $a$.

The following evaluation is defined on the full hom-set $\mathcal C(I,O)$. Let $\operatorname{Out}(f)$ be
the set of tuples labelling configurations of $f$ with one event on every
output component. Define
\[
 k_i(f)=\{u_i(a):a\in\operatorname{Out}(f)\}.
\]
Order subsets of $\mathbb R$ by the Hoare preorder
\[
 U\preceq_H V\quad\Longleftrightarrow\quad
 \forall u\in U\ \exists v\in V\ (u\leq v).
\]
This is reflexive and transitive: reflexivity uses $v=u$, and
transitivity composes the two witnesses with transitivity of $\leq$.
On singleton sets it restricts to the order on real numbers. This particular extension to
nondeterministic or incomplete outcomes is used only to make the context
total; the extension provides a total comparison context for this construction.
\end{construction}

\begin{theorem}
\label{thm:egt-normal}
The constructed concurrent open game has profile space $\prod_i A_i$, the ordinary
unilateral deviations, and exactly the pure Nash equilibria of $(A_i,u_i)$.
Its equilibrium assertion in the extended linear language is the ordinary
Nash inequality.
\end{theorem}
\begin{proof}
\leavevmode
\begin{enumerate}[label=\textup{(\arabic*)},leftmargin=*,itemsep=4pt,topsep=3pt]
  \item \textbf{Recovery of normal-form payoffs.} All factors except agent $i$'s factor in $G_i$ are singletons; consequently
the composite has $\Sigma_i=A_i$. Its realized outcome set at $a$ is
$\{a\}$, so $k_i(g_a)=\{u_i(a)\}$. Strict comparison of these singleton
values is strict comparison of the utilities.
  \item \textbf{Descent along the kernel pair.} The kernel pair $D_i$ fixes
exactly the other agents' actions. Substituting these facts into
Equation~\eqref{eq:egt-ne} gives
precisely the ordinary Nash inequalities.
\end{enumerate}
\end{proof}

\noindent The pure action profile and its realized concurrent strategy determine identical payoffs.
\begin{center}
\begin{tikzcd}[column sep=large,row sep=large]
\prod_i A_i\arrow[r,"g"]\arrow[d,"u_i"']
 &\mathcal C(I,O)\arrow[d,"k_i"]\\
 \mathbb R\arrow[r,"r\mapsto\{r\}"']&\mathcal P(\mathbb R)
\end{tikzcd}
\end{center}

The recovered pure-equilibrium predicate can be empty.
In \textbf{matching pennies}, both agents choose $H$ or $T$. Player $1$
receives $1$ when the choices agree and $-1$ otherwise; agent $2$ receives
the opposite payoff. At an agreeing profile agent $2$ improves by switching,
and at a disagreeing profile agent $1$ improves by switching.
Thus every pure profile has a profitable unilateral deviation.
A mixed strategy would instead be a probability distribution over $H,T$;
existence results using such strategies belong to the probabilistic
extension.

\subsubsection{Boolean games}

\begin{definition}[Boolean games and truth assignments]
\label{def:egt-boolean-games-and-truth-assignments}
A \textbf{finite Boolean game} consists of a finite partition
$V=\coprod_{i\in N}V_i$ of propositional variables, with policy sets
\[
 \Sigma_i=2^{V_i},\qquad
 \Sigma=\prod_i\Sigma_i\cong 2^V,
\]
and a propositional goal $\gamma_i\in\operatorname{Form}(V)$ for each
agent. Here a \textbf{valuation} is a propositional truth assignment $\sigma:V\to2$. Interpret the
formula by its satisfaction subobject
$\llbracket\gamma_i\rrbracket\hookrightarrow\Sigma$ and give
agent $i$ the payoff map
\[
 u_i:\Sigma\longrightarrow\{0,1\},\qquad
 u_i(\sigma)=
 \begin{cases}1,&\sigma\models\gamma_i,\\
               0,&\sigma\not\models\gamma_i.
 \end{cases}
\]
Thus satisfying a goal is strictly preferred to failing it. The
constructed game is the Boolean-game instance of
\citet{BonzonEtAl2006}, realized by the concurrent action protocol.
\end{definition}
Truth of a finite propositional formula at a valuation is decidable by
structural recursion, and quantification over the finitely many assignments
to $V_i$ is likewise decidable. The case distinction in the proof below
therefore does not invoke a general excluded-middle principle for the
assertion language.

\begin{corollary}
\label{cor:egt-boolean}
For a joint valuation $a$,
\[
 a\in\operatorname{NE}
 \quad\Longleftrightarrow\quad
 a\models\bigwedge_{i\in N}
       \bigl(\gamma_i\lor\neg\exists V_i\,\gamma_i\bigr),
\]
where the quantified variables are changed while all other variables retain
their values at $a$.
\end{corollary}
\begin{proof}
\leavevmode
\begin{enumerate}[label=\textup{(\arabic*)},leftmargin=*,itemsep=4pt,topsep=3pt]
  \item \textbf{Strategic comparison.} If $\gamma_i$ already holds, the utility is $1$ and cannot increase. If it
fails, a profitable deviation exists precisely when some assignment of
$V_i$ makes it true while the other variables remain fixed. Taking the
conjunction over agents gives the result.
\end{enumerate}
\end{proof}

The Boolean-game formula is a predicate on complete valuations, not
a linear operation on the move realizing a valuation. Consequently logical
goals may be tested against resource-sensitive interactions without
turning a truth value into a consumable resource.

\subsubsection{Equilibrium under composition and reparameterization}

\begin{theorem}[Selection after sequential and parallel wiring]
\label{thm:cgt-composite-selection}
Let $G:A\to B$ and $H:B\to C$ be concurrent open arenas. Fix a
comparison context $P$ on $\mathcal C(A,C)$ and a reference profile
$(\sigma,\tau)$. The feedback at the composed interface of owner $i$ is
\begin{align*}
 q_i^{H\circ G}((a,b))((a',b'))
 \quad\Longleftrightarrow\quad
 P_i\bigl(&h_{(\tau_{-i},b)}g_{(\sigma_{-i},a)},\\
         &h_{(\tau_{-i},b')}g_{(\sigma_{-i},a')}\bigr).
\end{align*}
For every candidate response $(\sigma',\tau')$,
\[
 \mathsf{BR}_{H\circ G}(P)((\sigma,\tau),(\sigma',\tau'))
 \quad\Longleftrightarrow\quad
 \bigwedge_i
 \varepsilon_{\Sigma_i^G\times\Sigma_i^H}
       ((\sigma'_i,\tau'_i),q_i^{H\circ G}).
\]
The same construction for parallel wiring replaces composition of
concurrent arrows by their tensor. Both laws respect reassociation,
unit insertion, and interchange. They hold for arbitrary comparisons
and repeated ownership of components.
\end{theorem}

\begin{proof}
\leavevmode
\begin{enumerate}[label=\textup{(\arabic*)},leftmargin=*,itemsep=4pt,topsep=3pt]
  \item \textbf{Ownerwise alternatives.} The ownerwise product in the composite arena gives the agency carrier
$\Sigma_i^G\times\Sigma_i^H$. Insert the fixed outsider coordinates
in each component, apply the two plays, and then apply the concurrent
composition map. Pulling $P_i$ back along both arguments of this
function gives the displayed $q_i^{H\circ G}$ by
Theorem~\ref{thm:cgt-coplay}.
  \item \textbf{Strategic comparison.} Applying the agent selection relation
to that costate is precisely the definition of the composite
best-response relation. This calculation uses all alternatives
$(a',b')$ at once. Equivalently,
Theorem~\ref{prop:egt-composite-deviation} identifies their paired
realizations with the component deviation pairs followed by concurrent
composition.
  \item \textbf{Base-change comparison.} The tensor calculation uses the same product of policy
coordinates and the tensor wiring function. Construction~\ref{con:fib-context-composition} gives the composite
reference context, and Theorem~\ref{thm:fib-wiring} gives coherence of
its cartesian lifts and feedback. Pulling back the internal selection
predicates of Proposition~\ref{prop:fib-selection-internal} along equal
costate terms therefore preserves those equations.
\end{enumerate}
\end{proof}


\begin{example}[One owner with complementary decisions]
Equilibrium must compare the complete policy of each owner, including the owner's coordinates in separate components. Consider one agent
choosing bits $x$ and $y$ in two components, with value $xy$. At $(0,0)$,
changing either coordinate alone gives no improvement, while the joint
change $(0,0)\mapsto(1,1)$ is profitable. Componentwise stationarity is
therefore strictly weaker than the correct unilateral-deviation condition.
\end{example}

\begin{construction}[Component deviation subrelations]
Under the canonical isomorphism of
Theorem~\ref{prop:egt-composite-deviation}, define the subrelations
\begin{align*}
 E_i^G&=D_i^G\times\Delta_{\Sigma^H},\\
 E_i^H&=\Delta_{\Sigma^G}\times D_i^H
\end{align*}
of $D_i^{H\circ G}$. The first consists of deviations that change only the
$G$-policy of owner $i$, and the second of deviations that change only its
$H$-policy.
\end{construction}

\begin{theorem}[Deviation-cover factorization]
\label{thm:egt-deviation-cover-factorization}
Suppose that for every owner $i$ the two component-deviation subrelations
cover the composite deviation relation,
\begin{equation}
\label{eq:egt-deviation-cover}
 D_i^{H\circ G}=E_i^G\cup E_i^H.
\end{equation}
Then, for every composite profile $(\sigma,\tau)$,
\[
 (\sigma,\tau)\in\operatorname{NE}(H\circ G,\mathsf{Better})
 \quad\Longleftrightarrow\quad
 \begin{cases}
 \sigma\in\operatorname{NE}
      (G,\mathsf{Better}^{G\mid\tau}),\\
 \tau\in\operatorname{NE}
      (H,\mathsf{Better}^{H\mid\sigma}).
 \end{cases}
\]
The analogous statement holds for tensor and its pullback comparison
contexts.
\end{theorem}

\begin{proof}
\leavevmode
\begin{enumerate}[label=\textup{(\arabic*)},leftmargin=*,itemsep=4pt,topsep=3pt]
  \item \textbf{Strategic comparison.} If the composite profile is an equilibrium, every $G$-only deviation is a
composite deviation with the $H$-profile held fixed. Its two realized
interactions are
\[
 h_\tau\circ g_\sigma
 \quad\text{and}\quad
 h_\tau\circ g_{\sigma'},
\]
so its comparison is exactly
$\mathsf{Better}^{G\mid\tau}_i(g_\sigma,g_{\sigma'})$. Hence no profitable
$G$-deviation exists. The same argument for $H$-only deviations gives the
second component equilibrium.
  \item \textbf{Reverse implication.} Conversely suppose both component equilibrium predicates hold and take any
composite $i$-deviation. By
Equation~\eqref{eq:egt-deviation-cover} it is either in $E_i^G$ or in
$E_i^H$. In the first case its $H$-coordinate is unchanged, so profitability
would contradict the $G$-equilibrium with comparison
$\mathsf{Better}^{G\mid\tau}$.
  \item \textbf{Strategic comparison.} In the second case profitability would
contradict the $H$-equilibrium with comparison
$\mathsf{Better}^{H\mid\sigma}$. Thus no profitable composite
$i$-deviation exists. The tensor proof is identical with whiskering by
composition replaced by tensoring with the fixed other component.
\end{enumerate}
\end{proof}


Disjoint active ownership is a classical sufficient
criterion for the cover equation.

\begin{corollary}[Disjoint active owners; classical specialization]
\label{thm:egt-disjoint}
In ordinary classical $\mathbf{Set}$, let
\[
 \operatorname{act}(G)=\{i:|\Sigma_i^G|>1\}.
\]
If
$\operatorname{act}(G)\cap\operatorname{act}(H)=\varnothing$, then
\[
 (\sigma,\tau)\in\operatorname{NE}(H\circ G,\mathsf{Better})
 \quad\Longleftrightarrow\quad
 \begin{cases}
 \sigma\in\operatorname{NE}
      (G,\mathsf{Better}^{G\mid\tau}),\\
 \tau\in\operatorname{NE}
      (H,\mathsf{Better}^{H\mid\sigma}).
 \end{cases}
\]
The analogous equivalence holds for tensor.
\end{corollary}
\begin{proof}
\leavevmode
\begin{enumerate}[label=\textup{(\arabic*)},leftmargin=*,itemsep=4pt,topsep=3pt]
  \item \textbf{Construction on policy coordinates.} Fix $i$. If the composite profile space is empty there is nothing to prove
for its deviation relation. Otherwise the two policy factors of owner $i$
are inhabited.
  \item \textbf{Ownerwise alternatives.} Classical disjointness of the active-owner sets implies that
at least one of $\Sigma_i^G,\Sigma_i^H$ has at most one element; being
inhabited, that factor is a singleton. Hence every composite $i$-deviation
changes only the other factor, so
$D_i^{H\circ G}=E_i^G\cup E_i^H$. Apply
Theorem~\ref{thm:egt-deviation-cover-factorization}.
\end{enumerate}
\end{proof}

\begin{corollary}[Selection by one agent across a composite interaction]
\label{cor:egt-agent-composite-selection}
For $G:A\to B$, $H:B\to C$, and a comparison context $P$ on
$\mathcal C(A,C)$, the equilibrium condition at $(\sigma,\tau)$ is
\[
 \bigwedge_i\ \forall a\in\Sigma_i^G\ \forall b\in\Sigma_i^H\;
 \neg P_i\bigl(h_\tau g_\sigma,
      h_{\tau[i\leftarrow b]}g_{\sigma[i\leftarrow a]}\bigr).
\]
For tensor composition, replace the two sequential composites by the
corresponding tensor products. These are selections on the complete agent
interfaces of Proposition~\ref{prop:egt-agent-composition}.
\end{corollary}
\begin{proof}
\leavevmode
\begin{enumerate}[label=\textup{(\arabic*)},leftmargin=*,itemsep=4pt,topsep=3pt]
  \item \textbf{Strategic comparison.} The alternatives of agent $i$ in the composite are exactly the pairs
$(a,b)\in\Sigma_i^G\times\Sigma_i^H$. Substitute its composite realization
in Theorem~\ref{thm:cgt-selection-equilibrium}. Its universal quantifier
over that product is the displayed pair of quantifiers. The tensor case is
the same substitution using the tensor realization.
\end{enumerate}
\end{proof}

\subsubsection{Additive reward composition}

The deviation-cover condition is one route to equilibrium
factorization. Another is a specified separable reward: if the evaluation
is additive across the two components, a joint change improves the sum
only if the component best-response conditions fail somewhere. This
argument applies even when one owner controls both components, because
its complete paired policy remains the unit of deviation.

\begin{construction}[Additive reward composition]
\label{con:egt-additive-reward}
For two games with separate retained outcome records, form their
parallel interaction and evaluate agent $i$ by
$v_i(s,t)=u_i^G(s)+u_i^H(t)$. This defines the additive reward
composition, with owner $i$ controlling both $s_i$ and $t_i$.
\end{construction}

\begin{proposition}[Additive equilibrium factorization]
\label{prop:egt-additive-equilibrium}
\[
 \operatorname{NE}(G\boxplus H)
   =\operatorname{NE}(G)\times\operatorname{NE}(H).
\]
\end{proposition}
\begin{proof}
\leavevmode
\begin{enumerate}[label=\textup{(\arabic*)},leftmargin=*,itemsep=4pt,topsep=3pt]
  \item \textbf{Strategic comparison.} A profitable component deviation gives a profitable joint
deviation by keeping the other coordinate fixed. Conversely, at two
component equilibria each component's utility weakly decreases under
its corresponding deviation, so their sum cannot increase under any
joint deviation. This is a specified economic
composition with a proof; it does not impose separability on an
arbitrary pre-existing institutional context.
\end{enumerate}
\end{proof}

\subsubsection{Fibred substitution and changes of strategic power}

\begin{construction}[The deviation comparison of a reparametrisation]
\label{con:fib-deviation-comparison}
For $r:G\to H$ and $i\in N$, put
\[
 E_i(r)=\Sigma^G\times_{\Sigma^H,p_1}D_i^H,
 \qquad
 \lambda_i:D_i^G\to E_i(r),\quad
 (\sigma,\tau)\mapsto(\sigma,(\bar r\sigma,\bar r\tau)).
\]
It is a map over $\Sigma^G$. Its paired realization in $H$ equals
the paired realization in $G$, because $g=h\bar r$.
These maps commute with composition of reparametrisations: factoring
through successive reference pullbacks or through their composite sends
$(\sigma,\tau)$ to the same pair of images.
\end{construction}

\begin{theorem}[Strategic transport and exact quotient transport]
\label{thm:fib-strategic-transport}
For every ownerwise reparametrisation $r:G\to H$ and common comparison
context on their concurrent-arrow sort,
\begin{equation}\label{eq:fib-ne-inclusion}
 \bar r^*\mathsf{NE}_H\leq\mathsf{NE}_G.
\end{equation}
For the ownerwise quotient presentations of
Theorem~\ref{thm:egt-surjection}, the comparison maps $\lambda_i$
are surjective and the displayed inclusion is equality.
For a reference context $s:W\to\Sigma^G$, by contrast, substitution
always gives the exact predicate
\begin{equation}\label{eq:fib-ne-reference}
 s^*\mathsf{NE}_G
 =\bigwedge_i\forall_{d_i^s}\neg(\delta_i^Ga_i)^*P_i,
 \quad
 D_i^s=W\times_{\Sigma^G,p_1}D_i^G,
\end{equation}
where $d_i^s$ and $a_i$ are the two pullback projections.
\end{theorem}

\noindent A reparametrization carries each source deviation to a deviation at the same target reference.
\begin{center}
\begin{tikzcd}[column sep=large,row sep=large]
D_i^G\arrow[r,"\lambda_i"]\arrow[d,"p_1"']
 &\Sigma^G\times_{\Sigma^H}D_i^H\arrow[d]\\
 \Sigma^G\arrow[r,equal]&\Sigma^G
\end{tikzcd}
\end{center}


\begin{proof}
\leavevmode
\begin{enumerate}[label=\textup{(\arabic*)},leftmargin=*,itemsep=4pt,topsep=3pt]
  \item \textbf{Strategic comparison.} An alternative in $D_i^G$ maps under $\lambda_i$ to an alternative of
$H$ at the image reference, with identical realized interactions. If none of the latter alternatives improves the comparison, neither
does an original alternative. Conjunction over $i$ proves
Equation~\eqref{eq:fib-ne-inclusion}.
  \item \textbf{Lifting along the surjection.} For an ownerwise quotient $r$, take an element
$(\sigma,(\bar r\sigma,\eta))$ of $E_i(r)$. Its outsiders already
agree with those of $\bar r\sigma$. Surjectivity of $r_i$ supplies
$b\in\Sigma_i^G$ with $r_i(b)=\eta_i$.
  \item \textbf{Lifting along the surjection: second comparison.} The profile
$\tau=(\sigma_{-i},b)$ is a lift. Thus every alternative in the
reference pullback has an original representative with the same
comparison. Universal quantification over these fibres gives equality. Only the displayed single alternative is lifted; no section of the
quotient is selected.
  \item \textbf{Beck--Chevalley for reference reindexing.} Finally, $D_i^s$ is the cartesian lift of the deviation context along
$s$. Beck--Chevalley from Theorem~\ref{thm:fib-comprehension} gives
$s^*\forall_{p_1}=\forall_{d_i^s}a_i^*$. Inverse image preserves
negation and conjunction, proving Equation~\eqref{eq:fib-ne-reference}.
\end{enumerate}
\end{proof}


\[
\begin{tikzcd}[column sep=large]
 D_i^G\arrow[r,"\lambda_i"]\arrow[dr,"p_1"']&
 E_i(r)\arrow[d]\arrow[r]&D_i^H\arrow[d,"p_1"]\\
 &\Sigma^G\arrow[r,"\bar r"']&\Sigma^H
\end{tikzcd}
\]

\begin{example}[A restriction is a change of game]
One owner has policies $a,b$ with values $0,1$. The larger game's
equilibrium predicate is $\{b\}$. Reparametrising along
$\{a\}\hookrightarrow\{a,b\}$ constructs the one-policy game,
whose equilibrium predicate is $\{a\}$; the inverse image of
$\{b\}$ is empty. Reference substitution along the same underlying
map instead retains the alternative $b$ in $D_i^s$, and correctly
returns the empty predicate. The two operations have the same reference
map but different deviation contexts.
\end{example}

\subsubsection{Changing the presentation of available strategies}

\begin{theorem}
\label{thm:egt-surjection}
Let $f_i:\Sigma_i^G\twoheadrightarrow\Sigma_i^H$ be surjective maps with
$g_\sigma=h_{f(\sigma)}$. In the same comparison context,
\[
 \operatorname{NE}(G,\mathsf{Better})
 =f^{-1}(\operatorname{NE}(H,\mathsf{Better})).
\]
In a value-presented context the same equation is written with $k$ in place
of $\mathsf{Better}$.
\end{theorem}
\begin{proof}
\leavevmode
\begin{enumerate}[label=\textup{(\arabic*)},leftmargin=*,itemsep=4pt,topsep=3pt]
  \item \textbf{Lifting deviations along surjective maps.} Apply \cref{lem:ownerwise-legal-lifting} to the ownerwise surjection
$f$, with full legal domain. It lifts every unilateral target deviation
at a given reference while fixing all outsider policies. Since both
realized endpoints agree under $g=h f$, the predicates of profitable
deviations agree.
  \item \textbf{Lifting along the surjection.} Universal negation gives the stated equality of Nash
predicates. The lift chooses preimages only for the finitely many coordinates varied by the deviation.
\end{enumerate}
\end{proof}

\begin{remark}[Restriction of available policies]
For injective parameter maps the useful implication is different. A profile
that is an equilibrium in a larger strategy space remains an equilibrium
when the available deviations are restricted and that profile is retained.
The converse fails because enlarging a strategy space can create a profitable
deviation. The $2$-cell language therefore distinguishes redundant
descriptions from genuine changes in strategic power.
\end{remark}

\begin{example}[Complementary commitments]
\label{ex:egt-commitments}
Two owners independently choose $0$ or $1$. Tensor their value-state
presentations and evaluate the completed record by
\[
 u(0,0)=(0,0),\quad u(0,1)=(0,-1),\quad
 u(1,0)=(-1,0),\quad u(1,1)=(1,1).
\]
The profile $(0,0)$ has only unilateral deviations giving the deviator
$-1$; at $(1,1)$ a unilateral withdrawal gives $0$. Both are equilibria.
At either off-diagonal profile the contributor improves by withdrawing.
Hence the equilibrium set is exactly $\{(0,0),(1,1)\}$.
\end{example}

\subsection{Situated agent presentations}
\label{agt-sec-situated}

\subsubsection{Declared presentations and reference states}

\begin{definition}[Situated presentation]
A situated presentation is displayed data
\[
 \ctx=(a,L\hookrightarrow\Sigma,s:W\to L,
       (o_i:W\to O_i)_{i\in N},(P_i)_{i\in N}).
\]
Here $a$ is an owned agent presentation, $L$ is its legal policy domain,
$s$ selects the reference profile, $o_i$ specifies observation, and $P_i$
is its comparison interface, typed in Section~\ref{agt-sec-equilibrium}.
The state space may retain histories, private inputs, policies, and
institutional records. Each observation is corestricted to its effective
image whenever local states are quantified.
\end{definition}

\begin{construction}[Situated and legal presentations generated by a protocol]
\label{agt-con-protocol-situation}
Let $U$ be the product of the external-input carriers of a finite
causal decision protocol, $V$ its node set, and $\Sigma$ its complete
owner profile space. The causal-height recursion gives the unique map
\[
 r:U\times\Sigma\longrightarrow\prod_{v\in V}M_v,
 \qquad W_{\mathrm{prot}}
 =\{(u,\sigma,a):a=r(u,\sigma)\}\cong U\times\Sigma.
\]
Let $s:W_{\mathrm{prot}}\to\Sigma$ be projection and
$z:W_{\mathrm{prot}}\to Z$ the specified outcome-record map. Nodewise
observations are the maps of Theorem~\ref{agt-thm-protocol-prefix}:
they are evaluated on their \emph{own causal prefix} and use only
actually available incoming wires. They need not coincide with one
owner's later, accumulated observation record. Any final owner
observation used in a situated presentation is the image of the
specified actually delivered end-of-protocol signal; it is not the
unrestricted product of all its earlier and later signals.

For an exhibited institutional acceptance predicate $Q\hookrightarrow Z$,
first form the accepted-\emph{state} subobject
$P=z^*Q\hookrightarrow W_{\mathrm{prot}}$. There are canonical,
distinct policy-domain constructions:
\begin{align*}
 L_{\exists}&=\exists_sP=\operatorname{im}(s|_P)
     \hookrightarrow\Sigma,\\
 L_{\forall}&=\forall_sP\hookrightarrow\Sigma.
\end{align*}
The restricted map $s|_P$ factors canonically into a surjection
$P\twoheadrightarrow L_{\exists}$ and its image inclusion; this gives
the situated presentation of \emph{accepted executions}. Pulling back
$s$ along $L_{\forall}\hookrightarrow\Sigma$ gives
$W_{\forall}\to L_{\forall}$, the presentation of policies whose
\emph{every} external-input execution passes the acceptance test.
For exhibited certificates $k:K\to W_{\mathrm{prot}}$, the policy
domain is $L_K=\operatorname{im}(sk)$, retaining $K$ as its witness
carrier. None of these subobjects is confused with the accepted-state
predicate $P$. Preference comparisons and institutional acceptance rules
remain separately supplied economic and normative inputs.
\end{construction}
\begin{proposition}[Universal characterization of protocol legality]
\label{agt-prop-protocol-legality}
$L_{\exists}$ consists exactly of profiles having at least one
accepted execution. The subobject $L_{\forall}$ is the greatest
$J\hookrightarrow\Sigma$ such that all protocol states with reference
in $J$ lie in $P$. These constructions are stable under substitution of
ordinary protocol parameters along cartesian squares.
\end{proposition}
\begin{proof}
\leavevmode
\begin{enumerate}[label=\textup{(\arabic*)},leftmargin=*,itemsep=4pt,topsep=3pt]
  \item \textbf{Adjunction and transposition.} The first assertion is the image formula
$\sigma\in L_{\exists}\iff\exists w\,(s(w)=\sigma\wedge P(w))$. For the second, the indexed adjunction gives
\[
 s^*J\subseteq P\iff J\subseteq\forall_sP.
\]
Taking $J=L_{\forall}$ yields its desired acceptance and greatestness. Beck--Chevalley for the image/existential and the right dependent
adjoint supplies substitution, with the same source pullback as the
protocol's reference map. No execution representative is selected.
\end{enumerate}
\end{proof}

\noindent The accepted execution space factors onto precisely the existentially legal profiles.
\begin{center}
\begin{tikzcd}[column sep=large,row sep=large]
P\arrow[r,two heads]\arrow[d,hook]
 &L_{\exists}\arrow[d,hook]\\
 W_{\mathrm{prot}}\arrow[r,"s"']&\Sigma
\end{tikzcd}
\end{center}

A situated presentation depends on declared owner and policy data.
It does not extract a new reusable owner presentation from an arbitrary
linear agent variable. Instead, an agent executes within a state context,
whose displayed signature determines which control and observation
tests are meaningful.

The need for the economic fields can be checked independently of
the categorical construction. Reverse the preference on the same
two-choice concurrent process: equilibrium changes, although realization
does not. Change the permitted subobject from $E$ to $\varnothing$ for an
unchanged execution span:
permission changes without changing the underlying process. Replace a
constant observation by an injective observation of the same worlds:
knowledge changes without changing execution. Preferences, permissions,
and nondelivered signals therefore cannot be reconstructed from a bare
operational arrow. Where the information is generated by a protocol,
Construction~\ref{agt-con-protocol-situation} instead derives its
associated maps by image, pullback, and quotient. No additional agency
equation is needed.

These constructions occupy three judgments: ordinary terms in
$\cT(W)$, linear terms in $\cL(W)$, and assertions in $\cP(W)$.
The entailment $W\mid\Phi\vdash\varphi$ means
$\bigcap_{\psi\in\Phi}\sem\psi\subseteq\sem\varphi$.
Logical conjunction, implication, and ordinary quantification act on
assertions; they do not supply a tensor diagonal to a linear resource.
The concrete predicate fibres are Boolean, while proofs using only the
adjunctions and Heyting operations remain constructive.

\subsection{Legal feedback, best response and equilibrium}
\label{agt-sec-equilibrium}

\subsubsection{Comparison feedback and its pullback}

The legal version of selection must compare only legal alternatives.
Let $H$ denote the retained outcome or interaction carrier, let
$g:L\to H$ be its realization, and let
$Q_i\hookrightarrow H\times H$ classify strict improvement in the
second coordinate. Pulling this relation back along the paired
realization supplies the improvement predicate on legal profiles.
\[
 P_i(\sigma,\tau)\iff Q_i(g\sigma,g\tau).
\]
The realization $g$ and comparison $Q_i$ may be indexed by outsider policies or by an ordinary
state context; the calculation is the same in each fibre. Evaluating a
component inside a larger interaction first supplies its contextual
realization into $H$. Inverse image of the joint comparison then gives
the component feedback, compatibly with further wiring.

\begin{proposition}[Compositional comparison feedback]
\label{agt-prop-feedback}
If $k:H\to Z$ is a further contextual realization and $Q\subseteq Z^2$,
then
\[
 (g\times g)^*(k\times k)^*Q
 =((kg)\times(kg))^*Q.
\]
Hence evaluating a component in a composite context agrees with successive
contextual evaluations, retaining the entire joint predicate.
\end{proposition}
\begin{proof}
\leavevmode
\begin{enumerate}[label=\textup{(\arabic*)},leftmargin=*,itemsep=4pt,topsep=3pt]
  \item \textbf{Equality of the two constructions.} Both sides are $((kg)\times(kg))^*Q$ by the functoriality of
inverse image proved in \cref{thm:cgt-coplay,thm:cgt-wiring-soundness}. Their common test at $(\sigma,\tau)$ is
$Q(kg\sigma,kg\tau)$, including all joint alternatives.
\end{enumerate}
\end{proof}

The backward comparison must remain joint: separate marginal
predicates would lose improvements that depend on complementary changes
at decisions controlled by the same owner.

\subsubsection{Selection and best response}

For a legal own-choice fibre
$S_i^L(u)=\{v:\Sigma_i\mid(u,v)\in L\}$, define the canonical selection
predicate
\[
 \operatorname{Sel}_i(u,v)
 \iff v\in S_i^L(u)\ \wedge\
 \forall v'\in S_i^L(u)\,
       \neg Q_i\bigl(g(u,v),g(u,v')\bigr).
\]
The predicate therefore retains exactly the legal choices admitting
no strictly improving legal alternative. Its construction uses only
comprehension, comparison feedback, and dependent universal
quantification, so it is an assertion in the established doctrine.

\begin{definition}[Best response and equilibrium]
For a legal reference profile $\sigma$, define
\begin{align*}
 \BR_i(\sigma,v)
 &\iff v\in S_i^L(\sigma_{-i})\ \wedge\
   \forall u\in S_i^L(\sigma_{-i})\,
   \neg Q_i\bigl(g(\sigma[i:=v]),g(\sigma[i:=u])\bigr),\\
 \NE(\sigma)
 &\iff\bigwedge_{i\in N}\BR_i(\sigma,\sigma_i).
\end{align*}
The same predicates at a situated state are obtained by substituting
$\sigma=s(w)$.
\end{definition}

\begin{proposition}[Equilibrium from owned alternatives]
\label{agt-prop-equilibrium}
\[
 \NE(\sigma)\iff
 \bigwedge_{i\in N}\forall\tau\,
 \bigl(\sigma D_i^L\tau\Rightarrow\neg P_i(\sigma,\tau)\bigr).
\]
It is therefore an assertion constructed from the agent's displayed control,
legality, realization, and comparison interfaces.
\end{proposition}
\begin{proof}
\leavevmode
\begin{enumerate}[label=\textup{(\arabic*)},leftmargin=*,itemsep=4pt,topsep=3pt]
  \item \textbf{Predicate calculation.} For fixed $\sigma$, the map $u\mapsto\sigma[i:=u]$ is a bijection from
$S_i^L(\sigma_{-i})$ to the target fibre of $D_i^L$ at $\sigma$. Under this bijection the two comparisons are the same pullback of $Q_i$. Dependent universal quantification along the bijection proves the equality.
\end{enumerate}
\end{proof}

An informed response is represented by an element of the observation-uniform legal policy carrier $\Pol_{o_i}(F_i)$, evaluated through its realized choices. For instance the assertion that $\pi$ is uniformly undominated is
\[
 \pi\in\Pol_{o_i}(F_i)\ \wedge\
 \forall w\ \forall\rho\in\Pol_{o_i}(F_i)\,
 \neg Q_i\bigl(g(s(w)[i:=\pi(o_iw)]),
                g(s(w)[i:=\rho(o_iw)])\bigr).
\]
The displayed quantifiers distinguish one uniform policy from a
collection of pointwise responses. Other economic criteria may change
the comparison relation or test whole policy outcomes, but they use the
same linear execution operations.

\begin{example}[Jointly controlled decisions]
One owner controls two bits $x,y$, and values the outcome by $xy$.
At $(0,0)$, changing either coordinate in isolation yields no improvement;
changing the complete owned policy to $(1,1)$ does. The product law
$\Sigma_i^{a\seq b}=\Sigma_i^a\times\Sigma_i^b$ makes the latter an
ordinary unilateral deviation by that owner. Consequently equilibrium of
the composite is not established by testing isolated decision positions.
\end{example}

\part{Information, institution, and epistemic dynamics}

\section{Observation, Perspective, and Knowledge}\label{sec:paper-6}

Information matters to agency before it enters a modal formula:
it determines which responses an agent can implement. Following the
local-state perspective of \citet{HalpernFagin1989KnowledgeAction},
we first represent available information by an effective observation
quotient. Its descent property classifies admissible responses; its
predicate adjunction then constructs knowledge. The same quotient also
acts on dependent linear families, distinguishing proof-relevant
perspective from the truth conditions of knowledge. Finally, joint
observations and their quotients distinguish distributed information
from information actually made common.

\subsection{Observation and admissible policies}
\label{sec:egt-information}

Before reasoning about what an agent knows, we must determine which
histories it can distinguish when choosing an action. The raw signal
is factored through its image to obtain an effective observation quotient.
Functions on this quotient are exactly the implementable observation-based
responses, by kernel-pair descent. The same quotient later induces the
knowledge operation of Section~\ref{sec:egt-logic}; the admissibility
condition is not an additional constraint placed on an arbitrary
history-dependent response.

\subsection{Observation-based responses of an agent}

\begin{construction}[Observation quotient and admissible policies]
Let $H$ be a set of histories available at a decision and let
$\widetilde o:H\to Y$ be the raw signal supplied by the protocol. Construct
the observation object by the image factorization in $\mathbf{Set}$,
\[
 H\xrightarrow{o}O\lhook\joinrel\longrightarrow Y,
 \qquad O=\widetilde o[H],
\]
so $o$ is the effective quotient obtained by corestricting the raw signal to
its image. Equivalently, if indistinguishability is supplied directly as an
equivalence relation $R\rightrightarrows H$, one may take its coequalizer
$H\twoheadrightarrow H/R$. Available policies with action set $A$ are then
functions $\pi:O\to A$. When $H,O,A$ are finite, these functions have the
concurrent realization of Proposition~\ref{prop:egt-functions}.
\end{construction}

For a concrete example take histories $H=\{h_1,h_2,h_3\}$, observations
$O=\{o,o'\}$ with $o(h_1)=o(h_2)=o$ and $o(h_3)=o'$, and actions
$A=\{0,1\}$. An admissible history response is determined by two choices,
one for $o$ and one for $o'$, so there are four such policies. The response
assigning $0$ at $h_1$ and $1$ at $h_2$ is inadmissible for this observation.
The raw signal's image removes unused signal labels without changing
any decision the agent can make.

\begin{proposition}[Admissibility as effective descent]
\label{prop:egt-observation}
Let
\[
 H\times_OH
 \mathrel{\substack{\xrightarrow{p_2}\\[-2pt]\xrightarrow[p_1]{}}}H
 \xrightarrow{o}O
\]
be the kernel pair of the constructed observation quotient. For every action
set $A$, precomposition gives an equalizer
\[
 A^O\xrightarrow{o^*}A^H
 \mathrel{\substack{\xrightarrow{p_2^*}\\[-2pt]
                     \xrightarrow[p_1^*]{}}}A^{H\times_OH}.
\]
Consequently a response $f:H\to A$ is implementable using only the observation
$o$ exactly when it equalizes the observation kernel pair, equivalently when
there is a unique $\pi:O\to A$ with $f=\pi o$.
\end{proposition}
\begin{proof}
\leavevmode
\begin{enumerate}[label=\textup{(\arabic*)},leftmargin=*,itemsep=4pt,topsep=3pt]
  \item \textbf{Descent along the kernel pair.} The image factorization makes $o$ surjective. By effective regularity of
$\mathbf{Set}$ it is the coequalizer of its kernel pair. A map
$f:H\to A$ equalizes the two projections precisely when
\[
 o(h)=o(h')\Longrightarrow f(h)=f(h').
\]
The coequalizer universal property therefore gives a unique factor
$\pi:O\to A$ with $f=\pi o$.
  \item \textbf{Representation and Yoneda decoding.} Applying the contravariant representable
$\mathbf{Set}(-,A)$ to this coequalizer is exactly the displayed equalizer. Thus ``admissible'' introduces no further policy datum: it names this descent
property.
\end{enumerate}
\end{proof}

\noindent Admissible response strategies factor through the effective observation quotient.
\begin{center}
\begin{tikzcd}[column sep=large,row sep=large]
H\arrow[r,two heads,"o"]\arrow[d,"f"']
 &O\arrow[dl,dashed,"\pi"]\\
 A&
\end{tikzcd}
\end{center}

\begin{figure}[htbp]
\centering
\begin{tikzcd}[column sep=large,row sep=large]
H \arrow[r,"o",two heads] \arrow[dr,"f"'] & O \arrow[d,dashed,"\pi"] \\
& A
\end{tikzcd}
\qquad
\begin{tikzcd}[column sep=large,row sep=large]
D(H) \arrow[r,"\widehat o"] \arrow[dr,"\widehat f"'] &
D(O) \arrow[d,"\widehat\pi"] \\
& D(A)
\end{tikzcd}
\caption{An information restriction is effective descent through the
constructed observation quotient; finite response protocols realize the
factorization concurrently.}
\end{figure}

The left square expresses the information restriction; for finite
data, Proposition~\ref{prop:egt-functions} realizes the same factorization
by concurrent response strategies. The policy responds after its actual
observation event. Other strategies on the arenas may stall or behave
differently; effective descent classifies the specified total response
family. In particular, a decision cannot depend on a signal whose event
has not occurred in its causal past.

\subsection{An individual agent's knowledge from observation}

\begin{definition}[Observations, indistinguishability, and knowledge]
\label{def:egt-observations-indistinguishability-and-knowledge}
An \textbf{observation} is information actually made available to an
agent. States yielding the same observation are
\textbf{indistinguishable} to that agent. It knows a predicate when the
predicate holds at every indistinguishable state.

Let $W$ be a set of economic states and $o_i:W\to O_i$ an observation
map. Its epistemic relation is the kernel-pair span
\[
 W\xleftarrow{p_1}W\times_{O_i}W\xrightarrow{p_2}W.
\]
Knowledge is the corresponding universal span modality:
\begin{equation}
\label{eq:egt-knowledge-span}
 K_iP
 =[\operatorname{ker}(o_i)]P
 =\forall_{p_1}p_2^*P.
\end{equation}
\end{definition}

For example, observing
whether a public total is zero need not reveal the individual contributions
that produced it. The resulting knowledge concerns everything forced by
that observation, not a probability assigned to the unobserved states.

\begin{definition}[Factivity and introspection]
\label{def:egt-factivity-introspection}
A knowledge operator $K_i$ is \textbf{factive} when $K_iP\leq P$ for every
predicate $P$. It has \textbf{positive introspection} when $K_iP\leq K_iK_iP$,
and \textbf{negative introspection} when $\neg K_iP\leq K_i\neg K_iP$.
These conditions say, respectively, that knowledge implies truth, that an
agent knows what it knows, and that it knows when it does not know.
\end{definition}

Because an observation kernel pair is an equivalence relation, it
also determines the familiar factivity and introspection conditions.
The same reasoning applies after epistemic action update when the
updated relation arises from the effective observation quotient of
Proposition~\ref{prop:egt-action-observation-quotient}.

\begin{proposition}[Factivity and introspection]
For an observation map, $K_iP\leq P$, $K_iP\leq K_iK_iP$, and
$\neg K_iP\leq K_i\neg K_iP$.
\end{proposition}
\begin{proof}
\leavevmode
\begin{enumerate}[label=\textup{(\arabic*)},leftmargin=*,itemsep=4pt,topsep=3pt]
  \item \textbf{Descent along the kernel pair.} The observation kernel pair is reflexive, so the reference state is an
alternative. Every related state has the same equivalence class. Truth of
$P$ throughout that class is therefore invariant throughout it, as is its
negation.
  \item \textbf{Predicate calculation.} These arguments use the implication and negation operations of the predicate doctrine.
\end{enumerate}
\end{proof}

Deviation and knowledge thus share a relational construction but not
an interpretation: the kernel pair of outsider projection records changes
within a coalition's control, while the kernel pair of observation records
indistinguishable states. The corresponding modalities quantify over
different alternatives.

\begin{corollary}[Observation quotients classify invariant predicates]
\label{prop:egt-observation-interior}
For a constructed observation quotient $o:W\twoheadrightarrow O$, put
$K_o=o^*\forall_o$ and $M_o=o^*\exists_o$. Knowledge is the
idempotent comonad $K_o$, and possibility is the idempotent monad $M_o$.
For a predicate $P$, the following are equivalent: it is the pullback
of a unique predicate on $O$; it is invariant on $\ker(o)$;
$K_oP=P$; and $M_oP=P$. Moreover $K_oP\leq P\leq M_oP$ and
\[
 \operatorname{Coalg}(K_o)\cong\mathsf{Pred}(O)
 \cong\operatorname{Alg}(M_o).
\]
\end{corollary}
\begin{proof}
\leavevmode
\begin{enumerate}[label=\textup{(\arabic*)},leftmargin=*,itemsep=4pt,topsep=3pt]
  \item \textbf{Adjunction and transposition.} The relation factorization identifies the knowledge and possibility
operators with the displayed composites. Apply
Theorem~\ref{thm:egt-invariant-adjunctions} and effective predicate descent.
\end{enumerate}
\end{proof}

The quotient plays two roles that are now visibly linked. Its maps
into action sets classify implementable responses by
Proposition~\ref{prop:egt-observation}; its descended predicates classify
observation-invariant assertions by
Corollary~\ref{prop:egt-observation-interior}. The operational and
epistemic interpretations therefore arise from one effective quotient,
not from separate information primitives.

\begin{proposition}[An agent's responses and observation-local assertions]
\label{prop:egt-agent-observation-interface}
For an effective observation quotient $o_i:W\twoheadrightarrow\Omega_i$ and
an action set $M$, there are canonical identifications
\[
 M^{\Omega_i}\cong
 \{f:W\to M:\ker(o_i)\subseteq\ker(f)\},
 \qquad
 \mathsf{Pred}(\Omega_i)\cong
 \{P\in\mathsf{Pred}(W):K_iP=P\}.
\]
The first identifies admissible responses; the second identifies assertions
local to the agent's current observation.
\end{proposition}
\begin{proof}
\leavevmode
\begin{enumerate}[label=\textup{(\arabic*)},leftmargin=*,itemsep=4pt,topsep=3pt]
  \item \textbf{Observation and knowledge.} The first bijection is precomposition with $o_i$ and its inverse is the unique
factorization of Proposition~\ref{prop:egt-observation}.
  \item \textbf{Lifting along the surjection.} For the second,
$o_i^*B$ is constant on every observation fibre, hence equals its own
knowledge predicate. If $K_iP=P$, every fibre meeting $P$ lies entirely in
$P$, so $P=o_i^{-1}(o_i[P])$. Surjectivity makes the representing subset
unique.
  \item \textbf{Observation and knowledge.} These maps are exactly the predicate descent and knowledge
identifications already constructed.
\end{enumerate}
\end{proof}

Observation-local assertions live on the common situated carrier
$W$ alongside assertions about other agents and later events. At a
particular decision, however, the relevant quotient is that decision's
actual history-to-signal map. Applying the same factorization at the
correct causal stage yields precisely the responses available there.

\subsection{Qualitative belief and common knowledge}

\begin{definition}[Possibility model and belief]
\label{def:egt-possibility-model-and-belief}
A \textbf{possibility model} records which states each agent considers
possible. An agent \textbf{believes} a predicate when it holds at every
state it considers possible. The actual state need not be among them, so
belief need not imply truth. If the possibility set is empty, universal
belief is vacuous; no nonemptiness condition is implicit.

A possibility model over a profile space $\Sigma$ consists of a set $W$,
a map $s:W\to\Sigma$, and relations $R_i\hookrightarrow W^2$.
Put $P_i(w)=\{w':wR_iw'\}$, $B_i=[R_i]$, and
$E(P)=\bigwedge_iB_iP$.
No relational laws are implicit. Observation maps construct the special
case $R_i=\ker(o_i)$.
\end{definition}

\begin{definition}[Common belief and common knowledge]
\label{def:egt-common-belief-and-common-knowledge}
\textbf{Common belief} requires belief by every agent, belief by every
agent that all agents believe, and all further finite iterations.
For observation relations, \textbf{common knowledge} also includes truth at
the current state. The two fixed points below distinguish iterated belief
from truth along all finite chains, including the chain of length zero.

Define
\begin{equation}\label{eq:egt-common-belief}
 \mathsf{CB}(P)=\nu X\,E(P\wedge X),
 \qquad C_P=\nu X\,(P\wedge E(X)).
\end{equation}
When the relations are observation kernel pairs, call $C_P$ the
\textbf{common-knowledge predicate} $\mathsf{CK}(P)$.
Let $R=\bigcup_iR_i$ and let $R^*$ be the reflexive transitive closure
constructed from finite $R$-paths, including paths of length zero.
\end{definition}

\begin{proposition}[Truth, common belief, and finite paths]
\label{prop:egt-true-common-belief}
For every finite owner list, including the empty list,
\[
 C_P=P\wedge\mathsf{CB}(P)=[R^*]P.
\]
In particular $\mathsf{CK}(P)\leq P$.
\end{proposition}
\begin{proof}
\leavevmode
\begin{enumerate}[label=\textup{(\arabic*)},leftmargin=*,itemsep=4pt,topsep=3pt]
  \item \textbf{Coinduction and bisimulation.} The fixed-point equation for $C_P$ gives $C_P\leq P$ and
$C_P\leq E(C_P)=E(P\wedge C_P)$, so coinduction gives
$C_P\leq P\wedge\mathsf{CB}(P)$.
  \item \textbf{Reverse implication.} Conversely, with
$B=\mathsf{CB}(P)$, the predicate $P\wedge B$ equals
$P\wedge E(P\wedge B)$ and is fixed for the operator defining $C_P$. For the path description, any post-fixed point of $X\mapsto P\wedge E(X)$
satisfies $P$ after every finite path by induction on path length.
  \item \textbf{Reverse implication: completion.} Conversely $[R^*]P\leq P$ by the zero-length path and is closed under
one further $R$-step, so it is post-fixed. Coinduction proves equality. For $N=\varnothing$, these formulas give
$\mathsf{CB}(P)=\top$ and $C_P=P$, as required.
\end{enumerate}
\end{proof}

\begin{theorem}[Common knowledge as a coreflection]
\label{thm:egt-common-coreflection}
For observation relations $R_i=\ker(o_i)$, the relation
$E=(\bigcup_iR_i)^*$ is an equivalence relation. Its quotient
$c:W\twoheadrightarrow W/E$ satisfies
\[
 \mathsf{CK}=c^*\forall_c,\qquad
 \operatorname{Coalg}(\mathsf{CK})
 =\bigcap_i\operatorname{Coalg}(K_i).
\]
Thus common knowledge is the greatest jointly observation-invariant
predicate below the given predicate.
\end{theorem}
\begin{proof}
\leavevmode
\begin{enumerate}[label=\textup{(\arabic*)},leftmargin=*,itemsep=4pt,topsep=3pt]
  \item \textbf{Fixed-point argument.} Reflexivity and transitivity hold by the path construction. Reversing each
edge of a finite path proves symmetry. A predicate is $E$-invariant iff
it is invariant under every $R_i$: one direction uses the inclusions and
the other is induction on path length.
  \item \textbf{Adjunction and transposition.} The finite-path characterization
gives $\mathsf{CK}=[E]=c^*\forall_c$. Apply
Theorem~\ref{thm:egt-invariant-adjunctions} to $c$ and to the image
factorization of each $o_i$. Their fixed predicates are exactly the
invariant predicates just identified.
  \item \textbf{Ownerwise alternatives.} If the owner list is empty, $E$ is
equality and the same proof gives the identity coreflection.
\end{enumerate}
\end{proof}

\noindent Each observation quotient factors through the equivalence quotient defining common knowledge.
\begin{center}
\begin{tikzcd}[column sep=large,row sep=large]
W\arrow[r,"o_i",two heads]\arrow[dr,"c"']
 &O_i\arrow[d,dashed,"q_i"]\\
 &W/E
\end{tikzcd}
\end{center}

\begin{corollary}[Common knowledge after epistemic action update]
\label{cor:egt-action-common-knowledge}
Let $R_i=\ker(o_i)$ and $S_i=\ker(b_i)$ in the action model of
Section~\ref{sec:egt-action-model-update}. On the updated carrier $U$,
put $c_{\mathcal E}:U\twoheadrightarrow U/{\sim}$, where $\sim$ is
generated by the union of its updated knowledge relations. Then
\[
 \mathsf{CK}_{\mathcal E}(Q)
 =\left[\left(\bigcup_iR_i^{\mathcal E}\right)^*\right]Q
 =c_{\mathcal E}^*\forall_{c_{\mathcal E}}Q.
\]
For the pointed action $e$, the corresponding assertion at the old
state sort is $[\mathcal E,e]\mathsf{CK}_{\mathcal E}(Q)$.
For a singleton public action, this recovers the common-knowledge
specialization of Section~\ref{sec:egt-dynamic-update}.
\end{corollary}
\begin{proof}
\leavevmode
\begin{enumerate}[label=\textup{(\arabic*)},leftmargin=*,itemsep=4pt,topsep=3pt]
  \item \textbf{Observation and knowledge.} Proposition~\ref{prop:egt-action-observation-quotient} constructs each
updated kernel relation. Their reflexive transitive union is an
equivalence relation, and Theorem~\ref{thm:egt-common-coreflection}
identifies its box with the displayed quotient coreflection. Equation~\eqref{eq:egt-action-dynamic-modalities} transports the
result through the chosen action component; the singleton case is
Theorem~\ref{thm:egt-action-model-semantics}.
\end{enumerate}
\end{proof}

\begin{theorem}[Categorical validity of agency operations]
\label{thm:egt-agency-soundness}
The ordinary dependent fibres and their subterminal fragment recover
$\mathsf P_G$, its quantifier adjunctions, and comprehension. The
relation-action and comprehension-update equations hold in these fibres,
as do the order-theoretic induction and coinduction properties of positive
fixed-point operators. For external parameters
$\Gamma$ and a fixed state sort $X$, pointwise parameterized fixed points
commute with ordinary substitution in $\Gamma$.
\end{theorem}
\begin{proof}
\leavevmode
\begin{enumerate}[label=\textup{(\arabic*)},leftmargin=*,itemsep=4pt,topsep=3pt]
  \item \textbf{Base-change comparison.} Theorem~\ref{thm:fib-comprehension} supplies the quantifier
adjunctions, equality, comprehension, and Beck--Chevalley. Construction~\ref{con:fib-dependent-types} supplies dependent types;
Proposition~\ref{prop:fib-span-typing} supplies the internal modal
typing.
  \item \textbf{Fixed-point argument.} The modal and update equations were proved using these operations,
and Theorem~\ref{thm:fib-linear-logical-compatibility} identifies their
operational tests with substitution of concurrent realization terms. Structural induction on polarity proves monotonicity of each fixed-point
body, so the fixed-point theorem applies.
  \item \textbf{Fixed-point argument: second comparison.} For a family $F_\gamma$ on
$\mathsf{Pred}(X)$, the operator on $\mathsf{Pred}(\Gamma\times X)$ acting
fiberwise has least fixed point with fiber $\mu F_\gamma$: it is fixed
fiberwise and lies below every fixed predicate fiberwise. The greatest case
is dual. Substitution in $\Gamma$ merely selects these same fibers.
  \item \textbf{Functoriality and substitution.} Changes to the state carrier are handled by the separate transport results,
not by this parameter-substitution assertion.
\end{enumerate}
\end{proof}

\begin{remark}[Relation to strategy logics]
Explicit quantification over complete policies places these assertions in
contact with strategy logics such as \citet{MogaveroEtAl2012}.
Here the interpretation is supplied by concurrent open-arena realizations and
the predicate doctrine. The soundness theorem concerns this interpretation;
it does not import model-checking results for a different class of models.
\end{remark}

\subsection{Dependent observation descent}
\label{sec:cgt-information}

The quotient already classifies individual observable responses and
predicates. To compare different presentations and information refinements,
we must also transport whole dependent families and their witnesses.
Effective descent provides this transport; Morita comparisons identify
presentations of the same observations, while kernel inclusions order
observation strength. The subsequent pooled-policy comparison separates
what multiple agents could compute from a delivered joint signal from
what they can compute independently.

\subsubsection{Descent in the dependent-type fibration}

\begin{construction}[The slice adjunctions]
\label{con:egt-slice-adjoints}
For a map $q:X\to Q$, regard an object of $\mathbf{Set}/X$ as a
family $(E_x)_{x\in X}$. Define
\[
 (q_!E)_u=\coprod_{q(x)=u}E_x,
 \qquad (q^*F)_x=F_{q(x)},
 \qquad(q_*E)_u=\prod_{q(x)=u}E_x.
\]
The universal properties of disjoint union and product give natural
bijections proving $q_!\dashv q^*\dashv q_*$. The empty sum and empty
product are included. On maps these functors act componentwise.
Write $T_q=q^*q_!$ and $G_q=q^*q_*$ for the induced monad and comonad
on $\mathbf{Set}/X$.
\end{construction}

\begin{definition}[Kernel groupoid and its actions]
\label{def:egt-presentation-groupoid}
For a surjection $q:X\twoheadrightarrow Q$, the \textbf{presentation
groupoid} $R_q$ has objects $X$ and an arrow $(x,y):y\to x$ exactly
when $q(x)=q(y)$. Its identity, inverse, and composition are
$(x,x)$, $(x,y)^{-1}=(y,x)$, and $(x,y)(y,z)=(x,z)$.
An action of $R_q$ is a family $E_x$ with maps
$\tau_{xy}:E_y\to E_x$ satisfying
\[
 \tau_{xx}=\operatorname{id},\qquad
 \tau_{xy}\tau_{yz}=\tau_{xz}.
\]
An equivariant map is a family of functions commuting with all these maps.
The resulting category is denoted $\operatorname{Act}(R_q)$.
\end{definition}

\begin{theorem}[Effective descent as algebra and coalgebra]
\label{thm:egt-descent-algebras}
There are canonical equivalences
\[
 \operatorname{Alg}(T_q)\simeq\operatorname{Act}(R_q)
 \simeq\mathbf{Set}/Q
 \simeq\operatorname{Coalg}(G_q).
\]
The comparison from $\mathbf{Set}/Q$ is pullback with its canonical
transport. It requires no chosen section of $q$.
\end{theorem}

\begin{proof}
\leavevmode
\begin{enumerate}[label=\textup{(\arabic*)},leftmargin=*,itemsep=4pt,topsep=3pt]
  \item \textbf{Unit and identity laws.} The unit of $T_q$ includes $E_x$ in the $x$-summand, and its
multiplication forgets the middle index in a twice-indexed sum. Consequently a $T_q$-algebra is exactly a family $\tau_{xy}$ with the
displayed identity and composition equations. These equations imply that
$\tau_{xy}$ and $\tau_{yx}$ are inverse. Algebra morphisms are exactly
equivariant maps.
  \item \textbf{Construction on states and events.} Given an action, define
\[
 \bar E_u=\left(\coprod_{q(x)=u}E_x\right)/\sim,
 \qquad (x,e)\sim(y,e')\ \Longleftrightarrow\ e=\tau_{xy}(e').
\]
The action laws prove that this is an equivalence relation. For each $x$,
the map $E_x\to\bar E_{q(x)}$, $e\mapsto[(x,e)]$, is injective by
$\tau_{xx}=\operatorname{id}$ and surjective because
$[(y,e')]=[(x,\tau_{xy}(e'))]$.
  \item \textbf{Descent along the kernel pair.} These bijections intertwine the action
with the identity transports on the pullback of $\bar E$. Equivariant
maps descend by $[(x,e)]\mapsto[(x,f_x(e))]$.
  \item \textbf{Reverse implication.} Conversely, pulling back a
family on $Q$ and applying this quotient returns the original family:
the map $[(x,e)]\mapsto e$ is a bijection on each fiber, with
surjectivity supplied by $q$. The displayed maps are natural, establishing
the equivalence.
  \item \textbf{Structure of perspective coalgebras.} A $G_q$-coalgebra is a family $c_x:E_x\to\prod_{q(y)=q(x)}E_y$. Its counit law says that the $x$-component of $c_x$ is the identity. Its coassociativity says that the component maps from $E_x$ to $E_y$
and then $E_z$ compose to the component from $E_x$ to $E_z$.
  \item \textbf{Equivalence with groupoid actions.} Thus its components are precisely the transports $\tau_{yx}$ above;
coalgebra morphisms are again equivariant maps.
  \item \textbf{Conclusion of the calculation: completion.} This proves the remaining
equivalence. The slice adjunctions follow the categorical-logical account of \citet{Jacobs1999}; the fiber formulas give the descent used here.
\end{enumerate}
\end{proof}


\[
\begin{tikzcd}[column sep=large]
 \mathbf{Set}/Q\arrow[r,shift left=1.3ex,"q^*"]&
 \mathbf{Set}/X\arrow[l,shift left=1.3ex,"q_!",swap]
 \arrow[l,bend left=42,"q_*"']
\end{tikzcd}
\qquad
\begin{tikzcd}[column sep=large]
 \operatorname{Alg}(T_q)\arrow[r,"\simeq"]&
 \operatorname{Act}(R_q)\arrow[r,"\simeq"]&\mathbf{Set}/Q
\end{tikzcd}
\]

The same observable information can have different raw-history
presentations. Suppose $X\twoheadrightarrow Q$ and
$Y\twoheadrightarrow Q$ present the quotient $Q$. The pullback of
pairs $(x,y)$ with the same image gives the comparison bibundle. Its
principal actions ensure that policies and dependent families descended
from $Q$ do not depend on a chosen history presentation; no additional
signal becomes available to an agent.

\begin{definition}[Principal comparison and Morita equivalence]
\label{def:egt-principal-comparison-and-morita-equivalence}
Fix groupoids $G$ and $H$ in $\mathbf{Set}$. A
$(G,H)$-\textbf{bibundle} consists of anchor maps
$G_0\xleftarrow{\alpha}P\xrightarrow{\beta}H_0$ and commuting
left and right actions, typed as
\[
 G_1\times_{s,\alpha}P\longrightarrow P,\quad(g,p)\mapsto gp,
 \qquad
 P\times_{\beta,t}H_1\longrightarrow P,\quad(p,h)\mapsto ph.
\]
Each action satisfies its groupoid unit and composition laws and
preserves the opposite anchor. It is \textbf{biprincipal} when both
anchors are surjective and both canonical maps
\[
 \begin{aligned}
 G_1\times_{s,\alpha}P&\xrightarrow{\ \cong\ }P\times_{\beta}P,
   &(g,p)&\longmapsto(gp,p),\\
 P\times_{\beta,t}H_1&\xrightarrow{\ \cong\ }P\times_{\alpha}P,
   &(p,h)&\longmapsto(p,ph)
 \end{aligned}
\]
are bijective. Thus, over either anchor, two points differ by a unique
arrow of the other groupoid. \textbf{Morita equivalence} here means an
equivalence implemented by such a bibundle. The balanced product
identifies $(p\cdot h,p')$ with $(p,h\cdot p')$ wherever defined.
All these groupoids are set-based; compare \citet{Landsman2001}.
\end{definition}
\begin{theorem}[Morita invariance of effective presentations]
\label{thm:egt-morita-presentations}
For surjections $X\xrightarrow qQ\xleftarrow rY$, the set
$P=X\times_QY$ is a biprincipal $(R_q,R_r)$-bibundle. Moreover
\[
 P\otimes_{R_r}P^{\mathrm{op}}\cong R_q,
 \qquad
 P^{\mathrm{op}}\otimes_{R_q}P\cong R_r,
 \qquad
 \operatorname{Act}(R_q)\simeq\mathbf{Set}/Q
 \simeq\operatorname{Act}(R_r).
\]
These equivalences preserve data descended from $Q$.
\end{theorem}
\begin{proof}
\leavevmode
\begin{enumerate}[label=\textup{(\arabic*)},leftmargin=*,itemsep=4pt,topsep=3pt]
  \item \textbf{Lifting along the surjection.} The anchors are the two projections. The actions replace the $X$ or $Y$
coordinate along the corresponding kernel arrow and therefore commute. The anchors are surjective by pullback stability of surjections.
  \item \textbf{Monoidal comparison.} Over a
fixed $y$, the unique arrow from $(x,y)$ to $(x',y)$ is $(x',x)$;
the other principality equation is identical with the roles reversed. The first balanced product sends the class of
$((x,y),(y,x'))$ to $(x,x')$. Each admissible pair $(x,x')$ has such a
representative because $r$ is surjective.
  \item \textbf{Monoidal comparison: completion.} Any two choices of $y$ differ
by the unique kernel arrow and hence give the same balanced class. This
proves the first bijection, including equivariance; the second is symmetric. The equivalences of action categories follow from
Theorem~\ref{thm:egt-descent-algebras}. Both identify pulled-back data
with the same family on $Q$, which proves the last assertion.
\end{enumerate}
\end{proof}

\noindent The comparison bibundle is the pullback of the two quotient presentations.
\begin{center}
\begin{tikzcd}[column sep=large,row sep=large]
X\times_QY\arrow[r]\arrow[d]\arrow[dr,phantom,"\lrcorner",very near start]
 &Y\arrow[d,"r",two heads]\\
 X\arrow[r,"q"',two heads]&Q
\end{tikzcd}
\end{center}

\[
\begin{tikzcd}[column sep=large]
 &X\times_QY\arrow[dl,two heads]\arrow[dr,two heads]&\\
 X\arrow[dr,two heads,"q"']&&Y\arrow[dl,two heads,"r"]\\
 &Q&
\end{tikzcd}
\]

\begin{remark}[What the comparison retains]
A semantic realization $Q\to\mathcal C(A,B)$, an observation on $Q$,
or a predicate on $Q$ pulls back to both presentations. Data not constant
on the quotient fibers remain on the larger presentation. In particular,
economic invariance also requires the deviation comparisons constructed in
Section~\ref{sec:egt-reduction}; it does not follow from equivalence of
unadorned groupoids alone.
\end{remark}

For a response term $a:W\to A$ and observation $o:W\to O$,
descent along $o$ says that $a$ is constant on the kernel of $o$. Its unique
factor is $\bar a$, and tests satisfy
$a^*P=o^*\bar a^*P$. This factorization through $o$ also holds for dependent
policies in the corresponding slice; constant types place it in the
state context of an owned arena. This identifies the mathematical
meaning of a policy's dependence on available information.

\begin{proposition}[Effective descent of dependent linear interfaces]
\label{prop:egt-linear-descent}
Let $q:X\twoheadrightarrow Y$. A family $A_x\in\widehat{\mathcal C}$
with natural isomorphisms $\theta_{x'x}:A_x\to A_{x'}$ whenever
$q(x)=q(x')$, satisfying identity and composition, descends to a family
$\bar A_y$, uniquely up to its canonical isomorphism. Equivariant
families of arrows descend as well.
\end{proposition}
\begin{proof}
\leavevmode
\begin{enumerate}[label=\textup{(\arabic*)},leftmargin=*,itemsep=4pt,topsep=3pt]
  \item \textbf{Predicate calculation.} For each $y,c$ form the set
\[
 \bar A_y(c)=
 \left(\coprod_{q(x)=y}A_x(c)\right)\big/\bigl(
 (x,a)\sim(x',\theta_{x'x,c}(a))\bigr).
\]
The identity and composition equations make this an equivalence relation. For $v:d\to c$ define restriction by
$[x,a]\mapsto[x,A_x(v)a]$; naturality of $\theta$ makes it
well-defined. The presheaf identities hold on representatives.
  \item \textbf{Functoriality and substitution.} For each
$x$ the map $a\mapsto[x,a]$ is bijective onto $\bar A_{q(x)}(c)$:
its inverse sends $[x',a']$ to $\theta_{xx',c}(a')$. This formula is independent of the representative and is natural in $c$. An equivariant arrow sends $[x,a]$ to the class of its image.
  \item \textbf{Descent along the kernel pair.} The same
formula proves uniqueness of descent and of its canonical comparison
with any other descended family. No representative of a fibre of $q$
is selected. Tensor and the other fibre constructions transport along
these natural isomorphisms; their substitution comparisons are those of
Theorem~\ref{thm:egt-dependent-model}.
\end{enumerate}
\end{proof}

\subsubsection{Observation quotients and local agreement}

\begin{theorem}[Testing the order of information]
\label{thm:egt-information-order}
For two observation maps $o_j:H\twoheadrightarrow O_j$, the following are
equivalent:
\begin{enumerate}[label=\textup{(\roman*)},leftmargin=*]
\item there is a map $q:O_2\to O_1$ with $o_1=qo_2$;
\item $\operatorname{ker}(o_2)\subseteq\operatorname{ker}(o_1)$;
\item for every action set $A$, every response implementable using $o_1$ is
      implementable using $o_2$;
\item the response $o_1:H\to O_1$, induced by the identity policy on $O_1$,
      is implementable using $o_2$.
\end{enumerate}
The factor $q$ is unique and surjective.
\end{theorem}
\begin{proof}
\leavevmode
\begin{enumerate}[label=\textup{(\arabic*)},leftmargin=*,itemsep=4pt,topsep=3pt]
  \item \textbf{Descent along the kernel pair.} Factorization implies kernel inclusion. Conversely, if
$\operatorname{ker}(o_2)\subseteq\operatorname{ker}(o_1)$, then $o_1$
equalizes the kernel pair of $o_2$. Since $o_2$ is the coequalizer of its
kernel pair, there is a unique $q$ with $o_1=qo_2$.
  \item \textbf{Verification of the universal property.} Surjectivity of $q$
follows from surjectivity of $o_1$: for $z\in O_1$, choose $h$ with
$o_1(h)=z$ and observe that $z=q(o_2(h))$.
  \item \textbf{Policy-level comparison.} If $o_1=qo_2$ and $\pi:O_1\to A$, then the same response is implemented
using $o_2$ by the policy $\pi q$, proving the third clause. The third
clause applied to $A=O_1$ and $\pi=\operatorname{id}_{O_1}$ gives the
fourth. The fourth clause is exactly the existence of a map $q$ with
$o_1=qo_2$, so the cycle of implications closes.
\end{enumerate}
\end{proof}

\noindent Information refinement is witnessed by factorization of effective observation maps.
\begin{center}
\begin{tikzcd}[column sep=large,row sep=large]
H\arrow[r,"o_2",two heads]\arrow[dr,"o_1"']
 &O_2\arrow[d,"q",two heads]\\
 &O_1
\end{tikzcd}
\end{center}

\begin{corollary}[Binary tests in the classical finite case]
If $H$ is finite and the finite observation sets are interpreted in ordinary
classical $\mathbf{Set}$, the preceding conditions are also equivalent to:
every binary-valued policy based on $o_1$ is implementable using $o_2$.
\end{corollary}
\begin{proof}
\leavevmode
\begin{enumerate}[label=\textup{(\arabic*)},leftmargin=*,itemsep=4pt,topsep=3pt]
  \item \textbf{Observation and knowledge.} Factorization gives the binary condition by precomposition. Conversely, if
the kernel inclusion fails, take $h,h'$ with equal $o_2$-observations and
distinct $o_1$-observations. The characteristic function of the singleton
$\{o_1(h)\}$ separates the two $o_1$-observations and therefore cannot factor
through $o_2$.
\end{enumerate}
\end{proof}

\begin{proposition}[Gluing observation-based policies]
Let
$H=\bigcup_\alpha H_\alpha$, and suppose local policies
$\pi_\alpha:o[H_\alpha]\to A$ are given. They glue to one admissible
policy iff
\[
 o(h)=o(h'),\ h\in H_\alpha,\ h'\in H_\beta
 \quad\Longrightarrow\quad
 \pi_\alpha(o(h))=\pi_\beta(o(h')).
\]
\end{proposition}
\begin{proof}
\leavevmode
\begin{enumerate}[label=\textup{(\arabic*)},leftmargin=*,itemsep=4pt,topsep=3pt]
  \item \textbf{Observation and knowledge.} Necessity follows by restriction of a common function. For sufficiency,
fix an observation $z\in O$. Because the patches cover $H$ and $o$ is
onto, there is a pair $(\alpha,h)$ with $h\in H_\alpha$ and $o(h)=z$.
  \item \textbf{Gluing observation-based policies.} The displayed agreement shows that all such pairs determine the same
value $\pi_\alpha(z)$. Thus there exists a unique value $a_z\in A$
with this property; define $\pi(z)=a_z$. This pointwise unique-value
construction gives the global function without choosing representatives,
and proves uniqueness.
  \item \textbf{Observation and knowledge.} Agreement only on the literal intersections
$H_\alpha\cap H_\beta$ would be insufficient: disjoint history patches
can still contain observationally indistinguishable histories.
\end{enumerate}
\end{proof}

\subsubsection{Information refinement and complete agent policies}

\begin{construction}[Information refinement and its presentation map]
Suppose $o_1=q\circ o_2$ for observation quotients
$o_j:H\twoheadrightarrow O_j$. Surjectivity of $o_1$ implies that
$q:O_2\twoheadrightarrow O_1$ is surjective. Precomposition gives
\[
 q^*:A^{O_1}\hookrightarrow A^{O_2},\qquad
 \pi\longmapsto\pi\circ q.
\]
It is injective because equal composites with a surjection are equal. The
resulting responses on $H$ are unchanged. Thus old policies embed in the more
informed policy space, preserving their realized behavior.

When $H,O_j,A$ are finite this embedding is itself a $2$-cell of open arenas. Define the single-owner observed decision arena
\[
 \mathsf{ObsDec}_i(o,A):D(H)\longrightarrow D(A)
\]
by
\[
 \Sigma_i=A^O,
 \qquad
 \Sigma_j=1\ (j\ne i),
 \qquad
 \pi\longmapsto\widehat\pi\circ\widehat o
             =\widehat{\pi o}.
\]
If $o_1=qo_2$, functoriality of the response protocols gives
\[
 \widehat{\pi o_1}
 =\widehat{(\pi q)o_2},
\]
so the injective map $q^*:A^{O_1}\hookrightarrow A^{O_2}$ is a
realization-preserving $2$-cell
\[
 \mathsf{ObsDec}_i(o_1,A)
 \Longrightarrow
 \mathsf{ObsDec}_i(o_2,A).
\]
Thus refinement of information is represented by ordinary parameter
reindexing in $\mathbf{Arena}_N(\mathcal C)$, while the observation wire itself
is the explicitly constructed map $\widehat o$.
\end{construction}

A finer observation gives an injection of old policies into a
larger observation-based policy carrier. Because equilibrium tests all
available alternatives, stability against the finer-information choices
implies stability against the original ones at any retained profile, but
not conversely. This is a change of the game, not merely a new reference
state: the new observations permit genuinely new contingent responses.

\begin{construction}[Pooled observation]
\label{con:egt-pooled-observation}
For a coalition $C$ of observers on a common history space, form the product
map
\[
 \langle o_i\rangle_{i\in C}:H\longrightarrow\prod_{i\in C}O_i
\]
and take its image factorization
\[
 H\twoheadrightarrow O_C\lhook\joinrel\longrightarrow
 \prod_{i\in C}O_i.
\]
This is the pooled observation quotient. Its kernel pair is
\[
 \ker(o_C)=\bigcap_{i\in C}\ker(o_i),
\]
because equality in the product is coordinatewise equality. Thus pooled
information is constructed as the effective quotient of the joint signal,
not assumed to fill the whole Cartesian product of observation values.
Pooling information remains distinct from allowing a coalition to coordinate
its choices; either operation may occur without the other.
\end{construction}

\begin{proposition}[Decentralized policies inside the pooled-information interface]
\label{prop:egt-pooled-policy-comparison}
For a nonempty coalition $C$, let
$o_C:H\twoheadrightarrow O_C$ be the pooled observation quotient and
$p_i:O_C\twoheadrightarrow O_i$ its coordinate maps. For action sets $M_i$
there is a canonical injection
\[
 J_C:\prod_{i\in C}M_i^{O_i}
      \lhook\joinrel\longrightarrow
      \left(\prod_{i\in C}M_i\right)^{O_C},
 \qquad
 J_C((\pi_i)_i)(o)=(\pi_i(p_i(o)))_i.
\]
Its image consists exactly of joint policies whose $i$th output is constant
on $\ker(p_i)$ for each $i$. Its composite with the observation $o_C$ is the
original tuple of decentralized responses on $H$.
\end{proposition}
\begin{proof}
\leavevmode
\begin{enumerate}[label=\textup{(\arabic*)},leftmargin=*,itemsep=4pt,topsep=3pt]
  \item \textbf{Lifting along the surjection.} Each $p_i$ is surjective because $p_io_C=o_i$ and $o_i$ is surjective. Consequently equality after precomposition with every $p_i$ implies equality
of every $\pi_i$, proving injectivity. If a joint policy is in the image,
its $i$th component factors through $p_i$ and is constant on its kernel pair.
  \item \textbf{Reverse implication.} Conversely such constancy gives a unique factor $\pi_i:O_i\to M_i$ by
effective descent. Their tuple is the unique preimage under $J_C$. Finally $p_io_C=o_i$ gives
$J_C((\pi_i)_i)o_C=(\pi_io_i)_i$.
\end{enumerate}
\end{proof}

\noindent Decentralized responses coincide with evaluation of the induced pooled policy.
\begin{center}
\begin{tikzcd}[column sep=large,row sep=large]
H\arrow[r,"o_C",two heads]\arrow[d,"(\pi_io_i)_i"']
 &O_C\arrow[d,"J_C((\pi_i)_i)"]\\
 \prod_iM_i\arrow[r,equal]&\prod_iM_i
\end{tikzcd}
\end{center}

Regrouping ownership combines the original decentralized policy
coordinates without giving their owners more information. Delivery of
the pooled signal changes the policy interface itself. The injection
above specifies exactly which pooled policies were already implementable
with separate observations. Finite collector wires realize the delivery,
and their synchronization barrier remains in the concurrent process.
Thus joint control and joint information are independent constructions.

\paragraph{Observation and implementable dependence.}
The quotient equalizer of \cref{prop:egt-observation} characterizes
observation-based responses; \cref{thm:egt-descent-algebras,prop:egt-linear-descent}
extend it to dependent families. Consequently the response
$\pi_i o_i$ uses only the effective signal actually displayed. No
separate admissibility predicate needs to be added to this response
construction.

\subsubsection{Local linear reasoning}
\label{agt-sec-perspective}

Knowledge is a predicate on the states compatible with an
observation. A linear argument may instead carry resources that vary
across those states. For an effective observation
$o:W\twoheadrightarrow O$, strong monoidal reindexing
$o^*:\cL(O)\to\cL(W)$ has right adjoint $\Pi_o^\ell$.
A local family $X$ is indexed by observations, whereas a world family
$A$ is indexed by states. Their adjunction is
\begin{equation}
 \cL(W)(o^*X,A)\cong\cL(O)(X,\Pi_o^\ell A),
 \qquad(\Pi_o^\ell A)_z=\prod_{o(w)=z}A_w.
 \label{agt-eq-perspective-adjunction}
\end{equation}
The right adjoint packages one compatible linear argument for every
world in an observation fibre. In judgment form it transposes
$O;\Delta\vdash u:\Pi_o^\ell A$ and
$W;o^*\Delta\vdash\widetilde u:A$. This is proof-relevant
perspective: the product retains the resources available across
indistinguishable worlds, rather than only the assertion that all such
worlds satisfy a predicate.

\begin{definition}[Proof-relevant perspective]
Define $K_o^\ell=o^*\Pi_o^\ell:\cL(W)\to\cL(W)$.
Its counit is projection at the actual state:
\[
 \epsilon_{A,w}:\prod_{v:o(v)=o(w)}A_v\longrightarrow A_w.
\]
Its comultiplication has, for every $v$ in that fibre, the component
identity into the same product:
\[
 \delta_{A,w}:\prod_{u:o(u)=o(w)}A_u
 \longrightarrow
 \prod_{v:o(v)=o(w)}\prod_{u:o(u)=o(v)}A_u.
\]
\end{definition}

\begin{proposition}[Perspective comonad]
\label{agt-prop-perspective-comonad}
The displayed counit and comultiplication make $K_o^\ell$ a comonad.
\end{proposition}
\begin{proof}
\leavevmode
\begin{enumerate}[label=\textup{(\arabic*)},leftmargin=*,itemsep=4pt,topsep=3pt]
  \item \textbf{Adjunction and transposition.} They are the counit and $o^*\eta\Pi_o^\ell$ of the adjunction
$o^*\dashv\Pi_o^\ell$. Explicitly, either counit composite projects back
to the original indexed family. Both coassociativity composites have, at
an index $(v,u,t)$, the projection to $A_t$. Product uniqueness proves
the equations, including empty products where applicable.
\end{enumerate}
\end{proof}

Unlike the predicate knowledge operator, the linear perspective
comonad need not be idempotent. Reapplying it can retain multiple
additive copies of world-indexed witness products, so its
comultiplication need not be invertible. In particular, the construction
introduces no contraction map $A\to A\otimes A$.

For a concrete distinction in this very model, take two indistinguishable
worlds and the constant family $\Lin 2$. Evaluation at the original unit
gives $(\Lin2)(I)=2$, since $\cC(I,I)$ is a singleton. The first perspective
product has four such values, whereas its iterated product has sixteen.
Thus its comultiplication cannot be an isomorphism. Evaluation at the tensor unit distinguishes the two presheaf objects directly.

\subsubsection{Compatibility with linear composition}

\begin{theorem}[Monoidal perspective comparison]
\label{agt-thm-perspective-monoidal}
There are coherent natural maps
\[
 \Pi_o^\ell A\otimes\Pi_o^\ell B
 \longrightarrow\Pi_o^\ell(A\otimes B),
 \qquad\I_O\longrightarrow\Pi_o^\ell\I_W.
\]
Consequently $K_o^\ell$ is lax symmetric monoidal. For world-indexed
agents there is, in particular, a natural operation
\[
 K_o^\ell\Ag_{\sig}(A,B)\otimes
 K_o^\ell\Ag_{\taus}(B,C)
 \longrightarrow
 K_o^\ell\Ag_{\sig\times\taus}(A,C),
\]
and a corresponding parallel-composition operation.
\end{theorem}
\begin{proof}
\leavevmode
\begin{enumerate}[label=\textup{(\arabic*)},leftmargin=*,itemsep=4pt,topsep=3pt]
  \item \textbf{Observation and knowledge.} At observation $z$, for each $w$ with $o(w)=z$, tensor the two projections
to obtain a map into $A_w\otimes B_w$. Product uniqueness gives the
comparison. The unit comparison is the family of identity maps on $\I$.
  \item \textbf{Associativity.} For associativity, each composite followed by the projection at $w$ is
the tensor of the same three projections, reassociated coherently. Units
and symmetry are checked by the same argument. Naturality follows by
postcomposing the projections with the given component maps.
  \item \textbf{Composition under linear perspective.} Applying $o^*$ gives the comparisons for $K_o^\ell$. Compose the first
comparison with $K_o^\ell(\mathsf{seq})$ to obtain agent composition;
replace $\mathsf{seq}$ by $\mathsf{par}$ for the parallel operation. Their coherence follows from the comparison coherence and
Theorem~\ref{agt-thm-agent-coherence}.
\end{enumerate}
\end{proof}

\noindent The perspective comparison is characterized by its projections at each observationally possible world.
\begin{center}
\begin{tikzcd}[column sep=large,row sep=large]
\Pi_o^\ell A\otimes\Pi_o^\ell B\arrow[r,"m_{A,B}"]\arrow[d,"\pi_w^A\otimes\pi_w^B"']
 &\Pi_o^\ell(A\otimes B)\arrow[d,"\pi_w^{A\otimes B}"]\\
 A_w\otimes B_w\arrow[r,equal]&A_w\otimes B_w
\end{tikzcd}
\end{center}

The lax monoidal map combines two resources that are each available
from the same perspective; it does not create two copies of one agent
or pool observations from different owners. Pooling requires its own
joint observation map and, when implemented, a communication protocol.
This separates composability of perspective-sensitive resources from
access to shared information.

\subsubsection{Perspective coherence and its unit}

The comparison $m_{A,B}:\Pi_oA\otimes\Pi_oB\to\Pi_o(A\otimes B)$
has component $\pi_w^A\otimes\pi_w^B$. Either associativity route,
followed by the $w$-projection, is
$\pi_w^A\otimes\pi_w^B\otimes\pi_w^C$ under tensor reassociation.
The unit route uses $\id_\I$ at every $w$ and then the tensor unitor.
This proves lax monoidal coherence directly. The counit is monoidal because
its projection at the actual world selects precisely these tensor
projections. The comultiplication is compatible because all its repeated
components are the same indexed-family map. These checks explain why
perspective and agent composition are compatible even though the lax
comparison need not be invertible.

\subsubsection{Knowledge as the predicate fragment}

For $P\subseteq W$, define
\[
 K_oP=o^*\forall_oP,
 \qquad
 K_oP(w)\iff\forall v\,(o(v)=o(w)\Rightarrow P(v)).
\]
The accessibility relation is $R_o=\kerpair(o)$. Since it is an equivalence
relation, the following properties are immediate consequences of its fibres.

\begin{proposition}[Epistemic laws]
\label{agt-prop-S5}
$K_o$ preserves finite conjunction and truth, is monotone and factive, and
satisfies $K_oK_oP=K_oP$ and
$\neg K_oP\subseteq K_o\neg K_oP$. It validates the normal modal
distribution rule.
\end{proposition}
\begin{proof}
\leavevmode
\begin{enumerate}[label=\textup{(\arabic*)},leftmargin=*,itemsep=4pt,topsep=3pt]
  \item \textbf{Adjunction and transposition.} The kernel pair $R_o=\ker(o)$ is an equivalence relation. \Cref{prop:egt-observation-interior,thm:egt-invariant-adjunctions}
identify $K_o=o^*\forall_o$ with the universal kernel-pair modality
and prove its factivity, idempotence, and invariant-predicate
characterization. Preservation of conjunction and truth follows from the right-adjoint
universal quantifier.
  \item \textbf{Predicate calculation.} The negation of a fibre-invariant predicate is
fibre-invariant, giving negative introspection. The normal modal
implication rule follows from meet preservation: since
$K_o((P\Rightarrow Q)\wedge P)\leq K_oQ$, one has
$K_o(P\Rightarrow Q)\wedge K_oP\leq K_oQ$, and hence
$K_o(P\Rightarrow Q)\leq(K_oP\Rightarrow K_oQ)$ by the Heyting
implication adjunction. No preservation-of-implication equality is asserted.
  \item \textbf{Restriction to the predicate fragment.} These are the same S5 equations on the ordinary
subterminal fragment, not claims about idempotence of $K_o^\ell$.
\end{enumerate}
\end{proof}

\begin{proposition}[Predicate reflection of linear perspective]
\label{agt-prop-perspective-reflection}
For $P\subseteq W$, let $\widehat P\in\cT(W)$ be the subterminal
ordinary family with fibre $1$ on $P$ and $0$ elsewhere.  In the canonical
completed concurrent model there is a natural isomorphism
\[
 \Val_W K_o^\ell\Lin_W\widehat P
 \cong\widehat{K_oP}.
\]
\end{proposition}
\begin{proof}
\leavevmode
\begin{enumerate}[label=\textup{(\arabic*)},leftmargin=*,itemsep=4pt,topsep=3pt]
  \item \textbf{Adjunction and transposition.} The scalar calculation $\cC(I,I)=1$ gives
$\Val\Lin X\cong X$ in this canonical model. As a right adjoint,
$\Val$ preserves products. Therefore at $w$ the left side is
\[
 \Val\left(\prod_{o(v)=o(w)}\Lin\widehat P_v\right)
 \cong\prod_{o(v)=o(w)}\Val\Lin\widehat P_v
 \cong\prod_{o(v)=o(w)}\widehat P_v.
\]
This subterminal product is inhabited exactly when every world in the
observation fibre satisfies $P$, which is $\widehat{K_oP}_w$.
  \item \textbf{Functoriality and substitution.} The maps preserve all reindexings and predicate inclusions, proving
naturality. The comparison uses the proved $\Val\Lin\cong1$ for the
chosen concurrent completion, not an equation in every abstract LNL
model.
\end{enumerate}
\end{proof}

The predicate reflection keeps exactly the truth information that
is constant across an observation fibre. The proof-relevant linear
perspective instead retains the associated product of witness objects.
Hence its non-idempotence is compatible with the idempotent S5 operator
on assertions.

\subsubsection{Distributed and common knowledge}

For nonempty $C\subseteq N$, factor the tuple of observations through its
image:
\[
 W\xrightarrow{q_C}O_C\hookrightarrow\prod_{i\in C}O_i.
\]
Its kernel is $R_C^{\cap}=\bigcap_{i\in C}R_i$. Define
$\DK_CP=[R_C^{\cap}]P$. This is knowledge of a joint signal. A
decentralized coalition does not automatically possess a policy on $O_C$;
Section~\ref{agt-sec-uniform} constructs the exact comparison.

Let $R_C^*$ be the relation generated by finite paths whose steps belong
to some $R_i$, $i\in C$. It is an equivalence relation. Let
$q_C^*:W\to W/R_C^*$ be its quotient, and define
$\CK_CP=(q_C^*)^*\forall_{q_C^*}P$.

\begin{proposition}[Common knowledge as a greatest fixed point]
\label{agt-prop-common}
\[
 \CK_CP=\nu X.\left(P\wedge\bigwedge_{i\in C}K_iX\right).
\]
\end{proposition}
\begin{proof}
\leavevmode
\begin{enumerate}[label=\textup{(\arabic*)},leftmargin=*,itemsep=4pt,topsep=3pt]
  \item \textbf{Finite-path characterization of knowledge.} Apply the finite-path construction in
\cref{prop:egt-true-common-belief,thm:egt-common-coreflection} to the
subfamily $(\ker(o_i))_{i\in C}$. Its equivalence closure consists
exactly of finite $C$-paths, including paths of length zero. The proof there identifies its box with the greatest post-fixed point
of $X\mapsto P\wedge\bigwedge_{i\in C}K_iX$. For $C=\varnothing$, the path relation is equality and both sides are $P$.
\end{enumerate}
\end{proof}

\subsubsection{Perspective and contextual adjunction rules}

For $f:V\to W$, the dependent-adjoint rules are the transpositions
\[
 \frac{V;\Delta\vdash t:f^*B}{W;\Sigma_f^\ell\Delta\vdash\overline t:B},
 \qquad
 \frac{V;f^*\Delta\vdash t:A}{W;\Delta\vdash\widehat t:\Pi_f^\ell A}.
\]
In the first rule $\Delta$ denotes its tensor object, and the premise is
the corresponding arrow from that object; this notation does not distribute
$\Sigma_f^\ell$ separately over every context entry. In the second,
$f^*$ preserves the tensor context. Both transpositions have inverse
elimination rules and their adjunction $\beta\eta$ equations.
Perspective uses the second rule at an observation $o$.

At the predicate level, for $P\subseteq V$, $Q\subseteq W$,
\[
 \exists_fP\subseteq Q\iff P\subseteq f^*Q,
 \qquad f^*Q\subseteq P\iff Q\subseteq\forall_fP.
\]
These supply existential and universal introduction and elimination.
Knowledge is the definition $K_o=o^*\forall_o$, and action modalities
are the definitions $\exists_s t^*$ and $\forall_s t^*$.
Their monotonicity rules follow from functoriality. Factivity,
introspection, and reduction equations are the derived results of the
main text, not operations acting on arbitrary linear resources.

Common knowledge has the unfolding equation
$\CK_CP=P\wedge\bigwedge_{i\in C}K_i\CK_CP$ and the rule
\[
 \frac{X\subseteq P\qquad X\subseteq K_iX\ \text{for all }i\in C}
      {X\subseteq\CK_CP}.
\]
The finite-path construction proves both. Under epistemic update the
relations in these rules are those of the updated model.

\section{Informed Choice and Collective Agency}\label{sec:paper-7}

Legality and information constrain agency in different ways. A legal
profile domain identifies permitted choices, whereas an observation
quotient identifies histories at which the same response must be used.
Worldwise feasibility alone does not produce an informed policy: the
chosen action must be legal at every world compatible with the observed
signal. We construct these uniform legal policies by dependent universal
images, then compare decentralized and pooled choices and derive the
coalitional effectivity and stability predicates. The distinction between
control and ability connects with coalition effectivity
\citep{Pauly2002}.

\subsection{Institutionally available agent interfaces}

For the legal inclusion $m:L\hookrightarrow\Sigma$, write
$L_S=\operatorname{im}(\pi_Sm)$ and $e_S:L\twoheadrightarrow L_S$.
These image-factorized coordinate maps specify the available outsider
indices used here; their full functorial construction is given in
Section~\ref{sec:egt-coalitions}.

\begin{definition}[Legal domains and legal choice fibres]
\label{def:egt-legal-domain}
A \textbf{legal domain} is a specified subobject $L\hookrightarrow\Sigma$
of profiles permitted by the institution. Against outsiders $\sigma_{-i}$,
the \textbf{legal choice fibre} is
\[
 S_i^L(\sigma_{-i})=\{a_i\in\Sigma_i:(a_i,\sigma_{-i})\in L\}.
\]
\end{definition}

Legality can couple owners' choices even when their unrestricted
policy space is a product. The set available to one owner may therefore
depend on the outsiders' coordinates. Pulling back the legal subobject
along a fixed outsider profile constructs this dependent menu, rather than
mistaking its coordinate image for an independently feasible choice.

\begin{construction}[Legal selection fibers]
\label{con:cgt-dependent-agency}
For a reference $\sigma\in L$ and owner $i$, put
\[
 S_i^L(\sigma_{-i})
 =\{a\in\Sigma_i:(\sigma_{-i},a)\in L\}.
\]
Let $j^L_{i,\sigma}:S_i^L(\sigma_{-i})\to L$ insert the fixed
outsider coordinates. If $m:L\hookrightarrow\Sigma$ is the inclusion,
the agent's legal feedback is the costate transported through
$\mathbb L(gmj^L_{i,\sigma})$. Connect the agent selection relation
$\varepsilon_{S_i^L(\sigma_{-i})}$ to this dependent interface.
For a coalition $C$, replace $i$ by $C$ and use the corresponding
joint choice fiber. All legally admissible changes of its members are
then made at one coalition interface.
\end{construction}

\begin{proposition}[Legal agency and logical equilibrium]
\label{prop:cgt-dependent-selection}
The conjunction of legal agent selections at a profile $\sigma\in L$
is exactly the legal Nash predicate constructed below. More generally,
the coalition interface has paired realization
$\delta_C^L:D_C^L\to\mathcal C(A,B)^2$, where
$D_C^L=\ker(e_{-C}:L\to L_{-C})$ and $L_{-C}$ is the image of the
outsider-coordinate map. Its feedback is $(\delta_C^L)^*P$.
\end{proposition}
\begin{proof}
\leavevmode
\begin{enumerate}[label=\textup{(\arabic*)},leftmargin=*,itemsep=4pt,topsep=3pt]
  \item \textbf{Ownerwise alternatives.} The fiber of the first projection $D_C^L\to L$ at $\sigma$ consists
exactly of the legal profiles with those outsider coordinates. Inserting
the alternative coalition coordinates is a bijection from
$S_C^L(\sigma_{-C})$ to that fiber. Both descriptions consequently test
$P(g_\sigma,g_\tau)$ on the same pairs.
  \item \textbf{Strategic comparison.} Universal quantification of its
negation is precisely the selection relation. For $C=\{i\}$, conjoin
over $i$ to obtain legal Nash equilibrium.
\end{enumerate}
\end{proof}

\noindent Legal choice fibres are pullbacks of the realizable outsider-coordinate map.
\begin{center}
\begin{tikzcd}[column sep=large,row sep=large]
S_C^L(u)\arrow[r,"j_{C,u}^L"]\arrow[d]\arrow[dr,phantom,"\lrcorner",very near start]
 &L\arrow[d,"e_{-C}"]\\
 1\arrow[r,"u"']&L_{-C}
\end{tikzcd}
\end{center}

\[
\begin{tikzcd}[column sep=large,row sep=large]
 S_C^L(\sigma_{-C}) \arrow[r,"j^L_{C,\sigma}"] \arrow[d]
 \arrow[dr,phantom,very near start,"\lrcorner"] &
 L \arrow[d,"e_{-C}"]\\
 1 \arrow[r,"\sigma_{-C}"'] & L_{-C}
\end{tikzcd}
\]
The pullback fixes outsiders and retains exactly the legal
alternative choices at that reference. Effectivity has a different
quantifier order: a coalition first chooses a realizable coordinate and
then requires the goal at every legal completion. The same legal map
supports both operations, but they answer different agency questions.

Coalitional power pulls an outcome goal back to legal profiles,
quantifies universally over completions of a proposed coalition choice,
and then existentially over realizable choices. The base of the final
quantifier is the image $L_C$, which removes choices without any legal
completion. Thus a formally vacuous universal predicate cannot be
mistaken for an enforceable goal.

\subsection{Institutional arrangements and collective agent interfaces}

\begin{definition}[Institutional arrangement on an agent process]
\label{def:egt-institutional-agent-arrangement}
An \textbf{institutional arrangement} on a situated agent process is its
displayed legal subobject, observation interfaces, and comparison or outcome
terms, together with the concurrent wiring that implements its specified
communication and outcome rules. A \textbf{collective agent interface} for a
nonempty coalition $C\subseteq N$ is obtained by regrouping its controlled
policy coordinates.
Its available legal carrier against outsiders $u$ is $S_C^L(u)$.
When the arrangement supplies a collective observation and comparison, these
are attached as the situated and evaluated interfaces of that collective
agent. The constituent agents' interfaces remain recorded by the original
indexing and the regrouping map. The empty coalition has the unit control
carrier $\Sigma_\varnothing=1$ and identity alternatives; at a legal
reference its legal choice fibre is a singleton. This is the nullary control
interface, requiring no nonempty block in a regrouping partition.
\end{definition}

For nonempty $C$, the partition with block $C$ and singleton outsiders has
control given by the regrouping of Theorem~\ref{thm:egt-owner-regrouping}.
A particular collective
comparison constructed from the members is
\[
 P_C(f,f')=\bigwedge_{i\in C}P_i(f,f').
\]
The collective selection interface tests legal joint deviations
against the displayed conjunction of members' comparisons. Its
construction uses the choice-fibre bijection of
Proposition~\ref{prop:cgt-dependent-selection}; different rules for
collective preference are not inferred from regrouping. This is the
comparison later used to distinguish collective stability from ability
to secure an outcome.

For independent arrangements, the legal domain of the composite is
$L_G\times L_H\hookrightarrow\Sigma^G\times\Sigma^H$.
Additional institutional rules are further comprehensions, and
simultaneous rules intersect. These operations remain coherent with
product reassociation. They also reveal why the legal domain of a
composite cannot in general be read off from its separate marginal
projections: a joint rule may couple otherwise available decisions.

\subsection{Quantified spans and their action}

\begin{definition}[Relation modalities]
\label{def:egt-relation-modalities}
For a relation presented by its jointly monic span
$X\xleftarrow{p}R\xrightarrow{q}Y$, put
\begin{equation}\label{eq:egt-span-modalities}
 [R]=\forall_pq^*,\qquad\langle R\rangle=\exists_pq^*.
\end{equation}
The converse is denoted $R^\smile$. Relational composition $R;S$ first
forms $R\times_Y S$ and then takes its image in the two external factors.
\end{definition}

\begin{proposition}[Internal typing of an alternative span]
\label{prop:fib-span-typing}
Let $s:X\to\Sigma^G$ be a reference context and
$X\xleftarrow{p}R\xrightarrow{q}Y$ a span of alternatives.
Its source leg is the context map
$p:(R,sp)\to(X,s)$; its target is the term
$(sp,q):(R,sp)\to\Delta_GY$. For $P\hookrightarrow Y$,
\[
 [R]P=\forall_p(sp,q)^*\Delta_GP,
 \qquad \langle R\rangle P=\exists_p(sp,q)^*\Delta_GP.
\]
Thus both modalities are operations of $\mathsf P_G$ even when the
target alternative changes the reference strategy.
\end{proposition}
\begin{proof}
\leavevmode
\begin{enumerate}[label=\textup{(\arabic*)},leftmargin=*,itemsep=4pt,topsep=3pt]
  \item \textbf{Predicate calculation.} Projection of $(sp,q)$ to $\Sigma^G$ is $sp$, so it is a slice
arrow. Equation~\eqref{eq:fib-constant-test} identifies its predicate
pullback with $q^*P$. The two quantifiers along $p$ then give precisely
Equation~\eqref{eq:egt-span-modalities}. No equality between the
reference of an alternative and the original reference is required.
\end{enumerate}
\end{proof}

\begin{theorem}[The adjunction of a relation]
\label{thm:egt-relation-adjunction}
For $R\subseteq X\times Y$, put
$R_!=\exists_qp^*:\mathsf{Pred}(X)\to\mathsf{Pred}(Y)$.
Then
\[
 R_!=\langle R^\smile\rangle\dashv[R]=\forall_pq^*.
\]
Consequently $[R]R_!$ is an idempotent monad on $\mathsf{Pred}(X)$
and $R_![R]$ is an idempotent comonad on $\mathsf{Pred}(Y)$.
\end{theorem}
\begin{proof}
\leavevmode
\begin{enumerate}[label=\textup{(\arabic*)},leftmargin=*,itemsep=4pt,topsep=3pt]
  \item \textbf{Relational adjunction and idempotence.} For $P\in\mathsf{Pred}(X)$ and $Q\in\mathsf{Pred}(Y)$,
\[
 \exists_qp^*P\leq Q
 \iff p^*P\leq q^*Q
 \iff P\leq\forall_pq^*Q.
\]
The unit and counit supply the monad and comonad inequalities;
monotonicity gives the reverse idempotence inequalities. The converse
relation interchanges $p$ and $q$, proving the first identity.
\end{enumerate}
\end{proof}

\begin{remark}[Direction of adjunction and qualitative belief]
The left adjoint of $[R]$ is the diamond of the converse relation.
When $R$ is symmetric this is $\langle R\rangle$ itself. An arbitrary
belief operator $[R]$ is not thereby a comonad: for a nonempty state set
and the empty relation, $[R]\bot=\top$, so no counit
$[R]\leq\operatorname{id}$ exists. The universally available comonad
is the composite $R_![R]$ on the indicated carrier.
\end{remark}

\begin{proposition}[Relation action and graph modalities]
\label{prop:egt-graph-modalities}
One has $[R;S]=[R][S]$ and
$\langle R;S\rangle=\langle R\rangle\langle S\rangle$.
For the graph $\Gamma_f$ of $f:X\to Y$,
\[
 [\Gamma_f]=\langle\Gamma_f\rangle=f^*,\qquad
 [\Gamma_f^\smile]=\forall_f,\qquad
 \langle\Gamma_f^\smile\rangle=\exists_f.
\]
In particular,
\begin{equation}\label{eq:egt-kernel-modality-factorization}
 [\ker f]=f^*\forall_f,\qquad
 \langle\ker f\rangle=f^*\exists_f.
\end{equation}
\end{proposition}
\begin{proof}
\leavevmode
\begin{enumerate}[label=\textup{(\arabic*)},leftmargin=*,itemsep=4pt,topsep=3pt]
  \item \textbf{Witness and legality comparison.} Expanding a composite box gives
$\forall y\,(R(x,y)\Rightarrow
\forall z\,(S(y,z)\Rightarrow P(z)))$, which is equivalent to the box
for the image relation. The existential calculation combines the two
witnesses. For a graph the sole forward candidate is $f(x)$; its converse
has exactly the fiber $f^{-1}(y)$. Finally
$\ker f=\Gamma_f;\Gamma_f^\smile$.
\end{enumerate}
\end{proof}

\[
\begin{tikzcd}[column sep=large]
 &R\times_Y S\arrow[dl]\arrow[dr]&\\
 R\arrow[d,"p"']\arrow[dr,"q"]&&S\arrow[dl,"r"']\arrow[d,"s"]\\
 X&Y&Z
\end{tikzcd}
\]

\subsection{Informed legal policies as sections of a dependent product}
\label{sec:egt-informed-legal-policies}

The individual choice fibre is not yet an informed policy. At the
same observation, several worlds may have different feasible menus, so
a single response must lie in their intersection. The right adjoint to
inverse image constructs that universally feasible menu over each
observation. Taking its sections then produces precisely the policies
that are simultaneously legal and observation-based.

\begin{construction}[The informed legal interface]
\label{con:egt-informed-legal-interface}
Let $o:W\twoheadrightarrow O$ be an effective observation quotient,
$M$ an action carrier, and $F\hookrightarrow W\times M$ the specified
feasibility predicate. Put $u=o\times\operatorname{id}_M$ and define
\[
 A_o=\forall_u F\hookrightarrow O\times M,\qquad
 A_o(z)=\{a\in M:\forall w\in o^{-1}(z)\;(w,a)\in F\}.
\]
The informed legal policy carrier is the set of sections
\[
 \mathsf{Pol}_o(F)
 =\operatorname{Sect}_O(A_o\to O)
 =\{\pi:O\to M:\forall w\in W\;(w,\pi(o(w)))\in F\}.
\]
Thus ordinary descent imposes dependence on observations, and dependent
universal quantification imposes legality at every world with that
observation. The predicate $F$ may describe physical availability,
institutional permission, or their intersection; these interpretations
remain explicit in the supplied data.
\end{construction}

\begin{proposition}[Universal characterization and information refinement]
\label{prop:egt-informed-legal-universal}
For every $\pi:O\to M$, writing $\Gamma_\pi$ for its graph,
\[
 \Gamma_\pi\subseteq A_o
 \quad\Longleftrightarrow\quad u^*\Gamma_\pi\subseteq F.
\]
In particular, $\mathsf{Pol}_o(F)$ is the maximal subobject of $M^O$
whose evaluated responses satisfy $F$.

If $o=q\,o'$ for surjective observation maps
$o':W\twoheadrightarrow O'$ and $q:O'\twoheadrightarrow O$, then
\[
 A_o=\forall_{q\times\operatorname{id}_M} A_{o'},\qquad
 q^*:\mathsf{Pol}_o(F)\lhook\joinrel\longrightarrow
             \mathsf{Pol}_{o'}(F),\quad \pi\longmapsto\pi q.
\]
Its image consists exactly of the fine-observation legal policies
constant on the fibres of $q$, and it preserves their responses on $W$.
\end{proposition}
\begin{proof}
\leavevmode
\begin{enumerate}[label=\textup{(\arabic*)},leftmargin=*,itemsep=4pt,topsep=3pt]
  \item \textbf{Adjunction and transposition.} The first equivalence is the adjunction
$u^*\dashv\forall_u$. The inverse-image graph consists of exactly
$(w,\pi(o(w)))$, so it is contained in $F$ exactly when that response
is feasible at every $w$. This proves both the section formula and
maximality.
  \item \textbf{Adjunction and transposition: second comparison.} Since $u=(q\times\operatorname{id}_M)
(o'\times\operatorname{id}_M)$, composition of right adjoints gives
the displayed equality for $A_o$. Equivalently, the intersection
of feasible menus over $o^{-1}(z)$ is the intersection of the
corresponding fine-observation menus over $q^{-1}(z)$. If $\pi$ is legal, $\pi q$ has exactly the same world response,
so is legal.
  \item \textbf{Verification of the universal property.} Surjectivity of $q$ makes precomposition injective. A fine policy constant on $\ker q$ descends uniquely to $\pi$;
its identical world response satisfies $F$, proving the image claim.
\end{enumerate}
\end{proof}

\begin{proposition}[Compatibility of admissibility, legality, and knowledge]
\label{prop:egt-uniform-knowledge-square}
For the effective observation $o:W\twoheadrightarrow O$ and
$F\hookrightarrow W\times M$ above, let
$F_a=\{w\in W:(w,a)\in F\}$. The following square of sets is a
pullback, in which both horizontal arrows and the right vertical arrow
are inclusions or injections:
\[
\begin{tikzcd}[column sep=large,row sep=large]
\mathsf{Pol}_o(F)\arrow[r,hook]\arrow[d,"o^*"']
  &M^O\arrow[d,"o^*",hook]\\
\{f:W\to M:\forall w\ F(w,f(w))\}\arrow[r,hook]
  &M^W.
\end{tikzcd}
\]
Furthermore, for each $a\in M$ the predicate obtained by pulling the
universally legal menu back to the history carrier is knowledge of
feasibility of that action:
\[
 (o\times\id_M)^*A_o(w,a)
 \quad\Longleftrightarrow\quad K_o(F_a)(w).
\]
Thus the same effective observation quotient simultaneously determines
which responses descend, which descended responses are uniformly legal,
and which individual-action feasibility predicates are known.
\end{proposition}
\begin{proof}
\leavevmode
\begin{enumerate}[label=\textup{(\arabic*)},leftmargin=*,itemsep=4pt,topsep=3pt]
  \item \textbf{Descent of responses.} By
\cref{prop:egt-observation}, precomposition with $o$ embeds $M^O$
into $M^W$ and its image is the set of functions constant on the
kernel pair of $o$. This is the right vertical map.
  \item \textbf{Legal pullback.} The inverse image under this embedding of
the lower-left set consists exactly of those $\pi$ for which
$F(w,\pi(o(w)))$ holds for every $w$. By
\cref{prop:egt-informed-legal-universal}, this set is
$\mathsf{Pol}_o(F)$; hence the square is cartesian.
  \item \textbf{Knowledge of the universally legal menu.} The definition
of $A_o$ gives
\[
 A_o(o(w),a)\iff
 \forall v\,(o(v)=o(w)\Rightarrow F(v,a)),
\]
which is exactly the kernel-pair knowledge predicate
$K_oF_a(w)$ of \cref{def:egt-observations-indistinguishability-and-knowledge}.
All identities are over the original observation and therefore commute
with substitution of a state reference by the relevant pullback square.
\end{enumerate}
\end{proof}

\noindent The dependent universal image identifies the observation-uniform legal actions.
\begin{center}
\begin{tikzcd}[column sep=large,row sep=large]
W\times M\arrow[r,"o\times\id"]\arrow[d,"\pi_W"']
 &O\times M\arrow[d,"\pi_O"]\\
 W\arrow[r,"o"']&O
\end{tikzcd}
\end{center}

For example, let $W=\{w_0,w_1\}$ have one observation and
$M=\{a,b\}$, with feasible menus $\{a\}$ and $\{b\}$ at the
two worlds. Then $A_o$ has empty fibre and there is no informed
legal policy. Under the fully revealing observation, the policy
$w_0\mapsto a$, $w_1\mapsto b$ is legal. The empty universal image exhibits the obstruction to a coarse-observation legal policy.

The construction applies directly to a situated agent. For
$s:W\to L$, take $M=\Sigma_i$ and
\[
 F_i(w,a)\iff(s(w)_{-i},a)\in L.
\]
The sections are complete-policy deviations that depend only on the
agent's observation and remain legal at every compatible reference
state. When the decision is a local action rather than a complete policy,
its own action set and enabling predicate must be used instead.
The distinction is fixed by the causal stage of the protocol.

\begin{proposition}[Decentralized legal coalition policies]
\label{prop:egt-decentralized-legal}
For observers $o_i:W\twoheadrightarrow O_i$, a nonempty coalition
$C$, and a joint feasibility predicate
$F_C\hookrightarrow W\times\prod_{i\in C}M_i$, let
$o_C:W\twoheadrightarrow O_C$ be the pooled quotient. Then
\[
 \mathsf{Pol}^{\mathrm{dec}}_C(F_C)
 =
 J_C^{-1}\!\left(
   \operatorname{Sect}_{O_C}
   \bigl(\forall_{o_C\times\operatorname{id}}F_C\to O_C\bigr)
          \right)
 \hookrightarrow\prod_{i\in C}M_i^{O_i},
\]
where $J_C$ is Proposition~\ref{prop:egt-pooled-policy-comparison}.
The inclusion into pooled legal policies is injective, and its
image is exactly those pooled legal policies whose $i$th output
descends along $p_i:O_C\to O_i$ for every $i$.
\end{proposition}
\begin{proof}
\leavevmode
\begin{enumerate}[label=\textup{(\arabic*)},leftmargin=*,itemsep=4pt,topsep=3pt]
  \item \textbf{Policy-level comparison.} The response at $w$ of $J_C((\pi_i)_i)$ is
$(\pi_i(o_i(w)))_i$. By the section characterization, this
pooled policy is legal exactly when that tuple lies in $F_C(w)$
at every $w$. This is precisely decentralized joint feasibility.
  \item \textbf{Characterization of decentralized legal policies.} The image characterization and injectivity follow by restricting
the injection and image characterization of $J_C$ to this
comprehension.
\end{enumerate}
\end{proof}

\noindent Decentralized legal policies embed in the globally pooled admissible-policy space.
\begin{center}
\begin{tikzcd}[column sep=large,row sep=large]
\mathsf{Pol}^{\mathrm{dec}}_C(F_C)\arrow[r,hook]\arrow[d,hook]
 &\mathsf{Pol}_{o_C}(F_C)\arrow[d,hook]\\
 \prod_iM_i^{O_i}\arrow[r,"J_C"',hook]
 &(\prod_i M_i)^{O_C}
\end{tikzcd}
\end{center}

Separate marginal feasibility does not imply joint legal
feasibility. If two binary actions must agree, the joint possibilities
$(0,0)$ and $(1,1)$ give both owners nonempty marginal menus, but the
product also admits the illegal tuple $(0,1)$. The dependent joint
predicate retains this coupling and is therefore needed for informed
coalition choice.

The resulting interface is the deterministic policy structure of
team decision problems: each agent chooses from its own information,
while a joint evaluation is attached to the combined response. The
comparison with \citet{Radner1962TeamDecisionProblems} concerns these
information and decision interfaces, not a probability or optimality
principle added to the present qualitative construction.

\subsection{Local states, knowledge, and a uniform ability witness}
\label{sec:egt-local-state-comparison}

\begin{proposition}[Local-state presentation of an observation frame]
\label{prop:egt-interpreted-system-comparison}
The map
\[
 \ell:W\longrightarrow W\times\prod_{i\in N}O_i,\qquad
 \ell(w)=(w,(o_i(w))_i)
\]
is injective. On its image, equality of agent $i$'s local coordinate
pulls back to $\ker o_i$. An observation policy $\pi_i:O_i\to M_i$
evaluates at the local coordinate exactly as the original response
$\pi_io_i$.
\end{proposition}
\begin{proof}
\leavevmode
\begin{enumerate}[label=\textup{(\arabic*)},leftmargin=*,itemsep=4pt,topsep=3pt]
  \item \textbf{Functoriality and substitution.} The first coordinate recovers $w$, proving injectivity. Equality of the $i$th local coordinates is the equation
$o_i(w)=o_i(v)$, and evaluation is substitution in $\pi_i$.
\end{enumerate}
\end{proof}

This yields the static local-state interface of interpreted-systems
semantics \citep{HalpernFagin1989KnowledgeAction}. The carrier $W$
may consist of histories or points of runs supplied by a protocol;
transition relations arise separately from executions or epistemic actions.
Knowing one's own currently selected policy, when modeled, is an
additional observation fact, not a consequence of ownership alone.

To distinguish knowledge of pointwise availability from possession
of one informed ability witness, let $\Pi$ be a specified policy
carrier and $S\hookrightarrow W\times\Pi$ a success predicate.
For example, $S(w,\pi)$ may already universally quantify legal
outsiders' responses. At a current world $w$, compare
\[
 \forall v\;(o_i(v)=o_i(w)\Rightarrow
                    \exists\pi\in\Pi\;S(v,\pi))
\]
with
\[
 \exists\pi\in\Pi\;\forall v\;(o_i(v)=o_i(w)\Rightarrow S(v,\pi)).
\]
The second formula exhibits one policy succeeding across the whole
information class; the first only says that each possible world admits
some successful policy. The implication cannot be reversed. For two
indistinguishable worlds, take policies $\pi_0,\pi_1$ with
$S(w_j,\pi_k)$ precisely when $j=k$. Each world has a successful
policy, but neither policy succeeds at both. The distinction between
knowledge that success is possible and knowing how to secure it is the
point of the comparison with \citet{JamrogaVanDerHoek2004KnowHow}.

\subsection{Legal coalition effectivity and stability}
\label{sec:egt-coalitions}

Regrouping gives coalitions a joint control interface, and legal
comprehension gives each controlled choice its permitted completions.
To determine power, we must now place the quantifiers in the order that
corresponds to a coalition choosing before its outsiders. Image
factorization removes unrealizable choices, while the same kernel-pair
relations continue to express legal deviations. We then compare
coalitional guarantees with preference-based stability.

The collective-control map is Theorem~\ref{thm:egt-owner-regrouping};
legal selection is Proposition~\ref{prop:cgt-dependent-selection}; and
the comparison of decentralized with pooled information is
Proposition~\ref{prop:egt-pooled-policy-comparison}. These supply
separate ingredients of the effectivity calculation. The resulting
relationship between dependence, preference, and coalitional power also
connects with \citet{ChenShiWang2022}.

\subsubsection{Institutional choices and power to secure a goal}

\begin{definition}[Coalition effectivity and enforcement]
\label{def:egt-effectivity-vocabulary}
Let $m:L\hookrightarrow\Sigma$ be a legal profile domain,
$e_C:L\twoheadrightarrow L_C$ the image-factorized projection to the
coalition's coordinates, $q:L\to Z$ an outcome, and
$P\hookrightarrow Z$ a goal. \textbf{Coalition effectivity} is the
predicate-valued operator
\[
 [C]_{L,q}:\mathsf{Pred}(Z)\longrightarrow\mathsf{Pred}(1),
 \qquad [C]_{L,q}(P)=
 \exists_{!:L_C\to1}\,\forall_{e_C:L\to L_C}\,q^*P.
\]
Equivalently, it holds just when some $c\in L_C$ satisfies
$\forall\ell\in L\,(e_C(\ell)=c\Rightarrow q(\ell)\in P)$.
A coalition \textbf{enforces} $P$ by such a realizable coordinate;
every legal completion must satisfy the goal. For fixed $q$, the
outcome subscript is suppressed, giving the notation $[C]_L$ used below. The image $L_C$ excludes
unrealizable choices and therefore excludes vacuous enforcement by them.
Equilibrium and institutional sanctions are evaluated through their respective additional interfaces.
Its state-dependent version is constructed by dependent quantification
in the public-state and announcement sections.
\end{definition}
Here ``enforce'' refers to a quantified guarantee, not to
rationality or preference satisfaction. A coalition enforces a goal when
one realizable choice yields it under every legal completion by outsiders.
The enforcing policy need not be an equilibrium choice, since the
ability predicate does not compare its desirability with other policies.

\begin{construction}[Legal profiles and realizable coordinate images]
Let
\[
 m:L\lhook\joinrel\longrightarrow
 \Sigma=\prod_{i\in N}\Sigma_i
\]
be a specified subobject of legal profiles. The ordinary unconstrained case is
$m=\operatorname{id}_\Sigma$. For every $S\subseteq N$, factor the coordinate
map through its image,
\begin{equation}
\label{eq:egt-legal-coordinate-image}
 L\xrightarrow{e_S}L_S
 \lhook\joinrel\xrightarrow{j_S}\Sigma_S,
 \qquad
 \pi_Sm=j_Se_S.
\end{equation}
Thus $L_S$ is the set of realizable $S$-coordinates and $e_S$ is a
surjection by construction. No legal choice is introduced by selecting a
representative of an equivalence class.
\end{construction}

The coordinate images inherit all coordinate restrictions by effective
descent.

\begin{proposition}[The realizable-coordinate functor]
\label{prop:egt-realizable-coordinate-functor}
For every inclusion $S\subseteq T\subseteq N$ there is a unique surjection
\[
 \rho_S^T:L_T\twoheadrightarrow L_S
\]
satisfying
\begin{equation}
\label{eq:egt-realizable-coordinate-restriction}
 \rho_S^T e_T=e_S.
\end{equation}
These maps obey
\[
 \rho_S^S=\operatorname{id}_{L_S},
 \qquad
 \rho_R^S\rho_S^T=\rho_R^T
 \quad(R\subseteq S\subseteq T).
\]
Hence realizable coordinates form a contravariant functor
\[
 L_{(-)}:\mathcal P(N)^{\mathrm{op}}\longrightarrow\mathbf{Set},
 \qquad S\longmapsto L_S,
\]
where the powerset is ordered by inclusion.
\end{proposition}
\begin{proof}
\leavevmode
\begin{enumerate}[label=\textup{(\arabic*)},leftmargin=*,itemsep=4pt,topsep=3pt]
  \item \textbf{Construction of the comparison.} Let $\pi_S^T:\Sigma_T\to\Sigma_S$ be the ordinary coordinate projection. If $e_T(\ell)=e_T(\ell')$, monicity of the image inclusion
$L_T\hookrightarrow\Sigma_T$ gives
$\pi_Tm(\ell)=\pi_Tm(\ell')$. Applying $\pi_S^T$ gives
$\pi_Sm(\ell)=\pi_Sm(\ell')$, and monicity of
$L_S\hookrightarrow\Sigma_S$ therefore gives
$e_S(\ell)=e_S(\ell')$.
  \item \textbf{Construction by quotient descent.} Thus $e_S$ equalizes the kernel pair of the
surjection $e_T$. Proposition~\ref{prop:egt-kernel-descent} supplies a unique
map $\rho_S^T$ satisfying
Equation~\eqref{eq:egt-realizable-coordinate-restriction}.
  \item \textbf{Lifting along the surjection.} It is surjective because every $s\in L_S$ has the form $s=e_S(\ell)$ and
then
\[
 s=\rho_S^T(e_T(\ell)).
\]
For $R\subseteq S\subseteq T$, both $\rho_R^S\rho_S^T$ and $\rho_R^T$
become $e_R$ after precomposition with the surjection $e_T$; cancellation of
that epimorphism gives equality. The identity equation is the same argument
with $S=T$.
\end{enumerate}
\end{proof}

\begin{construction}[Legal deviation kernel pairs]
The corresponding legal deviation relation is the kernel pair
\begin{equation}
\label{eq:egt-legal-deviation}
 D_C^L=L\times_{L_{-C}}L
 \mathrel{\substack{\xrightarrow{p_2}\\[-2pt]
                     \xrightarrow[p_1]{}}}L
\end{equation}
of $e_{-C}$. It relates exactly those legal profiles that agree on their
realizable outsider coordinate. When $L=\Sigma$ this is the deviation
relation $D_C$ constructed in Equation~\eqref{eq:egt-deviation-kernel}.
\end{construction}

\begin{proposition}[Functional dependence of legal policy coordinates]
\label{prop:egt-functional-dependence}
For $S\subseteq N$ and $j\in N$, form the kernel pair
$K_S=L\times_{L_S}L$ of $e_S$, with projections $p_1,p_2$.
The assertion that the $S$-coordinates locally determine owner $j$'s
coordinate is the predicate
\[
 \mathsf{FD}_{S\Rightarrow j}
 =\forall_{p_1:K_S\to L}
   \operatorname{Eq}\bigl((\pi_jm)p_1,(\pi_jm)p_2\bigr)
 \quad\in\mathsf{Pred}(L).
\]
It holds at $\ell$ exactly when every legal profile agreeing with
$\ell$ on $S$ also agrees on coordinate $j$. Moreover
$\mathsf{FD}_{S\Rightarrow j}=\top_L$ exactly when $\pi_jm$ factors
uniquely through $e_S:L\twoheadrightarrow L_S$.
This is the functional-dependence test associated with the preference-and-
dependence interpretation of \citet{ChenShiWang2022}, expressed in
the existing assertion doctrine.
\end{proposition}
\begin{proof}
\leavevmode
\begin{enumerate}[label=\textup{(\arabic*)},leftmargin=*,itemsep=4pt,topsep=3pt]
  \item \textbf{Descent along the kernel pair.} The fibre of $p_1$ at $\ell$ consists of all legal profiles having
$S$-coordinate $e_S(\ell)$; universal quantification of equality proves
the pointwise claim. Global truth says precisely that
$(\pi_jm)p_1=(\pi_jm)p_2$. Since $e_S$ is the coequalizer of its kernel
pair, Proposition~\ref{prop:egt-kernel-descent} gives the unique factor.
\end{enumerate}
\end{proof}

\begin{proposition}[Legal deviation is dynamic restriction]
\label{prop:egt-legal-deviation-restriction}
For every coalition $C$,
\begin{equation}
\label{eq:egt-legal-deviation-restriction}
 D_C^L\cong(m\times m)^*D_C.
\end{equation}
Thus the legal deviation relation is exactly the ordinary coalition-deviation
relation restricted to the legal state space. The image object $L_{-C}$ does
not introduce a new accessibility relation; it presents the realizable
outsider coordinate through which that restricted relation is the kernel pair.
\end{proposition}
\begin{proof}
\leavevmode
\begin{enumerate}[label=\textup{(\arabic*)},leftmargin=*,itemsep=4pt,topsep=3pt]
  \item \textbf{Descent along the kernel pair.} By definition
$D_C^L=\ker(e_{-C})$. Since
\[
 \pi_{-C}m=j_{-C}e_{-C}
\]
and $j_{-C}$ is monic,
\[
 e_{-C}(\ell)=e_{-C}(\ell')
 \Longleftrightarrow
 \pi_{-C}m(\ell)=\pi_{-C}m(\ell').
\]
Hence
$\ker(e_{-C})=\ker(\pi_{-C}m)$. Kernel pairs are pullbacks of a map along
itself and therefore are stable under change of base.
  \item \textbf{Pullback of deviation kernel pairs.} Pulling
$D_C=\ker(\pi_{-C})$ back along $m\times m$ gives
$\ker(\pi_{-C}m)$, proving the claimed isomorphism.
\end{enumerate}
\end{proof}

\begin{construction}[Legal domains from institutional data]
\label{con:egt-legal-domain}
A legal subobject need not be introduced by an uninterpreted legality
predicate when lower-level institutional data are available. For a set-indexed
family, if $c_\alpha:\Sigma\to Y_\alpha$ are typed constraint maps and
$A_\alpha\hookrightarrow Y_\alpha$ are their permitted-value subobjects,
define
\[
 L=\bigwedge_\alpha c_\alpha^*A_\alpha
 \hookrightarrow\Sigma.
\]
If legality instead means that a profile admits an institutional witness,
start from a witness object $W$ and a map $w:W\to\Sigma$ and define
\[
 L=\operatorname{im}(w)\hookrightarrow\Sigma.
\]
Finite combinations of such conditions are obtained by pullback and image
factorization. Thus the generic subobject $L\hookrightarrow\Sigma$ below is
supplied institutional data only when no lower-level construction has been
specified; whenever constraints or witnesses are supplied, legality is the
subobject constructed from them.
\end{construction}

For an open-arena realization $g:\Sigma\to\mathcal C(A,B)$, the legal
counterfactual semantics of coalition $C$ is
\begin{equation}
\label{eq:egt-legal-counterfactual}
 \delta_C^{G,L}
 =(g\circ m\circ p_1,g\circ m\circ p_2):
 D_C^L\longrightarrow\mathcal C(A,B)^2,
\end{equation}
\begin{construction}[Legal equilibrium]
For a singleton $i$, put
\[
 \mathsf{Prof}_i^L
 =(\delta_i^{G,L})^*\mathsf{Better}_i.
\]
The legal Nash predicate is then
\begin{equation}
\label{eq:egt-legal-ne}
 \mathsf{NE}_{G,L}
 =\bigwedge_{i\in N}
   \forall_{p_1:D_i^L\to L}\neg\mathsf{Prof}_i^L.
\end{equation}
\end{construction}

When $L=\Sigma$, these constructions are exactly the unrestricted
counterfactual semantics and Nash predicate already defined.

\begin{proposition}[Legal equilibrium as restricted modal stability]
\label{prop:egt-legal-ne-modal}
Let
\[
 \mathsf{PDev}_i
 =D_i\cap(g\times g)^*\mathsf{Better}_i
 \hookrightarrow\Sigma\times\Sigma.
\]
Then its legal restriction is
\[
 \mathsf{PDev}_i^L=(m\times m)^*\mathsf{PDev}_i,
\]
and
\begin{equation}
\label{eq:egt-legal-ne-modal}
 \mathsf{NE}_{G,L}
 =\bigwedge_i\neg\langle\mathsf{PDev}_i^L\rangle\top.
\end{equation}
Viewed again as a predicate on the ambient profile space through the dynamic
box $[!m]=\forall_m$,
\begin{equation}
\label{eq:egt-legal-ne-ambient-reduction}
 [!m]\mathsf{NE}_{G,L}
 =L\Rightarrow
   \bigwedge_i\neg\langle\mathsf{PDev}_i\rangle L.
\end{equation}
Thus restricting legality can create equilibria exactly by removing
profitable alternatives; legal equilibrium is not generally the pullback of
unrestricted equilibrium.
\end{proposition}

\begin{proof}
\leavevmode
\begin{enumerate}[label=\textup{(\arabic*)},leftmargin=*,itemsep=4pt,topsep=3pt]
  \item \textbf{Ownerwise alternatives.} Proposition~\ref{prop:egt-legal-deviation-restriction} gives
$D_i^L=(m\times m)^*D_i$, while functoriality of pullback gives
\[
 ((gm)\times(gm))^*\mathsf{Better}_i
 =(m\times m)^*(g\times g)^*\mathsf{Better}_i.
\]
Pullback preserves finite meets, hence the profitable-deviation relation
restricts as displayed. Proposition~\ref{prop:egt-ne-modal-safety} applied in
the state space $L$ gives
Equation~\eqref{eq:egt-legal-ne-modal}.
  \item \textbf{Witness and legality comparison.} Write $A=\exists_m\top$ for the ambient predicate represented by the legal
subobject and put
\[
 S=\bigwedge_i\neg\langle\mathsf{PDev}_i\rangle A
 \in\mathsf{Pred}(\Sigma).
\]
For every relation $R\hookrightarrow\Sigma\times\Sigma$, change of base along
$m$ gives
\begin{equation}
\label{eq:egt-restricted-diamond}
 m^*\langle R\rangle A
 =\langle(m\times m)^*R\rangle\top.
\end{equation}
To see this entirely in the relation calculus, let $\Gamma_m$ be the graph of
$m$.
  \item \textbf{Predicate calculation.} Proposition~\ref{prop:egt-graph-modalities} gives
\[
 m^*=\langle\Gamma_m\rangle,
 \qquad
 A=\exists_m\top=\langle\Gamma_m^{\smile}\rangle\top.
\]
Hence
\begin{align*}
 m^*\langle R\rangle A
 &=\langle\Gamma_m;R;\Gamma_m^{\smile}\rangle\top\\
 &=\langle(m\times m)^*R\rangle\top,
\end{align*}
because $\Gamma_m;R;\Gamma_m^{\smile}$ is exactly the restriction of $R$ to
the subobject $L$.
  \item \textbf{Legal Nash reduction by implication.} Therefore
\[
 m^*S
 =\bigwedge_i\neg\langle\mathsf{PDev}_i^L\rangle\top
 =\mathsf{NE}_{G,L}.
\]
For the monomorphism $m:L\hookrightarrow\Sigma$, each fibre has
one element over a point of $A=\exists_m\top$ and no elements
outside $A$. The universal and existential fibre formulas therefore
give $\forall_m Q=A\Rightarrow\exists_m Q$ for every predicate
$Q$ on $L$: outside $A$ the universal clause is vacuous, and on
$A$ its unique preimage determines both quantifiers. Frobenius gives
\begin{align*}
 [!m]\mathsf{NE}_{G,L}
 &=A\Rightarrow\exists_m\mathsf{NE}_{G,L}\\
 &=A\Rightarrow\exists_m m^*S\\
 &=A\Rightarrow(A\wedge S)\\
 &=A\Rightarrow S.
\end{align*}
Identifying the image predicate $A$ with the displayed legal predicate $L$
gives Equation~\eqref{eq:egt-legal-ne-ambient-reduction}.
\end{enumerate}
\end{proof}


\begin{definition}[Amalgamation of legal choices]
\label{def:egt-legal-amalgamation}
For disjoint coalitions $C,D\subseteq N$, use the restriction maps of
Proposition~\ref{prop:egt-realizable-coordinate-functor} to form the
\textbf{amalgamation comparison}
\[
 \lambda_{C,D}=
 \langle\rho_C^{C\cup D},\rho_D^{C\cup D}\rangle:
 L_{C\cup D}\longrightarrow L_C\times L_D.
\]
Two separately realizable coordinates $c\in L_C$ and $d\in L_D$
\textbf{amalgamate} precisely when the dependent witness fibre
\[
 \mathsf{Amal}_{C,D}(c,d)
 =\{u\in L_{C\cup D}:\lambda_{C,D}(u)=(c,d)\}
\]
is inhabited. All such pairs amalgamate precisely when
$\lambda_{C,D}$ is surjective. A witness is a single realizable joint
coordinate simultaneously extending both choices.
\end{definition}
Coupled legality can prevent individually realizable choices from
being jointly executable. If each of two agents may request a resource but
both requests cannot occur in one legal profile, both requests lie in
their respective coordinate images while their pair has no legal
completion. The amalgamation comparisons below identify this obstruction;
it disappears for unrestricted product domains.

\begin{definition}[Coalition effectivity specialized to an evaluated arena]
\label{def:egt-coalition-effectivity-on-a-legal-domain}
For an open-arena realization $g:\Sigma\to\mathcal C(A,B)$ with an
observable evaluation $o:\mathcal C(A,B)\to Z$, the outcome map in
Definition~\ref{def:egt-effectivity-vocabulary} is
$q=o\circ g\circ m:L\to Z$; it may also be supplied directly.
The already constructed effectivity operator, now written $[C]_L$, is
\begin{equation}
\label{eq:egt-legal-ability}
 [C]_L P
 =\exists_{!:L_C\to1}\,\forall_{e_C:L\to L_C}\,q^*P.
\end{equation}
This agrees with the elementwise enforcing criterion of
Definition~\ref{def:egt-effectivity-vocabulary}: a realizable coalition
coordinate $c\in L_C$ guarantees the goal under every legal completion.
Because $e_C$ is surjective, every quantified coalition choice has at least
one legal completion. Hence ability cannot be obtained vacuously from an
unrealizable coordinate. If $L=\Sigma$ and the profile space is inhabited,
then every factor is inhabited, $L_C\cong\Sigma_C$, and the fiber of $e_C$ is
canonically $\Sigma_{-C}$; Equation~\eqref{eq:egt-legal-ability} becomes
\begin{equation}
\label{eq:egt-ability}
 [C]P\quad\Longleftrightarrow\quad
 \exists s_C\in\Sigma_C\ \forall s_{-C}\in\Sigma_{-C},
       \ q(s_C,s_{-C})\in P.
\end{equation}
For an empty profile space the image construction leaves no realizable
coalition choice and therefore introduces no vacuous power.
\end{definition}

The coalition chooses before the outsiders' simultaneous choice is
known. Exchanging the existential and universal quantifiers would grant a
different kind of power. The image and dependent-quantifier construction
thus makes both the implementable choices and the order of control
explicit in the effectivity predicate.

For disjoint coalitions $C,D$, the realizable-coordinate functor constructs
the comparison
\begin{equation}
\label{eq:egt-choice-amalgamation-map}
 \lambda_{C,D}
 =\bigl\langle\rho_C^{C\cup D},\rho_D^{C\cup D}\bigr\rangle:
 L_{C\cup D}\longrightarrow L_C\times L_D.
\end{equation}
It measures exactly whether separately realizable choices of $C$ and $D$ can
be amalgamated into one legal joint choice.

\begin{theorem}[Coalition ability and legal amalgamation]
\label{thm:egt-superadditive}
Let $C\cap D=\varnothing$. If $\lambda_{C,D}$ is surjective, then for every
outcome map $q:L\to Z$ and predicates $P,Q\hookrightarrow Z$,
\[
 [C]_LP\ \wedge\ [D]_LQ
 \quad\Longrightarrow\quad
 [C\cup D]_L(P\cap Q).
\]
Conversely, if this implication holds for every $q,P,Q$, then
$\lambda_{C,D}$ is surjective. In particular the implication holds for the
ordinary full-product game.
\end{theorem}

\begin{proof}
\leavevmode
\begin{enumerate}[label=\textup{(\arabic*)},leftmargin=*,itemsep=4pt,topsep=3pt]
  \item \textbf{Choice of representatives and quotients.} Let $c\in L_C$ and $d\in L_D$ witness the two ability assertions. If
$\lambda_{C,D}$ is surjective, take $u\in L_{C\cup D}$ with
$\lambda_{C,D}(u)=(c,d)$. For any legal profile $\ell$ with
$e_{C\cup D}(\ell)=u$, functoriality of coordinate restriction gives
$e_C(\ell)=c$ and $e_D(\ell)=d$.
  \item \textbf{Joint coalition witness.} Hence $q(\ell)\in P$ by the first witness
and $q(\ell)\in Q$ by the second. Thus $u$ witnesses
$[C\cup D]_L(P\cap Q)$.
  \item \textbf{Reverse implication.} Conversely, take an arbitrary pair $(c,d)\in L_C\times L_D$. Let
\[
 q=(e_C,e_D):L\longrightarrow L_C\times L_D,
 \qquad
 P=\{c\}\times L_D,
 \qquad
 Q=L_C\times\{d\}.
\]
Then $c$ witnesses $[C]_LP$ and $d$ witnesses $[D]_LQ$. By the assumed
superadditivity there is $u\in L_{C\cup D}$ such that every legal completion
of $u$ is sent to $(c,d)$.
  \item \textbf{Verification of the universal property.} Surjectivity of
$e_{C\cup D}:L\twoheadrightarrow L_{C\cup D}$ supplies such a completion
$\ell$, and
\[
 \lambda_{C,D}(u)
 =\lambda_{C,D}e_{C\cup D}(\ell)
 =(e_C(\ell),e_D(\ell))
 =(c,d).
\]
Hence $\lambda_{C,D}$ is surjective.
  \item \textbf{Lifting along the surjection.} For $L=\Sigma$ a nonempty full product has
$L_S\cong\Sigma_S$, and for disjoint $C,D$ the map $\lambda_{C,D}$ is the
canonical product comparison and is surjective. If $\Sigma$ is empty, every
$L_S$ is empty and the same map is the unique surjection
$\varnothing\to\varnothing$. Thus the ordinary full-product implication
holds in either case.
\end{enumerate}
\end{proof}


For product choices the resulting implication recovers the familiar
superadditivity pattern of coalition effectivity
\citet{Pauly2002}, while the quantifier order parallels the
coalition modalities of alternating-time logic \citep{AlurEtAl2002}.
Here the image comparison identifies the exact amalgamation condition
needed when legality couples the choices.

\begin{construction}[Outsider-coordinate amalgamation map]
There is a parallel exact criterion for composing legal deviation relations.
Put $O=N\setminus(C\cup D)$. Proposition~\ref{prop:egt-realizable-coordinate-functor}
constructs the restriction maps
\[
 \rho_O^{-C}:L_{-C}\to L_O,
 \qquad
 \rho_O^{-D}:L_{-D}\to L_O.
\]
Form their pullback
\[
 P_{C,D}=L_{-C}\times_{L_O}L_{-D}
\]
and the canonical comparison
\begin{equation}
\label{eq:egt-deviation-amalgamation-map}
 a_{C,D}:L\longrightarrow P_{C,D},
 \qquad
 \ell\longmapsto(e_{-C}(\ell),e_{-D}(\ell)).
\end{equation}
\end{construction}

\begin{proposition}[Legal deviation amalgamation]
\label{prop:egt-legal-deviation-amalgamation}
For arbitrary $C,D\subseteq N$,
\[
 D_C^L;D_D^L=D_{C\cup D}^L
 \quad\Longleftrightarrow\quad
 a_{C,D}\text{ is surjective}.
\]
\end{proposition}
\begin{proof}
\leavevmode
\begin{enumerate}[label=\textup{(\arabic*)},leftmargin=*,itemsep=4pt,topsep=3pt]
  \item \textbf{Inclusion of successive legal deviations.} A successive $C$-change and $D$-change preserves every coordinate outside
$C\cup D$, so
$D_C^L;D_D^L\subseteq D_{C\cup D}^L$ always.
  \item \textbf{Forward implication.} Suppose first that $a_{C,D}$ is surjective and
$\sigma D_{C\cup D}^L\tau$. Then the pair
\[
 (e_{-C}(\sigma),e_{-D}(\tau))
\]
has equal images in $L_O$, because $\sigma$ and $\tau$ agree on the
realizable coordinate outside $C\cup D$. It therefore lies in $P_{C,D}$.
  \item \textbf{Choice of representatives and quotients.} Choose $\rho\in L$ above this point under the surjection $a_{C,D}$. Then
\[
 e_{-C}(\rho)=e_{-C}(\sigma),
 \qquad
 e_{-D}(\rho)=e_{-D}(\tau),
\]
so $\sigma D_C^L\rho D_D^L\tau$.
  \item \textbf{Reverse implication.} Conversely assume equality of the relations and take
$(x,y)\in P_{C,D}$. Since the image maps $e_{-C}$ and $e_{-D}$ are
surjective, choose legal profiles $\sigma,\tau$ with
$e_{-C}(\sigma)=x$ and $e_{-D}(\tau)=y$. Compatibility of $(x,y)$ over
$L_O$ says $\sigma D_{C\cup D}^L\tau$.
  \item \textbf{Lifting along the surjection.} By the assumed equality there is
$\rho$ with $\sigma D_C^L\rho D_D^L\tau$. Hence
$a_{C,D}(\rho)=(x,y)$, proving surjectivity.
\end{enumerate}
\end{proof}

\noindent The comparison for successive legal deviations is induced by the pullback of outsider-coordinate images.
\begin{center}
\begin{tikzcd}[column sep=large,row sep=large]
L_{-C}\times_{L_O}L_{-D}\arrow[r,"\mathrm{pr}_2"]\arrow[d,"\mathrm{pr}_1"']\arrow[dr,phantom,"\lrcorner",very near start]
 &L_{-D}\arrow[d,"\rho_O^{-D}"]\\
 L_{-C}\arrow[r,"\rho_O^{-C}"']&L_O
\end{tikzcd}
\end{center}

For the concrete dependent domain
\[
 L=\{(0,0),(1,1)\}\subseteq\{0,1\}^2,
\]
one has $L_{\{1\}}=L_{\{2\}}=\{0,1\}$, while
$\lambda_{\{1\},\{2\}}:L\to\{0,1\}^2$ is the diagonal inclusion and is
not surjective. Taking $q=\operatorname{id}_L$,
$P=\{(0,0)\}$, and $Q=\{(1,1)\}$, agent $1$ can force $P$ with its
realizable choice $0$ and agent $2$ can force $Q$ with its realizable choice
$1$, but their witnesses have no legal amalgamation and
$P\cap Q=\varnothing$. Likewise both singleton deviation relations are
identities whereas the two-player relation is $L\times L$.
The example separates legal dependence from the full-product case.

\subsubsection{Choice cells and the one-stage logic of agency}
\label{sec:egt-choice-cell-logic}

Effectivity asks whether some coalition choice guarantees a goal.
A related one-stage operation asks what is guaranteed by the coalition's
already selected choice. This distinction makes contact with the
seeing-to-it-that interpretation of agency
\citep{HortyBelnap1995DeliberativeStit}: choice and ability quantify
over different fibres even when they use the same legal domain.

\begin{proposition}[Choice-cell guarantees and effectivity]
\label{prop:egt-choice-cell-comparison}
For $e_C:L\twoheadrightarrow L_C$, put
\[
 Q_C=\ker e_C=D_{N\setminus C}^L.
\]
For a predicate $R\hookrightarrow L$, the guarantee of the current
coalition choice is
\[
 [Q_C]R=e_C^*\forall_{e_C}R.
\]
Some realizable coalition choice guarantees $R$ exactly when
$\exists_{!:L_C\to1}\forall_{e_C}R$ holds. At a legal reference
$\ell$, the predicate
\[
 [Q_C]R\ \wedge\ \neg[L\times L]R
\]
says that the current choice guarantees $R$ while $R$ is not
already inevitable among all legal profiles.
\end{proposition}
\begin{proof}
\leavevmode
\begin{enumerate}[label=\textup{(\arabic*)},leftmargin=*,itemsep=4pt,topsep=3pt]
  \item \textbf{Policy-level comparison.} The $Q_C$-successors of $\ell$ are exactly the legal profiles
$\ell'$ with $e_C(\ell')=e_C(\ell)$, which proves the first
equation by the fibre formula for $\forall_{e_C}$. Existentially quantifying over its image $L_C$ gives the
second assertion and excludes unrealizable choices. The universal relation $L\times L$ tests $R$ at every legal
profile, giving the final conjunction.
  \item \textbf{Complementary choice-cell relation.} Finally, fixing the
$C$-coordinates is exactly allowing changes of their complement,
which proves the equality with $D_{N\setminus C}^L$.
\end{enumerate}
\end{proof}

The current-choice modality ranges over outsiders' legal
completions; profitable-deviation modalities instead range over changes
to the owner's own policy with outsiders fixed. Interchanging these
relations would reverse the control interpretation. Their conjunction in
the one-stage deliberative clause therefore records two genuinely
different tests.

For unrestricted products, independently fixed choices of different
owners always have a common realization, given by their coordinate
tuple. Coupled legality need not have this property: with
$L=\{(0,0),(1,1)\}$, choosing $0$ for one owner and $1$ for the
other produces no joint legal profile. Compatibility of choices must
therefore be established from the legal domain, not from ownership
alone.

\subsubsection{Collective improvements and stability}

\begin{definition}[Blocking and cores]
\label{def:egt-blocking-and-core}
Let $L$ be a reference carrier. For each nonempty coalition
$C\subseteq N$, a specified \textbf{blocking-witness family} is a map
$p_C:B_C\to L$, whose fibre $B_C(\ell)=p_C^{-1}(\ell)$ contains the
exhibited $C$-blocks at $\ell$. The \textbf{blocking predicate} and
its corresponding \textbf{core} are
\[
 \mathsf{Block}_C=\exists_{p_C}\top\hookrightarrow L,
 \qquad
 \mathsf{Core}_{B}=
 \bigwedge_{\varnothing\ne C\subseteq N}
 \neg\mathsf{Block}_C\hookrightarrow L.
\]
For strict strong-Nash blocks, $B_C$ is the comprehension of
$\mathsf{Improve}_C^L\hookrightarrow D_C^L$, with $p_C$ induced by
the first deviation projection. For effectivity blocks it is the
family of enforcing coalition coordinates; for endowment blocks it is
the family of certified reallocations of the coalition's resources.
These are different cores because their witness and alternative
carriers are different.
\end{definition}
The distinctions used here are:
\begin{center}
\begin{tabular}{@{}p{.24\textwidth}p{.69\textwidth}@{}}
\toprule
Stability notion & Alternative ruled out\\
\midrule
Nash equilibrium & One agent strictly improves while outsiders retain
their policies.\\
Strong Nash stability & A nonempty coalition changes its policies,
with outsiders fixed, and every member strictly improves.\\
Strong-Pareto stability & A coalition changes its policies,
with outsiders fixed; all members weakly improve and at least one
strictly improves.\\
Effectivity core & A coalition chooses a coordinate whose every legal
completion strictly improves all its members.\\
Endowment core & A coalition reallocates its own initial resources so
that every member strictly improves, using a certified reallocation
of its endowment.\\
\bottomrule
\end{tabular}
\end{center}

In the endowment interpretation, the outsider comparison
concerns exclusion of outsiders' resources, leaving their
policy coordinates. These are consequently different stability predicates.

\begin{definition}[Strict coalitional improvement and strong Nash stability]
\label{def:egt-strict-coalitional-improvement-and-strong-nash-stability}
For a nonempty coalition $C$, put
\begin{equation}\label{eq:egt-coalition-improve}
 \mathsf{Improve}_C=\bigwedge_{i\in C}
                  (\delta_C^G)^*\mathsf{Better}^{A,B}_i
 \quad\hookrightarrow D_C.
\end{equation}
Define strong Nash stability by
\begin{equation}\label{eq:egt-strong-nash-doctrine}
 \mathsf{SNE}_G=\bigwedge_{\varnothing\ne C\subseteq N}
       \forall_{p_1:D_C\to\Sigma}\neg\mathsf{Improve}_C.
\end{equation}
For the legal domain $m:L\hookrightarrow\Sigma$, put
\[
 \mathsf{Improve}_C^L=\bigwedge_{i\in C}
                    (\delta_C^{G,L})^*\mathsf{Better}^{A,B}_i
 \quad\hookrightarrow D_C^L
\]
and
\begin{equation}\label{eq:egt-legal-strong-nash}
 \mathsf{SNE}_{G,L}=\bigwedge_{\varnothing\ne C\subseteq N}
       \forall_{p_1:D_C^L\to L}\neg\mathsf{Improve}_C^L.
\end{equation}
The Pareto-improvement variant is made into an explicit optimality
predicate in Construction~\ref{con:egt-pareto-optimality} below.
\end{definition}

\begin{construction}[Weak and strong Pareto optimality]
\label{con:egt-pareto-optimality}
Use the paper's real-valued presentation
$v_i(\ell)=k_i(g(m\ell))$ on the legal profile space $L$.
For every nonempty coalition $C$, the existing strict coalition-improvement
predicate is $\mathsf{Improve}_C^L\hookrightarrow D_C^L$. Construct the
Pareto-improvement predicate on that same dependent alternative space by
\[
 \mathsf{ParetoGain}_C^L(\ell,\ell')\ \Longleftrightarrow\
 \left(\bigwedge_{i\in C}v_i(\ell')\geq v_i(\ell)\right)
 \wedge
 \left(\bigvee_{i\in C}v_i(\ell')>v_i(\ell)\right).
\]
Define the weak and strong Pareto-optimality assertions on $L$ by
\[
 \mathsf{WPO}_{C,L}
  =\forall_{p_1:D_C^L\to L}\neg\mathsf{Improve}_C^L,
 \qquad
 \mathsf{SPO}_{C,L}
  =\forall_{p_1:D_C^L\to L}\neg\mathsf{ParetoGain}_C^L.
\]
For the empty coalition both predicates are $\top_L$. Weak Pareto
optimality rules out an alternative strictly improving every member;
strong Pareto optimality rules out an alternative weakly improving every
member and strictly improving at least one. For $C=N$ these are the
ordinary global Pareto predicates relative to $L$; for other coalitions
outsider coordinates remain fixed. Their fibres are subterminal ordinary
dependent objects. For preorder-valued preferences, the construction evaluates the supplied weak preorder and its strict part \citep{ChenShiWang2022}.
\end{construction}

\begin{proposition}[Pareto optimality and strategic stability]
\label{prop:egt-pareto-optimality}
In the real-valued instance, for every coalition $C$,
$\mathsf{SPO}_{C,L}\leq\mathsf{WPO}_{C,L}$. Moreover
\[
 \mathsf{NE}_{G,L}
 =\bigwedge_{i\in N}\mathsf{WPO}_{\{i\},L}
 =\bigwedge_{i\in N}\mathsf{SPO}_{\{i\},L},
 \qquad
 \mathsf{SNE}_{G,L}
 =\bigwedge_{\varnothing\ne C\subseteq N}\mathsf{WPO}_{C,L},
\]
and the all-coalition strong-Pareto condition satisfies
\[
 \bigwedge_{\varnothing\ne C\subseteq N}\mathsf{SPO}_{C,L}
 \ \leq\ \mathsf{SNE}_{G,L}.
\]
All these assertions commute with substitution of a reference context;
legal strategic quotients preserve and reflect them when the utility
comparisons descend.
\end{proposition}
\begin{proof}
\leavevmode
\begin{enumerate}[label=\textup{(\arabic*)},leftmargin=*,itemsep=4pt,topsep=3pt]
  \item \textbf{Strategic comparison.} A strict gain for every member is a weak gain for every member with at
least one strict gain; taking complements and universally quantifying
gives $\mathsf{SPO}_{C,L}\leq\mathsf{WPO}_{C,L}$. For a singleton
coalition the two improvement predicates coincide and are exactly the
profitable unilateral-deviation predicate. Their conjunction is legal
Nash equilibrium.
  \item \textbf{Base-change comparison.} Conjoining the weak Pareto predicates over all
nonempty coalitions gives the definition of strong Nash stability. The final inequality follows pointwise from the first. Substitution
retains the entire alternative fibre by Beck--Chevalley, and the
quotient claim follows from alternative lifting in
Theorem~\ref{thm:egt-legal-reduction} and preservation of the descended
comparisons.
\end{enumerate}
\end{proof}

\begin{proposition}[Coalitional interpretation and quotient invariance]
For real-valued utilities, $\mathsf{SNE}_G(\sigma)$ holds exactly when,
for every nonempty $C$, there is no $\tau_C$ such that
$v_i(\tau_C,\sigma_{-C})>v_i(\sigma)$ for all $i\in C$.
On $L$ the same assertion ranges over legal alternatives with the actual
outsider coordinates fixed. Legality-sensitive strategic reduction preserves
and reflects these predicates, and also the Pareto variant when its utility
comparisons descend.
\end{proposition}
\begin{proof}
\leavevmode
\begin{enumerate}[label=\textup{(\arabic*)},leftmargin=*,itemsep=4pt,topsep=3pt]
  \item \textbf{Lifting through the comparison.} The fiber of $p_1$ over the reference consists precisely of the stated
alternatives. Expand the right quantifier and the pulled-back comparisons. For quotient invariance, Theorem~\ref{thm:egt-legal-reduction} supplies
both forward preservation and lifting of every reduced legal alternative.
  \item \textbf{Strategic comparison.} The commuting counterfactual realizations preserve each comparison. A
profitable alternative therefore exists on one side exactly when it exists
on the other. Finite conjunction over the coalitions gives the result.
\end{enumerate}
\end{proof}

\begin{remark}[Carriers of collective alternatives]
Strong Nash stability fixes actual outsiders. Effectivity quantifies over
all legal completions of a coalition choice. Exchange economies use another carrier: certificates for reallocating
a coalition's own initial endowment. Its core predicate uses the same assertion doctrine
with that distinct form of collective control.
\end{remark}

\begin{definition}[Enforceable blocks and the effectivity core]
\label{def:egt-enforceable-blocks-and-the-effectivity-core}
For real-valued utilities, $\ell\in L$, and nonempty $C$, define
\[
 P_{\ell,C}(\ell')\quad\Longleftrightarrow\quad
                  \bigwedge_{i\in C}v_i(\ell')>v_i(\ell).
\]
An enforceable block at $\ell$ is a realizable coalition coordinate
$c\in L_C$ such that every legal completion $\ell'$ of $c$ satisfies
$P_{\ell,C}$. Thus existence of such a block is exactly
$[C]_L P_{\ell,C}$, using the identity outcome map of $L$.
Define the effectivity core by
\[
 \mathsf{Core}^{\mathrm{eff}}_{G,L}(\ell)
   =\bigwedge_{\varnothing\ne C\subseteq N}\neg[C]_L P_{\ell,C}.
\]
This is a notion of stability defined by effectivity.
Characteristic-function games and cores with transferable utility require
their own specified data.
\end{definition}

\begin{proposition}[Amalgamation turns enforceable blocks into deviations]
An enforcing coordinate $c\in L_C$ has a candidate with the actual outsiders
of $\ell$ exactly when
\[
 (c,e_{-C}(\ell))\in\operatorname{im}
 \bigl(\lambda_{C,-C}:L\to L_C\times L_{-C}\bigr).
\]
Whenever this condition holds, the candidate is a strict strong-Nash block.
In particular surjectivity of $\lambda_{C,-C}$ makes every enforceable block
a strong-Nash block. The condition holds on a full product.
\end{proposition}
\begin{proof}
\leavevmode
\begin{enumerate}[label=\textup{(\arabic*)},leftmargin=*,itemsep=4pt,topsep=3pt]
  \item \textbf{Image characterization of deviations.} Membership in the displayed image says exactly that there is an
$\ell'\in L$ with $e_C(\ell')=c$ and
$e_{-C}(\ell')=e_{-C}(\ell)$. The latter equality places
$(\ell,\ell')$ in $D_C^L$; the former invokes the enforcing property, so
every member strictly benefits. Surjectivity provides that witness for
every pair. On a full product the witness is the tuple
$(c,\ell_{-C})$.
\end{enumerate}
\end{proof}

\begin{example}[Authority and information in complementary commitments]
For Example~\ref{ex:egt-commitments}, the coalition of both owners can move
from $(0,0)$ to $(1,1)$ and strictly benefit both members. Thus $(0,0)$
fails strong Nash stability while $(1,1)$ satisfies it. Regrouping the owners
turns this joint deviation into a unilateral deviation of the new owner;
pooling their observations alone does not change its ownership. A collective
preference for a regrouped owner remains separate from that structural change.
\end{example}

\subsubsection{Substitution and three different changes of context}

\begin{proposition}[Stability of agent formation]
\label{agt-prop-agent-basechange}
For $f:V\to W$ there are coherent isomorphisms
\[
 f^*\Ag_{\sig}(A,B)
 \cong\Ag_{f^*\sig}(f^*A,f^*B).
\]
They preserve sequential and parallel composition, execution, and
ownerwise reparametrization.
\end{proposition}
\begin{proof}
\leavevmode
\begin{enumerate}[label=\textup{(\arabic*)},leftmargin=*,itemsep=4pt,topsep=3pt]
  \item \textbf{Fibrewise calculation.} At $v$, both sides are
$[\Lin\Sigma_{f(v)}\otimes A_{f(v)},B_{f(v)}]$. All the named operations
are defined by these pointwise tensor, evaluation, and ordinary profile
maps. Identity and composition of substitutions give the same evaluations.
\end{enumerate}
\end{proof}

Three substitutions that can look similar on state carriers have
different agency meanings. First, changing the reference state pulls back
only the source of an alternative span. For
$W\xleftarrow{s}E\xrightarrow{t}Y$ and $f:V\to W$, the new
source is the corresponding pullback.
\[
 V\xleftarrow{}V\times_WE\xrightarrow{t\pi_E}Y.
\]
The target remains all of $Y$, so the original alternatives stay
available. Beck--Chevalley then gives exact substitution of existential
and universal tests along the reference map.

Second, a policy reparametrization changes the available controlled
alternatives themselves. Its comparison is the realization-preserving
map and alternative-lifting construction of
Section~\ref{agt-sec-ownership}, not a source-only pullback.

Third, epistemic restriction changes both endpoints of
accessibility. For $j:U\hookrightarrow W$, the relation $R$ becomes
$(j\times j)^*R$, deleting excluded alternatives as well as excluded
references. Public announcements in Section~\ref{agt-sec-public}
use this operation, which is not interchangeable with the first two.

\subsection{Uniformly informed and legal policies}
\label{agt-sec-uniform}

\subsubsection{The legal local menu in agent syntax}

Let $o:W\twoheadrightarrow O$ be an effective observation and
$F\hookrightarrow W\times M$ a feasibility or permission relation.
The construction in \cref{con:egt-informed-legal-interface} gives
$A_o=\forall_{o\times1_M}F$. Its agent-calculus dependent type is
\begin{equation}
 \Pol_o(F)=\prod_{z:O}\{m:M\mid A_o(z,m)\},
 \label{agt-eq-uniform-policies}
\end{equation}
canonically identified with the sections of $A_o\to O$.

\begin{theorem}[Uniform legality]
\label{agt-thm-uniform}
For $\pi:O\to M$, membership in $\Pol_o(F)$ is equivalent to
$F(w,\pi(o(w)))$ for every $w$. Thus it represents exactly the
observation-based responses legal at all compatible worlds.
\end{theorem}
\begin{proof}
\leavevmode
\begin{enumerate}[label=\textup{(\arabic*)},leftmargin=*,itemsep=4pt,topsep=3pt]
  \item \textbf{Adjunction and transposition.} Under the sections isomorphism, membership means
$\Gamma_\pi\subseteq A_o$. The universal-image adjunction proves
$\Gamma_\pi\subseteq\forall_{o\times1}F$ if and only if
$(o\times1)^*\Gamma_\pi\subseteq F$; this is exactly the worldwise
criterion of \cref{prop:egt-informed-legal-universal}.
\end{enumerate}
\end{proof}

For complete-policy changes of owner $i$, take $M=\Sigma_i$ and
$F_i(w,u)\iff(s(w)_{-i},u)\in L$. The empty-menu counterexample and
its quantifier-order interpretation are given after
\cref{prop:egt-informed-legal-universal}.

\subsubsection{Refinement and decentralized coalitions}

\begin{proposition}[Refinement of information]
\label{agt-prop-info-refinement}
If $o=r\circ o'$ with both observations effective, precomposition
$\pi\mapsto\pi r$ injects $\Pol_o(F)$ into $\Pol_{o'}(F)$.
Its image is precisely the refined policies constant on the fibres of $r$.
\end{proposition}
\begin{proof}
\leavevmode
\begin{enumerate}[label=\textup{(\arabic*)},leftmargin=*,itemsep=4pt,topsep=3pt]
  \item \textbf{Lifting along the surjection.} This is the section-level specialization of the refinement and
effective-descent comparisons in \cref{prop:egt-informed-legal-universal}:
surjectivity of $r$ makes the map injective, and unchanged worldwise
responses preserve legality in both directions.
\end{enumerate}
\end{proof}

For joint legality $F_C\hookrightarrow W\times\prod_{i\in C}M_i$,
the decentralized type is
\[
 \Pol_C^{\mathrm{dec}}(F_C)=
 \left\{(\pi_i)\in\prod_{i\in C}M_i^{O_i}:
 \forall w\ F_C(w,(\pi_i(o_i(w)))_i)\right\}.
\]
The pooled observation image $O_C$ defines the map
$(\pi_i)_i\mapsto(z\mapsto(\pi_i(p_i z))_i)$.

\begin{proposition}[Exact decentralized comparison]
\label{agt-prop-decentralized}
This map into pooled legal policies is injective and has as image
exactly those pooled policies whose $i$th output descends along
$p_i:O_C\to O_i$, for each member of $C$.
\end{proposition}
\begin{proof}
\leavevmode
\begin{enumerate}[label=\textup{(\arabic*)},leftmargin=*,itemsep=4pt,topsep=3pt]
  \item \textbf{Policy-level comparison.} \Cref{prop:egt-pooled-policy-comparison,prop:egt-decentralized-legal}
give the injectivity, coordinatewise descent, and unchanged joint
feasibility test. Their explicit map is the one displayed above, so
this is an identification of dependent policy types, not a new
independence or marginal-legality principle.
\end{enumerate}
\end{proof}

Pooling can strictly enlarge the policy type, but executable access
to the pooled signal still requires its specified communication wiring.

\subsubsection{Collective evaluation and institutional power}

For a nonempty coalition $C$, one constructed collective comparison is
\[
 P_C^{\mathrm{all}}(\sigma,\tau)
 =\bigwedge_{i\in C}P_i(\sigma,\tau).
\]
The corresponding stability predicate says that no legal $C$-deviation
strictly improves every member. Weak Pareto improvement is a different
predicate, formed from the members' weak comparisons and at least one
strict comparison. A designated corporate comparison can also be supplied
as the institution's evaluation interface. Regrouping controls does not
choose between these evaluations.

\begin{corollary}[Delegation preserves the implemented stability test]
For the image comparison of Theorem~\ref{agt-thm-delegation-image}, pull back
the same parent comparison predicates. Stability against a family of parent
coalitions pulls back exactly to stability against the corresponding whole
delegate blocks. In particular parent Nash equilibrium pulls back to Nash
equilibrium after regrouping each delegate block as one owner.
\end{corollary}
\begin{proof}
\leavevmode
\begin{enumerate}[label=\textup{(\arabic*)},leftmargin=*,itemsep=4pt,topsep=3pt]
  \item \textbf{Lifting through the comparison.} Every subordinate deviation maps to a parent deviation with the same
comparison value. Every parent deviation lifts at the given reference
while retaining outsider blocks. Existence of an improving alternative is
therefore equivalent on the two sides. Negating and conjoining the specified
tests gives the result.
\end{enumerate}
\end{proof}

Institutional ability has a different quantifier pattern. Let
$e_C:L\twoheadrightarrow L_C=\im(L\to\Sigma_C)$ be the realizable
coalition-coordinate map, and let $z:L\to Z$ record outcomes. Define
\begin{equation}
 \mathsf{Can}_C(P)
 \iff\exists u\in L_C\ \forall\ell\in L\,
          (e_C(\ell)=u\Rightarrow P(z\ell)).
 \label{agt-eq-effectivity}
\end{equation}
Each quantified coordinate has an actual legal completion because $e_C$ is
surjective. Thus the formula gives no vacuous power to impossible choices.
When choices must be informed, the existential witness is a member of the
decentralized or pooled legal policy type, and the universal quantification
is taken over its compatible worlds and legal completions. This constructs
uniform ability directly, with the information constraint in its witness.

The resulting distinctions are now explicit: having an agent
interface, having a feasible execution, possessing an informed legal
policy, satisfying rational selection, and having coalition power are
different predicates or judgments on the displayed process. The maps and
quantifiers above specify their relations; none is inferred from the
others merely by terminology.

\section{Witnessed Institutions and Dynamic Agency}\label{sec:paper-8}

Legality so far restricts the policies available in one fixed
interaction. An institution may also change public records, create
entitlements, or certify that an action occurred. For such changes,
a relation between initial and final records is insufficient: different
executions can have the same endpoints yet carry different concurrent
processes and certificates. We therefore retain witnessed execution spans
and their realizers, show how they compose, and use their dependent sums
and products to interpret dynamic agency.

\subsection{Dependent witnesses and certified realizations}
\label{sec:egt-witnesses}

The crucial distinction is between a feasible outcome and an
exhibited execution establishing it. Existential images classify feasible
outcomes, while a dependent witness family retains the actual
certificates and concurrent process data. These two readings will
support the proof-irrelevant and proof-relevant institutional operations,
respectively.

\begin{construction}[Dependent witness families]
\label{con:egt-general-witness-family}
Let $A,B\in\mathsf L(\Gamma)$ and let $K\in\mathsf T(\Gamma)$.
Its context extension is
$p_K:\Gamma.K=\coprod_{\gamma\in\Gamma}K_\gamma\to\Gamma$.
A witness realization is an ordinary morphism
\[
 \rho:K\longrightarrow\mathsf{Val}_\Gamma[A,B]_D.
\]
Thus each witness $k\in K_\gamma$ exhibits a linear arrow
$\rho_\gamma(k):A_\gamma\to B_\gamma$.
A certified family is a section $\kappa:1_\Gamma\to K$; its realized
family is $g=\rho\kappa$. The predicate that a witness exists is
\[
 \|K\|=\exists_{p_K}\top\hookrightarrow\Gamma.
\]
The witness object $K$, its realization $\rho$, and any exhibited section
remain available separately from this predicate. For $f:\Delta\to\Gamma$,
pullback gives $f^*K$, $f^*\rho$, and $f^*\kappa$, and
\[
 f^*(\rho\kappa)=(f^*\rho)(f^*\kappa),\qquad
 f^*\|K\|=\|f^*K\|.
\]
The first equation is functoriality of substitution; the second is
Proposition~\ref{prop:egt-propositional-recovery}. Consequently realization
and existential feasibility share the constructed substitution while
retaining their different carriers. Allocation and production witnesses are particular instances of
this construction.
\end{construction}

For a constant interaction boundary and a displayed certificate carrier,
the compatibility of its tests with operational tests has the following
form. Here $K$ denotes the total certificate carrier and $U$ its
realization map.

\begin{proposition}[Certificate tests and operational tests]
\label{prop:fib-certificate-tests}
Let $K$ be a displayed certificate carrier,
$U:K\to\mathcal C(A,B)$ its concurrent realization, and
$\kappa:\Sigma\to K$ a certified arena with $g=U\kappa$.
For $Q\hookrightarrow\mathcal C(A,B)$ and $P\hookrightarrow K$,
\[
 g^*Q=\kappa^*U^*Q,
 \qquad \operatorname{Test}_P(\kappa)=\kappa^*P.
\]
Thus operational and certificate predicates belong to the same assertion
fibre, and the second class may retain distinctions unseen by $U$.
\end{proposition}
\begin{proof}
\leavevmode
\begin{enumerate}[label=\textup{(\arabic*)},leftmargin=*,itemsep=4pt,topsep=3pt]
  \item \textbf{Witness and legality comparison.} The first equality is functoriality of predicate substitution; the second
is semantic testing of the certificate term. Both have typed expressions
using $\Delta_GK$ and $\Delta_G\mathcal C(A,B)$. No inverse to $U$
is needed: the carrier $K$ and the term $\kappa$ are retained.
\end{enumerate}
\end{proof}

\subsection{Witnessed institutional executions}
\label{sec:egt-institutional-executions}

A legal domain specifies permitted profiles of one interaction;
institutional composition additionally requires matched changes of public
records. We therefore retain a witness linking each initial record to a
final record, together with its concurrent realization. This is the
record-sensitive interpretation of institutional power considered in
\citet{JonesSergot1996InstitutionalPower}. Distinct witnesses with the
same endpoints remain distinct, since subsequent actions may depend on
their certificates or traces.

\begin{definition}[A witnessed institutional execution]
\label{def:egt-witnessed-institution}
An institutional boundary is a pair $(X,A)$ of an ordinary record set
$X$ and a linear boundary $A\in\mathcal C$. A witnessed execution from
$(X,A)$ to $(Y,B)$ is
\[
 \mathcal E=(X\xleftarrow{s}E\xrightarrow{t}Y,\,
             \rho:E\to\mathcal C(A,B),\,P\hookrightarrow E).
\]
Here $E$ retains execution witnesses, $\rho(e)$ is the realized
concurrent interaction, and $P$ consists of permitted witnesses.
A map $\mathcal E\Rightarrow\mathcal E'$ is a function $f:E\to E'$
with $s'f=s$, $t'f=t$, $\rho'f=\rho$, and $f[P]\subseteq P'$.
The last condition preserves permission; reflection of permission
is a stronger property and is not included in every map.
\end{definition}

An owned arena $G=(\vec\Sigma,g)$ yields such a witness
presentation when execution data include a policy map
$\lambda:E\to\Sigma$, record maps, and a permission predicate
$P\subseteq\lambda^{-1}(L)$. Its realizer is $\rho=g\lambda$.
The witness may retain a policy tuple, a certificate, or hidden
communication. Restricting to permitted witnesses changes the
institutional action but does not identify distinct executions.

\begin{theorem}[Composition of witnessed institutions]
\label{thm:egt-institution-bicategory}
\label{agt-thm-institutions}
Witnessed executions and their maps form a symmetric monoidal
bicategory $\mathbf{Inst}(\mathcal C)$. For consecutive executions
$\mathcal E_1:(X,A)\to(Y,B)$ and
$\mathcal E_2:(Y,B)\to(Z,C)$, composition is
\[
 E_{21}=E_1\times_Y E_2,\qquad
 \rho_{21}(e_1,e_2)=\rho_2(e_2)\circ\rho_1(e_1),\qquad
 P_{21}=P_1\times_Y P_2.
\]
The source and target are $s_1(e_1)$ and $t_2(e_2)$.
The identity at $(X,A)$ has witness carrier $X$, both endpoint
maps equal to $\operatorname{id}_X$, constant realizer
$\operatorname{id}_A$, and every witness permitted.
The tensor has
\[
 (X,A)\otimes(X',A')=(X\times X',A\otimes A'),\qquad
 E_{\mathcal E\otimes\mathcal F}=E\times F,\qquad
 P_{\mathcal E\otimes\mathcal F}=P_{\mathcal E}\times P_{\mathcal F},
\]
with coordinatewise endpoints and realizer
$\rho_{\mathcal E}(e)\otimes\rho_{\mathcal F}(f)$.
\end{theorem}

\noindent Composing witnessed executions matches the output record of the first span to the input record of the second.
\begin{center}
\begin{tikzcd}[column sep=large,row sep=large]
E_1\times_Y E_2\arrow[r,"\mathrm{pr}_2"]\arrow[d,"\mathrm{pr}_1"']\arrow[dr,phantom,"\lrcorner",very near start]
 &E_2\arrow[d,"s_2"]\\
 E_1\arrow[r,"t_1"']&Y
\end{tikzcd}
\end{center}


\begin{proof}
\leavevmode
\begin{enumerate}[label=\textup{(\arabic*)},leftmargin=*,itemsep=4pt,topsep=3pt]
  \item \textbf{Secured operational interaction.} For composable witnessed executions, the pullback has the explicit
fibrewise description
\begin{equation}
\label{eq:institution-composite-witness-fibre}
 (E_1\times_YE_2)(x,z)
 \cong\coprod_{y\in Y}E_1(x,y)\times E_2(y,z),
\end{equation}
which sends a matching pair to its uniquely determined intermediate
record and its two constituent witnesses. Its inverse assembles that
pair, so it retains both realizer labels and both permissions.
  \item \textbf{Naturality of witnessed composition.} If $f:E_1\to E'_1$ and $g:E_2\to E'_2$ preserve the endpoints,
$(e_1,e_2)\mapsto(f(e_1),g(e_2))$ is a well-defined map of the
pullbacks. It preserves the composite realizer by the two realization
equations and sends permitted pairs to permitted pairs. This defines
horizontal composition of maps; vertical composition is composition
of functions.
  \item \textbf{Interchange of witnessed executions.} The interchange equation is pointwise the equality
of the two componentwise composite functions.
  \item \textbf{Secured operational interaction.} The associator sends $((e_1,e_2),e_3)$ to $(e_1,(e_2,e_3))$. Both sides are exactly the triples with matching consecutive records. Their realizers agree by associativity in $\mathcal C$, and their
permission predicates are both the conjunction of the three
permission predicates.
  \item \textbf{Witness and legality comparison.} The unitors delete the unique intermediate
identity witness, whose record is already fixed by the adjacent
endpoint. Pentagon and triangle coherence follow because each
route has the same tuple projections.
  \item \textbf{Parallel compatibility.} For tensor and horizontal composition, the comparison
\[
 (E_1\times E'_1)\times_{Y\times Y'}(E_2\times E'_2)
 \cong (E_1\times_YE_2)\times(E'_1\times_{Y'}E'_2)
\]
reorders a matching quadruple. The interchange law in $\mathcal C$
preserves its realizer, and the four permission tests agree on both
sides.
  \item \textbf{Secured operational interaction.} The monoidal associativity, unit, and symmetry components
are the graph spans of the corresponding bijections of ordinary
records, labelled by the corresponding structural isomorphisms
of $\mathcal C$, with all witnesses permitted. Their naturality
and coherence reduce to those bijections and the monoidal coherence
of $\mathcal C$. This supplies the stated symmetric monoidal
bicategory structure.
\end{enumerate}
\end{proof}


Products here combine ordinary institutional records, not copies
of arbitrary linear resources. The realized concurrent interactions
continue to compose by secured interaction and hiding. Two distinct
witnesses may therefore have identical source and target records while
remaining distinguishable by their strategies or certifications.

\begin{proposition}[Witness retention, permission, and institutional modalities]
\label{prop:egt-institution-modalities}
Write
$E(x,y)=\{e\in E:s(e)=x,\ t(e)=y\}$.
Composition retains the dependent witness family
\[
 E_{21}(x,z)\cong\coprod_{y\in Y}E_1(x,y)\times E_2(y,z).
\]
Taking inhabitedness gives ordinary relational composition.
For $Q\hookrightarrow Y$, the source predicates
\[
 \langle\mathcal E\rangle Q=\exists_s\,t^*Q,\qquad
 [\mathcal E]Q=\forall_s\,t^*Q
\]
express, respectively, that some witnessed execution reaches $Q$
and that every witnessed execution reaches $Q$. They obey
\[
 \langle\mathcal E_2\circ\mathcal E_1\rangle Q
   =\langle\mathcal E_1\rangle\langle\mathcal E_2\rangle Q,\qquad
 [\mathcal E_2\circ\mathcal E_1]Q
   =[\mathcal E_1][\mathcal E_2]Q.
\]
The same equations hold when each execution is restricted to its
permitted subobject. A guarantee requiring an actual execution is
\[
 (\exists_s\top)\wedge\forall_s t^*Q.
\]
Finally the violating witnesses of a composite are exactly
\[
 \begin{aligned}
 (E_1\times_YE_2)\setminus(P_1\times_YP_2)
 &=
 ((E_1\setminus P_1)\times_YE_2)\\
 &\quad{}\cup (E_1\times_Y(E_2\setminus P_2)).
 \end{aligned}
\]
\end{proposition}
\begin{proof}
\leavevmode
\begin{enumerate}[label=\textup{(\arabic*)},leftmargin=*,itemsep=4pt,topsep=3pt]
  \item \textbf{Witness and legality comparison.} A composite witness is a pair with one common intermediate record;
recording that value gives the displayed disjoint-union bijection. It has an inhabitant exactly when some intermediate record has
an inhabitant in each of the two fibres.
  \item \textbf{Fibrewise calculation.} At $x$, the composite diamond says that there are $e_1,e_2$
with $s_1(e_1)=x$, $s_2(e_2)=t_1(e_1)$, and $Q(t_2(e_2))$. This is the nested existential formula on the right. The box
replaces these existential tests by the two universal tests,
giving its equation, including intermediate records with no
outgoing witness.
  \item \textbf{Predicate calculation.} Categorically these are the dependent
quantifier composition and pullback equations. The restricted
permitted spans compose by the same pullback, so the argument
applies to them as well. The extra conjunct $\exists_s\top$
excludes a source with no execution, where the box alone is
vacuously true.
  \item \textbf{De Morgan law for matched execution witnesses.} The final equality is the pointwise equivalence
$\neg(P_1(e_1)\wedge P_2(e_2))$ with
$\neg P_1(e_1)\vee\neg P_2(e_2)$ on the pullback.
\end{enumerate}
\end{proof}

\noindent An institutional execution witness carries a source record and a target record.
\begin{center}
\begin{tikzcd}[column sep=large,row sep=large]
&E\arrow[dl,"s"']\arrow[dr,"t"]&\\
 X && Y
\end{tikzcd}
\end{center}

An authorization rule is represented by the enabling
comprehension and target record of an execution; the separate predicate
$P$ specifies whether that execution is permitted. Retaining the
witness allows the model to record prohibited occurrences for later
sanction or audit without treating permission as a property of the
underlying concurrent arrow. An institutional state change and its
normative status consequently remain distinct.

\subsection{Linear realization and dynamic modalities}
\label{agt-sec-institutions}

The witnessed executions of \cref{def:egt-witnessed-institution} compose
in the symmetric monoidal bicategory of
\cref{thm:egt-institution-bicategory}. We now construct their linear
execution maps and the dependent type operations induced by their witness
spans. These constructions retain both the realizer and the permission
predicate.

\subsubsection{The linear execution induced by witnesses}
\label{agt-sec-institution-run}

\begin{proposition}[Realized witnessed execution]
\label{agt-prop-witness-run}
For a witnessed execution
$\mathcal E=(X\xleftarrow{s}E\xrightarrow{t}Y,\rho,P)$,
let $\underline{yA}$ and $\underline{yB}$ denote the constant linear
families on $X$ and $Y$, and put $J_E=\Sigma_s^\ell\I_E$.
The concurrent realizer constructs a canonical linear map over $X$:
\begin{equation}
 \operatorname{run}_\rho:
 \underline{yA}\otimes J_E\longrightarrow
 \Sigma_s^\ell t^*\underline{yB}.
 \label{agt-eq-witness-run}
\end{equation}
These maps compose along the canonical pullback and coproduct-flattening
comparisons, with the realizer of a matching pair $(e,f)$ equal to
$y\rho_F(f)\,y\rho_E(e)$.  Restricting to $P$ gives the corresponding
permitted execution map.
\end{proposition}

\begin{proof}
\leavevmode
\begin{enumerate}[label=\textup{(\arabic*)},leftmargin=*,itemsep=4pt,topsep=3pt]
  \item \textbf{Fibrewise calculation.} Pointwise Yoneda gives the family morphism
$\rho^\sharp:s^*\underline{yA}\to t^*\underline{yB}$ with component
$y\rho(e)$ at $e$. By Frobenius and the tensor unitor, define
\[
 \underline{yA}\otimes\Sigma_s^\ell\I_E
 \xrightarrow{\cong}
 \Sigma_s^\ell(s^*\underline{yA}\otimes\I_E)
 \xrightarrow{\Sigma_s^\ell(\rho^\sharp\otimes1)}
 \Sigma_s^\ell(t^*\underline{yB}\otimes\I_E)
 \xrightarrow{\Sigma_s^\ell(\text{right unitor})}
 \Sigma_s^\ell t^*\underline{yB}.
\]
At $x\in X$, this is the copairing of the maps
$y\rho(e):yA\to yB$ over $e$ with $s(e)=x$, retaining each summand's
target-record index.
  \item \textbf{Base-change comparison.} Under a second execution, Beck--Chevalley and
successive dependent sums flatten the summands indexed by matching
$(e,f)$; the component map is
$y\rho_F(f)\circ y\rho_E(e)=y(\rho_F(f)\rho_E(e))$. The composite execution has exactly that realizer. Coproduct uniqueness
proves the comparison equation, and flattening triples proves associativity
coherence. Replacing $E$ by its permitted comprehension gives the last
assertion.
\end{enumerate}
\end{proof}


\subsubsection{Dynamic type and predicate operations}

For the underlying span $X\xleftarrow{s}E\xrightarrow{t}Y$, define
\[
 \langle E\rangle^\ell=\Sigma_s^\ell t^*,\qquad
 [E]^\ell=\Pi_s^\ell t^*.
\]
At $x$, the first retains a choice of execution and an output resource;
the second supplies an additive family indexed by every execution from
$x$. The predicate operations are
\[
 \langle E\rangle P=\exists_s t^*P,\qquad
 [E]P=\forall_s t^*P.
\]
They quantify over witnessed executions, including distinct executions with
the same source and target.

\begin{theorem}[Compositional dynamic modalities]
\label{agt-thm-dynamics}
For consecutive witnessed spans $E,F$,
\[
 [F\circ E]^\ell\cong[E]^\ell[F]^\ell,
 \qquad
 \langle F\circ E\rangle^\ell
 \cong\langle E\rangle^\ell\langle F\rangle^\ell.
\]
The predicate versions are equalities. Identity spans give identity
operations. The comparisons are associative and unital.
\end{theorem}
\begin{proof}
\leavevmode
\begin{enumerate}[label=\textup{(\arabic*)},leftmargin=*,itemsep=4pt,topsep=3pt]
  \item \textbf{Fibrewise calculation.} At $x$, the box on the composite is
\[
 \prod_{\substack{e:s_E(e)=x\\ f:s_F(f)=t_E(e)}}A_{t_F(f)}.
\]
Flattening a nested product gives exactly
$\prod_{e:s_E(e)=x}\prod_{f:s_F(f)=t_E(e)}A_{t_F(f)}$. The sum proof flattens the corresponding nested coproduct. Both include
empty fibres.
  \item \textbf{Representation and Yoneda decoding.} Three successive comparisons flatten the same triples of
witnesses, so product projections and coproduct injections prove coherence. The identity fibre has its unique record witness. Predicate quantifiers
obey the identical nested quantification formulas.
\end{enumerate}
\end{proof}

For a relation $R\subseteq X\times Y$, the dynamic diamond
$\exists_s t^*$ is left adjoint to $\forall_t s^*$, the box
associated with the converse direction. The forward diamond and
forward box are not generally an adjoint pair. Recording this variance
prevents an incorrect transport rule when sequential execution spans
are composed.

Passing from a witness span $E$ to its relational image retains
existential and universal truth tests on target predicates. It discards
proof-relevant multiplicity: distinct executions at one pair of
endpoints still contribute separate dependent summands or factors and
may retain different strategies or certificates. The institutional
type interpretation therefore uses the full span.

\subsubsection{Permission and collective institutional agency}

Restricting $E$ to $P$ gives permitted modalities. A guarantee of $Q$ with
an actual permitted execution is
\[
 (\exists_{s|P}\top)\wedge\forall_{s|P}(t|P)^*Q.
\]
The first conjunct excludes the vacuous truth of a universal test at a
record with no permitted execution. On a composite, violation is precisely
\[
 (\neg P_1\wedge\top)\ \vee\ (\top\wedge\neg P_2)
\]
as a predicate on $E_1\times_YE_2$. This records the stage at which
permission failed.

An institution can be exposed as a collective agent by retaining its
owner-indexed presentation and applying Section~\ref{agt-sec-ownership}.
Its observation and evaluation are then situated collective interfaces.
This does not make the institution's normative rules into a preference:
who may act, which actions are authorized, and which outcomes the
collective prefers remain separately displayed. The preceding bicategory
makes their witnessed executions compositional.

\section{Public Announcements and Epistemic Actions}\label{sec:paper-9}

Institutional executions change records; epistemic actions change
which states remain possible and what event an agent observes. Public
announcement is the case in which an asserted predicate restricts the
state carrier and both endpoints of each accessibility relation. A
general action model additionally retains an event coordinate and its
observation structure. We construct both updates through comprehension,
then identify the lifting conditions needed to preserve modal formulas
when abstract updates are compared with implemented communication.

\subsection{Communication and epistemic updates}
\label{sec:egt-communication}

An agent may announce information, but a stated formula is not yet
a delivered message. At the logical level, public update restricts the
situated state space and reindexes the epistemic relations. At the
operational level, a communicating agent must have a complete policy,
an enabled message, and a causal realization of its delivery. We first
establish the logical update and then compare it with the complete
communication construction.

The later policy-to-message comparison keeps each owner's
pre-existing observation and policy record alongside the public message
tuple. Its forward map and source-fixed lifts will transport knowledge,
distributed knowledge, coalition deviations, and the specified fixed-point
assertions. This is stronger than agreement on current messages alone:
modal formulas quantify over alternatives that must also be preserved.

\subsubsection{Public restriction by comprehension}
\label{sec:egt-dynamic-update}

Public announcement is the comprehensional update of public
communication logic, originating with Plaza and represented here by
\citet{Plaza2007PublicCommunications}. Its operation on the state set
is simple, but its epistemic meaning depends on restricting both
arguments of accessibility. We construct this restriction using the
predicate doctrine rather than adding an independent update primitive.

\begin{definition}[Public announcement]
\label{def:egt-public-announcement}
Fix a relation frame $M$ with state set $W$ and relations
$(R_i)_{i\in N}$. For a predicate $P\hookrightarrow W$, let
$a:U=\{P\}\hookrightarrow W$ be its comprehension. The \textbf{public announcement} of $P$ updates the
entire frame to
\[
 M|_P=\bigl(U,(R_i^P)_{i\in N}\bigr),\qquad
 R_i^P=(a\times a)^*R_i\hookrightarrow U^2.
\]
Its box and diamond are the dependent quantifiers
\[
 [!P]=\forall_a:\mathsf{Pred}(U)\to\mathsf{Pred}(W),\qquad
 \langle!P\rangle=\exists_a:\mathsf{Pred}(U)\to\mathsf{Pred}(W).
\]
For a fixed predicate $Q\hookrightarrow W$,
$[!P](a^*Q)=P\Rightarrow Q$. Writing $[!a]$ and
$\langle!a\rangle$ for the same operators when the inclusion $a$
is specified, epistemic formulas in the updated frame are evaluated
using the restricted relations $R_i^P$ on both endpoints.
Merely substituting reference states would instead retain the original
target alternatives and would not be this update.
\end{definition}
This changes the epistemic model, not merely the chosen reference
state. Selecting only references satisfying $P$ would leave excluded
states available as epistemic alternatives; restricting both endpoints
removes them. The corresponding reindexed observation fibres represent
the information available after the announcement.
\begin{construction}[Public update]
For a monomorphism $a:U\hookrightarrow W$, put $A=\exists_a\top$.
Restrict an operational family $t$ to $ta$ and an assertion to $a^*P$.
Define $[!a]=\forall_a$ and $\langle!a\rangle=\exists_a$.
For a relation $R$ on $W$, its restricted relation is
$R|_U=(a\times a)^*R$.
This restriction changes both endpoints, unlike a reference pullback.
\end{construction}

\begin{proposition}[Monic change of base and its nucleus]
\label{prop:egt-monic-nucleus}
Put $L=\exists_a$, $\rho=a^*$, and $J=\forall_a$.
Then $\rho L=\rho J=\operatorname{id}$ and $L\rho(P)=A\wedge P$.
Moreover
\begin{equation}\label{eq:egt-monic-right-adjoint}
 JQ=A\Rightarrow LQ,
\end{equation}
and
\begin{equation}\label{eq:egt-monic-nucleus}
 J\rho(P)=j_A(P)=A\Rightarrow P.
\end{equation}
Thus $L$ identifies $\mathsf{Pred}(U)$ with $\downarrow A$, and $J$
identifies it with the fixed predicates of the nucleus $j_A$.
\end{proposition}
\begin{proof}
\leavevmode
\begin{enumerate}[label=\textup{(\arabic*)},leftmargin=*,itemsep=4pt,topsep=3pt]
  \item \textbf{Base-change comparison.} The self-pullback of a monomorphism has identity projections, so
Beck--Chevalley gives $\rho L=\rho J=\operatorname{id}$. Frobenius gives
$L\rho(P)=A\wedge P$. Consequently
$P\leq A\Rightarrow LQ$ iff $L\rho P\leq LQ$ iff $\rho P\leq Q$.
  \item \textbf{Verification of the universal property.} Uniqueness of right adjoints proves the formula for $J$. Substitute
$Q=\rho P$. The unit of $\rho\dashv J$ gives inflationarity of $J\rho$;
$\rho J=\operatorname{id}$ gives idempotence, and both factors preserve
finite meets. The two claimed identifications follow from the same identities.
\end{enumerate}
\end{proof}

\noindent The inclusion of an announced submodel and its existential and universal predicate extensions form the update comparison.
\begin{center}
\begin{tikzcd}[column sep=large,row sep=large]
\mathsf{Pred}(U)\arrow[r,"\exists_m"]\arrow[d,"\forall_m"']
 &\mathsf{Pred}(W)\arrow[d,"m^*"]\\
 \mathsf{Pred}(W)\arrow[r,"m^*"']&\mathsf{Pred}(U)
\end{tikzcd}
\end{center}

\begin{corollary}[Update monad, support comonad, and restricted knowledge]
\label{cor:egt-update-adjoints}
For the inclusion $a:U\hookrightarrow W$, the endomaps
$L\rho=A\wedge(-)$ and $J\rho=A\Rightarrow(-)$ are respectively an
idempotent comonad and monad. Their coalgebra and algebra posets are
equivalent to $\mathsf{Pred}(U)$ via $L$ and $J$, respectively.
For observational knowledge $K=[\ker(o)]$ and $K^U=[\ker(oa)]$,
\[
 K^U=\rho KJ,\qquad JK^U=j_AKJ.
\]
\end{corollary}
\begin{proof}
\leavevmode
\begin{enumerate}[label=\textup{(\arabic*)},leftmargin=*,itemsep=4pt,topsep=3pt]
  \item \textbf{Adjunction and transposition.} The preceding monic adjunction identities identify both fixed-point
posets and give their units and counits. For $u\in U$, the assertion
$(\rho KJQ)(u)$ says that every $w$ with $o(w)=o(a(u))$ satisfies
$JQ$. Outside $U$ this is true; inside $U$ it is exactly $Q$.
  \item \textbf{Observation and knowledge.} This is
knowledge for the restricted observation $oa$. Applying $J$ and using
$J\rho=j_A$ gives the second equation.
\end{enumerate}
\end{proof}

\begin{example}[Restriction and knowledge need not commute]
Let $W=\{0,1\}$ have one observation class and let $A=P=\{0\}$.
Then $Kj_AP=W$, whereas $j_AKP=\{1\}$. The restriction formula is
$K^U=\rho KJ$; no commutation of $K$ with $j_A$ is used.
\end{example}

\begin{theorem}[Public-update reduction and composition]
\label{thm:egt-public-update-reduction}
\label{prop:egt-public-update-composition}
For predicates $P$ on $W$ and $Q,Q'$ on $U$,
\[
 [!a]a^*P=A\Rightarrow P,\qquad
 [!a](Q\wedge Q')=[!a]Q\wedge[!a]Q',
\]
\[
 [!a](Q\Rightarrow Q')
 =A\Rightarrow([!a]Q\Rightarrow[!a]Q'),\qquad
 [!a]\neg Q=A\Rightarrow\neg[!a]Q.
\]
For a relation $R$ on $W$,
\begin{equation}\label{eq:egt-public-update-modal-reduction}
 [!a][R|_U]Q=A\Rightarrow[R]([!a]Q).
\end{equation}
For $R=\ker o$, the restricted relation is $\ker(oa)$; hence
\begin{equation}\label{eq:egt-public-update-knowledge-reduction}
 [!a]K^UQ=A\Rightarrow K([!a]Q),
\end{equation}
and
\begin{equation}\label{eq:egt-announcement}
 [!a]K^U(a^*P)=A\Rightarrow K(A\Rightarrow P).
\end{equation}
For $V\xhookrightarrow bU\xhookrightarrow aW$,
$[!a][!b]=[!(ab)]$. If $B=\exists_b\top$ in $U$, then
\begin{equation}\label{eq:egt-composite-announcement}
 \exists_aB=A\wedge[!a]B.
\end{equation}
\end{theorem}
\begin{proof}
\leavevmode
\begin{enumerate}[label=\textup{(\arabic*)},leftmargin=*,itemsep=4pt,topsep=3pt]
  \item \textbf{Reindexed implication and negation.} The first equations follow from $J\rho=j_A$, preservation of meets by $J$,
and preservation of implication by $\rho$. For implication use
$Q\Rightarrow Q'=\rho(JQ\Rightarrow JQ')$; the negation case also uses
$J\bot=\neg A$ and the Heyting implication identities. Let $\Gamma_a$ be the graph of $a$.
  \item \textbf{Descent along the kernel pair.} Then
$R|_U=\Gamma_a;R;\Gamma_a^\smile$ and
$\Gamma_a^\smile;\Gamma_a$ is the partial identity on $A$. Its box is $A\Rightarrow(-)$, proving the modal equation by relation
composition. Kernel pairs restrict as stated.
  \item \textbf{Frobenius law for repeated announcements.} Finally right adjoints compose,
and Frobenius with $a^*\forall_a=\operatorname{id}$ gives
$\exists_aB=A\wedge\forall_aB$.
\end{enumerate}
\end{proof}

\noindent Repeated public restriction is computed by a composite of comprehension inclusions.
\begin{center}
\begin{tikzcd}[column sep=large,row sep=large]
U\cap Q\arrow[r,hook]\arrow[d,hook]
 &U\arrow[d,hook,"m_U"]\\
 Q\arrow[r,hook]&W
\end{tikzcd}
\end{center}

\begin{remark}[Fixed predicates and re-evaluated epistemic assertions]
\label{rem:egt-announcement-reevaluation}
For an inclusion $a:A\hookrightarrow W$, the equation
$\forall_a a^*P=(A\Rightarrow P)$ is an equation about a fixed predicate.
It does not imply that an epistemic formula remains true when its own
truth set is announced.

\end{remark}

\begin{example}[A truthful announcement that ceases to be true]
\label{ex:egt-announcement-success}
Take $W=\{x,y\}$, one indiscriminate agent, and $p=\{x\}$.
The formula $\varphi=p\wedge\neg Kp$ holds at $x$. Announcing its
truth set leaves $\{x\}$, where $Kp$ holds and $\varphi$ is false.
The factual predicate $p$ has not changed; its epistemic context has.
\end{example}

\begin{corollary}[Dynamic semantic testing]
\label{prop:egt-dynamic-semantic-testing}
For $\Phi=\operatorname{Test}_P(t_1,\ldots,t_n)$,
\begin{equation}\label{eq:egt-dynamic-test-reduction}
 [!a]\operatorname{Test}_P(t_1a,\ldots,t_na)=A\Rightarrow\Phi,
\end{equation}
and
\begin{equation}\label{eq:egt-dynamic-epistemic-test}
 [!a]K_i^U\operatorname{Test}_P(t_1a,\ldots,t_na)
 =A\Rightarrow K_i(A\Rightarrow\Phi).
\end{equation}
\end{corollary}
\begin{proof}
\leavevmode
\begin{enumerate}[label=\textup{(\arabic*)},leftmargin=*,itemsep=4pt,topsep=3pt]
  \item \textbf{Functoriality and substitution.} Apply substitution of semantic tests and the preceding reduction equations.
\end{enumerate}
\end{proof}

\begin{corollary}[Fixed points under public restriction]
\label{cor:egt-update-fixedpoints}
For $a:U\hookrightarrow W$ and monotone $F$ on $\mathsf{Pred}(U)$,
\[
 \forall_a(\mu F)=\mu(\forall_aFa^*),\qquad
 \forall_a(\nu F)=\nu(\forall_aFa^*).
\]
\end{corollary}
\begin{proof}
\leavevmode
\begin{enumerate}[label=\textup{(\arabic*)},leftmargin=*,itemsep=4pt,topsep=3pt]
  \item \textbf{Repeated-announcement fixed-point calculation.} Use $a^*\forall_a=\operatorname{id}$ in
Lemma~\ref{lem:egt-retraction-fixedpoints}.
\end{enumerate}
\end{proof}

\subsubsection{Epistemic frames and categorical update}
\label{sec:egt-frame-updates}

The comprehension inclusion for public restriction is a map of
state carriers equipped with an accessibility comparison. General action
updates have the same shape, but carry an event coordinate as well. To
relate their assertions before and after update, we must distinguish
ordinary reindexing of dependent families from exact lifting of
accessible alternatives.

\begin{definition}[Relation frames and their maps]
\label{def:egt-relation-frame-category}
An object $M=(W,(R_i)_{i\in N})$ of $\mathbf{Fr}_N$ is a set with an
$N$-indexed family of binary relations. A morphism
$f:(U,(T_i))\to(W,(R_i))$ is a function satisfying
\[
 T_i\subseteq(f\times f)^*R_i\qquad(i\in N).
\]
Write $|M|=W$ for the underlying set. S5 frames form the full subcategory
whose relations are equivalence relations. The definitions also permit the
paper's non-S5 belief relations.
\end{definition}

\begin{proposition}[Pullbacks and the existing indexed semantics]
\label{prop:egt-frame-pullbacks}
The category $\mathbf{Fr}_N$ has finite limits. For maps $f:M\to K$ and
$g:L\to K$, its pullback has underlying set $|M|\times_{|K|}|L|$
and relations
\[
 R_i^{M\times_KL}
 = (\pi_M\times\pi_M)^*R_i^M
   \cap(\pi_L\times\pi_L)^*R_i^L.
\]
These limits restrict to S5 frames. An epistemic refinement update of $M$
is an object $p:U\to M$ of $\mathbf{Fr}_N/M$. Its ordinary and linear families
are transported by the existing functors
\[
 |p|^*:\mathsf T(|M|)\to\mathsf T(|U|),\qquad
 |p|^*:\mathsf L(|M|)\to\mathsf L(|U|).
\]
Successive updates compose as morphisms. Pullback gives base change
between the slice categories, coherently up to their canonical isomorphisms.
\end{proposition}
\begin{proof}
\leavevmode
\begin{enumerate}[label=\textup{(\arabic*)},leftmargin=*,itemsep=4pt,topsep=3pt]
  \item \textbf{Predicate calculation.} The set-theoretic pullback carries the displayed relations so that both
projections preserve relations. Any compatible cone has a unique underlying
map to that pullback. Since both legs preserve each relation, this map
preserves their intersection.
  \item \textbf{Finite limits of equivalence frames.} This proves the universal property. A singleton
with reflexive relations is terminal; products and equalizers follow. Inverse images and intersections preserve equivalence relations.
  \item \textbf{Coherence of epistemic pullbacks.} The remaining statements follow by composition of the underlying functions
and the universal property of pullback. The canonical reassociation maps
preserve the projected tuples, so their associativity and unit diagrams
commute. The same uniqueness gives the coherence of iterated base change.
  \item \textbf{Base-change comparison.} The indexed categories of
Section~\ref{sec:egt-dependent-semantics}
and Beck--Chevalley comparisons apply to these underlying set pullbacks.
\end{enumerate}
\end{proof}

A public announcement can delete accessible alternatives, so its update projection is treated as a frame map with the separately stated boundedness criteria.
For a forward frame map $p$ one has the comparison
\[
 |p|^*[R_i]\,Q\ \subseteq\ [T_i]|p|^*Q.
\]
The forward frame-map comparison gives an inclusion for each
predicate $Q$: every $T_i$-successor maps to an $R_i$-successor.
Equality requires lifting abstract successors as well, which is the
bounded-morphism condition. The exact equations in
Section~\ref{sec:egt-modal-transport} make this distinction explicit,
in line with the categorical treatment of epistemic updates in
\citep{Kishida2017CategoriesDEL}.

\subsubsection{Epistemic action-model updates as dependent comprehension}
\label{sec:egt-action-model-update}

Public restriction is the singleton-event case of epistemic
product update \citet{BaltagMossSolecki1998PublicAnnouncements}.
With several possible events, an agent may distinguish some occurrences
but not others, and the events may change ordinary observables.
The update therefore retains enabled world--event pairs and constructs
their new accessibility from both the old state observations and the
observations of the action event.

\begin{definition}[Semantic epistemic action model]
\label{def:egt-epistemic-action-model}
Let $W$ be a state set with relations $R_i\hookrightarrow W^2$ for
$i\in N$. An \textbf{epistemic action model over $W$} consists of
an action-occurrence set $E$, relations $S_i\hookrightarrow E^2$, and
an enabling subobject $P\hookrightarrow W\times E$. Its fibre at
$e\in E$ is $P_e=\{w:P(w,e)\}\hookrightarrow W$.
The relations $R_i$ and $S_i$ may be arbitrary qualitative possibility
relations. If they are observation kernel pairs, they describe knowledge.
An updated observable of type $Z$ is a term $v^+:\{P\}\to Z$; the
unchanged-value case is $v^+=v p$ for $v:W\to Z$, where
$p:\{P\}\to W$ is the projection. A state-transition presentation may
instead provide $t:\{P\}\to W'$ and set $v^+=v't$ for $v':W'\to Z$.
\end{definition}

\begin{construction}[Epistemic product update]
\label{con:egt-action-product}
Take comprehension and its two projections,
\[
 j:U=\{P\}\lhook\joinrel\longrightarrow W\times E,
 \qquad p:U\to W,\quad q:U\to E.
\]
Write $R_i\boxtimes S_i$ for the product relation on $W\times E$,
\[
 (R_i\boxtimes S_i)((w,e),(w',e'))
 \ \Longleftrightarrow\ R_i(w,w')\wedge S_i(e,e').
\]
Define the updated relation on $U$ by
\begin{equation}\label{eq:egt-action-updated-relation}
 R_i^{\mathcal E}=(j\times j)^*(R_i\boxtimes S_i)
       =(p\times p)^*R_i\wedge(q\times q)^*S_i.
\end{equation}
For each $e$ write $a_e:P_e\hookrightarrow W$ for comprehension and
$j_e:P_e\to U$ for $w\mapsto(w,e)$. For a predicate $Q\hookrightarrow U$
define its dynamic box and diamond on $W$ by
\begin{equation}\label{eq:egt-action-dynamic-modalities}
 [\mathcal E,e]Q=\forall_{a_e}j_e^*Q,\qquad
 \langle\mathcal E,e\rangle Q=\exists_{a_e}j_e^*Q.
\end{equation}
The action model, its comprehension, and the updated relations are
ordinary dependent objects and subobjects in the model of
Theorem~\ref{thm:egt-dependent-model}. They can equally be constructed
in the uncompleted $\mathbf{Set}$ assertion doctrine.
\end{construction}

\noindent Epistemic product update is a comprehension of enabled world--event pairs.
\begin{center}
\begin{tikzcd}[column sep=large,row sep=large]
U\arrow[r,hook]\arrow[d,"p"']
 &W\times E\arrow[d,"\pi_1"]\\
 W\arrow[r,equal]&W
\end{tikzcd}
\end{center}

\begin{proposition}[Action models in an economic state context]
\label{prop:egt-action-state-tower}
For $s:W\to\Sigma^G$, the canonical state context after the action
update is $sp:U\to\Sigma^G$. Its updated relations, action-indexed
predicates, and postcondition terms belong to the ordinary dependent
fibres above $U$. For a linear family $A$ over $W$, its unchanged
interface after the action is $p^*A$; a transition-indexed residual
interface instead uses its exhibited post-state map $t:U\to W'$ and
the dependent family over $W'$. These reindexings commute with the
ordinary comparison tests under the substitution laws of
Theorem~\ref{thm:fib-linear-logical-compatibility}. If the action changes the economic reference profile, the changed profile is the specified post-update term $s^+:U\to\Sigma^G$; the composite $sp$ denotes the unchanged reference.
\end{proposition}
\begin{proof}
\leavevmode
\begin{enumerate}[label=\textup{(\arabic*)},leftmargin=*,itemsep=4pt,topsep=3pt]
  \item \textbf{Composition.} The composite $sp$ is an object of $\mathbf{Set}/\Sigma^G$. Every relation and postcondition on $U$ is a subobject or term in
$\mathsf T(U)$; the linear reindexings are index precomposition. Theorem~\ref{thm:fib-linear-logical-compatibility} gives the asserted
comparison-test equations.
  \item \textbf{Policy-level comparison.} A post-update profile is simply a term of
the constant profile family on $U$, which supplies its typed state
context without imposing any equality with $sp$.
\end{enumerate}
\end{proof}

\begin{theorem}[Action-model semantics, base change, and public specialization]
\label{thm:egt-action-model-semantics}
For every $Q\hookrightarrow U$ and $w\in W$,
\[
 \begin{aligned}
 [\mathcal E,e]Q(w)&\Longleftrightarrow P_e(w)\Rightarrow Q(w,e),\\
 \langle\mathcal E,e\rangle Q(w)&\Longleftrightarrow
                                      P_e(w)\wedge Q(w,e).
 \end{aligned}
\]
If $E=1$, $S_i=1^2$, and $P_e=A\hookrightarrow W$, the updated
relation and modalities are respectively $R_i|_A$, $[!a]$, and
$\langle!a\rangle$ of Section~\ref{sec:egt-dynamic-update}.
For $f:W'\to W$, pull back $R_i$ along $f\times f$ and $P$ along
$f\times1_E$. The resulting updated carrier $U'$ is canonically
$W'\times_WU$, its relations are inverse images of
$R_i^{\mathcal E}$, and the action boxes and diamonds commute with
this pullback. These comparisons are coherent under successive
substitutions.
\end{theorem}
\begin{proof}
\leavevmode
\begin{enumerate}[label=\textup{(\arabic*)},leftmargin=*,itemsep=4pt,topsep=3pt]
  \item \textbf{Predicate calculation.} The fibre of $a_e$ at $w$ is a singleton if $P_e(w)$ holds and empty
otherwise. The quantifier formulas of
Theorem~\ref{thm:fib-comprehension} give the two equivalences. For $E=1$ comprehension identifies $U$ with $A$, and
Equation~\eqref{eq:egt-action-updated-relation} reduces to the ordinary
restriction of $R_i$.
  \item \textbf{Base-change comparison.} The dynamic modalities are exactly the two
adjoints along $a$. The universal property of pullback identifies
$U'$ with $W'\times_WU$, while relation substitution and both
Beck--Chevalley comparisons give the remaining equations. Repeated
substitution identifies the same indexed tuples, proving coherence.
\end{enumerate}
\end{proof}

\noindent Changing the economic state context commutes with epistemic product update.
\begin{center}
\begin{tikzcd}[column sep=large,row sep=large]
U'\arrow[r]\arrow[d]\arrow[dr,phantom,"\lrcorner",very near start]
 &U\arrow[d,"p"]\\
 W'\arrow[r,"f"']&W
\end{tikzcd}
\end{center}

\begin{theorem}[Presentation of every semantic refinement by an action model]
\label{thm:egt-action-presentation}
Every object $p:(U,(T_i))\to M=(W,(R_i))$ of $\mathbf{Fr}_N/M$ is
isomorphic over $M$ to a semantic action-model update of $M$.
Conversely, every action-model update has such a projection.
\end{theorem}
\begin{proof}
\leavevmode
\begin{enumerate}[label=\textup{(\arabic*)},leftmargin=*,itemsep=4pt,topsep=3pt]
  \item \textbf{Causal execution.} For the forward direction take the event set to be $E=U$, its uncertainty
relations to be $S_i=T_i$, and its enabling subobject to be the graph
\[
 P=\{(w,u)\in W\times U:w=p(u)\}.
\]
The updated carrier is $P$, with canonical bijection
$u\mapsto(p(u),u)$. Under this bijection its $i$-relation is
\[
 (p\times p)^*R_i\cap T_i=T_i,
\]
because $p$ preserves relations. The first projection becomes $p$.
  \item \textbf{Reverse implication.} Conversely, the product-update relation
$(\pi_W\times\pi_W)^*R_i\cap(\pi_E\times\pi_E)^*S_i$
is contained in the first factor, so its projection is a frame map.
\end{enumerate}
\end{proof}

The presentation result is relative to a specified base model;
its action-event set need not be finite. For finite $W,U$, the
construction in \cref{con:egt-finite-action-enablement} realizes
its enabling relation and the map $P\to U$ by finite response
protocols. The quotient $U\to U/T_i$ classifies the relation $T_i$ on
the updated carrier $U$ and describes epistemic indistinguishability after the event, not an extra physical wire.
The witness data retain the underlying execution, so the same update can
be compared with an operational protocol when one is supplied.

\begin{construction}[Updates with state-dependent edge conditions]
\label{con:egt-edge-guard-update}
Let $P\subseteq W\times E$ be an enabling carrier and let
$G_i\subseteq P\times P$ be specified edge conditions. Put
\[
 T_i=(\pi_W\times\pi_W)^*R_i\cap G_i.
\]
This is another object of $\mathbf{Fr}_N/M$, so Theorem~\ref{thm:egt-action-presentation}
gives an action-model presentation. Ordinary product update is the case
$G_i=(\pi_E\times\pi_E)^*S_i$. State-dependent edge conditions
provide a semantic counterpart of arrow updates. The syntactic arrow-update
literature has its own expressivity and synthesis results
\citep{VanDitmarschEtAl2020ArrowUpdateSynthesis};
those results are not being inferred from this set-based presentation.
\end{construction}

\begin{construction}[Finite operational realization of action enablement]
\label{con:egt-finite-action-enablement}
For finite $W,E$ and an action-enabling predicate
$P\hookrightarrow W\times E$, let $U=\{P\}$, regarded as its finite
set of enabled pairs. Define a strategy
\[
 a_P:D(W)\longrightarrow D(U)
\]
with one negative input event $n_w$ for every $w\in W$ and one
positive event $o_{w,e}$ for every $(w,e)\in U$. Generate the causal
edges $n_w<o_{w,e}$. Different input branches conflict; within a
branch the positive events conflict, so at most one enabled pair is
emitted. Label the negative event by $w$ at the input boundary and
the positive event by $(w,e)$ at the output boundary. Thus the
completed outputs following an input $w$ are exactly the enabled
actions $(w,e)\in U$. For an updated value term $v^+:U\to Z$,
corestrict to its finite image $F=v^+[U]$ and compose $a_P$ with
$\widehat{\bar v^+}:D(U)\to D(F)$. If updated S5 observations are
supplied, their finite image quotients likewise have the response
strategies of Proposition~\ref{prop:egt-functions}.
\end{construction}

\begin{proposition}[Sound operational realization of finite action models]
\label{prop:egt-finite-action-enablement}
The map $a_P$ is a receptive courteous strategy. Its enabled output
relation is precisely $P$ and its composition with any finite-image
updated observable returns that observable on each completed updated
pair. Hence every finite semantic action-model enabling relation and
its specified finite-image postconditions admit explicit concurrent
operational realizations, while the action uncertainty relations
$S_i$ remain ordinary epistemic data until observation protocols are
specified. Countable action presentations admit the same carrier-based
construction; uncountable $\mathbf{Set}$-based action models belong to
the ordinary dependent semantics and are not claimed to be single
countable operational arenas.
\end{proposition}

\begin{proof}
\leavevmode
\begin{enumerate}[label=\textup{(\arabic*)},leftmargin=*,itemsep=4pt,topsep=3pt]
  \item \textbf{Causal execution.} The source has finite principal pasts. Its consistency relation is
hereditary and closed under adjoining causes, so it is an event
structure. Every source configuration maps injectively to a
configuration of $D(W)^\perp\parallel D(U)$.
  \item \textbf{Fibrewise calculation.} At the empty source
configuration every negative input $w$ has exactly the lift $n_w$;
after one input all other negative inputs conflict in the arena. All new immediate causal edges are negative-to-positive. Thus the
map is receptive and courteous.
  \item \textbf{Causal execution.} Every source configuration containing
an output has $n_w$ and precisely one event $o_{w,e}$, with $P(w,e)$;
conversely, each
$(w,e)\in P$ exhibits such a configuration. If no action is enabled at
$w$, the maximal source configuration $\{n_w\}$ has no positive
output.
  \item \textbf{Composition.} Composing with the response strategy for the
corestriction of $v^+$ gives the value $v^+(w,e)$ by
Equation~\eqref{eq:egt-value-state-naturality} and composition of
finite response strategies. The same proof uses countably many
branches and outputs when the carriers are countable; every
individual causal past remains finite.
\end{enumerate}
\end{proof}


\begin{proposition}[Knowledge reduction for epistemic actions]
\label{prop:egt-action-knowledge-reduction}
Put $K_i=[R_i]$ on $W$ and $K_i^{\mathcal E}=[R_i^{\mathcal E}]$
on $U$. For every $Q\hookrightarrow U$,
\begin{equation}\label{eq:egt-action-knowledge-reduction}
 [\mathcal E,e]K_i^{\mathcal E}Q
 =P_e\Rightarrow
   \bigwedge_{f:\,S_i(e,f)}K_i([\mathcal E,f]Q).
\end{equation}
The empty meet is $\top$. Ordinary conjunction and unchanged atomic
predicates have the reductions
\[
 \begin{aligned}
 [\mathcal E,e](Q\wedge Q')
   &=[\mathcal E,e]Q\wedge[\mathcal E,e]Q',\\
 [\mathcal E,e]p^*V&=P_e\Rightarrow V.
 \end{aligned}
\]
The same predicate doctrine interprets belief modalities for arbitrary
$R_i$, without asserting factivity.
\end{proposition}
\begin{proof}
\leavevmode
\begin{enumerate}[label=\textup{(\arabic*)},leftmargin=*,itemsep=4pt,topsep=3pt]
  \item \textbf{Fibrewise calculation.} At a state $w\in P_e$, the left side requires $Q(w',f)$ for every
$(w',f)\in U$ satisfying $wR_iw'$ and $eS_if$. Reordering the
quantifiers yields
\[
 \forall f\;(eS_if\Rightarrow
   \forall w'\;(wR_iw'\Rightarrow(P_f(w')\Rightarrow Q(w',f)))).
\]
The inner assertion is $K_i([\mathcal E,f]Q)(w)$. At states outside
$P_e$ both sides are true. Preservation of meets by the right adjoint
and $a_e^*p^*V=a_e^*V$ prove the other equations.
\end{enumerate}
\end{proof}

\begin{proposition}[Observation quotients after action update]
\label{prop:egt-action-observation-quotient}
If $R_i=\ker(o_i)$ and $S_i=\ker(b_i)$ for maps
$o_i:W\to O_i$ and $b_i:E\to L_i$, then
\[
 R_i^{\mathcal E}
 =\ker\bigl(\langle o_ip,b_iq\rangle:U\to O_i\times L_i\bigr).
\]
Its image factorization is the effective updated observation quotient.
In particular S5 knowledge, the invariant-predicate adjunctions, and
the common-knowledge quotient construction remain available after the
update, including when $U$ is empty.
\end{proposition}
\begin{proof}
\leavevmode
\begin{enumerate}[label=\textup{(\arabic*)},leftmargin=*,itemsep=4pt,topsep=3pt]
  \item \textbf{Descent along the kernel pair.} Two updated pairs have the same displayed observation exactly when
their old observations and action observations agree. This is the
conjunction of their two kernel-pair relations, which is
Equation~\eqref{eq:egt-action-updated-relation}. Image factorization
and Corollary~\ref{prop:egt-observation-interior} give the effective
quotient.
  \item \textbf{Observation and knowledge.} Equivalence relations and their generated common-knowledge
relation restrict to the new carrier as constructed, with the empty
case interpreted in the ordinary predicate doctrine.
\end{enumerate}
\end{proof}

\begin{example}[Private epistemic information cannot be reduced to announcement]
Let $W=\{w_0,w_1\}$, let two owners initially have the indiscrete
observation equivalence, and take $E=\{e_0,e_1\}$ with
$P_{e_k}=\{w_k\}$. Give owner $1$ the diagonal action relation and
owner $2$ the indiscrete action relation. Then
$U=\{(w_0,e_0),(w_1,e_1)\}$. Owner $1$'s updated relation is equality,
while owner $2$'s still relates both points. Neither original state is
removed. Hence this is a genuine private-information update, not the
public restriction of a subset of $W$.
\end{example}

\begin{construction}[Sequential semantic action update]
\label{con:egt-action-sequential}
Let $\mathcal E=(E,(S_i),P)$ update a relation model on $W$ to
$U=\{P\}$. Let $\mathcal F=(F,(T_i),Q)$ be an action model over
$U$, with $Q\hookrightarrow U\times F$. Define the composite
$\mathcal E\diamond\mathcal F$ over $W$ with action set $E\times F$,
action relations $S_i\boxtimes T_i$, and enabling predicate
\begin{equation}\label{eq:egt-action-composite-precondition}
 P\diamond Q\,(w,e,f)
 \Longleftrightarrow
 P(w,e)\wedge Q((w,e),f).
\end{equation}
The expression involving $Q$ is evaluated only in the fibre where
$P(w,e)$ holds, equivalently by the image under comprehension of the
second enabling object. Updated value terms on the two-stage carrier
are transported by the same canonical reassociation.
\end{construction}

\begin{theorem}[Coherent composition of epistemic actions]
\label{thm:egt-action-composition}
There is a canonical isomorphism of relation models and their retained
updated ordinary terms
\begin{equation}\label{eq:egt-action-composition}
 (W\odot\mathcal E)\odot\mathcal F
 \cong W\odot(\mathcal E\diamond\mathcal F),
 \qquad ((w,e),f)\longmapsto(w,(e,f)).
\end{equation}
It satisfies associativity and the singleton everywhere-enabled unit
laws, including coherent substitution in the reference state. The
corresponding dynamic modalities satisfy
\[
 [\mathcal E,e][\mathcal F,f]Q
   =[\mathcal E\diamond\mathcal F,(e,f)]Q
\]
under the indicated identification of final predicates.
\end{theorem}
\begin{proof}
\leavevmode
\begin{enumerate}[label=\textup{(\arabic*)},leftmargin=*,itemsep=4pt,topsep=3pt]
  \item \textbf{Predicate calculation.} The two carriers are defined by the identical conjunction of the two
enablement tests after reassociation. The two-stage relation between
$((w,e),f)$ and $((w',e'),f')$ is
\[
 R_i(w,w')\wedge S_i(e,e')\wedge T_i(f,f'),
\]
which is the relation of the combined action model. Updated terms are
the same functions on these identified points.
  \item \textbf{Adjunction and transposition.} Reassociation of
triples and quadruples obeys Cartesian coherence, including the unit
comparison. For the dynamic box, two successive right adjoints along
the inclusions of the action components agree with the right adjoint
along their composite by uniqueness of right adjoints; the elementwise
precondition is Equation~\eqref{eq:egt-action-composite-precondition}.
\end{enumerate}
\end{proof}

\noindent Two consecutive epistemic action updates associate by matching the intermediate state.
\begin{center}
\begin{tikzcd}[column sep=large,row sep=large]
(W\odot\mathcal E)\odot\mathcal F\arrow[r,"\cong"]\arrow[d,"p_1"']
 &W\odot(\mathcal E\diamond\mathcal F)\arrow[d,"p"]\\
 W\odot\mathcal E\arrow[r,"p"']&W
\end{tikzcd}
\end{center}

\begin{remark}[Scope of epistemic action models]
\label{rem:egt-action-scope}
The construction is fully semantic: $P$, $R_i$, $S_i$, and updated
observable terms are ordinary typed data. Syntactic precondition or
postcondition languages are not required. The update changes the
state carrier as well as both endpoints of each accessibility relation;
it is not the reference pullback of
Proposition~\ref{prop:egt-reference-bc}. A state map that forgets
accessible alternatives need not preserve knowledge formulas merely
because it preserves the underlying updated carrier. The bounded
morphism and alternative-lifting results of
Section~\ref{sec:egt-modal-transport} address that distinction.
\end{remark}

\begin{corollary}[Action models transport across bounded state quotients]
\label{cor:egt-action-bounded-quotient}
Let $f:X\twoheadrightarrow Y$ be a surjective bounded morphism from
relations $R_i$ to $\bar R_i$. Give $Y$ an epistemic action model with
action set $E$, action relations $S_i$, and enabling subobject
$\bar P\hookrightarrow Y\times E$. Pull it back along $f\times1_E$
to form the action model over $X$. The induced map on updated carriers
\[
 f^+:X\odot\mathcal E\longrightarrow Y\odot\bar{\mathcal E},
 \qquad(x,e)\longmapsto(f(x),e)
\]
is a surjective bounded morphism for the updated relations. For every
$e\in E$ and $\bar Q\hookrightarrow Y\odot\bar{\mathcal E}$,
\[
 \begin{aligned}
 f^*[\bar{\mathcal E},e]\bar Q
   &=[\mathcal E,e](f^+)^*\bar Q,\\
 f^*\langle\bar{\mathcal E},e\rangle\bar Q
   &=\langle\mathcal E,e\rangle(f^+)^*\bar Q.
 \end{aligned}
\]
Hence it preserves and reflects every modal fixed-point assertion with
atomic tests descended through $f^+$, including updated knowledge
assertions and their dynamic action tests.
\end{corollary}
\begin{proof}
\leavevmode
\begin{enumerate}[label=\textup{(\arabic*)},leftmargin=*,itemsep=4pt,topsep=3pt]
  \item \textbf{Lifting along the surjection.} For an enabled $(y,e)$, surjectivity of $f$ supplies $x$ with $f(x)=y$;
the pullback enabling predicate makes $(x,e)$ enabled. Thus $f^+$ is
surjective. Forward preservation of accessibility is inherited from
$f$, with the unchanged action relation $S_i$.
  \item \textbf{Construction of the comparison.} Given a reduced
updated alternative $(y',e')$ accessible from $(f(x),e)$, the back
clause of $f$ supplies $x'$ with $xR_ix'$ and $f(x')=y'$. Since
$\bar P(y',e')$, its pullback $P(x',e')$ holds. Thus $(x',e')$ is the
required updated alternative.
  \item \textbf{Base-change comparison.} The action-component comprehensions
of the two models form pullback squares along $f$; both
Beck--Chevalley identities of Theorem~\ref{thm:fib-comprehension}
give the dynamic box and diamond equations. Apply
Theorem~\ref{thm:egt-modal-adequacy} to $f^+$.
\end{enumerate}
\end{proof}

\subsubsection{Lifting alternatives and preserving fixed points}
\label{sec:egt-modal-transport}

A map between epistemic models may preserve actual states and
forward accessibility while still losing alternatives quantified over
by a modality. Exact modal transport therefore needs two comparisons:
every concrete successor has an abstract image, and every abstract
successor at an image state lifts to a concrete successor at the same
reference. These are the forward and back conditions of a bounded
morphism. With both, the modal and fixed-point predicates descend
without deleting relevant alternatives.

\begin{definition}[Bounded morphism]
\label{def:egt-bounded-morphism}
A map $q:X\to Y$ between families of relation frames is \textbf{bounded} if,
for every relation index $a$,
\[
 xR_ax'\Longrightarrow q(x)\bar R_aq(x'),
\]
and every $q(x)\bar R_ay'$ has a lift $xR_ax'$ with $q(x')=y'$.
For interpreted atomic predicates require $P_X=q^*P_Y$.
The second clause is equivalently fibrewise coverage by the canonical comparison to the cartesian lift,
\[
 R_a\longrightarrow X\times_{Y,\bar p_1}\bar R_a,
 \qquad(x,x')\longmapsto(x,(q(x),q(x'))).
\]
\end{definition}

\[
\begin{tikzcd}[column sep=large]
 R_a\arrow[r,two heads]\arrow[dr,"p_1"']&
 X\times_Y\bar R_a\arrow[d]\arrow[r]&\bar R_a\arrow[d,"\bar p_1"]\\
 &X\arrow[r,"q"']&Y
\end{tikzcd}
\]

\begin{theorem}[Modal fixed-point adequacy]
\label{thm:egt-modal-adequacy}
Let $q:X\to Y$ be a bounded morphism. Interpret every free predicate
variable by a valuation $\rho$ on $Y$ and by its componentwise pullback
$q^*\rho$ on $X$; interpret each designated atomic predicate on $X$ as
the pullback of its interpretation on $Y$. Then every formula of the
positive modal fixed-point language satisfies the naturality equation
\[
 \llbracket\varphi\rrbracket_X(q^*\rho)
 =q^*\llbracket\varphi\rrbracket_Y(\rho).
\]
This includes nested least and greatest fixed points. Validity is
preserved; if $q$ is surjective it is also reflected.
\end{theorem}

\noindent Forward preservation of accessibility is a commuting square on relation witnesses.
\begin{center}
\begin{tikzcd}[column sep=large,row sep=large]
R_i\arrow[r,"f\times f"]\arrow[d,hook]
 &S_i\arrow[d,hook]\\
 X\times X\arrow[r,"f\times f"']
 &Y\times Y
\end{tikzcd}
\end{center}


\begin{proof}
\leavevmode
\begin{enumerate}[label=\textup{(\arabic*)},leftmargin=*,itemsep=4pt,topsep=3pt]
  \item \textbf{Construction of the comparison.} Fix the valuation $\rho$ on the target and its pulled-back valuation
$q^*\rho$ on the source. Induct on formulas in the assertion fibres. The variable and atomic cases follow from these interpretations by
definition.
  \item \textbf{Lifting through the comparison.} Reindexing preserves the remaining logical operations by
Theorem~\ref{thm:fib-comprehension}. The comparison to the cartesian
alternative context has the displayed fibrewise lifts; quantifying those
fibres gives
$[R_a]q^*=q^*[\bar R_a]$ and
$\langle R_a\rangle q^*=q^*\langle\bar R_a\rangle$ by respectively
universal and existential quantification over the corresponding fibers.
  \item \textbf{Fixed-point argument.} For a fixed-point body the induction hypothesis gives
$F_Xq^*=q^*F_Y$. Since $q^*\mu F_Y$ is fixed,
$\mu F_X\leq q^*\mu F_Y$. Put $Q=\forall_q\mu F_X$.
  \item \textbf{Adjunction and transposition.} Then $q^*Q\leq\mu F_X$, whence
$q^*F_YQ=F_Xq^*Q\leq\mu F_X$. Adjunction gives $F_YQ\leq Q$;
leastness and the counit give
$q^*\mu F_Y\leq q^*Q\leq\mu F_X$. For greatest fixed points put $Q=\exists_q\nu F_X$.
  \item \textbf{Unit and identity laws.} The unit gives $\nu F_X\leq q^*Q$, hence
$\nu F_X=F_X\nu F_X\leq q^*F_YQ$. Adjunction gives $Q\leq F_YQ$, and coinduction gives
$\nu F_X\leq q^*Q\leq q^*\nu F_Y$. The reverse inequality holds because $q^*\nu F_Y$ is fixed.
  \item \textbf{Fixed-point argument.} For nested fixed points, first apply structural induction to the body
with a designated free variable $p$: for every target candidate
$U\in\mathsf{Pred}(Y)$ its source instance uses $q^*U$, giving
$F_X(q^*U)=q^*F_Y(U)$ with all other free variables interpreted by
$q^*\rho$. The two adjunction arguments above then establish the
fixed-point comparison for the binder of $p$.
  \item \textbf{Descent along the kernel pair.} Inducting on binding depth
proves the displayed equation for nested least and greatest fixed
points, without treating arbitrary source valuations as descending. Finally, a surjective $q$ reflects the top predicate.
\end{enumerate}
\end{proof}


\subsection{Public announcement in the agent calculus}
\label{agt-sec-public}

\subsubsection{Restricting a situated model}

Let $M$ be a situated model on $W$, and let $\varphi$ be an assertion
interpreted in $M$. Form its comprehension
\[
 j:U=\{w\in W\mid M,w\models\varphi\}\hookrightarrow W.
\]
The updated model $M|\varphi$ is carried by $U$. Its valuations
and linear families are restricted, and observations become $o_i j$,
corestricted to their images when appropriate. Accessibility is
$R_i^U=(j\times j)^*R_i$, so excluded worlds disappear from both
positions of the relation. Strong monoidality of linear restriction
preserves agent types and their composition; only the available
epistemic context changes.

The dynamic type operation is $\Pi_j^\ell$. At a state removed by
the announcement it is the empty additive product $\Top$; at a
retained state it is the post-update fibre. The tensor unit $\I$
cannot be substituted for $\Top$: these are different limits and
resource interpretations.

\begin{theorem}[Perspective after announcement]
\label{agt-thm-announcement-perspective}
For a linear family $A$ over $U$, there is a natural isomorphism
\[
 K_{o j}^\ell A\cong j^*K_o^\ell\Pi_j^\ell A.
\]
It is coherent with successive public restrictions and the monoidal
perspective comparisons.
\end{theorem}
\begin{proof}
\leavevmode
\begin{enumerate}[label=\textup{(\arabic*)},leftmargin=*,itemsep=4pt,topsep=3pt]
  \item \textbf{Adjunction and transposition.} Right adjoints compose: $\Pi_{oj}^\ell\cong\Pi_o^\ell\Pi_j^\ell$. Also $(oj)^*=j^*o^*$. Their composite gives the displayed isomorphism.
  \item \textbf{Fibrewise calculation.} At $u\in U$, the right side takes a product over all old alternatives,
with terminal factors for alternatives outside $U$. Its nonterminal factors
are exactly $A_v$ for $v\in U$ with $o(v)=o(u)$. Product projections prove
the comparison and its coherence. Tensor comparisons on either side are
obtained by tensoring these same remaining projections.
\end{enumerate}
\end{proof}

\subsubsection{Dynamic formulas and reduction}

The interpretation of a public-announcement formula is
\begin{equation}
 \sem{[!\varphi]\psi}_M
 =\forall_j\sem{\psi}_{M|\varphi}.
 \label{agt-eq-public-semantics}
\end{equation}
The formula $\psi$ is always interpreted in the updated model,
using its restricted accessibility relations. This rule continues to
apply when the postcondition contains knowledge, another announcement,
or a common-knowledge fixed point. Otherwise the reduction would
silently evaluate an epistemic operator against the old alternatives.

\begin{theorem}[Public-announcement reduction laws]
\label{agt-thm-public-reduction}
For an atom $p$ whose valuation is restricted by the update,
\begin{align*}
 [!\varphi]p&\leftrightarrow(\varphi\Rightarrow p),\\
 [!\varphi](\psi\wedge\chi)
 &\leftrightarrow[!\varphi]\psi\wedge[!\varphi]\chi,\\
 [!\varphi]\neg\psi
 &\leftrightarrow\bigl(\varphi\Rightarrow\neg[!\varphi]\psi\bigr),\\
 [!\varphi]K_i\psi
 &\leftrightarrow\bigl(\varphi\Rightarrow K_i([!\varphi]\psi)\bigr),\\
 [!\varphi][!\psi]\chi
 &\leftrightarrow[!(\varphi\wedge[!\varphi]\psi)]\chi.
\end{align*}
\end{theorem}
\begin{proof}
\leavevmode
\begin{enumerate}[label=\textup{(\arabic*)},leftmargin=*,itemsep=4pt,topsep=3pt]
  \item \textbf{Fibrewise calculation.} At each stage interpret the postformula in the restricted model and
let $j:U\hookrightarrow W$ be the current announcement comprehension. The atomic, Boolean, and knowledge equations follow, respectively,
from inverse image, the monic adjunction identities, and
\eqref{eq:egt-public-update-modal-reduction} of
\cref{thm:egt-public-update-reduction}.
  \item \textbf{Iteration of announcement reduction.} Thus they apply inductively to
postformulas involving knowledge as well as fixed pre-update
predicates. For successive announcements the second comprehension
inside $U$ composes with $j$ to the inclusion classified in $W$ by
$\varphi\wedge[!\varphi]\psi$; the right adjoints of the two
inclusions compose. This gives the final equation and retains the
re-evaluation of epistemic formulas after the first update.
\end{enumerate}
\end{proof}

The reduction equations follow from comprehension and the
right-adjoint interpretation of announcement; compare
\citet{Plaza2007PublicCommunications,BaltagMossSolecki1998PublicAnnouncements}.
They express the logical operation independently of a physical test or
message that might implement it.

\begin{example}[Re-evaluating an epistemic announcement]
Take two indistinguishable worlds, with $p$ true at just one. At that world
$p\wedge\neg K_ip$ is true, so its announcement is enabled. Afterward only
the $p$-world remains and $K_ip$ holds. Thus the announced epistemic formula
itself becomes false. The update recomputes epistemic formulas in the restricted model; the example shows the possible failure of $[!\varphi]\varphi$.
\end{example}

\subsubsection{The effect on policies and resources}

Public restriction can change which observation-based policies are
uniformly legal. For $F\subseteq W\times M$, restricting to $U$
removes some worlds from surviving observation fibres. With
$O_U=\im(oj)$ and $r:O_U\hookrightarrow O$, every former legal
policy restricts to one over the new quotient, giving the map below.
\[
 \Pol_o(F)\longrightarrow\Pol_{oj}((j\times1)^*F),\qquad
 \pi\longmapsto\pi r.
\]
The restriction map need not be injective: policies distinguished
only by observations absent from $U$ become indistinguishable. New
uniformly legal actions may also appear after a counterexample world is
removed. By contrast, linear execution resources are merely reindexed;
a public update does not create a resource or authorization token.

So far the update is logical. To regard it as a physical public
announcement, an enabled test or message must have execution witnesses
whose successful image is exactly $U$. The finite communication protocols
in the next part provide such realizations and make it possible to compare
the epistemic update with the complete policies that produced it.

\begin{proposition}[Observation-based implementability of a public test]
\label{agt-prop-test-implementation}
Let $o:W\twoheadrightarrow O$ be an effective observation and
$P\subseteq W$.  The following are equivalent:
\begin{enumerate}[label=(\roman*),nosep]
\item $P$ is the inverse image $o^{-1}(B)$ of a predicate $B\subseteq O$;
\item $P$ is constant on the kernel pair of $o$;
\item its characteristic function factors as $1_P=\chi o$ for
      $\chi:O\to2$.
\end{enumerate}
For finite $O$, these equivalent descriptions give a receptive courteous
protocol accepting the observed input and emitting a success event precisely
at worlds in $P$.  Its successful-execution witness image is the
comprehension $P\hookrightarrow W$.
\end{proposition}
\begin{proof}
\leavevmode
\begin{enumerate}[label=\textup{(\arabic*)},leftmargin=*,itemsep=4pt,topsep=3pt]
  \item \textbf{Descent along the kernel pair.} The first and third descriptions are equivalent by taking the
characteristic predicate. Either makes $P$ constant on the kernel pair.
  \item \textbf{Reverse implication.} Conversely kernel-pair constancy lets $1_P$ descend uniquely along the
effective quotient $o$, producing $\chi$ and $B=\chi^{-1}(\{1\})$. For finite $O$, realize $\chi$ by
$\widehat\chi:D(O)\to D(2)$. Connect it to the gate
$D(2)\to D(1)$ whose source accepts both negative bit inputs but
emits its sole positive success output only following input $1$.
  \item \textbf{Lifting through the comparison.} It is receptive because both inputs lift uniquely, and courteous because
its only added dependency is negative-to-positive. The composite emits
success exactly when $\chi(o(w))=1$; the image of successful witnesses
is therefore $P$. The absent success event is not reinterpreted as a
false operational scalar.
  \item \textbf{Observation and knowledge.} A test of a predicate not constant on the
available observation fibre requires a different actually delivered signal,
not an implicit truth oracle.
\end{enumerate}
\end{proof}

\subsection{General epistemic action models and delivery}
\label{agt-sec-actions}

\subsubsection{General action models}

Let an action presentation have event set $E$, event relations $S_i$, and
an enabling predicate $\mathsf{pre}\subseteq W\times E$. Its updated
carrier and relations are
\begin{align}
 U&=\{(w,e):\mathsf{pre}(w,e)\},
 \label{agt-eq-actioncarrier}\\
 (w,e)R_i^+(v,f)&\iff wR_iv\ \wedge\ eS_if.
 \label{agt-eq-actionrelation}
\end{align}
The updated relation is the standard dynamic-epistemic product
update \citep{BaltagMossSolecki1998PublicAnnouncements}. In the
observation-generated case, $R_i=\kerpair(o_i)$ and
$S_i=\kerpair(b_i)$, so the effective updated signal is the image
of $(w,e)\mapsto(o_i(w),b_i(e))$. This construction uses the
information actually delivered by the event.

The projection $r:U\to W$ forgets the event, and $r^*$
transports linear families and agent types across that change of state.
An implemented event also retains a witness family and concurrent
realizer, as in Section~\ref{agt-sec-institutions}. Neither an abstract
event relation nor a semantic update alone supplies an operational
communication wire.

For an event $e$, let $j_e:W_e\hookrightarrow W$ be its enabling inclusion,
and $k_e:W_e\to U$ send $w$ to $(w,e)$. Define
\[
 [\mathcal A,e]\psi=\forall_{j_e}k_e^*\sem\psi_{M\otimes\mathcal A}.
\]
Its type-level version is $\Pi_{j_e}^\ell k_e^*$.

\begin{proposition}[Action reduction and composition]
\label{agt-prop-action}
For a finite event presentation,
\[
 [\mathcal A,e]K_i\psi
 \leftrightarrow
 \mathsf{pre}_e\Rightarrow
 \bigwedge_{f:eS_if}K_i([\mathcal A,f]\psi).
\]
Successive action updates compose by retaining pairs of enabled events,
and this composition is associative up to the canonical tuple bijection.
\end{proposition}
\begin{proof}
\leavevmode
\begin{enumerate}[label=\textup{(\arabic*)},leftmargin=*,itemsep=4pt,topsep=3pt]
  \item \textbf{Fibrewise calculation.} At an enabled $(w,e)$, knowledge quantifies over exactly the pairs $(v,f)$
with $wR_iv$, $eS_if$, and $\mathsf{pre}(v,f)$. Grouping first by $f$
gives the displayed conjunction and its guarded boxes. At a disabled $w$
both sides are true.
  \item \textbf{Causal execution.} For a second action with precondition $Q$ on $U$,
the sequential carrier consists of triples $(w,e,f)$ with
$\mathsf{pre}(w,e)\wedge Q((w,e),f)$. Its relation is the conjunction
of the original world relation and the two event relations. Reassociating
triples preserves these tests, the reindexed linear families, and, for
implemented actions, the composite concurrent realizer. Fourfold tuple
reassociation proves the coherence diagrams.
\end{enumerate}
\end{proof}

\begin{proposition}[Public restriction as singleton action update]
\label{prop:public-singleton-action}
For a state predicate $P\hookrightarrow W$, let the action carrier be
$E=1$, its sole event have enabling predicate $P$, and every action
relation be equality on $1$. Its product update is canonically the
comprehension $j:\{P\}\hookrightarrow W$, with relations
$R_i^+= (j\times j)^*R_i$. The pointed action box and diamond are
$\forall_j$ and $\exists_j$; in particular the corresponding
knowledge-reduction law is
\[
 [!P]K_i\psi\ \longleftrightarrow\
 P\Rightarrow K_i([!P]\psi).
\]
\end{proposition}
\begin{proof}
\leavevmode
\begin{enumerate}[label=\textup{(\arabic*)},leftmargin=*,itemsep=4pt,topsep=3pt]
  \item \textbf{Monoidal comparison.} The enabled carrier is
$\{(w,*):P(w)\}\cong\{w:P(w)\}$ by projection. Because both action coordinates are $*$, the product-update relation
reduces to $wR_iv$ on the retained endpoints.
  \item \textbf{Communication and delivery.} Pointed action
quantification along its enabling inclusion is therefore exactly the
comprehension quantification of
\cref{def:egt-public-announcement}; the displayed reduction is the
singleton instance of \cref{agt-prop-action} (equivalently
\cref{thm:egt-public-update-reduction}).
\end{enumerate}
\end{proof}

\subsubsection{Institutional and epistemic composition coherence}

For three execution spans, both bracketings have witness carrier
\[
 \{(e,f,g):t_E(e)=s_F(f),\ t_F(f)=s_G(g)\}.
\]
For three composed executions, both bracketings have the same
matched triple of witnesses. Permission is
$P_E(e)\wedge P_F(f)\wedge P_G(g)$, and the concurrent realizer is
$\rho_G(g)\rho_F(f)\rho_E(e)$. Reassociating the witnesses
preserves these fields. The epistemic comparison replaces record
matching by enabled event matching and uses the same dependent-sum and
-product coherence. Thus repeated updates do not require a new
composition law.
Thus the extraction of action events from execution witnesses is
completed only after the finite causal protocol and session adequacy
results of Section~\ref{sec:paper-10}; see
Construction~\ref{agt-con-witnessed-action-data}.


\part{Operational and strategic communication}

\section{True-Concurrent Communication Between Agents}\label{sec:paper-10}

An epistemic action describes which events agents distinguish,
but it does not by itself implement communication. We now build finite
protocols out of response strategies, data wires, and secured causal
interaction. Complete owner policies specify responses at every
reachable or off-path decision, while the event structure records the
causal prerequisites of actual delivery. Finite binary sessions give a
second presentation of the same discipline through dual endpoints,
residual typing, and witnessed execution. These implementations supply
the operational content needed for the epistemic comparisons of the
following section.

\subsection{Agents implemented by causal decision protocols}

\begin{definition}[Decision generators and generated presentations]
\label{def:egt-decision-generators-and-generated-presentations}
Use the locally fully faithful
nonstrategic embedding $\mathbf j_N$ of
Equation~\eqref{eq:egt-nonstrategic-embedding}. For finite $O,A$ and $i\in N$,
define the singly owned decision presentation
\[
 \mathsf{Dec}_i(O,A):D(O)\longrightarrow D(A)
\]
by
\[
 \Sigma_i=A^O,\qquad \Sigma_j=1\ (j\ne i),\qquad
 \pi\longmapsto\widehat\pi.
\]
Let $\mathbf{DecArena}_N(\mathcal C)$ be the smallest locally full symmetric
monoidal sub-bicategory of $\mathbf{Arena}_N(\mathcal C)$ that contains the
image of $\mathbf j_N$ and all finite decision presentations, is closed under
horizontal composition and tensor, is closed under the cartesian parameter
reindexings of Equation~\eqref{eq:egt-owned-reindexing}, and contains every
presentation invertibly $2$-isomorphic to one already generated. Thus a
restriction or duplication of policy descriptions remains in the generated
class only through an exhibited parameter map; it is not silently absorbed
into the word ``decision.'' This gives an explicit meaning to ``generated by
singly owned decisions and nonstrategic wiring.''
\end{definition}

\begin{definition}[Finite causal decision protocol]
\label{def:egt-finite-causal-decision-protocol}
A \textbf{finite causal decision protocol} consists of a finite acyclic
directed graph of decision nodes $V$, an ownership map $i:V\to N$,
finite nonempty action sets $A_v$ for $v\in V$, a finite set $P_0$
of external ports with finite value sets $(X_p)_{p\in P_0}$, and
specified typed input wires. Each incoming wire $e$ of a node $v$
has source either an external port $p$ (carrier $X_e=X_p$) or an
earlier decision node $u$ (carrier $X_e=A_u$); the wires from nodes
respect the directed acyclic graph. If $I_v$ is the finite set of
incoming wires to $v$, its input carrier and raw observation map are
\[
 H_v=\prod_{e\in I_v}X_e,\qquad
 \widetilde o_v:H_v\to Y_v,\qquad
 H_v\xrightarrow{o_v}O_v\hookrightarrow Y_v,
 \quad O_v=\operatorname{im}(\widetilde o_v).
\]
The effective observation quotient $o_v$ determines the node policies
$\pi_v:O_v\to A_v$ and the complete owner policy spaces
\[
 \Sigma_i=\prod_{v:\,i(v)=i}A_v^{O_v}.
\]
The protocol also specifies an outcome function on its full institutional
record (external data together with all node actions),
\[
 z:\left(\prod_{p\in P_0}X_p\right)\times
    \left(\prod_{v\in V}A_v\right)\longrightarrow Z.
\]
The value of $z$ can be corestricted to its finite image for the
concurrent response protocol. An action at $v$ is
$a_v=\pi_v(o_v((x_e)_{e\in I_v}))$. The graph determines which
observations are available at each decision; recording an action for
an institution does not create an additional observation wire.
The generated arena is constructed and its scheduling independence
proved below.
\end{definition}

A finite causal decision protocol makes the origin of each
observation explicit. Nodes respond to the values arriving along their
incoming wires; independent nodes remain concurrent. The owned
realization assigns a complete contingent function to every decision
node, so a policy includes responses at nodes not encountered during a
particular execution. The ensuing construction connects these policies
to secured event-structure strategies.

\begin{lemma}[Decision-event path reflection in a finite protocol]
\label{lem:egt-protocol-event-path}
Fix the external inputs and a complete owner-policy profile of a finite
causal decision protocol. Form its secured interaction with all wired
occurrences still displayed, before hiding the internal channels. In a
completed configuration let $d_v$ be the positive action-output
occurrence at decision node $v$. If $u\ne v$ and $d_u<d_v$ in the
induced causal order, then the protocol DAG contains a nonempty
directed path from $u$ to $v$. Consequently the chosen action outputs
of incomparable nodes are concurrent in this completed interaction.
Hiding preserves this path restriction on any distinct decision-output
occurrences retained as visible events.
\end{lemma}
\begin{proof}
\leavevmode
\begin{enumerate}[label=\textup{(\arabic*)},leftmargin=*,itemsep=4pt,topsep=3pt]
  \item \textbf{Local causal generators.} Expand each response and
collector in the protocol into its source event structure. On a
completed input branch, a response has only its negative-input to
positive-output dependency. A collector has exactly its individual
input-to-tuple-output dependencies, with no dependency between its
distinct inputs. A copying strategy has an input-to-each-output
dependency, with no order between copied outputs. Disposal has no
output. These descriptions are the explicit constructions of
\cref{prop:egt-functions,prop:egt-data-comonoid,prop:egt-finite-collectors}.
  \item \textbf{Paths through secured matching.} In secured
synchronization, each matched internal event is identified with its
input/output occurrence on the adjoining components. The secured
order is the transitive closure of the component orders, and is
acyclic by the secured-matching condition. Hence every strict causal
comparison in the unhidden completed interaction is witnessed by a
finite chain of local generator edges and matched wire occurrences.
The other occurrence of a matched channel cannot insert an edge from
its output backwards to its input.
  \item \textbf{Reflection into the protocol DAG.} Mark each generator
edge that crosses from a decision output towards the next decision
input by its specified protocol wire. Edges inside copying or
collection boxes only progress towards their outputs. Traversing a
chain from $d_u$ to $d_v$, each passage through a decision component
other than the starting one must enter that component through one of
its specified incoming wires and leave through its response output.
Removing the segments within components therefore gives a directed
walk from $u$ to $v$ in the protocol DAG, which is a directed path
because the graph is acyclic. External inputs are sources and the
final institutional collector is a sink, so neither can introduce a
return path between decision nodes.
  \item \textbf{Incomparability and hiding.} If $u$ and $v$ are
incomparable, neither $d_u<d_v$ nor $d_v<d_u$ can hold. Both
occurrences belong to the completed configuration and are therefore
consistent, so they are concurrent. Hiding restricts the transitive
secured order to visible occurrences; a newly visible comparison
between retained decision outputs would already be a comparison in
the expanded interaction, which has just been excluded.
\end{enumerate}
\end{proof}

\begin{theorem}[Coherence and scheduling independence of finite causal decision protocols]
\label{thm:egt-causal-protocol}
Every finite causal decision protocol determines an open arena in
$\mathbf{DecArena}_N(\mathcal C)$. For each external input and policy profile it
has a unique completed action assignment satisfying
\[
 a_v=\pi_v\bigl(o_v((a_u)_{u\to v},\text{incoming external data})\bigr).
\]
Every evaluation along a linear extension of the causal order yields that
assignment. In the concurrent realization, incomparable decision components
remain unordered; a linear extension is only a reading of the same
symmetric-monoidal diagram.
\end{theorem}

\noindent The causal protocol records outputs of independent nodes before collecting their joint result.
\begin{center}
\begin{tikzcd}[column sep=large,row sep=large]
D(M_v)\otimes D(M_w)\arrow[r,"c_{M_v,M_w}"]\arrow[d,"\widehat f\otimes\widehat g"']
 &D(M_v\times M_w)\arrow[d,"\widehat{f\times g}"]\\
 D(N_v)\otimes D(N_w)\arrow[r,"c_{N_v,N_w}"']&D(N_v\times N_w)
\end{tikzcd}
\end{center}


\begin{proof}
\leavevmode
\begin{enumerate}[label=\textup{(\arabic*)},leftmargin=*,itemsep=4pt,topsep=3pt]
  \item \textbf{Monoidal comparison.} Construct a progressive symmetric-monoidal diagram directly from the protocol
DAG. A data value with $k$ recipients is routed through the iterated comonoid
of Proposition~\ref{prop:egt-data-comonoid}; zero recipients uses disposal. At a node $v$, put $H_v=\prod_{e\in I_v}X_e$.
  \item \textbf{Ownerwise alternatives.} Its implemented
component, evaluated at the owner's complete policy coordinate $\pi_v$,
is the composite
\[
 \bigotimes_{e\in I_v}D(X_e)
 \xrightarrow{c_v}D(H_v)
 \xrightarrow{\widehat{o_v}}D(O_v)
 \xrightarrow{\widehat{\pi_v}}D(A_v).
\]
Here $c_v$ is the finite collector, including $c_0$ when $I_v$ is
empty; the last arrow is the response realized by the singly owned
decision generator.
  \item \textbf{Causal execution.} Copy the
resulting action to exactly its prescribed successor and institutional-record
wires. After all required records are available, use a final collector and
the response strategy for the corestriction of the outcome function to its
finite image. Nodes without incoming wires
use the nullary collector $c_0:I\to D(1)$.
  \item \textbf{Ownerwise alternatives.} Every box in this diagram is either
nonstrategic wiring from $\mathcal C$ or a singly owned decision generator, so
its interpretation lies in $\mathbf{DecArena}_N(\mathcal C)$; grouping the node
parameters by owner gives exactly the displayed $\Sigma_i$.
  \item \textbf{Fixed-point argument.} Because the graph is finite and acyclic, causal height is well founded. At a
minimal node all arguments of its policy equation are external inputs, so its
action is uniquely determined. Inductively, once all nodes of smaller height
have unique actions, the right side of the equation at $v$ has a unique value.
  \item \textbf{Independence from execution ordering.} This constructs a unique completed assignment. The argument mentions no
choice of linear extension, so evaluation along every such extension yields
the same assignment.
  \item \textbf{Causal execution.} The causal-height filtration used above gives a layered factorization of the
progressive monoidal diagram: all boxes at one height are tensored, and the
successive layers are composed.
  \item \textbf{Monoidal comparison.} Within one finite layer, any enumeration is
identified with any other by symmetry; reassociation and exchanges of
incomparable boxes are identified by symmetric-monoidal coherence, while the
constructed data comonoids identify the alternative copying trees by
coassociativity and cocommutativity. Hence every reading along a linear
extension of the causal order denotes the same layered arrow.
  \item \textbf{Monoidal comparison: second comparison.} Any two linear extensions of a finite partial order differ by
adjacent swaps of incomparable nodes: move the first node of one reading
past its incomparable predecessors in the other, and iterate. Each such swap is justified by monoidal interchange, not by imposing an
edge in the event structure. Incomparable decision boxes remain in
parallel parts of the diagram, so scheduling independence inserts no
causal edge between them. More precisely,
\cref{lem:egt-protocol-event-path} shows directly on the secured,
unhidden event-level interaction that a comparison of two decision
outputs would require a directed path of dependency wires. This proves
the stronger event-order assertion independently of the coherence
argument for completed values.
  \item \textbf{Causal execution.} A linear extension is
therefore only an external reading of the constructed arrow, not a scheduler
included in the strategy.
\end{enumerate}
\end{proof}


\begin{figure}[htbp]
\centering
\begin{tikzcd}[column sep=large,row sep=large]
\text{incoming data at }v \arrow[r,"o_v"] & O_v
 \arrow[r,"\pi_v"] & A_v \arrow[d,"\text{specified wires}"] \\
&& \text{later observations and outcome records}
\end{tikzcd}
\caption{A policy receives its constructed observation quotient. Wiring can
retain additional records without making them available to that policy.}
\end{figure}

The protocol construction now identifies complete economic policies
with families of operational response functions. Each function uses the
effective observation at its own node, and its realization is receptive
and courteous. When several nodes belong to one owner, the resulting
policy is their product; ownership neither duplicates the process nor
orders otherwise independent events. This provides a concrete
implementation of the agent representation.

\begin{remark}[Protocol coherence and incentives]
The linear-logical account of multiparty protocol compatibility in
\citet{CarboneEtAl2015,CarboneEtAl2016} concerns the operational layer.
In the present construction, the causal protocol specifies the agents' implemented interactions,
while preference comparisons and profitable deviations determine incentive
properties. The latter are constructed in the following sections.
\end{remark}

\subsection{Causal decision and session protocols}
\label{agt-sec-protocols}

\subsubsection{Decision generators and causal diagrams}

For finite $O,M$, the singly owned decision generator is
\[
 \Dec_i(O,M):D(O)\to D(M),\qquad
 \Sigma_i=M^O,\quad\Sigma_j=1\ (j\ne i),
\]
with policy $\pi$ realized by $\widehat\pi$. The response strategy has an
input occurrence for each $o$ and an output labelled $\pi(o)$ dependent
on that input. The finite-data construction proves
$\widehat{\id}=\id$ and $\widehat g\,\widehat f=\widehat{gf}$.
It supplies data copying and disposal on $D(O)$ and collectors
\[
 c_{(X_e)}:\bigotimes_eD(X_e)\to D\left(\prod_eX_e\right).
\]
A collector output requires all its inputs, without ordering their arrivals.

\begin{definition}[Finite causal decision syntax]
A finite causal decision diagram is a finite directed acyclic graph of
nodes $v$, owned by $i(v)$. Each incoming wire has a finite carrier $X_e$.
A node has an effective observation map
$o_v:\prod_{e\to v}X_e\twoheadrightarrow O_v$ and a finite output carrier
$M_v$. External ports and outcome-record functions are specified as typed
wires. Its complete owner policy family is
\[
 \Sigma_i=\prod_{v:i(v)=i}M_v^{O_v}.
\]
Terms are formed by composing and tensoring decision generators,
nonstrategic data wiring, and exhibited ownerwise parameter maps.
\end{definition}

\begin{construction}[The interpreted finite-data generator signature]
\label{agt-con-finite-generators}
For each finite $X$, take the representable linear object
$d_X=yD(X)$.  Section~\ref{sec:paper-1} constructs, for every finite
function $f:X\to Y$, the arrow $y\widehat f:d_X\to d_Y$, as well
as the copying, disposal, distributor, and collector arrows with their
concurrent realizations.  Via Theorem~\ref{agt-thm-representation}, the
family $\pi\mapsto y\widehat\pi$ determines a state
\[
 \delta_i:\I\longrightarrow\Ag_{\sig}(d_O,d_M),\qquad
 \Sigma_i=M^O,\quad\Sigma_j=1\ (j\ne i).
\]
These are named generators with \emph{exhibited interpretations}, not
arrows obtained from $\Lin$, tensor, and the other abstract LNL
constructors alone.  Their event-structure equations, including data
comonoid coherence and collector synchronization, were proved in Section~\ref{sec:paper-1} and are used below in that interpretation.  Equality
of concurrent interpretations is not automatically a definitional equality
between freely declared generator names.
\end{construction}

\begin{lemma}[A physically token-guarded finite response]
\label{agt-lem-token-guard}
For finite $O,M$ and $f:O\to M$, the composite
\[
 D(1)\otimes D(O)\xrightarrow{c_{1,O}}D(1\times O)
 \xrightarrow{\widehat f}D(M)
\]
(where $f$ is regarded as a map $1\times O\to M$) is receptive and
courteous. Every source configuration containing its positive output
labelled $f(o)$ also contains the authorization input event and the
input event labelled $o$ in the causal past of that output. Thus the
execution cannot deliver its result without receiving both inputs.
\end{lemma}
\begin{proof}
\leavevmode
\begin{enumerate}[label=\textup{(\arabic*)},leftmargin=*,itemsep=4pt,topsep=3pt]
  \item \textbf{Secured operational interaction.} In the finite collector the tuple output $(\ast,o)$ requires the
negative inputs on both components. The response strategy adds the
negative-to-positive dependency from this tuple input to its output
$f(o)$. Secured composition synchronizes the tuple events; hiding
them retains the induced dependencies from both original inputs to the
final output.
  \item \textbf{Secured operational interaction: completion.} Downward closure of configurations proves the claim. Receptivity and courtesy hold for both factors by the preceding constructions,
and are preserved under their secured composition.
\end{enumerate}
\end{proof}
\begin{remark}[Typing, issuance, and unforgeability]
A single linear variable of type $D(1)$ cannot be contracted by the core
rules; however the full concurrent category contains the state
$\mathsf{pt}_1(\ast):I\to D(1)$. Thus this arena alone is not an
unforgeable capability. An institutional issuance process is constructed
by composing an actual finite observation-based permission gate
(Proposition~\ref{agt-prop-test-implementation}) with the token input of
the above guarded response and retaining the issuance witness. Its output
is causally dependent on successful issuance. The issued event's causal predecessor and its execution witness record
which gate produced it.  Authenticity of issuers is a property of the
institution's exhibited certificate and permission interface, not of the
mere type $D(1)$: the independent state
$\mathsf{pt}_1(\ast):I\to D(1)$ is available in the full operational
category.  The bearer can therefore prove provenance only when the
provenance witness or its verified interface is retained.
\end{remark}

\begin{theorem}[Definability and scheduling independence]
\label{agt-thm-protocol}
Every finite causal decision diagram has an agent term over the interpreted
generator signature of Construction~\ref{agt-con-finite-generators}, with its
displayed complete policy family and the concurrent interpretation constructed
in the preceding constructions. Each input and profile determine a unique completed
assignment of outputs. Different linear readings have the same interpreted concurrent
interaction and add no order between independent decisions.  Equality of
these interpretations is not silently made a definitional equality between
unrelated generator terms in the freely presented syntax.
\end{theorem}

\begin{proof}
\leavevmode
\begin{enumerate}[label=\textup{(\arabic*)},leftmargin=*,itemsep=4pt,topsep=3pt]
  \item \textbf{Causal execution.} The generator signature maps each displayed $\Dec_i(O_v,M_v)$ to the
same singly owned decision presentation used in the constructive DAG
factorization in the proof of \cref{thm:egt-causal-protocol}. Its routing maps, nullary collectors, finite copying/disposal, and
ownerwise product coordinates are exactly those of that construction.
  \item \textbf{Causal execution: completion.} The causal-height recursion therefore gives the same unique output
assignment, while adjacent incomparable-node swaps and collector
coherence give the same concurrent interaction for every linear reading. This establishes the interpretation of the agent term, without identifying
freely declared primitive terms merely because their interpretations
happen to agree.
\end{enumerate}
\end{proof}


\begin{theorem}[Prefix realizations and non-anticipation]
\label{agt-thm-protocol-prefix}
For a downward-closed set $D\subseteq V$ of decision nodes, let $U_D$
be the product of external-input coordinates feeding nodes of $D$, and
$\Sigma_D=\prod_{v\in D}M_v^{O_v}$ (regrouped by owner when required).
The causal-height recursion constructs a unique prefix realization
\[
 r_D:U_D\times\Sigma_D\longrightarrow\prod_{v\in D}M_v.
\]
For downward-closed $D\subseteq E$, these maps obey
\begin{equation}
 \operatorname{pr}_D r_E
 =r_D(\operatorname{pr}_{U_D}\times\operatorname{pr}_{\Sigma_D}).
 \label{agt-eq-protocol-nonanticipation}
\end{equation}
Every node's output therefore depends only on external inputs and policy
coordinates at its causal predecessors and itself. For a node $v$, the
actual incoming tuple and declared effective observation determine
\[
 U_v\times\Sigma_{D_v^-}\xrightarrow{\operatorname{in}_v}
 H_v:=\prod_{e\to v}X_e\xrightarrow{o_v}O_v,
 \qquad
 U_v\times\Sigma_{D_v^-}\twoheadrightarrow
 O_v^{\mathrm{reach}}\lhook\joinrel\longrightarrow O_v,
\]
where the second factorization is the image of $o_v\operatorname{in}_v$,
$D_v^-$ is the strict causal past, and $U_v$ includes the external
ports feeding it or $v$. Execution evaluates the existing complete
policy $\pi_v:O_v\to M_v$ on that image; it does not replace the full
policy carrier by one containing only reachable observations.
\end{theorem}
\begin{proof}
\leavevmode
\begin{enumerate}[label=\textup{(\arabic*)},leftmargin=*,itemsep=4pt,topsep=3pt]
  \item \textbf{Fibrewise calculation.} At a minimal node of $D$, both sides of
\eqref{agt-eq-protocol-nonanticipation} evaluate the same external
inputs under its same policy coordinate. For a nonminimal $v\in D$,
each source of an incoming node wire is already in $D$, because $D$ is
downward closed. By induction on causal height all those values agree
between the full and prefix realizations.
  \item \textbf{Observation and knowledge.} Applying the same raw
observation $o_v$ and the same policy $\pi_v$ gives the same output. Extensionality of finite products gives the equation. The map
$\operatorname{in}_v$ is formed by projecting the required external
ports and the already recursively calculated predecessor outputs.
  \item \textbf{Monoidal comparison.} Image factorization of its composite with the declared $o_v$ gives
$O_v^{\mathrm{reach}}\hookrightarrow O_v$ directly. All maps are built
by recursion, products, and image factorization in $\Set$, so the
comparisons commute with pullback substitution of ordinary parameters.
\end{enumerate}
\end{proof}

\noindent A prefix response factors through the observation supplied by its causal predecessors.
\begin{center}
\begin{tikzcd}[column sep=large,row sep=large]
H_v\arrow[r,"o_v",two heads]\arrow[d,"f_v"']
 &O_v\arrow[dl,dashed,"\pi_v"]\\
 M_v&
\end{tikzcd}
\end{center}

\begin{example}[One agent informs another]
Compose $\Dec_i(O,M)$ with $\Dec_j(M,A)$. At policies
$\alpha:O\to M$ and $\beta:M\to A$, the first agent sends $\alpha(o)$
and the second responds with $\beta(\alpha(o))$. Their composite is
$\widehat\beta\,\widehat\alpha=\widehat{\beta\alpha}$.
Hiding the message boundary retains both owners' policy coordinates.
Recording or forwarding the message uses an explicit output wire.
\end{example}

\subsubsection{Finite session boundaries}

Finite session types provide a second view of these communications:
they prescribe how an endpoint can send, receive, and continue after a
message. The treatment of channel references in \citet{agt-Harper} supplies programming-language context; the present
construction makes the linear protocol discipline and its concurrent
realization explicit. We use the dependent finite grammar below.
\begin{equation}
 S::=\End\mid !_{m:M}S_m\mid ?_{m:M}S_m,
 \label{agt-eq-session-grammar}
\end{equation}
where $M$ is a finite ordinary message type and each continuation is already
a finite session. The nondependent notations $!M.S$ and $?M.S$ are special
cases. Internal choice $\oplus\{\ell:S_\ell\}$ is a send of a finite label;
external choice $\&\{\ell:S_\ell\}$ is its receive dual. The prefix $!$
here is not the exponential $\Lin\Val$. Duality is defined recursively by
\[
 \overline{\End}=\End,\quad
 \overline{!_{m:M}S_m}=?_{m:M}\overline{S_m},\quad
 \overline{?_{m:M}S_m}=!_{m:M}\overline{S_m}.
\]
Sending $m$ consumes the current endpoint of type
$!_{m:M}S_m$ and returns its residual endpoint $S_m$.
Its partner consumes the corresponding receive endpoint and continues
at $\overline{S_m}$. An ordinary message may be stored, but the old
endpoint does not remain available as a second linear resource.

\begin{construction}[Session arenas]
\label{agt-con-sessionarena}
Let $H_S$ be the finite tree of legal message histories of $S$, with empty
history at its root. A nonterminal history $h$ has an extension $hm$ for
each element of its current message carrier $M_h$. Events are nonempty
histories, ordered by prefix. A finite event set is consistent when its
elements lie on a single branch. The event $hm$ is positive when the
protocol at $h$ sends, and negative when it receives. Write
$\mathsf A_S$ for this polarized event structure.
\end{construction}

Predecessor sets are finite, consistency is downward closed, and adjoining
a predecessor preserves consistency. Thus the construction is an event
structure. Directly, $\mathsf A_{\overline S}=\mathsf A_S^\perp$ and
$\mathsf A_{\End}=I$. The continuation guard is an actual arena order.
In particular $\mathsf A_{!M.S}$ is not defined as the unguarded tensor
$D(M)\otimes\mathsf A_S$: that tensor alone would allow independent activity
in the two components.

\subsubsection{Complete policies of communicating agents}

Fix endpoint owners $a,b$.  Let $H(S)$ be the prefix-closed set of
session histories, including the empty history $\epsilon$.  A
\emph{decision position} is a nonterminal history, whereas an event of
$\mathsf A_S$ is a \emph{nonempty extension} of a history.  These are
different indexing sets.  Put
\[
 \mathcal D_a(S)=\{h\in H(S):S|h\ne\End\text{ and }S|h
                    \text{ begins with a send}\},
\]
and define $\mathcal D_b(S)$ by replacing send with receive; the send
moves of $a$ are exactly the receive moves of $b$.  At $h$ let $M_h$ be
the current finite message carrier.  Define the complete policy spaces
\begin{equation}
 \Sigma_i^S=\prod_{h\in\mathcal D_i(S)}M_h,
 \qquad i\in\{a,b\}.
 \label{agt-eq-session-policy}
\end{equation}
These products retain every legal decision position, including the
positions inconsistent with the choices made at other positions by the
same policy.  The empty product is the unit policy type.  For an
ordinary private value $u_i:O_i$, the corresponding policy carrier is
$\prod_{h\in\mathcal D_i(S)}M_h^{O_i}$, and its executed choice is
$\alpha_h(u_i)$.  When the private value is delivered physically, the
source strategy additionally records its negative input event, as
constructed below.

For $\alpha\in\Sigma_a^S$, recursively construct the prefix-closed
subtree $P_\alpha$ by retaining the uniquely selected outgoing event at
an $a$-send position and every outgoing event at a $b$-send position.
The inclusion $P_\alpha\hookrightarrow\mathsf A_S$ denotes a strategy
$p_\alpha:I\to\mathsf A_S$.  Dually, the policy $\beta$ of owner $b$
prunes its positive events in the opposite-polarity arena, giving
$q_\beta:\mathsf A_S\to I$.

\begin{lemma}[Endpoint realization]
\label{agt-lem-endpoint}
The two inclusions are receptive courteous strategies. Their complete
policy families define closed displayed agent terms
\[
 p:\Ag_{(\Sigma_a^S,1)}(yI,y\mathsf A_S),\qquad
 q:\Ag_{(1,\Sigma_b^S)}(y\mathsf A_S,yI)
\]
with owner list $(a,b)$.
\end{lemma}
\begin{proof}
\leavevmode
\begin{enumerate}[label=\textup{(\arabic*)},leftmargin=*,itemsep=4pt,topsep=3pt]
  \item \textbf{Causal execution.} Each pruned subtree is down-closed, locally injective, and preserves
single-branch consistency and polarity. Given an enabled negative
arena event after a retained source configuration, it belongs to the
next receive position and was retained, uniquely, by the recursive
construction. All causal edges of a pruned tree are inherited prefix
edges, so courtesy holds.
  \item \textbf{Packaging policy families as agent terms.} Applying
Theorem~\ref{agt-thm-representation} to the complete policy families
packages them as the displayed closed agent terms.
\end{enumerate}
\end{proof}

\subsubsection{Physically received private information}
\label{agt-sec-session-input}

\begin{construction}[Receptive endpoint with a private input]
\label{agt-con-private-endpoint}
For finite $O_a$, take the complete policy
$\alpha\in\prod_{h\in\mathcal D_a(S)}M_h^{O_a}$.  We construct
\[
 p_\alpha^{O_a}:D(O_a)\longrightarrow\mathsf A_S.
\]
Write $B_-$ for the set of nonempty session-event histories whose entire
prefix contains only negative moves of $\mathsf A_S$.  For each
$u\in O_a$, let $P_{\alpha,u}$ be the prefix-closed session subtree
obtained by applying $\alpha_h(u)$ at every positive decision position
and retaining all negative branches.  Its events in $B_-$ are shared
between all input values.

The source event set has: a negative event $e_u$ for each $u\in O_a$;
a common event $b_k$ for each $k\in B_-$; and an event $(u,k)$ for
each $k\in P_{\alpha,u}\setminus B_-$.  The $e_u$ are mutually
conflicting.  Order common events by their session prefix relation.
Order the tagged events, and the common-to-tagged dependencies, by the
same prefix relation.  Every first tagged event (which is necessarily
positive) has the additional predecessor $e_u$.  Distinct tagged values
are inconsistent.  A finite set is consistent exactly when its session
labels belong to one branch, it uses at most one input value and one
tag, and the input value, if present with a tag, equals that tag.
The source's finite-consistency family includes down-closures of all
these compatible finite sets.  Label $e_u$ in $D(O_a)^\perp$ and label
$b_k,(u,k)$ by the session event $k$ in $\mathsf A_S$.
The added prerequisite makes the private input causally prior to the
first positive session event without imposing an order between the
private input and an initial negative-only session prefix.

Apply the identical construction with exchanged session polarities for
owner $b$ to obtain
\[
 q_\beta^{O_b}:\mathsf A_S\otimes D(O_b)\longrightarrow I,
 \qquad\Sigma_b=\prod_{h\in\mathcal D_b(S)}M_h^{O_b}.
\]
The private-input event is negative in both arrow arenas.
\end{construction}

\begin{lemma}[Input causality, receptivity, and courtesy]
\label{agt-lem-private-endpoint}
These are receptive courteous strategies.  Session inputs available
before the endpoint's first send can precede the private-input event;
each input-dependent positive response has that event in its causal
past.  Every complete policy coordinate remains present in the policy
carrier, including off-path coordinates.
\end{lemma}

\begin{proof}
\leavevmode
\begin{enumerate}[label=\textup{(\arabic*)},leftmargin=*,itemsep=4pt,topsep=3pt]
  \item \textbf{Causal execution.} Principal causal downsets are finite because $S$ is finite and a tagged
suffix acquires only the one extra private predecessor. Cause extension
preserves the stated consistency: predecessors of a tagged event have
the same tag or are common prefix events, and its input predecessor is
$e_u$. Hence the source is an event structure.
  \item \textbf{Causal execution: second comparison.} Its labelling preserves
configurations and is locally injective: conflicting input values or
tags cannot occur in one configuration and the common and tagged event
sets have disjoint session labels.
  \item \textbf{Lifting through the comparison.} For receptivity, an enabled negative private-input move has a unique
$e_u$ lift. It is available after any common negative-only prefix,
and no local positive move can have occurred without some $e_u$. Before the first local positive move, every enabled negative session
move has its unique common-event lift, whether the private input is
already present or not.
  \item \textbf{Lifting along the surjection: second comparison.} After the first local positive move its value
tag $u$ is fixed, and the corresponding negative continuation has a
unique tagged lift. These cases exhaust negative arena moves.
  \item \textbf{Secured operational interaction.} An immediate source edge inherited from the session prefix tree is an
arena dependency. The only newly inserted immediate dependencies go
from the negative input $e_u$ to a first tagged positive event. Thus
courtesy holds.
  \item \textbf{Communication and delivery.} The dual endpoint is proved by the same argument with
the session polarities reversed. Neither construction restricts the
original product of contingent response maps.
\end{enumerate}
\end{proof}


\subsubsection{Connection, reduction, and fidelity}

Connection is linear composition. In provider--client notation its rule is
\begin{equation}
 \frac{\Gamma;\Delta\vdash P::c:S\qquad
       \Gamma;\Theta,c:S\vdash Q::z:T}
      {\Gamma;\Delta,\Theta\vdash(\nu c)(P\mid Q)::z:T}.
 \label{agt-eq-session-cut}
\end{equation}
Here the session boundary is interpreted by
$\sem S_\ell=y\mathsf A_S$ and its dual by
$\sem{\overline S}_\ell=y\mathsf A_S^\perp$.
The first premise denotes a linear arrow
$\sem\Delta\to y\mathsf A_S$; the second denotes
$\sem\Theta\otimes y\mathsf A_S\to\sem T_\ell$.
The latter occurrence is an \emph{input} linear variable: it is already
polarity-reversed when interpreted in the arrow arena.  Thus the
endpoint duality $\mathsf A_{\overline S}=\mathsf A_S^\perp$
is enforced by the input/output placement, and one does not replace
that input variable by a variable of the distinct linear object
$y\mathsf A_S^\perp$ in this derived cut rule.
The conclusion is their tensor-and-composition composite in $\cV$.
This gives a derived cut judgment from the core linear rules, not an
unrestricted process language. The finite operational language consists
only of the constructed endpoint prefix trees, complete-policy registers,
and a current history. Its $\Send$ and $\Recv$ notations name the
corresponding selected positive branch and receptive negative branch of
these trees, not independent freely asserted primitive processes.
Matching prefixes reduce by
\begin{equation}
 (\nu c)(\Send_c(m);P\mid\Recv_c(x);Q)
 \longrightarrow(\nu c)(P\mid Q[m/x]).
 \label{agt-eq-send-reduction}
\end{equation}
The connection theorem is therefore an operational statement, not
just a typing convention: dual session processes synchronize their
matching communication events, and hiding produces the completed
conversation. The residual theorem checks that each consumed endpoint
is replaced by exactly the continuation required by the received
message.

\subsubsection{Operational configurations and residual session endpoints}
\label{agt-sec-session-residual}

A runtime configuration is $(S,h,\alpha,\beta)$, where $h$ is a legal
history compatible with the complete policies and $S|h$ is its residual
session. At a nonterminal history the next sender is uniquely determined
by the polarity.  Its register gives the next label $m_h$, and the
transition is
\begin{equation}
 (S,h,\alpha,\beta)\longrightarrow
 (S,hm_h,\alpha,\beta).
 \label{agt-eq-runtime-transition}
\end{equation}
There is no transition from $\End$.  At a send node of $a$ one has
$m_h=\alpha_h$ and at a send node of $b$ one has $m_h=\beta_h$.
For physically received private values $u_a,u_b$, evaluate the corresponding
register entries as $\alpha_h(u_a)$ or $\beta_h(u_b)$; the causal input
events are prerequisites of the first relevant sends.  Local prefix
notation projects this communication transition to the matching sender
and receiver reduction~\eqref{agt-eq-send-reduction}.  Neither transition
alters the complete register.

\begin{lemma}[Residual session arena]
\label{agt-lem-residual-arena}
Let $h$ be a history of $S$, viewed as a configuration of $\mathsf A_S$.
The residual arena consisting of events compatible with and strictly
after $h$, with the events of $h$ removed and the inherited causal order,
is canonically isomorphic to $\mathsf A_{S|h}$.  Residual formation
commutes with duality.
\end{lemma}
\begin{proof}
\leavevmode
\begin{enumerate}[label=\textup{(\arabic*)},leftmargin=*,itemsep=4pt,topsep=3pt]
  \item \textbf{Causal execution.} Every remaining event is a nonempty extension $hk$ of $h$, and
$hk\mapsto k$ is a bijection to the nonempty histories of the syntactic
continuation $S|h$. Prefix order, single-branch consistency, and the
send/receive polarity at a continuation node are unchanged by deleting
$h$. Prepending $h$ is the inverse.
  \item \textbf{Causal execution: completion.} Reversing polarities before or after
this bijection gives the same event labels, proving the duality equation.
\end{enumerate}
\end{proof}

\begin{construction}[Restriction of complete session policies]
\label{agt-con-session-policy-residual}
For a legal session history $h$, concatenation gives a bijection
\[
 \mathcal D_i(S|h)\ \xrightarrow{\ k\mapsto hk\ }
 \{hk\in\mathcal D_i(S):k\in H(S|h)\}.
\]
In contrast to the residual \emph{event} bijection, the residual
\emph{decision} positions include $k=\epsilon$ when $S|h$ begins
with a decision by $i$.  Define the ordinary product projection
\[
 \boxed{\rho_{i,h}:\Sigma_i^S\longrightarrow\Sigma_i^{S|h},
 \qquad\rho_{i,h}(\alpha)_k=\alpha_{hk}.}
\]
This supplies the coordinates needed by the residual endpoint; the
original complete policy register remains unchanged and still contains
its earlier and counterfactual coordinates.  For a further legal history
$k$ of $S|h$, product extensionality gives
\begin{equation}
 \rho_{i,\epsilon}=\id,\qquad
 \rho_{i,hk}=\rho_{i,k}^{S|h}\rho_{i,h}.
 \label{agt-eq-session-policy-residual}
\end{equation}
Indeed both sides have component $\alpha_{hkl}$ at a residual decision
position $l$ of $S|hk$.
\end{construction}

\begin{construction}[Residual status of a physical private input]
\label{agt-con-session-input-status}
For owner $i$ define the ordinary local input-status carrier
$\mathsf{InputStatus}_i=1+O_i$.  Its left summand means the physical
input is still outstanding; a right-summand value $u$ means its negative
input occurrence belongs to the completed source history.  Let $x$ be
a configuration of the strategy of
Construction~\ref{agt-con-private-endpoint} and let $h$ be its session
projection, so $h$ is a prefix of $S$.

When the input is outstanding, no local positive session event has
occurred, because every first such event requires the missing private
input.  Residualizing at $x$ leaves the private-input strategy associated
with $\rho_{i,h}(\alpha)$ and the external input interface $D(O_i)$.
When $e_u\in x$, the input has already occurred.  Define the fixed-value
residual session policy by
\[
 \alpha^{u,h}_k=(\rho_{i,h}\alpha)_k(u)=\alpha_{hk}(u).
\]
Residualizing at $x$ and removing its completed private-input occurrence
leaves the fixed-value endpoint strategy over $\mathsf A_{S|h}$, with no
second private-input port.  For owner $b$ use the identical construction
in the dual session arena.
\end{construction}

\begin{lemma}[Coherence of input-status residuals]
\label{agt-lem-session-input-residual}
The two residual descriptions of
Construction~\ref{agt-con-session-input-status} are isomorphisms of
strategies over their respective residual session arenas.  Receiving a
private input changes the status from outstanding to the delivered value
$u$; receiving a session message leaves the status unchanged; a local
send is possible only in the received-value status.  These residual
isomorphisms commute with successive residual formation and session
duality.
\end{lemma}

\begin{proof}
\leavevmode
\begin{enumerate}[label=\textup{(\arabic*)},leftmargin=*,itemsep=4pt,topsep=3pt]
  \item \textbf{Causal execution.} An outstanding-input configuration has only common negative-prefix
session events. Deleting that prefix identifies the remaining common
negative events and tagged suffixes, by $hk\mapsto k$, with the source
constructed from $\rho_{i,h}\alpha$; its private input events $e_u$
remain present. If the input $e_u$ has arrived, delete it and the
completed session prefix.
  \item \textbf{Causal execution: second comparison.} The remaining tagged events $(u,hk)$ are
identified with the fixed-$u$ pruning of $S|h$, while any common
negative-prefix suffix has the identical image under this bijection. The identifications preserve order, consistency, polarity, and labelling,
and the additional private-input cause has already been discharged.
  \item \textbf{Causal execution.} A transition receiving $e_u$ selects the second description without
changing $h$. A negative session transition extends $h$ through the
same common or tagged branch and leaves the input status unchanged. Every positive session event is tagged and therefore can occur only in
the received-value case.
  \item \textbf{Policy-level comparison.} For successive session prefixes the resulting
policy components are
$\alpha_{hkl}(u)$ by
Equation~\eqref{agt-eq-session-policy-residual}; the input is removed
exactly once. Thus both orders of residualization yield the same
source event structure over the suffix arena. Reversal of session
polarities exchanges the two endpoint arguments without changing the
input-status construction.
\end{enumerate}
\end{proof}


\begin{lemma}[Residual endpoint strategy isomorphisms]
\label{agt-lem-residual-endpoints}
For a history $h$ compatible with the complete policy profile, residual
formation of its two fixed-policy endpoint strategies satisfies
\[
 \operatorname{Res}_h(p_\alpha)\cong p_{\rho_{a,h}(\alpha)},
 \qquad
 \operatorname{Res}_h(q_\beta)\cong q_{\rho_{b,h}(\beta)},
\]
using the identification
$\mathsf A_S/h\cong\mathsf A_{S|h}$ of
Lemma~\ref{agt-lem-residual-arena}. These isomorphisms commute with
further residual formation and with polarity reversal.
\end{lemma}
\begin{proof}
\leavevmode
\begin{enumerate}[label=\textup{(\arabic*)},leftmargin=*,itemsep=4pt,topsep=3pt]
  \item \textbf{Causal execution.} A remaining session event is exactly an extension $hk$, $k\ne\varnothing$. The bijection $hk\mapsto k$ preserves and reflects its prefix order,
branchwise consistency, and polarity. At a positive endpoint decision position $k$ (including $k=\epsilon$
when the residual begins with a local send) both prunings retain the
same chosen output, since the residual policy has value
$\rho_{i,h}(\alpha)_k=\alpha_{hk}$.
  \item \textbf{Fibrewise calculation.} At a negative endpoint node both
retain every enabled input. The resulting restriction is therefore an
isomorphism of the two strategy event structures over the residual arena. The indexed equation~\eqref{agt-eq-session-policy-residual} proves
compatibility of successive residuals, and the event bijection is
unchanged by polarity reversal.
\end{enumerate}
\end{proof}
\begin{proposition}[Derived finite-prefix subject reduction]
\label{agt-prop-session-subject-reduction}
For the above finite protocol-prefix notation and its closed complete
policy registers, a matching send and receive at history $h$ preserves
the external linear typing judgment, changes both connected endpoint
types to the dual pair $S|hm$ and $\overline{S|hm}$, and preserves the
full registers. The reduction~\eqref{agt-eq-send-reduction} is interpreted
by residual formation of the secured interaction, followed by the
canonical residual endpoint isomorphisms.
\end{proposition}

\begin{proof}
\leavevmode
\begin{enumerate}[label=\textup{(\arabic*)},leftmargin=*,itemsep=4pt,topsep=3pt]
  \item \textbf{Fibrewise calculation.} At the nonterminal history $h$, the fixed sender entry selects
$m\in M_h$; the receptive endpoint uniquely accepts that same input. The session syntax selects the continuation $S|hm$, and polarity
reversal selects its dual. The preceding lemma identifies the resulting
endpoint strategies with the respective prunings of these continuations.
  \item \textbf{Functoriality and substitution.} Substitution of $m$ into a dependent continuation is ordinary
substitution, so it cannot duplicate or drop unrelated linear variables. The core linear substitution lemma validates the cut before and after
that step. The projections $\rho_{i,hm}$ supply the first
post-prefix decision as well as all later decisions.
  \item \textbf{Causal execution.} Its full register
is unchanged; only those projections determine its residual use. For endpoints with physical private inputs, the additional input-status
isomorphisms of Lemma~\ref{agt-lem-session-input-residual}
verify that an input is neither used before arrival nor requested a
second time.
  \item \textbf{Secured operational interaction.} Secured interaction of the
matching prefix and hiding give precisely the residual composed
strategies, by the interaction and hiding construction of Section~\ref{sec:paper-1}.
  \item \textbf{Session subject reduction.} This proves subject reduction for the specified finite
prefix language without claiming a process calculus with arbitrary
channel creation or mobility.
\end{enumerate}
\end{proof}


\begin{theorem}[Finite session preservation, progress, and adequacy]
\label{agt-thm-session}
For every $(\alpha,\beta)\in\Sigma_a^S\times\Sigma_b^S$:
\begin{enumerate}[label=(\roman*),nosep]
\item every communication step leaves dual residual endpoint types and
preserves the complete-policy profile;
\item every nonterminal runtime configuration has a unique next
communication, and execution reaches $\End$ in finitely many steps;
\item the configurations of the secured interaction of $p_\alpha$ and
$q_\beta$ are exactly the prefixes of this execution. Its unique maximal
configuration records the completed transcript.
\end{enumerate}
\end{theorem}

\begin{proof}
\leavevmode
\begin{enumerate}[label=\textup{(\arabic*)},leftmargin=*,itemsep=4pt,topsep=3pt]
  \item \textbf{Fibrewise calculation.} At the current history $h$, exactly one endpoint sends. Its complete-policy
coordinate selects $m\in M_h$. The other endpoint retains every enabled
input and therefore receives this event uniquely.
  \item \textbf{Causal execution.} The continuations
are indexed by the same $m$ and are dual by
Lemma~\ref{agt-lem-residual-arena}. Their typed reduction and the
unchanged complete register follow from
Proposition~\ref{agt-prop-session-subject-reduction}. The finite height
of the residual syntax tree strictly decreases.
  \item \textbf{Policy-level comparison.} A reached node with an
empty message carrier would
require a component of the given complete profile in the empty set, which
is impossible. Hence execution terminates specifically at $\End$.
  \item \textbf{Secured operational interaction.} A matching configuration of the endpoint maps identifies session events
bearing the same history label. The maps are locally injective, and each
endpoint's configurations lie on one root-to-leaf branch. At each
nonterminal node the fixed sender pruning retains precisely the event
selected by its complete policy, while receptivity gives its unique
matching occurrence on the receiver.
  \item \textbf{Fixed-point argument.} Induction on the prefix length
therefore identifies a matching with exactly one initial segment of
the calculated run; conversely, each calculated initial segment gives
such a matching. Both endpoint causal orders restrict to prefix order,
so their union has no oriented cycle and every such matching is secured.
  \item \textbf{Secured operational interaction.} The secured-matching configuration bijection of
Theorem~\ref{thm:egt-pullback} now identifies the interaction
configurations with these prefixes, \emph{before hiding}. The complete
run gives the unique maximal secured witness. Hiding can erase the
internal transcript events from the composite scalar; the completed
witness or the explicit recorder of Lemma~\ref{agt-lem-session-recorder}
retains them for subsequent evaluation.
\end{enumerate}
\end{proof}


\begin{corollary}[Input-indexed session adequacy]
\label{agt-cor-session-input}
For fixed complete input-indexed policies $\alpha,\beta$ and supplied
private values $u_a\in O_a,u_b\in O_b$, compose the strategies of
Construction~\ref{agt-con-private-endpoint} with the value states
$\mathsf{pt}_{O_a}(u_a)$ and $\mathsf{pt}_{O_b}(u_b)$.
The projection of their secured interaction to the shared session arena
has exactly the prefixes of the execution of
\eqref{agt-eq-runtime-transition}, with next labels
$\alpha_h(u_a)$ and $\beta_h(u_b)$.  Its maximal session transcript
is unique.  The physical input events may be concurrent with earlier
received session moves; no ordering between the two private inputs is
added.
\end{corollary}
\begin{proof}
\leavevmode
\begin{enumerate}[label=\textup{(\arabic*)},leftmargin=*,itemsep=4pt,topsep=3pt]
  \item \textbf{Fibrewise calculation.} At the first local send of each owner its private input is already in
its source causal past, by Lemma~\ref{agt-lem-private-endpoint}. Lemma~\ref{agt-lem-session-input-residual} then identifies the remaining
endpoint with the fixed-value pruning at each later history, without
reissuing the private input. Every negative session continuation remains
uniquely receivable.
  \item \textbf{Fixed-point argument.} Induction on the current session prefix therefore gives the same unique
next synchronized label as the closed case, evaluated at the supplied
input. After forgetting private-input occurrences, its session matchings
are exactly the calculated prefixes. The input dependencies cannot form
a cycle: their edges point into the relevant positive moves and never
back towards an input.
  \item \textbf{Securedness of synchronized messages.} Hence the matchings are secured. The two input
events have no dependency between them, and finite-height termination
is unchanged.
\end{enumerate}
\end{proof}

The theorem establishes adequacy at the level of witnesses. After hiding
the entire session, $q_\beta p_\alpha$ has boundary $I\to I$, and the
concurrent category has a single scalar. Its transcript must therefore be
retained as an execution witness, or delivered on an explicit output wire,
when it determines later observations, records, or evaluation.

\begin{lemma}[A physical transcript output]
\label{agt-lem-session-recorder}
The finite closed session has a recorded implementation with boundary
$I\to D(T_S)$, where $T_S$ is its finite set of terminating transcripts.
For the input-indexed session there is a recorded implementation with
boundary $D(O_a)\otimes D(O_b)\to D(T_S)$.  Its complete policy carriers
are canonically the original ones.  At each complete profile, its visible
output is the calculated transcript, and a completed secured matching
retains the internal communication history witnessing that output.
\end{lemma}

\begin{proof}
\leavevmode
\begin{enumerate}[label=\textup{(\arabic*)},leftmargin=*,itemsep=4pt,topsep=3pt]
  \item \textbf{Fibrewise calculation.} At each terminating history $h$, adjoin to the session arena a positive
record event $r_h$ whose immediate cause is the final session event of
$h$ (or no cause if $h=\epsilon$). Its full causal past is therefore
exactly the events of $h$. Record events at different leaves remain
inconsistent, and their labels distinguish the corresponding leaves.
  \item \textbf{Ownerwise alternatives.} Call the extended arena $\mathsf A_S^r$. These new final positive moves
have a prescribed singleton choice; they contribute unit factors rather
than new owner policies.
  \item \textbf{Policy-level comparison.} Extend the first endpoint by producing $r_h$ at every leaf consistent
with its policy. Extend the dual endpoint by accepting every enabled
negative occurrence $r_h$ and then producing a positive output $d_h$
labelled by $h\in T_S$. Its only new immediate dependency is the
negative-to-positive edge $r_h<d_h$.
  \item \textbf{Secured operational interaction.} The two source maps preserve
prefix order and conflict, are locally injective, and accept every
enabled negative move uniquely. Their new dependencies satisfy courtesy;
thus both extended endpoints are receptive courteous strategies.
  \item \textbf{Policy-level comparison.} For a fixed complete policy profile, the secured interaction of the
extended endpoints has precisely the session prefixes established by
Theorem~\ref{agt-thm-session}, together with a final matched $r_h$ and
visible output $d_h$ on the unique completed branch. Before hiding, the
completed matching explicitly records the internal session events and
the final record occurrence.
  \item \textbf{Secured operational interaction.} After hiding the $\mathsf A_S^r$
interface, the output event $d_h$ remains with the induced visible
causal order, but the matched session events themselves are no longer
events of the composite strategy. This distinction proves the claim
about the retained execution witness and the visible output separately.
  \item \textbf{Causal execution.} For physical private inputs use the endpoints of
Construction~\ref{agt-con-private-endpoint}, extending their session
components by the same terminal recorder. Their reception events
remain independent until the protocol requires them, and the extra
causal predecessors of input-dependent positive events persist in the
secured interaction.
  \item \textbf{Secured operational interaction.} For the input-indexed version, first copy each physical private input
by its designated data comonoid $\Delta_{O_a}$ or $\Delta_{O_b}$. Send one copy to the corresponding private-input endpoint and the
other two copies to $c_{O_a,O_b}:D(O_a)\otimes D(O_b)\to D(O_a\times O_b)$.
  \item \textbf{Communication and delivery.} After the session recorder supplies $D(T_S)$, form the final output by
\[
 D(O_a\times O_b)\otimes D(T_S)
 \xrightarrow{c_{O_a\times O_b,T_S}}
 D((O_a\times O_b)\times T_S)
 \xrightarrow{\widehat{\operatorname{pr}_{T_S}}}D(T_S).
\]
This makes the final visible record causally dependent on both physical
private inputs and the completed transcript, without ordering the two
private input occurrences.
  \item \textbf{Ownerwise alternatives.} The entire construction uses only the
existing copying, collector, and finite response arrows. Its fixed data wiring preserves the displayed owner policy factors. On composing
with the selected value states,
Lemma~\ref{agt-lem-session-input-residual} and
Corollary~\ref{agt-cor-session-input} identify the same unique
completed transcript.
  \item \textbf{Policy-level comparison.} No policy coordinate is created by the
recorder: its outputs are functions of the existing terminal history.
\end{enumerate}
\end{proof}


\subsubsection{Composing conversations}

Sequential continuation is tree grafting. Given a continuation $T_h$ at
each terminating leaf $h$ of $S$, its decision nodes are those of $S$ and
the tagged nodes of every $T_h$. The complete owner carriers satisfy
\[
 \Sigma_i^{S[T]}\cong
 \Sigma_i^S\times\prod_{h\in\operatorname{Leaves}(S)}\Sigma_i^{T_h}.
\]
Substituting a continuation at each completed leaf gives the
complete policy space of the composite conversation. The policy at an
off-path leaf remains part of that product, even when the chosen execution
never visits it. Consequently sequential composition retains the
counterfactual choices needed by later strategic evaluation.

This compositional account of finite sessions is the operational
basis for the communication witnesses used below: typing records which
messages can be exchanged, while secured histories record which
messages were actually delivered. Neither alone determines a truthful
announcement or the recipient's epistemic update.

This supplies request--reply, labelled choice, private delivery, and finite
multi-round conversations. The general logical correspondence between session
typing and cut reduction is developed in \citet{agt-CairesEtAl}; the finite
calculus here proves its operational results directly while retaining
ownership and off-path policies. Recursive protocols and unrestricted channel
mobility are further syntax, not operations silently added by the connection
rule.


\subsection{Epistemic action events from completed communication}
\label{sec:egt-implemented-action-events}

The completed secured matchings and the transcript-output wires from
this section now provide the event and observation data required by the
abstract epistemic action semantics of Section~\ref{sec:paper-9}.
The construction keeps distinct execution witnesses even when their
observable action-event labels coincide.

\begin{construction}[Action data from completed execution witnesses]
\label{agt-con-witnessed-action-data}
For a finite implemented protocol, construct $\Xi$ as the dependent
carrier of its completed secured matchings, retaining the source world,
profile, and completed labelled causal history.  Concretely, over input
and complete-policy space $W\times\Sigma$, it is the dependent sum
of the sets of completed secured matching configurations associated
with the specified protocol realizers.  The deterministic
protocol and session adequacy theorems just proved construct those fibres; for
general witnessed finite interactions they retain all completed
matchings together with their individual witnesses. Its source projection is
$s:\Xi\to W$.  Each recipient's specified actual input ports determine
a local-record set $V_i$ of finite labelled causal subhistories and a
restriction map
\[
 d_i:\Xi\longrightarrow V_i.
\]
Here the local record includes exactly the events delivered to $i$
and any designated retained local policy register, with their labels
and induced causal order; independent receptions are not assigned a total
schedule.  The finite protocol has finitely many event occurrences and
finite message carriers, so its collections of completed labelled local
subhistories form finite sets.  If an externally visible public
record is required, use the final visible output to construct its
record map $d_0:\Xi\to V_0$; otherwise take $V_0=1$ and its unique
record map.  Construct the tuple and its image factorization
\[
 \widetilde d=\langle d_0,(d_i)_{i\in N}\rangle:
 \Xi\xrightarrow{q} E\lhook\joinrel\xrightarrow{\iota}
 V_0\times\prod_{i\in N}V_i.
\]
Here $q$ is the surjective image map and $\iota$ is its inclusion.
For each owner define $b_i=\operatorname{pr}_{V_i}\iota:E\to V_i$,
so that $d_i=b_iq$.  Thus the event labels and
their recipient observations are \emph{constructed from delivery}, not
separately specified event equivalence relations.  Explicitly,
\[
 q(\xi)=q(\xi')
 \quad\Longleftrightarrow\quad
 d_0(\xi)=d_0(\xi')\text{ and }
 d_i(\xi)=d_i(\xi')\text{ for all }i.
\]
This is the coarsest event-label identification preserving the retained
public record and every actually delivered recipient record.

The enabling predicate is the image
\[
 U=\im\langle s,q\rangle\hookrightarrow W\times E,
 \qquad (w,e)\in U\iff
 \exists\xi\in\Xi\ (s\xi=w\ \wedge\ q\xi=e).
\]
The effective updated observation of owner $i$ is the image of
\[
 o_i^+:U\to O_i\times V_i,
 \qquad o_i^+(w,e)=(o_i(w),b_i(e)).
\]
Write $R_i=\kerpair(o_i)$ and construct
$S_i=\kerpair(b_i)$.  The witness carrier $\Xi$ remains displayed above
$U$ together with each witness's concurrent realization and any
institutional certificate.  The image is used for the epistemic event
label, not to identify distinct execution witnesses.
\end{construction}


\begin{proof}
\leavevmode
\begin{enumerate}[label=\textup{(\arabic*)},leftmargin=*,itemsep=4pt,topsep=3pt]
  \item \textbf{Secured operational interaction.} The local record restrictions are functions because every completed
secured history has a determined finite labelled subhistory at each
designated recipient port. The image construction has the universal
property that every map out of $\Xi$ constant on $\kerpair(q)$
descends uniquely through the surjection $q$. The coordinate maps
$b_i$ exist by factorization through the image and satisfy
$b_iq=d_i$ by definition.
  \item \textbf{Monoidal comparison.} The product-update kernel is calculated on the enabled carrier:
\begin{align*}
 (w,e)\mathrel{\kerpair(o_i^+)}(v,f)
 &\Longleftrightarrow o_i(w)=o_i(v)\ \wedge\ b_i(e)=b_i(f)\\
 &\Longleftrightarrow wR_iv\ \wedge\ eS_if.
\end{align*}
Corestriction to the effective image does not change this kernel pair. The first line is equality in a Cartesian product, and the second is
the definition of the two component kernels.
  \item \textbf{Witness and legality comparison.} This is exactly
Equation~\eqref{agt-eq-actionrelation}. Nothing in the factorization
forgets the separate elements of $\Xi$: each remains an exhibited
witness in the institution's dependent execution carrier.
\end{enumerate}
\end{proof}


\section{Strategic Communication and the Creation of Knowledge}\label{sec:paper-11}

The previous section constructed causal protocols and their
completed transcripts. We can now use actual delivery records to
construct epistemic action events: taking the image of the completed
execution map gives the event carrier, and each recipient's delivered
projection determines its observation relation. To compare this update
with strategic choice, however, we must retain complete announcement
policies rather than only the messages sent at the current world.
The following split-surjection and alternative-lifting results identify
exactly which strategic and epistemic assertions survive evaluation,
including the information gained by coalition disclosure.

\subsection{Public inspection and the disclosure of pooled information}
\label{sec:egt-public-inspection}

The image in Construction~\ref{con:egt-pooled-observation} is a joint
signal. Disclosing its value is an epistemic action that can be described
simultaneously at every reference state. For observation relations
$R_i=\ker(o_i)$ on $W$ and a coalition $C$, write
\[
 R_C=\bigcap_{i\in C}R_i,\qquad
 \mathsf{DK}_C(P)=[R_C]P,\qquad
 \mathsf{CK}_C(P)=\left[\left(\bigcup_{i\in C}R_i\right)^*\right]P.
\]
The first predicate is distributed knowledge. The second is common
knowledge within $C$; for $C=N$ it is the previously defined
$\mathsf{CK}$. Empty intersection is $W^2$, while the reflexive
transitive closure of the empty union is equality. These conventions
also define both operations for the empty coalition.

\begin{construction}[Public inspection of a signal]
\label{con:egt-public-inspection}
For $q:W\to O$, let $E=\operatorname{im}(q)$, enable precisely the pair
$(w,q(w))$, and equip every agent's action set with equality. The
updated carrier is canonically $W$ and
\[
 R_i^q=R_i\cap\ker(q).
\]
If $R_i=\ker(o_i)$, the updated observation is
$\langle o_i,q\rangle$. At a designated world $w$, the component
$q^{-1}(q(w))$ is the announcement of the observed value. This fibre is closed under all
updated accessibility relations. Its inclusion is bounded, so
Theorem~\ref{thm:egt-modal-adequacy} identifies its modal interpretation
with that of the full-carrier update at $w$.
Retaining every possible outcome at once gives the full-carrier
presentation of public inspection
\citep{VanEijckGattingerWang2017KnowingValues}.
\end{construction}

\begin{proposition}[Inspection is an idempotent coreflection]
\label{prop:egt-inspection-coreflection}
In the poset of relation families on a fixed carrier $W$, ordered by
componentwise inclusion, let $\mathcal R_q$ be the subposet with
$T_i\subseteq\ker(q)$ for every $i$. Its inclusion $\iota$ has
right adjoint
\[
 \mathsf{Inspect}_q(R)_i=R_i\cap\ker(q),
 \qquad\iota\dashv\mathsf{Inspect}_q.
\]
Write $\mathsf I_q=\iota\mathsf{Inspect}_q$ for the associated
idempotent comonad on the full relation poset. For two fixed signals
$q,r$ on $W$, these endofunctors satisfy
\[
 \mathsf I_r\mathsf I_q
 =\mathsf I_{\langle q,r\rangle}
 =\mathsf I_q\mathsf I_r.
\]
These statements also hold in the poset of S5 relation families.
\end{proposition}
\begin{proof}
\leavevmode
\begin{enumerate}[label=\textup{(\arabic*)},leftmargin=*,itemsep=4pt,topsep=3pt]
  \item \textbf{Adjunction and transposition.} For $T\in\mathcal R_q$, $T\subseteq R$ if and only if
$T\subseteq R\cap\ker(q)$. This is the adjunction. Its associated endofunctor $\mathsf I_q$ has counit
$R\cap\ker(q)\subseteq R$ and is idempotent.
  \item \textbf{Predicate calculation.} The composition law follows from
$\ker\langle q,r\rangle=\ker(q)\cap\ker(r)$. Meets preserve
equivalence relations, proving the S5 version.
\end{enumerate}
\end{proof}

\begin{theorem}[Public pooling turns distributed knowledge into common knowledge]
\label{thm:egt-public-pooling}
Let $R_i=\ker(o_i)$ and let $C\subseteq N$ be nonempty. Form the
image factorization of the tuple of coalition observations,
\[
 W\xrightarrow{q_C}O_C\lhook\joinrel\longrightarrow
 \prod_{i\in C}O_i,
 \qquad R_C=\ker(q_C)=\bigcap_{i\in C}R_i.
\]
After public inspection of $q_C$, the common-knowledge relation of all
owners is exactly $R_C$. Consequently, for every fixed factual predicate
$A\subseteq W$,
\[
 \mathsf{CK}_N^{\,q_C}(A)=\mathsf{DK}_C(A),\qquad
 \mathsf{DK}_C(A):=[R_C]A.
\]
\end{theorem}
\begin{proof}
\leavevmode
\begin{enumerate}[label=\textup{(\arabic*)},leftmargin=*,itemsep=4pt,topsep=3pt]
  \item \textbf{Predicate calculation.} Each updated relation is $R_i^+=R_i\cap R_C$. If $i\in C$ this
is $R_C$; for other $i$ it is contained in $R_C$. Since $C$ is nonempty,
\[
 \bigcup_{i\in N}R_i^+=R_C.
\]
The relation $R_C$ is an equivalence relation, hence equals its own
reflexive transitive closure. Taking relational boxes proves the equation.
\end{enumerate}
\end{proof}

\begin{corollary}[Private pooling within a coalition]
\label{cor:egt-private-pooling}
Use the same event set and enabling relation, but take $S_i$ to be
equality for $i\in C$ and the universal relation for $i\notin C$.
Then members have updated relation $R_C$ and outsiders retain $R_i$.
For fixed $A$, common knowledge within $C$ after this update equals
$\mathsf{DK}_C(A)$ before it.
\end{corollary}
\begin{proof}
\leavevmode
\begin{enumerate}[label=\textup{(\arabic*)},leftmargin=*,itemsep=4pt,topsep=3pt]
  \item \textbf{Ownerwise alternatives.} Product update intersects member relations with $\ker(q_C)$ and leaves
outsider relations unchanged. The union of the member relations is $R_C$.
\end{enumerate}
\end{proof}

This private construction is the resolution of distributed knowledge studied
by \citet{AgotnesWang2016ResolvingDistributedKnowledge}.
Public inspection additionally
reveals the pooled value to outsiders. In both equations $A$ is retained as a
fixed subobject. An epistemic formula re-evaluated after the update may
denote a different subobject; Remark~\ref{rem:egt-announcement-reevaluation} makes this explicit.

\begin{proposition}[A connectivity obstruction for general signals]
\label{prop:egt-inspection-connectivity}
Let $E=(\bigcup_iR_i)^*$ and
$E^q=(\bigcup_i(R_i\cap\ker(q)))^*$. Then
\[
 E^q\subseteq E\cap\ker(q).
\]
For the quotient $c:W\to W/E$, there is a canonical surjection
\[
 W/E^q\longrightarrow\operatorname{im}\langle c,q\rangle.
\]
It is an isomorphism exactly when every nonempty intersection of an
$E$-class with a $q$-fibre is connected by the original $R_i$-edges
that stay inside that intersection.
\end{proposition}
\begin{proof}
\leavevmode
\begin{enumerate}[label=\textup{(\arabic*)},leftmargin=*,itemsep=4pt,topsep=3pt]
  \item \textbf{Lifting along the surjection.} Every generating updated edge preserves both the $E$-class and $q$. Thus both are constant on $E^q$-classes. This gives the factor map,
which is surjective because its target is the image.
  \item \textbf{Injectivity of the event quotient.} It is injective
exactly when equality of those two retained values implies $E^q$
relatedness. The definition of reflexive transitive closure identifies
this with the stated connectivity condition.
\end{enumerate}
\end{proof}

\begin{example}
On $W=\{a,b,c\}$ let $R_1$ have classes $\{a,b\},\{c\}$ and
$R_2$ have classes $\{a\},\{b,c\}$. Put $q(a)=q(c)=0$ and
$q(b)=1$. Then $E=W^2$, but every updated individual relation is
equality. Hence $E^q$ is equality, whereas $E\cap\ker(q)$ still
relates $a$ to $c$. Common knowledge cannot generally be updated by
intersecting the old common-knowledge relation with the signal kernel.
\end{example}

\begin{construction}[Finite operational disclosure]
\label{con:egt-finite-disclosure}
For the observation maps $o_i:W\to O_i$, take their image
factorizations $W\twoheadrightarrow F_i\hookrightarrow O_i$.
When these images are finite and $C$ is nonempty, the collector and
iterated data comultiplication form the strategy
\[
 \bigotimes_{i\in C}D(F_i)
 \xrightarrow{c_C}D\!\left(\prod_{i\in C}F_i\right)
 \xrightarrow{\Delta^{(N)}}
 \bigotimes_{j\in N}D\!\left(\prod_{i\in C}F_i\right).
\]
The collector law of Proposition~\ref{prop:egt-finite-collectors}
sends the supplied value states to their tuple state. The comonoid law
of Proposition~\ref{prop:egt-data-comonoid} places that tuple at every
recipient's output. Equality of tuples has kernel $R_C$, so this realizes
the disclosed signal of Theorem~\ref{thm:egt-public-pooling} on the
supplied observation values. Retaining each recipient's old observation
gives the updated map $\langle o_j,q_C\rangle$.
The collector waits for all coalition inputs; a decision that uses the
disclosure is wired after the corresponding output occurrence.
\end{construction}

The ordinary signal map specifies which information is disclosed. The
displayed strategy specifies how already supplied signal occurrences are
collected and delivered. This maintains the distinction between semantic
information and its causal availability at a decision interface.

\subsection{Agent-labelled truthful public actions}

\begin{construction}[The public announcement performed by an agent]
\label{con:egt-agent-authored-announcement}
Fix an observed agent presentation on $W$, with effective quotients
$\bar o_j:W\twoheadrightarrow\Omega_j$. For a sender $i$, take the action
set $E_i=\mathcal P(\Omega_i)$ and enable $B\in E_i$ at $w$ precisely when
$\bar o_i(w)\in B$. Every observer's action relation is equality on $E_i$.
The content and updated carrier are
\[
 A_{i,B}=\bar o_i^{-1}(B),\qquad
 \mathcal U_i=\{(w,B):w\in A_{i,B}\}.
\]
The sender identity is part of the action interface. Write
$\mathsf{say}_i(B)$ for the corresponding pointed public action. If
$a_{i,B}:A_{i,B}\hookrightarrow W$ and
$j_{i,B}:A_{i,B}\to\mathcal U_i$ are the inclusion and
$w\mapsto(w,B)$, its modalities are
\[
 [\mathsf{say}_i(B)]Q=\forall_{a_{i,B}}j_{i,B}^*Q,
 \qquad
 \langle\mathsf{say}_i(B)\rangle Q=\exists_{a_{i,B}}j_{i,B}^*Q.
\]
\end{construction}

\begin{proposition}[Authored announcement and public restriction]
\label{prop:egt-agent-announcement-restriction}
The content $A_{i,B}$ is invariant under the sender's observation relation,
so $K_iA_{i,B}=A_{i,B}$. The component of $\mathcal U_i$ at event $B$ is
isomorphic as an observed frame to $M|_{A_{i,B}}$. Consequently an
agent-labelled announcement has exactly the public-restriction semantics of
its content, with epistemic formulas evaluated in that restricted frame.
\end{proposition}
\begin{proof}
\leavevmode
\begin{enumerate}[label=\textup{(\arabic*)},leftmargin=*,itemsep=4pt,topsep=3pt]
  \item \textbf{Observation and knowledge.} If $\bar o_i(v)=\bar o_i(w)$, membership of that common value in $B$ is the
same at $v$ and $w$. This proves observation invariance and the knowledge
equation. The product-update relation at $B$ is
\[
 (w,B)R_j^+(v,B)
 \iff wR_jv\ \text{and}\ w,v\in A_{i,B},
\]
because the action relation is equality.
  \item \textbf{Isomorphism of restricted action models.} Thus $w\mapsto(w,B)$ is a frame
isomorphism from the restriction onto that component, preserving the
restricted factual terms. The displayed dynamic modalities are precisely
the inclusion quantifiers of the public restriction. Frame isomorphism and
the defining logical operations give the assertion for re-evaluated
epistemic formulas and their fixed points.
\end{enumerate}
\end{proof}

This realizes a sender-labelled public test for the
observation-invariant message family. Joint communication consists of
one such component per owner, but its strategic interpretation requires
complete policies specifying messages at every observation. We
construct those carriers next.

\begin{construction}[A public test as an institutional agent action]
\label{con:egt-institutional-public-test}
For any predicate $P\hookrightarrow W$, form its characteristic term
$\chi_P:W\to2$ in the paper's ordinary base and corestrict it to
$\Omega_P=\operatorname{im}(\chi_P)$. Extend the owner list by a fresh
institutional identity $\ast$, retaining all original world observations.
Give this agent observation $\chi_P:W\twoheadrightarrow\Omega_P$ and the
truthful message interface $\mathcal P(\Omega_P)$. The complete policy
$\eta_P(o)=\{o\}$ discloses the test result. Its successful public event is
$B_P=\Omega_P\cap\{1\}$; if $P$ is empty, that event has empty enabling
carrier. The original agents publicly receive the disclosed result.
The ordinary test and its concurrent realization have distinct types:
\[
 W\xrightarrow{\chi_P}2,
 \qquad
 D(\Omega_P)\xrightarrow{\widehat{\chi}}D(2)
 \xrightarrow{\text{success gate}}D(1),
\]
where $\chi:\Omega_P\to2$ is the inclusion of the image of $\chi_P$.
The first arrow is the \emph{specified ordinary signal input} and not
an event-structure sensor inferred from an arbitrary predicate on $W$.
The data strategy implements the second, finite-stage operation after
that input is actually supplied.
\end{construction}

\begin{proposition}[Public comprehension realized by a test agent]
\label{prop:egt-institutional-public-test}
The successful branch of this institutional action has enabling carrier
$P$, and its frame on the original agent indices is exactly $M|_P$.
Its dynamic modalities on retained predicates are the paper's public
restriction modalities. The disclosure component has a finite concurrent
realization, with its input supplied by the ordinary term $\chi_P$.
\end{proposition}

\begin{proof}
\leavevmode
\begin{enumerate}[label=\textup{(\arabic*)},leftmargin=*,itemsep=4pt,topsep=3pt]
  \item \textbf{Policy-level comparison.} The event $B_P$ is enabled precisely when
$\chi_P(w)\in\Omega_P\cap\{1\}$, equivalently $w\in P$. At such a
world the complete policy $\eta_P$ outputs exactly $B_P=\{1\}$; if no
such world exists both enabling carriers are empty. Equality of public
events restricts each original relation $R_j$ to $P\times P$.
  \item \textbf{Predicate calculation.} Projection
of the successful component to $P$ is therefore a frame isomorphism for
those indices and the retained factual terms. Its inclusion into $W$ is
the same comprehension map, proving both dynamic quantifier equations.
  \item \textbf{Observation and knowledge.} The observation image has at most two elements. Its singleton response and
public broadcast are the finite data constructions of the operational
semantics. The term $\chi_P$ supplies their ordinary reference input; the
implementation realizes the finite disclosure interface.
  \item \textbf{Observation and knowledge: completion.} Predicate
evaluation remains the ordinary semantic test that constructed $\chi_P$. The original agents' pre-action observations are retained, and the extra
institutional observation is explicitly part of this extended presentation.
\end{enumerate}
\end{proof}


The implementation records who sends the test, which result is
delivered, and which event succeeds. Thus the logical update has an
observable cause rather than an unexplained restriction of worlds.
The finite disclosure protocol acts on the test-result interface; the
underlying world remains an ordinary parameter.

\subsection{The full truthful message family and coalition announcement games}
\label{sec:egt-coalition-announcements}

A truthful public message is controlled by its sender, but its
content is constrained by that sender's observation. Power objects and
comprehension generate the entire semantic family of such messages.
The coalition effectivity quantifiers then test whether a group can
secure a post-announcement goal against truthful outsiders. This
relates strategic control of disclosure to the joint-signal operation
of Section~\ref{sec:egt-public-inspection}.

\begin{construction}[Truthful message families]
\label{con:egt-announcement-game}
Let $M=(W,(\ker(o_i))_{i\in N})$. Factor each observation as
$W\xrightarrow{\bar o_i}\Omega_i\hookrightarrow O_i$, where
$\Omega_i=\operatorname{im}(o_i)$. Define the ordinary message object
$\mathsf{Msg}_i=\mathcal P(\Omega_i)$ and its admissible fibre
\[
 \mathsf{Msg}_i(w)=\{B_i\subseteq\Omega_i:\bar o_i(w)\in B_i\}.
\]
The message $B_i$ announces the predicate
$A_i=\bar o_i^{-1}(B_i)$. This predicate is observation-invariant;
when $B_i$ is admissible at $w$, $K_iA_i(w)$ holds. Conversely every
observation-invariant predicate is the inverse image of a unique subset
of $\Omega_i$, by Proposition~\ref{prop:egt-predicate-descent}.

For $\mathsf{Msg}=\prod_i\mathsf{Msg}_i$ and $B\in\mathsf{Msg}$, put
\[
 A_B=\bigcap_i\bar o_i^{-1}(B_i),\qquad M_B=M|_{A_B}.
\]
The updated carriers form the ordinary dependent family
$\pi:\mathcal U=\{(B,v):v\in A_B\}\to\mathsf{Msg}$.
Its $i$-relation relates $(B,v)$ and $(B',v')$ precisely when
$B=B'$ and $v\ker(o_i)v'$. Thus all epistemic operations in a fibre
are those of $M_B$. A goal is a specified subobject
$\Psi\hookrightarrow\mathcal U$.

The legal message family over the reference state is
\[
 L=\{(w,B):B_i\in\mathsf{Msg}_i(w)\text{ for all }i\}
 \longrightarrow W.
\]
For $C\subseteq N$, let
$L_C=\{(w,B_C):B_i\in\mathsf{Msg}_i(w)\ (i\in C)\}$,
and write $e_C:L\to L_C$ and $u_C:L_C\to W$ for the projections.
These are the realizable coalition-coordinate objects in each fibre:
every partial tuple has the canonical completion $B_j=\Omega_j$ for
$j\notin C$. The map $g:L\to\mathcal U$, $(w,B)\mapsto(B,w)$,
evaluates the goal in the announced context. Define
\begin{equation}\label{eq:egt-announcement-ability}
 \mathsf{AnnAbility}_C(\Psi)
 =\exists_{u_C}\,\forall_{e_C}\,g^*\Psi
 \quad\in\mathsf{Pred}(W).
\end{equation}
In particular, at $w$ this asserts that some truthful coalition message
tuple succeeds against every simultaneous truthful outsider tuple.
\end{construction}

The construction uses power objects, products, projection, and
comprehension. Equation~\eqref{eq:egt-announcement-ability} is the
fibrewise instance of Equation~\eqref{eq:egt-legal-ability}; its
quantifiers follow the coalition-announcement pattern. Standard
announcement logics restrict messages to a specified epistemic syntax
\citep{AgotnesEtAl2010GroupAnnouncements,
GalimullinAlechina2018CoalitionGroupAnnouncement}; our full semantic
family instead contains all observation-invariant subobjects. Group
announcement logics with distributed knowledge are discussed by
\citet{AgotnesEtAl2022GroupDistributed}. The results below state their
own factual-goal and full-family scope rather than identifying these
different languages.

\begin{theorem}[Exact announcement ability for factual common knowledge]
\label{thm:egt-coalition-announcement}
For $P\subseteq W$, let $P^+=\{(B,v)\in\mathcal U:v\in P\}$
and take $\Psi_P=\mathsf{CK}_N^{\mathcal U}(P^+)$, using the
fibre-preserving relations of Construction~\ref{con:egt-announcement-game}.
For every nonempty coalition $C$,
\begin{equation}\label{eq:egt-announcement-distributed-knowledge}
 \mathsf{AnnAbility}_C(\Psi_P)=\mathsf{DK}_C(P).
\end{equation}
The sufficient coalition policy is the uniform observation-admissible
map $o\mapsto\{o\}$ for each member. For $C=\varnothing$,
$\mathsf{AnnAbility}_{\varnothing}(\Psi_P)=\mathsf{CK}_N^M(P)$.
\end{theorem}

\noindent A truthful announcement is represented as an informed section over the observation quotient.
\begin{center}
\begin{tikzcd}[column sep=large,row sep=large]
W\arrow[r,"o_C",two heads]\arrow[d,"m"']
 &O_C\arrow[d,"\pi"]\\
 M\arrow[r,equal]&M
\end{tikzcd}
\end{center}


\begin{proof}
\leavevmode
\begin{enumerate}[label=\textup{(\arabic*)},leftmargin=*,itemsep=4pt,topsep=3pt]
  \item \textbf{Fibrewise calculation.} At $w\in\mathsf{DK}_C(P)$, let every member announce
$B_i=\{\bar o_i(w)\}$. For every admissible outsider tuple,
\[
 w\in A_B\subseteq R_C[w]\subseteq P.
\]
The fact therefore holds throughout the updated carrier and is common
knowledge there. Each chosen message is a function of the owner's
observation alone, proving uniform observation-admissibility by
Proposition~\ref{prop:egt-observation}.
  \item \textbf{Reverse implication.} Conversely, let a coalition tuple establish the displayed common knowledge
against every outsider tuple. Take $B_j=\Omega_j$ for outsiders. Truthfulness and observation invariance give $R_C[w]\subseteq A_B$.
  \item \textbf{Observation and knowledge.} For any $i\in C$, each $v\in R_C[w]$ is $R_i$-accessible from $w$
inside $A_B$. Common knowledge in this updated model therefore gives
$v\in P$. This is $\mathsf{DK}_C(P)(w)$.
  \item \textbf{Ownerwise alternatives.} For the empty coalition, the outsider tuple of all tautological messages
leaves the original model unchanged, proving necessity.
  \item \textbf{Reverse implication.} Conversely every
path in a restriction containing $w$ is an original path. Hence common
knowledge of the fixed fact is preserved under all these restrictions. For $N=\varnothing$, the common-knowledge relation is equality and the
same argument gives the predicate $P$.
\end{enumerate}
\end{proof}


\begin{proposition}[Finite message realizations]
\label{prop:egt-finite-announcement-messages}
When all observation images $\Omega_i$ are finite, the full message
objects, singleton disclosure maps, and truthfulness tests in
Construction~\ref{con:egt-announcement-game} have finite data realizations.
\end{proposition}
\begin{proof}
\leavevmode
\begin{enumerate}[label=\textup{(\arabic*)},leftmargin=*,itemsep=4pt,topsep=3pt]
  \item \textbf{Message-policy and membership protocols.} Each $\mathcal P(\Omega_i)$ is finite. The map
$\Omega_i\to\mathcal P(\Omega_i)$, $o\mapsto\{o\}$, and the
membership function $\Omega_i\times\mathcal P(\Omega_i)\to2$
have response strategies by Proposition~\ref{prop:egt-functions}. Collectors form joint messages, and the data comonoids broadcast them.
  \item \textbf{Observation and knowledge.} The public enabling subobject is the conjunction of the membership tests,
so its finite presentation is
Construction~\ref{con:egt-finite-action-enablement} on these signal and
message carriers. Observational distinctions on an arbitrary original
state carrier remain in the ordinary dependent family over that carrier.
\end{enumerate}
\end{proof}

The construction specifies a full semantic message family. A chosen
communication language is represented by its own message object and
content map; its effectivity is then obtained by the same dependent
quantifiers. The equality in
Theorem~\ref{thm:egt-coalition-announcement} is established for the
constructed full family, whose singleton policy is explicitly available.

\subsection{Complete truthful policies of communicating agents}
\label{sec:egt-complete-announcement-policies}

Strategic announcement choice must specify a response at every
possible observation, not just at the current world. Contingent
announcements in \citet{AgotnesVanDitmarsch2011AnnouncementGames}
provide a precedent for this policy-based treatment. Here the
observation-invariant semantic message carrier is fixed by the preceding
construction. We compare its complete truthful policies with current
public messages while retaining private policy records and concurrent
execution witnesses. This comparison is not asserted for every
formula-restricted announcement language or for expected utility.

\begin{construction}[The complete truthful announcement interface]
\label{con:egt-complete-announcement-policies}
Use the effective observation images and message carriers of
Construction~\ref{con:egt-announcement-game}. Define
\[
 \mathsf{AnnPol}_i
 =\prod_{o\in\Omega_i}\{B\subseteq\Omega_i:o\in B\},
 \qquad
 \mathsf{AnnPol}=\prod_{i\in N}\mathsf{AnnPol}_i.
\]
A policy $\pi_i$ is a truthful message response at every possible
observation, including observations different from the current one.
It is an ordinary subobject of $\mathsf{Msg}_i^{\Omega_i}$. Its evaluation is
\[
 e_i:\Omega_i\times\mathsf{AnnPol}_i\to\mathsf{Msg}_i,
 \qquad e_i(o,\pi_i)=\pi_i(o).
\]
At $w$, the legal output of the agent is $\pi_i(\bar o_i(w))$.
The agent's policy carrier is $\mathsf{AnnPol}_i$, with its canonical
selection interface; a goal or comparison is connected to the resulting
message or updated state as in the preceding agent calculus.
\end{construction}

\begin{lemma}[Explicit extension of a truthful message]
\label{lem:egt-truthful-policy-extension}
\label{agt-lem-message-extension}
For every $o\in\Omega_i$, the evaluation map
\[
 e_{i,o}:\mathsf{AnnPol}_i\twoheadrightarrow
        \{B\subseteq\Omega_i:o\in B\}
\]
has the explicit section $B\mapsto E_{i,B}$, where
\[
 E_{i,B}(x)=B\cup\{x\}.
\]
In particular every legal message at $w$ is the current output of an
exhibited complete truthful policy.
\end{lemma}
\begin{proof}
\leavevmode
\begin{enumerate}[label=\textup{(\arabic*)},leftmargin=*,itemsep=4pt,topsep=3pt]
  \item \textbf{Split evaluation of public messages.} Every value of $E_{i,B}$ contains $x$, so the function belongs to
$\mathsf{AnnPol}_i$. Since $o\in B$, evaluation at $o$ is
$E_{i,B}(o)=B\cup\{o\}=B$. Hence $e_{i,o}(E_{i,B})=B$. The formula is defined
for every indicated pair and uses no choice of a representative.
\end{enumerate}
\end{proof}

\noindent Truthful messages extend to complete policies that evaluate to the chosen message at the given observation.
\begin{center}
\begin{tikzcd}[column sep=large,row sep=large]
\{B\subseteq\Omega_i:o\in B\}\arrow[r,"E_{i,-}"]\arrow[dr,"\id"']
 &\AnnPol_i\arrow[d,"e_{i,o}",two heads]\\
 &\{B\subseteq\Omega_i:o\in B\}
\end{tikzcd}
\end{center}

\begin{construction}[The policy-state model of public communication]
\label{con:egt-announcement-policy-model}
Put $\mathcal X=W\times\mathsf{AnnPol}$ and define the joint evaluation
\[
 b(w,\pi)=(\pi_i(\bar o_i(w)))_{i\in N}.
\]
The map to the paper's public-message carrier, retaining its original
coordinate order, is
\[
 F:\mathcal X\to\mathcal U,
 \qquad F(w,\pi)=(b(w,\pi),w).
\]
Give agent $j$ the pre-communication observation
$\widetilde o_j^0(w,\pi)=(\bar o_j(w),\pi_j)$. Use the public action set
$\mathsf{Msg}$ with equality as every action relation and enabling predicate
$B=b(w,\pi)$. Its updated carrier is canonically $\mathcal X$, since each
source has exactly one enabled joint message. Its post-communication
observation is
\[
 \widetilde o_j(w,\pi)
    =(\bar o_j(w),\pi_j,b(w,\pi)),
 \qquad R_j^{\mathcal X}=\ker(\widetilde o_j).
\]
Thus each agent retains its own complete policy, its original observation,
and the delivered public messages. This is an explicit product-update
instance of the earlier action-model construction.

The strategic relations hold the reference world fixed:
\[
 \begin{aligned}
 (w,\pi)D_C^{\mathcal X}(v,\pi')
   &\iff w=v\ \text{and}\ \pi_{-C}=\pi'_{-C},\\
 (B,w)D_C^{\mathcal U}(B',v)
   &\iff w=v\ \text{and}\ B_{-C}=B'_{-C}.
 \end{aligned}
\]
The latter are the legal message deviations in each fibre over $W$.
\end{construction}

\begin{theorem}[Exact policy-to-message comparison]
\label{thm:egt-announcement-policy-comparison}
\label{agt-thm-message-comparison}
The map $F:\mathcal X\to\mathcal U$ is a split surjection and a bounded
morphism for every individual epistemic relation, every nonempty
intersection of those relations, and every coalition-deviation relation.
More explicitly, for $\varnothing\ne J\subseteq N$ it is bounded from
$\bigcap_{j\in J}R_j^{\mathcal X}$ to
$\bigcap_{j\in J}R_j^{\mathcal U}$, and for every $C\subseteq N$ it is
bounded from $D_C^{\mathcal X}$ to $D_C^{\mathcal U}$.
\end{theorem}
\begin{proof}
\leavevmode
\begin{enumerate}[label=\textup{(\arabic*)},leftmargin=*,itemsep=4pt,topsep=3pt]
  \item \textbf{Ownerwise alternatives.} A section sends $(B,w)\in\mathcal U$ to
\[
 \bigl(w,(E_{i,B_i})_{i\in N}\bigr).
\]
It is well-defined by truthfulness of $B$ at $w$ and is a section by
Lemma~\ref{lem:egt-truthful-policy-extension}. This includes the empty owner
list, when both carriers are canonically $W$.
  \item \textbf{Observation and knowledge.} Forward preservation of the epistemic relations follows from their
observations: equality of $\widetilde o_j$ implies equality of the original
$j$-observation and of the public tuple. For the back condition, let
$x=(w,\pi)$ and suppose
\[
 F(x)\Bigl(\bigcap_{j\in J}R_j^{\mathcal U}\Bigr)(B,v),
 \qquad J\ne\varnothing.
\]
Then $B=b(w,\pi)$ and $\bar o_j(v)=\bar o_j(w)$ for every $j\in J$.
  \item \textbf{Construction on policy coordinates.} Define a profile at $v$ by
\[
 \pi'_j=\pi_j\quad(j\in J),\qquad
 \pi'_k=E_{k,B_k}\quad(k\notin J).
\]
For a retained coordinate,
$\pi_j(\bar o_j(v))=\pi_j(\bar o_j(w))=B_j$; for every other coordinate
the extension lemma gives the same output. Hence $F(v,\pi')=(B,v)$. For each $j\in J$ the old observation, its own policy, and the public tuple
are unchanged, so $xR_j^{\mathcal X}(v,\pi')$.
  \item \textbf{Lifting intersections along the report map.} This proves the simultaneous
back condition and, in particular, its individual case. For an empty
intersection the universal relation has the same lifting property by
surjectivity alone.
  \item \textbf{Ownerwise alternatives.} For strategic relations, retaining the outsider policies at fixed $w$
retains their evaluated messages, which proves the forward clause. Suppose $F(w,\pi)D_C^{\mathcal U}(B',w)$. Keep
$\pi'_j=\pi_j$ for $j\notin C$ and set
$\pi'_i=E_{i,B'_i}$ for $i\in C$.
  \item \textbf{Ownerwise alternatives: completion.} The outsider outputs are already $B'_{-C}$, and the changed outputs are
$B'_C$ by the extension lemma. Therefore
$F(w,\pi')=(B',w)$ and
$(w,\pi)D_C^{\mathcal X}(w,\pi')$. This proves the back clause for every
coalition, including the empty and grand coalitions.
\end{enumerate}
\end{proof}

\noindent The complete-policy comparison is a split surjection with an explicit section.
\begin{center}
\begin{tikzcd}[column sep=large,row sep=large]
\mathcal U\arrow[r,"s"]\arrow[dr,"\id"']
 &\mathcal X\arrow[d,"F",two heads]\\
 &\mathcal U
\end{tikzcd}
\end{center}

\begin{lemma}[Dependent successor-fibre exactness of truthful communication]
\label{lem:egt-announcement-successor-fibres}
For each nonempty observer set $J\subseteq N$, write
$K_J=\bigcap_{j\in J}R_j$ for the joint epistemic relation; for each
coalition $C\subseteq N$, retain $D_C$ for strategic deviations.
For either $R=K_J$ or $R=D_C$, with the corresponding superscripts
$\mathcal X$ and $\mathcal U$, the canonical map of dependent
successor families
\[
 \lambda_R:R^{\mathcal X}\longrightarrow
 \mathcal X\times_{\mathcal U,p_1}R^{\mathcal U},
 \qquad (x,x')\longmapsto(x,(Fx,Fx'))
\]
is surjective. Consequently for every $Q\hookrightarrow\mathcal U$,
\[
 F^*[R^{\mathcal U}]Q=[R^{\mathcal X}]F^*Q,
 \qquad
 F^*\langle R^{\mathcal U}\rangle Q
   =\langle R^{\mathcal X}\rangle F^*Q.
\]
The same statement for an empty epistemic intersection follows from
the split surjection $F$.
\end{lemma}
\begin{proof}
\leavevmode
\begin{enumerate}[label=\textup{(\arabic*)},leftmargin=*,itemsep=4pt,topsep=3pt]
  \item \textbf{Fibrewise calculation.} At a fixed $x\in\mathcal X$, a point of the displayed pullback is a
specified $R^{\mathcal U}$-successor of $Fx$. The epistemic and strategic
back-lifts constructed in
Theorem~\ref{thm:egt-announcement-policy-comparison} provide an
$R^{\mathcal X}$-successor of that very $x$ mapping to it. Thus each
fibre of $\lambda_R$ is inhabited, which is surjectivity in
$\mathbf{Set}$.
  \item \textbf{Predicate calculation.} Forward preservation ensures well-definedness. The source and target existential successor predicates therefore agree on
inverse-image tests. Their universal successor predicates agree as well:
a target test is true on every successor if and only if its pullback is
true on every source successor, because the successor map is onto.
  \item \textbf{Modal preservation for empty coalitions.} This proves both displayed modal equations. For an empty epistemic
intersection the relation is universal; the split surjection supplies
the requisite lift from any fixed source state.
\end{enumerate}
\end{proof}

\noindent Alternative lifting is performed over the fixed policy reference.
\begin{center}
\begin{tikzcd}[column sep=large,row sep=large]
D_C^{\mathcal X}\arrow[r,"\lambda_C",two heads]\arrow[d,"p_1"']
 &\mathcal X\times_{\mathcal U}D_C^{\mathcal U}\arrow[d]\\
 \mathcal X\arrow[r,equal]&\mathcal X
\end{tikzcd}
\end{center}

\begin{corollary}[Knowledge, strategic predicates, and dynamic transport]
\label{cor:egt-announcement-policy-logic}
For predicates and comparisons retained on $\mathcal U$, inverse image
along $F$ preserves the individual knowledge, distributed knowledge,
common-knowledge, coalition-deviation, and positive modal fixed-point
operations. It preserves and reflects validity. For a comparison
$Q_i\hookrightarrow\mathcal U^2$, use
$(F\times F)^*Q_i$ on $\mathcal X^2$. The Nash predicates defined using the
displayed deviation relations satisfy
\[
 \mathsf{NE}_{\mathcal X}=F^*\mathsf{NE}_{\mathcal U},
\]
and the corresponding strict coalitional stability predicates satisfy the
same equality. Public restrictions by retained assertions and action models
pulled back from $\mathcal U$ have the same exact dynamic comparison.
\end{corollary}

\begin{proof}
\leavevmode
\begin{enumerate}[label=\textup{(\arabic*)},leftmargin=*,itemsep=4pt,topsep=3pt]
  \item \textbf{Communication and delivery.} The explicit comparison
$\lambda_R:R^{\mathcal X}\twoheadrightarrow
 \mathcal X\times_{\mathcal U,p_1}R^{\mathcal U}$
of Lemma~\ref{lem:egt-announcement-successor-fibres} gives the two
inverse-image equations
\[
 F^*[R^{\mathcal U}]Q=[R^{\mathcal X}]F^*Q,
 \qquad
 F^*\langle R^{\mathcal U}\rangle Q
    =\langle R^{\mathcal X}\rangle F^*Q.
\]
Here the source reference is not chosen afresh: it is fixed in the
pullback, and every target successor lifts over that same reference.
  \item \textbf{Lifting through the comparison.} Theorem~\ref{thm:egt-modal-adequacy} proves the positive fixed-point
comparison. The intersection lifting in the theorem supplies the distributed
knowledge case directly. Common knowledge uses the same individual
relations and its previously constructed greatest fixed point.
  \item \textbf{Fibrewise calculation.} At a source $x$, every concrete $i$-deviation maps to a target deviation
with comparison $Q_i(Fx,Fx')$, and every target deviation lifts at that same
source. Universal absence of improving alternatives is therefore equivalent
on the two carriers. For a coalition use its lifting clause and the specified
collective comparison; conjunction over coalitions gives the claimed
stability equality.
  \item \textbf{Base-change comparison.} For $V\hookrightarrow\mathcal U$, the restriction
$F^{-1}(V)\to V$ is surjective. A lift of an alternative in $V$ belongs to
$F^{-1}(V)$ by its image, so it remains bounded. Comprehension and
Beck--Chevalley give the public-restriction comparison.
  \item \textbf{Lifting along the surjection.} For a pulled-back
action model apply Corollary~\ref{cor:egt-action-bounded-quotient}. Surjectivity reflects validity in each case.
\end{enumerate}
\end{proof}


The retained language has its predicates on $\mathcal U$, pulled back to
$\mathcal X$. Predicates concerning additional private policy records are
also available on $\mathcal X$ and remain attached to that presentation.
The comparison theorem relates the specified interfaces and their retained
tests; it keeps the whole policy family needed to lift alternatives.

\begin{example}[A public policy record carries additional information]
\label{ex:egt-public-policy-record}
Take two worlds $0,1$. Agent $1$ observes the world, while agent $2$ has one
observation. The truthful policy
$\pi_1(0)=\{0,1\}$, $\pi_1(1)=\{1\}$ sends the tautological message at
$0$. Its content alone leaves both worlds possible to agent $2$. If the
complete policy $\pi_1$ is also publicly delivered, observing that message
identifies world $0$. In the policy-state model of
Construction~\ref{con:egt-announcement-policy-model}, only agent $1$ retains
its own policy as a private record. At world $1$ the alternative policy
$E_{1,\{0,1\}}$ gives the same public message, supplying the lift required
for agent $2$. The example exhibits the informational effect of an additional
policy-record wire and explains the role of the full policy family in $F$.
\end{example}

\subsection{Concurrent implementation and retained policy witnesses}

\begin{proposition}[Concurrent realization of complete announcing agents]
\label{prop:egt-complete-announcement-realization}
\label{agt-prop-ann-implementation}
When the observation images $\Omega_i$ are finite, the complete truthful
policy presentation and its post-communication observations have an
implementation in $\mathbf{DecArena}_N(\mathcal C)$. Its complete owner
policy carriers are $\mathsf{AnnPol}_i$. At reference input
$(\bar o_i(w))_i$ and policy profile $\pi$, its delivered observation to
agent $j$ is
\[
 (\bar o_j(w),\pi_j,b(w,\pi)).
\]
The report decisions are concurrent, and the joint message is delivered
after its component messages are collected.
\end{proposition}

\begin{proof}
\leavevmode
\begin{enumerate}[label=\textup{(\arabic*)},leftmargin=*,itemsep=4pt,topsep=3pt]
  \item \textbf{Realization of announcement policies.} Each $\mathsf{Msg}_i=\mathcal P(\Omega_i)$ and $\mathsf{AnnPol}_i$ is finite. Restrict the decision generator
$\mathsf{Dec}_i(\Omega_i,\mathsf{Msg}_i)$ along the exhibited inclusion
\[
 \mathsf{AnnPol}_i\hookrightarrow\mathsf{Msg}_i^{\Omega_i}.
\]
It realizes policy $\pi_i$ by the response strategy $\widehat{\pi_i}$. Tensor these components and apply the finite collector to their message
outputs.
  \item \textbf{Broadcast of the joint report.} The resulting joint output is $b(w,\pi)$. Broadcast this output
using the iterated data comonoid, one copy for each agent.
  \item \textbf{Observation and knowledge.} Copy the original input observation to its reporting component and its
owner's retained-observation wire. For a private policy record, use the
singly owned nullary decision generator with response carrier
$\mathsf{AnnPol}_i$; at a fixed policy it realizes the constant
$\ulcorner\pi_i\urcorner:I\to D(\mathsf{AnnPol}_i)$, using the nullary
collector and the response from $1$.
  \item \textbf{Ownerwise alternatives.} Compose these two occurrences of owner
$i$ and reparametrize their product carrier by the diagonal assigning the
same $\pi_i$ to the reporting and record components. This is a permitted
ownerwise reparametrization in the generated calculus. It leaves exactly
$\mathsf{AnnPol}_i$ as the complete policy carrier.
  \item \textbf{Collection of observation outputs.} Finally collect, at each owner's output, its retained observation, private
policy record, and public tuple. The full boundaries can be taken as
\[
 A_{\mathrm{ann}}=\bigotimes_iD(\Omega_i),\qquad
 B_{\mathrm{ann}}=
   \bigotimes_jD(\Omega_j\times\mathsf{AnnPol}_j\times\mathsf{Msg}).
\]
Every component is a decision generator, nonstrategic data wiring, or an
exhibited parameter reindexing, so the process lies in
$\mathbf{DecArena}_N(\mathcal C)$.
  \item \textbf{Causal execution.} Evaluating each response and collector
gives the displayed record. The reporting components are tensored; the
message collector is their explicit delivery barrier. Coherence and
scheduling independence follow from Theorem~\ref{thm:egt-causal-protocol}.
  \item \textbf{Observation and knowledge.} The original carrier $W$ remains an ordinary reference index: its tuple of
observations supplies the finite input, so no finite realization of $W$ is
required.
\end{enumerate}
\end{proof}


\begin{construction}[Policy witnesses over a public message state]
\label{con:egt-announcement-policy-witnesses}
For $u\in\mathcal U$ retain the ordinary dependent family
\[
 K_u=F^{-1}(u).
\]
Its total carrier has the canonical bijection
\[
 \coprod_{u\in\mathcal U}K_u\cong\mathcal X,
 \qquad(u,x)\longmapsto x,
\]
whose inverse is $x\mapsto(F(x),x)$. In the finite realization, let
$g_{\mathrm{ann}}:\mathsf{AnnPol}\to
 \mathcal C(A_{\mathrm{ann}},B_{\mathrm{ann}})$ be the process just
constructed. Associate to $x=(w,\pi)\in K_u$ the realized arrow
$g_{\mathrm{ann}}(\pi)$. This is a witness realization of
Construction~\ref{con:egt-general-witness-family}; its dependent linear
transpose exists by the ordinary--linear adjunction. The family retains
every implementing policy and its private records, while $F$ gives the
exact logical comparison on retained message tests. Reconstructing its total
carrier recovers all of $\mathcal X$ and the attached realizations.
\end{construction}

\subsection{Coalition ability through complete communication policies}

\begin{theorem}[Complete-policy announcement ability]
\label{thm:egt-complete-policy-announcement-ability}
For every goal $\Psi\hookrightarrow\mathcal U$ and coalition $C\subseteq N$,
\[
 \mathsf{AnnAbility}_C(\Psi)(w)
 \iff
 \exists\pi_C\in\prod_{i\in C}\mathsf{AnnPol}_i\;
 \forall\pi_{-C}\in\prod_{j\notin C}\mathsf{AnnPol}_j\;
       \Psi\bigl(F(w,(\pi_C,\pi_{-C}))\bigr).
\]
This is the paper's pointwise effectivity predicate, expressed through
the complete policy interfaces of the announcing agents.
\end{theorem}
\begin{proof}
\leavevmode
\begin{enumerate}[label=\textup{(\arabic*)},leftmargin=*,itemsep=4pt,topsep=3pt]
  \item \textbf{Ownerwise alternatives.} If $B_C$ is a winning legal coalition message at $w$, choose the explicit
policies $E_{i,B_i}$ for its members. Every outsider policy
evaluates to a legal outsider message, so every resulting joint message
satisfies $\Psi$.
  \item \textbf{Reverse implication.} Conversely, let $\pi_C$ satisfy the right-hand side and let
$B_C=(\pi_i(\bar o_i(w)))_{i\in C}$. Given any legal outsider tuple
$B_{-C}$, extend it by
$\pi_j=E_{j,B_j}$ for $j\notin C$. The resulting policy tuple
evaluates to $(B_C,B_{-C})$, so that legal joint message satisfies $\Psi$.
  \item \textbf{Monoidal comparison.} This is exactly the fibrewise existential--universal quantifier defining
$\mathsf{AnnAbility}_C$. The formulas also cover empty products.
\end{enumerate}
\end{proof}

\begin{corollary}[Uniform disclosure by complete agent policies]
\label{cor:egt-uniform-agent-disclosure}
For $P\hookrightarrow W$, put
$\Psi_P=\mathsf{CK}_N^{\mathcal U}(P^+)$ as in
Theorem~\ref{thm:egt-coalition-announcement}. For every nonempty coalition
$C$, define the fixed complete policies
\[
 \eta_i(o)=\{o\}\quad(i\in C).
\]
Then, at every $w$,
\[
 \begin{aligned}
 \mathsf{DK}_C(P)(w)
 &\iff \forall\pi_{-C}\;
          \Psi_P(F(w,(\eta_C,\pi_{-C})))\\
 &\iff \exists\pi_C\;\forall\pi_{-C}\;
          \Psi_P(F(w,(\pi_C,\pi_{-C}))).
 \end{aligned}
\]
The same $\eta_C$ is used at every world. For $C=\varnothing$, complete-policy
announcement ability for this goal is $\mathsf{CK}_N^M(P)$.
\end{corollary}

\begin{proof}
\leavevmode
\begin{enumerate}[label=\textup{(\arabic*)},leftmargin=*,itemsep=4pt,topsep=3pt]
  \item \textbf{Observation and knowledge.} If $R_C[w]\subseteq P$, singleton coalition messages give
$w\in A_B\subseteq R_C[w]\subseteq P$ against every truthful outsider
output. Hence the fixed fact holds throughout the updated carrier and is
common knowledge.
  \item \textbf{Comparison of existential witnesses.} This proves the first implication with the exhibited
complete policies. Those policies are a particular existential witness,
giving the second implication.
  \item \textbf{Reverse implication.} Conversely, take any successful complete coalition policies and give each
outsider the constant tautological policy $o\mapsto\Omega_j$. At $w$ their
joint output $B$ satisfies $R_C[w]\subseteq A_B$ by truthfulness of the
coalition messages. For $i\in C$, every $v\in R_C[w]$ is $R_i$-accessible
from $w$ inside that updated carrier.
  \item \textbf{Observation and knowledge.} Common knowledge of $P^+$ gives
$v\in P$, so $R_C[w]\subseteq P$. This proves the converse.
  \item \textbf{Observation and knowledge: second comparison.} For the empty coalition, tautological outsider policies leave the original
model unchanged and establish necessity of original common knowledge.
  \item \textbf{Reverse implication.} Conversely every path in a public restriction containing $w$ is an original
path, so a fixed fact that was common knowledge remains common knowledge. If $N$ is empty, the common-knowledge relation is equality and the same
statement reduces to $P$.
\end{enumerate}
\end{proof}


The general ability predicate in
Theorem~\ref{thm:egt-complete-policy-announcement-ability} is evaluated
pointwise at $w$. For the factual common-knowledge goal, the corollary
constructs a stronger witness: one coalition policy independent of $w$
that succeeds at every world satisfying distributed knowledge.
The distinction is between a response chosen at the current reference
and a single complete policy valid throughout the specified information
class.

\subsection{Messages change observations}

To turn a delivered message into epistemic information, the state
carrier must retain enough evidence to identify what was sent.
Let $X$ contain the reference world, the complete policy profile, and
any execution witness determining $m:X\to M$. A recipient $j$
keeps its previous observation and receives the message through the
updated observation map below.
\[
 o_j^+(x)=(o_j(x),m(x)),\qquad
 R_j^+=R_j\cap\kerpair(m).
\]
For deterministic private delivery, each recipient's event
observation $b_i$ records the delivered value, whereas a nonrecipient's
$b_i$ is constant on event labels. Outsiders therefore keep their prior
relations. Public delivery makes the event relation equality for every
recipient. These are instances of
\eqref{agt-eq-actioncarrier}--\eqref{agt-eq-actionrelation}.
In finite sessions, the retained transcript of
Theorem~\ref{agt-thm-session} determines the message $m$, while its
recipient projection specifies exactly which part of the transcript
enters a particular agent's observation.

Session typing certifies the permitted order and shape of the
exchange, not the truth of the communicated proposition. Truthfulness
is expressed by the enabling predicate relating a message to the world.
Thus well-typed reception and informative disclosure remain distinct
properties of one implemented communication.

\subsection{Pooling and collective knowledge}

Let $q_C:W\twoheadrightarrow O_C$ be the effective joint signal of a
nonempty coalition. In the finite protocol of
Lemma~\ref{agt-lem-completed-delivery}, a public delivery uses $J=N$
and a private delivery to the coalition uses $J=C$. At the \emph{completed
execution boundary}, the state carrier is the graph of the deterministic
completed transcript, canonically isomorphic to $W$. A designated
recipient's new observation is the pair of its earlier observation and
the actually received tuple. Thus for public delivery
$R_i^+=R_i\cap\kerpair(q_C)$ for every owner; private pooling makes
that same change for members of $C$ and leaves outsiders' individual
relations unchanged. These statements concern the completed state
carrier, not an intermediate configuration where some recipients have
not yet received the value.

\begin{lemma}[Physical completed-delivery boundary]
\label{agt-lem-completed-delivery}
Let $C\ne\varnothing$ be finite with finite effective observation
carriers $\Omega_i$, and let $J$ be the designated recipient set. Put
$T_C=\prod_{i\in C}\Omega_i$.  The previously constructed true-concurrent
data strategies give a receptive courteous arrow
\[
 \bigotimes_{i\in C}D(\Omega_i)\longrightarrow D(T_C)
\]
with a final visible tuple record.  In its \emph{secured interaction
witness before hiding}, the final record causally follows each
recipient's reception of the tuple, while distinct recipients' reception
events remain unordered.  After hiding, the external input events remain
causal predecessors of the final output.  At supplied observations the
final value is $\widetilde q_C(w)=(o_i(w))_{i\in C}$; its effective
observation quotient is the image
$W\twoheadrightarrow O_C\hookrightarrow T_C$.
\end{lemma}

\begin{proof}
\leavevmode
\begin{enumerate}[label=\textup{(\arabic*)},leftmargin=*,itemsep=4pt,topsep=3pt]
  \item \textbf{Observation and knowledge.} First use the finite collector $c_C$ to form the tuple from the
independent observation inputs. Apply the designated data comonoid to
supply one copy to each recipient and retain one for the final output. The recipient $j$ is connected to the receptive response strategy
$\widehat {!}:D(T_C)\to D(1)$ for the unique function $!:T_C\to1$.
  \item \textbf{Observation and knowledge: second comparison.} Its negative tuple input is the recipient's actual reception event;
its positive acknowledgement follows that input. Collect the retained
tuple and all the acknowledgements with
\[
 c_{T_C,(1)_{j\in J}}:
 D(T_C)\otimes\bigotimes_{j\in J}D(1)
 \longrightarrow D(T_C\times 1^J),
\]
and then apply the response for the projection $T_C\times1^J\to T_C$. The empty recipient case uses the nullary product $1^\varnothing=1$.
  \item \textbf{Observation and knowledge.} To verify the actual causal claim, inspect the secured interaction
\emph{before hiding}. Write $e_i$ for original input occurrences,
$r_j$ for each matched recipient reception, $a_j$ for its
acknowledgement, and $f$ for the final record.
  \item \textbf{Observation and knowledge: second comparison.} The collector and
copying dependencies followed by acknowledgement and final collection
supply, for every $i\in C$ and $j\in J$,
\[
 e_i<r_j<a_j<f.
\]
There is no path from $r_j$ to $r_k$ or conversely for distinct
recipients: the copies share predecessors, but their only later common
dependency is the final collector output. The causal graph progresses
from original inputs through collection, copying, acknowledgements, and
the final output.
  \item \textbf{Secured operational interaction.} It contains no directed cycle; every matching used is
therefore secured. The components are receptive courteous, and their
secured composite is receptive courteous by the closure theorem of Section~\ref{sec:paper-1}.
  \item \textbf{Observation and knowledge.} Upon hiding the internal copy, reception, and acknowledgement events,
the induced visible causal order still contains $e_i<f$ for every $i$,
but the internal receptions themselves are recorded by the completed
secured witness, not by events in the hidden composite. Data copying,
collectors, and functoriality of the finite response maps ensure that
the visible label is exactly the input tuple.
  \item \textbf{Quotient for pooled observation.} Image factorization of
the tuple map on worlds gives $O_C$, without selecting a retraction of
its inclusion into $T_C$.
\end{enumerate}
\end{proof}


\begin{theorem}[Communication realizes pooled information]
\label{agt-thm-pooling}
For a fixed factual predicate $P\subseteq W$, public inspection satisfies
\[
 \CK_N^+P=\DK_CP.
\]
Private pooling satisfies $\CK_C^+P=\DK_CP$ and leaves each outsider's
individual relation unchanged.
\end{theorem}
\begin{proof}
\leavevmode
\begin{enumerate}[label=\textup{(\arabic*)},leftmargin=*,itemsep=4pt,topsep=3pt]
  \item \textbf{Ownerwise alternatives.} Write $R_C=\kerpair(q_C)=\bigcap_{i\in C}R_i$. For $i\in C$ the updated
relation is exactly $R_C$, and for every other owner under public inspection
it is contained in $R_C$. Because $C$ is nonempty, the union of all updated
relations is $R_C$.
  \item \textbf{Ownerwise alternatives: completion.} This is already an equivalence relation, so its finite
path closure is $R_C$. Taking boxes gives the first equality. Under private
pooling the members again all have relation $R_C$, while outsiders have
$R_i$ by their universal event relation. The same path calculation within
$C$ proves the second equality.
\end{enumerate}
\end{proof}

\noindent At completed delivery, the joint signal becomes a factor
of every designated recipient's observation.
\begin{center}
\begin{tikzcd}[column sep=large,row sep=large]
W\arrow[r,"o_C",two heads]\arrow[d,"o_i"']
 &O_C\arrow[dl,"p_i"]\\
 O_i&
\end{tikzcd}
\end{center}

Lemma~\ref{agt-lem-completed-delivery} supplies the physical finite
implementation and its completion witness. The final collector is a
barrier for completed delivery, while the independent recipient events
remain unordered. The member policies can then use the actual delivered
tuple. Owner regrouping separately exposes their combined control;
composition coherence for the policy parameters and data wiring makes
the two operations compatible. At intermediate asynchronous histories,
knowledge must be evaluated on the corresponding history-indexed
observations determined by the actual communication history.

\subsection{Complete truthful communication policies in the agent calculus}
\label{sec:agt-announcement-internal-language}

The complete response type and the policy-to-message comparison of
\cref{con:egt-complete-announcement-policies,con:egt-announcement-policy-model}
have the following ordinary-family expressions in the agent calculus:
\[
 \begin{aligned}
 \AnnPol_i &=\prod_{z\in\Omega_i}\{B\subseteq\Omega_i:z\in B\},\\
 X &=W\times\prod_i\AnnPol_i,\\
 U &=\{(B,w):\forall i\ \bar o_i(w)\in B_i\}.
 \end{aligned}
\]
Evaluation is $F(w,\pi)=(b(w,\pi),w)$. The post-communication observations
are $o_j^X(w,\pi)=(\bar o_j(w),\pi_j,b(w,\pi))$ and
$o_j^U(B,w)=(\bar o_j(w),B)$. Each sender retains its own complete
policy as a private record, and deviations keep the world and outsider
policy coordinates fixed.

For enumerated finite observations the message power set has the ordinary
finite function type $2^{\Omega_i}$, with $2=1+1$. The complete truthful
policy carrier is represented by
\[
 \prod_{z:\Omega_i}\sum_{b:2^{\Omega_i}}
       \operatorname{Id}_2(b(z),1).
\]
Finite enumeration, equality, and membership give the displayed policy
semantics. For infinite sets, $\mathcal P(\Omega_i)$ denotes the power set
in the set-based semantics.

By \cref{lem:egt-truthful-policy-extension}, the term
$E_{i,B}(x)=B\cup\{x\}$ is truthful at every observation and satisfies
$E_{i,B}(z)=B$ whenever $z\in B$. It supplies the section
$(B,w)\mapsto(w,(E_{i,B_i})_i)$ of $F$.
The simultaneous epistemic back-lifts and the outsider-preserving strategic
back-lifts of \cref{thm:egt-announcement-policy-comparison} apply to these
same families. Thus inverse image along $F$ commutes with the boxes and
diamonds of the designated relations, their epistemic intersections,
distributed knowledge, and common knowledge, as proved in
\cref{lem:egt-announcement-successor-fibres,cor:egt-announcement-policy-logic}.
The lifts keep the reference complete policy fixed.

For comparisons on $U$ pulled back along $F\times F$,
\cref{cor:egt-announcement-policy-logic} gives
\[
 \NE_X=F^*\NE_U,
\]
together with the analogous equalities for the designated coalitional
stability tests. Public restrictions by predicates on $U$ and action
updates pulled back from $U$ retain this exact comparison. Assertions about
private policy registers remain available on $X$; the comparison applies
to the predicates retained on $U$.

\subsection{Concurrent implementation and coalition announcement ability}

For finite $\Omega_i$, the implementation of
\cref{prop:egt-complete-announcement-realization} realizes the internal
families above by decision generators, policy reparametrization, collectors,
and data copying. At $(w,\pi)$ the delivered observation to owner $j$ is
$(\bar o_j(w),\pi_j,b(w,\pi))$. The reporting decisions are concurrent,
and the joint public delivery follows their collection. The possibly
infinite reference carrier $W$ remains an ordinary index; the physical
interfaces encode the displayed finite observations.

For a goal $\Psi\subseteq U$, coalition announcement ability at $w$ is
the existence of a truthful coalition message tuple that achieves $\Psi$
against every truthful outsider tuple. The extension lemma gives exactly
\begin{equation}
 \mathsf{AnnCan}_C(\Psi)(w)
 \iff\exists\pi_C\ \forall\pi_{-C}\,
 \Psi(F(w,(\pi_C,\pi_{-C}))).
 \label{agt-eq-ann-ability}
\end{equation}
Indeed, extend any winning coalition message explicitly to policies; every
outsider policy evaluates to a legal outsider message. Conversely extend
each possible outsider message while retaining the proposed coalition
policies. These two constructions prove both implications.

\begin{theorem}[Uniform disclosure of distributed knowledge]
\label{agt-thm-disclosure}
For nonempty $C$ and a fixed factual predicate $P\subseteq W$, let
$P^+(B,w)=P(w)$ and $\eta_i(z)=\{z\}$. Then
\begin{align*}
 \DK_CP(w)
 &\iff\forall\pi_{-C}\,
       \CK_N^U(P^+)(F(w,(\eta_C,\pi_{-C})))\\
 &\iff\exists\pi_C\ \forall\pi_{-C}\,
       \CK_N^U(P^+)(F(w,(\pi_C,\pi_{-C}))).
\end{align*}
The same complete disclosure policies $\eta_C$ work at every world.
\end{theorem}
\begin{proof}
\leavevmode
\begin{enumerate}[label=\textup{(\arabic*)},leftmargin=*,itemsep=4pt,topsep=3pt]
  \item \textbf{Ownerwise alternatives.} The policies $\eta_i(z)=\{z\}$ are complete and truthful. At $w$ their
messages cut the compatible component down to a subset of
$R_C[w]=\bigcap_{i\in C}\ker(o_i)[w]$, proving the forward direction
from $\DK_CP(w)$ independently of outsider policies.
  \item \textbf{Reverse implication.} Conversely, let outsiders use the complete tautological policies
$z\mapsto\Omega_j$. Every $v\in R_C[w]$ survives the coalition's
truthful messages and is reachable from $w$ by an individual $i\in C$
epistemic step. Common knowledge implies $P(v)$, hence $\DK_CP(w)$.
  \item \textbf{Policy-level comparison.} This is the complete-policy realization of the exact disclosure result
\cref{cor:egt-uniform-agent-disclosure}, with evaluation justified by
\cref{agt-thm-message-comparison}. The nonempty condition on $C$ is
used only for the final individual-step argument.
\end{enumerate}
\end{proof}

\noindent Policy evaluation retains the underlying world while updating public messages.
\begin{center}
\begin{tikzcd}[column sep=large,row sep=large]
\mathcal X\arrow[r,"F"]\arrow[d,"\pi_W"']
 &\mathcal U\arrow[d,"\pi_W"]\\
 W\arrow[r,equal]&W
\end{tikzcd}
\end{center}

\subsection{The running institution}\label{sec:agt-running-institution}

Return to the two private bits $x,y$ and the complete policy carriers
$\Sigma_i=2^2$ of the introduction. A policy profile specifies the
functions $\pi_1,\pi_2$ independently of the current world, and at
$w=(x,y)$ its response realization produces the action tuple
\[
 (a,b)=(\pi_1(x),\pi_2(y)).
\]
The legal coupling $\pi_1(1)\leq\pi_2(1)$ has $12$ profiles: for
the entries at observation $1$ it allows $(0,0),(0,1),(1,1)$,
while both observation-$0$ responses vary freely. At the world
$w=(0,0)$, let $\pi_1^0$ and $\pi_2^0$ be the constant-zero
policies. The reference $(\pi_1^0,\pi_2^0)$ is legal and yields
$u_1=u_2=0$. Holding either outsider policy fixed at zero, a change
of the current action from $0$ to $1$ gives the deviator payoff $-1$;
changes only at the off-path observation $1$ give payoff $0$. Thus
there is no strictly profitable unilateral legal deviation at this
reference. By contrast, the joint deviation to policies
$\pi_i^*(0)=1$, $\pi_i^*(1)=0$ is legal and yields the action tuple
$(1,1)$, with both utilities equal to $1$. This calculates the
unilateral and grand-coalition deviation fibres using the same
owner-indexed complete-policy space. It also illustrates why
changing two decisions under one regrouped owner requires a
whole-block deviation test.

Initially each owner sees only its own bit, while the joint image distinguishes the pair.
A collector receives both signals and delivers the tuple to the two
owners. Theorem~\ref{agt-thm-pooling} identifies facts about $(x,y)$
that thereby become common knowledge. For example, put
$P=\{(0,0),(1,1)\}\hookrightarrow W$, the assertion $x=y$.
Before delivery, $\ker(o_1)\cap\ker(o_2)$ is equality on $W$,
so $\DK_{\{1,2\}}P(0,0)$ holds; but the union of the two
observation relations connects all four worlds, so
$\CK(P)(0,0)$ fails. After both owners actually receive $(x,y)$,
each updated observation is injective and
$\CK^{+}(P)=P$. In particular, $\CK^{+}(P)(0,0)$ holds.
This equality is asserted only at the completed-delivery boundary,
not while either copy of the public tuple remains in transit. The recipients' new response
policies may depend on the delivered pair; the exact enlargement is
given by Propositions~\ref{agt-prop-info-refinement}
and~\ref{agt-prop-decentralized}.

Regroup the owners to form one control interface while retaining
their individual comparisons; a chosen collective comparison then
evaluates joint decisions. Let the institution record their messages
and commitment. Its execution witness retains the policy profile and
transcript, its target map records the commitment, and its permitted
subobject enforces the rule. The token-guarded executor of
Lemma~\ref{agt-lem-token-guard} receives an actual issuer-produced
authorization event. Final execution depends causally on that event;
the bare type $D(1)$ does not certify the token's authority.
Composition of witnessed spans and of their concurrent realizations
now describes the institutional transaction. A public announcement of
the commitment is a further information-changing action.

\subsection{A concrete finite check of linearity and information}

The finite-data barrier can be tested on two independent bits.
The collector $c:D(2)\otimes D(2)\to D(2\times2)$ waits for both;
the distributor $d:D(2\times2)\to D(2)\otimes D(2)$ sends the
completed pair. Although $cd=\id$ on the paired interface,
$dc$ is not the identity on two independent channels: each output
inherits both input prerequisites. The distinction is operational,
not merely one of value types.

Initially the two observation kernel pairs $R_1$ and $R_2$
distinguish the first and second bit, respectively. Their intersection
is equality on the four possible worlds, so private pooling makes facts
about the entire pair common knowledge once delivery is complete.
For a material resource $R$, this logical change does not supply a
linear contraction $R\to R\otimes R$. The example thereby
separates information sharing from resource duplication.

\subsection{A running example}\label{sec:agt-running-example}

Consider the same owners with private signals $x$ and $y$. They may
exchange messages directly or submit them to an institution that
redistributes a joint observation. Let $K$ be the dependent carrier of
accepted completed executions
\[
 K=\{(x,y,\pi_1,\pi_2)\in W\times L:
           \pi_1(x)=1\ \wedge\ \pi_2(y)=1\}.
\]
Its projection to the full institutional record retains the source
world and complete owner-policy profile. The issuer receives the two
recorded positive action outputs through a collector, checks the
predicate defining $K$, and only on an accepted completed record
produces its designated positive authorization event $t$.
The token-guarded executor receives the matching authorization input
and produces the final execution event $e$ with $t<e$ in the
completed interaction witness. The legal joint policy
$(\pi_1^*,\pi_2^*)$ at $(0,0)$ belongs to $K$, so this witness
carrier is inhabited. The type $D(1)$ alone does not certify an
issuer-produced $t$; the actual occurrence and its causal
predecessors are retained in the institution's witness. A later
decision requires such a linear authorization token. This example tracks the successive questions
answered by the construction: can the messages be delivered, what does
each recipient learn, which subsequent policies become available, and
which actual authorization event enables the final action?

\paragraph{From completed communication to continuing agency.}
The preceding constructions now connect complete agent policies to
causally realized communication and to the epistemic information created
by delivery. The remaining part uses the same owner and observation
interfaces after a history: residual policies describe continuing
execution, modal fixed points describe repeated goals, and exact
reductions identify the distinctions that must be retained for later
strategic tests.

\part{Continuing agency and exactness}

\section{Continuing Agents and Temporal Agency}\label{sec:paper-12}

An agent's control does not end when one history has been executed.
The remaining interaction is a residual strategy, and the remaining
policy is the restriction of an original complete contingent plan.
Keeping both permits agents to continue without reconstructing
control from an arbitrary interleaving of events. We construct
residual controllers and temporal fixed-point objectives, relate
witnessed transitions to epistemic updates, and study visible
refinements that preserve the designated alternatives. Finite
backward induction and subgame perfection then check the sequential
interpretation.

\subsection{Fixed points in assertion fibres and their transport}

The assertion fibres are complete lattices, so the fixed-point
construction of \citet{Tarski1955} applies to monotone modal
operators. We write $\mu$ for the least and $\nu$ for the greatest
fixed predicate. These operations express distinct temporal tasks:
least fixed points capture progress towards a goal, while greatest
fixed points capture its persistence. Their transport must respect
the witnessed successor relations on which the operators depend.

\begin{construction}[Predicate fixed points]
\label{con:egt-fixedpoint}
For a monotone $F:\mathsf{Pred}(X)\to\mathsf{Pred}(X)$, put
\[
 \mu F=\bigwedge\{P:F(P)\leq P\},\qquad
 \nu F=\bigvee\{P:P\leq F(P)\}.
\]
The set-indexed joins and meets are pointwise existential and universal
quantification in the chosen predicate doctrine.
\end{construction}

\begin{theorem}[Predicate induction and coinduction; recalled]
\label{prop:egt-fixedpoint}
$\mu F$ and $\nu F$ are respectively the least and greatest fixed points.
They satisfy $F(P)\leq P\Rightarrow\mu F\leq P$ and
$P\leq F(P)\Rightarrow P\leq\nu F$.
The same holds on finite products of predicate lattices.
\end{theorem}
\begin{proof}[Reference]
\leavevmode
\begin{enumerate}[label=\textup{(\arabic*)},leftmargin=*,itemsep=4pt,topsep=3pt]
  \item \textbf{Fixed-point argument.} Apply the complete-lattice fixed-point theorem of \citet{Tarski1955}. The indicated products and pointwise predicate lattices are complete.
\end{enumerate}
\end{proof}

\begin{definition}[Positive modal fixed-point language]
\label{def:egt-positive-modal-fixed-point-language}
Fix a state set $X$, atomic predicates $A_p\hookrightarrow X$,
relations $R_a\hookrightarrow X^2$, and predicate variables $Z$ drawn
from a set $\mathsf{Var}$. Formulas have the grammar
\[
 \varphi::=p\mid Z\mid\top\mid\bot\mid
 (\varphi\wedge\varphi)\mid(\varphi\vee\varphi)\mid
 (\varphi\Rightarrow\varphi)\mid[R_a]\varphi\mid
 \langle R_a\rangle\varphi\mid\mu Z.\varphi\mid\nu Z.\varphi.
\]
An environment $\eta:\mathsf{Var}\to\mathsf{Pred}(X)$ gives the
interpretation $\llbracket\varphi\rrbracket_{X,\eta}\in\mathsf{Pred}(X)$
by interpreting atoms as $A_p$, variables as $\eta(Z)$, and connectives
and modalities by the already-defined operations. A bound variable must
occur \textbf{positively} in its binder: antecedents of implication
reverse polarity, consequents preserve it, and the other constructors
preserve polarity. Hence
$F_{\varphi,\eta}(U)=\llbracket\varphi\rrbracket_{X,\eta[Z\mapsto U]}$
is monotone, and the binder clauses are
\[
 \llbracket\mu Z.\varphi\rrbracket_{X,\eta}=\mu F_{\varphi,\eta},
 \qquad
 \llbracket\nu Z.\varphi\rrbracket_{X,\eta}=\nu F_{\varphi,\eta}.
\]
The two fixed points are those of
Construction~\ref{con:egt-fixedpoint}; they are not extra syntactic
actions on concurrent strategies.
\end{definition}
\begin{lemma}[Transport across a retraction]
\label{lem:egt-retraction-fixedpoints}
Let $\rho:L\to M$ and $J:M\to L$ be monotone maps between complete lattices
with $\rho J=\operatorname{id}_M$. For monotone $F:M\to M$, put
$\widehat F=JF\rho$. Then
\[
 \mu\widehat F=J\mu F,\qquad\nu\widehat F=J\nu F.
\]
\end{lemma}
\begin{proof}
\leavevmode
\begin{enumerate}[label=\textup{(\arabic*)},leftmargin=*,itemsep=4pt,topsep=3pt]
  \item \textbf{Fixed-point argument.} If $Fx=x$, then $\widehat F(Jx)=Jx$. If $\widehat Fy=y$, then
$F(\rho y)=\rho y$ and $J\rho y=y$. Thus $J$ and $\rho$ restrict to
inverse order isomorphisms between the two fixed-point posets. They preserve
their least and greatest elements.
\end{enumerate}
\end{proof}

\subsection{Residual strategies and continuing agents}
\label{sec:egt-continuation}

To interpret action after a history, we must retain both the
strategy's remaining event structure and the owner's contingent
responses at later decision positions. Pointed residuals provide the
operational data. Restriction of a complete policy gives the corresponding
future response family, while epistemic action models record information
from actually completed rounds. The resulting continuation interface
connects temporal objectives to the same control and observation data
already used for one-stage agency.

Each completed public round can be read as an enabled epistemic
action event from Section~\ref{sec:egt-action-model-update}. Its
state map records the concurrent transition, while its observation
records the information actually delivered at completion. Intermediate
histories remain represented by operational configurations. Thus
successive public updates do not impose a total schedule on events
that were independent within a round.

A continuing state can carry commitments or goals as ordinary
records. Their persistence, discharge, and revision require explicit
update maps; current preference selection alone does not determine
such transitions. This is the point of contact with intention accounts
such as \citet{CohenLevesque1990Intention}. Here continued agency is
established for the displayed record and residual operations rather
than by assigning an intention law to every strategy.

\subsubsection{Complete policies and their residual controllers}
\label{sec:egt-residual-controllers}

A complete local policy specifies a response to every finite
observation history. Its residuals form a deterministic coalgebra,
in the sense of \citet{Rutten2000UniversalCoalgebra}, whose state
records the remaining response function. The controller realizes the
policy by updating that state after each observed input. The same
coalgebraic perspective also gives a criterion for when the full
response family admits a finite-state implementation.

\begin{definition}[Local response controller]
\label{def:egt-local-controller}
For a sequential local observation interface $O$ and action interface
$M$, a controller is a coalgebra
\[
 \delta:S\longrightarrow(M\times S)^O.
\]
Write $\delta(s)(o)=(a(s,o),n(s,o))$. A controller map $f:S\to S'$
preserves both components:
$a'(f(s),o)=a(s,o)$ and $n'(f(s),o)=f(n(s,o))$.
A pointed controller also specifies its initial state $s_0$.

Let $O^+$ and $O^*$ denote nonempty and arbitrary finite words,
respectively. A complete local policy is an element
$b\in\mathcal B=M^{O^+}$: its value at a word is the action
after the final observation of that word. Set
\[
 \zeta:\mathcal B\to(M\times\mathcal B)^O,\qquad
 \zeta(b)(o)=(b(o),b_o),\qquad b_o(v)=b(ov).
\]
\end{definition}

\begin{theorem}[Final behaviour and the generated residual controller]
\label{thm:egt-residual-controller}
The coalgebra $(\mathcal B,\zeta)$ is final.
For a complete policy $b$, put
\[
 b_u(v)=b(uv)\quad(u\in O^*,\ v\in O^+),\qquad
 S_b=\{b_u:u\in O^*\}\subseteq\mathcal B.
\]
It is a pointed subcoalgebra with initial state $b=b_{\varnothing}$.
Every pointed controller implementing $b$ has a unique surjective
pointed controller map from its reachable states onto $S_b$.
Consequently $b$ has a finite-state implementation if and only if
$S_b$ is finite. In that case $|S_b|$ is the least possible number
of reachable states of an implementation.
\end{theorem}

\noindent The residual controller coalgebra sends a policy state to its one-step response and successor state.
\begin{center}
\begin{tikzcd}[column sep=large,row sep=large]
S\arrow[r,"\zeta"]\arrow[d,"f"']
 &F(S)\arrow[d,"Ff"]\\
 T\arrow[r,"\theta"']&F(T)
\end{tikzcd}
\end{center}


\begin{proof}
\leavevmode
\begin{enumerate}[label=\textup{(\arabic*)},leftmargin=*,itemsep=4pt,topsep=3pt]
  \item \textbf{Fixed-point argument.} For a controller $\delta$, define its behaviour recursively by
\[
 \operatorname{beh}_\delta(s)(o)=a(s,o),\qquad
 \operatorname{beh}_\delta(s)(ov)
   =\operatorname{beh}_\delta(n(s,o))(v)
       \quad(v\in O^+).
\]
Induction on word length defines a unique function
$S\to M^{O^+}$. The two equations assert exactly that this
function is a coalgebra map into $(\mathcal B,\zeta)$.
  \item \textbf{Reverse implication.} Conversely, every such map must satisfy those equations;
induction proves uniqueness. This proves finality, including
empty carriers whenever the displayed coalgebra exists.
  \item \textbf{Closure of controller derivatives.} For a derivative $b_u$, its output on $o$ is $b(uo)$ and
its next behaviour is $b_{uo}$, which lies in $S_b$. Hence restriction of $\zeta$ is well defined. Its behaviour
map is its inclusion in $\mathcal B$, by the same recursion.
  \item \textbf{Construction of the comparison.} Let $\delta$ with initial state $s_0$ implement $b$, so that
$\operatorname{beh}_\delta(s_0)=b$. Define
$s_{\varnothing}=s_0$ and $s_{uo}=n(s_u,o)$. The reachable carrier is $\{s_u:u\in O^*\}$.
  \item \textbf{Lifting along the surjection.} Induction gives
$\operatorname{beh}_\delta(s_u)=b_u$. The restricted behaviour map therefore lands in $S_b$,
is surjective, and preserves initial state, outputs, and
transitions. A pointed controller map must send $s_u$ to
$b_u$ by induction on $u$, so it is unique, even when several
words reach the same state.
  \item \textbf{Verification of the universal property.} Surjectivity implies that a
finite reachable implementation has finite $S_b$ and at
least its cardinality. Conversely, finite $S_b$ is itself
a pointed implementation with that many states.
\end{enumerate}
\end{proof}


The behavioural quotient identifies controller states exactly when
they prescribe the same future responses for every finite local
observation word. It therefore preserves the local input--output policy.
Preferences, permissions, and observations of additional internal
records need not descend along that quotient unless their corresponding
tests are retained.

\begin{proposition}[Concurrent realization of a finite controller step]
\label{prop:egt-controller-realization}
For finite $S,O,M$, the uncurried response function
$\delta:S\times O\to M\times S$ has the one-step realization
\[
 \begin{aligned}
 D(S)\otimes D(O)
 &\xrightarrow{\operatorname{col}}D(S\times O)\\
 &\xrightarrow{\widehat\delta}D(M\times S)
 \xrightarrow{\operatorname{dist}}D(M)\otimes D(S).
 \end{aligned}
\]
Connecting the returned state to the next local step realizes,
for every finite observation word, the output sequence of the
pointed controller. If the controller implements $b$, its final
output on a nonempty word $u$ is $b(u)$.
\end{proposition}
\begin{proof}
\leavevmode
\begin{enumerate}[label=\textup{(\arabic*)},leftmargin=*,itemsep=4pt,topsep=3pt]
  \item \textbf{Fixed-point argument.} The finite-data collector and distributor and
Proposition~\ref{prop:egt-functions} realize the displayed ordinary
function on value inputs. Thus the value pair $(s,o)$ yields
exactly $(a(s,o),n(s,o))$. Induction on the number of connected
steps shows that the state after $u$ is $s_u$ and that the output
at its last step is
$\operatorname{beh}_\delta(s_0)(u)$.
  \item \textbf{Observation and knowledge.} The implementation equation makes this $b(u)$. Each finite composite is an arrow of the existing concurrent
semantics. The collector waits for the local state and observation;
its synchronization is part of that arrow and is not an
isomorphism identifying $D(S\times O)$ with
$D(S)\otimes D(O)$.
\end{enumerate}
\end{proof}

The words used in the controller are sequential observations at
one specified local interface. Concurrent events belonging to other
interfaces retain their own partial orders, so the implementation
does not choose a global interleaving. The constructed realization
covers each finite connected segment, not an unspecified recursive
process beyond the declared local protocol.

Petri games provide another setting for causally informed
strategies of distributed processes
\citep{FinkbeinerOlderog2017PetriGames}. The comparison here concerns
implementation: which complete local policies have controllers
respecting the observation available at each step? The residual
coalgebra answers that question without conflating global concurrency
with the chosen local sequence of observations.

\subsubsection{Maintaining and eventually achieving collective goals}

\begin{definition}[Public states, histories, and winning regions]
\label{def:egt-public-state-winning-region}
For the finite public-state interaction specified below, let $W$ be the
set of publicly observed round states, $A_i(w)$ the finite nonempty
set of agent $i$'s available actions at $w$, and
$\delta_w:\prod_i A_i(w)\to W$ its deterministic transition.
The \textbf{history carrier} is the set of finite legal histories
\[
 \mathsf{Hist}=\coprod_{n\geq0}
 \{(w_0,a^0,w_1,\ldots,a^{n-1},w_n):
        \delta_{w_k}(a^k)=w_{k+1}\ (k<n)\},
\]
with $\operatorname{last}:\mathsf{Hist}\to W$.
A \textbf{history-dependent pure policy} of owner $i$ is a section
of the pulled-back action bundle, equivalently a family
$\pi_i(h)\in A_i(\operatorname{last}(h))$ for every $h\in\mathsf{Hist}$.
Let $\mathsf{Play}$ be the carrier of infinite legal
transition-consistent action/state sequences. For initial $w$ and a
complete policy profile $\pi$, deterministic recursion gives
$\operatorname{play}(w,\pi)\in\mathsf{Play}$.
For a \textbf{temporal objective} $\Phi\hookrightarrow\mathsf{Play}$,
the \textbf{winning region} of coalition $C$ is
\[
 \operatorname{Win}_C(\Phi)
 =\{w\in W:\exists\pi_C\ \forall\pi_{-C},\;
            \operatorname{play}(w,\pi_C,\pi_{-C})\in\Phi\}
 \hookrightarrow W.
\]
All agents observe the completed public history, not their opponents'
current simultaneous choices. The predecessor construction below
derives this ability using dependent action bundles.
\end{definition}
Temporal ability requires more than one-stage effectivity: a
coalition must choose at successive public histories while outsiders
may respond at each stage. The predecessor operator uses dependent
quantifiers in the existing predicate doctrine. The ordinary
classical $\mathbf{Set}$ instance supplies the duality between
failed coalition ability and a counterpolicy, expressed by
$\neg\exists a_C\,\forall a_{-C}P\Rightarrow
 \forall a_C\,\exists a_{-C}\,\neg P$.
The fixed-point construction below retains the quantifier order when
moving from one round to repeated play.

\begin{construction}[Finite public-state interaction and controllable predecessor]
A finite public-state interaction has a finite state set $W$, a finite nonempty action set
$A_i(w)$ for each agent at each state, and a transition function
\[
 \delta_w:\prod_i A_i(w)\to W.
\]
Players observe the finite public history before each simultaneous
round. A coalition policy assigns its action tuple to each such
history. Opponents have policies of the same kind. This is an
economic policy space in the preceding sense; one round is a causal
decision protocol with concurrent action inputs and a deterministic
transition rule. Finite prefixes are obtained by composition of these
rounds, retaining the current state as the interface.

To expose the categorical structure when action sets depend on the current
state, form the action bundles
\begin{align*}
 E_C&=\{(w,a_C):a_C\in\prod_{i\in C}A_i(w)\},\\
 E&=\{(w,a_C,a_{-C}):a_i\in A_i(w)\text{ for all }i\}.
\end{align*}
Let
\[
 E\xrightarrow{\pi_C}E_C\xrightarrow{\pi_W}W
\]
forget first the outsider actions and then the coalition actions, and let
$\delta:E\to W$ be the transition map. For a state predicate $X$, define
\begin{equation}
\label{eq:egt-pre-doctrine}
 \operatorname{Pre}_C(X)
 =\exists_{\pi_W}\,\forall_{\pi_C}\,\delta^*X.
\end{equation}
Pointwise this is exactly
\[
 \operatorname{Pre}_C(X)(w)
 \quad\Longleftrightarrow\quad
 \exists a_C\ \forall a_{-C},\
                  X(\delta_w(a_C,a_{-C})).
\]
The operator is monotone. It is the state-indexed version of
Equation~\eqref{eq:egt-ability} and is constructed entirely from
reindexing and the two quantifier adjoints of the assertion doctrine.
\end{construction}

\begin{proposition}[Epistemic product update of a public concurrent round]
\label{prop:egt-public-round-action-update}
For the finite public-state interaction above, let the action-occurrence
set be the disjoint union
\[
 E^{\mathrm{round}}=\coprod_{w\in W}\prod_i A_i(w),
\]
so an occurrence is a tagged pair $(w,a)$. Give it enabling predicate
$P(u,(w,a))\iff u=w$. Take equality as the prior public-state
relation and as the action relation representing the publicly revealed
completed tuple. Then its product-update carrier is canonically
the original joint-action bundle $E$ used to define
Equation~\eqref{eq:egt-pre-doctrine}, with the transition
postcondition
\[
 t(u,(u,a))=\delta_u(a).
\]
Under this identification the controllable predecessor is
$\exists_{\pi_W}\forall_{\pi_C}t^*X$. This action-model presentation
uses a completed joint action as one public update, while the
underlying operational realization uses the independent concurrent
choice events and the collector's synchronization barrier.
\end{proposition}
\begin{proof}
\leavevmode
\begin{enumerate}[label=\textup{(\arabic*)},leftmargin=*,itemsep=4pt,topsep=3pt]
  \item \textbf{Monoidal comparison.} Comprehension consists precisely of the pairs $(u,(u,a))$, so
forgetting the redundant first copy of $u$ is a bijection with the
original action bundle $E$. Both action and old-state relations are
equality, hence their restricted product is equality on the updated
carrier. The postcondition is the original transition map.
  \item \textbf{Ownerwise alternatives.} The
quantifier expression is Equation~\eqref{eq:egt-pre-doctrine} under
the carrier bijection. The circuit of
Theorem~\ref{thm:egt-causal-protocol} realizes each action tuple by
independent owner-controlled response protocols, followed by the
finite collector and transition response; no causal order between
simultaneous decisions is added.
\end{enumerate}
\end{proof}

\begin{definition}[Temporal objectives]
\label{def:egt-temporal-objectives}
Let $\mathsf{Play}$ be the carrier of infinite legal plays
$\lambda=(w_0,a^0,w_1,a^1,\ldots)$ of the public-state interaction.
For state predicates $P,Q\hookrightarrow W$, define the play predicates
\[
 \begin{aligned}
 \llbracket P\,\mathsf U\,Q\rrbracket
   &=\{\lambda:\exists n\in\mathbb N,\;
         w_n\in Q\ \wedge\ \forall k<n\;(w_k\in P)\},\\
 \llbracket\mathsf G P\rrbracket
   &=\{\lambda:\forall n\in\mathbb N,\;w_n\in P\}.
 \end{aligned}
\]
The first requires reaching $Q$ after finitely many rounds while
maintaining $P$ beforehand; the second requires $P$ at every state.
The coalition-enforcement predicate for either objective is
$\operatorname{Win}_C(\llbracket\cdot\rrbracket)\hookrightarrow W$
of Definition~\ref{def:egt-public-state-winning-region}, quantifying
one coalition policy against every outsider policy.
\end{definition}
\begin{theorem}[Temporal ability and fixed points]
\label{thm:egt-temporal-ability}
For a finite public-state interaction,
\begin{align*}
 \operatorname{Win}_C(P\,\mathsf U\,Q)
   &=\mu X\bigl(Q\cup(P\cap\operatorname{Pre}_C(X))\bigr),\\
 \operatorname{Win}_C(\mathsf G P)
   &=\nu X\bigl(P\cap\operatorname{Pre}_C(X)\bigr).
\end{align*}
Every state in either winning set has a policy depending only on the
current state. For the until objective, the construction also gives
a finite bound on the time to reach $Q$.
\end{theorem}

\noindent The controllable predecessor is the composite of universal quantification over opponents and existential quantification over coalition choices.
\begin{center}
\begin{tikzcd}[column sep=large,row sep=large]
E\arrow[r,"\pi_C"]\arrow[dr,"\pi_W\pi_C"']
 &E_C\arrow[d,"\pi_W"]\\
 &W
\end{tikzcd}
\end{center}


\begin{proof}
\leavevmode
\begin{enumerate}[label=\textup{(\arabic*)},leftmargin=*,itemsep=4pt,topsep=3pt]
  \item \textbf{Fixed-point argument.} For until, iterate $X_0=\varnothing$ and
$X_{n+1}=Q\cup(P\cap\operatorname{Pre}_C(X_n))$. Because $W$ is finite this increasing sequence reaches its least fixed
point $L$ after finitely many steps. Give each state in $L$ its least
positive index of entry.
  \item \textbf{Witness and legality comparison.} States of index $1$ lie in $Q$. At every
other state take an action tuple witnessing predecessor membership in
the preceding stage. There are only finitely many such states, so these
finitely many witnesses assemble by finite recursion into a
state-dependent coalition policy.
  \item \textbf{Ownerwise alternatives.} Against every outsider action it moves
to a state of smaller index. Until $Q$ is reached, the current state lies
in $P$. The index decreases at each step, so this state-dependent
policy reaches $Q$ in fewer steps than the starting index.
  \item \textbf{Ownerwise alternatives: second comparison.} To prove the converse, fix any coalition policy at a state outside
$L$. Such a state is not in $Q$. If it is outside $P$, the objective
has already failed.
  \item \textbf{Fibrewise calculation.} At a state in $P\setminus L$, fixed-point equality implies
$w\notin\operatorname{Pre}_C(L)$. Expanding the predecessor quantifiers and
using the classical finite specialization fixed above gives
\[
 \forall a_C\ \exists a_{-C},\qquad
       \delta_w(a_C,a_{-C})\notin L.
\]
For the finitely many state/action pairs, assemble one such outsider
response for each pair by finite recursion.
  \item \textbf{Ownerwise alternatives.} Against the fixed coalition
policy, this gives an outsider policy on public histories: evaluate that fixed policy at
the history and use the selected response. It keeps the play outside
$L$ until $P$ is left, if that happens. The play never reaches $Q$;
in either case the until objective fails. Thus no state outside $L$
is winning.
  \item \textbf{Policy-level comparison.} For safety let $S=\nu X(P\cap\operatorname{Pre}_C(X))$. For the finitely many states in $S$, assemble predecessor witnesses
that stay in $S$ into a state-dependent policy by finite recursion. Since $S\subseteq P$, this policy enforces safety.
  \item \textbf{Reverse implication.} Conversely let $Y$ be the set of states from which
some policy enforces safety. At $w\in Y$, the first coalition action
must send every outsider completion to another state of $Y$. Indeed the winning continuation after that first history, with its
prefix fixed as a parameter, is a valid policy from the successor
state.
  \item \textbf{Coinduction and bisimulation.} Also $w\in P$. Hence
$Y\subseteq P\cap\operatorname{Pre}_C(Y)$. Coinduction gives
$Y\subseteq S$ and proves the converse.
\end{enumerate}
\end{proof}


The counterpolicy constructed in the until argument responds to
a fixed coalition policy: the conclusion is
$\forall\sigma_C\exists\sigma_{-C}$, not
$\exists\sigma_{-C}\forall\sigma_C$.
Both policies use completed public histories as inputs. Exchanging
these quantifiers would claim a single outsider policy defeating all
coalition policies, which is not what the construction proves.

A one-round matching game already separates these statements.
Two agents simultaneously choose bits, and success means equality.
Neither can enforce equality against the other's arbitrary choice,
and neither can enforce inequality against every alternative. Thus the
negation of one coalition's effectivity is not the other coalition's
positive ability without further analysis of the quantifiers.

\begin{corollary}[Translation of alternating-time state formulas]
\label{cor:egt-atl}
On these interactions, the state language generated by propositional connectives and the operators
$\langle\!\langle C\rangle\!\rangle\mathsf X$,
$\langle\!\langle C\rangle\!\rangle\mathsf U$, and
$\langle\!\langle C\rangle\!\rangle\mathsf G$
has an exact interpretation in the assertion extension:
next is $\operatorname{Pre}_C$, until is the displayed least fixed
point, and always is the displayed greatest fixed point.
\end{corollary}
\begin{proof}
\leavevmode
\begin{enumerate}[label=\textup{(\arabic*)},leftmargin=*,itemsep=4pt,topsep=3pt]
  \item \textbf{Fixed-point argument.} Induct on state formulas. Atomic and propositional cases are the
predicate interpretation. The next case is the definition of controllable
predecessor.
  \item \textbf{Ownerwise alternatives.} Once the immediate subformulas have been interpreted as
state predicates, Theorem~\ref{thm:egt-temporal-ability} supplies the
until and always cases. This also allows nested coalition operators,
since their interpretations are again state predicates.
\end{enumerate}
\end{proof}

This gives the specified connection to temporal and alternating-time
reasoning \citep{AlurEtAl2002}. The linear semantics constructs and
composes the rounds; ownerwise policy quantifiers specify control;
and least and greatest fixed points express progress and persistence.
The resulting policies and counterpolicies are exhibited by the
finite-history construction rather than inferred from the modal
notation alone.

\subsubsection{What remains after a history}

\begin{definition}[Residual games, continuation policies, and off-path histories]
\label{def:egt-residual-game}
For a receptive courteous strategy
$s:S\to A^\perp\parallel B$ and a finite source configuration
$x\in\operatorname{Conf}(S)$ with $s(x)=a_x\amalg b_x$, its
\textbf{pointed operational residual} is the restriction
\[
 s/x:S/x\longrightarrow(A^\perp\parallel B)/s(x)
 \cong (A/a_x)^\perp\parallel(B/b_x).
\]
It records the remaining events of that particular strategy after the
actually exhibited source history $x$; its validity is proved in
Proposition~\ref{prop:egt-residual-strategy}.

For a complete-policy game with a total family of histories
$\mathsf{Hist}$, a \textbf{continuation game} or \textbf{subgame} at
$h\in\mathsf{Hist}$ is the owned open arena
$G_h=(\vec\Sigma^{\,h},g_h):A_h\to B_h$ together with the restriction
\[
 r_h:\Sigma\to\Sigma^h,\qquad
 g_h:\Sigma^h\to\mathcal C(A_h,B_h),\qquad
 g^h=g_hr_h.
\]
A continuation policy is an element of $\Sigma_i^h$; its restriction
from a complete policy is $r_{h,i}(\sigma_i)$. If
$\mathsf{Reach}(\sigma)\subseteq\mathsf{Hist}$ denotes the histories
on the play of $\sigma$, then
$\mathsf{Off}(\sigma)=\mathsf{Hist}\setminus\mathsf{Reach}(\sigma)$
is the \textbf{off-path history} set. A subtree game is defined even at
such a history; it is not asserted to be the residual of the particular
strategy realized by $\sigma$. The comparison at reached histories is
constructed below.
\end{definition}

\begin{construction}[The continuing agent at a history]
\label{con:egt-continuing-agent}
For the continuation presentation $G_h$ and ownerwise restriction $r_h$,
the continuing agent of identity $i$ has interface
\[
 \mathsf{Ag}_i(G_h)
 =\bigl(i,\Sigma_i^h,\varepsilon_{\Sigma_i^h}\bigr),
 \qquad r_{h,i}:\Sigma_i\to\Sigma_i^h.
\]
Its situated observations are the displayed observation terms on the
history or continuation carrier. Its comparison is $P_{i,h}$ on the
continuation-arrow sort. Pulling that comparison back along $g_h$ and then
$r_h$ gives the complete-policy comparison
\[
 (r_h\times r_h)^*(g_h\times g_h)^*P_{i,h}.
\]
Functoriality of inverse image identifies this with feedback through
$g^h=g_hr_h$. Thus the history changes the agent's continuation interface
while retaining its identity and the complete-policy restriction map.
Pointed operational residuals supply the exhibited implementation at reached
histories; the full continuation family retains the off-path interfaces used
by sequential rationality.
\end{construction}

The next notation packages a continuation at a history, but a
visible configuration alone may not determine the full residual.
For a general concurrent strategy, the source configuration must also
be retained. The pointed-history object supplies precisely that extra
information.

Write $r_h:\Sigma\to\Sigma^h$ for restriction of complete policies to
the future of a history, and $g_h:\Sigma^h\to\mathcal C(A_h,B_h)$
for residual play. The history-indexed open play map is $g^h=g_hr_h$.
For every continuation comparison $P_h$, its complete-policy feedback is
\[
 (g^h)^\sharp P_h=(r_h\times r_h)^*(g_h\times g_h)^*P_h.
\]
The equality is the functorial action of coplay on residual
realizations. The operational construction first produces
continuations at pointed histories; the finite complete-policy trees
then organize all continuation choices, including histories not
reached in the current play. This retains the alternatives needed
for subgame-perfect evaluation.

\begin{construction}[Residual event structure and residual strategy]
For a finite configuration $x$ of an event structure $E$, define its
residual $E/x$ as follows. Its events are the $e\notin x$ such that
$x\cup\{e\}\in\operatorname{Con}_E$; its order is the restriction of $\leq_E$;
and its finite consistency relation is
\[
 Z\in\operatorname{Con}_{E/x}\quad\Longleftrightarrow\quad
 x\cup Z\in\operatorname{Con}_E.
\]
Residual polarities are inherited.
\end{construction}

\begin{lemma}
\label{lem:egt-residual}
The construction $E/x$ is an event structure, and for finite sets of
residual events $z$,
\[
 z\in\operatorname{Conf}(E/x)
 \quad\Longleftrightarrow\quad
 x\cup z\in\operatorname{Conf}(E).
\]
If $x\subseteq y$ are configurations, there is a canonical isomorphism
\[
 (E/x)/(y\setminus x)\cong E/y.
\]
\end{lemma}
\begin{proof}
\leavevmode
\begin{enumerate}[label=\textup{(\arabic*)},leftmargin=*,itemsep=4pt,topsep=3pt]
  \item \textbf{Causal execution.} Finite pasts and the consistency axioms are inherited: adjoining a cause
of a consistent event preserves consistency with $x$. Every event in a
residual past is compatible with $x$, by the same axiom. If $z$ is a
residual configuration, the causes of an event of $z$ either lie in $x$
or are residual causes and lie in $z$.
  \item \textbf{Residual configuration equivalence.} Thus $x\cup z$ is down-closed,
and it is consistent by definition. Conversely, down-closure and
consistency of $x\cup z$ give the corresponding properties in the
residual.
  \item \textbf{Causal execution.} The events of either iterated residual are exactly the events outside
$y$ compatible with $y$. On these events the orders are both the
restriction of the original order, and consistency on both sides says
that adjoining $y$ is consistent in $E$. The identity on remaining events
is the asserted isomorphism.
\end{enumerate}
\end{proof}

\begin{figure}[htbp]
\centering
\begin{tikzcd}[column sep=large,row sep=large]
E \arrow[r,"\text{after }x"] \arrow[dr,"\text{after }y"']
 & E/x \arrow[d,"\text{after }y\setminus x"] \\
 & E/y
\end{tikzcd}
\caption{Continuation is independent of whether a history is exposed in
one step or in successive prefixes. The arrows denote residual operations.}
\end{figure}

\begin{proposition}[Residual strategies]
\label{prop:egt-residual-strategy}
If $s:S\to A$ is a receptive courteous strategy and
$x\in\operatorname{Conf}(S)$, restriction of the labelling map gives a
receptive courteous strategy
\[
 s/x:S/x\longrightarrow A/s(x).
\]
\end{proposition}

\begin{proof}
\leavevmode
\begin{enumerate}[label=\textup{(\arabic*)},leftmargin=*,itemsep=4pt,topsep=3pt]
  \item \textbf{Causal execution.} First the restricted map is well defined. Let $e$ be an event of $S/x$. Then $x\cup\{e\}$ is consistent.
  \item \textbf{Causal execution: second comparison.} Adjoining the finitely many causes of $e$
one at a time shows that $x\cup\downarrow e$ is a finite configuration of
$S$. Its image under $s$ is therefore a configuration of $A$. Local
injectivity on this configuration implies $s(e)\notin s(x)$, while
$s(x)\cup\{s(e)\}$ is consistent. Hence $s(e)$ is an event of $A/s(x)$.
  \item \textbf{Causal execution.} If $z\in\operatorname{Conf}(S/x)$, then
$x\cup z\in\operatorname{Conf}(S)$ by
Lemma~\ref{lem:egt-residual}. Thus $s(x\cup z)$ is a configuration of $A$. Local injectivity separates the labels in $s(x)$ from those in $s(z)$, so
$s(z)$ is a configuration of $A/s(x)$. Local injectivity of $s/x$ is
inherited from $s$.
  \item \textbf{Secured operational interaction.} For receptivity, let $z$ be a residual configuration and let $a$ be a
negative event enabled after $(s/x)(z)$ in $A/s(x)$. Then
\[
 s(x)\cup s(z)\cup\{a\}
\]
is a configuration of $A$. Receptivity of $s$ at $x\cup z$ gives a unique
event $e$ with $s(e)=a$ and
$x\cup z\cup\{e\}\in\operatorname{Conf}(S)$.
  \item \textbf{Lifting through the comparison.} Since $a\notin s(x)$, local
injectivity implies $e\notin x$; hence $e$ lies in $S/x$ and is the unique
residual lift.
  \item \textbf{Secured operational interaction.} For courtesy, suppose $e\lessdot_{S/x}e'$. If some $d$ satisfied
$e<d<e'$ in $S$, then $d\notin x$ because $x$ is down-closed while
$e\notin x$. Since $e'$ is compatible with $x$ and $d\leq e'$, closure of
consistency under adjoining causes makes $d$ compatible with $x$, so $d$
would be an event of $S/x$.
  \item \textbf{Secured operational interaction: completion.} This contradicts immediacy there. Therefore
$e\lessdot_Se'$. Whenever courtesy for $s$ requires
$s(e)\lessdot_A s(e')$, that immediate edge also remains immediate in the
residual arena: an event strictly between the two residual labels would
already be an event strictly between them in $A$. Hence the courtesy
condition restricts to the residual strategy.
\end{enumerate}
\end{proof}


The source history $x$ cannot be discarded: different source
configurations may have the same visible history but distinct future
interactions. Residualization also does not claim that every arbitrary
future policy extends to the original complete policy space. It
constructs the continuation of a specified strategy and the exact
restriction maps that are available in the presentation.

\begin{construction}[The dependent object of pointed operational histories]
\label{con:egt-pointed-histories}
For an open arena $G=(\vec\Sigma,g):A\to B$, take all triples
\[
 (\sigma,s:S\to A^\perp\parallel B,x),
 \qquad [s]=g(\sigma),\quad x\in\operatorname{Conf}(S),
\]
using the set of countable strategy presentations on subsets of $\mathbb N$ from
Convention~\ref{conv:egt-base}. Quotient by strategy isomorphisms preserving
$\sigma$ and carrying the marked configuration to the marked configuration.
Write $H_G$ for this quotient. It has canonical projections to $\Sigma$ and
to $\operatorname{Conf}(A)\times\operatorname{Conf}(B)$, with the latter
defined by $s[x]=a\amalg b$. Residualization is a term in the pulled-back
dependent arrow family:
\[
 [\sigma,s,x]\longmapsto[s/x]\in\mathcal C(A/a,B/b).
\]
Indeed an isomorphism $f:S\cong S'$ carrying $x$ to $x'$ restricts to an
isomorphism $S/x\cong S'/x'$ over the same residual boundary. Hence this
term is independent of the representative. Proposition~\ref{prop:egt-residual-strategy}
supplies its strategy conditions. The isomorphisms
$(S/x)/(y\setminus x)\cong S/y$ of Lemma~\ref{lem:egt-residual} are
identities on remaining events, so their iterated comparison maps satisfy
coherence. No representative strategy is selected for each $\sigma$.

The projection $H_G\to\Sigma$ retains histories actually realized by each
strategy. It need not be a product with a common history set. Complete-policy
trees additionally specify choices at unreached nodes; their subtree
construction below supplies a total continuation family over
$\Sigma\times\mathsf{Hist}$.
\end{construction}

\begin{construction}[The dependent type of residual realizations]
\label{con:fib-history-type}
For a full continuation family, as constructed below for finite
complete-policy trees, let $\mathsf{Hist}$ be its history set, and write
$G_h:A_h\to B_h$ for the residual arena,
with policy restriction
$r_h:\Sigma\to\Sigma^{G_h}$. The families $\mathbf A_h=yA_h$ and $\mathbf B_h=yB_h$ are objects
of $\mathsf L(\mathsf{Hist})$. Their ordinary value family is
\[
 \mathsf{Val}_{\mathsf{Hist}}[\mathbf A,\mathbf B]_D
 \cong\bigl(\mathcal C(A_h,B_h)\bigr)_{h\in\mathsf{Hist}}.
\]
Its context extension is exactly
$\coprod_h\mathcal C(A_h,B_h)\to\mathsf{Hist}$.
Thus the dependent residual interfaces and the existing continuation
carrier are the linear and ordinary parts of the same constructed model.
The realization section sends $(\sigma,h)$ to $(h,g_hr_h\sigma)$;
comparison terms take values in the corresponding dependent family of
squares. Write $\mathsf{NE}_{G_h}$ for the residual comparison
predicate and form the comprehension
\[
 E=\{(\sigma,h):r_h\sigma\in\mathsf{NE}_{G_h}\}
 \hookrightarrow\Sigma\times\mathsf{Hist}.
\]
For $\pi:\Sigma\times\mathsf{Hist}\to\Sigma$, universal continuation
equilibrium is the internal predicate $\forall_\pi E$.
The quantifier formula proves that it holds exactly when every residual
profile is an equilibrium. The history type contains all histories, so
the formula includes histories not reached by the reference profile.
The subsequent finite decision and concurrent-round constructions specify
these residual games and their solution predicates explicitly.
\end{construction}

\subsubsection{Operational action updates and pointed-history transport}
\label{sec:egt-operational-action-update}

The epistemic action-model construction can now obtain its
enabling data directly from event-structure steps. The updated states
are pointed transitions, so they retain the causal prerequisites and
consistency of the move as well as its visible label. Keeping the
source configuration prevents two different internal histories from
being identified merely because they expose the same event.

\begin{construction}[An action model generated by a concurrent strategy]
\label{con:egt-strategy-generated-action}
For a receptive courteous strategy $s:S\to A^\perp\parallel B$, put
$W_s=\operatorname{Conf}(S)$ and $E_s=|S|$. Use the enabling predicate
\[
 P_s(x,e)\quad\Longleftrightarrow\quad
   e\notin x\ \wedge\ x\cup\{e\}\in\operatorname{Conf}(S).
\]
Its comprehension is the enabled-event span
$\mathsf{Step}(S)$ of Construction~\ref{con:egt-enabled-event-span}.
The state map after an action is the target
\[
 t_s:\mathsf{Step}(S)\longrightarrow W_s,
 \qquad (x,e)\longmapsto x\cup\{e\}.
\]
For $x\in W_s$, write $s(x)=a_x\amalg b_x$ and form the
representable dependent boundary families
$\mathbf A_x=y(A/a_x)$ and $\mathbf B_x=y(B/b_x)$.
Define the pointed residual section
\[
 \rho_s:1_{W_s}\longrightarrow
       \mathsf{Val}_{W_s}[\mathbf A,\mathbf B]_D,
 \qquad \rho_s(x)=[s/x].
\]
Proposition~\ref{prop:egt-residual-strategy} gives the component's
strategy conditions, and pointwise Yoneda gives its displayed type.
For a history observable $v:W_s\to Z$, define its operational
postcondition by $v^+=vt_s$. No choice of total order on independent
events is involved. Action uncertainty relations will be constructed
from the actual observation data below.
\end{construction}

\begin{theorem}[Operational realization of one-step product update]
\label{thm:egt-operational-action-correspondence}
The carrier of the semantic action-model product update determined by
$P_s$ is canonically
\begin{equation}\label{eq:egt-operational-action-carrier}
 W_s\odot(E_s,P_s)\cong\mathsf{Step}(S).
\end{equation}
The residual realization over this carrier is
\begin{equation}\label{eq:egt-operational-updated-residual}
 (x,e)\longmapsto [s/(x\cup\{e\})]
 \in\mathcal C(A/a_{x,e},B/b_{x,e}),
\end{equation}
where $s(x\cup\{e\})=a_{x,e}\amalg b_{x,e}$. It is the reindexing
$t_s^*\rho_s$ of the pointed residual family. For an enabled $e$ after
$x$ there is also a canonical residual comparison
\[
 (s/x)/\{e\}\cong s/(x\cup\{e\}).
\]
All comparisons respect isomorphisms of strategy presentations.
\end{theorem}
\begin{proof}
\leavevmode
\begin{enumerate}[label=\textup{(\arabic*)},leftmargin=*,itemsep=4pt,topsep=3pt]
  \item \textbf{Equality of the two constructions.} Both sides of the carrier bijection consist of exactly the pairs
$(x,e)$ satisfying $P_s(x,e)$, and their source and action projections
are the same. The target $t_s$ is the configuration reached by that
step. Proposition~\ref{prop:egt-residual-strategy} constructs the
residual arrow at $x\cup\{e\}$; the iterated-residual isomorphism is
Lemma~\ref{lem:egt-residual}.
  \item \textbf{Secured operational interaction.} The corresponding dependent section is
Construction~\ref{con:fib-history-type} on the appropriate pointed
history base, and its reindexing is
Proposition~\ref{prop:egt-dependent-realization-transport}. Every
construction is by restriction of labels, configurations, and causal
order, so it commutes with source isomorphisms preserving marked
histories.
\end{enumerate}
\end{proof}

\noindent A witnessed event step determines the corresponding epistemic action transition.
\begin{center}
\begin{tikzcd}[column sep=large,row sep=large]
\mathsf{Step}(S)\arrow[r,"\mathrm{record}"]\arrow[d,"p"']
 &U\arrow[d,"\pi_W"]\\
 \operatorname{Conf}(S)\arrow[r,"\mathrm{state}"']&W
\end{tikzcd}
\end{center}

\begin{construction}[Action observations and their effective quotients]
\label{con:egt-operational-action-observations}
For each owner $i$, take its already designated observation on source
histories $o_i:W_s\to O_i$. If the report of the new event is a map
$b_i:E_s\to L_i$, construct the retained observation
\[
 o_i^+:\mathsf{Step}(S)\to O_i\times L_i,
 \qquad o_i^+(x,e)=(o_i(x),b_i(e)).
\]
Use $R_i=\ker(o_i)$ and $S_i=\ker(b_i)$ in the action model.
Then the product-updated relation is $\ker(o_i^+)$ by
Proposition~\ref{prop:egt-action-observation-quotient}.

A protocol may instead supply a history-dependent fresh observation
$\omega_i:\mathsf{Step}(S)\to L_i$. To retain both the preceding
information and this new observation, construct
\[
 o_i^+=\langle o_ip,\omega_i\rangle:
 \mathsf{Step}(S)\to O_i\times L_i.
\]
This too has an exact action-model presentation: use the
\emph{pointed-action set} $E_s^\bullet=\mathsf{Step}(S)$, make an action
$\alpha=(x,e)$ enabled exactly at its source $x$, and put
$S_i^\bullet=\ker(\omega_i)$ on pointed actions. Product update then has
carrier canonically isomorphic to $\mathsf{Step}(S)$, with epistemic
relation $\ker\langle o_ip,\omega_i\rangle$.
\end{construction}

\begin{proposition}[Pointed action presentation of retained observations]
\label{prop:egt-pointed-action-observations}
For the history-dependent construction above, let
\[
 P^\bullet(w,\alpha)\quad\Longleftrightarrow\quad
 w=p(\alpha),\qquad \alpha\in E_s^\bullet.
\]
Its action-model update and resulting kernel relations identify with
$\mathsf{Step}(S)$ and $\ker(o_i^+)$, respectively. State-dependent action observations are constructed from semantic precondition predicates and effective observation quotients.
\end{proposition}
\begin{proof}
\leavevmode
\begin{enumerate}[label=\textup{(\arabic*)},leftmargin=*,itemsep=4pt,topsep=3pt]
  \item \textbf{Monoidal comparison.} The comprehension of $P^\bullet$ is the graph of
$p:E_s^\bullet\to W_s$, canonically isomorphic to $E_s^\bullet$. Under that isomorphism the old-world projection is $p$, while the
new-action projection is the identity. The product-updated relation is
\[
 (p\times p)^*\ker(o_i)\wedge\ker(\omega_i)
 =\ker\langle o_ip,\omega_i\rangle.
\]
The latter kernel has its effective image factorization as in
Corollary~\ref{prop:egt-observation-interior}.
\end{enumerate}
\end{proof}

\begin{construction}[Owner-parameterized pointed action families]
\label{con:egt-owned-pointed-actions}
For an owned interaction $G=(\vec\Sigma,g)$, use the pointed-history
carrier $H_G$ of Construction~\ref{con:egt-pointed-histories}. Form the
set of classes
\[
 \mathsf{Step}(G)=
 \{[\sigma,s,x,e]:[s]=g(\sigma),\ 
       x\in\operatorname{Conf}(S),\ e\notin x,
       x\cup\{e\}\in\operatorname{Conf}(S)\},
\]
where strategy isomorphisms preserve $\sigma$, $x$, and the marked event
$e$. Its source and target maps send a class to
$[\sigma,s,x]$ and $[\sigma,s,x\cup\{e\}]$, respectively. The
post-step residual is the class of $s/(x\cup\{e\})$. For
$r:G\Rightarrow H$, the realization equation $g=h\Pi_Nr$ sends these
classes to the corresponding classes for $H$, changing only their policy
index by $\Pi_Nr$. These maps commute with both endpoints and the
residual realization. Pulling back the reference along
$W\to\Sigma^G$ retains the whole fibre of pointed operational
histories; it does not restrict the alternative policy sets.
\end{construction}

\begin{proposition}[Causal interchange of iterated action updates]
\label{prop:egt-operational-action-interchange}
Let $\mathsf{Path}(S)$ be the set of pairs consisting of a finite initial
configuration $x$ and a finite sequence $(e_1,\ldots,e_n)$ of events
successively enabled from $x$ (including $n=0$). Let $\equiv$ be
generated by
exchanging two adjacent occurrences whenever they are concurrent and
both orders yield configurations. Then endpoint assignment induces
\begin{equation}\label{eq:egt-action-trace-quotient}
 \mathsf{Path}(S)/{\equiv}
 \ \cong
 \{(x,y)\in\operatorname{Conf}(S)^2:x\subseteq y\}.
\end{equation}
Residualization after a path depends only on its class, with resulting
arrow $s/y$. This quotient removes arbitrary linear readings of
independence but does not identify epistemic records that genuinely
observe their execution order: such records remain on the path carrier
until their effective descent is established.
\end{proposition}
\begin{proof}
\leavevmode
\begin{enumerate}[label=\textup{(\arabic*)},leftmargin=*,itemsep=4pt,topsep=3pt]
  \item \textbf{Causal execution.} A sequence yields $y=x\cup\{e_1,\ldots,e_n\}$, unchanged by an
interchange. Conversely, given $x\subseteq y$, any linear extension
of the finite causal order on $y\setminus x$ provides a sequence of
enabled events, since $y$ is consistent and down-closed. Any two
linear extensions are connected by swaps of adjacent incomparable
occurrences; both swapped steps lie within $y$ and are compatible.
  \item \textbf{Coherence of iterated residuals.} This proves the bijection. Iterated residuals are canonically
isomorphic to $s/y$ by Lemma~\ref{lem:egt-residual}; their comparison
maps are identities on the remaining events and hence coherent.
\end{enumerate}
\end{proof}

\subsubsection{Secured visible actions, hiding, and epistemic witnesses}
\label{sec:egt-secured-action-update}

The comparison of communication update with concurrent composition
has to be made before hiding removes private events. Secured
synchronization \citep{CastellanEtAl2017} retains the causal witness,
while epistemic product update
\citep{BaltagMossSolecki1998PublicAnnouncements} retains the
observed event. We construct their correspondence and then identify
the additional information required to descend it to visible
executions. One visible event can depend on several hidden histories,
so forgetting those witnesses too early can change knowledge.

\begin{construction}[Visible projection and causal lifts]
\label{con:egt-visible-causal-lift}
Let $U$ be an event structure with a designated set $V\subseteq|U|$ of
visible events. Write $H=U\mathord{\downarrow}V$ for the event structure
with event set $V$, the restricted causal order, and consistency inherited
from $U$. For $z\in\operatorname{Conf}(H)$ put
\begin{equation}\label{eq:egt-visible-causal-lift}
 \widehat z=\bigcup_{v\in z}\downarrow_Uv.
\end{equation}
The visible projection of a configuration is
$\operatorname{vis}(X)=X\cap V$; its fibre is the dependent witness set
\begin{equation}\label{eq:egt-hidden-history-witness}
 K_z=\{X\in\operatorname{Conf}(U):X\cap V=z\}.
\end{equation}
\end{construction}

\begin{lemma}[Causal lifting and residualization across hiding]
\label{lem:egt-hiding-residual}
For every $z\in\operatorname{Conf}(H)$, $\widehat z$ is a
configuration of $U$, is the least such configuration with visible
projection $z$, and satisfies $\widehat z\cap V=z$. In particular
$\operatorname{vis}:\operatorname{Conf}(U)\twoheadrightarrow
\operatorname{Conf}(H)$ is a surjection with canonical section
$z\mapsto\widehat z$.
Put $V_z=V\cap|U/\widehat z|$, the visible events that remain in
the residual. There is a canonical isomorphism
\begin{equation}\label{eq:egt-hiding-residual}
 (U\mathord{\downarrow}V)/z
 \ \cong
 (U/\widehat z)\mathord{\downarrow}V_z.
\end{equation}
For $z\subseteq z'$ the residual extension corresponds to the finite
configuration $\widehat{z'}\setminus\widehat z$ of $U/\widehat z$.
\end{lemma}

\begin{proof}
\leavevmode
\begin{enumerate}[label=\textup{(\arabic*)},leftmargin=*,itemsep=4pt,topsep=3pt]
  \item \textbf{Causal execution.} Consistency of $z$ and repeated adjoining of causal predecessors show
that $\widehat z$ is consistent. It is finite because every principal
past is finite, and it is down-closed by construction. Any configuration
with visible part $z$ contains all these predecessors.
  \item \textbf{Lifting along the surjection.} Since $z$ is
down-closed in the restricted visible order, $\widehat z\cap V=z$. Hence $\widehat z$ is the least lift and the visible projection is
surjective.
  \item \textbf{Causal execution.} An event $v\in V\setminus z$ belongs to the left residual exactly when
$z\cup\{v\}$ is consistent in $U$. Closure of consistency under
adjoining causes makes this equivalent to
$\widehat z\cup\{v\}\in\operatorname{Con}_U$, which is precisely
membership in the right residual.
  \item \textbf{Causal execution: completion.} For every finite set $Z$ of these
remaining visible events, the same cause-extension argument gives
\[
 z\cup Z\in\operatorname{Con}_U
 \quad\Longleftrightarrow\quad
 \widehat z\cup Z\in\operatorname{Con}_U.
\]
Both residuals inherit the same restricted causal order, proving the
isomorphism. Monotonicity of $z\mapsto\widehat z$ and
Lemma~\ref{lem:egt-residual} give the final comparison.
\end{enumerate}
\end{proof}


\begin{proposition}[Secured interaction at a pointed history]
\label{prop:egt-pointed-secured-interaction}
Let $U=S\mathbin{\circledast_B}T$ be the secured interaction of
$s:S\to A^\perp\parallel B$ and $t:T\to B^\perp\parallel C$.
For $X\in\operatorname{Conf}(U)$, write $x$ and $y$ for its projections
to $S$ and $T$, and $b$ for their common synchronized $B$-history.
There is a canonical event-structure isomorphism, over the residual
three-interface arena,
\begin{equation}\label{eq:egt-secured-pointed-interaction}
 U/X\ \cong\ (S/x)\mathbin{\circledast_{B/b}}(T/y).
\end{equation}
It is natural under marked secured-history isomorphisms and compatible
with further residualization.
\end{proposition}

\begin{proof}
\leavevmode
\begin{enumerate}[label=\textup{(\arabic*)},leftmargin=*,itemsep=4pt,topsep=3pt]
  \item \textbf{Secured matching of interaction configurations.} By Theorem~\ref{thm:egt-pullback}, configurations of the secured
interaction are secured matchings of the augmented component maps. The configuration $X$ determines a fixed secured matching of the
component histories $x,y$, including the private interface events. A configuration extending $X$ gives a matching of the additional
residual configurations of $S/x$ and $T/y$, with the already matched
$B$-labels removed.
  \item \textbf{Reverse implication.} Conversely, the union of this fixed matching
with a matching of the residual component configurations is a matching
of the original augmented maps. The causality relation has no edges
from a new pair to an old pair, because $x,y$ are down-closed. Consequently a cycle in the union exists exactly when a cycle exists
in the residual matching; security is preserved and reflected.
  \item \textbf{Secured operational interaction.} These constructions are inverse and inclusion-preserving on the
secured-matching configuration posets. Taking their prime secured
histories identifies the corresponding event structures, including
causality, consistency, and boundary projections. The identification
uses the same fixed prefix under further residualization, so the
comparison maps compose by the identities of
Lemma~\ref{lem:egt-residual}.
\end{enumerate}
\end{proof}


\begin{construction}[Secured weak-step relation]
\label{con:egt-secured-weak-steps}
For the visible projection above construct
\[
 \begin{aligned}
 \mathsf{WStep}(U,V)=\{(X,Y):\;&X,Y\in\operatorname{Conf}(U),\\
 &X\subseteq Y,\quad
 \operatorname{vis}(Y)=\operatorname{vis}(X)\cup\{v\}
 \text{ for one }v\in V\setminus\operatorname{vis}(X)\}.
 \end{aligned}
\]
The two projections to $\operatorname{Conf}(U)$ are the source and
post-step internal history. The comparison with the visible action-update
carrier is
\begin{equation}\label{eq:egt-secured-weak-comparison}
 \Theta:\mathsf{WStep}(U,V)\to\mathsf{Step}(H),\qquad
 (X,Y)\longmapsto(\operatorname{vis}(X),v),
\end{equation}
where $v$ is the unique newly visible event. Each fibre $\Theta^{-1}(z,v)$
is a dependent set of secured witnesses realizing the visible action.
\end{construction}

\begin{theorem}[Secured interaction, epistemic update, and hiding]
\label{thm:egt-secured-del-comparison}
Let $U=S\mathbin{\circledast_B}T$ be the secured interaction of two
concurrent strategies and let $V$ consist of their exterior visible
events. The source event structure of their composite is
$H=U\mathord{\downarrow}V$. The action-model carrier generated from
$H$ is $\mathsf{Step}(H)$. The comparison
$\Theta:\mathsf{WStep}(U,V)\to\mathsf{Step}(H)$ is surjective, with
canonical section
\begin{equation}\label{eq:egt-secured-weak-section}
 (z,v)\longmapsto
 (\widehat z,\widehat{z\cup\{v\}}).
\end{equation}
For every visible step $(z,v)$, residualizing after the visible update
is canonically isomorphic to the visible part of the secured residual:
\begin{equation}\label{eq:egt-secured-del-residual}
 H/(z\cup\{v\})
 \ \cong
 \bigl(U/\widehat{z\cup\{v\}}\bigr)
 \mathord{\downarrow}\bigl(V\cap|U/\widehat{z\cup\{v\}}|\bigr).
\end{equation}
The visible residual strategy has the same external labels on both
sides. If $x,y$ are the component projections of
$\widehat{z\cup\{v\}}$, then by
Proposition~\ref{prop:egt-pointed-secured-interaction} this residual
strategy is also the concurrent composite of the residual component
strategies $s/x$ and $t/y$, with the common synchronized $B$-past
removed. The comparison and the residual isomorphisms are invariant
under arena-preserving isomorphisms of the two component strategies.
\end{theorem}

\begin{proof}
\leavevmode
\begin{enumerate}[label=\textup{(\arabic*)},leftmargin=*,itemsep=4pt,topsep=3pt]
  \item \textbf{Causal execution.} Every $(X,Y)$ in the weak-step carrier projects to a visible source
configuration followed by exactly one enabled visible event, so $\Theta$
is well-defined.
  \item \textbf{Reverse implication.} Conversely, Lemma~\ref{lem:egt-hiding-residual}
constructs the two configurations in
Equation~\eqref{eq:egt-secured-weak-section}; the first is included in
the second, and their visible parts differ by precisely $v$. This is a
section and proves surjectivity. Apply
Equation~\eqref{eq:egt-hiding-residual} at $z\cup\{v\}$ to obtain
Equation~\eqref{eq:egt-secured-del-residual}.
  \item \textbf{Secured operational interaction.} The hidden composite's
source is the visible restriction of $U$, and its label map is the
restriction to the exterior arena. Therefore the identity on visible
remaining events also preserves their arena labels. Proposition~\ref{prop:egt-pointed-secured-interaction} identifies the
secured residual at the canonical lift with interaction of the
component residual strategies.
  \item \textbf{Secured operational interaction: completion.} Hiding those matched internal events
gives precisely the displayed exterior residual. The causal-closure
construction commutes with isomorphisms of the secured interaction and
with the canonical source comparisons of its pullback presentation.
\end{enumerate}
\end{proof}


\begin{remark}[Internal histories are not erased by epistemic hiding]
\label{rem:egt-hidden-epistemic-witnesses}
The configuration projection
$\operatorname{vis}:\operatorname{Conf}(U)\to\operatorname{Conf}(H)$
preserves the secured weak-step relation of
Construction~\ref{con:egt-secured-weak-steps} and the visible enabled-step
relation, but it need not be bounded: a visible step need not lift from
each internal history over its source. The surjection $\Theta$ compares
the corresponding step carriers; its surjectivity asserts only that every
visible step has some witness. An internal history may include a hidden
occurrence that disables a visible
move, even though that move remains available at its coarse visible
projection. For example, take a hidden event $h$ conflicting with a
visible event $v$, both initially enabled. At internal history $\{h\}$,
the event $v$ is impossible, while the empty visible history still permits
$v$. The minimal section in
Equation~\eqref{eq:egt-secured-weak-section} and the complete witness
fibres $K_z$ are therefore both essential. We do not identify arbitrary
private observations of $K_z$ with predicates on visible histories.
\end{remark}

\begin{proposition}[Exact epistemic descent of secured weak steps]
\label{prop:egt-secured-action-descent}
For an owner $i$, let
$\omega_i:\mathsf{WStep}(U,V)\to O_i$ be a specified observation of
secured weak steps. The following are equivalent:
\begin{enumerate}[label=\textup{(\roman*)},leftmargin=*]
\item $\omega_i$ is constant on the fibres of $\Theta$;
\item its two pullbacks to the kernel pair of $\Theta$ agree;
\item there is a unique $\bar\omega_i:\mathsf{Step}(H)\to O_i$
with $\omega_i=\bar\omega_i\Theta$.
\end{enumerate}
Under these equivalent descriptions, $\Theta$ is a surjective bounded
morphism from $\ker(\omega_i)$ to $\ker(\bar\omega_i)$. For a finite
family of owners satisfying the same descent test, it preserves and
reflects all modal fixed-point predicates with descended atomic tests.
Observations that do not descend remain on the dependent witness carrier
$\Theta^{-1}(z,v)$ and have their own epistemic update there.
\end{proposition}
\begin{proof}
\leavevmode
\begin{enumerate}[label=\textup{(\arabic*)},leftmargin=*,itemsep=4pt,topsep=3pt]
  \item \textbf{Descent along the kernel pair.} The equivalence is effective descent along the surjection $\Theta$,
using Proposition~\ref{prop:egt-kernel-descent}. Since
$\ker(\omega_i)=(\Theta\times\Theta)^*\ker(\bar\omega_i)$,
forward preservation of the kernel relation follows immediately. For a reduced successor $z'$ of $\Theta(X,Y)$ with the same
$\bar\omega_i$-value, surjectivity of $\Theta$ supplies any preimage
of $z'$.
  \item \textbf{Descent along the kernel pair: completion.} Its $\omega_i$-value equals that of $(X,Y)$, so it supplies
the bounded-morphism back lift. Theorem~\ref{thm:egt-modal-adequacy}
then gives the modal fixed-point assertion. For non-descending
observations, the witness family is the ordinary fibre of $\Theta$;
no image-level observation term is falsely asserted.
\end{enumerate}
\end{proof}

The results give two complementary readings. The visible action
model is constructed from the hidden composite strategy $H$ through
its least causal lift; the internal action model retains every secured
history and its private observations. Their comparison identifies
which information descends to visible histories and which must remain
with the witnesses. Hiding alone is not an epistemic quotient.

\subsubsection{Canonical exact refinements of visible information}
\label{sec:egt-exact-visible-refinement}

The preceding descent result identifies the observations that the
visible quotient can represent. Private observations may nevertheless
distinguish histories with the same visible projection. We first
retain those distinctions by an image refinement. To preserve modal
and temporal reasoning as well, the refinement must also lift the
specified relations in both directions; a relation-sensitive quotient
will supply that exactness.

\begin{proposition}[The coarsest observation-retaining refinement]
\label{prop:egt-observation-retaining-refinement}
Let $q:X\twoheadrightarrow Y$ be a surjection and $o_i:X\to O_i$
specified observations. The image factorization
\[
 X\xrightarrow{r_0}Z_0:=\operatorname{im}\langle q,(o_i)_{i\in N}\rangle
 \lhook\joinrel\longrightarrow Y\times\prod_iO_i
\]
is the coarsest quotient through which $q$ and every $o_i$ factor.
Moreover $r_0$ is a surjective bounded morphism from each $\ker(o_i)$
to $\ker(\bar o_i)$ on $Z_0$.
\end{proposition}
\begin{proof}
\leavevmode
\begin{enumerate}[label=\textup{(\arabic*)},leftmargin=*,itemsep=4pt,topsep=3pt]
  \item \textbf{Descent along the kernel pair.} Its kernel pair is $\ker(q)\cap\bigcap_i\ker(o_i)$. If a quotient
$s:X\twoheadrightarrow Z$ retains all these maps, then
$\ker(s)\subseteq\ker(r_0)$, giving the unique factor
$h:Z\twoheadrightarrow Z_0$ with $r_0=hs$. Forward preservation of observation kernels is immediate.
  \item \textbf{Lifting along the surjection.} For a back lift
from $r_0(x)$ to $z$ with the same $i$-observation, choose a preimage
$x'$ of $z$. Then $o_i(x')=o_i(x)$, as required.
\end{enumerate}
\end{proof}

The image quotient alone can still fail to lift visible
transitions. To retain both observations and transitions, take a set
$\mathcal A$ of relation labels with $R_a\subseteq X^2$ for
$a\in\mathcal A$, together with the specified observation and test
maps $t_\lambda:X\to V_\lambda$. The greatest-bisimulation
construction identifies histories only when they are indistinguishable
by all of these designated interfaces.
\[
 \ell=\langle q,(o_i)_i,(t_\lambda)_\lambda\rangle:X\to L,
 \qquad L=Y\times\prod_iO_i\times\prod_\lambda V_\lambda.
\]
An atomic predicate can be retained by its characteristic map to $2$.

\begin{construction}[The bisimulation predicate]
\label{con:egt-retained-bisimulation}
For $B\subseteq X^2$ define $\mathcal F(B)$ by
\begin{align*}
 (x,x')\in\mathcal F(B)\quad\Longleftrightarrow\quad
 &\ell(x)=\ell(x')\ \text{ and, for each }a\in\mathcal A,\\
 &\forall z\;(xR_az\Rightarrow
       \exists z'\;(x'R_az'\wedge zBz')),\\
 &\forall z'\;(x'R_az'\Rightarrow
       \exists z\;(xR_az\wedge zBz')).
\end{align*}
The last two lines are simultaneous clauses. The map $\mathcal F$ is
monotone. Define $B_*:=\nu\mathcal F$ in the existing predicate lattice
$\mathcal P(X^2)$.
\end{construction}

\begin{definition}[Exact quotient refinements over retained data]
\label{def:egt-exact-quotient-category}
Let $\mathbf Q_\ell(X)$ have objects consisting of a relation frame
$(Z,(\bar R_a))$, a surjective bounded map $s:X\to Z$, and a map
$\bar\ell:Z\to L$ with $\ell=\bar\ell s$. An arrow from
$(s_1,\bar\ell_1)$ to $(s_2,\bar\ell_2)$ is a bounded map
$h:Z_1\to Z_2$ satisfying $s_2=hs_1$. It is automatically
surjective, and $\bar\ell_1=\bar\ell_2h$ follows by cancellation
of the surjection $s_1$. Identities and compositions are inherited from
frame maps; successive back lifts prove boundedness of composition.
Thus the category retains the visible map, observations, and tests
as well as the selected relations.
\end{definition}

\begin{theorem}[Coarsest exact quotient refining the visible quotient]
\label{thm:egt-exact-visible-refinement}
The relation $B_*$ is an equivalence relation. Let $r:X\twoheadrightarrow
Z=X/B_*$ and give $Z$ the direct-image relations
\[
 u\bar R_av\quad\Longleftrightarrow\quad
 \exists x,x'\;(r(x)=u,\ r(x')=v,\ xR_ax').
\]
Then:
\begin{enumerate}
\item $q$, all observations, and all retained tests factor uniquely through $r$.
\item $r$ is a surjective bounded morphism for every $R_a$.
\item For every other surjective bounded morphism $s:X\to Z'$ through
which $\ell$ factors, there is a unique surjective bounded morphism
$h:Z'\to Z$ satisfying $r=hs$.
\end{enumerate}
Thus $(r,\bar\ell)$ is terminal in $\mathbf Q_\ell(X)$.
It preserves and reflects the modal fixed-point predicates built
from the retained atoms.
\end{theorem}

\noindent The retained visible observation factors through the greatest-bisimulation quotient.
\begin{center}
\begin{tikzcd}[column sep=large,row sep=large]
X\arrow[r,"r",two heads]\arrow[dr,"q"']
 &Z\arrow[d,"\bar q"]\\
 &Y
\end{tikzcd}
\end{center}


\begin{proof}
\leavevmode
\begin{enumerate}[label=\textup{(\arabic*)},leftmargin=*,itemsep=4pt,topsep=3pt]
  \item \textbf{Coinduction and bisimulation.} A relation $B$ with $B\subseteq\mathcal F(B)$ is a label-preserving
bisimulation. The diagonal is such a relation. The converse of such a
relation is again such a relation.
  \item \textbf{Witness and legality comparison.} The relational composite of two such
relations is also one: to match a successor, first use the witness for the
first relation and then the witness for the second; the reverse direction
is identical, and label equalities compose. Greatestness therefore gives
\[
 \Delta_X\subseteq B_*,\qquad B_*^{\mathrm{op}}\subseteq B_*,
 \qquad B_*;B_*\subseteq B_*.
\]
These are reflexivity, symmetry, and transitivity. Since
$B_*\subseteq\ker(\ell)$, effective quotient gives the factorizations.
  \item \textbf{Lifting through the comparison.} Forward preservation of relations by $r$ is built into $\bar R_a$. For back lifting suppose $r(x)\bar R_av$. By definition choose
$yR_az$ with $r(y)=r(x)$ and $r(z)=v$.
  \item \textbf{Coinduction and bisimulation.} As $xB_*y$, the bisimulation
clause gives $xR_az'$ with $z'B_*z$. Then $r(z')=v$.
  \item \textbf{Universal property of the bounded quotient.} For the universal property, $\ker(s)$ preserves labels. We construct
$h$ from the quotient universal property and verify both bounded clauses,
including the existence of a set map and both bounded-morphism clauses. If $s(x)=s(y)$
and $xR_az$, forward preservation gives
$s(y)R_a^{Z'}s(z)$; back lifting along $s$ gives a matching
$yR_az'$ with $s(z')=s(z)$.
  \item \textbf{Lifting along the surjection.} Interchanging $x,y$ gives the other
clause. Thus $\ker(s)$ is a label-preserving bisimulation and is
contained in $B_*$. Quotient universality gives the unique surjective
map $h:Z'\twoheadrightarrow Z$ with $hs=r$.
  \item \textbf{Lifting along the surjection: second comparison.} Its compatibility with the
retained labels is forced by the equality
$\bar\ell_s s=\ell=\bar\ell_r r$ and surjectivity of $s$.
  \item \textbf{Lifting through the comparison.} To see that $h$ preserves a relation $uR_a^{Z'}v$, take $s(x)=u$ and
lift to $xR_ax'$ with $s(x')=v$. Forward preservation by $r$ gives
$h(u)\bar R_ah(v)$. For the back condition, if $h(u)\bar R_av$,
take $s(x)=u$ and use boundedness of $r$ to obtain $xR_ax'$ with
$r(x')=v$.
  \item \textbf{Choice of representatives and quotients.} Put $u'=s(x')$. Then $uR_a^{Z'}u'$ and $h(u')=v$. Modal fixed-point preservation and reflection are exactly Theorem~\ref{thm:egt-modal-adequacy}, since $r$ is surjective and the atoms
descend.
\end{enumerate}
\end{proof}


\begin{corollary}[Compatibility with knowledge and action update]
\label{cor:egt-exact-refinement-action}
If $R_i=\ker(o_i)$ is among the retained relations, then
$\bar R_i=\ker(\bar o_i)$. If the only relations are these observation
kernels, then $B_*=\ker(\ell)$, so the theorem reduces to the image
construction, with any additional tests included in that image.

For an action set $E$ with action relations $(S_a)_{a\in\mathcal A}$,
an arbitrary enabling subobject $P\subseteq X\times E$ is retained
by including among the tests the function
\[
 t_P:X\to\mathcal P(E),\qquad t_P(x)=\{e:(x,e)\in P\}.
\]
Its factor $\bar t_P:Z\to\mathcal P(E)$ defines
$\bar P(z,e)\Longleftrightarrow e\in\bar t_P(z)$, and
$P=(r\times1_E)^*\bar P$. The resulting map on product updates
$\{P\}\to\{\bar P\}$, $(x,e)\mapsto(r(x),e)$, is again
surjective and bounded for the updated relations. Hence modal
fixed-point adequacy continues through this update and through finite
successions for which the enabling data are retained at each stage.
\end{corollary}
\begin{proof}
\leavevmode
\begin{enumerate}[label=\textup{(\arabic*)},leftmargin=*,itemsep=4pt,topsep=3pt]
  \item \textbf{Observation and knowledge.} Equality of quotient observations is equality of the original observations
of any representatives, which is exactly $R_i$; this proves the first claim.
  \item \textbf{Coinduction and bisimulation.} For the second, if $\ell(x)=\ell(x')$, every $R_i$-successor of $x$
is also an $R_i$-successor of $x'$. It matches itself under
$\ker(\ell)$. Thus this entire kernel is a bisimulation.
  \item \textbf{Lifting along the surjection.} The added test makes $P$ constant on $r$-fibres for each event. Therefore
the updated map is well-defined and surjective. A target successor
$(v,f)$ of $(r(x),e)$ gives $r(x)\bar R_av$, $eS_af$, and enabled
$(v,f)$.
  \item \textbf{Lifting along the surjection: completion.} Boundedness supplies $xR_ax'$ with $r(x')=v$; descent of
enablement makes $(x',f)$ enabled. This is the required back lift. Forward preservation is immediate. Iterate this argument for the last claim.
\end{enumerate}
\end{proof}

The greatest-bisimulation construction is related to coalgebraic
behavioural equivalence \citep{Rutten2000UniversalCoalgebra}.
Its present role is more specific: it preserves a chosen visible map,
private observations, and the designated epistemic and secured
transition relations simultaneously. The resulting refinement is
therefore exact for the stated tests, without declaring every internal
piece of the residual arena observable.

For the example in Remark~\ref{rem:egt-hidden-epistemic-witnesses}, take
$X=\operatorname{Conf}(U)$ and $q=\operatorname{vis}$, with the
secured weak-step relation and the chosen history observations. The empty
history and the hidden history $\{h\}$ have the same visible projection,
but only the former can match the visible $v$-step. They cannot be
$B_*$-equivalent. The refined carrier therefore retains exactly the
distinction that the original visible projection loses. The map from the
refined quotient to visible configurations still preserves the visible-step
relation: each quotient edge is represented by a secured weak step, whose
visible image is an enabled step. For the map
$\Theta$ between step carriers, apply the theorem on those carriers with
their explicitly specified relations; do not identify it with the
configuration projection.

\begin{corollary}[Comparison of retained data and invariance under isomorphism]
\label{cor:egt-exact-refinement-comparison}
Adding retained tests to $\ell$ gives a canonical surjective bounded
map from the more discriminating exact quotient to the previous one.
An isomorphism of the source relation frames that intertwines all retained
data induces a unique isomorphism of their exact quotients. These comparison
maps compose coherently.
\end{corollary}
\begin{proof}
\leavevmode
\begin{enumerate}[label=\textup{(\arabic*)},leftmargin=*,itemsep=4pt,topsep=3pt]
  \item \textbf{Coinduction and bisimulation.} The quotient retaining more tests is itself an exact quotient retaining
the original tests. The terminal property of
Theorem~\ref{thm:egt-exact-visible-refinement} supplies the comparison. An intertwining isomorphism transports label-preserving bisimulations in
both directions, hence transports their greatest relation exactly.
  \item \textbf{Lifting along the surjection.} It therefore induces inverse quotient maps. All composition and identity
equations follow from uniqueness after precomposition with the relevant
surjective quotient map.
\end{enumerate}
\end{proof}

\begin{construction}[Dependent operational witnesses over the refined quotient]
\label{con:egt-refined-history-witnesses}
For $r:X\twoheadrightarrow Z=X/B_*$, retain the ordinary witness
family $K_z=r^{-1}(z)$. The comparison
\[
 X\xrightarrow{\cong}Z.K,\qquad x\longmapsto(r(x),x),
\]
has inverse $(z,x)\mapsto x$. Thus every ordinary or linear family over
$X$, including the pointed residual arenas and their realized strategies,
is transported exactly to the total witness context $Z.K$.
The refined state quotient carries the preserved modal predicates; its
dependent witnesses retain the individual operational residuals.
This is an isomorphism of contexts, so all reindexing and realization
comparisons follow from
Proposition~\ref{prop:egt-dependent-realization-transport}.
\end{construction}

The quotient concerns the chosen visible and modal tests, not an
identification of residual linear arenas. Complete linear data remain
available on the underlying witness family; Proposition~\ref{prop:egt-linear-descent}
then specifies when an interface itself descends. This separates exact
logical abstraction from isomorphism of all operational structure.

\subsubsection{Finite perfect-information decisions}

\begin{definition}[Finite perfect-information tree and complete policies]
\label{def:egt-finite-perfect-information-tree-and-complete-policies}
A \textbf{finite perfect-information tree} is a finite prefix-closed
set $H$ of histories with root $\varepsilon$, a set
$H_{\mathrm{nt}}\subseteq H$ of nonterminal histories, and leaf set
$Z=H\setminus H_{\mathrm{nt}}$. Each $h\in H_{\mathrm{nt}}$ has an
owner $i(h)\in N$ and a finite nonempty successor set $A(h)$, with
the extension $ha\in H$ for each $a\in A(h)$ giving precisely the
children of $h$. For each $i$ a leaf payoff has type
$u_i:Z\to\mathbb R$.
A \textbf{complete policy} of owner $i$ is an element of
\[
 \Sigma_i=\prod_{h\in H_{\mathrm{nt}}:\,i(h)=i}A(h),
 \qquad \Sigma=\prod_i\Sigma_i,
\]
including its actions at nodes not on the realized path. Following the
prescribed successor from the root defines
$\operatorname{play}:\Sigma\to Z$. The \textbf{subtree} $T_h$ is
rooted at a history $h$ and has as nodes the extensions
$H_h=\{k\in H:h\preceq k\}$, with inherited ownership, action sets,
and leaf payoffs; its policy restrictions are the coordinate projections
$r_h:\Sigma\to\Sigma^h$ constructed below. An agent observes its
current history and therefore knows its decision node; perfect
information says nothing about shared preferences.
\end{definition}
\begin{construction}[Retained-history realization]
There is a fixed-polarity event-structure realization of the complete
history. Let $\mathsf{Hist}(T)$ have the nonempty tree histories as events,
ordered by prefix extension, with finite consistent sets precisely the
chains, and give every event positive polarity. For a profile $\sigma$,
let $P_\sigma$ be the chain of nonempty prefixes of
$\operatorname{play}(\sigma)$ and let $S_\sigma$ be the induced chain event
structure. Inclusion gives
\[
 t_\sigma:S_\sigma\longrightarrow\mathsf{Hist}(T),
\]
regarded as a strategy $I\to\mathsf{Hist}(T)$. Receptivity is vacuous
because the arena has no negative events. Every immediate dependency in
$S_\sigma$ is the corresponding prefix dependency in the arena, so
courtesy holds. Hence
\[
 \sigma\longmapsto t_\sigma
\]
is an owned family on one common polarized arena; no event changes
polarity with the identity of the agent who owns its decision.
\end{construction}

\begin{construction}[Leaf observation, including the terminal root]
If the root is a leaf $z$, then $\mathsf{Hist}(T)=I$, its unique
complete-policy profile has $t_\sigma=\operatorname{id}_I$, and put
$r=\mathsf{pt}(z):I\to D(\{z\})$.
For a nonterminal root every leaf is a nonempty history. Define $r:\mathsf{Hist}(T)\to D(Z)$ to have one negative input event $\bar h$
for every history $h$ and one positive output $o_z$ for every leaf $z$.
The negative events inherit the tree order. Add the dependency
$\bar z<o_z$. A finite set is consistent precisely when its negative
history events form a chain, it contains at most one output, and, if it
contains $o_z$, every negative history in the set is a prefix of $z$.
These clauses make the source an event structure: adjoining a cause of
$o_z$ stays on the same prefix chain, while inherited causes of negative
history events remain on that chain. The labelling sends $\bar h$ to the
input occurrence of $h$ and $o_z$ to the output value $z$. A source
configuration containing $o_z$ contains $\bar z$ and all its prefixes, so
its image is a configuration of
$\mathsf{Hist}(T)^\perp\parallel D(Z)$; local injectivity is immediate.
Thus the labelling is a total map. Every enabled negative history
extension is accepted exactly once. The only dependency not already
present on the input arena is $\bar z<o_z$, which is
negative-to-positive. Thus $r$ is receptive and courteous. Synchronizing
$r$ with $t_\sigma$ exposes exactly the prefix chain $P_\sigma$; only the
leaf $\operatorname{play}(\sigma)$ occurs on that chain. Hiding the
history events therefore gives
\[
 r\circ t_\sigma
 \cong
 \mathsf{pt}_Z(\operatorname{play}(\sigma)):I\to D(Z).
\]
Thus the retained-history and outcome-only presentations are related by an
explicit strategy with the interface polarities fixed.
The staged public-state circuit constructed below gives a decentralized
factorization of the same play function. When specialized to a
perfect-information tree, every nonowner has a singleton action at the
current node. To compare its flat outcome interface with the prefix-ordered
history arena, we now construct an explicit strategy. Merely copying public
state values into independent data ports would not itself impose the
history arena's prefix order.

\end{construction}

\begin{construction}[Expansion of a leaf into its prefix history]
\label{con:egt-history-expansion}
Let $Z$ be the leaves of the finite tree $T$. For $z\in Z$, write $P_z$
for its chain of nonempty prefixes and let
$t_z:I\to\mathsf{Hist}(T)$ be inclusion of that chain.
Construct
\[
 e_T:D(Z)\longrightarrow\mathsf{Hist}(T)
\]
as follows. For each $z$ its source has a negative event $n_z$ and a
positive event $p_{z,h}$ for every $h\in P_z$. Within this branch,
$n_z$ precedes every positive event and the positive events are ordered
by prefix. A finite set is consistent exactly when it is contained in
one branch; events in distinct branches conflict. The labelling sends
$n_z$ to the input value $z$ and $p_{z,h}$ to the history event $h$.
If the root is terminal, $P_z$ is empty, so the unique branch has only
its negative input and $e_T:D(\{z\})\to I$ is disposal.
\end{construction}

\begin{proposition}[Comparison with the retained-history realization]
\label{prop:egt-history-expansion}
The displayed map $e_T$ is a receptive courteous strategy. For every leaf
$z$ and complete policy profile $\sigma$,
\[
 e_T\circ\mathsf{pt}_Z(z)\cong t_z,
 \qquad
 e_T\circ\mathsf{pt}_Z(\operatorname{play}(\sigma))
       \cong t_\sigma.
\]
These comparisons include the terminal-root case.
\end{proposition}

\begin{proof}
\leavevmode
\begin{enumerate}[label=\textup{(\arabic*)},leftmargin=*,itemsep=4pt,topsep=3pt]
  \item \textbf{Causal execution.} The source is an event structure: each branch is a finite chain and its
causes remain in that branch, so the consistency axioms hold. A nonempty configuration contains one $n_z$ and an initial segment of
the prefix chain $P_z$. Its image is consequently a configuration of
$D(Z)^\perp\parallel\mathsf{Hist}(T)$, and the labels are injective on
that configuration.
  \item \textbf{Local injectivity across conflicting branches.} Repeated history labels on different branches cause
no failure of local injectivity because those branches conflict.
  \item \textbf{Fibrewise calculation.} At the empty configuration every negative input $z$ has the unique lift
$n_z$. Once an input has occurred, the other negative inputs conflict
with it already in $D(Z)^\perp$, so there are no further enabled negative
extensions.
  \item \textbf{Receptivity and courtesy of prefixes.} This proves receptivity. The only new immediate dependency
is from $n_z$ to the first positive prefix event. All positive-to-positive
immediate dependencies are the corresponding prefix edges of
$\mathsf{Hist}(T)$. Hence courtesy holds.
  \item \textbf{Causal execution.} Synchronizing with $\mathsf{pt}_Z(z)$ selects the $z$-branch and
hides its input event. The remaining labelled source is exactly the
chain $P_z$, giving the first isomorphism. Substituting the leaf
$\operatorname{play}(\sigma)$ gives the second.
  \item \textbf{Causal execution: completion.} For a terminal root there
are no remaining events, and both sides are the unique arrow $I\to I$.
\end{enumerate}
\end{proof}


The comparison recovers prefix-ordered records from terminal
history values, giving equality for the constructed closed strategies.
It does not identify arbitrary strategies on a flat data arena with
arbitrary strategies on a history arena. The causal record retained by
the protocol is what makes the reconstruction possible.

\subsubsection{Backward induction, with all ties retained}

\begin{definition}[Backward induction and subgame-perfect equilibrium]
\label{def:egt-backward-induction-spe}
For a finite perfect-information tree, write $\Sigma^h$ for the
complete profile carrier of its subtree $T_h$, and define the
\textbf{backward-induction solution family}
$\mathsf{BI}(h)\hookrightarrow\Sigma^h$ recursively.
At a leaf $z$, $\mathsf{BI}(z)=\{*\}=\Sigma^z$.
At a nonterminal history $h$ owned by $i(h)$, use the canonical
policy-coordinate identification
\[
 \Sigma^h\cong A(h)\times\prod_{b\in A(h)}\Sigma^{hb}
\]
and put
\[
 \begin{split}
 \mathsf{BI}(h)=\{(a,(s^{hb})_{b\in A(h)}):\;&
       s^{hb}\in\mathsf{BI}(hb)\ \text{for every }b,\\
 &u_{i(h)}(\operatorname{play}_{ha}(s^{ha}))
      =\max_{b\in A(h)}
        u_{i(h)}(\operatorname{play}_{hb}(s^{hb}))\}.
 \end{split}
\]
Thus every maximizing action and every admissible choice of its
continuation profiles is retained, including off-path continuations.
A profile is a \textbf{subgame-perfect equilibrium} (SPE) when its
restriction is Nash in every subtree. The independent internal
predicate $\mathsf{SPE}=\bigwedge_h r_h^*\mathsf{NE}_{G_h}$ is typed
in Definition~\ref{def:egt-subgame-perfection};
Theorem~\ref{thm:egt-backward} proves
$\mathsf{BI}(\varepsilon)=\mathsf{SPE}$ for these finite trees.
\end{definition}
This rules out a prescribed future action that ceases to be optimal when
its decision is actually reached.

These backward-induction characterizations concern finite trees.
The choice of a maximizing child is justified in the ordinary
classical specialization with real-valued payoffs in $\mathbb R$.
The categorical construction supplies all finite continuation choices;
ordinary maximization selects an optimal continuation while retaining
all ties in the resulting policy set.

\begin{construction}[Subtree policies, restriction, and reference-based extension]
For each history $h$, let $T_h$ be the rooted subtree whose nodes are the
histories extending $h$. Construct its ownerwise policy spaces by
\[
 \Sigma_i^h
 =\prod_{\substack{k\succeq h\\ i(k)=i}}A(k),
\]
where the product ranges over the nonterminal nodes of $T_h$ owned by $i$.
The canonical product projections give restriction maps
\[
 r_{h,i}:\Sigma_i\longrightarrow\Sigma_i^h,
 \qquad
 r_h=\prod_i r_{h,i}:\Sigma\longrightarrow\Sigma^h.
\]
Applying the retained-history and outcome construction above to $T_h$ gives
an owned continuation game $G_h$ on the parameter family
$\vec\Sigma^{\,h}$. Thus ``the game after $h$'' is a constructed subtree
game, with continuation data carried by its dependent policy family.

Every residual alternative has an explicit extension to a complete policy
relative to a reference profile. For $\sigma_i\in\Sigma_i$, define
\[
 \operatorname{ext}_{h,\sigma_i}:\Sigma_i^h\longrightarrow\Sigma_i
\]
coordinatewise by
\[
 \operatorname{ext}_{h,\sigma_i}(\tau_i)(k)=
 \begin{cases}
  \tau_i(k),&k\succeq h,\\
  \sigma_i(k),&k\not\succeq h.
 \end{cases}
\]
Then
\[
 r_{h,i}\operatorname{ext}_{h,\sigma_i}
 =\operatorname{id}_{\Sigma_i^h}.
\]
No global section of $r_{h,i}$ is selected once and for all; the displayed
extension is constructed from the reference policy whose off-subtree
coordinates are to be retained.

When $h$ lies on the realized history of a profile $\sigma$, this subtree
construction agrees with the operational residual. Regard $T_h$ as rooted at
$h$; the nonempty histories of $T_h$ are the nonempty suffixes after $h$ and
therefore correspond to the strict extensions of $h$ in $T$. If $P_h$ is the
prefix configuration containing all nonempty prefixes through $h$, deletion
of the executed prefix gives canonical isomorphisms
\[
 \mathsf{Hist}(T)/P_h\cong\mathsf{Hist}(T_h),
 \qquad
 t_\sigma/P_h\cong t^h_{r_h(\sigma)}.
\]
The first sends a strict extension of $h$ to its relative suffix in the rooted
subtree, and the second restricts that map to the realized continuation chain;
order, chain consistency, polarities, and labels are inherited. Off-path
subgames are still defined by the subtree construction even though they are
not residuals of the particular realized strategy.
\end{construction}

\begin{definition}[Subgame perfection]
\label{def:egt-subgame-perfection}
A profile is \textbf{subgame perfect} exactly when every one of these reindexed
Nash predicates holds. In the assertion doctrine,
\begin{equation}
\label{eq:egt-spe-doctrine}
 \mathsf{SPE}
 =\bigwedge_h r_h^*\mathsf{NE}_{G_h}.
\end{equation}
Subgame perfection is the conjunction of equilibrium predicates over all constructed continuation games. The backward-induction theorems below
determine this predicate for finite causal trees.
\end{definition}

\begin{theorem}
\label{thm:egt-backward}
For every finite perfect-information tree, all pure subgame-perfect
profiles are generated by the following recursion. At a leaf there is
the unique empty policy profile. At a node owned by $i$, choose a
subgame-perfect profile in every child subtree, and choose any child
whose resulting payoff to $i$ is maximal. Retain the chosen profiles in
all child subtrees, including those not selected at the node.
In particular a pure subgame-perfect profile exists.
\end{theorem}

\begin{proof}
\leavevmode
\begin{enumerate}[label=\textup{(\arabic*)},leftmargin=*,itemsep=4pt,topsep=3pt]
  \item \textbf{Fixed-point argument.} Proceed by induction on tree height. The leaf case is immediate. At a nonterminal root, let the selected child be $a$ and the chosen
subgame-perfect continuation in child $b$ be $\sigma^b$.
  \item \textbf{Strategic comparison.} If a deviating agent $j$ is not the root's owner $i$, the selected child
does not change. The deviation cannot increase its payoff within child
$a$, because $\sigma^a$ is a Nash equilibrium there. If the deviator is
$i$, it can choose another child $b$ and also change its entire policy
inside that child.
  \item \textbf{Strategic comparison: second comparison.} The Nash property of $\sigma^b$ bounds the resulting
payoff by $u_i(\operatorname{play}(\sigma^b))$, and the maximizing choice
of $a$ bounds this by the current payoff. Thus no agent has a profitable
complete-policy deviation at the root. All proper subgames already have
the required property by induction.
  \item \textbf{Reverse implication.} Conversely, a subgame-perfect profile restricts to a subgame-perfect
profile in every child. Its chosen child must maximize the root owner's
payoff among those fixed continuations: otherwise changing only the
root decision is profitable. This is exactly one of the profiles in
the stated recursion.
  \item \textbf{Finite maxima and continuation witnesses.} Finally each finite nonempty set of child payoffs has a maximum, and the
finitely many continuation witnesses are assembled by induction over the
child set. No choice over an infinite family is involved. This proves
existence.
\end{enumerate}
\end{proof}


The later node may belong to the same owner as an earlier one.
Accordingly, optimality is tested against changes to that owner's
complete contingent policy rather than one local action at a time.
The recursion therefore avoids the componentwise-equilibrium error
identified in Section~\ref{sec:egt-equilibrium}. The construction
relates to selection-function accounts of backward induction
\citep{EscardoOliva2010}, but preserves the explicitly concurrent
structure of completed stages.

\begin{remark}[Finite and infinite continuation]
The continuation arguments here use finite trees and their residual games.
For a coalgebraic treatment of subgame-perfect equilibria in infinite games
without discounting, see \citet{AbramskyWinschel2013}. Such an infinite-game
analysis is separate from the finite backward decompositions proved here.
\end{remark}

\subsubsection{Subgame perfection with concurrent public rounds}

\begin{definition}[Finite public game tree with simultaneous rounds]
\label{def:egt-finite-public-game-tree-with-simultaneous-rounds}
A \textbf{finite public game tree} has a finite prefix-closed history
set $H$, its nonterminal subset $H_{\mathrm{nt}}$, leaves
$Z=H\setminus H_{\mathrm{nt}}$, and, at every
$h\in H_{\mathrm{nt}}$, finite nonempty action sets $A_i(h)$ for all
agents. Its dependent joint-action carrier and successor map are
\[
 \mathsf{Act}=
 \coprod_{h\in H_{\mathrm{nt}}}\prod_{i\in N}A_i(h)
 \xrightarrow{\ \operatorname{succ}\ }H,\qquad
 \operatorname{succ}(h,a)=ha.
\]
The $ha$ are precisely the distinct children of $h$; each leaf has
payoffs given by $u_i:Z\to\mathbb R$. Complete owner policies are
\[
 \Sigma_i=\prod_{h\in H_{\mathrm{nt}}}A_i(h),\qquad
 \Sigma=\prod_i\Sigma_i.
\]
All agents see the public history $h$ before choosing their current
components of $a$, but no agent sees the others' current choices
before making its own. The joint tuple is revealed only when the
child $ha$ is reached. The \textbf{stage game} at $h$ is the one-round
normal-form game with its payoffs determined by a specified complete
family of child continuation profiles, as typed in
Definition~\ref{def:egt-stage-game-induced-by-continuation-profiles}.
Its Nash set need not be nonempty. The one-round interaction is a
concurrent reporting protocol followed by the continuation interface.
\end{definition}
The following construction gives a fixed interface for the entire tree.

\begin{construction}[Concurrent circuit for a finite public tree]

\label{con:egt-public-tree-circuit}
Let its height be $d$. At layer $k$, let $H_k$ contain its histories
of depth $k$, together with absorbing copies of leaves reached
earlier and an error symbol $\bot$. A leaf keeps its label in every
subsequent layer. For agent $i$ take the layer policy set
\[
 S_i^k=\prod_{\substack{h\in H_k\\h\text{ nonterminal}}}A_i(h)
\]
and the finite response set
\[
 B_i^k=\{(h,a):h\in H_k\text{ nonterminal},\ a\in A_i(h)\}
       \amalg\{\mathsf{idle}\}.
\]
The policy $s_i^k$ realizes the function $f_{s_i^k}:H_k\to B_i^k$
which sends a nonterminal $h$ to $(h,s_i^k(h))$ and a leaf or
$\bot$ to $\mathsf{idle}$. The map
\[
 \iota_i^k:S_i^k\longrightarrow(B_i^k)^{H_k},
 \qquad
 s_i^k\longmapsto f_{s_i^k},
\]
is injective, since the value of $f_{s_i^k}$ at every nonterminal history
recovers the chosen action there. The resulting singly owned response family
is therefore the cartesian reindexing
\[
 (\iota_i^k)^*\mathsf{Dec}_i(H_k,B_i^k)
\]
of the full finite decision presentation. In particular it belongs to
$\mathbf{DecArena}_N(\mathcal C)$ by construction. Use the iterated data
comonoid to copy the
public $H_k$ input to all agents and to one transition record; if
layer-labelled data records are desired, retain one additional copied
value at each layer as well. Apply the independent singly owned response families and collect
their responses together with the transition copy of the public state.
The next-state function sends a tuple of correctly tagged responses at
$h$ to the child $ha$, preserves a leaf when all responses are idle, and
sends every other tuple to $\bot$. Set $\bot$ absorbing. This is a
total finite function, so its response strategy exists by
Proposition~\ref{prop:egt-functions}. Every remaining box in the layer is
nonstrategic wiring, copying, collecting, or one of the displayed reindexed
decision families. Hence each layer, and therefore the complete finite tree
circuit, is an open arena in $\mathbf{DecArena}_N(\mathcal C)$.

Start with the root state and compose the $d$ layers. Induction on the
layers proves that $\bot$ is never reached by an available policy: all
responses to a genuine nonterminal state have its tag, and all responses
to a retained leaf are idle. The same induction shows that the state at
each layer is the prescribed public history or the already reached leaf.
The optional copied values record the successive public histories in
separate data ports; hiding them leaves the leaf-only realization.
After the last layer, the reachable output states lie in the leaf set
$Z\subseteq H_d$. Define the nonstrategic leaf projection
$p_Z:D(H_d)\to D(Z)$ as follows. For every $h\in H_d$ it has one
negative input $n_h$; for every $z\in Z$ it has one positive output
$o_z$ with $n_z<o_z$. Distinct input branches conflict, and an input
outside $Z$ has no positive continuation. This strategy is total as
an event-structure map and receptive to every negative input, but is
value-partial on the unreachable complement $H_d\setminus Z$.
Its only added immediate dependencies are negative-to-positive, so it is
courteous. The preceding reachability induction ensures that its
composite with the circuit is value-total on all policy profiles.
It satisfies $p_Z\circ\mathsf{pt}_{H_d}(z)=\mathsf{pt}_Z(z)$.
Composing the layer circuit with $\mathbf j_N(p_Z)$ gives, within
$\mathbf{DecArena}_N(\mathcal C)$, the leaf-only family
$\ell_\sigma:I\to D(Z)$. The finite-function response and collector
equations identify each internal value with its uniquely prescribed value
on that profile; after hiding, the sole output event is its leaf. Hence
\[
 \ell_\sigma\cong
 \mathsf{pt}_Z(\operatorname{play}(\sigma)).
\]
Regard each completed public round as one tree step when forming
$\mathsf{Hist}(T)$. Applying the explicit comparison of
Proposition~\ref{prop:egt-history-expansion} now gives
\[
 e_T\circ\ell_\sigma\cong t_\sigma.
\]
Thus the circuit and the retained-history presentation are related by
a constructed operational strategy, not by identifying independent
copied data ports with prefix-ordered history events. Its ownerwise parameter product
$\prod_k S_i^k$ is canonically the complete contingent policy space,
with no choices added at leaves or at the error state. The error
symbol merely makes the transition functions total on unreachable
tuples; it creates no additional strategic option. Within each round
the response components are tensored, so the agents' simultaneous choices remain causally independent.
\end{construction}

\begin{proposition}[Public trees as iterated epistemic action updates]
\label{prop:egt-public-tree-action-update}
For the finite public tree of
Construction~\ref{con:egt-public-tree-circuit}, let $W_k$ consist of
its valid public histories at layer $k$, including absorbing copies of
earlier leaves but excluding the circuit's unreachable error symbol.
At a nonterminal $h\in W_k$ put
$A(h)=\prod_iA_i(h)$; at a leaf put $A(h)=1$. Form the tagged action
set $E_k=\coprod_{h\in W_k}A(h)$ and enable $(h,a)$ only at $h$.
Take the publicly observed completed action as the action label.
The semantic product-update carrier is then canonically
$\coprod_{h\in W_k}A(h)$, and its deterministic successor map is a
bijection with $W_{k+1}$. Iterating these updates from the root recovers
the full set of reachable complete public histories, including
histories off the path of any particular policy profile. On the
product with the complete-policy context $\Sigma$, the residual
realization at the successor history $ha$ is $g_{ha}r_{ha}$, with
$r_{ha}=r_{h,ha}r_h$ for the subtree restriction maps. Thus history
substitution and owner-policy restriction act on their distinct indices
and commute in the constructed continuation family.
\end{proposition}
\begin{proof}
\leavevmode
\begin{enumerate}[label=\textup{(\arabic*)},leftmargin=*,itemsep=4pt,topsep=3pt]
  \item \textbf{Lifting along the surjection.} The enabling comprehension pairs $h$ with a single tagged legal action
$(h,a)$, so it is isomorphic to $E_k$. A child is uniquely the history
$ha$, and an absorbing leaf has its sole idle successor. Hence the
successor function is injective and surjective onto $W_{k+1}$.
  \item \textbf{Fixed-point argument.} Induction on layers gives the iterated carrier correspondence. The
underlying concurrent circuit realizes the same child map by the
finite response and collector equations in
Construction~\ref{con:egt-public-tree-circuit}. The residual fibre
after $h$ is the constructed subtree game.
  \item \textbf{Policy-level comparison.} Its policy restriction to
$ha$ is the product projection $r_{h,ha}$; composing it with $r_h$
gives $r_{ha}$, by equality of the selected policy coordinates. Proposition~\ref{prop:egt-dependent-realization-transport} applies
separately to the history-indexed linear family and shows that the
resulting component is $g_{ha}r_{ha}$. This keeps off-path fibres: no
update is restricted to the actual reference policy's unique path.
\end{enumerate}
\end{proof}

\begin{definition}[Stage game induced by continuation profiles]
\label{def:egt-stage-game-induced-by-continuation-profiles}
At a nonterminal public history $h$, put
$A(h)=\prod_{i\in N}A_i(h)$. A specified family of complete
child-continuation profiles has type
\[
 s=(s^{ha})_{a\in A(h)}\in\prod_{a\in A(h)}\Sigma^{ha}.
\]
It induces the \textbf{stage game at $(h,s)$}, with one-round utility
maps
\[
 U_i^{h,s}:A(h)\longrightarrow\mathbb R,\qquad
 U_i^{h,s}(a)=u_i\bigl(\operatorname{play}_{ha}(s^{ha})\bigr).
\]
Different continuation families may induce different stage games.
Subgame perfection compares the current action against all alternative
current actions while keeping this entire child-continuation family
fixed, and separately checks Nash equilibrium in every child subgame.
The backward-decomposition theorem retains $s$ as the indexing
parameter of each stage game's Nash set.
\end{definition}
\begin{theorem}[Concurrent backward decomposition]
\label{thm:egt-concurrent-backward}
A profile at $h$ is subgame perfect iff every child continuation is
subgame perfect and its current joint action is a pure Nash
equilibrium of the stage game $(A_i(h),U_i^{h,s})$. Consequently all
pure subgame-perfect profiles are obtained by the recursion
\[
 \operatorname{SPE}(h)=
 \coprod_{(s^{ha})_a\in\prod_a\operatorname{SPE}(ha)}
       \operatorname{NE}\bigl(A_i(h),U_i^{h,s}\bigr),
\]
where the displayed elements retain both the current action and
every chosen child continuation.
\end{theorem}
\begin{proof}
\leavevmode
\begin{enumerate}[label=\textup{(\arabic*)},leftmargin=*,itemsep=4pt,topsep=3pt]
  \item \textbf{Ownerwise alternatives.} Subgame perfection implies the property in every child. A change
only to agent $i$'s current action is among its admissible complete
policy deviations; holding its child policies fixed gives exactly
the stage-game inequality. This proves necessity.
  \item \textbf{Forward implication.} For sufficiency, take a complete deviation of agent $i$. Its new
current action, combined with the other agents' unchanged actions,
selects a child $hb$. In that child the deviator may also change
its entire continuation.
  \item \textbf{Strategic comparison.} The Nash property of the chosen child
profile bounds its resulting payoff by $U_i^{h,s}(b)$. The current
stage Nash inequality bounds this by the payoff of the original
joint action. Thus the complete deviation is not profitable at $h$.
  \item \textbf{Fixed-point argument.} Every proper subgame already has the required property. Induction
on height proves the theorem and the displayed decomposition.
\end{enumerate}
\end{proof}

\noindent Child-continuation policy restriction composes with deeper restriction.
\begin{center}
\begin{tikzcd}[column sep=large,row sep=large]
\Sigma_h\arrow[r,"r_{h,ha}"]\arrow[dr,"r_{h,hab}"']
 &\Sigma_{ha}\arrow[d,"r_{ha,hab}"]\\
 &\Sigma_{hab}
\end{tikzcd}
\end{center}

A simultaneous stage need not have a pure Nash equilibrium, so
the resulting subgame-perfect set may be empty. When each stage has
only one nontrivial owner, its Nash condition becomes the nonempty
argmax condition and the construction recovers
Theorem~\ref{thm:egt-backward}. The same continuation restriction maps
cover both cases, distinguishing finite perfect-information choice
from simultaneous concurrent rounds.

\begin{corollary}[One-round deviations suffice]
For a finite public game tree, a profile is subgame perfect iff,
at every history, no agent can improve by changing only its action
in that round and retaining its original continuation everywhere.
\end{corollary}
\begin{proof}
\leavevmode
\begin{enumerate}[label=\textup{(\arabic*)},leftmargin=*,itemsep=4pt,topsep=3pt]
  \item \textbf{Fixed-point argument.} The forward implication is immediate. For the reverse, induct upward
from the leaves. The induction gives subgame perfection of every
child profile, while the one-round inequalities are the stage Nash
condition. Theorem~\ref{thm:egt-concurrent-backward} proves the
inductive step.
\end{enumerate}
\end{proof}

\section{Epistemic Rationality and Strategic Dynamics}\label{sec:paper-13}

Continuation and epistemic update supply the temporal and
information structures needed to evaluate rationality across states.
An improving alternative is not automatically an executed transition;
when improvement dynamics are studied, their steps must be given
explicitly. Likewise, rationality depends on the opponent policies an
agent considers possible, not just on the realized profile. We first
construct improvement relations and termination ranks, then examine
common belief, elimination, and truthful reporting as predicates on
these situated games.

\subsection{Improvement dynamics and termination}
\label{sec:egt-improvement}

Numerical utility differences can be evaluated on the existing
deviation spans. Comprehension of positive gain then gives a relation
of profitable changes, whose sinks are Nash equilibria. For a finite
nonempty policy carrier, acyclicity permits an integer rank decreasing
along every improvement step. This relates equilibrium to termination
without identifying an arbitrary execution with an improvement path.

\begin{notation}[Utilities on profiles]
In the real-valued results below, $u_i:\Sigma\to\mathbb R$ denotes the
induced profile utility $v_i^G=k_i g$ of a value presentation. This is the
strategic-form notation. A utility evaluated on a typed outcome will
instead have its domain displayed explicitly, as in the reporting
contexts of Section~\ref{sec:egt-reporting}.
\end{notation}

\begin{construction}[The comprehension of profitable alternatives]
\label{con:fib-profitable-comprehension}
On the dependent deviation context $p_1:D_i\to\Sigma$, put
$\gamma_i=u_ip_2-u_ip_1$ and take the comprehension
$J_i=\{\gamma_i>0\}\hookrightarrow D_i$. Its paired endpoint map
is jointly monic and presents $\mathsf{PDev}_i\subseteq\Sigma^2$.
Hence in the assertion fibre over $1_G$,
\[
 \mathsf{NE}_G
 =\bigwedge_i\neg\exists_{p_1:J_i\to\Sigma}\top.
\]
Indeed the existential image consists exactly of profiles with a
positive-gain alternative. The term $\gamma_i$ is retained before
comprehension, so the construction preserves the magnitudes used by
potential and integrability arguments.
\end{construction}

\subsubsection{Profitable changes and termination}

\begin{definition}[Improvement paths and ordinal potentials]
\label{def:egt-improvement-paths-and-ordinal-potentials}
Put $P=\bigcup_{i\in N}\mathsf{PDev}_i\hookrightarrow\Sigma^2$,
the directed relation of strictly profitable unilateral changes.
A \textbf{finite improvement path} of length $n$ is a tuple
$(\sigma_0,\ldots,\sigma_n)\in\Sigma^{n+1}$ with
$\sigma_kP\sigma_{k+1}$ for every $k<n$. An \textbf{infinite
improvement path} is a map $\sigma:\mathbb N\to\Sigma$ with
$\sigma(k)P\sigma(k+1)$ for every $k$. A path is \textbf{maximal} when
it is infinite or its final vertex has no outgoing $P$-edge.

A \textbf{generalized ordinal potential} is a map
$\Phi:\Sigma\to\mathbb R$ satisfying
\[
 \sigma P\tau\Longrightarrow\Phi(\sigma)<\Phi(\tau).
\]
In the numerical setting, an \textbf{exact potential} additionally
satisfies $\Phi(\tau)-\Phi(\sigma)=u_i(\tau)-u_i(\sigma)$
for every $(\sigma,\tau)\in D_i$ and every owner $i$.
Exact potentials can be analyzed through the same gain interface; neither kind
of potential is identified with welfare or an adjustment scheduler.
\end{definition}
\begin{theorem}[Finite improvement and equilibrium]
\label{thm:egt-finite-improvement}
Let the profile space $\Sigma$ be finite and nonempty. The following are
equivalent: $P$ has no directed cycle; every
improvement path is finite; and an integer-valued generalized ordinal
potential exists. Under these conditions every maximal improvement
path ends at a Nash equilibrium, and at least one such equilibrium
exists. A potential can be constructed with values in
$\{0,\ldots,|\Sigma|-1\}$.
\end{theorem}
\begin{proof}
\leavevmode
\begin{enumerate}[label=\textup{(\arabic*)},leftmargin=*,itemsep=4pt,topsep=3pt]
  \item \textbf{Infinite paths and directed cycles.} A directed cycle can be repeated to make an infinite path. Conversely
an infinite path in a finite graph repeats a vertex and therefore
contains a directed cycle. An increasing potential excludes such a
cycle, since returning to its starting vertex would strictly increase
the same real number.
  \item \textbf{Construction of a ranking function.} If the graph is acyclic, define $h(\sigma)$ to be the largest length
of a directed path ending at $\sigma$. Every such path is simple,
so this integer exists and is at most $|\Sigma|-1$. For an edge
$\sigma P\tau$, a longest path to $\sigma$ cannot already contain
$\tau$, since that would give a cycle.
  \item \textbf{Strategic comparison.} Appending the edge proves
$h(\tau)\geq h(\sigma)+1$. Thus $h$ is the required potential. Every maximal path is finite and ends at a vertex with no profitable
deviation, exactly Equation~\eqref{eq:egt-ne}.
  \item \textbf{Strategic comparison: completion.} For existence, take a
profile on which the integer $h$ is maximal. It cannot have an outgoing
improvement edge, since every such edge raises $h$ by at least one. Thus it is a Nash equilibrium.
\end{enumerate}
\end{proof}

\noindent Profitable deviations give a relation whose maximal points are precisely the equilibrium states.
\begin{center}
\begin{tikzcd}[column sep=large,row sep=large]
\mathsf{PDev}_i\arrow[r,hook]\arrow[d,"p_1"']
 &L\times L\arrow[d,"\pi_1"]\\
 L\arrow[r,equal]&L
\end{tikzcd}
\end{center}

This is also a logical termination certificate. The integer
$\rho=|\Sigma|-1-h$ satisfies
\[
 P(\sigma,\tau)\Longrightarrow\rho(\tau)<\rho(\sigma).
\]
The rank is a certificate for paths that actually take profitable
steps. A process permitted to idle at a non-equilibrium state can still
run forever; no fairness principle turns arbitrary execution into an
improvement path. The termination result is therefore about the
displayed improvement relation, not every operational transition.

\subsection{Epistemic rationality and elimination}
\label{sec:egt-epistemic}

Situated rationality combines the realized profile
$s:W\to\Sigma^G$, supplied comparison feedback, and the owners'
possibility relations. An agent's response is tested against the
opponents it considers possible, which need not be the actual opponents.
To connect such a predicate to elimination of strategies, we must keep
the rectangular closure of these possibilities explicit and retain the
witnesses used by best-response and dominance tests.

Public restriction is already the comprehension update of
Section~\ref{sec:egt-communication}; its agent-labelled executions
supply particular announcements. Iterated rationality elimination
uses this same operation, but specifies which predicates are publicly
retained at each stage. We compare this update with rationality under
common belief and with survival under elimination, following the
qualitative epistemic-game analyses of
\citet{VanBenthemGheerbrant2010,VanBenthemPacuitRoy2011}.

\subsubsection{Belief, common belief, and survival under elimination}

\begin{definition}[Pointwise rationality and rectangular closure]
\label{def:egt-rationality-vocabulary}
For $S\hookrightarrow\Sigma=\prod_{j\in N}\Sigma_j$, put
\[
 \operatorname{Rect}(S)=(\pi_j[S])_{j\in N},\qquad
 \operatorname{Prod}(\operatorname{Rect}(S))=
 \prod_{j\in N}\pi_j[S]\hookrightarrow\Sigma.
\]
The latter is the \textbf{rectangular closure}: it recombines the
coordinates of represented profiles independently. Given a
possibility model $s:W\to\Sigma$ with possibility fibres
$P_i(w)=\{w':wR_iw'\}$, form
\[
 Q_i(w)=\prod_{j\in N}s_j[P_i(w)]\hookrightarrow\Sigma.
\]
For numerical utility maps $u_i:\Sigma\to\mathbb R$, owner $i$ is
\textbf{pointwise rational} at $w$ precisely when
\[
 \exists t\in Q_i(w)\ \forall a_i\in\Sigma_i,\qquad
 u_i(s_i(w),t_{-i})\geq u_i(a_i,t_{-i}).
\]
Thus rationality uses a possible pure opponent tuple, not expected
utility or simultaneous optimality against every possible tuple.
The universal property of $\operatorname{Rect}$ and the internally
typed rationality predicate are given below.
\end{definition}
Pointwise rationality asks whether a policy is a pure best response
to some opponent profile considered possible. That differs from
expected-utility maximization and from one response optimal at every
possible opponent profile. In particular, the profiles $(a,b)$ and
$(a',b')$ generate the rectangular opponent set
$\{a,a'\}\times\{b,b'\}$, including the crossed combinations.
The rationality predicate below uses that rectangular closure so that
individual beliefs are not silently treated as a correlated joint
opponent policy.

Fix a possibility model as in Section~\ref{sec:egt-logic}, over a real-valued strategic game. The
rectangular closure used below has a universal property.

\begin{proposition}[Rectangular closure]
\label{prop:egt-rectangular-reflection}
For $S\subseteq\Sigma=\prod_j\Sigma_j$ put
\[
 \operatorname{Rect}(S)=(\pi_j[S])_{j\in N},
 \qquad
 \operatorname{Prod}(L)=\prod_jL_j.
\]
With rectangular restrictions ordered coordinatewise,
\begin{equation}
\label{eq:egt-rectangular-reflection}
 \operatorname{Rect}(S)\leq L
 \quad\Longleftrightarrow\quad
 S\subseteq\operatorname{Prod}(L).
\end{equation}
Hence $\operatorname{Rect}\dashv\operatorname{Prod}$, and
$\operatorname{Prod}(\operatorname{Rect}(S))$ is the closure of $S$ under
independent recombination of its represented owner coordinates.
\end{proposition}
\begin{proof}
\leavevmode
\begin{enumerate}[label=\textup{(\arabic*)},leftmargin=*,itemsep=4pt,topsep=3pt]
  \item \textbf{Adjunction and transposition.} The inequality $\operatorname{Rect}(S)\leq L$ says
$\pi_j[S]\subseteq L_j$ for every owner $j$. This holds exactly when every
$s\in S$ has each coordinate in $L_j$, namely when
$S\subseteq\prod_jL_j$. The adjunction follows.
  \item \textbf{Rectangular-closure unit.} Its unit
$S\subseteq\operatorname{Prod}(\operatorname{Rect}(S))$ is the asserted
rectangular closure.
\end{enumerate}
\end{proof}

\begin{corollary}[Rectangular algebras]
The endomap $\operatorname{Prod}\operatorname{Rect}$ is an idempotent
monad on $\mathsf{Pred}(\Sigma)$. Its algebra poset consists of the
subsets that are products of coordinate subsets, ordered by inclusion.
\end{corollary}
\begin{proof}
\leavevmode
\begin{enumerate}[label=\textup{(\arabic*)},leftmargin=*,itemsep=4pt,topsep=3pt]
  \item \textbf{Adjunction and transposition.} The adjunction gives an extensive monotone map, and the two triangle
inequalities give idempotence. Its fixed subsets are exactly the image
of $\operatorname{Prod}$, because any right-adjoint value is fixed by
the induced posetal monad and every fixed value has that form.
\end{enumerate}
\end{proof}

\begin{definition}[Pointwise rationality]
\label{def:egt-pointwise-rationality}
For agent $i$ define
\[
 Q_i(w)
 =\operatorname{Prod}\!\left(\operatorname{Rect}(s[P_i(w)])\right)
 =\prod_{j\in N}s_j[P_i(w)].
\]
Rectangularity is the closure under independent recombination of strategy coordinates represented at states related to $w$. Player $i$ is \textbf{pointwise rational} at $w$ when
some $t\in Q_i(w)$ makes its current strategy a global best response:
\begin{equation}
\label{eq:egt-point-rat}
 \exists t\in Q_i(w)\ \forall a_i\in\Sigma_i,\qquad
 u_i(s_i(w),t_{-i})\geq u_i(a_i,t_{-i}).
\end{equation}
This is rationality relative to a possible pure opponent profile, not expected
utility relative to a probability distribution. Let $\operatorname{Rat}$ be
the states where every agent is pointwise rational and put
\[
 E(X)=\bigcap_iB_iX.
\]
\end{definition}

\begin{proposition}[Pointwise rationality after private action update]
\label{prop:egt-action-point-rationality}
Let $s:W\to\Sigma$ be the world realization of the qualitative
possibility model, and apply an epistemic action model
$\mathcal E=(E,(S_i),P)$. Use the unchanged reference-profile term
$sp:U\to\Sigma$. At $(w,e)\in U$, define the old-world projection of
the updated possibility set by
\[
 V_i^{\mathcal E}(w,e)=
 \{w'\in W:\exists f\in E,\ 
       R_i(w,w')\wedge S_i(e,f)\wedge P_f(w')\}.
\]
The updated pointwise-rationality predicate of owner $i$ is exactly
\begin{equation}\label{eq:egt-updated-point-rationality}
 \exists t\in\prod_{j\in N}s_j[V_i^{\mathcal E}(w,e)]\ 
 \forall b_i\in\Sigma_i,\qquad
 u_i(s_i(w),t_{-i})\geq u_i(b_i,t_{-i}).
\end{equation}
Thus the same selection predicate is tested against the product-updated
possibility fibre, which can change under private action information
even when the economic policy and utility maps do not change.
\end{proposition}
\begin{proof}
\leavevmode
\begin{enumerate}[label=\textup{(\arabic*)},leftmargin=*,itemsep=4pt,topsep=3pt]
  \item \textbf{Policy-level comparison.} Equation~\eqref{eq:egt-action-updated-relation} says that the
updated possible states at $(w,e)$ are exactly the enabled pairs
$(w',f)$ with $R_i(w,w')$ and $S_i(e,f)$. Their image under the
unchanged profile map $sp$ is the set of profiles $s(w')$ with
$w'\in V_i^{\mathcal E}(w,e)$. Taking rectangular closure gives the
product in Equation~\eqref{eq:egt-updated-point-rationality}, as in
Definition~\ref{def:egt-pointwise-rationality}.
  \item \textbf{Strategic comparison.} The global best-response
inequality uses the same current policy $s_i(w)$ and original
alternatives $\Sigma_i$, proving the formula even when the updated
possibility fibre is empty.
\end{enumerate}
\end{proof}

\begin{construction}[Rationality and common belief]
\begin{equation}\label{eq:egt-true-common-rationality}
 C=\nu X(\operatorname{Rat}\wedge E(X))
  =\operatorname{Rat}\wedge\mathsf{CB}(\operatorname{Rat}).
\end{equation}
The equality is Proposition~\ref{prop:egt-true-common-belief}.
\end{construction}

By Theorem~\ref{prop:egt-fixedpoint}, $C\subseteq\operatorname{Rat}$ and
$P_i(w)\subseteq C$ whenever $w\in C$. Equivalently, $C$ consists of states
that are rational themselves and whose states reachable by every finite
positive-length chain of supplied belief relations are rational. The
factorization above makes explicit that no factivity axiom has been hidden in
the word ``belief.''

\begin{definition}[Local strict dominance and iterated elimination]
\label{def:egt-strict-dominance}
For a restriction $L=(L_j)_{j\in N}$ with $L_j\subseteq\Sigma_j$, a policy
$a_i\in L_i$ is \textbf{strictly dominated within $L$} if there is
$b_i\in L_i$ such that the opponent product $\prod_{j\ne i}L_j$ is inhabited
and
\[
 u_i(b_i,a_{-i})>u_i(a_i,a_{-i})
 \quad\text{for every }a_{-i}\in\prod_{j\ne i}L_j.
\]
\textbf{Iterated elimination} repeatedly removes such dominated policies,
using the remaining policy sets at each step and intersections at limits.
\end{definition}

The tested policy and its possible dominator are both drawn from
the current restriction. The explicit opponent witness also prevents
vacuous strict dominance when the opponent carrier is empty. This is
a different test from being a best response to at least one possible
opponent profile; changing their quantifiers changes the resulting
elimination operator.

\begin{construction}[Pure strict-dominance elimination]
Define iterated elimination of locally strictly dominated pure strategies by
$L_i^0=\Sigma_i$ and
\begin{align*}
 L_i^{\alpha+1}=\{a_i\in L_i^\alpha:
  \ &\text{there is no }b_i\in L_i^\alpha\text{ such that}\\[-2pt]
  &\bigl(\exists a_{-i}\in\prod_{j\ne i}L_j^\alpha\bigr)\ \wedge\\[-2pt]
  &\forall a_{-i}\in\prod_{j\ne i}L_j^\alpha,
       \ u_i(b_i,a_{-i})>u_i(a_i,a_{-i})\},
\end{align*}
with intersections at limit ordinals. The explicit witness of an opponent
profile prevents strict dominance from becoming vacuously true after a
restriction has no realizable opponent tuple. On every inhabited rectangular
restriction this is exactly the ordinary definition. Each stage is determined
extensionally by the preceding restriction; no strategy is selected for
deletion. For finite games the sequence stabilizes after finitely many strict
deletions. The survival theorem below only uses transfinite induction over
these recursively defined stages, so no stabilization claim for arbitrary
strategy sets is needed.
\end{construction}

\begin{theorem}
\label{thm:egt-rationality}
Every strategy played at a state in $C$ survives every stage of this
elimination procedure.
\end{theorem}
\begin{proof}
\leavevmode
\begin{enumerate}[label=\textup{(\arabic*)},leftmargin=*,itemsep=4pt,topsep=3pt]
  \item \textbf{Construction of the comparison.} Put $S_i=s_i[C]$. We prove $S_i\subseteq L_i^\alpha$ for every $i$ by
transfinite induction. The initial and limit cases are immediate.
  \item \textbf{Fixed-point argument.} For the
successor step, take $a_i=s_i(w)$ with $w\in C$. Pointwise rationality gives
$t\in Q_i(w)$ satisfying Equation~\eqref{eq:egt-point-rat}. Since
$P_i(w)\subseteq C$, each coordinate $t_j$ belongs to $S_j$, hence by
induction to $L_j^\alpha$.
  \item \textbf{Strategic comparison.} If a surviving $b_i$ strictly dominated $a_i$
against every surviving opponent profile, then in particular
\[
 u_i(b_i,t_{-i})>u_i(a_i,t_{-i}),
\]
contradicting the global best-response inequality, which quantifies over every
original $b_i$. Thus $a_i$ is not removed at the successor stage.
\end{enumerate}
\end{proof}

This reconstructs a qualitative rationality-to-survival comparison
of the kind considered in \citet{ZvesperApt2010}. Global best
response and local strict dominance have deliberately different
quantifier patterns, and the theorem relies on those exact definitions.
The resulting implication uses the pointwise rationality predicate
of \eqref{eq:egt-point-rat}, not an expected-payoff or arbitrary
best-response variant.

Epistemic conditions leading to Nash equilibrium are studied in
\citet{AumannBrandenburger1995}. Here the comparison keeps actual
and possible outsiders distinct at every step, so the epistemic
assumptions are represented by the supplied possibility relation
rather than folded into the outcome comparison.

\subsubsection{Common belief under public restriction}
\label{sec:egt-relativized-common-belief}

\begin{construction}[Relativized common belief]
For $a:U\hookrightarrow W$ put $A=\exists_a\top$,
$E_W(X)=\bigwedge_i[R_i]X$, and
$E_U(X)=\bigwedge_i[R_i|_U]X$. For $Q$ on $U$, put
$\widehat Q=\forall_aQ$ and
\[
 \mathsf{RCB}_{A,\widehat Q}
 =\nu X\;\bigl(A\Rightarrow
 E_W(A\Rightarrow(\widehat Q\wedge X))\bigr).
\]
\end{construction}

\begin{theorem}[Common belief after restriction]
\label{thm:egt-common-belief-update}
\begin{equation}\label{eq:egt-common-belief-update}
 [!a]\mathsf{CB}_U(Q)=\mathsf{RCB}_{A,\widehat Q}.
\end{equation}
\end{theorem}
\begin{proof}
\leavevmode
\begin{enumerate}[label=\textup{(\arabic*)},leftmargin=*,itemsep=4pt,topsep=3pt]
  \item \textbf{Construction of the comparison.} Let $\rho=a^*$, $J=\forall_a$, and $j_A=J\rho$. The public-update equation gives $JE_U(X)=j_AE_W(JX)$. For $F(X)=E_U(Q\wedge X)$, the operator $JF\rho$ is therefore
\[
 Y\longmapsto j_AE_W(JQ\wedge j_AY)
 =j_AE_W(j_A(\widehat Q\wedge Y)),
\]
since $\widehat Q$ is fixed by the meet-preserving nucleus $j_A$. Apply Corollary~\ref{cor:egt-update-fixedpoints} to $\nu F$.
\end{enumerate}
\end{proof}

\noindent Restriction of the epistemic model is obtained by intersecting each accessibility relation with the announced subspace.
\begin{center}
\begin{tikzcd}[column sep=large,row sep=large]
U\times U\arrow[r,hook]\arrow[d,"\pi_1"']
 &W\times W\arrow[d,"\pi_1"]\\
 U\arrow[r,hook,"m"']&W
\end{tikzcd}
\end{center}

\subsubsection{An exact epistemic characterization}

\begin{definition}[Global pure best responses and point rationalizability]
\label{def:egt-global-pure-best-responses-and-point-rationalizability}
A \textbf{point-rationalizable} policy survives repeated removal of
policies that cannot be a best response to any remaining pure opponent
tuple. In the global operator defined here, a candidate is tested against all of
the agent's original alternatives; only the possible opponents are
restricted at each step. This distinction from local dominance elimination
determines the exact epistemic characterization. No distribution over
opponents is used in this notion of rationalizability.

The preceding survival theorem can be strengthened to an exact
characterization when the elimination operator matches the rationality
predicate. For a finite game, let a restriction be a tuple
$L=(L_i)_{i\in N}$ with $L_i\subseteq\Sigma_i$. Define
\[
 \mathcal B_i(L)=\{a_i\in\Sigma_i:
   \exists t_{-i}\in\prod_{j\ne i}L_j\
   \forall b_i\in\Sigma_i,\quad
        u_i(a_i,t_{-i})\geq u_i(b_i,t_{-i})\}.
\]
The opponent witness ranges over the current restriction, but the
comparison ranges over every original alternative of the agent. Throughout
this paper, ``point rationalizability'' refers to the greatest fixed point
of this displayed global-best-response operator with pure opponent witnesses.
The operator $\mathcal B=(\mathcal B_i)_i$ is monotone on the
product lattice $\prod_i\mathcal P(\Sigma_i)$. Call its greatest
fixed point $R=\nu\mathcal B$ the point-rationalizable restriction.
\end{definition}

\begin{construction}[Descending best-response iteration]
Starting at the full restriction, put
\[
 L^0=\Sigma,
 \qquad
 L^{n+1}=\mathcal B(L^n).
\]
The sequence is descending: $\mathcal B(\Sigma)\subseteq\Sigma$, and if
$L^n\subseteq L^{n-1}$ then monotonicity gives
$L^{n+1}=\mathcal B(L^n)\subseteq\mathcal B(L^{n-1})=L^n$.
The product lattice is finite, so the sequence stabilizes at some $L^*$
with $L^*=\mathcal B(L^*)$. If $F=\mathcal B(F)$ is any fixed point, then
$F\subseteq L^0$, and induction gives
$F=\mathcal B(F)\subseteq\mathcal B(L^n)=L^{n+1}$. Hence
$F\subseteq L^*$, so $L^*$ is the greatest fixed point $R=\nu\mathcal B$.
No separate nonemptiness premise is needed. The fixed-point equation itself
supplies exactly the extension property required by the epistemic
construction.
\end{construction}

\begin{lemma}[Surviving strategies extend to surviving profiles]
\label{lem:egt-rationalizable-extension}
Let $R=\mathcal B(R)$. For every owner $i$ and every $a_i\in R_i$, there is a
profile $w\in\prod_jR_j$ with $w_i=a_i$. Consequently, if any surviving
individual strategy exists then the full product $\prod_jR_j$ is nonempty.
\end{lemma}
\begin{proof}
\leavevmode
\begin{enumerate}[label=\textup{(\arabic*)},leftmargin=*,itemsep=4pt,topsep=3pt]
  \item \textbf{Strategic comparison.} Since $a_i\in R_i=\mathcal B_i(R)$, the definition of $\mathcal B_i$
provides a tuple
\[
 t_{-i}\in\prod_{j\ne i}R_j
\]
against which $a_i$ is a global best response. The profile
$w=(a_i,t_{-i})$ lies in $\prod_jR_j$ and has the required $i$th coordinate. No selection from a family of nonempty sets is involved: the needed opponent
tuple is already the witness carried by membership in the fixed point.
\end{enumerate}
\end{proof}

\begin{theorem}[Rationality with common belief and point rationalizability]
\label{thm:egt-epistemic-exact}
In every possibility model satisfying the definitions of
Section~\ref{sec:egt-epistemic}, every strategy played at a state satisfying
pointwise rationality together with common belief of pointwise rationality
belongs to $R$. Equivalently, the conclusion holds at every state in the
predicate $C$ of Equation~\eqref{eq:egt-true-common-rationality}.
Conversely every profile in $\prod_iR_i$ occurs in a finite epistemic model
in which agents know their own strategies and pointwise rationality is
common knowledge at every state. Thus the surviving strategies have an exact
epistemic realization.
\end{theorem}

\noindent Pointwise rationality uses the actual observation kernel pair on epistemic states.
\begin{center}
\begin{tikzcd}[column sep=large,row sep=large]
W\times_{O_i}W\arrow[r,hook]\arrow[d,"p_1"']
 &W\times W\arrow[d,"\pi_1"]\\
 W\arrow[r,equal]&W
\end{tikzcd}
\end{center}


\begin{proof}
\leavevmode
\begin{enumerate}[label=\textup{(\arabic*)},leftmargin=*,itemsep=4pt,topsep=3pt]
  \item \textbf{Witness and legality comparison.} For the forward direction, let $C$ be the rationality-and-common-belief
predicate constructed above and put $S_i=s_i[C]$. For $a_i=s_i(w)$ with $w\in C$, rationality supplies an opponent
witness drawn from $s_j[P_i(w)]$ for each $j\ne i$. Possibility closure gives $s_j[P_i(w)]\subseteq S_j$.
  \item \textbf{Coinduction and bisimulation.} The rationality inequality is global in the agent's alternatives,
so $a_i\in\mathcal B_i(S)$. Therefore $S\subseteq\mathcal B(S)$. Coinduction in the product lattice gives $S\subseteq R$.
  \item \textbf{Reverse implication.} For the converse take the finite state space
\[
 W=\prod_iR_i.
\]
It may be empty only when there is no surviving individual strategy to
represent, by Lemma~\ref{lem:egt-rationalizable-extension}. Let the strategy
map be inclusion into
$\Sigma$ and define $o_i(w)=w_i$.
  \item \textbf{Choice of representatives and quotients.} Take the epistemic relation to be the
constructed kernel pair
\[
 R_i^{\mathrm{ep}}=\ker(o_i),
\]
so its possibility fiber is
\[
 P_i(w)=\{t\in W:t_i=w_i\}.
\]
Player $i$ therefore knows its own strategy. Since
$R_i=\mathcal B_i(R)$, the current strategy $w_i$ is a global best response
to some tuple in $R_{-i}$. Combining that tuple with $w_i$ gives a state in
$P_i(w)$ and hence a witness for Equation~\eqref{eq:egt-point-rat}.
  \item \textbf{Common knowledge of rationality.} Thus all
agents are rational at every state. Every epistemically accessible state is
again in $W$, so rationality is common knowledge everywhere. Every profile in
$W$ is represented by construction.
  \item \textbf{Policy-level comparison.} Lemma~\ref{lem:egt-rationalizable-extension}
shows in addition that every surviving individual strategy occurs as a
coordinate of one of those profiles. Thus the same constructed model realizes
all surviving behavior without a separate nonemptiness premise.
\end{enumerate}
\end{proof}


The converse constructs a finite state space and observation maps
in which every surviving profile satisfies the specified common-knowledge
rationality condition. It also identifies the relevant elimination
operator: global never-best-response elimination with pure opponent
witnesses. This is the exact characterization compared with
\citet{ZvesperApt2010}, not a claim about every iterative deletion
procedure.

\begin{example}[Pure best responses versus pure strict dominance]
Pure-witness rationalizability need not coincide with elimination
of pure strictly dominated actions. For example, give agent 1
the rows $a,b,c$, agent 2 the columns $L,R$, and utilities
\[
\begin{array}{c|cc}
 & L&R\\\hline
a&2&0\\
b&0&2\\
c&11/10&11/10
\end{array}
\]
for agent 1; give agent 2 utility zero everywhere. No row strictly
dominates $c$, but $c$ is not a best response to either pure column
and is removed by $\mathcal B$. This exact difference is why the
development has not called all these elimination procedures one notion
of rationality. Probability distributions over opponent choices give a further
best-response operator in probabilistic models of games.
\end{example}

The inclusion of a public comprehension has the quantifier
adjoints of Theorem~\ref{thm:fib-comprehension}. When an update also
changes the policy carriers, the counterfactual relation must be
transported by Construction~\ref{con:fib-deviation-comparison}.
The rectangular-image and possibility comparisons therefore occur at
the correct updated strategic context, rather than merely at its
reference world.

\subsubsection{Re-evaluation and repeated public restriction}
\label{sec:egt-announcement-iteration}

Iterated public announcements are successive comprehensions of
the epistemic state carrier, as in
\citet{Plaza2007PublicCommunications}. The relevant rationality or
survival predicate is re-evaluated after each restriction. Keeping this
change of model explicit distinguishes public elimination from simply
repeating a predicate on the original state space.

\begin{proposition}[Finite stabilization of re-evaluated announcements]
\label{prop:egt-announcement-stabilization}
For every submodel carrier $A\subseteq W$, let $\Phi_A\subseteq A$
be a specified updated goal predicate, for example the denotation of a fixed
epistemic formula in $M|_A$. Define
\[
 F_\Phi(A)=\Phi_A\quad\text{viewed as a subset of }W,
 \qquad A_0=W,\quad A_{n+1}=F_\Phi(A_n).
\]
If $W$ is finite, this chain stabilizes after at most $|W|$ strict deletions.
At the stable carrier the goal holds everywhere. If $F_\Phi$ is monotone,
the stable carrier is its greatest fixed point. Without monotonicity that
last conclusion need not hold.
\end{proposition}
\begin{proof}
\leavevmode
\begin{enumerate}[label=\textup{(\arabic*)},leftmargin=*,itemsep=4pt,topsep=3pt]
  \item \textbf{Fixed-point argument.} By construction $F_\Phi(A)\subseteq A$. Every strict step removes at
least one element; thus there are at most $|W|$ strict steps. Equality
at one step is the fixed-point equation.
  \item \textbf{Fixed-point argument: second comparison.} If $F_\Phi$ is monotone,
every fixed point $B$ is contained in $A_0$ and induction gives
$B=F_\Phi(B)\subseteq F_\Phi(A_n)=A_{n+1}$.
  \item \textbf{Greatest fixed point and counterexample.} This proves greatestness. For failure without monotonicity, use the two-world model of Example~\ref{ex:egt-announcement-success} and
$\Phi_A=\llbracket Kp\rrbracket_{M|_A}$. Then
$F_\Phi(W)=\varnothing$ while $F_\Phi(\{x\})=\{x\}$.
\end{enumerate}
\end{proof}

The operator $\mathcal B$ records one such elimination step, and
$\nu\mathcal B$ denotes the states surviving every finite iteration.
Its greatest-fixed-point meaning is justified by the descending
comprehensions, not by treating a public update as an ordinary
reference substitution.

\subsubsection{Public elimination as an epistemic update}

For each current policy carrier, the institutional test agent of
Construction~\ref{con:egt-institutional-public-test} uses the current
elimination predicate as its observable test. By
Proposition~\ref{prop:egt-institutional-public-test}, success of that
public test realizes the next comprehension carrier. Thus the
fixed-point elimination sequence also has an agent-labelled execution
at each step; its predicate is evaluated anew on the surviving carrier.

\begin{construction}[Successive rationality announcements]
Begin with state
space $W_0=\prod_i\Sigma_i$ and observations $o_i(s)=s_i$.
At a stage with $W_n=\prod_iL_i^n$, announce the assertion
\[
 \operatorname{Rat}_n(s)\quad\Longleftrightarrow\quad
   \forall i\ \exists t_{-i}\in L_{-i}^n\ \forall b_i\in\Sigma_i,
       \ u_i(s_i,t_{-i})\geq u_i(b_i,t_{-i}).
\]
This is pointwise rationality in the restricted observation model:
the agent's own action is known, while every surviving opponent
tuple is possible. Public announcement restricts the state space
and its observation equivalence relations as in
Equation~\eqref{eq:egt-announcement}.
\end{construction}

\begin{proposition}[Update equals elimination]
\label{prop:egt-public-elimination}
Repeated announcements of these current rationality predicates give
\[
 W_{n+1}=\prod_i\mathcal B_i(L^n).
\]
They stabilize at the canonical point-rationalizable epistemic model
with state space $\prod_iR_i$.
\end{proposition}
\begin{proof}
\leavevmode
\begin{enumerate}[label=\textup{(\arabic*)},leftmargin=*,itemsep=4pt,topsep=3pt]
  \item \textbf{Descent along the kernel pair.} Within the current product domain, the announcement says independently
for each coordinate that $s_i\in\mathcal B_i(L^n)$. Thus the
new domain is $\prod_i(L_i^n\cap\mathcal B_i(L^n))$. Along the descending iteration already proved,
$\mathcal B_i(L^n)\subseteq L_i^n$, giving the displayed equality.
  \item \textbf{Observation and knowledge.} Observation relations remain equality of the appropriate coordinate,
restricted to the new domain. Each announcement is the singleton action
model with enabling predicate $\operatorname{Rat}_n$ and a trivial
action relation; their iterated product update is canonically the
successive comprehension by
Theorem~\ref{thm:egt-action-composition}. Finite stabilization of the
restrictions therefore proves stabilization of the models as well.
\end{enumerate}
\end{proof}

\noindent An elimination step is the pullback of the surviving-strategy inclusion along the profile map.
\begin{center}
\begin{tikzcd}[column sep=large,row sep=large]
W_{n+1}\arrow[r,hook]\arrow[d,"s_{n+1}"']
 &W_n\arrow[d,"s_n"]\\
 \Sigma_{n+1}\arrow[r,hook]&\Sigma_n
\end{tikzcd}
\end{center}

The resulting public updates realize the successive eliminations
through the announcement and product-update semantics, while retaining
the rationality predicate appropriate to each restricted game. This
explains the link with the elimination result: repeatedly announcing one
fixed subset of the original model would not, in general, implement the
same sequence.

\begin{proposition}[Actual-opponent rationality]
At a pointwise rational state $w$, if every state considered possible by
$i$ has opponents' tuple $s_{-i}(w)$, for every $i$, then $s(w)$ is Nash.
Conversely every Nash profile occurs in a one-state observation model with
these properties.
\end{proposition}
\begin{proof}
\leavevmode
\begin{enumerate}[label=\textup{(\arabic*)},leftmargin=*,itemsep=4pt,topsep=3pt]
  \item \textbf{Observation and knowledge.} The rationality witness has the actual opponents' tuple, so its inequalities
are the Nash inequalities. Conversely take a singleton state with the given
profile and the unique observation maps. Its only possible tuple is the
given profile, and its Nash inequalities establish pointwise rationality.
\end{enumerate}
\end{proof}

\subsection{Reporting agents, incentive compatibility and truthfulness}
\label{sec:egt-reporting}

The reporting problem has two levels of meaning: the type and
message carriers are ordinary data in $\mathsf T$, while their
transmission and receipt are implemented by linear processes in
$\mathsf L$. A report need not coincide with the sender's true type;
that relation is enforced, when appropriate, by the incentive
compatibility predicate.

\subsubsection{Type-indexed reporting interfaces}

\begin{definition}[Economic types and type-dependent valuations]
\label{def:egt-economic-type}
An \textbf{economic type} of owner $i$ is a point $t_i\in T_i$ of
its specified ordinary type carrier. For an outcome or allocation
carrier $Z$, a \textbf{type-dependent valuation} has type
\[
 v_i:T_i\times Z\longrightarrow\mathbb R,\qquad
 (t_i,z)\longmapsto v_i(t_i,z).
\]
If an outcome is presented by a resource certificate in $K$ with
allocation map $a:K\to Z$, the corresponding evaluation of a
certificate is the composite
$T_i\times K\xrightarrow{1\times a}T_i\times Z
\xrightarrow{v_i}\mathbb R$; a certificate need not be identified with
its allocation. The true type is held fixed when reports are varied.
Economic types are ordinary parameters of valuation, distinct from
resource types and from dependent semantic types.
\end{definition}
\begin{definition}[Messages, reports, and truthful reporting]
\label{def:egt-messages-reports}
For an economic type carrier $T_i$ and an allowed message carrier
$M_i$, a \textbf{message} is $m_i\in M_i$. A \textbf{report} is a
message interpreted by the institution as a declaration of information
about type. A \textbf{reporting policy} has type
\[
 s_i:T_i\longrightarrow M_i;
\]
its value $s_i(t_i)$ is the message sent at true type $t_i$.
In a \textbf{direct mechanism} $M_i=T_i$, and the \textbf{truthful
reporting policy} is $\operatorname{tr}_i=\operatorname{id}_{T_i}:
T_i\to T_i$. For finite carriers, its concurrent response realization
has type $\widehat{s_i}:D(T_i)\to D(M_i)$.
A general mechanism may instead use messages that do not denote types;
no probabilistic inference from a message to a true type is implied.
\end{definition}
Each reporting policy is a type-indexed function
$s_i:T_i\to M_i$. It chooses the message associated with each
possible type, including types not realized in the current execution.
No prior or probabilistic inference from a message to a true type is
introduced by the protocol.

\begin{construction}[Concurrent reporting protocol]
For finite message sets, a reporting rule has a direct concurrent
realization even when its ambient outcome set is infinite. Tensor the
agents' message-selection components and collect the chosen messages. The
collector has one negative event for each possible input on each component;
distinct values on a component conflict. For each full tuple it has an
output event labelled by that tuple, dependent on all its input events. An
input can occur as soon as its component is ready, in any order. Its only
new dependencies are negative-to-positive, and it accepts every enabled
input, so it is a receptive courteous strategy. If
$m:\prod_iM_i\to Z$ is the rule, construct its finite image
$F_m=m[\prod_iM_i]$ and corestrict it to
$\bar m:\prod_iM_i\to F_m$. Proposition~\ref{prop:egt-functions} realizes
$\widehat{\bar m}$ on the finite outcome arena $D(F_m)$. Thus the
construction preserves concurrency of reporting while explicitly waiting
for the data the rule needs; no operational presentation of unused elements of
$Z$ is required.
\end{construction}

\subsubsection{Dominant strategies and direct mechanisms}

\begin{definition}[Mechanism, selected dominance, and direct truthfulness]
\label{def:egt-mechanism-selected-dominance-and-direct-truthfulness}
A reporting strategy is \textbf{dominant} when, for every true type and
every outsider message tuple, its prescribed message is at least as good
as every alternative message. The weak inequality permits ties.
A direct mechanism is \textbf{dominant-strategy incentive compatible}
(DSIC) when truthful reporting has this property.
Nash equilibrium fixes the outsiders' actual policies in its comparison;
dominance requires the comparison for every outsider message tuple.
A dominant truthful strategy can coexist with other equilibrium reports.

Fix finite type sets $T_i$, an outcome set $Z$, and utilities
$u_i:T_i\times Z\to\mathbb R$. A mechanism is a family of finite message
sets $M_i$ and a rule $m:\prod_iM_i\to Z$. Its selected strategy is a
family $s_i:T_i\to M_i$. The selection is dominant if
\begin{equation}
\label{eq:egt-dominant}
 \forall i,\theta_i,m_i,m_{-i},\qquad
 u_i\bigl(\theta_i,m(s_i(\theta_i),m_{-i})\bigr)
 \geq
 u_i\bigl(\theta_i,m(m_i,m_{-i})\bigr).
\end{equation}
Only the finite image of the outcome rule is needed for its operational
presentation. No prior distribution over types occurs in this definition.

A direct rule $r:\prod_iT_i\to Z$ is \textbf{dominant-strategy incentive
compatible} (DSIC) if truthful reporting is dominant. In the assertion
language this is the formula
\[
 \forall i,\theta_i,\widehat\theta_i,\widehat\theta_{-i},\qquad
 u_i\bigl(\theta_i,r(\theta_i,\widehat\theta_{-i})\bigr)
 \geq
 u_i\bigl(\theta_i,r(\widehat\theta_i,\widehat\theta_{-i})\bigr).
\]
It is an explicit comparison of two realizations of the same typed
reporting protocol. The report changes; the true type in the utility
comparison does not.
\end{definition}

\begin{construction}[The dependent type of dominant selections]
\label{con:egt-dominant-selection-type}
For each owner, form the comprehension
\[
\begin{aligned}
 \mathsf{Dom}_i(m)=\{(t,a)\in T_i\times M_i:\quad
 &\forall\beta\in M_{-i}\;\forall b\in M_i,\\[-2pt]
 &u_i(t,m(a,\beta))\geq u_i(t,m(b,\beta))\}.
\end{aligned}
\]
with projection $d_i:\mathsf{Dom}_i(m)\to T_i$. Its section set is
$\{s_i:T_i\to\mathsf{Dom}_i(m):d_i s_i=\operatorname{id}_{T_i}\}$.
Define the dependent product
\[
 \mathsf{Sel}(m)
 =\prod_i\mathsf{Sect}(d_i)
 \cong\prod_i\prod_{t\in T_i}\mathsf{Dom}_i(m)_t.
\]
Evaluation gives the selected messages uniformly over $\mathsf{Sel}(m)$;
write $s_i(t)\in M_i$ for the message component. Its defining inequality is
exactly Equation~\eqref{eq:egt-dominant}. Substituting
$b=s_i(\widehat t)$ and $\beta=s_{-i}(\widehat\theta_{-i})$ proves DSIC of
$r_s=m\prod_i s_i$. Thus a selected implementation is precisely an element
of this constructed witness type, and directization is defined uniformly
on that type.

The type can be empty. In the two-player bit-matching game, agent $1$
receives $1$ for equal bits and $0$ otherwise. Message $0$ fails to dominate
against opponent message $1$, and message $1$ fails against opponent message
$0$. With a singleton true-type set, $\mathsf{Dom}_1(m)$ is empty.
In a concrete truthful mechanism, its incentive proof exhibits the
identity report map as an inhabitant of this predicate.
\end{construction}

\begin{proposition}[Dominance as uniform agent selection]
\label{prop:cgt-dominance-selection}
For a true type $t\in T_i$ and fixed outsider messages $\beta$, put
\[
 q_{i,t,\beta}(a)(b)
 \quad\Longleftrightarrow\quad
 u_i(t,m(a,\beta))<u_i(t,m(b,\beta)).
\]
Then selected dominance is
$\forall i,t,\beta\;\varepsilon_{M_i}(s_i(t),q_{i,t,\beta})$.
If $r=m\prod_i s_i$, the direct mechanism's feedback is the pullback
of the mechanism's feedback along the selected report maps in both
comparison arguments.
\end{proposition}
\begin{proof}
\leavevmode
\begin{enumerate}[label=\textup{(\arabic*)},leftmargin=*,itemsep=4pt,topsep=3pt]
  \item \textbf{Strategic comparison.} Expanding selection gives the absence of every strictly improving
message. In the classical real order this is exactly
Equation~\eqref{eq:egt-dominant}. For directization the play equation is
$r=m\prod_i s_i$; applying
Theorem~\ref{thm:cgt-coplay} gives its two-argument feedback equation.
  \item \textbf{Direct-message alternative comparison.} Hence every alternative direct report is compared through its exhibited
message, with the true type fixed in the continuation.
\end{enumerate}
\end{proof}

\subsubsection{Dominance and direct truthfulness in the assertion doctrine}

\begin{proposition}[Dominance in the internal reporting language]
\label{prop:fib-dominance}
For owner $i$, put $\Gamma_i=T_i\times M_{-i}$ and let
$\rho_i:\Gamma_i\times M_i\to\Gamma_i$ project the alternative message.
Define the bad-alternative predicate
\[
 B_i(t,\beta,a)\iff
 u_i(t,m(s_i(t),\beta))<u_i(t,m(a,\beta)).
\]
Selected dominance is
$\bigwedge_i\forall_{!:\Gamma_i\to1}\forall_{\rho_i}\neg B_i$.
Substituting $a=s_i(\widehat t)$ and
$\beta=s_{-i}(\widehat\beta)$ gives direct truthfulness of
$r=m\prod_i s_i$.
\end{proposition}
\begin{proof}
\leavevmode
\begin{enumerate}[label=\textup{(\arabic*)},leftmargin=*,itemsep=4pt,topsep=3pt]
  \item \textbf{Adjunction and transposition.} The adjoint formulas expand the displayed predicate to all true types,
all outsider messages, and all alternative messages, with no improving
alternative. This is Equation~\eqref{eq:egt-dominant}. Universal
elimination substitutes the two displayed report terms, leaving the
true type $t$ fixed in both values.
  \item \textbf{Lifting along the surjection.} The resulting inequality is exactly
DSIC for $r$. The terms take values in constant message types; neither
surjectivity nor a choice of representatives is used.
\end{enumerate}
\end{proof}

Direct truthfulness fixes the true type and tests all outsider messages
and alternative reports. The selected factorization also determines a
universal comparison, as follows.

\begin{construction}[The category of selected implementations]
\label{con:egt-selected-implementation-category}
Fix the type carriers $T_i$, outcome carrier $Z$, and utilities
$u_i:T_i\times Z\to\mathbb R$ of the reporting interface. For a selected
dominant implementation $(M,m,s)$, put $r=m\circ\prod_i s_i$.
Define $\mathbf{Impl}_r$ to have as objects the selected dominant
implementations $(M',m',s')$ of this same rule: each $M'_i$ is a finite
message carrier, $s'_i:T_i\to M'_i$ is dominant for $m'$, and
$m'\circ\prod_i s'_i=r$. A morphism
\[
 f:(M',m',s')\longrightarrow(M'',m'',s'')
\]
is an ownerwise family $f_i:M'_i\to M''_i$ satisfying
\[
 m''\circ\prod_i f_i=m',\qquad f_i\circ s'_i=s''_i.
\]
Identities are coordinatewise identity maps, and composition is
coordinatewise function composition. Substitution in the two displayed
equations proves closure under composition; associativity and the unit
laws follow from those of functions.
\end{construction}

\begin{proposition}[Initiality of the truthful direct implementation]
\label{prop:egt-direct-implementation-initial}
The truthful direct implementation $(T,r,\mathrm{id})$ is initial in
$\mathbf{Impl}_r$. Its unique morphism to $(M',m',s')$ is the selected
policy family $(s'_i)_i$.
\end{proposition}
\begin{proof}
Proposition~\ref{prop:fib-dominance} proves that $r$ is DSIC, so truthful
reporting is dominant and $(T,r,\mathrm{id})$ is an object of
$\mathbf{Impl}_r$. For every object $(M',m',s')$, the equations
$m'\circ\prod_i s'_i=r$ and $s'_i\circ\mathrm{id}_{T_i}=s'_i$
make $(s'_i)_i$ a morphism from this direct implementation. Any such
morphism $f$ must satisfy $f_i\circ\mathrm{id}_{T_i}=s'_i$, hence
$f_i=s'_i$ for every owner. This proves existence and uniqueness.
\end{proof}

\begin{proposition}[Reporting agents as observation-based response interfaces]
\label{prop:egt-reporting-agent-interface}
At a reporting interface where agent $i$ observes its economic type $t_i$,
the complete reporting carrier is $M_i^{T_i}$. Its policy $s_i$ realizes the
response $t_i\mapsto s_i(t_i)$, and for finite carriers its concurrent
realization is $\widehat{s_i}$. Connecting the reporting agents to the
institutional outcome rule gives the reporting protocol above. Uniform
selection over true types and outsider reports is the dominance predicate
of Proposition~\ref{prop:cgt-dominance-selection}.
\end{proposition}
\begin{proof}
\leavevmode
\begin{enumerate}[label=\textup{(\arabic*)},leftmargin=*,itemsep=4pt,topsep=3pt]
  \item \textbf{Choice of representatives and quotients.} Take the observation quotient to be the identity of $T_i$ in the response
interface of Proposition~\ref{prop:egt-agent-observation-interface}. Its policies are
exactly the maps $T_i\to M_i$, and the finite decision generator realizes
each by its response strategy. Tensor and the collector connect these
responses to the same rule.
  \item \textbf{Ownerwise alternatives.} The selection comparison keeps $t_i$ fixed and
varies only the report; quantification over all true types and outsider
reports is exactly the displayed dominance construction. The internal
identification is proved in Proposition~\ref{prop:cgt-dominance-selection}.
\end{enumerate}
\end{proof}

\section{Reconstruction and Exact Abstraction of Agent Systems}\label{sec:paper-14}

The preceding parts built agency using displayed ownership,
concurrent realization, observation, legality, and execution witnesses.
It remains to identify exactly which of those data can be recovered
from the categorical and logical presentations and which distinctions
may be forgotten without changing designated tests. We first reconstruct
the full displayed record, then construct exact reductions relative to
specified observations and deviations. Finally, the calculus is checked
for soundness and its first-layer equational completeness is established
by a classifying model. Representation, operational adequacy,
definability, and completeness are kept separate throughout.

\subsection{Reconstructing the complete agent presentation}
\label{sec:egt-abstraction}
\label{sec:egt-reduction}

The fibred semantics retains a structured agent rather than merely
its truth predicates: owner coordinates, concurrent realizations,
comparisons, observations, legal alternatives, continuations, and
execution witnesses. We first reconstruct this displayed data and its
composition. Exact reductions will then retain the tests and
source-fixed alternative lifts on which strategic and epistemic
assertions depend. This also recovers the policy witnesses of the
communication comparison in
Section~\ref{sec:egt-complete-announcement-policies}
and the observations and transitions of the visible refinement.

\subsubsection{Reconstruction of the structured agent presentation}

Reconstruction concerns the displayed economic presentation.
Ownership, realization, observation, and comparison data must be
recovered together; a collection of truth predicates alone would
not identify who controls a given deviation. The constructions below
exhibit inverse maps between the declared presentation and its
fibred semantic record.

\begin{theorem}[Equivalence of agent presentations and fibred economic presentations]
\label{thm:cgt-lossless}\label{thm:fib-reconstruction}
Fix the operational category $\mathcal C$ of
Section~\ref{sec:egt-interaction}. Its owner-indexed presentations
and the constructed fibred open-arena presentations are equivalent,
with inverse constructions in indexed-family notation and canonical
equivalences in slice notation. The fibred presentation retains
\[
 \vec\Sigma^G,\quad
 \mathsf{Assn}_G\longrightarrow\mathsf{St}(G),\quad
 \Delta_G,\quad\widehat g:1_G\longrightarrow
             \Delta_G\mathcal C(A,B),
\]
the designated owner projections and deviation contexts, and canonical
coplay. It retains every ownerwise reparametrisation, including its
deviation comparison maps, and commutes with sequential and parallel
composition.

Equivalently, by Proposition~\ref{prop:egt-agent-presentation}, these are
the agent presentations with their designated interfaces. The correspondence
recovers the complete-policy composition law and the control regroupings.

The correspondence extends to evaluated, observed, legal,
history-indexed, certified, and epistemically action-decorated games by
retaining their evaluation terms, observation and action-relation terms,
enabling comprehensions, updated observable terms, dependent continuation
families, and certificate carriers with their structure maps. Public
inspection, the full dependent announcement games and their complete-policy
realizations, and the exact visible
refinements are recovered from these retained data. Every interpreted
formula of the agency language has the same denotation.
Player selection, best response, Nash equilibrium, coalition stability,
continuation equilibrium, and truthful implementation are preserved
and reflected on their original carriers.
\end{theorem}

\par\noindent The proof is given in Appendix~\ref{app:metatheory}.

The recovered object includes the full declared presentation.
The hyperdoctrine alone would not determine ownership or concurrent
play, since distinct realizations can share one profile carrier.
Keeping the underlying arena and its distinguished terms makes the
reconstruction reversible. The operational arrow remains taken up to
arena-preserving strategy isomorphism; the reductions to follow are
additional quotients relative to explicitly chosen economic tests.

\begin{corollary}[Recovery from the completed dependent semantics]
\label{cor:egt-completed-reconstruction}
Embedding an original game and its dependent arena families into the
completed model preserves and reflects its concurrent realizations and
all the economic predicates in Theorem~\ref{thm:fib-reconstruction}.
\end{corollary}
\begin{proof}
\leavevmode
\begin{enumerate}[label=\textup{(\arabic*)},leftmargin=*,itemsep=4pt,topsep=3pt]
  \item \textbf{Representation and Yoneda decoding.} For each representable boundary, recover a realized map $\alpha:yA\to yB$
as $\alpha_c(u)u^{-1}$ for any local isomorphism $u:Jc\to A$, by Theorem~\ref{thm:egt-day-preservation}. This is inverse to embedding and commutes with composition, tensor,
and parameter substitution.
  \item \textbf{Predicate calculation.} Ordinary carriers and their maps are
recovered by $\mathsf{Val}\mathsf{Lin}\cong1$; the subterminal fragment,
its quantifiers, and its complete predicate lattices are the same ones. Apply the preceding reconstruction proof to these recovered maps. The argument applies fibrewise to dependent residual arenas and retains
the specified certificate objects and their realization maps.
\end{enumerate}
\end{proof}

\begin{corollary}[Recovery of action-updated concurrent realizations]
\label{cor:egt-action-completed-reconstruction}
For an original concurrent interaction with a specified epistemic
action model, embedding its boundaries by Yoneda preserves the action
enablement comprehension and all ordinary updated epistemic relations.
For an operationally generated action, its post-action residual is
recovered pointwise from the updated dependent linear internal-hom
value family. Secured weak-step witnesses and their observable
predicates are recovered on their own dependent carriers.
\end{corollary}
\begin{proof}
\leavevmode
\begin{enumerate}[label=\textup{(\arabic*)},leftmargin=*,itemsep=4pt,topsep=3pt]
  \item \textbf{Representation and Yoneda decoding.} The epistemic update data and their pullbacks are objects and arrows
of $\mathsf T$, unchanged by the embedding. Yoneda's full faithfulness
recovers each residual strategy from its linear representative;
Proposition~\ref{prop:egt-dependent-realization-transport} commutes
this identification with action-target reindexing.
  \item \textbf{Observation and knowledge.} The witness carrier
of Construction~\ref{con:egt-secured-weak-steps} is ordinary and is
retained before any existential image or observational descent.
\end{enumerate}
\end{proof}

\begin{corollary}[Reconstruction of agents and their communication witnesses]
\label{cor:egt-agent-reconstruction}
The reconstruction equivalence recovers the agent interfaces of
Definition~\ref{def:egt-agent-interface}, their complete-policy composition,
regrouped control, legal choice fibres, observations, and continuation
interfaces. For the complete-policy announcement presentation it also
recovers $\mathsf{AnnPol}_i$, the source $\mathcal X$, the comparison
$F:\mathcal X\to\mathcal U$, and the dependent family of policy witnesses
and concurrent realizations over $\mathcal U$. The retained modal,
fixed-point, dynamic, selection, and coalition predicates have the comparison
equations established in Section~\ref{sec:egt-communication}.
\end{corollary}
\par\noindent The proof is given in Appendix~\ref{app:metatheory}.

\begin{corollary}[Reconstruction of the enriched agent calculus]
\label{cor:egt-enriched-agent-reconstruction}
Retain, in addition to the preceding agent presentation, its displayed
delegation maps, feasibility predicates, institutional execution spans
with realizers and permissions, and local controller data whenever
these are part of the presentation. The reconstruction of these
ordinary carriers and maps and of their linear realization values
also reconstructs the following derived objects and comparisons:
\begin{enumerate}
 \item substituted agent processes, their implemented-policy images,
       and the deviation comparison of
       Theorem~\ref{thm:egt-delegation-image};
 \item informed legal policy carriers and the decentralized-to-pooled
       comparison of Proposition~\ref{prop:egt-decentralized-legal};
 \item composites of witnessed institutions, their permission
       predicates, and their box and diamond operations;
 \item the behaviour of each retained local controller and the
       residual controller generated by each retained complete local policy.
\end{enumerate}
The identifications are canonical up to the same carrier isomorphisms
as the reconstructed presentation.
\end{corollary}
\par\noindent The proof is given in Appendix~\ref{app:metatheory}.

The extra displayed data are genuine parts of the enriched
presentation: they are not claimed to be recoverable from a bare
concurrent arrow. Reconstruction recovers what the presentation
retains and the universal constructions derived from it.
All these constructions can also be made fibrewise over an
ordinary base. Reindexing then selects the same fibrewise products,
images, sections, pullbacks, and controller equations, so the
displayed comparisons commute with base substitution.

\subsubsection{Strategic reduction relative to semantic tests}

\begin{definition}[Semantic tests and contextual policy equivalence]
\label{def:egt-semantic-policy-tests}
A \textbf{semantic-test family} on
$\Sigma=\prod_{i\in N}\Sigma_i$ consists of set-indexed functional
tests $T_\lambda:\Sigma\to X_\lambda$ ($\lambda\in\Lambda$) and
predicate tests $P_\mu\hookrightarrow\Sigma$ ($\mu\in M$).
For owner $i$, transpose these along
$\Sigma\cong\Sigma_{-i}\times\Sigma_i$ into
\[
 \begin{split}
 \kappa_i^{T,P}:\Sigma_i\longrightarrow
 &\left(\prod_{\lambda\in\Lambda}X_\lambda^{\Sigma_{-i}}\right)
  \times\left(\prod_{\mu\in M}\mathsf{Pow}(\Sigma_{-i})\right),\\
 \kappa_i^{T,P}(a_i)=
 &\Bigl(
  (\sigma_{-i}\mapsto T_\lambda(\sigma_{-i},a_i))_\lambda,
  (\{\sigma_{-i}:P_\mu(\sigma_{-i},a_i)\})_\mu
  \Bigr).
 \end{split}
\]
Here $\mathsf{Pow}(X)$ is the power object representing predicates on
$X$. Two policies are \textbf{contextually equivalent} for the tests
precisely when
\[
 a_i\approx_i^{T,P}b_i
 \quad\Longleftrightarrow\quad
 \kappa_i^{T,P}(a_i)=\kappa_i^{T,P}(b_i).
\]
Thus $\approx_i^{T,P}$ is the kernel pair of the contextual-test map.
Its image factorization gives the exact test-induced quotient below;
equivalence compares policies against every fixed outsider profile.
\end{definition}
The quotient identifies precisely the distinctions not visible to
the retained test family. ``Behaviour'' therefore has a relative meaning:
agreement of current payoffs need not imply agreement of residual,
legal, or epistemic observations tested later.

A strategic abstraction should forget a difference in policies
only when the designated assertions cannot detect it. Some of those
tests are functions recording observable outputs; others are
predicates recording comparisons or permissions. The following image
factorization reduces both kinds of test simultaneously, making the
meaning of the resulting equivalence explicit.

\begin{construction}[Contextual test quotient and its comparison category]
Let
\[
 T_\lambda:\Sigma=\prod_i\Sigma_i\longrightarrow X_\lambda
 \qquad(\lambda\in\Lambda)
\]
be a set-indexed family of functional semantic tests, and let
\[
 P_\mu\hookrightarrow\Sigma
 \qquad(\mu\in M)
\]
be a set-indexed family of predicate tests. A predicate is treated by pullback to owner contexts and transposition along the power-object representation. Thus the construction remains valid in the
intuitionistic reading of the assertion doctrine.

We use the predicate-descent form of
Proposition~\ref{prop:egt-predicate-descent}, placed beside functional descent
above. Because $\mathbf{Set}$ is Cartesian closed and its subobject doctrine has
power objects, the contextual tests can be assembled without replacing a
predicate by a chosen Boolean-valued characteristic function. Write
$\mathsf{Pow}(X)$ for the power object representing subobjects of $X$.
For each owner $i$, exponential transposition gives
\[
 \widehat T_{\lambda,i}:\Sigma_i\longrightarrow
 X_\lambda^{\Sigma_{-i}},
 \qquad
 \widehat T_{\lambda,i}(a_i)(\sigma_{-i})
 =T_\lambda(\sigma_{-i},a_i),
\]
while power-object transposition of
$P_\mu\hookrightarrow\Sigma_{-i}\times\Sigma_i$ gives
\[
 \widehat P_{\mu,i}:\Sigma_i\longrightarrow
 \mathsf{Pow}(\Sigma_{-i}),
\]
characterized by
\[
 \sigma_{-i}\in\widehat P_{\mu,i}(a_i)
 \quad\Longleftrightarrow\quad
 P_\mu(\sigma_{-i},a_i).
\]
Define the canonical contextual-test map
\begin{equation}
\label{eq:egt-contextual-test-map}
 \kappa_i^{T,P}
 =\left\langle
   (\widehat T_{\lambda,i})_{\lambda\in\Lambda},
   (\widehat P_{\mu,i})_{\mu\in M}
  \right\rangle:
 \Sigma_i\longrightarrow
 \left(\prod_{\lambda\in\Lambda}X_\lambda^{\Sigma_{-i}}\right)
 \times
 \left(\prod_{\mu\in M}\mathsf{Pow}(\Sigma_{-i})\right).
\end{equation}
Take its image factorization in $\mathbf{Set}$,
\begin{equation}
\label{eq:egt-test-image-factorization}
 \Sigma_i\xrightarrow{q_i^{T,P}}\bar\Sigma_i^{T,P}
 \lhook\joinrel\longrightarrow
 \left(\prod_{\lambda\in\Lambda}X_\lambda^{\Sigma_{-i}}\right)
 \times
 \left(\prod_{\mu\in M}\mathsf{Pow}(\Sigma_{-i})\right).
\end{equation}
Thus $\bar\Sigma_i^{T,P}$ is the image of the map that records exactly the
functional behavior and predicate behavior visible in every opponent context.
Elementwise,
\begin{equation}
\label{eq:egt-test-equivalence}
\begin{split}
 q_i^{T,P}(a_i)=q_i^{T,P}(b_i)
 \quad\Longleftrightarrow\quad
 &\forall\sigma_{-i}\;\forall\lambda\in\Lambda,
 \quad
 T_\lambda(\sigma_{-i},a_i)=T_\lambda(\sigma_{-i},b_i),\\
 &\forall\sigma_{-i}\;\forall\mu\in M,
 \quad
 P_\mu(\sigma_{-i},a_i)\leftrightarrow
 P_\mu(\sigma_{-i},b_i).
\end{split}
\end{equation}
Writing $\mathcal K_i$ for the displayed target of $\kappa_i^{T,P}$,
the image construction and its terminal factorization have the forms
\[
\begin{tikzcd}[column sep=large]
 \Sigma_i\arrow[r,two heads,"q_i^{T,P}"]
 \arrow[dr,"\kappa_i^{T,P}"']&
 \bar\Sigma_i^{T,P}\arrow[d,hook]\\
 &\mathcal K_i
\end{tikzcd}
\qquad
\begin{tikzcd}[column sep=large]
 \Sigma_i\arrow[r,two heads,"f_i"]
 \arrow[dr,"q_i^{T,P}"']&\Theta_i\arrow[d,dashed,"\bar f_i"]\\
 &\bar\Sigma_i^{T,P}.
\end{tikzcd}
\]
The dashed factorization is proved in the theorem below.
No representatives of equivalence classes are chosen: the reduced policy set
is the image object of the displayed universal map.

To state the universal property, let $\mathsf{Quot}_{T,P}(\vec\Sigma)$ have
as objects ownerwise surjections
\[
 f:\vec\Sigma\twoheadrightarrow\vec\Theta
\]
through which every $T_\lambda$ factors and along which every $P_\mu$
descends: there are maps $\widetilde T_\lambda$ and predicates
$\widetilde P_\mu$ with
\[
 T_\lambda=\widetilde T_\lambda\prod_i f_i,
 \qquad
 P_\mu=\left(\prod_i f_i\right)^*\widetilde P_\mu.
\]
A morphism $h:f\to f'$ is an ownerwise map between quotient parameter
objects satisfying $f'=hf$. Surjectivity of $f$ makes the descended maps and
predicates commute with $h$ automatically.
\end{construction}

\begin{theorem}[Test-induced strategic reduction]
\label{thm:egt-test-quotient}
Let $q^{T,P}=\prod_iq_i^{T,P}$. Every functional test descends uniquely to
\[
 \bar T_\lambda:\prod_i\bar\Sigma_i^{T,P}\longrightarrow X_\lambda,
 \qquad
 T_\lambda=\bar T_\lambda q^{T,P},
\]
and every predicate test descends uniquely to
\[
 \bar P_\mu\hookrightarrow\prod_i\bar\Sigma_i^{T,P},
 \qquad
 P_\mu=(q^{T,P})^*\bar P_\mu.
\]
The reduced family has no distinct policies of one owner with the same values
under every functional test and the same truth values under every predicate
test against every reduced opponent profile. Moreover, if
$f_i:\Sigma_i\twoheadrightarrow\Theta_i$ are ownerwise surjections through
which every $T_\lambda$ factors and along which every $P_\mu$ descends, then
there is a unique ownerwise map
\[
 \bar f:\vec\Theta\longrightarrow\bar{\vec\Sigma}^{\,T,P}
 \qquad\text{with}\qquad
 q^{T,P}=\bar f f.
\]
Hence $q^{T,P}$ is terminal in $\mathsf{Quot}_{T,P}(\vec\Sigma)$: it is the
coarsest ownerwise effective quotient preserving all chosen functional and
predicate tests.
\end{theorem}

\noindent The contextual-test map admits its canonical image factorization.
\begin{center}
\begin{tikzcd}[column sep=large,row sep=large]
\Sigma_i\arrow[r,"q_i^{T,P}",two heads]\arrow[dr,"\kappa_i^{T,P}"']
 &\bar\Sigma_i^{T,P}\arrow[d,hook]\\
 &\mathcal K_i
\end{tikzcd}
\end{center}


\begin{proof}
\leavevmode
\begin{enumerate}[label=\textup{(\arabic*)},leftmargin=*,itemsep=4pt,topsep=3pt]
  \item \textbf{Lifting along the surjection.} Each $q_i^{T,P}$ is surjective by image factorization, so their finite product
$q^{T,P}$ is surjective by Proposition~\ref{prop:egt-finite-surjections}. Suppose two profiles have the same $q^{T,P}$-image. Replace their finitely
many owner coordinates one at a time.
  \item \textbf{Fibrewise calculation.} At the $i$th replacement,
Equation~\eqref{eq:egt-test-equivalence} says, against the currently fixed
opponent profile, that every $T_\lambda$ retains its value and every
$P_\mu$ retains its truth value. Hence every $T_\lambda$ is constant on the
fibers of $q^{T,P}$, while every $P_\mu$ is fiber-invariant.
  \item \textbf{Predicate calculation.} Functional
descent follows from Proposition~\ref{prop:egt-kernel-descent} and predicate
descent from Proposition~\ref{prop:egt-predicate-descent}.
  \item \textbf{Separation by contextual tests.} For the separation statement, write
$\bar\Sigma_{-i}=\prod_{j\ne i}\bar\Sigma_j^{T,P}$ and
\[
 \mathcal K_i=
 \left(\prod_{\lambda}X_\lambda^{\Sigma_{-i}}\right)
 \times
 \left(\prod_{\mu}\mathsf{Pow}(\Sigma_{-i})\right)
\]
for the original contextual-test object, and let
\[
 \bar{\mathcal K}_i=
 \left(\prod_{\lambda}X_\lambda^{\bar\Sigma_{-i}}\right)
 \times
 \left(\prod_{\mu}\mathsf{Pow}(\bar\Sigma_{-i})\right)
\]
be its reduced analogue.
  \item \textbf{Conclusion: second comparison.} The quotient map
$q_{-i}:\Sigma_{-i}\twoheadrightarrow\bar\Sigma_{-i}$ induces a canonical map
\[
 \theta_i:\bar{\mathcal K}_i\longrightarrow\mathcal K_i
\]
by precomposition with $q_{-i}$ on every functional component and inverse
image along $q_{-i}$ on every predicate component. Let
\[
 \bar\kappa_i:\bar\Sigma_i\longrightarrow\bar{\mathcal K}_i
\]
be the contextual-test map formed from the descended tests.
  \item \textbf{Lifting along the surjection.} By construction,
\[
 \kappa_i^{T,P}=\theta_i\bar\kappa_iq_i^{T,P}.
\]
If
\[
 \kappa_i^{T,P}=m_iq_i^{T,P}
\]
is its image factorization, epimorphicity of $q_i^{T,P}$ gives
\[
 m_i=\theta_i\bar\kappa_i.
\]
Since the image inclusion $m_i$ is monic, $\bar\kappa_i$ is monic. Hence two
reduced policies with the same values under all reduced functional tests and
the same reduced contextual predicate behavior are equal. No representatives
of the reduced policies are selected in this argument.
  \item \textbf{Policy-level comparison.} For the universal property, suppose $f_i(a_i)=f_i(b_i)$. Against every
original opponent profile the two full profiles have the same $f$-image. Factorization of the functional tests and descent of the predicate tests then
give both equalities in Equation~\eqref{eq:egt-test-equivalence}.
  \item \textbf{Factorization of the representation map.} Thus
$q_i^{T,P}$ equalizes the kernel pair of $f_i$. Since $f_i$ is the
coequalizer of that kernel pair, there is a unique
$\bar f_i:\Theta_i\to\bar\Sigma_i^{T,P}$ with
$q_i^{T,P}=\bar f_i f_i$. Taking the ownerwise product gives $\bar f$.
\end{enumerate}
\end{proof}


\begin{corollary}[Canonical strategic reduction]
\label{thm:egt-strategic-quotient}
Take the single functional test
\[
 T=g:\Sigma\longrightarrow\mathcal C(A,B)
\]
and no predicate tests. Then owner $i$'s contextual-behavior map is
\[
 \kappa_i^g:\Sigma_i\longrightarrow
 \mathcal C(A,B)^{\Sigma_{-i}},
 \qquad
 a_i\longmapsto
 \bigl(\sigma_{-i}\mapsto g_{(\sigma_{-i},a_i)}\bigr),
\]
and its image gives the canonical behavioral reduction $\bar G$. For every coalition, the quotient map is a bounded morphism of deviation
relations. For every ownerwise surjective realization-preserving $2$-cell
$f:G\Rightarrow H$, there is a unique realization-preserving $2$-cell
$\bar f:H\Rightarrow\bar G$ with $q=\bar f f$.
\end{corollary}

\begin{proof}
\leavevmode
\begin{enumerate}[label=\textup{(\arabic*)},leftmargin=*,itemsep=4pt,topsep=3pt]
  \item \textbf{Descent along the kernel pair.} The realization descends by Theorem~\ref{thm:egt-test-quotient}. Write the
realizations of $G,H,\bar G$ as $g,h,\bar g$, respectively. The universal
property supplies a unique ownerwise parameter map $\bar f$ with
$q=\bar f f$.
  \item \textbf{Lifting along the surjection.} Since $f$ is realization-preserving and $q$ is the quotient
presentation map,
\[
 h f=g=\bar g q=\bar g\,\bar f\,f.
\]
The total profile map $f=\prod_i f_i$ is surjective by
Proposition~\ref{prop:egt-finite-surjections}, hence is an epimorphism in
$\mathbf{Set}$. Cancelling it gives
\[
 h=\bar g\,\bar f,
\]
so $\bar f:H\Rightarrow\bar G$ is itself realization-preserving. Uniqueness
as a $2$-cell follows from uniqueness of the ownerwise parameter map.
  \item \textbf{Descent along the kernel pair.} Every coalition alternative descends to the quotient. To lift a reduced
alternative, keep the original outsider coordinates and lift the finite tuple
of coalition coordinates along $\prod_{i\in C}q_i$. This product
is an effective quotient by Proposition~\ref{prop:egt-finite-surjections}, so
such a lift exists without choosing sections.
  \item \textbf{Lifting through the comparison.} Realizations and every
comparison or full-product outcome predicate that factors through realization
have the same truth value before and after this lifting. Thus the quotient has the claimed lifting property, and
Theorem~\ref{thm:egt-modal-adequacy} applies to all descended atoms.
\end{enumerate}
\end{proof}


\begin{proposition}[Reduction and composition]
\label{prop:egt-reduction-comparison}
Let $q_G:G\Rightarrow\bar G$ and $q_H:H\Rightarrow\bar H$ be canonical
behavioral reductions. There is a unique ownerwise surjective,
realization-preserving comparison
\[
 \bar H\circ\bar G\longrightarrow\overline{H\circ G}
\]
whose composite with $q_H\circ q_G$ is the canonical reduction of
$H\circ G$. The analogous comparison exists for tensor. Such a comparison
need not be invertible.
\end{proposition}
\begin{proof}
\leavevmode
\begin{enumerate}[label=\textup{(\arabic*)},leftmargin=*,itemsep=4pt,topsep=3pt]
  \item \textbf{Lifting along the surjection.} The ownerwise products of the two quotient maps are surjective and preserve
the composite realization. Terminality of the composite's test quotient
supplies the unique comparison. It is surjective at each owner because its
composite with a surjection is the surjective canonical quotient.
  \item \textbf{Ownerwise alternatives.} For noninvertibility, let one owner choose $0$ or $1$ in
$G:I\to D(\{0,1\})$, and let $H$ be the nonstrategic response induced by
the constant function $\{0,1\}\to1$. The reduction of $G$ retains two
policies, whereas every composite realization is the same value state in
$D(1)$ and its behavioral reduction has one policy. Thus the comparison
has a two-element domain and a singleton codomain at that owner.
\end{enumerate}
\end{proof}

\[
\begin{tikzcd}[column sep=large]
 H\circ G\arrow[r,Rightarrow,"q_H\circ q_G"]
 \arrow[dr,Rightarrow,"q_{H\circ G}"']&
 \bar H\circ\bar G\arrow[d,Rightarrow,dashed]\\
 &\overline{H\circ G}
\end{tikzcd}
\]

\subsubsection{The reflection and reduced composition}

\begin{definition}[Surjective presentation category]
\label{def:egt-surjective-presentation-category}
For fixed $A,B$, let $\mathsf{Pres}_{\mathrm{epi}}(A,B)$ have open arenas as objects and ownerwise surjective realization-preserving
$2$-cells as morphisms. Let $\mathsf{Red}(A,B)$ be its full subcategory
on presentations whose contextual-behavior maps $\kappa_i^g$ are
injective. Write $I$ for its inclusion.
\end{definition}

\begin{theorem}[Strategic reduction is a reflection]
\label{thm:egt-reduction-reflection}
Behavioral reduction defines a functor
\[
 \operatorname{Red}:\mathsf{Pres}_{\mathrm{epi}}(A,B)
 \longrightarrow\mathsf{Red}(A,B),\qquad
 \operatorname{Red}\dashv I.
\]
Its unit is $q_G:G\to\bar G$. For every morphism $f:G\to H$,
$\operatorname{Red}(f):\bar G\to\bar H$ is an isomorphism.
\end{theorem}

\noindent The quotient comparison satisfies the universal property of a reflection.
\begin{center}
\begin{tikzcd}[column sep=large,row sep=large]
G\arrow[r,"q_G"]\arrow[dr,"f"']
 &\overline G\arrow[d,dashed,"\bar f"]\\
 &H
\end{tikzcd}
\end{center}


\begin{proof}
\leavevmode
\begin{enumerate}[label=\textup{(\arabic*)},leftmargin=*,itemsep=4pt,topsep=3pt]
  \item \textbf{Lifting along the surjection.} Write $g=h\prod_i f_i$. For two $i$-policies $a,b$, one has
\[
 a\sim_G b\quad\Longleftrightarrow\quad f_i(a)\sim_H f_i(b).
\]
For the forward implication, lift every $H$-opponent profile along the
surjection $\prod_{j\ne i}f_j$ and use the realization equation.
  \item \textbf{Reverse implication.} For the reverse implication substitute the image of each $G$-opponent
profile. Hence $q_{H,i}f_i$ descends uniquely through $q_{G,i}$ to a
bijection $\bar f_i$. The descended realization equation follows by
cancelling the surjective total parameter map.
  \item \textbf{Lifting along the surjection.} The equations
$\bar f q_G=q_Hf$ prove identity and composition preservation by
epimorphic cancellation.
  \item \textbf{Lifting along the surjection: second comparison.} If $H$ is reduced, $a\sim_G b$ implies $f_i(a)=f_i(b)$. Thus $f$ factors uniquely as $\widetilde f q_G$ with
$\widetilde f:\bar G\to H$. It is surjective because $f$ is, and
preserves realization by cancelling $q_G$.
  \item \textbf{Reverse implication.} Conversely every such
$\widetilde f$ gives a morphism by precomposition with $q_G$. These inverse operations are natural and give the adjunction. Reduction
is reduced by Theorem~\ref{thm:egt-test-quotient}, so the stated codomain
is correct.
  \item \textbf{Policy-level comparison.} The bijections constructed in the first paragraph prove the
last assertion, including the empty-policy cases.
\end{enumerate}
\end{proof}


\[
\begin{tikzcd}[column sep=large]
 G\arrow[r,two heads,"q_G"]\arrow[d,"f"']&
 \bar G\arrow[d,"\operatorname{Red}(f)","\cong"']\\
 H\arrow[r,two heads,"q_H"']&\bar H
\end{tikzcd}
\qquad
\begin{tikzcd}[column sep=large]
 G\arrow[r,two heads,"q_G"]\arrow[dr,"f"']&
 \bar G\arrow[d,dashed,"\widetilde f"]\\
 &H\quad\text{reduced}
\end{tikzcd}
\]

\begin{construction}[Composition after reduction]
For reduced $G:A\to B$ and $H:B\to C$, put
\[
 H\star G=\operatorname{Red}(H\circ G),
 \qquad H\mathbin{\widehat\otimes}G
 =\operatorname{Red}(H\otimes G).
\]
On $2$-cells apply the original horizontal composition or tensor and
then the reduction functor. Retain the original arena objects and the
reduced identity presentations.
\end{construction}

\begin{theorem}[Coherence of reduced composition]
\label{thm:egt-reduced-coherence}
The preceding construction is a symmetric monoidal bicategory, with
hom-categories $\mathsf{Red}(A,B)$. Its associators, unitors, and
monoidal coherence maps are the canonical comparisons induced from
those of open arenas. In particular,
\[
 \operatorname{Red}(\bar H\circ\bar G)
 \cong\operatorname{Red}(H\circ G),
\]
although the unreduced comparison of
Proposition~\ref{prop:egt-reduction-comparison} need not be invertible.
\end{theorem}

\begin{proof}
\leavevmode
\begin{enumerate}[label=\textup{(\arabic*)},leftmargin=*,itemsep=4pt,topsep=3pt]
  \item \textbf{Composition.} Composition and tensor of ownerwise surjections are ownerwise
surjections. The preceding theorem therefore sends every quotient map
inserted in a composite to an isomorphism after reduction. For example,
reduce the map $q_H\circ q_G:H\circ G\to\bar H\circ\bar G$;
its image is the displayed isomorphism, with the orientation reversed
if necessary.
  \item \textbf{Monoidal comparison.} For three arrows, both $(K\star H)\star G$ and
$K\star(H\star G)$ are consequently canonically isomorphic to the
reduction of an unreduced triple composite. Transport its original
associator through these isomorphisms. Define the unitors and tensor
comparisons in the same way. Naturality follows from the naturality of
$q$ and of the original structural maps.
  \item \textbf{Lifting along the surjection.} For each pentagon, triangle, symmetry, or monoidal interchange diagram,
precompose the two candidate composites with the product of the
unreduced parameter quotients. This is surjective at every owner. The resulting diagram is the corresponding coherence diagram in
$\mathbf{Arena}_N(\mathcal C)$, followed by the unique reduction map,
and therefore commutes.
  \item \textbf{Lifting along the surjection: completion.} Cancel the surjections ownerwise. This proves
all the asserted coherence equations and their compatibility, without
choosing representatives of reduced policies.
\end{enumerate}
\end{proof}


\begin{remark}[The morphisms used by the reflection]
The proof uses the lifting of opponent contexts along the displayed
surjections. Arbitrary parameter maps remain available in the original
open-arena bicategory; the reflection just proved is on its surjective
presentation categories. A change that introduces new opponent policies
is represented there before any new contextual reduction is performed.
\end{remark}

\begin{remark}[Strategy reduction and epistemic state spaces]
For $s:W\to\Sigma$, behavioral reduction replaces it by
$qs:W\to\bar\Sigma$ and preserves $gs=\bar gqs$. It need not identify
states of $W$. Observations and belief relations on $W$ may be retained
unchanged. If a separate state quotient is taken, its relations and atoms
are subject to the bounded-morphism theorem independently. When action
enablement descends along such a bounded state quotient, the updated
state quotient is again bounded by
Corollary~\ref{cor:egt-action-bounded-quotient}; this is not inferred
from equality of play maps alone.
\end{remark}

\subsubsection{Reduction of dependent legal agency}

If an institution distinguishes legal from illegal policies,
legality must be retained by its reduction. Two policies may yield the
same immediate output but admit different legal completions, and hence
different deviations. The construction therefore descends both the
outcome map and the legal domain with its dependent alternative fibres.

\begin{construction}[Legality-sensitive strategic reduction]
The predicate-sensitive strategic reduction now gives the reduction that is
appropriate when institutional legality must survive abstraction. Apply
Theorem~\ref{thm:egt-test-quotient} to the functional test
\[
 g:\Sigma\to\mathcal C(A,B)
\]
and the predicate test $L\hookrightarrow\Sigma$. Write the resulting
ownerwise quotient as
\[
 r_i:\Sigma_i\twoheadrightarrow\bar\Sigma_i,
 \qquad
 r=\prod_i r_i:\Sigma\twoheadrightarrow\bar\Sigma.
\]
Then there are a unique realization $\bar g$ and a unique descended legal
subobject $\bar L\hookrightarrow\bar\Sigma$ such that
\begin{equation}
\label{eq:egt-legal-reduction-descent}
 g=\bar g r,
 \qquad
 L=r^*\bar L.
\end{equation}
Let
\[
 r_L:L\twoheadrightarrow\bar L
\]
be the restriction of $r$. It is surjective: every point of $\bar L$ has an
$r$-preimage by surjectivity, and
Equation~\eqref{eq:egt-legal-reduction-descent} places every such preimage in
$L$.

Construct the realizable-coordinate functor $\bar L_{(-)}$ of the reduced
legal domain in the same way as Proposition~\ref{prop:egt-realizable-coordinate-functor}.
For every $S\subseteq N$, the composite
$\bar e_Sr_L:L\to\bar L_S$ is constant on the fibers of $e_S$, because equal
original $S$-coordinates have equal reduced $S$-coordinates. Effective
descent therefore gives a unique map
\begin{equation}
\label{eq:egt-legal-coordinate-quotient}
 r_S^L:L_S\longrightarrow\bar L_S,
 \qquad
 r_S^Le_S=\bar e_Sr_L.
\end{equation}
It is surjective: for $\bar s\in\bar L_S$, take
$\bar\ell\in\bar L$ with $\bar e_S(\bar\ell)=\bar s$, lift
$\bar\ell$ along $r_L$ to $\ell\in L$, and then
$r_S^L(e_S(\ell))=\bar s$.

These quotient maps are compatible with every coordinate restriction. If
$S\subseteq T$, then
\begin{equation}
\label{eq:egt-legal-coordinate-naturality}
 \bar\rho_S^T r_T^L=r_S^L\rho_S^T.
\end{equation}
Indeed both sides become $\bar e_Sr_L$ after precomposition with the
surjection $e_T$. Hence
\[
 r^L_{(-)}:L_{(-)}\Longrightarrow\bar L_{(-)}
\]
is a levelwise surjective natural transformation of realizable-coordinate
functors.
\end{construction}

\begin{theorem}[Exact descent of legal alternatives]
\label{thm:egt-legal-reduction}
For every $S\subseteq N$, the commutative square
\[
\begin{tikzcd}[column sep=large,row sep=large]
L \arrow[r,"r_L",two heads] \arrow[d,"e_S"'] &
\bar L \arrow[d,"\bar e_S"] \\
L_S \arrow[r,"r_S^L"',two heads] & \bar L_S
\end{tikzcd}
\]
has a surjective comparison map
\begin{equation}
\label{eq:egt-legal-regular-pushout}
 L\twoheadrightarrow L_S\times_{\bar L_S}\bar L,
 \qquad
 \ell\longmapsto(e_S(\ell),r_L(\ell)).
\end{equation}
Thus, in $\mathbf{Set}$, each displayed square is a regular pushout of
surjections. Consequently, for every coalition $C$, $r_L$ is a surjective
bounded morphism from the legal deviation relation $D_C^L$ to
$D_C^{\bar L}$, and the counterfactual realizations commute with the induced
map of deviation objects. Legal Nash equilibrium is therefore preserved and
reflected by $r_L$; the same is true of every modal fixed-point assertion
built from legal deviation relations and atomic predicates that descend
through $r_L$.
\end{theorem}

\noindent Exact legal reduction is a Cartesian square of legal profiles over their policy spaces.
\begin{center}
\begin{tikzcd}[column sep=large,row sep=large]
L\arrow[r,"r_L",two heads]\arrow[d,hook]
 &\bar L\arrow[d,hook]\\
 \Sigma\arrow[r,"r"',two heads]&\bar\Sigma
\end{tikzcd}
\end{center}


\begin{proof}
\leavevmode
\begin{enumerate}[label=\textup{(\arabic*)},leftmargin=*,itemsep=4pt,topsep=3pt]
  \item \textbf{Choice of representatives and quotients.} Take a compatible pair
\[
 (s,\bar\ell)\in L_S\times_{\bar L_S}\bar L.
\]
Surjectivity of $e_S$ gives $\ell_0\in L$ with $e_S(\ell_0)=s$, and
surjectivity of $r_L$ gives $\ell_1\in L$ with
$r_L(\ell_1)=\bar\ell$.
  \item \textbf{Witness and legality comparison.} Compatibility says
\[
 r_S^L(e_S(\ell_0))
 =\bar e_S(r_L(\ell_1)).
\]
Using Equation~\eqref{eq:egt-legal-coordinate-quotient} on the left gives
\[
 \bar e_S(r_L(\ell_0))=\bar e_S(r_L(\ell_1)).
\]
Since $\bar L_S\hookrightarrow\bar\Sigma_S$ is monic, the reduced
$S$-coordinates of $\ell_0$ and $\ell_1$ coincide. Form the ambient profile
$\tau\in\Sigma$ by taking its $S$-coordinates from $\ell_0$ and its
coordinates outside $S$ from $\ell_1$.
  \item \textbf{Ownerwise alternatives.} Ownerwise equality of the reduced
$S$-coordinates gives
\[
 r(\tau)=r(\ell_1)=\bar\ell.
\]
Because $L=r^*\bar L$ and $\bar\ell\in\bar L$, the profile $\tau$ belongs to
$L$. It also satisfies $e_S(\tau)=s$. Hence $\tau$ maps to the chosen
compatible pair, proving surjectivity of
Equation~\eqref{eq:egt-legal-regular-pushout}.
  \item \textbf{Lifting through the comparison.} For the bounded-morphism lifting property take $S=-C$. If
\[
 r_L(\ell)\;D_C^{\bar L}\;\bar\tau,
\]
then
$\bar e_{-C}(r_L\ell)=\bar e_{-C}(\bar\tau)$. By
Equation~\eqref{eq:egt-legal-coordinate-quotient}, the pair
$(e_{-C}(\ell),\bar\tau)$ lies in
$L_{-C}\times_{\bar L_{-C}}\bar L$.
  \item \textbf{Verification of the universal property.} Surjectivity of the comparison map
supplies $\tau\in L$ with
\[
 e_{-C}(\tau)=e_{-C}(\ell),
 \qquad
 r_L(\tau)=\bar\tau.
\]
Thus $\ell D_C^L\tau$, giving the required lift. Forward preservation follows
from the commutative square.
  \item \textbf{Strategic comparison.} The induced map on legal deviation objects sends
$(\ell,\tau)$ to $(r_L\ell,r_L\tau)$. Equation
\eqref{eq:egt-legal-reduction-descent} gives
\[
 \delta_C^{\bar G,\bar L}(r_L\ell,r_L\tau)
 =\delta_C^{G,L}(\ell,\tau).
\]
Hence profitable legal alternatives are preserved and lifted with the same
two realized interactions. Equation~\eqref{eq:egt-legal-ne} is therefore
preserved and reflected.
  \item \textbf{Fixed-point argument.} The fixed-point statement is the bounded-morphism
argument of Theorem~\ref{thm:egt-modal-adequacy}, applied to the legal state
sort $L$ and its legal deviation relations.
\end{enumerate}
\end{proof}


\begin{theorem}[Legal ability under strategic reduction]
\label{thm:egt-legal-ability-reduction}
Let $r_L:L\twoheadrightarrow\bar L$ be the legality-sensitive reduction of
Theorem~\ref{thm:egt-legal-reduction}, and suppose an outcome map
$q:L\to Z$ descends to $\bar q:\bar L\to Z$, so $q=\bar q r_L$. Then for
every coalition $C$ and predicate $P\hookrightarrow Z$,
\[
 [C]_L P
 \quad\Longleftrightarrow\quad
 [C]_{\bar L}P.
\]
Thus the legality-sensitive strategic quotient preserves and reflects
coalition ability as well as legal strategic stability.
\end{theorem}
\begin{proof}
\leavevmode
\begin{enumerate}[label=\textup{(\arabic*)},leftmargin=*,itemsep=4pt,topsep=3pt]
  \item \textbf{Choice of representatives and quotients.} Put
\[
 A=\forall_{e_C:L\to L_C}q^*P,
 \qquad
 \bar A=\forall_{\bar e_C:\bar L\to\bar L_C}\bar q^*P.
\]
We first prove
\[
 A=(r_C^L)^*\bar A.
\]
Suppose $A(c)$ and let $\bar\ell\in\bar L$ satisfy
$\bar e_C(\bar\ell)=r_C^L(c)$. By the surjective comparison map in
Theorem~\ref{thm:egt-legal-reduction}, there is $\ell\in L$ with
$e_C(\ell)=c$ and $r_L(\ell)=\bar\ell$. Hence $P(q(\ell))$, and therefore
$P(\bar q(\bar\ell))$.
  \item \textbf{Reverse implication.} Conversely, suppose $\bar A(r_C^L(c))$ and let $\ell\in L$ satisfy
$e_C(\ell)=c$. Then
\[
 \bar e_C(r_L\ell)=r_C^L(c),
\]
so $P(\bar q(r_L\ell))$, equivalently $P(q(\ell))$. This proves the
predicate equality.
  \item \textbf{Preservation of coalition effectivity.} Finally $r_C^L$ is surjective, so existential quantification over $L_C$ is
preserved and reflected:
\[
 \exists_{!:L_C\to1}A
 \quad\Longleftrightarrow\quad
 \exists_{!:\bar L_C\to1}\bar A.
\]
Using Equation~\eqref{eq:egt-legal-ability} on both sides gives the result.
\end{enumerate}
\end{proof}

In particular, an observable outcome map of the form
$q=o\circ g\circ m$ descends automatically: if
$\bar m:\bar L\lhook\joinrel\longrightarrow\bar\Sigma$ is the descended
legal inclusion, then
\[
 \bar q=o\circ\bar g\circ\bar m
\]
satisfies $q=\bar q r_L$ by
Equation~\eqref{eq:egt-legal-reduction-descent}.

\subsubsection{Continuation-sensitive open-game reduction}

\begin{definition}[Continuation tests]
\label{def:egt-continuation-test}
For a complete history-indexed continuation family with history set
$\mathsf{Hist}$, use its restriction and realization maps
\[
 r_h:\Sigma\to\Sigma^h,\qquad
 g_h:\Sigma^h\to\mathcal C(A_h,B_h),\qquad
 g^h=g_hr_h:\Sigma\to\mathcal C(A_h,B_h)
 \quad(h\in\mathsf{Hist}).
\]
A \textbf{functional continuation test} is the map $g^h$, which
observes the continuation of a complete policy after $h$, whether
reached or off-path. The family $(g^h)_{h\in\mathsf{Hist}}$ induces the
ownerwise contextual equivalence
\[
 a_i\approx_i^{\mathrm{cont}}b_i
 \quad\Longleftrightarrow\quad
 \forall h\in\mathsf{Hist}\ \forall\sigma_{-i}\in\Sigma_{-i},\qquad
 g^h(a_i,\sigma_{-i})=g^h(b_i,\sigma_{-i}).
\]
Comparison predicates $P_h\hookrightarrow\mathcal C(A_h,B_h)^2$
are transported along the same maps; identical continuation
realizations give identical comparison tests. The superscript on $\approx_i^{\mathrm{cont}}$ is suppressed
in the continuation-quotient construction below. Additional factual
or certificate observations are retained by including them in the
general semantic-test family.
\end{definition}
A root outcome test sees only the realized initial interaction.
Subgame perfection evaluates policies after later histories as well.
A reduction preserving root outcomes may therefore identify policies
whose continuation incentives differ. Retaining all designated
continuation tests yields a quotient compatible with policy restriction;
the following example shows why that extra information is necessary.

\begin{example}[Root behavior does not determine sequential rationality]
Root-outcome equivalence alone does not preserve sequential rationality.
Consider a one-player tree: the root can stop for payoff $2$, or continue to
a choice between payoffs $1$ and $0$. The complete policies ``stop, then
choose $1$'' and ``stop, then choose $0$'' have the same root behavior and
are both root Nash equilibria. Only the first is subgame perfect. The
canonical root strategic quotient therefore identifies policies that the
continuation predicates in Equation~\eqref{eq:egt-spe-doctrine} still need
to distinguish.
\end{example}

\begin{construction}[Continuation test quotient]
For each history $h$, let
\[
 g^h:\Sigma\longrightarrow\mathcal C(A_h,B_h)
\]
be the realization of the residual game after restricting complete policies.
Use the family $(g^h)_h$ as the semantic tests of
Theorem~\ref{thm:egt-test-quotient}. This is the family in
Definition~\ref{def:egt-continuation-test}. The resulting ownerwise equivalence is
\[
 a_i\approx_i b_i
 \quad\Longleftrightarrow\quad
 \forall h\;\forall s_{-i},\qquad
 g^h(a_i,s_{-i})=g^h(b_i,s_{-i}).
\]
\end{construction}

\begin{lemma}[Coherent reduction of continuation games]
\label{lem:egt-reduced-continuation-diagram}
For either finite complete-policy tree construction, order histories by
extension and write $r_{hk,i}:\Sigma_i^h\to\Sigma_i^k$ for coordinate
restriction when $h\preceq k$. At each history define
\[
 a_i\approx_i^h b_i
 \quad\Longleftrightarrow\quad
 \forall k\succeq h\;\forall\tau_{-i}\in\Sigma_{-i}^h,\qquad
 g_k\bigl(r_{hk}(a_i,\tau_{-i})\bigr)
 =g_k\bigl(r_{hk}(b_i,\tau_{-i})\bigr).
\]
Let $q_i^h:\Sigma_i^h\twoheadrightarrow\bar\Sigma_i^h$ be its quotient.
There are unique surjective maps
$\bar r_{hk,i}:\bar\Sigma_i^h\to\bar\Sigma_i^k$ satisfying
\begin{equation}
\label{eq:egt-reduced-continuation-square}
 \bar r_{hk,i}q_i^h=q_i^k r_{hk,i}.
\end{equation}
They satisfy $\bar r_{hh,i}=1$ and
$\bar r_{k\ell,i}\bar r_{hk,i}=\bar r_{h\ell,i}$. Every $g_h$ descends
uniquely to $\bar g_h$ with $g_h=\bar g_hq^h$, where $q^h=\prod_iq_i^h$.
At the root, this is the continuation test quotient above.
\end{lemma}

\begin{proof}
\leavevmode
\begin{enumerate}[label=\textup{(\arabic*)},leftmargin=*,itemsep=4pt,topsep=3pt]
  \item \textbf{Lifting along the surjection.} The coordinate projections satisfy $r_{hh,i}=1$ and
$r_{k\ell,i}r_{hk,i}=r_{h\ell,i}$. They are surjective: the finitely many
omitted coordinates have nonempty action sets, so finite induction supplies
their product, and adjoining those coordinates extends any subtree policy. Equivalently, the paper's reference-based extension gives a section after
fixing the reference policy. Neither argument selects sections globally.
  \item \textbf{Lifting along the surjection: second comparison.} Suppose $a_i\approx_i^h b_i$. For $\ell\succeq k\succeq h$ and an
opponent profile $\eta_{-i}\in\Sigma_{-i}^k$, lift $\eta_{-i}$ to
$\tau_{-i}\in\Sigma_{-i}^h$.
  \item \textbf{Stability under continuation restriction.} Restriction composition gives
\begin{align*}
 g_\ell\bigl(r_{k\ell}(r_{hk,i}a_i,\eta_{-i})\bigr)
 &=g_\ell\bigl(r_{h\ell}(a_i,\tau_{-i})\bigr)\\
 &=g_\ell\bigl(r_{h\ell}(b_i,\tau_{-i})\bigr)\\
 &=g_\ell\bigl(r_{k\ell}(r_{hk,i}b_i,\eta_{-i})\bigr).
\end{align*}
Thus $r_{hk,i}a_i\approx_i^k r_{hk,i}b_i$, and the quotient universal
property supplies the unique map in
Equation~\eqref{eq:egt-reduced-continuation-square}.
  \item \textbf{Lifting along the surjection.} It is surjective because
$q_i^k r_{hk,i}$ is surjective. Moreover,
\[
 \bar r_{k\ell,i}\bar r_{hk,i}q_i^h
 =q_i^\ell r_{k\ell,i}r_{hk,i}
 =q_i^\ell r_{h\ell,i}
 =\bar r_{h\ell,i}q_i^h.
\]
Cancelling the epimorphism $q_i^h$ proves composition coherence; the identity
law follows in the same way. The test at $h$ occurs in the defining family.
  \item \textbf{Ownerwise alternatives.} Replacing the finitely many owner coordinates successively therefore proves
that $g_h$ is constant on the fibres of $q^h$, giving its unique descent. The root tests are exactly $(g^k)_k$.
\end{enumerate}
\end{proof}


\begin{theorem}[Continuation-sensitive reduction]
\label{thm:egt-continuation-quotient}
The test-induced quotient for the family $(g^h)_h$ is the coarsest
ownerwise quotient preserving all continuation realizations and has the
universal factorization property of Theorem~\ref{thm:egt-test-quotient}.
For the finite complete-policy trees, the reduced games and restriction maps
of Lemma~\ref{lem:egt-reduced-continuation-diagram} preserve and reflect
subgame-perfect profiles. Writing $o$ for the root,
\begin{equation}
\label{eq:egt-spe-quotient}
 \mathsf{SPE}_G
 =(q^o)^*\left(\bigwedge_h
      \bar r_{oh}^*\mathsf{NE}_{\bar G_h}\right).
\end{equation}
\end{theorem}
\begin{proof}
\leavevmode
\begin{enumerate}[label=\textup{(\arabic*)},leftmargin=*,itemsep=4pt,topsep=3pt]
  \item \textbf{Reduced continuations and test descent.} Theorem~\ref{thm:egt-test-quotient} gives descent and the universal property. Lemma~\ref{lem:egt-reduced-continuation-diagram} constructs the reduced
continuation diagram. For each $h$, an original unilateral alternative maps
to a reduced alternative.
  \item \textbf{Reverse implication.} Conversely any reduced unilateral alternative
lifts through $q_i^h$, while the reference opponents remain fixed. Since
$g_h=\bar g_hq^h$, both realized interactions in the comparison are
preserved. Hence, for each continuation comparison,
\[
 \mathsf{NE}_{G_h}=(q^h)^*\mathsf{NE}_{\bar G_h}.
\]
The product of the ownerwise descent squares gives
$q^h r_{oh}=\bar r_{oh}q^o$.
  \item \textbf{Functoriality and substitution.} Functoriality of reindexing and its preservation
of meets therefore yield
\[
 \bigwedge_h r_{oh}^*\mathsf{NE}_{G_h}
 =(q^o)^*\left(\bigwedge_h
    \bar r_{oh}^*\mathsf{NE}_{\bar G_h}\right),
\]
which is Equation~\eqref{eq:egt-spe-quotient}.
\end{enumerate}
\end{proof}

\noindent Continuation-sensitive quotient maps commute with
restriction to subsequent histories, so the reduced game retains the
designated sequential comparisons.
\begin{center}
\begin{tikzcd}[column sep=large,row sep=large]
\Sigma_i^h\arrow[r,"q_i^h",two heads]\arrow[d,"r_{hk,i}"']
 &\bar\Sigma_i^h\arrow[d,"\bar r_{hk,i}"]\\
 \Sigma_i^k\arrow[r,"q_i^k"',two heads]
 &\bar\Sigma_i^k
\end{tikzcd}
\end{center}

Thus the two strategic reductions in the paper are not different
abstraction mechanisms. They are the same test-induced quotient with
different observation families: root realization for ordinary strategic
stability, and all residual realizations for sequential stability.

\[
\begin{tikzcd}[column sep=large]
 \Sigma_i^h\arrow[r,"\operatorname{ext}_{h,\sigma_i}"]
 \arrow[dr,equal]&\Sigma_i\arrow[d,"r_{h,i}"]\\
 &\Sigma_i^h
\end{tikzcd}
\qquad
\begin{tikzcd}
 \Sigma\arrow[r,two heads,"q"]\arrow[dr,"g^h"']&\bar\Sigma
 \arrow[d,dashed,"\bar g^h"]\\
 &\mathcal C(A_h,B_h)
\end{tikzcd}
\]

\subsection{Soundness and classifying-model completeness}
\label{agt-sec-metatheory}

\subsubsection{The two-layer soundness theorem}

The soundness statement connects the four levels already
constructed: the ordinary judgments of Section~\ref{agt-sec-calculus},
the displayed agent operations, the finite protocols of
Section~\ref{agt-sec-protocols}, and the contextual assertions in
$\cT(W)$, $\cL(W)$, and $\cP(W)$. Ordinary equality is
interpreted through diagonal comprehension, while linear variables
remain governed by tensor-context splitting. Agent abstraction is
internal-hom transposition; observed responses are effective
coequalizer factors; legal policies are dependent universal images;
and institutional composition matches witnesses by pullback.
Primitive comparisons, permissions, and undelivered observations are
interpreted by their declared typed data, not inferred from the
operational arrows. The inference rules and their proof obligations
are displayed in the appendices.

\begin{theorem}[Soundness of the two-layer agency calculus]
\label{agt-thm-soundness}
The completed true-concurrent semantics of Section~\ref{sec:paper-1} interprets
every well-formed type, term, display, and contextual assertion of this
calculus. Derivable term equations are equalities of morphisms; derivable
assertion entailments are inclusions of predicates. Interpretation preserves
ordinary and linear substitution, agent composition, owner regrouping,
delegation, finite protocol execution, perspective, witnessed institutional
composition, and epistemic update.
\end{theorem}
\par\noindent The proof is given in Appendix~\ref{app:metatheory}.

The two semantic layers remain distinct: linear terms retain
resources and execution witnesses, whereas assertions retain truth
values. Soundness does not identify their units, products, or
notions of equality.

\subsubsection{Reconstruction of the displayed economic model}

\begin{theorem}[Exact reconstruction]
\label{agt-thm-reconstruction}
Encoding a displayed owned presentation on original arenas as an agent
state and decoding it by Theorem~\ref{agt-thm-representation} recovers that
presentation. The construction extends to situated presentations,
regroupings, delegation maps, witnessed institutions, and complete finite
communication presentations, retaining every displayed ordinary carrier
and map. Both round trips are identity up to their canonical isomorphisms.
\end{theorem}

\noindent Agent reconstruction is the inverse of policy-indexed currying and Yoneda decoding.
\begin{center}
\begin{tikzcd}[column sep=large,row sep=large]
\Set(\Sigma,\mathcal C(A,B))\arrow[r,"\mathrm{encode}"]\arrow[dr,"\id"']
 &\mathcal V(\I,\Ag_{\vec\Sigma}(yA,yB))\arrow[d,"\mathrm{decode}"]\\
 &\Set(\Sigma,\mathcal C(A,B))
\end{tikzcd}
\end{center}

\par\noindent The proof is given in Appendix~\ref{app:metatheory}.

Exact reconstruction depends on retaining the ownership signature
and execution witnesses as displayed data. The underlying linear
agent object does not recover who controls its policy factors, and
the predicate doctrine cannot recover a hidden execution from its
truth value. The reconstruction theorem recovers precisely the data
that were retained.

\begin{center}
\small
\begin{tabular}{p{0.39\textwidth}p{0.53\textwidth}}
\toprule
Type-theoretic construction & Recovered semantic construction\\
\midrule
Agent state & Owned complete-policy realization\\
Linear sequential and parallel operations & Composition and tensor of concurrent realizers\\
Ownerwise parameter substitution & Realization-preserving policy maps\\
Regrouping and delegation & Collective control and implemented parent policies\\
Linear dependent products & Families of perspective-dependent resources\\
Predicate observation box & Knowledge from observation fibres\\
Uniform policy section & Observation-based legality at every compatible world\\
Session interaction witness & Secured communication history\\
Institutional span composition & Matching records, realizers, and permissions\\
Public comprehension & Restriction of both epistemic endpoints\\
Complete communication policies & Truthful reporting and delivered private/public observations\\
\bottomrule
\end{tabular}
\end{center}

\subsubsection{The classifying model of the first layer}

Fix a typed primitive signature $\mathcal K$ for ordinary sorts and
operations, linear boundary sorts, and linear generators. All primitive
declarations have specified arities. The first-layer theory
$\mathbb T_{\mathrm{Ag}}(\mathcal K)$ is generated by the rules of
Section~\ref{agt-sec-calculus} and Section~\ref{agt-app-rules}; agent types,
composition, regrouping, and delegation are the displayed definitions above.
Finite data and session presentations use the exhibited interpreted
generators of Construction~\ref{agt-con-finite-generators}.  A primitive
signature declares operation symbols with typed arities; it does not silently
assert that unrelated symbols obey equations proved only in one semantic
interpretation.  Consequently the following completeness statement concerns
the equational first-layer theory, not operational progress or equality in
one fixed concurrent model.

A split indexed linear--nonlinear model consists of an ordinary
comprehension model for the stated extensional dependent type constructors
(with equality comprehensions given by pullbacks of diagonals),
an indexed symmetric monoidal closed linear category with additives and
dependent linear sums and products, and a substitution-stable adjunction
$\Lin\dashv\Val$ with strong monoidal $\Lin$. Its indexed adjoints obey
the base-change and Frobenius comparisons expressed by the rules. A model
of $\mathcal K$ interprets its declared constants at their arities. These
are the semantic structures specified by the type constructors, rather
than an identification of all ordinary contexts with sets.

\begin{construction}[Syntactic indexed model]
\label{agt-con-syntactic}
Let $\mathcal B_{\mathrm{syn}}$ have ordinary contexts as objects and
ordinary substitutions modulo derivable equality as arrows. Its terminal
object is the empty context and its comprehension is context extension.
Over $\Gamma$, let $\cL_{\mathrm{syn}}(\Gamma)$ have linear types modulo
conversion as objects and arrows $A\to B$ represented by
\[
 \Gamma;x:A\linturn t:B
\]
modulo derivable equality and bound-variable renaming. Composition is linear
substitution. Reindexing substitutes ordinary parameters in types and terms.
The ordinary type fibre is defined similarly using ordinary terms.
\end{construction}

\begin{proposition}[Diagonal construction of extensional equality]
\label{agt-prop-equality-diagonal}
In an ordinary extensional comprehension model, equality between
sections $a,b:\Gamma\to\Gamma.X$ is the diagonal pullback
\[
\begin{array}{ccc}
 \Gamma.\operatorname{Eq}_X(a,b)&\longrightarrow&\Gamma.X\\
 \downarrow e&&\downarrow\Delta_X\\
 \Gamma&\xrightarrow{\langle a,b\rangle}&
 \Gamma.X\times_\Gamma\Gamma.X.
\end{array}
\]
A term of this equality family yields the equation $a=b$; two terms
of the same equality family are equal. These properties are stable
under ordinary substitution. The corresponding extensional equality conversion in the syntax
\emph{presents} this universal property: the categorical argument proves
its soundness, not a derivation of equality reflection from weaker
syntactic constructors.  It is part of the chosen \emph{ordinary}
foundation, not a rule for agents or linear resources, and it does not
follow from intensional identity types or from dependent sums and
products alone.
\end{proposition}
\begin{proof}
\leavevmode
\begin{enumerate}[label=\textup{(\arabic*)},leftmargin=*,itemsep=4pt,topsep=3pt]
  \item \textbf{Extensional equality elimination.} The diagonal and its pullback $e$ are monomorphisms. The pullback
equations give $ae=be$. A term $p$ of the equality family is a
section $\bar p$ of $e$, so
$a=ae\bar p=be\bar p=b$.
  \item \textbf{Functoriality and substitution.} For two such sections $p,q$, monicity gives $p=q$ from
$ep=eq=\id$. Pullback stability of diagonals proves substitution
stability. This proof derives the categorical consequences of the
chosen extensional equality interpretation; it does not identify
intensional equality with extensional conversion.
\end{enumerate}
\end{proof}

\begin{lemma}[General-map dependent adjoints in the syntax]
\label{agt-lem-general-syntactic-adjoints}
For every definable ordinary substitution $f:V\to W$, the ordinary
dependent sums, products, and extensional equality construct left and
right adjoints to reindexing on the syntactic linear fibres.  They obey
Beck--Chevalley for pullback squares and Frobenius for tensor.
\end{lemma}
\par\noindent The proof is given in Appendix~\ref{app:metatheory}.

\begin{lemma}[Equality-fibre base change and Frobenius coherence]
\label{agt-lem-fibre-coherence}
For a pullback of definable ordinary substitutions
\[
\begin{tikzcd}[column sep=large,row sep=large]
V'\arrow[r,"a"]\arrow[d,"f'"']\arrow[dr,phantom,"\lrcorner",very near start]
 &V\arrow[d,"f"]\\
W'\arrow[r,"b"']&W,
\end{tikzcd}
\]
there is a substitution-natural equivalence of equality fibres
\[
 \chi_{w'}:\operatorname{Fib}_{f'}(w')
 \xrightarrow{\ \cong\ }\operatorname{Fib}_f(b(w')).
\]
Its induced comparison maps give
\[
 b^*\Sigma_f^\ell\cong\Sigma_{f'}^\ell a^*,\qquad
 b^*\Pi_f^\ell\cong\Pi_{f'}^\ell a^*.
\]
The comparisons compose coherently under iterated pullbacks, and
intertwine the Frobenius maps
$\Sigma_f^\ell(f^*B\otimes A)\cong B\otimes\Sigma_f^\ell A$.
\end{lemma}
\par\noindent The proof is given in Appendix~\ref{app:metatheory}.

\begin{lemma}[The syntax carries the required structure]
\label{agt-lem-syntactic-model}
The construction is a split indexed linear--nonlinear model of
$\mathbb T_{\mathrm{Ag}}(\mathcal K)$.
\end{lemma}
\par\noindent The proof is given in Appendix~\ref{app:metatheory}.

\begin{theorem}[Initiality and equational completeness]
\label{agt-thm-completeness}
Every interpretation of $\mathcal K$ in a split indexed linear--nonlinear
model extends uniquely, up to its specified coherent structure
comparisons, to an interpretation of the first-layer syntax. For well-typed terms $t,u$,
\[
 \mathbb T_{\mathrm{Ag}}(\mathcal K)\vdash t=u
 \quad\Longleftrightarrow\quad
 \sem t_{\mathcal M}=\sem u_{\mathcal M}
 \text{ in every such model }\mathcal M.
\]
The same assertion holds for its derived agent operations.
\end{theorem}

\par\noindent The proof is given in Appendix~\ref{app:metatheory}.

\subsubsection{Operational and modal completeness in the fixed semantics}

The classifying-model result concerns derivability relative to all split
indexed linear--nonlinear models. The fixed concurrent model admits stronger
comparison results of a different kind: original arenas separate its open
agent morphisms, while its retained transition relations admit a logical
characterization by modal tests. These statements concern equality and
expressivity in the specified semantics, rather than derivability of every
fixed-model equation in the freely presented syntax.

\begin{theorem}[Original-arena probes separate open agent morphisms]
\label{thm:agt-concurrent-probe-separation}
For original arenas $A,B$ and an ordinary complete-policy carrier $\Sigma$,
put $H_\Sigma(A,B)=[\Lin\Sigma\otimes yA,yB]\in\cV$. For every
$K\in\mathcal C_0$ there is a natural bijection
\begin{equation}\label{eq:agt-concurrent-probe}
 \cV(y(JK),H_\Sigma(A,B))
 \cong \Set\bigl(\Sigma,\cC(JK\otimes A,B)\bigr).
\end{equation}
For any $T\in\cV$ and $f,g:T\to H_\Sigma(A,B)$, the following are
equivalent:
\begin{enumerate}[label=\textup{(\roman*)},leftmargin=*]
\item $f=g$ in the fixed presheaf category;
\item for every $K\in\mathcal C_0$, every $z:y(JK)\to T$, and every
$\sigma\in\Sigma$, the concurrent arrows $JK\otimes A\to B$
decoded from $fz$ and $gz$ at $\sigma$ are equal in $\cC$.
\end{enumerate}
The comparison is stable under ordinary parameter substitution, by
applying it in each indexed fibre. For $T=yR$, it reduces to pointwise
equality of the families of actual concurrent arrows $R\otimes A\to B$.
\end{theorem}
\begin{proof}
\leavevmode
\begin{enumerate}[label=\textup{(\arabic*)},leftmargin=*,itemsep=4pt,topsep=3pt]
  \item \textbf{Closedness, copowers, and Yoneda decoding.} By closedness,
the copower description of $\Lin$, and strong monoidality of restricted
Yoneda, there are natural isomorphisms
\[
\begin{aligned}
 \cV(y(JK),H_\Sigma(A,B))
 &\cong\cV(y(JK)\otimes\Lin\Sigma\otimes yA,yB)\\
 &\cong\prod_{\sigma\in\Sigma}
             \cV(y(JK\otimes A),yB)\\
 &\cong\prod_{\sigma\in\Sigma}\cC(JK\otimes A,B).
\end{aligned}
\]
The second isomorphism distributes Day tensor over the copower and
identifies the resulting representables; the third is restricted Yoneda.
Each comparison is natural in the arena probe and in the policy
coordinate.
  \item \textbf{Separation by original arenas.} Yoneda identifies a map
$z:y(JK)\to T$ with an element of $T(K)$. Two natural transformations
$f,g:T\to H_\Sigma(A,B)$ are equal exactly when their component maps
agree on every such element, for every $K\in\cC_0$. The displayed
bijection identifies this equality componentwise with equality of the
decoded concurrent arrows at all $\sigma\in\Sigma$, proving both
implications.
  \item \textbf{Representable contexts and substitution.} For $T=yR$,
choose $K\in\cC_0$ and an isomorphism $JK\cong R$; evaluation at the
corresponding Yoneda element reduces the comparison to equality of the
families $R\otimes A\to B$. In a dependent fibre all comparisons
are pointwise, so ordinary reindexing preserves them. Ownerwise
reparametrization simply precomposes the displayed policy family.
\end{enumerate}
\end{proof}

\begin{corollary}[Effective semantic equality of finite causal agents]
\label{cor:agt-finite-agent-equality}
Consider two closed agent states with the same displayed finite policy
carrier $\Sigma$ and boundaries $A,B$, generated by the explicitly
presented finite causal protocols of
\cref{thm:egt-causal-protocol,agt-thm-protocol}. Their denotations in
$\cV$ are equal if and only if, for every $\sigma\in\Sigma$,
their realized finite strategies are isomorphic over
$A^\perp\parallel B$. Equality of these explicitly encoded finite
interpretations is decidable: a successful comparison supplies an
arena-preserving event-structure isomorphism for every policy
coordinate. The procedure also tests an exhibited finite ownerwise
reparametrization.
\end{corollary}
\begin{proof}
\leavevmode
\begin{enumerate}[label=\textup{(\arabic*)},leftmargin=*,itemsep=4pt,topsep=3pt]
  \item \textbf{Recovery of concurrent arrows.} For closed agent states,
\cref{agt-thm-representation} and restricted Yoneda identify semantic
equality with equality of the realized arrows in $\cC(A,B)$ at each
complete policy. By the definition of $\cC$, equality of two such
arrows is precisely isomorphism of their strategy event structures over
the common polarized arrow arena.
  \item \textbf{Finite comparison procedure.} In the finite causal
protocols, the policy sets, event carriers, event labels, orders, and
consistency tables are finite explicitly specified data. For each
policy, enumerate the finitely many bijections between the two event
carriers and test preservation of order, consistency, polarity, and
arena labelling. The procedure returns an isomorphism exactly when one
exists and terminates over the finite policy carrier.
  \item \textbf{Ownerwise reparametrization.} For a displayed finite map
$r:\Theta\to\Sigma$, compare the source arrow at $\theta$ with the
target arrow at $r(\theta)$. This is the defining realization equation
for reparametrization and again requires finitely many comparisons.
\end{enumerate}
\end{proof}

Let $X$, $\ell:X\to L$, and $(R_a\subseteq X^2)_{a\in\mathcal A}$
be the retained carrier, observable data, and designated relations of
Construction~\ref{con:egt-retained-bisimulation}. Write $B_*=\nu\mathcal F$
for its greatest retained-label bisimulation. The predicates
$P_v=\ell^{-1}\{v\}\hookrightarrow X$, $v\in\ell[X]$, are
inverse images of ordinary equality predicates in the already
constructed doctrine. Let $\mathsf{ML}_{\ell,R}$ be the Boolean
\emph{semantic test language} built from these equality-fibre tests,
$\top$, Boolean operations, and
$\langle a\rangle Q=\exists_{p_1}p_2^*Q$ for $R_a$.
As usual, $[a]Q=\neg\langle a\rangle\neg Q$. Each $P_v$ may be
regarded as a parameterized equality test at the retained value $v$;
a closed syntactic formula requires a term naming that value. This
distinction will be maintained in the expressive-completeness result.

\begin{theorem}[Modal characterization of the exact retained refinement]
\label{thm:egt-hennessy-milner-agency}
If each successor set $\{y:xR_ay\}$ is finite, then, for $x,y\in X$,
\[
 x\,B_*\,y
 \quad\Longleftrightarrow\quad
 \bigl(\forall\varphi\in\mathsf{ML}_{\ell,R}\bigr)
 \bigl(x\models\varphi\Longleftrightarrow y\models\varphi\bigr).
\]
Thus the coarsest exact quotient of
\cref{thm:egt-exact-visible-refinement} identifies precisely the
states indistinguishable by these retained Boolean modal tests. This
is the Hennessy--Milner separation argument \citep{HennessyMilner1980,HennessyMilner1985},
applied to the paper's displayed observation, strategic, and execution
relations.
\end{theorem}
\begin{proof}
\leavevmode
\begin{enumerate}[label=\textup{(\arabic*)},leftmargin=*,itemsep=4pt,topsep=3pt]
  \item \textbf{Invariance under retained bisimulation.} The relation
$B_*$ preserves $\ell$ and matches $R_a$-successors in both directions.
Therefore each equality-fibre test $P_v$ is invariant. Boolean
operations preserve invariance, and the two successor clauses show
that the diamond (hence also the box) preserves it. Induction on
modal tests proves the forward implication.
  \item \textbf{Separating an unmatched successor.} Let $M$ relate states
satisfying exactly the same tests. The tests $P_v$ give
$M\subseteq\ker(\ell)$. Suppose $xMy$ and $xR_ax'$, but no
$a$-successor $y'$ of $y$ satisfies $x'My'$. For each of the finitely
many such $y'$ choose a modal test $\varphi_{y'}$ true at $x'$ and
false at $y'$. Boolean negation permits this orientation. Then
\[
 \psi=\bigwedge_{yR_ay'}\varphi_{y'}
\]
is true at $x'$ and false at every $a$-successor of $y$; the
empty conjunction is $\top$. Thus $x\models\langle a\rangle\psi$
but $y\not\models\langle a\rangle\psi$, contradicting $xMy$.
  \item \textbf{Coinduction.} Interchanging $x,y$ gives the reverse
successor clause. Consequently $M\subseteq\mathcal F(M)$.
Coinduction for the greatest fixed point gives $M\subseteq B_*$,
completing the converse.
\end{enumerate}
\end{proof}

The characteristic formulas below belong to the semantic language of
retained equality tests: $P_{\ell(x)}$ is available as a predicate
interpreting equality with the displayed label $\ell(x)$. Their finite
expressivity is relative to those named semantic tests. A separately
presented first-layer signature denotes such a formula only when it
contains a term naming the corresponding retained label; no general
surjectivity from freely generated terms onto semantic labels is used.

\begin{theorem}[Finite-depth characteristic formulas and full expressivity]
\label{thm:egt-finite-modal-completeness}
Suppose $X$ and $\mathcal A$ are finite. Define
\[
 E_0=\ker(\ell),\qquad E_{n+1}=\mathcal F(E_n).
\]
If $X$ is nonempty, this sequence stabilizes at $B_*$ after at most
$|X|-1$ strict refinement steps. In particular, $B_*$ is definable
without a fixed-point binder by finitely many instances of ordinary
equality, the designated relation tests, and their quantifier adjoints.

More explicitly, for each $x\in X$ put
\begin{align*}
 \chi_x^0&=P_{\ell(x)},\\
 \chi_x^{n+1}
 &=P_{\ell(x)}\ \wedge
 \bigwedge_{a\in\mathcal A}
 \left(
   \bigwedge_{xR_az}\langle a\rangle\chi_z^n
   \ \wedge\
   [a]\Bigl(\bigvee_{xR_az}\chi_z^n\Bigr)
 \right).
\end{align*}
Empty conjunctions and disjunctions mean $\top$ and $\bot$.
For all $x,y,n$,
\[
 y\models\chi_x^n\quad\Longleftrightarrow\quad (x,y)\in E_n.
\]
If $m$ is the first stable index, then $\chi_x^m$ defines the
$B_*$-class of $x$. Every $B_*$-invariant predicate
$P\hookrightarrow X$ is consequently denoted by the finite modal
expression $\bigvee_{x\in P}\chi_x^m$.
In particular, evaluation induces an isomorphism of Boolean algebras
\[
 \mathsf{ML}_{\ell,R}/{\equiv_X}\ \cong\ \mathsf{Pred}(X/B_*),
\]
where $\equiv_X$ denotes equality of denotations in this fixed
finite model. This is expressive completeness of the semantic test
language, not proof-theoretic completeness of the freely generated
first-layer syntax.
For $X=\varnothing$ the analogous statement holds with the unique
predicate and $\bot$.
\end{theorem}
\begin{proof}
\leavevmode
\begin{enumerate}[label=\textup{(\arabic*)},leftmargin=*,itemsep=4pt,topsep=3pt]
  \item \textbf{Finite bisimulation refinement.} For an equivalence
relation $E$, let $[z]_E$ denote its equivalence class. Then
$\mathcal F(E)$ is the kernel equivalence of the map
\[
  x\longmapsto
  \left(\ell(x),\left(\{[z]_E:xR_az\}\right)_{a\in\mathcal A}\right).
\]
Indeed, equality of the reached class sets expresses both successor
matching clauses. Hence $\mathcal F(E)$ is an equivalence relation.
Since $E_1\subseteq E_0$ and $\mathcal F$ is monotone, the sequence
$E_n$ descends. On nonempty finite $X$ every strict refinement splits
at least one class, so there are at most $|X|-1$ strict refinements.
  \item \textbf{Identification with the greatest fixed point.} If
$E_m=E_{m+1}$, then $E_m$ is a fixed point and therefore
$E_m\subseteq B_*$. Conversely $B_*\subseteq E_0$ by retained-label
preservation, and monotonicity inductively gives $B_*\subseteq E_n$
for all $n$. Hence $E_m=B_*$. An explicit first-order definition is
obtained recursively from $E_0(x,y)\equiv\ell(x)=\ell(y)$ by
\[
\begin{aligned}
 E_{n+1}(x,y)\ \equiv\ {}&\ell(x)=\ell(y)\ \wedge\\
 &\bigwedge_{a\in\mathcal A}\left[
   \forall x'\,(xR_ax'\Rightarrow
      \exists y'\,(yR_ay'\wedge E_n(x',y')))\right.\\
 &\hspace{24mm}\left.\wedge\
   \forall y'\,(yR_ay'\Rightarrow
      \exists x'\,(xR_ax'\wedge E_n(x',y')))
 \right].
\end{aligned}
\]
The construction uses the existing equality, predicate, and quantifier
operations and gives $\mathcal F(E_n)$ at each stage.
  \item \textbf{Characteristic modal tests.} We prove
$y\models\chi_x^n\iff(x,y)\in E_n$ by induction on $n$. At $n=0$
it says precisely $\ell(x)=\ell(y)$. For the successor case, the
diamond conjunction states that every $a$-successor of $x$ has an
$E_n$-matching successor of $y$; the box states that every
$a$-successor of $y$ has an $E_n$-matching successor of $x$.
These are exactly the clauses defining $E_{n+1}$. At stabilization,
$\chi_x^m$ therefore denotes the $B_*$-class of $x$.
  \item \textbf{Classification of invariant predicates.} A predicate
$P\subseteq X$ is $B_*$-invariant precisely when it is a union of
$B_*$-classes. Therefore $\bigvee_{x\in P}\chi_x^m$ denotes $P$.
All modal tests are $B_*$-invariant by the previous theorem, so their
interpretations descend to predicates on $X/B_*$. Passing to tests
modulo equality of denotations gives an injective Boolean-algebra map
to $\mathsf{Pred}(X/B_*)$, and the characteristic formulas make it
surjective. When $X=\varnothing$, both Boolean algebras have one
element, represented by the denotation of $\bot$.
\end{enumerate}
\end{proof}

\begin{corollary}[Finite-model elimination of modal fixed points]
\label{cor:egt-finite-fixedpoint-elimination}
On a fixed finite retained relation frame, every modal fixed-point
predicate built from the retained atoms and relations has the same
denotation as a finite Boolean modal expression without $\mu$ or $\nu$.
For any monotone $F:\mathsf{Pred}(X)\to\mathsf{Pred}(X)$ given by the
displayed assertion constructors, its least and greatest fixed points
are also obtained by iteration from $\bot$ and $\top$, respectively,
after at most $|X|$ strict changes.
\end{corollary}
\begin{proof}
\leavevmode
\begin{enumerate}[label=\textup{(\arabic*)},leftmargin=*,itemsep=4pt,topsep=3pt]
  \item \textbf{Descent of modal fixed-point predicates.}
By \cref{thm:egt-exact-visible-refinement}, the quotient by $B_*$ is
surjective and bounded for each retained relation, and the designated
atoms descend. The modal fixed-point transport theorem
\ref{thm:egt-modal-adequacy} consequently makes every such fixed-point
denotation $B_*$-invariant. The characteristic-formula theorem then
expresses it by a finite Boolean modal test without fixed-point binders.
  \item \textbf{Finite monotone iteration.} For any monotone $F$ on
$\mathsf{Pred}(X)$, the iteration $\bot\leq F(\bot)\leq
F^2(\bot)\leq\cdots$ is increasing, while
$\top\geq F(\top)\geq F^2(\top)\geq\cdots$ is decreasing.
Each strict change adds or removes at least one state, so a chain
stabilizes after at most $|X|$ strict changes. At stabilization the
value is a fixed point; induction and coinduction identify it with
$\mu F$ and $\nu F$, respectively.
\end{enumerate}
\end{proof}

\begin{remark}[Why the fixed-model result is not unrestricted syntactic
completeness]
\label{rem:agt-fixed-completeness-scope}
These results separate all denotationally distinct open agent maps and
characterize every finite-state bisimulation-invariant predicate, but
they do not infer definitional equality of arbitrary primitive terms
from equality in the fixed concurrent model. In the interpreted
finite-data signature of
Construction~\ref{agt-con-finite-generators}, the primitive name for
the response $\widehat{\mathrm{id}_X}:D(X)\to D(X)$ denotes copycat,
which is also the interpretation of the structural identity. When that
name is freely declared, the existing rules need not equate its term
with the identity term. Indeed, the split indexed linear--nonlinear
model with both ordinary and linear fibres $\Set^{\Gamma_{\mathrm{disc}}}$,
Cartesian tensor, and $\Lin=\Val=1$ may interpret $d_X$ as a two-element
set and the free response name as a nonidentity endomap. Soundness
then rules out derivability of their equation in the free theory.
The modal Boolean-algebra isomorphism above is likewise modulo
\emph{equality of denotations in the specified finite model}, not an
unproved assertion about equations derivable in all indexed models.
In particular, the concrete assertion fibres are Boolean whereas the
listed logical inference rules are formulated through Heyting operations;
a proof-theoretic completeness claim would also need to derive the
corresponding fixed-model Boolean validities from the stated rules.
\end{remark}

\subsubsection{The different meanings of completeness}

Several completeness results must be distinguished. The representation
theorem \ref{agt-thm-representation} identifies semantic agent states
with owned concurrent realizations. The finite causal protocol theorem
\ref{agt-thm-protocol} gives definability for that generated fragment;
Theorem~\ref{agt-thm-session} retains its finite session realizations
and witnesses. The classifying-model theorem
\ref{agt-thm-completeness} establishes derivability from equality in
all the designated first-layer models. In contrast,
\cref{thm:agt-concurrent-probe-separation} separates arrows in the
single fixed presheaf model, and
\cref{thm:egt-hennessy-milner-agency,thm:egt-finite-modal-completeness}
characterize modal indistinguishability and finite-state expressivity
relative to the retained tests. Neither conclusion equates all
fixed-model denotational equalities with derivable syntactic equations.

Classifying-model completeness ranges over a class of indexed
linear--nonlinear models. Definability and equality reflection in a
fixed model are different questions. For instance, two unrelated
primitive scalar constants can remain distinct in the free syntax
while both necessarily denote the unique element of $\cC(I,I)$.
Remark~\ref{rem:agt-fixed-completeness-scope} exhibits the same
obstruction using a named finite response and copycat. The operational
separation and modal results characterize fixed-model semantic equality
at their specified interfaces; they do not prove derivability of those
equalities. Full completeness in particular causal-invariant linear
fragments \citep{agt-CastellanYoshida} is not transferred here to all
receptive courteous strategies.

What has been established is the required sequence of exhibited
interfaces: first-class agent types with composed owner policies,
controlled and legally informed alternatives, witnessed institutional
and communication processes, perspective compatible with linear
composition, residual temporal behaviour, and exact reconstruction.
Finite causal guards certify dependence on actual authorization events;
normative authority and economic preference remain separately displayed.
Each comparison is expressed through its designated operational,
strategic, and epistemic interfaces.

\section*{Conclusion}
\addcontentsline{toc}{section}{Conclusion}

The construction connects two accounts of agency that otherwise
leave complementary information unspecified. Event-structure strategies
record what can occur concurrently and what depends causally on earlier
events. Owner-indexed complete policies record which contingent
alternatives an agent controls. Their representation in the completed
linear semantics is compositional, and the same displayed factors
support deviations, preference-sensitive selection, observational
admissibility, legal power, and witnessed institutional action.
Communication and continuation are not merely added modalities: finite
protocols and residual strategies exhibit their execution data, while
alternative-lifting and refinement results identify the information
preserved by the relevant abstractions. The resulting calculus is sound,
and its first-layer equations are complete for their classifying-model
semantics. Original concurrent arenas also separate open-agent arrows,
while finite-state retained bisimulation has an exact characterization
by modal tests and finite characteristic formulas. Thus operational
interaction, economic control, and epistemic change can be reasoned
about through distinct but explicitly connected categorical interfaces.

\paragraph{Scope and interfaces.}\label{sec:egt-boundary}
The conclusions depend on retaining the displayed information appropriate
to each question. Yoneda recovers the original concurrent arrows;
owner factors determine control; legal and observation quotients
determine admissible alternatives; and execution spans preserve the
witnesses needed for subsequent action. The exact comparison and
reconstruction statements of
\Cref{thm:egt-announcement-policy-comparison,thm:fib-reconstruction,agt-thm-reconstruction}
identify what survives each passage. The following remarks describe how these retained interfaces connect
to economic and computational extensions.

\begin{remark}[Retained economic interfaces]
The linear semantics determines how exhibited interactions compose.
Ownership, observations, comparison predicates, technologies, and
institutional rules are explicitly typed economic data. Concurrent
strategies determine event enablement and residualization; observation
wires and retained pointed histories determine the displayed action
uncertainty. The resulting predicates express equilibrium, ability,
rationalizability, feasibility, coalition blocking, and truthful
implementation for the specified economic presentation. The
Cartesian-frame, local-state, and choice-cell comparisons identify their
respective interfaces. Persistent intentions and corporate judgements
can be represented through explicit commitment and aggregation data
retained by the same presentation.
\end{remark}

\begin{remark}[Probability and intensional interfaces]
The policy sets and private policy records used here are semantic values;
their finite realizations are extensional response protocols. Models of
open-source games additionally represent program descriptions, permitting
self-reference and inspection of another agent's strategy description.
Interactive diagonalization is studied by \citet{AbramskyZvesper2010}, and
provability-based equilibria by \citet{BaraszEtAl2014}. Rational synthesis,
as studied by \citet{FismanEtAl2009}, asks for implementations satisfying
strategic specifications.

The present theory uses pure policies and qualitative information.
Probabilistic extensions require distributions and chance kernels, as in
Bayesian games \citep{Harsanyi1967} and stochastic games
\citep{Shapley1953}. The disclosure and action-model constructions here
have the formula semantics and reduction laws of
Appendix~\ref{app:rules-assertions}. Observational distinctions that do
not descend to the original visible quotient are retained on the exact
refinement and its dependent operational witness family.

The completion, linearization, and dependent adjoints are set-theoretic
categorical constructions, and least and greatest fixed points are
order-theoretic operations on predicates. Finite path maxima construct
the potentials of \cref{thm:egt-finite-improvement}; directization constructs
$r_s=m\prod_i s_i$ with the universal comparison of
\cref{prop:egt-direct-implementation-initial}.
\end{remark}


\clearpage
\appendix

\section{Natural-Deduction Presentation of the Logic of Agency}
\label{app:rules}

This appendix specifies the ordinary, linear, and assertion judgments of
the agency calculus. All semantic and operational constructions and their
mathematical proofs are developed in the main text. The soundness and
classifying-model proofs for these rules are in
Appendix~\ref{app:metatheory}.

\subsection{Ordinary and linear agent type theory}
\label{app:rules-types}


The proof-theoretic fragment is interpreted through the
linear--nonlinear adjunction of \citet{Benton1994MixedLNL} and the
indexed dependent setting of \citet{Vakar2015LinearFrameworks}.
The rule tables preserve the distinction between ordinary variables,
which may be reused, and linear variables, which must be accounted for
by context splitting.

The judgments below are interpreted in the indexed ordinary and
linear categories already constructed. Their complete inference rules
make explicit the context-splitting conditions used in the soundness
induction; the displayed agent operations are derived from those rules
rather than introduced as unrelated primitives.

\subsubsection{Judgments of the agent calculus}
\label{agt-sec-calculus}

\paragraph{Ordinary parameters and linear resources}

The ordinary fragment is extensional dependent type theory with finite
products and coproducts, dependent sums and products, and equality. Ordinary
types and terms have judgments
\[
 \Gamma\ordturn X\ \type,
 \qquad \Gamma\ordturn u:X.
\]
An ordinary context is built by comprehension $\Gamma,x:X$. Linear families
may depend on ordinary variables:
\[
 \Gamma\vdash A\ \ltype,
 \qquad \Gamma;\Delta\linturn t:A.
\]
Here $\Delta=x_1:A_1,\ldots,x_n:A_n$ is a linear context whose
families are all indexed by $\Gamma$. Each $A_k$ is indexed by the ordinary context. A later linear type does not
depend on a preceding linear variable. Ordinary dependence is therefore
available in type formation, while resource occurrences remain governed
by tensor-context use.

The linear type constructors are
\[
 \I,\quad A\otimes B,\quad A\multimap B,\quad
 A\mathbin{\&}B,\quad\Top,\quad A\oplus B,\quad\botell,
 \quad\Sigma_x^\ell A,\quad\Pi_x^\ell A,\quad\Lin X.
\]
The ordinary family $\Val A$ records monoidal states of the
linear family $A$. The adjunction $\Lin\dashv\Val$ connects the
two kinds of judgments, and $!A$ abbreviates $\Lin\Val A$, the
resulting exponential. These rules do not supply an unrestricted
contraction on arbitrary linear terms.

\paragraph{The multiplicative and adjunction rules}

The essential multiplicative rules are
\[
 \frac{}{\Gamma;x:A\linturn x:A},\qquad
 \frac{\Gamma;\Delta\linturn t:A\quad
       \Gamma;\Theta\linturn u:B}
      {\Gamma;\Delta,\Theta\linturn t\otimes u:A\otimes B},
\]
\[
 \frac{\Gamma;\Delta,x:A\linturn t:B}
      {\Gamma;\Delta\linturn\lambda x.t:A\multimap B},
 \qquad
 \frac{\Gamma;\Delta\linturn f:A\multimap B\quad
       \Gamma;\Theta\linturn t:A}
      {\Gamma;\Delta,\Theta\linturn f(t):B}.
\]
Each multiplicative rule partitions its linear premises into
disjoint contexts. Tensor elimination consumes its argument and binds
both components, while the unit is introduced only with the appropriate
empty linear context. These conditions implement the resource-sensitive
meaning of the connectives.

Ordinary terms are introduced into $\Lin X$ by
\[
 \frac{\Gamma\ordturn u:X}{\Gamma;\varnothing\linturn
        \mathsf{return}(u):\Lin X}.
\]
Elimination of $\Lin X$ binds an ordinary $x:X$ inside the continuation:
\begin{equation}
 \frac{\Gamma;\Delta\linturn s:\Lin X\quad
       \Gamma,x:X;\Theta\linturn t:A}
      {\Gamma;\Delta,\Theta\linturn
       \mathsf{let}\ \mathsf{return}(x)=s\ \mathsf{in}\ t:A}.
 \label{agt-eq-Felim}
\end{equation}
The displayed instance has $\Theta$ and $A$ pulled back from
$\Gamma$. Dependent result types are expressed by the indexed
sum $\Sigma^\ell$. Elimination binds an ordinary parameter that may
be reused in its continuation, while the supplied linear term $s$ is
consumed exactly once.

The value rules are
\[
 \frac{\Gamma;\varnothing\linturn t:A}
      {\Gamma\ordturn\mathsf{thunk}(t):\Val A},\qquad
 \frac{\Gamma\ordturn v:\Val A}
      {\Gamma;\varnothing\linturn\mathsf{force}(v):A}.
\]
The value rules express the hom-set adjunction by their computation
and uniqueness laws. The introduction rule requires an empty linear
context: a term depending on a linear resource cannot thereby be
turned into reusable ordinary information.

Definitional equality consists of the displayed $\beta$ and
$\eta$ equations for tensor, unit, internal hom, additives,
dependent sums and products, and the adjunction. Its commuting
conversions make eliminators natural in their continuations. The
soundness proof checks each family of equations against the
corresponding categorical universal property.

\paragraph{Indexed linear abstraction}

For $\Gamma\ordturn X\ \type$ and
$\Gamma,x:X\vdash A\ \ltype$, the dependent product introduction rule is
\[
 \frac{\Gamma,x:X;p^*\Delta\linturn t:A}
      {\Gamma;\Delta\linturn\Lambda x.t:\Pi_x^\ell A},
 \qquad p:\Gamma.X\to\Gamma.
\]
Application selects one ordinary index from the dependent
product. The same linear resource may be mentioned in alternative
components of an additive product, but it cannot be extracted twice
as simultaneous linear occurrences. The distinction is essential for
the agent abstraction derived from the dependent product.

\begin{lemma}[Substitution]
\label{agt-lem-substitution}
Ordinary substitution preserves all formation, typing, and equality
judgments. If
\[
 \Gamma;\Delta\linturn s:A,
 \qquad \Gamma;\Theta,x:A\linturn t:B,
\]
then
$\Gamma;\Theta,\Delta\linturn t[s/x]:B$.
The substitution uses the resources of $\Delta$ once in the multiplicative
sense and preserves definitional equality.
\end{lemma}
\begin{proof}
\leavevmode
\begin{enumerate}[label=\textup{(\arabic*)},leftmargin=*,itemsep=4pt,topsep=3pt]
  \item \textbf{Fixed-point argument.} Proceed by induction on the derivation into which substitution is made. For a variable, substitute $s$ if it is $x$, and retain the variable otherwise. At a multiplicative rule, $x$ belongs to exactly one premise; apply the
induction to that premise and retain the other disjoint context.
  \item \textbf{Fibrewise calculation.} At an
additive introduction or case rule, the same context is present in the
alternative premises, so substitute in each alternative, retaining one
shared context in the conclusion. Under an ordinary binder, rename its
variable to avoid capture and pull $s$ back along comprehension. Under a
linear binder, rename that variable and apply the induction to the body.
  \item \textbf{Fixed-point argument.} The $\Val$ introduction case has empty linear context, so it cannot contain
a free $x$. The $\Lin$ eliminator uses the two induction steps for its split
premises. These exhaust the constructors.
  \item \textbf{Functoriality and substitution.} Each defining equation is stable
under the same capture-avoiding substitutions, proving the equality claim.
\end{enumerate}
\end{proof}

\begin{corollary}[Preservation for the core reductions]
Every well-typed $\beta$ reduction and every typed commuting conversion
preserves its judgment.
\end{corollary}
\begin{proof}
\leavevmode
\begin{enumerate}[label=\textup{(\arabic*)},leftmargin=*,itemsep=4pt,topsep=3pt]
  \item \textbf{Secured operational interaction.} The internal-hom and dependent elimination cases are instances of
Lemma~\ref{agt-lem-substitution}. Tensor and unit reduction consume precisely
the components introduced by their matching constructors. Additive reduction
selects one branch with its declared context.
  \item \textbf{Adjunction and transposition.} Adjunction reduction substitutes
the ordinary value supplied by $\mathsf{return}$ or removes a matching
$\mathsf{thunk}$--$\mathsf{force}$ pair. Commuting conversions rearrange the
same context splits and bindings.
\end{enumerate}
\end{proof}

\begin{example}[Linear occurrences and ordinary descriptions]
Given $a:R\multimap A$, $b:R\multimap B$, and one linear occurrence
$x:R$, the expression $a(x)\otimes b(x)$ is not derivable: tensor
introduction would place $x$ in both split premises. An exhibited map
$R\to R\otimes R$, or a supplied input $R\otimes R$, makes the needed
resource available. An ordinary identifier of $R$ supplies neither.
There is an important distinction for an actual value $v:\Val R$:
ordinary contraction together with the value and tensor rules derives
\[
 \Gamma,v:\Val R;\varnothing\linturn
 \mathsf{force}(v)\otimes\mathsf{force}(v):R\otimes R.
\]
This constructs two states from a reusable ordinary state description;
it does not contract a linear variable. Consequently linear typing alone
does not establish unforgeability of an authorization token. The
causally token-guarded protocol of Lemma~\ref{agt-lem-token-guard}
provides the operational dependency needed in the finite-data fragment.
\end{example}

\paragraph{Complete rule tables and formation conventions}
\label{agt-app-rules}

\paragraph{Binding and formation conventions}

The rules are schematic under ordinary substitution and type
conversion. Bound variables are chosen fresh and substitutions avoid
capture. Multiplicative context concatenation requires disjoint
variable sets; a rule over $\Gamma,x:X$ pulls back any context or target
originally over $\Gamma$ along its projection. A resource may appear
in alternative additive premises, but this does not allow simultaneous
duplication. These are the binding conventions
used in the later soundness induction.

Ordinary formation is generated from the declared sorts by
$1,0,X+Y,\Sigma_{x:X}Y,\Pi_{x:X}Y$, and extensional equality.
The ordinary rules are variable, weakening, contraction, exchange,
substitution, unit introduction, empty elimination, sum injections and case,
dependent pairing and projections, abstraction and application, and equality
reflexivity and substitution. Their $\beta\eta$ equations are the
evaluation and uniqueness equations for these constructors.
Ordinary equality is \emph{extensional}: the ordinary foundation presents
equality by comprehension of the diagonal. Its eliminator and conversion rule
are those of this extensional ordinary type theory.  Semantically, a section
of a diagonal pullback forces its two compared sections to be equal, and
monicity makes its equality witness unique.  The proof is given in
Proposition~\ref{agt-prop-equality-diagonal}.  This is a genuine feature of
the selected ordinary foundation, not a consequence of products, sums, or
dependent products and not an additional law of agency.  No equality
reflection for arbitrary linear resources is introduced.

The mixed formation rules are
\[
 \frac{\Gamma\ordturn X\ \type}{\Gamma\vdash\Lin X\ \ltype},
 \qquad
 \frac{\Gamma\vdash A\ \ltype}{\Gamma\ordturn\Val A\ \type},
\]
\[
 \frac{\Gamma\vdash A,B\ \ltype}
      {\Gamma\vdash A\otimes B,\ A\multimap B,\ A\mathbin{\&}B,
                   \ A\oplus B\ \ltype},
\]
\[
 \frac{\Gamma\ordturn X\ \type\qquad\Gamma,x:X\vdash A\ \ltype}
      {\Gamma\vdash\Sigma_x^\ell A,\ \Pi_x^\ell A\ \ltype}.
\]
The nullary constructors are $\I,\Top,\botell$. For every displayed
signature $(\sig;A,B)$, $\Ag_{\sig}(A,B)$ abbreviates
$[\Lin(\prod_i\Sigma_i)\otimes A,B]$ with that signature retained.

\paragraph{The remaining linear introduction and elimination rules}

Tensor elimination and the multiplicative unit are
\[
 \frac{\Gamma;\Delta\linturn s:A\otimes B\qquad
       \Gamma;\Theta,x:A,y:B\linturn t:C}
      {\Gamma;\Delta,\Theta\linturn
       \mathsf{let}\ x\otimes y=s\ \mathsf{in}\ t:C},
\]
\[
 \frac{}{\Gamma;\varnothing\linturn *:\I},\qquad
 \frac{\Gamma;\Delta\linturn s:\I\qquad\Gamma;\Theta\linturn t:C}
      {\Gamma;\Delta,\Theta\linturn
       \mathsf{let}\ *=s\ \mathsf{in}\ t:C}.
\]
Additive products have
\[
 \frac{\Gamma;\Delta\linturn t:A\qquad\Gamma;\Delta\linturn u:B}
      {\Gamma;\Delta\linturn\langle t,u\rangle:A\mathbin{\&}B},
 \qquad
 \frac{\Gamma;\Delta\linturn t:A\mathbin{\&}B}
      {\Gamma;\Delta\linturn\pi_1t:A},
\]
and the analogous second projection. The nullary additive product has a
unique term $\Gamma;\Delta\linturn\langle\rangle:\Top$ at every
linear context. This is the terminal-object rule, not weakening into an
arbitrary target.

For additive sums, either injection uses its premise's context, and
\[
 \frac{\Gamma;\Delta\linturn s:A\oplus B\qquad
       \Gamma;\Theta,x:A\linturn t:C\qquad
       \Gamma;\Theta,y:B\linturn u:C}
      {\Gamma;\Delta,\Theta\linturn
       \mathsf{case}\ s\ \mathsf{of}\ (x.t\mid y.u):C}.
\]
A term of $\botell$ eliminates into any target, with an additional linear
context allowed in that impossible branch. Semantically this is the unique
map $0\otimes X\to C$, using $0\otimes X\cong0$.

For dependent linear sums, with $p:\Gamma.X\to\Gamma$,
\[
 \frac{\Gamma\ordturn u:X\qquad\Gamma;\Delta\linturn t:A[u/x]}
      {\Gamma;\Delta\linturn(u,t):\Sigma_x^\ell A},
\]
\[
 \frac{\Gamma;\Delta\linturn s:\Sigma_x^\ell A\qquad
       \Gamma,x:X;p^*\Theta,z:A\linturn t:p^*B}
      {\Gamma;\Delta,\Theta\linturn
       \mathsf{let}\ (x,z)=s\ \mathsf{in}\ t:B}.
\]
Dependent product abstraction is in Section~\ref{agt-sec-calculus}; its
application is
\[
 \frac{\Gamma;\Delta\linturn f:\Pi_x^\ell A\qquad
       \Gamma\ordturn u:X}
      {\Gamma;\Delta\linturn f[u]:A[u/x]}.
\]
The ordinary--linear rules are those of
Section~\ref{agt-sec-calculus}. Each primitive linear generator
consumes the tensor of its declared input resources, with the indicated
context split. No additional copying rule is attached to a primitive
generator.

The defining equations include
\begin{align*}
 (\lambda x.t)(u)&=t[u/x],&\lambda x.f(x)&=f,\\
 \pi_1\langle t,u\rangle&=t,&\langle\pi_1s,\pi_2s\rangle&=s,\\
 (\Lambda x.t)[u]&=t[u/x],&\Lambda x.f[x]&=f,\\
 \mathsf{force}(\mathsf{thunk}(t))&=t,&
 \mathsf{thunk}(\mathsf{force}(v))&=v.
\end{align*}
Tensor, dependent-sum, additive-sum, unit, and $\Lin$ computation
laws identify elimination immediately following
the corresponding introduction. Their uniqueness counterparts say
that a value is determined by its projections or case components; in
particular, maps out of sums are determined by all branches.
For the mixed ordinary--linear eliminator, the fundamental equations
(with the context conditions of~\eqref{agt-eq-Felim}) include
\[
\begin{aligned}
 \mathsf{let}\ \mathsf{return}(x)=\mathsf{return}(u)
   \ \mathsf{in}\ t&=t[u/x],\\
 \mathsf{let}\ \mathsf{return}(x)=z
   \ \mathsf{in}\ \mathsf{return}(x)&=z
       \qquad(z:\Lin X).
\end{aligned}
\]
The remaining commuting conversions express naturality of the
eliminators in their linear continuations. If $E_s(t)$ eliminates a
constructor whose parameter is $s$ and continuation is $t$, and $k$
follows the result, then eliminating into the
composite continuation agrees with eliminating first and composing
with $k$, with identical context splitting. These equations include the coproduct and $\Lin$ commuting
conversions used by the classifying-model construction. Each equation
is checked under the same typed substitution rules. The value equations
retain the empty linear context required by $\mathsf{thunk}$. This completes the formal
first-layer equations used by the classifying model.


\subsection{Strategic, epistemic, institutional, and temporal assertions}
\label{app:rules-assertions}


The assertion calculus is generated from the quantifier adjunctions
in \citet{Jacobs1999} and the complete-lattice fixed-point
construction of \citet{Tarski1955}. The rules below record which
modalities are derived from which maps, so that ordinary logical
reasoning and linear execution remain separate judgments.

There are two coordinated languages: assertions over ordinary
state contexts and linear terms over those contexts. For a situated
agent with complete policy carrier $\Sigma$, legal carrier $L$,
reference map $s:W\to L$, observation quotients
$o_i:W\twoheadrightarrow O_i$, and comparison predicates $P_i$, the
clauses of the predicate language are interpreted in $\mathsf{Pred}(W)$ using exactly those carriers
and their typed relations. Its formulas therefore inherit the available alternatives
from the declared game rather than treating control as an untyped
modal primitive.

\subsubsection{Assertion and strategic judgments}
Write $W\mid\Phi\vdash\varphi$ when
$\bigcap_{\psi\in\Phi}\llbracket\psi\rrbracket\subseteq
\llbracket\varphi\rrbracket$.  Ordinary equality, conjunction,
disjunction, implication, quantifiers, and substitution have the extensional
subobject interpretation.  For a coalition $C$ and the kernel pair
$D_C^L=L\times_{L_{-C}}L$, the alternatives from a reference are its
fibres under the first projection.  The existential and universal
counterfactual operations are
\[
\langle C\rangle Q=\exists_{p_1}p_2^*Q,
\qquad
[C]Q=\forall_{p_1}p_2^*Q.
\]
Selection and best response are obtained by pulling back
strict-improvement comparisons along paired concurrent realizations.
The predicate $\mathsf{NE}$ universally excludes profitable unilateral
alternatives. Coalition effectivity uses a different order of quantifiers:
it chooses a realizable coalition coordinate and tests every legal
completion. The formal rules retain that distinction.

\subsubsection{Knowledge, perspective, and common knowledge}
For $o_i:W\twoheadrightarrow O_i$ put
\[
K_iP=o_i^*\forall_{o_i}P,\qquad
D_CP=q_C^*\forall_{q_C}P,
\]
where $q_C$ is the effective image of the tuple of the members'
observations.  Let $R_C^*$ be the finite-path equivalence closure of
$\bigcup_{i\in C}\ker(o_i)$.  Then
\[
\mathsf{CK}_C P=[R_C^*]P
=\nu X.\left(P\wedge\bigwedge_{i\in C}K_iX\right).
\]
The knowledge rules follow from the observation quotient and its
kernel pair. The proof-relevant linear operation
$K_{o_i}^{\ell}=o_i^*\Pi_{o_i}^{\ell}$ is separately interpreted
as a comonad. Its predicate reflection yields the stated knowledge
operator, without identifying their witness products.

\subsubsection{Institutional and epistemic action formulas}
For a witnessed span $X\xleftarrow{s}E\xrightarrow{t}Y$, the dynamic
clauses are
\[
\langle E\rangle Q=\exists_s t^*Q,\qquad
[E]Q=\forall_s t^*Q.
\]
A public announcement of $\varphi$ has updated carrier
$j:U=\{w: \varphi(w)\}\hookrightarrow W$ with relations
$(j\times j)^*R_i$ and box $[!\varphi]\psi=\forall_j
\llbracket\psi\rrbracket_{M|\varphi}$.  A general epistemic
action $(E,(S_i),\mathsf{pre})$ has updated carrier
$U=\{(w,e):\mathsf{pre}(w,e)\}$ and relation
$(w,e)R_i^+(v,f)$ exactly when $wR_iv$ and $eS_if$.
All postformulas are evaluated in the updated model.

\subsubsection{Fixpoint and temporal rules}
Least and greatest fixed points, denoted by $\mu$ and $\nu$,
are evaluated in the complete
predicate lattice of the current context. Induction uses the least
prefixpoint property, while coinduction uses the greatest postfixpoint
property. Temporal maintenance and achievement apply these operations
to the controllable predecessor supplied by the finite public
transition game. The formulas do not create transition witnesses; those
are supplied by the operational construction.

\subsubsection{Formal reductions and soundness obligations}
Public announcement, epistemic action, substitution, and modal
reduction are consequences of comprehension, pullback, and the
dependent adjunctions. The displayed equations fix the interpretation
of these operations for the soundness proof; no independent reduction
principle is added at the level of concurrent arrows.
\[
[!\varphi]K_i\psi\leftrightarrow
(\varphi\Rightarrow K_i[!\varphi]\psi),\qquad
[!\varphi][!\psi]\chi\leftrightarrow
[!(\varphi\wedge[!\varphi]\psi)]\chi.
\]
Each assertion rule is sound when it denotes an inclusion of
subobjects, and each typed term equation is sound when it denotes
equality of arrows. Appendix~\ref{app:metatheory} verifies these obligations
by induction on derivations.

\subsubsection{Types, formulas, and contexts}

Fix a displayed agent process $G=(\vec\Sigma,g):A\to B$, a legal
comprehension $m:L\hookrightarrow\Sigma$, and a reference context
$s:W\to L$.  Its sorts include ordinary state and policy families,
linear families over those states, the concurrent arrow carrier
$\mathcal C(A,B)$, and dependent witness families.  An ordinary term is
interpreted as a section of an ordinary family over its declared
context; a linear term is interpreted as a morphism from the tensor of
its linear context.  A formula in context $W$ denotes a subobject
$\llbracket\varphi\rrbracket\hookrightarrow W$.  For each ordinary
context map $f:V\to W$ there is a substitution operation
$f^*: \mathsf{Pred}(W)\to\mathsf{Pred}(V)$.

The logical grammar may be given schematically by
\[
\begin{aligned}
\varphi,\psi ::= {}&\top\mid\bot\mid P(t_1,\ldots,t_n)
 \mid t=u\mid(\varphi\wedge\psi)
 \mid(\varphi\vee\psi)\mid(\varphi\Rightarrow\psi)\\
 &\mid\exists_{f}\varphi\mid\forall_f\varphi
 \mid\langle C\rangle\varphi\mid[C]\varphi
 \mid K_i\varphi\mid\mathsf{DK}_C\varphi\mid\mathsf{CK}_C\varphi\\
 &\mid\langle E\rangle\varphi\mid[E]\varphi
 \mid[!\varphi]\psi\mid[\mathcal A,e]\psi
 \mid\mu X.\varphi(X)\mid\nu X.\varphi(X).
\end{aligned}
\]
Each predicate atom is typed at a declared ordinary carrier.
The coalition variable $C$ refers to actual owner coordinates; $E$
is an exhibited witness span; and $(\mathcal A,e)$ is an epistemic
event with enabling data. Fixed-point binders are formed only at
positive occurrences of the bound predicate, making the semantic
operator monotone.

For a state context $s:W\to\Sigma$, logical equality is the diagonal
of $W\times_\Sigma W$.  If a state is represented by a term
$t:V\to W$, an assertion is tested by the inverse image
$t^*\llbracket\varphi\rrbracket$.  In particular, assertions about
ownership, legality, or observations are not forced to factor through the
realized concurrent arrow: they may use the full displayed state record.

\subsubsection{Predicate inference from adjunctions}

The order structure on each predicate fibre supplies the inference
relation $\Phi\vdash_W\varphi$ by inclusion.  The ordinary connective
rules are those of its Heyting algebra, so for example
\[
R\wedge P\leq Q\quad\Longleftrightarrow\quad
R\leq(P\Rightarrow Q).
\]
The dependent quantifier introduction and elimination rules are the
following equivalences, applicable to every ordinary context map
$f:V\to W$:
\[
\boxed{
\exists_f P\leq Q\ \Longleftrightarrow\ P\leq f^*Q,
\qquad
f^*Q\leq P\ \Longleftrightarrow\ Q\leq\forall_f P.}
\]
The rules for conjunction, implication, and quantification follow
from the Heyting structure and adjunctions. Their stability under
ordinary substitution is the Beck--Chevalley comparison. In
particular, for the following pullback square, the two directions of
quantifier transport are fixed by the same fibre bijection.
\[
\begin{tikzcd}
V'\ar[r,"a"]\ar[d,"f'"']&V\ar[d,"f"]\\
W'\ar[r,"b"']&W
\end{tikzcd}
\]
one has
\[
b^*\exists_f\cong\exists_{f'}a^*,\qquad
b^*\forall_f\cong\forall_{f'}a^*,\qquad
\exists_f(P\wedge f^*Q)=\exists_fP\wedge Q.
\]
The dependent sum and product underlie ordinary existential and
universal quantification, with subterminal reflection converting
witness-bearing sums into truth values. The assertion rules therefore
reuse the indexed type operations rather than duplicating their
semantic structure.

\subsubsection{Counterfactual and selection rules}

Let $\pi_{-C}:\Sigma\to\Sigma_{-C}$ be the outsider projection.  Its
kernel pair $D_C=\Sigma\times_{\Sigma_{-C}}\Sigma$ has source and
target $p_1,p_2$.  The relation modalities are
\[
\begin{split}
\langle C\rangle P(\sigma)&\iff
\exists\tau\,[\pi_{-C}(\tau)=\pi_{-C}(\sigma)\wedge P(\tau)],\\
[C]P(\sigma)&\iff
\forall\tau\,[\pi_{-C}(\tau)=\pi_{-C}(\sigma)
 \Rightarrow P(\tau)].
\end{split}
\]
They are monotone, preserve the respective finite joins and meets,
and satisfy $\langle\varnothing\rangle P=[\varnothing]P=P$ on an
unrestricted profile carrier.  On unrestricted products the relation
identity $D_C;D_D=D_{C\cup D}$ gives
\[
\langle C\rangle\langle D\rangle P
 =\langle C\cup D\rangle P,
\qquad
[C][D]P=[C\cup D]P.
\]
With a legal carrier, replace the full product kernel pair by
$D_C^L=L\times_{L_{-C}}L$. Sequentially composed legal changes must
retain an intermediate profile lying in the same legal subobject.
The unrestricted identity for coalitional composition need not hold
without that witness.

For each owner $i$, the paired realization is
$\delta_i=(gp_1,gp_2):D_i\to\mathcal C(A,B)^2$.  If
$P_i\hookrightarrow\mathcal C(A,B)^2$ compares the second interaction
favorably with the first, define the strict-improvement test
$\delta_i^*P_i$.  The selected legal best-response predicate says that
no improvement exists against the fixed outsider profile:
\[
\mathsf{BR}_i(\sigma,\sigma_i)
\Longleftrightarrow
\forall\tau\,(
\sigma D_i^L\tau\Rightarrow\neg P_i(g\sigma,g\tau)).
\]
Equilibrium conjoins the ownerwise absence-of-improvement tests.
Under reparametrization its exact transport depends on lifting the
relevant deviations at a fixed reference, as well as preserving both
realized endpoints. Numerical utilities give a particular comparison
presentation; they do not alter this logical rule.

\subsubsection{Epistemic and informed-ability rules}

An observation quotient $o_i:W\twoheadrightarrow O_i$ induces the
relation $R_i=\ker(o_i)$.  The knowledge operator has the explicit
universal-fibre interpretation
\[
K_iP(w)\iff\forall v\in W
\,[o_i(v)=o_i(w)\Rightarrow P(v)].
\]
It preserves truth and conjunction, and the defining equivalence
relation gives
\[
K_i(P\Rightarrow Q)\leq(K_iP\Rightarrow K_iQ),
\quad K_iP\leq P,\quad K_iP=K_iK_iP,
\quad\neg K_iP\leq K_i\neg K_iP.
\]
These are the normality, factivity, and introspection properties of
an observation equivalence relation. The linear perspective
$K_{o_i}^\ell=o_i^*\Pi_{o_i}^\ell$ instead carries proof-relevant
witness products, satisfying its comonad and lax-monoidal laws.
The predicate laws are not read as equations of arbitrary linear
resources.

For a nonempty coalition $C$ the joint quotient $q_C$ is the image of
$\langle o_i\rangle_{i\in C}$, and
$\mathsf{DK}_CP=q_C^*\forall_{q_C}P$.  The common-knowledge operator
is the greatest fixed point
\[
\mathsf{CK}_CP
=\nu X.\left(P\wedge\bigwedge_{i\in C}K_iX\right).
\]
Its unfolding law and its coinduction rule have the forms
\[
\mathsf{CK}_CP=P\wedge\bigwedge_{i\in C}K_i\mathsf{CK}_CP,
\qquad
\frac{X\leq P\quad X\leq K_iX\ (i\in C)}{X\leq\mathsf{CK}_CP}.
\]
The greatest fixed point is computed using finite paths in the
union of the owners' observation relations. The joint quotient
classifies distributed knowledge but does not by itself deliver a
usable pooled signal. Section~\ref{sec:paper-11} supplies its
communication realization.

For a worldwise feasibility or permission predicate $F\hookrightarrow
W\times M$, an informed legal policy is a section of
$\forall_{o\times1_M}F$ over $O$.  The judgement
$\pi:\mathsf{Pol}_o(F)$ is equivalent to the uniformly quantified
permission formula $\forall w\,F(w,\pi(o(w)))$.
An individual or coalition guarantee uses existential quantification
over a realizable choice coordinate, followed by universal quantification
over its actual legal completions.  A separate existence conjunct is
retained for executable guarantees while a
vacuous universal assertion.

\subsubsection{Dynamic execution and public restriction}

For a witnessed span $X\xleftarrow{s}E\xrightarrow{t}Y$, put
\[
\langle E\rangle P=\exists_s t^*P,
\qquad [E]P=\forall_s t^*P.
\]
The composition of witnessed spans by pullback gives, for a second
span $F$,
\[
[F\circ E]P=[E][F]P,
\qquad\langle F\circ E\rangle P
=\langle E\rangle\langle F\rangle P.
\]
Identity spans induce identity modalities and pullback
composition supplies their sequential laws. The proof-relevant
versions use $\Sigma_s^\ell t^*$ and $\Pi_s^\ell t^*$, so two
distinct witnesses sharing one endpoint pair still contribute
separate summands or factors. Restricting $E$ to its permission predicate gives the permitted
modalities before either construction is applied.

For a public announcement $\varphi$ interpreted in model $M$,
construct its comprehensional inclusion $j:U\hookrightarrow W$ and
restrict every accessibility relation on both endpoints.  The box is
\[
\llbracket[!\varphi]\psi\rrbracket_M
=\forall_j\llbracket\psi\rrbracket_{M|\varphi}.
\]
The postcondition belongs to the restricted epistemic model, not
to the original state space. The displayed announcement reductions
follow from the dependent-product and comprehension equations and
must therefore use the updated accessibility relation.
\[
\begin{aligned}
[!\varphi]p&\leftrightarrow(\varphi\Rightarrow p),\\
[!\varphi](\psi\wedge\chi)
 &\leftrightarrow[!\varphi]\psi\wedge[!\varphi]\chi,\\
[!\varphi]\neg\psi
 &\leftrightarrow(\varphi\Rightarrow\neg[!\varphi]\psi),\\
[!\varphi]K_i\psi
 &\leftrightarrow(\varphi\Rightarrow K_i[!\varphi]\psi),\\
[!\varphi][!\psi]\chi
 &\leftrightarrow[!(\varphi\wedge[!\varphi]\psi)]\chi.
\end{aligned}
\]
At a removed state the first announcement makes the universal
condition vacuous. At a retained state, a knowledge reduction ranges
only over accessible worlds that also survive the restriction.
This is why updating both endpoints of the relation differs from
substituting the reference state alone.

\subsubsection{General epistemic events and their reduction}

For a finite event set $E$, relations $S_i$, and enabling predicate
$\mathsf{pre}\hookrightarrow W\times E$, the product-update model has
carrier $U=\{(w,e):\mathsf{pre}(w,e)\}$ with
\[
(w,e)R_i^+(v,f)
\iff wR_iv\wedge eS_if.
\]
For a designated event $e$ let $j_e:W_e\hookrightarrow W$ be its
precondition comprehension and $k_e:W_e\to U$ its insertion.  Define
\[
[\mathcal A,e]\psi
=\forall_{j_e}k_e^*\llbracket\psi\rrbracket_{M\otimes\mathcal A}.
\]
The derived finite-event knowledge reduction is
\[
[\mathcal A,e]K_i\psi
\leftrightarrow
\left(\mathsf{pre}_e\Rightarrow
\bigwedge_{f:eS_if}K_i[\mathcal A,f]\psi\right).
\]
The finite-event knowledge reduction quantifies over enabled,
epistemically compatible event alternatives through guarded boxes.
Successive product updates retain their event pairs and the
conjunction of enabling predicates. The canonical reassociation of
these pairs proves compatibility with repeated updates.

\subsubsection{Temporal induction and rationality}

For a finite public transition game with controllable predecessor
$\operatorname{Pre}_C$, the formulas for maintenance and bounded-until
ability are respectively
\[
\nu X.(P\wedge\operatorname{Pre}_C(X)),
\qquad
\mu X.(Q\vee(P\wedge\operatorname{Pre}_C(X))).
\]
For a monotone operator $F$ on a predicate lattice, its fixed-point
inference principles are
\[
\frac{F(P)\leq P}{\mu F\leq P},
\qquad
\frac{P\leq F(P)}{P\leq\nu F},
\]
The unfolding equations $F(\mu F)=\mu F$ and
$F(\nu F)=\nu F$ express the least and greatest fixed points.
Finite-state strategy extraction is a separate operational theorem:
its witnesses come from the displayed predecessor construction.
Likewise, a profitable-deviation relation becomes a temporal step
only when the corresponding transition is specified.

\subsubsection{Compositionality and exact transport of formulas}

The recursive interpretation of assertions commutes with ordinary
substitution by Beck--Chevalley and reindexing of the designated maps.
For semantic arrow tests, one has
\[
\operatorname{Test}_{Q}(mtu)
=u^*\operatorname{Test}_{Q}(mt)
=\operatorname{Test}_{m^*Q}(tu)
\]
For a wiring map $m$, a realized-arrow tuple $t$, and a reference
substitution $u$, semantic testing commutes with composition. Relation
modalities also require their alternatives to be preserved: a
comparison map must send related pairs forward and lift each abstract
successor at an image state. Under these two conditions the specified
modal fixed points transport exactly, rather than merely one-way.

The soundness induction sends an assertion entailment
$W\mid\Phi\vdash\varphi$ to inclusion of subobjects and each
linear equation to equality of its denotations. The cases use the
universal properties established in the main text and this appendix:
adjunction and Beck--Chevalley for quantifiers, kernel pairs for
alternatives, comprehension for updates, and fixed-point
induction and coinduction for temporal operators.

\subsubsection{The formal logic of agency}
\label{sec:egt-logic}

The agent calculus supplies control, observation, and state
contexts with their intended carriers. Relation spans turn these
structures into predicate modalities: owner projections describe
controlled changes, observation kernels describe indistinguishable
states, and witnessed executions describe possible action. The
different meanings arise from their displayed relations, not from
different untyped modal operators.

\paragraph{The carriers of agent assertions}

Assertions must be interpreted over their correct carrier. Legal
profile predicates lie in $\mathsf{Pred}(L)$, whereas situated
epistemic assertions lie in $\mathsf{Pred}(W)$. For a reference
$s:W\to L$, the corresponding lifted deviation relation has a
source on the situated-state carrier and a target in the legal profile
space, with type
\[
 W\xleftarrow{d_C}W\times_{L,p_1}D_C^L
       \xrightarrow{q_C}L.
\]
Its box therefore sends $\mathsf{Pred}(L)$ into
$\mathsf{Pred}(W)$. The knowledge relation $\ker(o_i)$ instead
relates states within $W$ and induces an endomap of
$\mathsf{Pred}(W)$. Fixed-point expressions require compatible state
sorts, which the typed span calculus enforces.


\clearpage
\section{Soundness, Reconstruction, and Equational Completeness}
\label{app:metatheory}

This appendix contains the proofs that reconstruct the displayed agent
presentations from their indexed semantics, establish soundness of the
natural-deduction calculus, and construct the syntactic classifying model.
The operational, game-theoretic, and exact-reduction proofs now appear
with their respective results in the main text. The categorical setting
is the indexed semantics of \citet{Jacobs1999,Vakar2015LinearFrameworks};
model-class equational completeness is kept distinct from full
completeness in the fixed concurrent strategy category.

\subsection{Reconstruction of semantic presentations}

\begin{proof}[Proof of \cref{thm:cgt-lossless}]
\leavevmode
\begin{enumerate}[label=\textup{(\arabic*)},leftmargin=*,itemsep=4pt,topsep=3pt]
  \item \textbf{Construction of the state-context tower.} From $\vec\Sigma^G$ and $g$ construct the tower of
Theorem~\ref{thm:fib-tower}, its constant types, and the section
$\widehat g(\sigma)=(\sigma,g_\sigma)$.
  \item \textbf{Reverse implication.} Conversely retain the
displayed owner carriers and recover $g=\pi_2\widehat g$. The canonical backward map is the inverse-image operation along that
same $g$, so forgetting and rebuilding it changes nothing. The family and slice presentations correspond by the canonical maps
between a displayed domain and the disjoint union of its fibres;
these preserve predicates and all defining maps.
  \item \textbf{Lifting through the comparison.} An ownerwise arrow is recovered as the same family $r_i$. Its equation
$g=h\bar r$ is the equation identifying the pulled-back realization
section. Its cartesian context lift is the displayed pullback and its
deviation comparison is forced by the two endpoint maps.
  \item \textbf{Recovery of original concurrent arrows.} Thus arrows
and their composition are recovered as well. Coherence with wiring is
Theorem~\ref{thm:fib-wiring}; joint alternative contexts agree by
Theorem~\ref{thm:egt-counterfactual}.
  \item \textbf{Witness and legality comparison.} For each additional datum retain its carrier and its map, interpreted
as a term using constant or dependent types. Comprehension recovers
the same subset, and a history-indexed type keeps all its fibres. Certificate objects and their arrows remain the specified witnesses,
with their individual identities and causal histories.
  \item \textbf{Monoidal comparison.} Action-model product updates
are the same enabling comprehensions and relation pullbacks by
Theorem~\ref{thm:egt-action-model-semantics}; their composition
comparisons are those of Theorem~\ref{thm:egt-action-composition}. The pointed operational update and its linear residual section are
recovered by Theorem~\ref{thm:egt-operational-action-correspondence}.
  \item \textbf{Descent along the kernel pair.} For a hidden composite the secured witness carrier and visible
comparison are retained as in
Theorem~\ref{thm:egt-secured-del-comparison}, without identifying
non-descending private observations. Image factorization recovers pooled
signals, and intersection with their kernels recovers public inspection.
  \item \textbf{Observation and knowledge.} The actual observation images, their power objects, and the truthful-message
comprehensions reconstruct the announcement game of
Construction~\ref{con:egt-announcement-game}; its effectivity uses the
same dependent quantifiers. For an exact visible refinement, the retained
labels and source relations define the same monotone endomap on
$\mathcal P(X^2)$, so its greatest fixed point and effective quotient
are recovered.
  \item \textbf{Causal execution.} Corollary~\ref{cor:egt-exact-refinement-comparison}
supplies the comparison under isomorphic presentations, and
Construction~\ref{con:egt-refined-history-witnesses} retains all
dependent residual families. Numerical evaluations remain the original
numerical functions.
  \item \textbf{Fixed-point argument.} Atomic predicates have the same pullbacks by
Equation~\eqref{eq:fib-constant-test}. The quantifier formulas of
Theorem~\ref{thm:fib-comprehension} range over the same fibres, so
structural induction proves equality for first-order operations and
the relation modalities built from them. A positive fixed-point body
defines the same monotone endomap on the same predicate lattice;
its least and greatest fixed points therefore agree.
  \item \textbf{Strategic comparison.} Selection and
best response agree by Proposition~\ref{prop:fib-selection-internal};
legal selection agrees by Proposition~\ref{prop:cgt-dependent-selection}. All the remaining predicates are the specified combinations of these
operations and retained terms, proving preservation and reflection.
\end{enumerate}
\end{proof}

\begin{proof}[Proof of \cref{cor:egt-agent-reconstruction}]
\leavevmode
\begin{enumerate}[label=\textup{(\arabic*)},leftmargin=*,itemsep=4pt,topsep=3pt]
  \item \textbf{Ownerwise alternatives.} An agent interface is a record of retained owner carriers and the canonical
selection predicate; its contextual realization is a rearrangement of the
retained realization map. Products reconstruct complete-policy composition
and control regrouping, while comprehension and the outsider projection
reconstruct the legal choice fibres. The retained observation and comparison
terms give its situated and evaluated interfaces. History-indexed families
retain every continuation carrier and restriction map.
  \item \textbf{Observation and knowledge.} For communication, the retained observation image $\Omega_i$ determines the
power object $\mathcal P(\Omega_i)$. Exponentiation and the truthful-policy
comprehension reconstruct $\mathsf{AnnPol}_i$ and its evaluation. Products
therefore reconstruct $\mathcal X$ and $F$, including the private policy
records.
  \item \textbf{Lifting through the comparison.} The fibre family $K_u=F^{-1}(u)$ recovers $\mathcal X$ through the
canonical bijection $\coprod_u K_u\cong\mathcal X$; its realization term is
retained as in Construction~\ref{con:egt-announcement-policy-witnesses}. The maps in the lifting proof are the explicit extension policies, so the
bounded comparison is reconstructed as well. Its logical consequences follow
from Theorem~\ref{thm:egt-announcement-policy-comparison} and its corollaries.
  \item \textbf{Policy-level comparison.} This retains the complete policy presentation together with its logical
comparison to the message presentation.
\end{enumerate}
\end{proof}

\begin{proof}[Proof of \cref{cor:egt-enriched-agent-reconstruction}]
\leavevmode
\begin{enumerate}[label=\textup{(\arabic*)},leftmargin=*,itemsep=4pt,topsep=3pt]
  \item \textbf{Lifting through the comparison.} The retained implementation maps give the same coordinate products,
composites, and image factorizations, hence the same substituted
realizations and legal inverse images. The lifting proof uses
precisely the fibres of those image maps, so is transported along
their canonical bijections.
  \item \textbf{Adjunction and transposition.} The observation and feasibility maps reconstruct
$\forall_{o\times\operatorname{id}}F$ through its adjoint
characterization. Sections, exponentials, and inverse images
then reconstruct the informed and decentralized legal policy
carriers. A bijection of the displayed data preserves both
evaluation and the descent equations defining the image of $J_C$.
  \item \textbf{Witness and legality comparison.} An institutional composite is a pullback with its composite
realization term and an intersection of inverse-image permission
predicates. These are reconstructed from the retained endpoints,
realizers, and permission inclusions. The dependent quantifiers
on its source and target consequently reconstruct its modalities.
  \item \textbf{Representation and Yoneda decoding.} Linear realizers on representable boundaries are recovered by
Yoneda as in the preceding reconstruction results.
  \item \textbf{Transport of controller behaviour.} Finally controller behaviour is uniquely determined by its
output and next-state equations. Transport of a controller
therefore transports that behaviour by induction on observation
words. For a retained complete policy, its derivatives and their
image $S_b$ are determined by word concatenation and evaluation.
  \item \textbf{Observation and knowledge.} Bijections of the observation and action carriers induce
bijections of words and behaviours, taking derivatives to
derivatives and preserving the pointed residual controller.
\end{enumerate}
\end{proof}


\subsection{Soundness and reconstruction of the agency calculus}

\begin{proof}[Proof of \cref{agt-thm-soundness}]
\leavevmode
\begin{enumerate}[label=\textup{(\arabic*)},leftmargin=*,itemsep=4pt,topsep=3pt]
  \item \textbf{Adjunction and transposition.} Formation and the core terms are interpreted by
Proposition~\ref{agt-prop-core-sound}. Agent operations are composites of
evaluation, tensor, currying, and ordinary parameter maps, so
Theorem~\ref{agt-thm-agent-coherence} and
Proposition~\ref{agt-prop-delegation} validate their equations. Regrouping
uses the product bijections and kernel-pair comparison of
Theorem~\ref{agt-thm-regroup}.
  \item \textbf{Observation and knowledge.} Decision protocols are interpreted by
Theorem~\ref{agt-thm-protocol}, with non-anticipation and stagewise
observations proved in Theorem~\ref{agt-thm-protocol-prefix}. The
acceptance domains of protocol-generated situated contexts are the
image and dependent quantifier constructions of
Proposition~\ref{agt-prop-protocol-legality}.
  \item \textbf{Witness and legality comparison.} Finite endpoint terms
and their reductions are interpreted by Lemma~\ref{agt-lem-endpoint},
Proposition~\ref{agt-prop-session-subject-reduction}, and
Theorem~\ref{agt-thm-session}, retaining execution witnesses when the
external interface is hidden. Lemma~\ref{agt-lem-session-recorder} supplies
the explicit physical output for a retained transcript.
  \item \textbf{Causal execution.} Lemma~\ref{agt-lem-private-endpoint},
Lemma~\ref{agt-lem-session-input-residual}, and
Corollary~\ref{agt-cor-session-input} cover physical private-input
occurrences, their status-indexed residuals, and their causal
constraints.
  \item \textbf{Adjunction and transposition.} For contextual terms, Theorem~\ref{agt-thm-indexed} validates the dependent
adjoints and their substitution laws. Theorem~\ref{agt-thm-agent-dependent-product}
constructs the agent abstraction as the corresponding dependent closed
product, constructed using the linear--nonlinear adjunction. Predicate introduction
and elimination rules are the corresponding order adjunctions in $\Sub(W)$.
  \item \textbf{Adjunction and transposition: second comparison.} Perspective
uses \eqref{agt-eq-perspective-adjunction}; its coherence and epistemic
laws are Theorem~\ref{agt-thm-perspective-monoidal} and
Proposition~\ref{agt-prop-S5}. Uniform policies and comparison feedback
are subobjects defined by the proved graph adjunction and inverse-image
equations. Hence their rules are sound comprehensions and quantifications.
  \item \textbf{Witness and legality comparison.} For dynamic rules, Theorems~\ref{agt-thm-institutions} and
\ref{agt-thm-dynamics} validate composition and its modalities. Theorems~\ref{agt-thm-announcement-perspective} and
\ref{agt-thm-public-reduction}, followed by
Proposition~\ref{agt-prop-action}, validate public and general epistemic
updates.
  \item \textbf{Observation and knowledge.} Proposition~\ref{agt-prop-test-implementation} distinguishes a
logical public comprehension from the exact tests implementable using
an actual finite observation wire.
  \item \textbf{Observation and knowledge: second comparison.} Construction~\ref{agt-con-witnessed-action-data}
constructs an implemented epistemic action and its event-observation
kernels from the image of actual recipient records in retained witnesses, and Lemma~\ref{agt-lem-completed-delivery}
supplies the synchronization barrier for completed pooling. Proposition~\ref{agt-prop-witness-run} constructs the separate realized
map of witnessed execution.
  \item \textbf{Coinduction and bisimulation.} Common-knowledge unfolding and coinduction are sound by the finite
path proof of Proposition~\ref{agt-prop-common}. Induction on a complete
derivation therefore establishes each judgment and equation, with the same
substitution comparisons at every constructor.
\end{enumerate}
\end{proof}

\begin{proof}[Proof of \cref{agt-thm-reconstruction}]
\leavevmode
\begin{enumerate}[label=\textup{(\arabic*)},leftmargin=*,itemsep=4pt,topsep=3pt]
  \item \textbf{Representation and Yoneda decoding.} Theorem~\ref{agt-thm-owned-equivalence} proves the mutually inverse
encoding and decoding functors on owned presentations and their ownerwise
morphisms. Concretely, encoding copairs the arrows $yg_\sigma$ and
curries their map $\Lin\Sigma\otimes yA\to yB$, retaining the owner
display. Decoding evaluates at each $\sigma$ and applies Yoneda full
faithfulness to recover $g_\sigma:A\to B$.
  \item \textbf{Frobenius and coproduct comparison.} Coproduct and internal-hom
$\beta\eta$ give both inverse equations, while naturality proves the
reparametrization and composition comparisons.
  \item \textbf{Descent along the kernel pair.} The owner family is retained, hence the original control projections and
their kernel pairs are recovered. The legal subobject, reference map,
observations, and comparisons are retained ordinary data.
  \item \textbf{Observation and knowledge.} For
protocol-generated data, Theorem~\ref{agt-thm-protocol-prefix} reconstructs
the decision-local observation from its actual causal input tuple, while
Proposition~\ref{agt-prop-protocol-legality} reconstructs acceptance
images and universally legal profile domains from the retained record
test. Neither construction substitutes a future record for a decision's
local observation.
  \item \textbf{Adjunction and transposition.} Their derived
legal choices, perspective operations, selection predicates, and modalities
are reconstructed by the same pullbacks and adjoints. Regrouping and
delegation retain their coordinate maps, so their executions decode to
the original precompositions.
  \item \textbf{Representation and Yoneda decoding.} For witnessed institutions the map $\rho:E\to\cC(A,B)$ is equivalent,
pointwise over $E$, to the family morphism
$\rho^\sharp:s^*\underline{yA}\to t^*\underline{yB}$ of
Proposition~\ref{agt-prop-witness-run}. Encoding and decoding each
component by Yoneda are inverse and natural in maps of witnesses. Endpoints and permission remain displayed, while composition uses the
same pullback and the same composite realizer.
  \item \textbf{Policy-level comparison.} For a finite session the
\emph{displayed} session type and endpoint policies reconstruct the two
pruned trees; their secured-matching synchronization reconstructs the
transcript witness. The witness is retained, not inferred from the
hidden scalar, and all off-path coordinates remain in the policy family.
  \item \textbf{Secured operational interaction.} For implemented action updates, the retained recipient-port records
reconstruct $d_i$, their tuple image $E$, and the derived event relation
$\kerpair(b_i)$ as in
Construction~\ref{agt-con-witnessed-action-data}; the original
completed matching witnesses remain available above the image. For truthful communication, the retained $\Omega_i$ reconstructs
$\pow(\Omega_i)$, the complete-policy subtype, evaluation, and the
explicit extension map.
  \item \textbf{Observation and knowledge.} These observations cover all the displayed fields, and each comparison
commutes with them by its construction.
\end{enumerate}
\end{proof}


\subsection{The classifying model and equational completeness}

\begin{proof}[Proof of \cref{agt-lem-general-syntactic-adjoints}]
\leavevmode
\begin{enumerate}[label=\textup{(\arabic*)},leftmargin=*,itemsep=4pt,topsep=3pt]
  \item \textbf{Secured operational interaction.} Encode each finite ordinary context telescope by its iterated dependent
sum, retaining the canonical packing and unpacking substitutions. Thus
$f$ may be written as a map of the corresponding ordinary total types.
  \item \textbf{Functoriality and substitution.} Over $w:W$ form its equality fibre
\[
 \operatorname{Fib}_f(w)
 =\sum_{v:V}\operatorname{Id}_W(f(v),w).
\]
The maps
$v\mapsto(f(v),v,\operatorname{refl}_{f(v)})$ and
$(w,v,p)\mapsto v$ give inverse context substitutions between $V$
and $W.\operatorname{Fib}_f$; the reverse composite uses the diagonal equality property of
Proposition~\ref{agt-prop-equality-diagonal} to identify $w=f(v)$
and its equality witness with reflexivity.
  \item \textbf{Dependent sum and product expressions.} For a linear family $A$ over
$V$ put
\[
 (\Sigma_f^\ell A)(w)
 =\Sigma^\ell_{(v,p):\operatorname{Fib}_f(w)}A(v),\qquad
 (\Pi_f^\ell A)(w)
 =\Pi^\ell_{(v,p):\operatorname{Fib}_f(w)}A(v).
\]
These are expressions of the existing dependent constructors, not new
quantifier symbols. Given $\alpha:\Sigma_f^\ell A\to B$, define
$\beta_v:A(v)\to B(f(v))$ by applying $\alpha$ at the fibre witness
$(v,\operatorname{refl})$.
  \item \textbf{Reverse implication.} Conversely, for $\beta:A\to f^*B$,
map an element $((v,p),a)$ to the transport along $p$ of
$\beta_v(a)$. Extensional equality elimination and the dependent-sum
$\beta\eta$ equations make the two operations inverse.
  \item \textbf{Adjunction and transposition.} For the right adjunction, evaluate a map
$\alpha:B\to\Pi_f^\ell A$ at $(v,\operatorname{refl})$ to obtain
$\beta_v:B(f(v))\to A(v)$.
  \item \textbf{Reverse implication.} Conversely, $\beta:f^*B\to A$
determines the dependent abstraction
\[
 \alpha_w(b)=\Lambda(v,p).\beta_v(p^{-1}_*b).
\]
The equality-transport and dependent-product $\beta\eta$ equations again
prove the two inverse equations. Thus
\[
 \Sigma_f^\ell\dashv f^*\dashv\Pi_f^\ell.
\]
For a pullback of ordinary substitutions, its two equality-fibre
presentations are equivalent by inserting or deleting the redundant
pullback coordinate.
  \item \textbf{Base-change comparison.} Substituting that equivalence into the displayed
sum and product constructs both Beck--Chevalley maps and their inverses. Every iterated comparison has the same projections to the original
fibre witnesses, proving its cocycle by dependent $\eta$.
The explicit comparison maps and the two inverse substitutions are
verified in \cref{agt-lem-fibre-coherence}; in particular their
composite-substitution equations hold at arbitrary linear contexts.
  \item \textbf{Linear Frobenius comparison maps.} Finally, the map
$\Sigma_f^\ell(f^*B\otimes A)\to B\otimes\Sigma_f^\ell A$
eliminates the fibre witness and moves the independent $B$-factor
outside its dependent sum. Its inverse eliminates the sum in the right
tensor factor and pairs $B$ with the selected fibre component. Tensor and
dependent-sum $\beta\eta$ make the composites identities, giving
Frobenius and its substitution coherence.
\end{enumerate}
\end{proof}

\begin{proof}[Proof of \cref{agt-lem-fibre-coherence}]
\leavevmode
\begin{enumerate}[label=\textup{(\arabic*)},leftmargin=*,itemsep=4pt,topsep=3pt]
  \item \textbf{Construction and inverse of the fibre comparison.}
In the equality-fibre presentation a point of
$\operatorname{Fib}_{f'}(w')$ is a pair $(v',p')$ with
$p':f'(v')=w'$. Define
\[
 \chi_{w'}(v',p')=(a(v'),b_*(p')),
\]
where the second component is the equality
$f(a(v'))=b(w')$ obtained by transporting $p'$ through the
commuting square. Conversely, a pair
$(v,p:f(v)=b(w'))$ defines the point
$v'=(w',v,p)$ of the ordinary pullback comprehension, with
$f'(v')=w'$ by the first projection; send it to
$(v',\operatorname{refl}_{w'})$. The pullback's defining property
and extensional equality elimination identify the two composites
with the identity substitutions, including their equality witnesses.
  \item \textbf{Dependent sums, products, and naturality.} After
substitution along $b$, the fibre of $\Sigma_f^\ell A$ at $w'$ is
$\Sigma^\ell_{(v,p):\operatorname{Fib}_f(b(w'))}A_v$.
The fibre of $\Sigma_{f'}^\ell a^*A$ is the analogous dependent sum
indexed by $\operatorname{Fib}_{f'}(w')$. Reindexing the summands
along $\chi_{w'}$ and its inverse gives the two comparison maps;
the same argument using dependent products gives the right-adjoint
comparison. Their inverse equations follow from the $\Sigma^\ell$
and $\Pi^\ell$ uniqueness laws. No particular representative of an
equality class is selected.
  \item \textbf{Iterated base change.} For two composable pullback
squares, either route sends a fibre witness to its projection in the
original context $V$ together with the equality obtained by
successively applying the two base substitutions. Functoriality of
equality transport identifies this with the direct transported
equality. By the uniqueness laws of the ordinary pullback and the
linear dependent sums and products, the associated comparisons
commute on arbitrary linear terms, and not only on closed states.
  \item \textbf{Frobenius comparison and compatibility.} At $w:W$,
the Frobenius map is the copower-distribution comparison
\[
 \Sigma^\ell_{(v,p):\operatorname{Fib}_f(w)}
      (B_w\otimes A_v)
 \longrightarrow
 B_w\otimes
 \Sigma^\ell_{(v,p):\operatorname{Fib}_f(w)}A_v.
\]
On a summand indexed by $(v,p)$, it sends $(b,a)$ to
$(b,\operatorname{in}_{(v,p)}a)$; its inverse eliminates the
dependent sum in the right-hand tensor factor and retains the same
$b$. The tensor and dependent-sum $\beta\eta$ laws prove the inverse
equations. Both base-change routes use $\chi_{w'}$ to transport
exactly the same fibre witness and leave the independent $B$-term
unchanged. Thus Frobenius commutes with Beck--Chevalley, including
its composite-substitution coherence.
\end{enumerate}
\end{proof}

\begin{proof}[Proof of \cref{agt-lem-syntactic-model}]
\leavevmode
\begin{enumerate}[label=\textup{(\arabic*)},leftmargin=*,itemsep=4pt,topsep=3pt]
  \item \textbf{Monoidal comparison.} The substitution lemma makes composition well-defined and associative;
variables give identities. Ordinary substitution gives functorial
reindexing and comprehension. Tensor on objects is the type constructor;
on arrows it is tensor introduction after tensor elimination on the input.
  \item \textbf{Secured operational interaction.} Its associators, symmetry, and unitors are the corresponding resource
rearrangements. Their equations follow from tensor and unit $\beta\eta$.
  \item \textbf{Frobenius and coproduct comparison.} Currying and application give mutually inverse bijections
\[
 \cL_{\mathrm{syn}}(\Gamma)(A\otimes B,C)
 \cong\cL_{\mathrm{syn}}(\Gamma)(A,B\multimap C).
\]
Additive pairing, projections, injections, and case analysis prove the
product and coproduct universal properties, including their nullary cases.
  \item \textbf{Adjunction and transposition.} Dependent abstraction and application prove the right-adjoint bijection
for comprehension projections; dependent pairing and elimination prove
the left-adjoint bijection. The uniqueness equations supply naturality. Lemma~\ref{agt-lem-general-syntactic-adjoints} constructs their extension
to general ordinary substitutions. The equality-fibre comparison of
\cref{agt-lem-fibre-coherence} verifies Beck--Chevalley, Frobenius,
and the two composite-substitution coherences, so these constructions
are coherent on open terms and not merely fibrewise on closed values.
  \item \textbf{Term-level linear--nonlinear adjunction.} Finally, the mutually inverse maps for $\Lin\dashv\Val$ send a linear
term $z:\Lin X\vdash t:A$ to
$x:X\vdash\mathsf{thunk}(t[\mathsf{return}(x)/z]):\Val A$, and send
an ordinary term $x:X\vdash v:\Val A$ to
$z:\Lin X\vdash\mathsf{let}\ \mathsf{return}(x)=z\ \mathsf{in}\
\mathsf{force}(v):A$. Their composites are the adjunction $\beta\eta$
equations.
  \item \textbf{Monoidal comparison.} For strong monoidality, the comparison
$\Lin X\otimes\Lin Y\to\Lin(X\times Y)$ eliminates the two
linearized ordinary values and returns their pair. Its inverse eliminates
the ordinary pair $p:X\times Y$ and forms
$\mathsf{return}(\pi_1p)\otimes\mathsf{return}(\pi_2p)$. Their composites are identities by the $\Lin$, tensor, and ordinary
product $\beta\eta$ equations.
  \item \textbf{Monoidal comparison: completion.} The same computation gives
$\Lin1\cong\I$ and the associativity, unit, and symmetry coherence
of this strong monoidal structure. All constructions are stable under ordinary substitution. The declared constants interpret themselves.
\end{enumerate}
\end{proof}

\begin{proof}[Proof of \cref{agt-thm-completeness}]
\leavevmode
\begin{enumerate}[label=\textup{(\arabic*)},leftmargin=*,itemsep=4pt,topsep=3pt]
  \item \textbf{Construction of the comparison.} Define the interpretation recursively on contexts, type constructors,
terms, and substitutions. Each constructor must be interpreted by the
corresponding chosen operation; hence the recursion is unique. The
universal-property equations make it respect conversion and derivable
equality, as in Proposition~\ref{agt-prop-core-sound}. This proves initiality
and the forward implication.
  \item \textbf{Reverse implication.} For the converse, use the syntactic model of
Lemma~\ref{agt-lem-syntactic-model}. Interpretation there returns precisely
the term's own derivable-equality class. Thus, if two interpretations
coincide in every model, they coincide in this model and are related by
the derivable-equality judgment.
  \item \textbf{Monoidal comparison.} This applies to open terms because the
hom-set of the syntactic fibre is \emph{defined} by their equality
classes, not by evaluation at global points. Open linear contexts are
represented by their tensor
and unpacked by tensor elimination; its $\beta\eta$ equations identify this
encoding with the original judgment.
  \item \textbf{Conservativity of displayed agent types.} The displayed agent judgment is conservative by
Definition~\ref{agt-def-displayed-agents}: erasure recovers the underlying
term and the tag is recovered unchanged. Agent operations are constructed
terms with specified displayed maps, so the same syntactic-model argument
applies to their typed equations. In split strict presentations the interpretation is unique on the
nose.
  \item \textbf{Secured operational interaction.} For models and morphisms preserving the chosen structures by
specified coherent comparisons, the same initiality argument yields
the corresponding bi-initial statement.
\end{enumerate}
\end{proof}


% Bibliography maintained directly inside this .tex document.
% Citation keys match those in the mathematical text; no BibTeX or Biber step is needed.
\begin{thebibliography}{99}
\setlength{\itemsep}{5pt}
\bibitem[Abramsky and Zvesper(2010)]{AbramskyZvesper2010}
Samson Abramsky and Jonathan Zvesper (2010).
\newblock \emph{From Lawvere to Brandenburger--Keisler: Interactive Forms of Diagonalization and Self-Reference}.
\newblock \url{https://arxiv.org/abs/1006.0992}.

\bibitem[Abramsky and Winschel(2017)]{AbramskyWinschel2013}
Samson Abramsky and Viktor Winschel (2017).
\newblock \emph{Coalgebraic Analysis of Subgame-perfect Equilibria in Infinite Games without Discounting}.
\newblock \emph{Mathematical Structures in Computer Science}, \textbf{27}(5), 751--761. First published online in 2015.
\newblock DOI: \url{https://doi.org/10.1017/S0960129515000365}.

\bibitem[Ågotnes and van Ditmarsch(2011)]{AgotnesVanDitmarsch2011AnnouncementGames}
Thomas Ågotnes and Hans van Ditmarsch (2011).
\newblock \emph{What Will They Say? Public Announcement Games}.
\newblock \emph{Synthese}, \textbf{179}, 57--85.
\newblock DOI: \url{https://doi.org/10.1007/s11229-010-9838-8}.

\bibitem[Ågotnes, Alechina, and Galimullin(2022)]{AgotnesEtAl2022GroupDistributed}
Thomas Ågotnes, Natasha Alechina, and Rustam Galimullin (2022).
\newblock \emph{Logics with Group Announcements and Distributed Knowledge: Completeness and Expressive Power}.
\newblock \emph{Journal of Logic, Language and Information}, \textbf{31}, 141--166.
\newblock DOI: \url{https://doi.org/10.1007/s10849-022-09355-0}.

\bibitem[Ågotnes et~al.(2010)]{AgotnesEtAl2010GroupAnnouncements}
Thomas Ågotnes, Philippe Balbiani, Hans van Ditmarsch, and Pablo Seban (2010).
\newblock \emph{Group Announcement Logic}.
\newblock \emph{Journal of Applied Logic}, \textbf{8}(1), 62--81.
\newblock DOI: \url{https://doi.org/10.1016/j.jal.2008.12.002}.

\bibitem[Ågotnes and Wáng(2016)]{AgotnesWang2016ResolvingDistributedKnowledge}
Thomas Ågotnes and Yì N. Wáng (2016).
\newblock \emph{Resolving Distributed Knowledge}.
\newblock \url{https://arxiv.org/abs/1606.07515}.

\bibitem[Alur, Henzinger, and Kupferman(2002)]{AlurEtAl2002}
Rajeev Alur, Thomas A. Henzinger, and Orna Kupferman (2002).
\newblock \emph{Alternating-time Temporal Logic}.
\newblock \emph{Journal of the ACM}, \textbf{49}(5), 672--713.
\newblock DOI: \url{https://doi.org/10.1145/585265.585270}.

\bibitem[Apt and Zvesper(2010)]{ZvesperApt2010}
Krzysztof R. Apt and Jonathan A. Zvesper (2010).
\newblock \emph{The Role of Monotonicity in the Epistemic Analysis of Strategic Games}.
\newblock \emph{Games}, \textbf{1}(4), 381--394.
\newblock DOI: \url{https://doi.org/10.3390/g1040381}.

\bibitem[Atkey et~al.(2021)]{AtkeyEtAl2021}
Robert Atkey, Bruno Gavranović, Neil Ghani, Clemens Kupke, Jérémy Ledent, and Fredrik Nordvall Forsberg (2021).
\newblock \emph{Compositional Game Theory, Compositionally}.
\newblock In \emph{Proceedings of the Third International Conference on Applied Category Theory (ACT 2020)}, Electronic Proceedings in Theoretical Computer Science, vol. 333, pp. 198--214.
\newblock DOI: \url{https://doi.org/10.4204/EPTCS.333.14}.

\bibitem[Aumann and Brandenburger(1995)]{AumannBrandenburger1995}
Robert J. Aumann and Adam Brandenburger (1995).
\newblock \emph{Epistemic Conditions for Nash Equilibrium}.
\newblock \emph{Econometrica}, \textbf{63}(5), 1161--1180.

\bibitem[Baltag, Coecke, and Sadrzadeh(2007)]{agt-BaltagCoeckeSadrzadeh}
Alexandru Baltag, Bob Coecke, and Mehrnoosh Sadrzadeh (2007).
\newblock \emph{Epistemic Actions as Resources}.
\newblock \emph{Journal of Logic and Computation}, \textbf{17}(3), 555--585.
\newblock DOI: \url{https://doi.org/10.1093/logcom/exm015}.

\bibitem[Baltag, Moss, and Solecki(1998)]{BaltagMossSolecki1998PublicAnnouncements}
Alexandru Baltag, Lawrence S. Moss, and Sławomir Solecki (1998).
\newblock \emph{The Logic of Public Announcements, Common Knowledge, and Private Suspicions}.
\newblock In \emph{Proceedings of the 7th Conference on Theoretical Aspects of Rationality and Knowledge}.

\bibitem[Barasz et~al.(2014)]{BaraszEtAl2014}
Mihaly Barasz, Paul Christiano, Benja Fallenstein, Marcello Herreshoff, Patrick LaVictoire, and Eliezer Yudkowsky (2014).
\newblock \emph{Robust Cooperation in the Prisoner's Dilemma: Program Equilibrium via Provability Logic}.
\newblock \url{https://arxiv.org/abs/1401.5577}.

\bibitem[Barber et~al.(1998)]{agt-BarberEtAl}
Andrew Barber, Philippa Gardner, Masahito Hasegawa, and Gordon Plotkin (1998).
\newblock \emph{From Action Calculi to Linear Logic}.
\newblock In \emph{Computer Science Logic (CSL 1997)}, Lecture Notes in Computer Science, vol. 1414, pp. 78--97.

\bibitem[Benton(1995)]{Benton1994MixedLNL}
Nick Benton (1995).
\newblock \emph{A Mixed Linear and Non-Linear Logic: Proofs, Terms and Models}.
\newblock In \emph{Computer Science Logic (CSL 1994)}, Lecture Notes in Computer Science, vol. 933, pp. 121--135, Springer.

\bibitem[Bonzon et~al.(2006)]{BonzonEtAl2006}
Elise Bonzon, Marie-Christine Lagasquie-Schiex, Jérôme Lang, and Bruno Zanuttini (2006).
\newblock \emph{Boolean Games Revisited}.
\newblock In \emph{Proceedings of the 17th European Conference on Artificial Intelligence}, pp. 265--269.

\bibitem[Caires and Pfenning(2010)]{agt-CairesEtAl}
Luís Caires and Frank Pfenning (2010).
\newblock \emph{Session Types as Intuitionistic Linear Propositions}.
\newblock In \emph{21st International Conference on Concurrency Theory (CONCUR 2010)}, Lecture Notes in Computer Science, vol. 6269, pp. 222--236.
\newblock DOI: \url{https://doi.org/10.1007/978-3-642-15375-4_16}.

\bibitem[Capucci(2023)]{Capucci2023}
Matteo Capucci (2023).
\newblock \emph{Diegetic Representation of Feedback in Open Games}.
\newblock In \emph{Proceedings of the Fifth International Conference on Applied Category Theory (ACT 2022)}, Electronic Proceedings in Theoretical Computer Science, vol. 380, pp. 145--158.
\newblock DOI: \url{https://doi.org/10.4204/EPTCS.380.9}.

\bibitem[Capucci et~al.(2022a)]{CapucciEtAl2022Cybernetics}
Matteo Capucci, Bruno Gavranović, Jules Hedges, and Eigil Fjeldgren Rischel (2022a).
\newblock \emph{Towards Foundations of Categorical Cybernetics}.
\newblock In \emph{Proceedings of the Fourth International Conference on Applied Category Theory (ACT 2021)}, Electronic Proceedings in Theoretical Computer Science, vol. 372, pp. 235--248.
\newblock DOI: \url{https://doi.org/10.4204/EPTCS.372.17}.

\bibitem[Capucci et~al.(2022b)]{CapucciEtAl2022Agency}
Matteo Capucci, Neil Ghani, Jérémy Ledent, and Fredrik Nordvall Forsberg (2022b).
\newblock \emph{Translating Extensive Form Games to Open Games with Agency}.
\newblock In \emph{Proceedings of the Fourth International Conference on Applied Category Theory (ACT 2021)}, Electronic Proceedings in Theoretical Computer Science, vol. 372, pp. 221--234.
\newblock DOI: \url{https://doi.org/10.4204/EPTCS.372.16}.

\bibitem[Carbone et~al.(2015)]{CarboneEtAl2015}
Marco Carbone, Fabrizio Montesi, Carsten Schürmann, and Nobuko Yoshida (2015).
\newblock \emph{Multiparty Session Types as Coherence Proofs}.
\newblock In \emph{26th International Conference on Concurrency Theory (CONCUR 2015)}, LIPIcs, vol. 42, pp. 412--426.
\newblock DOI: \url{https://doi.org/10.4230/LIPIcs.CONCUR.2015.412}.

\bibitem[Carbone et~al.(2017)]{CarboneEtAl2016}
Marco Carbone, Fabrizio Montesi, Carsten Schürmann, and Nobuko Yoshida (2017).
\newblock \emph{Multiparty Session Types as Coherence Proofs}.
\newblock \emph{Acta Informatica}, \textbf{54}(3), 243--269.
\newblock DOI: \url{https://doi.org/10.1007/s00236-016-0285-y}.

\bibitem[Castellan and Yoshida(2019)]{agt-CastellanYoshida}
Simon Castellan and Nobuko Yoshida (2019).
\newblock \emph{Causality in Linear Logic: Full Completeness and Injectivity}.
\newblock In \emph{Foundations of Software Science and Computation Structures (FoSSaCS 2019)}, Lecture Notes in Computer Science, pp. 150--168.
\newblock DOI: \url{https://doi.org/10.1007/978-3-030-17127-8_9}.

\bibitem[Castellan, Clairambault, and Winskel(2019)]{CastellanClairambaultWinskel2019}
Simon Castellan, Pierre Clairambault, and Glynn Winskel (2019).
\newblock \emph{Thin Games with Symmetry and Concurrent Hyland--Ong Games}.
\newblock \emph{Logical Methods in Computer Science}, \textbf{15}(1), article 18.
\newblock DOI: \url{https://doi.org/10.23638/LMCS-15(1:18)2019}.

\bibitem[Castellan et~al.(2017)]{CastellanEtAl2017}
Simon Castellan, Pierre Clairambault, Silvain Rideau, and Glynn Winskel (2017).
\newblock \emph{Games and Strategies as Event Structures}.
\newblock \emph{Logical Methods in Computer Science}, \textbf{13}(3), article 35.
\newblock DOI: \url{https://doi.org/10.23638/LMCS-13(3:35)2017}.

\bibitem[Chen, Shi, and Wang(2022)]{ChenShiWang2022}
Qian Chen, Chenwei Shi, and Yiyan Wang (2022).
\newblock \emph{Reasoning about Dependence, Preference and Coalitional Power}.
\newblock \url{https://arxiv.org/abs/2209.08213}.

\bibitem[Cohen and Levesque(1990)]{CohenLevesque1990Intention}
Philip R. Cohen and Hector J. Levesque (1990).
\newblock \emph{Intention Is Choice with Commitment}.
\newblock \emph{Artificial Intelligence}, \textbf{42}(2-3), 213--261.
\newblock DOI: \url{https://doi.org/10.1016/0004-3702(90)90055-5}.

\bibitem[Day(1970)]{Day1970ClosedCategories}
Brian Day (1970).
\newblock \emph{On Closed Categories of Functors}.
\newblock In \emph{Reports of the Midwest Category Seminar IV}, Lecture Notes in Mathematics, vol. 137, pp. 1--38, Springer.
\newblock DOI: \url{https://doi.org/10.1007/BFb0060438}.

\bibitem[Escardó and Oliva(2010)]{EscardoOliva2010}
Martín Escardó and Paulo Oliva (2010).
\newblock \emph{Selection Functions, Bar Recursion and Backward Induction}.
\newblock \emph{Mathematical Structures in Computer Science}, \textbf{20}(2), 127--168.
\newblock DOI: \url{https://doi.org/10.1017/S0960129509990351}.

\bibitem[Finkbeiner and Olderog(2017)]{FinkbeinerOlderog2017PetriGames}
Bernd Finkbeiner and Ernst-Rüdiger Olderog (2017).
\newblock \emph{Petri Games: Synthesis of Distributed Systems with Causal Memory}.
\newblock \emph{Information and Computation}, \textbf{253}, 181--203.
\newblock DOI: \url{https://doi.org/10.1016/j.ic.2016.07.006}.

\bibitem[Fisman, Kupferman, and Lustig(2009)]{FismanEtAl2009}
Dana Fisman, Orna Kupferman, and Yoad Lustig (2009).
\newblock \emph{Rational Synthesis}.
\newblock \url{https://arxiv.org/abs/0907.3019}.

\bibitem[Frittella et~al.(2016)]{agt-FrittellaEtAl}
Sabine Frittella, Giuseppe Greco, Alexander Kurz, Alessandra Palmigiano, and Vlasta Sikimić (2016).
\newblock \emph{Multi-type Display Calculus for Dynamic Epistemic Logic}.
\newblock \emph{Journal of Logic and Computation}, \textbf{26}(6), 2017--2065.
\newblock DOI: \url{https://doi.org/10.1093/logcom/exu068}.

\bibitem[Galimullin and Alechina(2018)]{GalimullinAlechina2018CoalitionGroupAnnouncement}
Rustam Galimullin and Natasha Alechina (2018).
\newblock \emph{Coalition and Group Announcement Logic}.
\newblock \url{https://arxiv.org/abs/1810.02769}.

\bibitem[Garrabrant, Herrmann, and Lopez-Wild(2021)]{GarrabrantEtAl2021CartesianFrames}
Scott Garrabrant, Daniel A. Herrmann, and Josiah Lopez-Wild (2021).
\newblock \emph{Cartesian Frames}.
\newblock \url{https://arxiv.org/abs/2109.10996}.

\bibitem[Ghani et~al.(2018)]{GhaniEtAl2018}
Neil Ghani, Jules Hedges, Viktor Winschel, and Philipp Zahn (2018).
\newblock \emph{Compositional Game Theory}.
\newblock In \emph{Proceedings of the 33rd Annual ACM/IEEE Symposium on Logic in Computer Science (LICS)}, pp. 472--481.
\newblock DOI: \url{https://doi.org/10.1145/3209108.3209165}.

\bibitem[Girard(1987)]{Girard1987}
Jean-Yves Girard (1987).
\newblock \emph{Linear Logic}.
\newblock \emph{Theoretical Computer Science}, \textbf{50}, 1--102.
\newblock DOI: \url{https://doi.org/10.1016/0304-3975(87)90045-4}.

\bibitem[Halpern and Fagin(1989)]{HalpernFagin1989KnowledgeAction}
Joseph Y. Halpern and Ronald Fagin (1989).
\newblock \emph{Modelling Knowledge and Action in Distributed Systems}.
\newblock \emph{Distributed Computing}, \textbf{3}(4), 159--177.

\bibitem[Harper(2016)]{agt-Harper}
Robert Harper (2016).
\newblock \emph{Practical Foundations for Programming Languages}.
\newblock 2nd edition. Cambridge University Press.

\bibitem[Harsanyi(1967)]{Harsanyi1967}
John C. Harsanyi (1967).
\newblock \emph{Games with Incomplete Information Played by Bayesian Players, Part I}.
\newblock \emph{Management Science}, \textbf{14}(3), 159--182.

\bibitem[Hedges(2020)]{Hedges2020}
Jules Hedges (2020).
\newblock \emph{The Game Semantics of Game Theory}.
\newblock \url{https://arxiv.org/abs/1904.11287}.

\bibitem[Hennessy and Milner(1980)]{HennessyMilner1980}
Matthew Hennessy and Robin Milner (1980).
\newblock \emph{On Observing Nondeterminism and Concurrency}.
\newblock In \emph{Automata, Languages and Programming (ICALP 1980)}, Lecture Notes in Computer Science, vol. 85, pp. 299--309, Springer.
\newblock DOI: \url{https://doi.org/10.1007/3-540-10003-2_79}.

\bibitem[Hennessy and Milner(1985)]{HennessyMilner1985}
Matthew Hennessy and Robin Milner (1985).
\newblock \emph{Algebraic Laws for Nondeterminism and Concurrency}.
\newblock \emph{Journal of the ACM}, \textbf{32}(1), 137--161.
\newblock DOI: \url{https://doi.org/10.1145/2455.2460}.

\bibitem[Horty and Belnap(1995)]{HortyBelnap1995DeliberativeStit}
John Horty and Nuel Belnap (1995).
\newblock \emph{The Deliberative STIT: A Study of Action, Omission, Ability, and Obligation}.
\newblock \emph{Journal of Philosophical Logic}, \textbf{24}, 583--644.

\bibitem[Jacobs(1999)]{Jacobs1999}
Bart Jacobs (1999).
\newblock \emph{Categorical Logic and Type Theory}.
\newblock Studies in Logic and the Foundations of Mathematics, vol. 141. Elsevier.

\bibitem[Jamroga and van der Hoek(2004)]{JamrogaVanDerHoek2004KnowHow}
Wojciech Jamroga and Wiebe van der Hoek (2004).
\newblock \emph{Agents that Know How to Play}.
\newblock \emph{Fundamenta Informaticae}, \textbf{63}(2-3), 185--219.

\bibitem[Jones and Sergot(1996)]{JonesSergot1996InstitutionalPower}
Andrew J. I. Jones and Marek Sergot (1996).
\newblock \emph{A Formal Characterisation of Institutionalised Power}.
\newblock \emph{Journal of the IGPL}, \textbf{4}(3), 427--443.

\bibitem[Kishida(2017)]{Kishida2017CategoriesDEL}
Kohei Kishida (2017).
\newblock \emph{Categories for Dynamic Epistemic Logic}.
\newblock In \emph{Proceedings of the Sixteenth Conference on Theoretical Aspects of Rationality and Knowledge (TARK 2017)}, Electronic Proceedings in Theoretical Computer Science, vol. 251, pp. 353--372.
\newblock DOI: \url{https://doi.org/10.4204/EPTCS.251.26}.

\bibitem[Landsman(2001)]{Landsman2001}
N. P. Landsman (2001).
\newblock \emph{Operator Algebras and Poisson Manifolds Associated to Groupoids}.
\newblock \emph{Communications in Mathematical Physics}, \textbf{222}, 97--116.
\newblock DOI: \url{https://doi.org/10.1007/s002200100496}.

\bibitem[List and Pettit(2011)]{ListPettit2011GroupAgency}
Christian List and Philip Pettit (2011).
\newblock \emph{Group Agency: The Possibility, Design, and Status of Corporate Agents}.
\newblock Oxford University Press.

\bibitem[Mogavero et~al.(2012)]{MogaveroEtAl2012}
Fabio Mogavero, Aniello Murano, Giuseppe Perelli, and Moshe Y. Vardi (2012).
\newblock \emph{What Makes ATL* Decidable? A Decidable Fragment of Strategy Logic}.
\newblock In \emph{23rd International Conference on Concurrency Theory (CONCUR 2012)}, Lecture Notes in Computer Science, vol. 7454, pp. 193--208.
\newblock DOI: \url{https://doi.org/10.1007/978-3-642-32940-1_15}.

\bibitem[Pauly(2002)]{Pauly2002}
Marc Pauly (2002).
\newblock \emph{A Modal Logic for Coalitional Power in Games}.
\newblock \emph{Journal of Logic and Computation}, \textbf{12}(1), 149--166.
\newblock DOI: \url{https://doi.org/10.1093/logcom/12.1.149}.

\bibitem[Plaza(2007)]{Plaza2007PublicCommunications}
Jan Plaza (2007).
\newblock \emph{Logics of Public Communications}.
\newblock \emph{Synthese}, \textbf{158}(2), 165--179.
\newblock DOI: \url{https://doi.org/10.1007/s11229-007-9168-7}.

\bibitem[Porello and Troquard(2014)]{agt-PorelloTroquard}
Daniele Porello and Nicolas Troquard (2014).
\newblock \emph{A Resource-Sensitive Logic of Agency}.
\newblock In \emph{21st European Conference on Artificial Intelligence (ECAI 2014)}, pp. 723--728.

\bibitem[Radner(1962)]{Radner1962TeamDecisionProblems}
Roy Radner (1962).
\newblock \emph{Team Decision Problems}.
\newblock \emph{The Annals of Mathematical Statistics}, \textbf{33}(3), 857--881.

\bibitem[Rutten(2000)]{Rutten2000UniversalCoalgebra}
Jan J. M. M. Rutten (2000).
\newblock \emph{Universal Coalgebra: A Theory of Systems}.
\newblock \emph{Theoretical Computer Science}, \textbf{249}(1), 3--80.
\newblock DOI: \url{https://doi.org/10.1016/S0304-3975(00)00056-6}.

\bibitem[Selinger(2011)]{Selinger2011}
Peter Selinger (2011).
\newblock \emph{A Survey of Graphical Languages for Monoidal Categories}.
\newblock In \emph{New Structures for Physics}, Lecture Notes in Physics, vol. 813, pp. 289--355.
\newblock DOI: \url{https://doi.org/10.1007/978-3-642-12821-9_4}.

\bibitem[Shapley(1953)]{Shapley1953}
Lloyd S. Shapley (1953).
\newblock \emph{Stochastic Games}.
\newblock \emph{Proceedings of the National Academy of Sciences}, \textbf{39}(10), 1095--1100.
\newblock DOI: \url{https://doi.org/10.1073/pnas.39.10.1095}.

\bibitem[Tarski(1955)]{Tarski1955}
Alfred Tarski (1955).
\newblock \emph{A Lattice-Theoretical Fixpoint Theorem and Its Applications}.
\newblock \emph{Pacific Journal of Mathematics}, \textbf{5}, 285--309.

\bibitem[Vákár(2015)]{Vakar2015LinearFrameworks}
Matthijs Vákár (2015).
\newblock \emph{A Categorical Semantics for Linear Logical Frameworks}.
\newblock \url{https://arxiv.org/abs/1501.05016}.

\bibitem[van Benthem and Gheerbrant(2010)]{VanBenthemGheerbrant2010}
Johan van Benthem and Amélie Gheerbrant (2010).
\newblock \emph{Game Solution, Epistemic Dynamics and Fixed-Point Logics}.
\newblock \emph{Fundamenta Informaticae}, \textbf{100}(1-4), 19--41.

\bibitem[van Benthem, Pacuit, and Roy(2011)]{VanBenthemPacuitRoy2011}
Johan van Benthem, Eric Pacuit, and Olivier Roy (2011).
\newblock \emph{Toward a Theory of Play: A Logical Perspective on Games and Interaction}.
\newblock \emph{Games}, \textbf{2}(1), 52--86.
\newblock DOI: \url{https://doi.org/10.3390/g2010052}.

\bibitem[van Ditmarsch et~al.(2020)]{VanDitmarschEtAl2020ArrowUpdateSynthesis}
Hans van Ditmarsch, Wiebe van der Hoek, Barteld Kooi, and Louwe B. Kuijer (2020).
\newblock \emph{Arrow Update Synthesis}.
\newblock \emph{Information and Computation}, \textbf{275}, article 104544.
\newblock DOI: \url{https://doi.org/10.1016/j.ic.2020.104544}.

\bibitem[van Eijck, Gattinger, and Wang(2017)]{VanEijckGattingerWang2017KnowingValues}
Jan van Eijck, Malvin Gattinger, and Yanjing Wang (2017).
\newblock \emph{Knowing Values and Public Inspection}.
\newblock In \emph{Logic and Its Applications}, Lecture Notes in Computer Science, vol. 10119, pp. 77--90.
\newblock DOI: \url{https://doi.org/10.1007/978-3-662-54069-5_7}.

\bibitem[Winskel(2023)]{Winskel2023}
Glynn Winskel (2023).
\newblock \emph{Making Concurrency Functional}.
\newblock In \emph{38th Annual ACM/IEEE Symposium on Logic in Computer Science (LICS)}.
\newblock DOI: \url{https://doi.org/10.1109/LICS56636.2023.10175727}.

\end{thebibliography}
\end{document}
